\documentclass[11pt,openany]{book}
\setcounter{secnumdepth}{3}
\usepackage[margin=1.1in]{geometry}
% Korean edition: build with XeLaTeX (see the Makefile)
\usepackage[hangul]{kotex}
\nocompresspunctuations % closing quotes at natural width: no gap between a closing quote and a following particle
\setmainhangulfont{UnBatang}[Extension=.ttf,BoldFont=UnBatangBold,ItalicFont=UnDotum,BoldItalicFont=UnDotumBold,SlantedFont=UnDotum,BoldSlantedFont=UnDotumBold] % \emph in Dotum: the Un fonts have no italic
\setsanshangulfont{UnDotum}[Extension=.ttf,BoldFont=UnDotumBold,ItalicFont=UnDotum,BoldItalicFont=UnDotumBold,SlantedFont=UnDotum,BoldSlantedFont=UnDotumBold]
\usepackage{amsmath,amssymb,amsthm}
\usepackage{graphicx}
\usepackage{tikz}
\usetikzlibrary{arrows.meta,decorations.markings,decorations.pathmorphing,decorations.pathreplacing,calc}
\input{figures/icons}
\usepackage[unicode,colorlinks=true,linkcolor=blue!60!black,citecolor=blue!60!black]{hyperref}
\usepackage{microtype}
\setlength{\emergencystretch}{3em} % long names (licences, references) in narrow lines
\usepackage{empheq}
\usepackage{array}
\usepackage{booktabs}
\usepackage{multirow}
\usepackage{longtable}
\usepackage{circuitikz}
\definecolor{keyyellow}{rgb}{1.0,0.97,0.78}
% key equations: \begin{empheq}[box=\keybox]{equation} ... \end{empheq} -- light-yellow box, number stays outside
\newcommand{\keybox}[1]{{\setlength{\fboxsep}{7pt}\colorbox{keyyellow}{#1}}}
\usepackage[most]{tcolorbox}
% key statements (propositions etc.): the same light-yellow background, breakable across pages
\newtcolorbox{keyblock}{enhanced,breakable,colback=keyyellow,colframe=keyyellow,boxrule=0pt,arc=0pt,left=6pt,right=6pt,top=2pt,bottom=2pt,boxsep=0pt}
\renewcommand{\chaptermark}[1]{\markboth{\small 제\thechapter 장\quad #1}{}}
\renewcommand{\sectionmark}[1]{\markright{\small\thesection.\ #1}}
\renewcommand{\topfraction}{0.95}
\renewcommand{\textfraction}{0.05}
\renewcommand{\floatpagefraction}{0.75}

% part pages: title at the top third, and text may follow on the same page (used for the part introductions)
\makeatletter
\renewcommand\part{\if@openright\cleardoublepage\else\clearpage\fi\thispagestyle{plain}\if@twocolumn\onecolumn\@tempswatrue\else\@tempswafalse\fi\null\vspace{0.18\textheight}\secdef\@part\@spart}
\renewcommand\@endpart{\par\vspace{3em}\if@tempswa\twocolumn\fi}
\makeatother
\newtheorem{theorem}{정리}[chapter]
\newtheorem{proposition}{명제}[chapter]
\newtheorem{lemma}{보조정리}[chapter]
\theoremstyle{remark}
\newtheorem{remark}{주의}[chapter]
% exercises: \begin{exercise}[\lvlB] ... \end{exercise}, numbered "Exercise 4.3"; difficulty marks
\theoremstyle{definition}
\newtheorem{exercise}{연습문제}[chapter]
\newcommand{\lvlA}{$\star$}
\newcommand{\lvlB}{$\star\star$}
\newcommand{\lvlC}{$\star\star\star$}
% answer to an exercise in the Answers appendix: \answer{ex:NN:k}{text}
\newcommand{\answer}[2]{\par\noindent\textbf{\ref{#1}.}~#2\par\smallskip}
\newcommand{\dd}{\mathrm{d}}
% neuron type code: first four letters of the primary mediator
\newcommand{\nt}[1]{\textsf{\textbf{#1}}}
\newcommand{\mV}{\,\text{mV}}
\newcommand{\ms}{\,\text{ms}}
\newcommand{\mM}{\,\text{mM}}
\newcommand{\um}{\,\text{\textmu m}}
\newcommand{\ion}[2]{\mathrm{#1}^{#2}}
\newcommand{\Na}{\ion{Na}{+}}
\newcommand{\K}{\ion{K}{+}}
\newcommand{\Cl}{\ion{Cl}{-}}
\newcommand{\Ca}{\ion{Ca}{2+}}
\newcommand{\conc}[1]{[\mathrm{#1}]}

\title{\Huge 20와트 컴퓨터\\[16pt] \Large 뇌는 어떻게 계산하는가: 엔지니어를 위한 정량적 접근}
\hypersetup{pdftitle={20와트 컴퓨터. 뇌는 어떻게 계산하는가: 엔지니어를 위한 정량적 접근},pdfauthor={Konstantin Kladko}}
\author{\Large 콘스탄틴 클라드코 (Konstantin Kladko)}
% cover: a salad of the book's figures around the title card (figures/bookcover.tex)
\newcommand{\covertitle}{20와트 컴퓨터}
\newcommand{\coversubtitle}{뇌는 어떻게 계산하는가: 엔지니어를 위한 정량적 접근}
\newcommand{\coverauthor}{콘스탄틴 클라드코 (Konstantin Kladko)}
\date{}

\begin{document}
\frontmatter
\input{figures/bookcover}
\thispagestyle{empty}
\vspace*{\fill}
\noindent{\small \copyright~2026 Konstantin Kladko(콘스탄틴 클라드코).\\[4pt]
이 책의 본문, 그림, 연습문제와 답은 크리에이티브 커먼즈 저작자표시-비영리-동일조건변경허락 4.0 국제 라이선스(CC BY-NC-SA 4.0)에 따라 배포된다: \url{https://creativecommons.org/licenses/by-nc-sa/4.0/deed.ko}. 저작자를 표시하고 2차적 저작물을 같은 라이선스로 배포하는 한, 비영리 목적으로 자유롭게 복제하고 수업에 사용하며 번역하고 개작할 수 있다.\\[4pt]
모델과 빌드 프로그램은 MIT 라이선스에 따라 배포된다.\\[4pt]
소스, 모델, 최신판: \url{https://github.com/kladkogex/20-watt-computer}.}
\clearpage

\tableofcontents
\chapter*{이 책을 읽는 방법}
\addcontentsline{toc}{chapter}{이 책을 읽는 방법}

이 책은 처음부터 순서대로 읽지 않아도 된다. 1--4장은 공통 기초로, 뉴런의 유형, \nt{GLUT}와 \nt{GABA}가 무엇을 계산하는지, 그리고 이들이 어떤 언어로 말하는지를 다룬다. 그다음부터 책은 여러 갈래로 나뉜다. 그림~\ref{fig:roadmap}의 지도는 어느 장이 어느 장에 기대는지 보여 준다. 장 $A$에서 장 $B$로 가는 화살표는 $B$가 $A$의 개념과 식을 사용한다는 뜻이다. 지도에는 주요 의존 관계만 표시했다. 앞 장을 가리키는 사소한 참조는 읽는 데 지장을 주지 않는다.

\begin{figure}[h]
\centering
\begin{tikzpicture}[>=Stealth,scale=0.95,
  ch/.style={draw,thick,rounded corners=3pt,align=center,font=\scriptsize,text width=1.85cm,minimum height=0.62cm,inner sep=2pt},
  base/.style={ch,fill=blue!8}, bus/.style={ch,fill=orange!14}, sum/.style={ch,fill=gray!18},
  io/.style={ch,fill=green!10}, sort/.style={ch,fill=cyan!8}, lab/.style={ch,fill=violet!10},
  dep/.style={->,thick,gray!60!black}]
% foundation
\node[base] (c1) at (0,0) {1 뉴런의 유형};
\node[base] (c2) at (2.3,0) {2 \nt{GLUT}};
\node[base] (c3) at (4.6,0) {3 \nt{GLUT}와 \nt{GABA}};
\node[base] (c4) at (6.9,0) {4 모스 부호};
% broadcast channels and the synthesis
\node[bus] (c5) at (4.6,-1.5) {5 \nt{SERO}};
\node[bus] (c6) at (7.3,-1.5) {6 \nt{DOPA}};
\node[bus] (c7) at (9.2,-0.9) {7 도파민 이론};
\node[bus] (c8) at (9.2,-2.1) {8 \nt{NORA}};
\node[bus] (c9) at (11.5,-2.1) {9 \nt{ACET}};
\node[sum] (c10) at (13.8,-0.9) {10 기계 전체};
% inputs and outputs
\node[io] (c12) at (2.3,-4.1) {12 인식};
\node[io] (c11) at (4.6,-4.1) {11 입력};
\node[io] (c13) at (6.9,-4.1) {13 근육};
\node[io] (c14) at (9.2,-3.05) {14 통증};
\node[io] (c15) at (9.2,-4.1) {15 운동\\학습};
\node[io] (c16) at (9.2,-5.15) {16 몸의\\화학};
% lab, sorting, sleep
\node[lab] (c17) at (11.5,-4.1) {17 디버깅};
\node[lab] (c21) at (13.8,-4.1) {21 살아 있는 기계};
\node[lab] (c22) at (13.8,-5.15) {22 DishBrain};
\node[sort] (c18) at (11.5,-5.15) {18 계측기\\뉴런};
\node[sort] (c19) at (13.8,-6.2) {19 가지돌기\\개수};
\node[bus] (c20) at (11.5,-6.2) {20 수면};
% dependencies
\draw[dep] (c1) -- (c2); \draw[dep] (c2) -- (c3); \draw[dep] (c3) -- (c4);
\draw[dep] (c1) -- (c5); \draw[dep] (c4) -- (c6); \draw[dep] (c5) -- (c6);
\draw[dep] (c6) -- (c7); \draw[dep] (c6) -- (c8); \draw[dep] (c8) -- (c9);
\draw[dep] (c7) -- (c10); \draw[dep] (c9) -- (c10);
\draw[dep] (c4) -- (c11); \draw[dep] (c11) -- (c12); \draw[dep] (c11) -- (c13); \draw[dep] (c5) -- (c13);
\draw[dep] (c13) -- (c14); \draw[dep] (c13) -- (c15); \draw[dep] (c13) -- (c16);
\draw[dep] (c15) -- (c17); \draw[dep] (c9) -- (c17); \draw[dep] (c17) -- (c21); \draw[dep] (c21) -- (c22);
\draw[dep] (c16) -- (c18); \draw[dep] (c18) -- (c19); \draw[dep] (c16) -- (c20);
% legend
\begin{scope}[yshift=-7.25cm]
\foreach \s/\t [count=\i] in {base/기초, bus/채널과 모드, sum/종합, io/입력과 출력, sort/뉴런 분류, lab/실험실}{
  \pgfmathtruncatemacro{\col}{mod(\i-1,3)}\pgfmathtruncatemacro{\row}{(\i-1)/3}
  \node[\s,text width=0.3cm,minimum height=0.32cm] at ({\col*4.8+0.2},{-\row*0.55}) {};
  \node[anchor=west,font=\scriptsize] at ({\col*4.8+0.55},{-\row*0.55}) {\t};}
\end{scope}
\end{tikzpicture}
\caption{장 지도. 화살표는 한 장에서 그 장에 기대는 장으로 향한다. 장 사이의 나머지 참조는 안내일 뿐 의존 관계가 아니다.}
\label{fig:roadmap}
\end{figure}

\section*{이 책을 지나는 경로}

\begin{center}
\footnotesize
\setlength{\tabcolsep}{4pt}
\renewcommand{\arraystretch}{1.2}
\begin{tabular}{>{\raggedright}p{3.0cm}>{\raggedright}p{3.9cm}>{\raggedright\arraybackslash}p{6.4cm}}
\toprule
경로 & 대상 & 장 순서 \\
\midrule
한 쿼터 강의 & 강의자, 10주 & 1주: 1--2장; 2주: 3장; 3주: 4장; 4주: 5--6장; 5주: 7--8장; 6주: 9장; 7주: 10장; 8주: 11--12장; 9주: 13장과~15장; 10주: 17장 \\
AI와 머신러닝 & 신경망을 알고 뇌가 어떻게 학습하는지 이해하고 싶은 사람 & 1--4, 5(훑어보기), 6, 7, 8, 9, 11(훑어보기), 12, 13(훑어보기), 15 \\
전자공학과 하드웨어 & 회로 설계자, 뉴로모픽 칩 개발자 & 1--4, 5, 11, 13, 16, 18, 19, 17(9장과 15장은 훑어보기) \\
뉴로인터페이스와 살아 있는 계산기 & 뇌-컴퓨터 인터페이스나 살아 있는 뉴런을 쓰는 칩을 만드는 사람 & 1, 2, 4, 11, 13, 17(9장과 15장은 훑어보기), 21, 22(12장은 훑어보기) \\
빠른 개관 & 주말 시간이 있는 엔지니어 & 1, 2, 3, 6, 10; 6장에서 4--5장을 가리키는 참조는 문맥으로 이해할 수 있다 \\
\bottomrule
\end{tabular}
\end{center}

“훑어보기”란 장의 도입부, 노란 상자 안의 명제들, 마지막 요약 표만 읽어도 충분하다는 뜻이다.

각 장 끝에는 연습문제가 있다. 어림 계산, 유도, 모델링, 설계 문제이며, 간단한 답은 부록~\ref{sec:answers}에 모아 두었다. 모든 기호는 부록~\ref{sec:notation}에, 각 장의 주요 명제를 뒷받침하는 실험은 부록~\ref{sec:supplement}에 정리했다. 책의 수치는 \texttt{models/} 폴더의 작은 프로그램들로 계산했다. 직접 실행하고 고쳐 볼 수 있다.

\mainmatter

\chapter{뉴런의 유형}
\label{sec:neurontypes}

\section{블랙박스로서의 뉴런}

뉴런 860억 개가 든 자루를 받았다고 하자~\cite{Azevedo2009}. 이것을 몇 무더기로 나누라고 하면 어떻게 하겠는가?

낯선 칩이 가득 든 자루를 받은 엔지니어라면 패키지 모양으로 분류하지 않을 것이다. 그는 이렇게 물을 것이다. 이 부품에는 입력이 몇 개이고 출력이 몇 개이며, 출력으로 무엇을 내보내는가? 뉴런도 그렇게 회로도의 블록으로 그려 보자(그림~\ref{fig:blackbox}).

\begin{figure}[t]
\centering
\begin{tikzpicture}[>=Stealth,x=1cm,y=1cm]
% the box
\node[draw,very thick,minimum width=2.6cm,minimum height=2.6cm,fill=blue!6] (N) at (0,0) {};
\node at (0,0.55) {\small 뉴런};
\node at (0,-0.15) {$\displaystyle u=\sum_i w_i x_i$};
\node at (0,-0.8) {\small $y=\Theta(u-\theta)$};
% inputs
\foreach \k/\y/\lab in {1/1.05/{x_1},2/0.55/{x_2},3/0.05/{x_3}}{
  \draw[->,thick,blue!60!black] (-4.2,\y) -- (-1.3,\y);
  \node[left] at (-4.2,\y) {$\lab$};
  \node[above,blue!60!black] at (-2.1,\y-0.06) {\scriptsize $w_{\k}$};}
\node at (-4.45,-0.4) {$\vdots$};
\node at (-2.1,-0.4) {$\vdots$};
\draw[->,thick,blue!60!black] (-4.2,-1.05) -- (-1.3,-1.05);
\node[left] at (-4.2,-1.05) {$x_n$};
\node[above,blue!60!black] at (-2.1,-1.11) {\scriptsize $w_{n}$};
\draw[decorate,decoration={brace,amplitude=5pt},gray!60!black] (-4.9,-1.2) -- (-4.9,1.2);
\node[left,align=right,gray!40!black] at (-5.1,0) {\small 입력\\[-2pt]\small $n\sim10^3$--$10^5$개};
% single output wire, then fan-out
\draw[very thick,red!70!black] (1.3,0) -- (3.2,0);
\foreach \x in {1.6,2.2,2.34,2.9}{\draw[thick,red!70!black] (\x,0) -- (\x,0.4);}
\node[above right,font=\tiny\itshape,gray!30!black,inner sep=1pt] at (1.42,0.45) {딸깍\dots\ 딸깍-딸깍\dots};
\node[below,red!70!black,align=center] at (2.25,-0.05) {\scriptsize 출력 $y$ 하나};
\fill[red!70!black] (3.2,0) circle (0.06);
\foreach \y in {1.3,0.65,0,-0.65,-1.3}{
  \draw[->,thick,red!70!black] (3.2,0) -- (5.0,\y);}
\draw[decorate,decoration={brace,amplitude=5pt},gray!60!black] (5.3,1.45) -- (5.3,-1.45);
\node[right,align=left,gray!40!black] at (5.5,0) {\small $y$의 복사본이\\[-2pt]\small 다른 뉴런\\[-2pt]\small $10^3$--$10^4$개로};
\end{tikzpicture}
\caption{엔지니어의 눈으로 본 뉴런. 입력 $x_i$가 많고 각각 자기 가중치 $w_i$를 가진다. 안에서는 가중합 $u$를 문턱값 $\theta$와 비교한다($\Theta$는 계단 함수로, 인수가 양수이면 $1$, 아니면 $0$이다). 출력 $y$는 하나이며, 스파이크의 열이 가지를 치며 복제되어 수천 개의 수신자에게 전달된다.}
\label{fig:blackbox}
\end{figure}

블록의 출력은 숫자가 아니라 짧은 전압 스파이크의 열이다. “딸깍”, 쉼, “딸깍-딸깍”, 긴 쉼. 스파이크는 모두 높이가 같으므로 모든 정보는 스파이크가 \emph{언제} 일어나는지에 담겨 있다. 블록 안에는, 거칠게 1차 근사로 말하면, 가산기와 비교기가 있다. 입력에 가중치를 곱해 더하고, 그 합이 문턱값을 넘으면 블록은 스파이크를 낸다. 다음 장에서는 이것을 연속 시간에서 제대로 동작하는 모델로 바꿀 것이다.

출력은 하나지만 전선은 가지를 치고, 가지마다 \emph{똑같은} 신호 복사본이 흐른다. 전자공학에서는 이것을 팬아웃(fan-out)이라고 부른다. 대뇌 피질 뉴런의 팬아웃은 수천 정도이다. 칩 안의 논리 게이트는 보통 다른 게이트 몇 개만 구동한다.

그렇다면 출력 전선을 따라 수신자에게 \emph{무엇이} 전달되는가? 스파이크는 각 가지의 끝까지 달려가 거기서 화학 신호로 바뀐다. 가지 끝은 세포 사이의 좁은 틈으로 특정 물질, 즉 \emph{신경전달물질}을 조금 뿜어내고, 이 물질이 수신자의 입력 전류를 바꾼다. 바로 여기에 분류 기준이 있다. \emph{뉴런의 출력이 어떤 신경전달물질을 내보내는가}이다.

\section{출력 하나, 신호 유형 하나}

분류가 의미를 가지려면 뉴런마다 출력 유형이 하나뿐이어야 한다. 정말 그런가? 실험에 따르면 거의 그렇다.

\begin{keyblock}
\begin{proposition}[뉴런 하나, 출력 유형 하나]
\label{prop:dale}
거의 모든 뉴런에는 주된 신경전달물질이 하나 있고, 뉴런은 출력의 모든 가지에서 그것을 내보낸다. 따라서 뉴런의 유형은 주된 신경전달물질로 정해진다~\cite{Strata1999}.
\end{proposition}
\end{keyblock}

뉴런의 유형은 \emph{주된 신경전달물질 영어 이름의 처음 네 글자}로 표시하기로 하자. 예를 들어 주된 신경전달물질이 글루탐산(glutamate)인 뉴런은 \nt{GLUT} 뉴런이다. 모든 유형은 표~\ref{tab:types}에 모아 두었다.

\begin{table}[t]
\centering
\footnotesize
\setlength{\tabcolsep}{3pt}
\renewcommand{\arraystretch}{1.1}
\begin{tabular}{>{\raggedright}p{1.0cm}>{\raggedright}p{2.25cm}>{\raggedright}p{1.95cm}>{\raggedright}p{1.8cm}>{\centering}p{0.7cm}>{\raggedright}p{1.55cm}>{\raggedright}p{1.8cm}>{\raggedright\arraybackslash}p{3.2cm}}
\toprule
유형 & 주된 신경전달물질 & 버스 & 속도 & 부호 & 뇌 속 뉴런 수 & 뉴런당 출력 수 & 계산에서의 역할 \\
\midrule
\nt{GLUT} & 글루탐산 & 주소 & 빠름과 느림 & $+$ & $\sim8\cdot10^{10}$ & $\sim10^4$ & 주된 양의 가중치 \\
\nt{GABA} & GABA & 주소 & 빠름과 느림 & $-$ & $\sim3\cdot10^{9}$ & $10^3$--$10^4$ & 주된 음의 가중치 \\
\nt{GLYC} & 글리신 & 주소 & 빠름 & $-$ & $10^6$--$10^7$ & $10^2$--$10^3$ & 척수와 뇌줄기의 음의 가중치; 글루탐산 수용체 중 하나의 두 번째 열쇠 \\
\nt{ACET} & 아세틸콜린 & 둘 다 & 빠름과 느림 & $\pm$ & $\sim10^6$ & 근육: $10$--$10^3$; 뇌: $\sim10^5$ & 근육으로 가는 출력; 뇌에서는 학습률(가설) \\
\nt{DOPA} & 도파민 & 브로드캐스트 & 느림 & $\pm$ & $\sim5\cdot10^{5}$ & $10^5$--$10^6$ & 보상 예측 오차 $\delta$ \\
\nt{SERO} & 세로토닌 & 브로드캐스트 & 느림(수용체 한 종류는 빠름) & $\pm$ & $\sim3\cdot10^{5}$ & $\sim10^5$ & 계획 지평(가설) \\
\nt{HIST} & 히스타민 & 브로드캐스트 & 느림 & $+$ & $\sim6\cdot10^{4}$ & 추정치 없음 & 전역 신호 “깨어 있어라” \\
\nt{NORA} & 노르아드레날린 & 브로드캐스트 & 느림 & $\pm$ & $\sim5\cdot10^{4}$ & $\sim10^5$ & 이득(가설) \\
\nt{ADRE} & 아드레날린 & 브로드캐스트 & 느림 & $\pm$ & $\sim10^4$ & 추정치 없음 & 뉴런이 적다; 뇌에서의 역할은 불분명 \\
\nt{ASPA} & 아스파르트산 & 주소 & 빠름 & $+$ & 추정치 없음 & 추정치 없음 & 약한 양의 가중치; 역할은 논쟁 중 \\
\bottomrule
\end{tabular}
\caption{뉴런의 유형. 뇌 속 뉴런 수가 많은 순서로 나열했다. \emph{버스}: 주소 버스는 신경전달물질을 특정 수신자로 가는 좁은 틈에 내보내는 방식이고, 브로드캐스트는 소수의 원천 뉴런 무리가 넓은 영역에 내보내는 방식이다(\ref{sec:buses}절). \emph{속도}: 빠름은 이온성 수용체를 거쳐 밀리초 단위, 느림은 대사성 수용체를 거쳐 수백 밀리초 이상이다. \emph{부호}는 통상적인 부호이다. 최종 부호는 수신자가 정하며(\ref{sec:receiver}절), $\pm$는 두 부호가 모두 나타난다는 뜻이다. \emph{뇌 속 뉴런 수}는 사람의 뇌 전체에 있는 해당 유형 뉴런의 수를 자릿수 수준으로 나타낸 것이다. 뉴런은 모두 약 $8.6\cdot10^{10}$개이다~\cite{Azevedo2009}. \nt{GLUT} 뉴런의 거의 전부, 약 $7\cdot10^{10}$개는 소뇌의 작은 입력 뉴런이다. \nt{GLUT}와 \nt{GABA}의 수는 이 집계와 피질 속 \nt{GABA} 뉴런의 비율에서 얻었다~\cite{Hendry1987}. 수가 적은 유형 가운데 직접 센 것은 \nt{HIST}~\cite{Airaksinen1991}와 \nt{DOPA} 뉴런의 두 무리 중 큰 쪽~\cite{Pakkenberg1991}뿐이므로, \nt{DOPA}의 값은 하한이다. \nt{SERO}는 두 부분 집계의 합이다~\cite{Baker1990,Baker1991}. \nt{NORA}, \nt{GLYC}, \nt{ACET}, \nt{ADRE}의 수는 직접 세지 않은 거친 자릿수 추정치이다. \emph{뉴런당 출력 수}는 뉴런 하나가 가진 시냅스 말단, 즉 출력 전선의 가지에서 신경전달물질을 내보내는 자리(그림~\ref{fig:blackbox})의 수를 자릿수 수준으로 나타낸 것이다. 브로드캐스트 유형의 말단은 직접 세지 않았다. \nt{DOPA}의 값은 출력 전선 하나의 가지가 표적 영역의 얼마만큼을 덮는지에서 유도했고~\cite{Matsuda2009}, 나머지 유형의 추정은 더 거칠다. \nt{ACET}에는 두 경우가 있다. 근육을 제어하는 뉴런과 뇌 속의 브로드캐스트 뉴런이다.}
\label{tab:types}
\end{table}

% the pie is a table-type float with a figure caption, so it can never float ahead of Table~\ref{tab:types}
\begin{table}[tb]
\makeatletter\def\@captype{figure}\makeatother
\centering
\definecolor{vizglut}{HTML}{2A78D6}
\definecolor{vizgaba}{HTML}{E34948}
\definecolor{vizrest}{HTML}{8A8986}
\begin{tikzpicture}[>=Stealth]
% pie: GLUT 96.4 %, GABA 3.6 % (13.0 degrees), the rest 0.006 % (a hairline at 0 degrees)
\fill[vizglut] (0,0) -- (13:1.9) arc (13:360:1.9) -- cycle;
\fill[vizgaba] (0,0) -- (0:1.9) arc (0:13:1.9) -- cycle;
\draw[white,line width=1pt] (0,0) -- (13:1.9);
\draw[vizrest,line width=1.2pt] (0,0) -- (0:1.9);
\node[white,align=center,font=\small] at (190:0.95) {\nt{GLUT}\\$96.4\,\%$};
\draw[gray!60!black] (6.5:1.75) -- (2.05,2.1);
\node[anchor=west,font=\small] at (2.05,2.1) {\nt{GABA} $3.6\,\%$};
% zoom box for the remaining types
\draw[densely dashed,vizrest] (1.9,0) -- (3.1,1.65);
\draw[densely dashed,vizrest] (1.9,0) -- (3.1,-2.2);
\draw[vizrest,rounded corners=3pt] (3.1,-2.2) rectangle (10.1,1.65);
\node[anchor=west,font=\scriptsize] at (3.2,1.4) {나머지 전부 합쳐서: $\approx0.006\,\%$; 로그 눈금};
\foreach \code/\lg/\y/\val/\lx in {GLYC/6.477/1.0/{$10^6$--$10^7$}/4.05,ACET/6.0/0.6/{$\sim10^6$}/3.0,DOPA/5.699/0.2/{$\sim5\cdot10^5$}/2.699,SERO/5.477/-0.2/{$\sim3\cdot10^5$}/2.477,HIST/4.778/-0.6/{$\sim6\cdot10^4$}/1.778,NORA/4.699/-1.0/{$\sim5\cdot10^4$}/1.699,ADRE/4.0/-1.4/{$\sim10^4$}/1.0}{
  \node[anchor=east,font=\scriptsize] at (4.35,\y) {\nt{\code}};
  \fill[vizrest] (4.45,\y-0.13) rectangle ({4.45+\lg-3},\y+0.13);
  \node[anchor=west,font=\scriptsize] at ({4.5+\lx},\y) {\val};}
\draw[vizrest,thick] (7.45,1.0) -- (8.45,1.0);
\draw[vizrest,thick] (7.45,0.92) -- (7.45,1.08);
\draw[vizrest,thick] (8.45,0.92) -- (8.45,1.08);
\draw[gray!60!black] (4.45,-1.72) -- (8.45,-1.72);
\foreach \e in {3,4,5,6,7}{
  \draw[gray!60!black] ({4.45+\e-3},-1.72) -- ({4.45+\e-3},-1.66);
  \node[below,font=\scriptsize] at ({4.45+\e-3},-1.72) {$10^{\e}$};}
\end{tikzpicture}
\caption{표~\ref{tab:types}의 추정치로 본 뇌 속 뉴런 유형의 비율. 원은 거의 전부 \nt{GLUT}가 차지한다. \nt{GABA}는 좁은 부채꼴이고, 나머지 유형은 모두 합쳐도 경계선 위의 머리카락 한 올이다. 오른쪽은 그 머리카락을 확대해 뉴런 수를 로그 눈금으로 나타낸 것이다. \nt{GLYC}의 막대는 범위 $10^6$--$10^7$의 가운데까지 그렸고, 선분은 범위 전체를 보여 준다. \nt{ASPA}는 추정치가 없어 그림에 넣지 않았다.}
\label{fig:typepie}
\end{table}

이 무더기들이 얼마나 불균등한지는 그림~\ref{fig:typepie}에 나와 있다. 뉴런 천 개 가운데 $964$개가 \nt{GLUT}, $36$개가 \nt{GABA}이고, 나머지 유형을 모두 합쳐도 십만 개 중 여섯 개 정도이다.

GABA는 원래 네 글자로 된 약칭이다. 주된 신경전달물질이 되는 일이 거의 없는 물질, 즉 신경펩타이드, 기체, 지질, 퓨린은 따로 코드를 받지 않는다. 이들은 부록~\ref{sec:othermediators}에서 다룬다. 명제~\ref{prop:dale}의 “거의”라는 말은 중요하다. 많은 뉴런이 주된 신경전달물질과 함께 보조 물질을 내보낸다~\cite{Yao2023}. 어떤 뉴런은 두 번째 빠른 신경전달물질을, 그것도 반대 부호의 것을 내보낸다. \nt{DOPA} 뉴런의 일부는 GABA도 내보낸다~\cite{Tritsch2012}. 또 어떤 뉴런은 같은 출력의 서로 다른 가지에서 서로 다른 신경전달물질을 내보낸다~\cite{Zhang2015}. 이 모든 것이 측정되었으므로 “뉴런 하나, 신경전달물질 하나”는 예외가 있는 규칙이지 법칙이 아니다. 그래도 첫 번째 분류로서는 검증을 견뎌 낸다. 발현하는 유전자에 따라 뉴런을 5000개가 넘는 무리로 나눈 생쥐 뇌 세포 아틀라스에서도, 뉴런의 큰 부류는 여전히 무엇보다 주된 신경전달물질에 따라 갈린다~\cite{Yao2023}.

이 규칙에는 뇌를 인공 신경망과 곧바로 구별해 주는 결과가 있다. 인공 신경망에서는 같은 뉴런이 한 수신자에게는 양의 가중치를, 다른 수신자에게는 음의 가중치를 가질 수 있다. 가중치 $w_{ij}$는 행렬 속의 숫자일 뿐이다. 뇌에서는 작용의 부호가 주로 신경전달물질로 정해지고, 뉴런의 신경전달물질은 하나이다. 따라서 뉴런 $j$가 한 수신자를 흥분시키면 다른 모든 수신자도 흥분시킨다. 뉴런 $j$에서 나가는 가중치 행렬의 \emph{열} 전체가 같은 부호를 가진다.
\begin{empheq}[box=\keybox]{equation}
\operatorname{sign} w_{ij} \;=\; s_j \quad\text{모든 수신자 } i\text{에 대해}, \qquad s_j\in\{+1,-1\}.
\label{eq:dalesign}
\end{empheq}
뉴런은 흥분성($s_j=+1$)과 억제성($s_j=-1$)으로 나뉘며, 이것은 연결이 아니라 뉴런 자체의 성질이다. 흥분성 뉴런에서 음의 연결이 필요하면 중개자를 거쳐야 한다. 흥분성 뉴런이 억제성 뉴런을 흥분시키고, 억제성 뉴런이 표적을 억제한다. 한 단계만큼 지연된 음의 가중치이다.

알려진 신경전달물질은 백 가지가 넘지만, 뇌가 하는 일의 거의 전부는 여섯 가지 남짓에 기대고 있으며, 이들은 전혀 다른 두 부류로 나뉜다.

\section{두 가지 버스: 주소 버스와 브로드캐스트}
\label{sec:buses}

컴퓨터 네트워크에는 메시지를 전달하는 방법이 두 가지 있다. 첫째는 \emph{주소 지정} 방식, 즉 유니캐스트(unicast)이다. 패킷이 한 송신자에서 특정 수신자로 빠르게 가고, 다른 누구도 그것을 보지 못한다. 둘째는 \emph{브로드캐스트}(broadcast)이다. 송신자가 메시지 하나를 네트워크의 모든 노드에 한꺼번에 보낸다. 뉴런은 두 방식을 모두 쓰며, 신경전달물질은 정확히 이 기준으로 나뉜다(그림~\ref{fig:phone_radio}).

\begin{figure}[t]
\centering
\begin{tikzpicture}[>=Stealth,scale=0.95]
% ---- unicast: point-to-point wiring
\node[above] at (2.0,3.3) {\small \textbf{주소 버스}: 글루탐산과 GABA};
\foreach \i in {0,...,4}{
  \fill[blue!60!black] (0.2+0.9*\i,2.5) circle (0.1);
  \fill[blue!60!black] (0.2+0.9*\i,0.6) circle (0.1);}
\foreach \a/\b in {0/0,0/1,1/1,1/3,2/2,3/2,3/4,4/4,2/0}{
  \draw[->,blue!60!black] (0.2+0.9*\a,2.38) -- (0.2+0.9*\b,0.72);}
\fill[red!70!black] (1.1,1.55) circle (0.09);
\pic[scale=1.2] at (1.75,1.9) {letter};
\draw[->,red!70!black,thick] (1.1,1.55) -- (0.28,0.7);
\draw[->,red!70!black,thick] (1.1,1.55) -- (1.95,0.72);
\node[align=center,below] at (2.0,0.2) {\small 뉴런 수십억 개,\\[-2pt]\small 저마다 자기 수신자가 있다,\\[-2pt]\small 시간: 밀리초};
% ---- broadcast: one small source, global targets
\begin{scope}[xshift=7.2cm]
\node[above] at (1.8,3.3) {\small \textbf{브로드캐스트}: 도파민, \dots};
\foreach \i in {0,...,8}{\fill[blue!60!black] (-0.2+0.5*\i,2.75) circle (0.08);}
\fill[orange!85!black] (1.8,0.35) ellipse (0.35 and 0.18);
\pic[scale=1.5] at (2.7,0.75) {mast};
\foreach \x in {-0.2,0.3,0.8,1.3,1.8,2.3,2.8,3.3,3.8}{
  \draw[->,orange!85!black] (1.8,0.5) .. controls (1.8,1.3) and (\x,1.6) .. (\x,2.62);}
\node[align=center,below] at (1.8,0.1) {\small 뉴런 수천--수십만 개,\\[-2pt]\small 모두에게 같은 메시지,\\[-2pt]\small 시간: 수백 밀리초 이상};
\end{scope}
\end{tikzpicture}
\caption{두 가지 전송 방식. 왼쪽은 주소 버스로, 봉투에 주소를 적은 우편과 비슷하다. 뉴런 하나와 그 수신자 둘을 빨간색으로 표시했다. 오른쪽은 브로드캐스트로, 라디오 방송국과 같다. 작은 원천 뉴런 무리가 있고, 모두가 같은 메시지를 받는다.}
\label{fig:phone_radio}
\end{figure}

\emph{주소 버스}는 글루탐산과 GABA이다. 압도적 다수의 뉴런이 이 둘을 내보낸다. 각 출력은 특정 세포로 이어지고, 신경전달물질은 접촉부의 좁은 틈 안에서만, 그것도 빠르게 작용한다. 1밀리초 만에 수신자의 이온 채널이 열리고, 몇 밀리초면 모든 것이 끝난다. 이것은 \emph{데이터} 버스이다. 뇌가 계산하는 내용이 이 버스로 전달된다.

\emph{브로드캐스트 버스}는 도파민, 세로토닌, 노르아드레날린, 그리고 부분적으로 아세틸콜린이다. 이들을 내보내는 것은 아주 작은 뉴런 무리이지만, 각 뉴런의 출력은 믿기 어려울 만큼 넓게 가지를 친다. 신경전달물질은 좁은 틈보다는 세포 사이 공간으로 방출되는 경우가 많고, 거기서 퍼져 나가 주변의 모든 세포에 작용한다. 작용은 느리다. 채널을 직접 여는 대신 수신자 안에서 화학 반응의 연쇄를 시작한다. 이 버스로는 데이터가 아니라 넓은 영역에 공통인 \emph{제어 파라미터}가 전달되며, 이것이 네트워크의 작동 모드를 바꾼다. 이득, 학습률, “잘했다”는 신호 같은 것이다. 그래서 이런 신경전달물질을 \emph{신경조절물질}이라고 부른다. 오랫동안 이런 채널 하나하나가 뇌 전체에 대해 숫자 하나라고 여겨졌다. 최근의 측정은 그림을 다듬었다. 채널은 전선 한 가닥이 아니라 병렬 선로 몇 가닥을 묶은 작은 다발이다. 한 신경전달물질의 원천 뉴런들은 하위 시스템으로 나뉘어 저마다 자기 수신자와 자기 메시지를 가지며, 그 자리에서 얼마나 방출할지는 국소 신호로도 조정된다~\cite{Ren2018,Poe2020}. 이런 선로의 수는 수십억이 아니라 몇 개에서 수십 개이므로, 채널은 여전히 브로드캐스트이다.

이 장에서 그림 하나만 가져가야 한다면 바로 이것이다. 뇌는 주소 지정 방식으로 데이터를 전달하는 거대한 네트워크이고, 그 위에 제어 신호를 싣는 브로드캐스트 채널 몇 개가 걸려 있다. 프로그래머라면 익숙한 구조를 알아볼 것이다. 계산이 있고, 그 위에 각각 네트워크의 넓은 영역에 공통인 제어 변수 몇 개가 있다.

\section{\nt{GLUT}: 양의 가중치}

글루탐산은 뇌의 주된 흥분성 신경전달물질이다. 뉴런의 출력이 글루탐산을 내보내면 수신자에서 채널이 열려 나트륨이 안으로 들어오고, 수신자의 막전압이 올라가 스스로 스파이크를 낼 가능성이 커진다. 그림~\ref{fig:blackbox}의 언어로 말하면 이것은 \emph{양의} 가중치 $w_i>0$을 가진 입력이다.

글루탐산은 누가 내보내는가? 데이터를 먼 거리로 전달하는 거의 모든 뉴런이다. 피질 뉴런의 4분의 3에서 5분의 4, 그리고 뇌의 여러 영역을 잇는 뉴런 대부분이 여기에 속한다. 뇌의 네트워크는 무엇보다 글루탐산 연결의 네트워크이다.

공학적 세부 사항 하나. 방출 직후 틈 속 글루탐산 농도는 수천 배로 치솟아 약 1밀리몰에 이르고, 약 1밀리초의 시간 상수로 떨어진다~\cite{Clements1992}. 글루탐산이 틈 밖으로 퍼져 나가고, 특수한 펌프 단백질이 그것을 다시 세포 안으로 퍼 올리기 때문이다. 왜 그럴까? 통신 선로에서 기호 하나를 보낼 때마다 신호를 0으로 되돌리는 것과 같은 이유이다. 신경전달물질을 치우지 않으면 다음 스파이크는 “통화 중”인 선로에 도착하고, 몇 밀리초 간격의 스파이크들이 서로 뭉개진다.

\section{\nt{GABA}: 음의 가중치}

모든 가중치가 양수인 네트워크를 상상해 보자. 평균적으로 스파이크 하나가 다음 스파이크를 하나보다 많이 낳으면, 활동은 지수적으로 늘어나 몇 단계 만에 네트워크 전체를 뒤덮는다. 전자공학자라면 이득이 1보다 큰 양의 되먹임 루프를 알아볼 것이다. 회로는 포화로 치닫는다. 이를 막으려면 음의 가중치가 필요하다. 그 주된 원천은 감마-아미노뷰티르산, 줄여서 GABA이다.

GABA는 수신자에서 염화 이온을 통과시키는 채널을 연다. 염소의 평형 전위는 휴지 전위 근처나 그보다 약간 아래에 있으므로, 열린 염소 채널은 막을 휴지 전위 근처에 붙잡아 두고 전압이 문턱값까지 오르지 못하게 한다. 이것은 \emph{음의} 가중치 $w_i<0$을 가진 입력이다. 피질 뉴런의 5분의 1에서 4분의 1이 GABA 뉴런이다~\cite{Hendry1987}. 이들의 출력은 짧고, 가장 가까운 이웃을 억제한다.

뇌에서 억제는 단순한 뺄셈보다 훨씬 많은 일을 한다. 억제는 네트워크 전체를 동작점에 붙잡아 두는 음의 되먹임으로 작동한다. 정상적으로 작동하는 피질에서는 뉴런으로 들어오는 흥분성 전류와 억제성 전류가 거의 정확히 균형을 이룬다고 여겨진다. 이 체제는 이론으로 예측되었고~\cite{vanVreeswijk1996} 살아 있는 피질의 전류 측정에서도 보이며~\cite{Haider2006}, 뉴런은 늘 문턱값 근처에 머문다. 불안정한 점 근처에서 강한 음의 되먹임으로 안정화된 회로는 작은 신호에 빠르고 민감하게 반응한다. 오래된 공학 기법이다.

\section{\nt{DOPA}: 오차 신호}

첫 번째 브로드캐스트 채널은 도파민이다. 사람의 도파민 뉴런은 50만 개 정도, 대략 17만 개 중 하나이다. 이것은 두 무리 중 큰 쪽에서 센 수이므로~\cite{Pakkenberg1991} 하한이다. 이들의 출력은 행동을 고르는 영역들로 퍼져 나간다.

이 채널은 무엇을 전달하는가? 보상을 받는 동물에서 도파민 뉴런의 스파이크를 기록한 실험이 답을 준다~\cite{Schultz1997}. 동물에게 이따금 예고 없이 주스 한 방울을 주면, 도파민 뉴런은 짧은 스파이크 버스트로 응답한다. 그다음에는 주스를 주기 전에 전등을 켜는데, 언제나 주스보다 1초 앞서 켠다. 동물이 이 관계를 익히면 뉴런은 \emph{전등}에 버스트로 응답하기 시작하고, 제때 정확히 온 주스 자체에는 전혀 응답하지 않는다. 그리고 전등을 켠 뒤 주스를 주지 않으면, 주스가 와야 했던 순간에 뉴런은 \emph{잠잠해지고} 발화율이 평소보다 떨어진다.

세 경우를 모두 설명하는 규칙은 이것이다(그림~\ref{fig:rpe}). 뉴런은 얼마나 좋은지가 아니라 \emph{기대보다 얼마나 좋은지 또는 나쁜지}를 알린다.

\begin{figure}[t]
\centering
\begin{tikzpicture}[>=Stealth,x=3.4cm,y=1cm]
\foreach \y/\lab in {2.8/{예고 없는 주스},1.4/{전등 뒤의 주스},0/{전등, 그러나 주스 없음}}{
  \draw[gray!70!black] (0,\y) -- (3,\y);
  \draw[densely dotted,gray!60!black] (0,\y+0.3) -- (3,\y+0.3);
  \node[left,align=right,font=\small] at (-0.03,\y+0.35) {\lab};}
% row 1: juice without a cue -> burst at the juice
\draw[very thick,orange!85!black] plot[domain=0:3,samples=120] (\x,{2.8+0.3+0.75*exp(-((\x-2.08)/0.07)^2)});
\pic[scale=0.8] at (2,2.58) {juice};
% row 2: cue then juice -> burst at the cue, nothing at the juice
\draw[very thick,orange!85!black] plot[domain=0:3,samples=120] (\x,{1.4+0.3+0.75*exp(-((\x-1.08)/0.07)^2)});
\pic[scale=0.85] at (1,1.15) {lamp}; \pic[scale=0.85] at (2,1.15) {juice};
% row 3: cue, juice omitted -> burst at the cue, dip when the juice was due
\draw[very thick,orange!85!black] plot[domain=0:3,samples=120] (\x,{0.3+0.75*exp(-((\x-1.08)/0.07)^2)-0.28*exp(-((\x-2.1)/0.12)^2)});
\pic[scale=0.85] at (1,-0.25) {lamp}; \pic[scale=0.85] at (2,-0.25) {nojuice};
\node[right,font=\scriptsize] at (2.36,0.1) {움푹 꺼짐};
\node[right,font=\scriptsize] at (2.13,3.75) {깜짝!};
\node[left,font=\scriptsize] at (0.97,2.25) {전등에 대한 버스트};
\node[above,font=\scriptsize] at (2,1.75) {반응 없음};
\draw[->,gray!70!black] (0,-0.75) -- (3.05,-0.75) node[right,black,font=\small] {$t$};
\draw[gray!60!black] (1,-0.8) -- (1,-0.7) node[below=2pt,font=\scriptsize] {전등, $1\,\text{s}$};
\draw[gray!60!black] (2,-0.8) -- (2,-0.7) node[below=2pt,font=\scriptsize] {주스, $2\,\text{s}$};
\end{tikzpicture}
\caption{세 실험에서 \nt{DOPA} 뉴런의 발화율(개략도, \cite{Schultz1997}에 따름). 점선은 일정한 배경 발화율이다. 전등은 신호, 물방울은 주스, 줄 그은 물방울은 약속했지만 주지 않은 주스이다. 응답은 예측 오차~\eqref{eq:rpe}이다.}
\label{fig:rpe}
\end{figure}

강화학습~\cite{SuttonBarto2018}을 아는 프로그래머라면 이 양을 바로 알아볼 것이다. 행동을 고르는 법을 배우는 에이전트는 추정치 $V(s)$, 즉 상태 $s$에 있을 때 기대하는 보상의 양을 저장한다. 에이전트는 매 단계 뒤에 기대와 실제로 받은 것을 비교하고,
\begin{empheq}[box=\keybox]{equation}
\delta_t \;=\; r_t \;+\; \gamma\,V(s_{t+1}) \;-\; V(s_t) ,
\label{eq:rpe}
\end{empheq}
추정치를 고친다. $V(s_t)\leftarrow V(s_t)+\alpha\,\delta_t$. 여기서 $r_t$는 받은 보상, $\gamma<1$은 할인율(미래의 보상이 즉각적인 보상보다 얼마나 덜 가치 있는가), $\alpha$는 학습률이다. 양 $\delta_t$를 \emph{시간차 오차}라고 한다.

이 양은 도파민 뉴런과 똑같이 행동한다. 예고 없는 주스에서는 $V$가 아직 0이고 $r>0$이므로 $\delta>0$이다. 학습된 주스에서는 전등이 보상을 예측했으므로 주스 직전의 $V(s_t)$가 이미 $r$과 같고 $\delta=0$이다. 주스가 오지 않으면 $V(s_t)=r$인데 $r_t=0$이므로 $\delta<0$이다. 그리고 버스트가 전등으로 옮겨 가는 것은 바로 전등의 순간에 기대가 한꺼번에 올라가기 때문이다. $\gamma V(s_{t+1})-V(s_t)>0$. 여기서 단계 $t,t+1,\dots$는 알고리즘의 약속일 뿐이다. 몸에는 단계를 세는 시계가 없으며, 같은 양을 연속 시간에서 \ref{sec:dofa}장에서 얻을 것이다.

따라서 도파민은 스칼라 오차 신호 $\delta$를 그것으로 학습해야 하는 모든 곳에 브로드캐스트하는 것이다. 1차 근사로는 뇌 전체에 전선 하나, 매 순간 숫자 하나이다. 이 전선이 실제로 얼마나 다발에 가까운지는 \ref{sec:dofaalt}장에서 살펴본다.

\section{나머지 브로드캐스트 채널}

나머지 신경조절물질에 대해서는 그 역할을 같은 전역 파라미터의 언어로 옮기는 작업 가설들만 있다~\cite{Doya2002}. 어디까지나 가설로 받아들이기 바란다.

\emph{\nt{NORA}, 노르아드레날린: 이득.} 노르아드레날린 뉴런은 도파민 뉴런보다도 적어서 거칠게 추정하면 수만 개 정도이지만, 그 출력은 사실상 뇌 전체로 퍼진다. 이들은 뜻밖의 일이나 중요한 일이 생기면 급격히 활성화된다. 노르아드레날린은 수신 뉴런의 전달 특성의 기울기를 바꾼다. 약한 입력은 더 약해지고 강한 입력은 더 강해진다. 측정해 보면 수신자의 배경 스파이크가 입력에 대한 응답보다 더 강하게 억제되어 신호 대 잡음비가 올라간다~\cite{Foote1975NE}. 간단한 모델에서는 이것을 숫자 하나로 기술한다. 뉴런이 입력 전류 $u$에 발화율 $f=F\bigl(g\cdot u\bigr)$로 응답한다면, 노르아드레날린은 계수 $g$를 키운다~\cite{AstonJones2005}(자세한 내용은 \ref{sec:nora}장). $g$가 큰 네트워크는 “대비가 뚜렷하게” 작동하며 결정을 더 빨리 내린다. $g$가 작으면 “흐릿하게” 작동하며, 탐색 모드의 강화학습 에이전트처럼 여러 선택지를 더 많이 시도한다.

\emph{\nt{ACET}, 아세틸콜린: 학습률.} 아세틸콜린은 네트워크가 현재 입력에 얼마나 귀 기울이고 저장된 자기 데이터에 얼마나 귀 기울일지를 조절한다. 아세틸콜린 수준이 높으면 감각 기관에서 오는 입력은 강해지고 내부의 “회상” 연결은 약해지며, 동시에 가중치가 바뀌기 쉬워진다~\cite{Hasselmo2006}. 거칠게 말해 “쓰기/읽기” 스위치이고, 그와 함께 학습률 $\alpha$이다.

\emph{\nt{SERO}, 세로토닌: 계획 지평.} 세로토닌이 무엇을 전달하는지는 가장 덜 알려져 있다. 한 가설은 세로토닌을 \eqref{eq:rpe}의 할인율 $\gamma$와 연결한다. 세로토닌 수준이 높으면 지금 작은 보상을 움켜쥐기보다 큰 보상을 기다릴 준비가 되어 있다는 것이다. 세로토닌 뉴런은 고르고 느리게 작동한다. 깨어 있을 때는 초당 스파이크 몇 개를 내고, 수면의 한 단계에서는 거의 잠잠해진다~\cite{JacobsAzmitia1992}.

여섯 가지 신경전달물질을 모두 표~\ref{tab:transmitters}에 모았다.

\begin{table}[t]
\centering
\small
\begin{tabular}{>{\raggedright}p{2.4cm}>{\raggedright}p{3.0cm}>{\raggedright}p{2.0cm}>{\raggedright\arraybackslash}p{5.2cm}}
\toprule
뉴런 유형 & 뉴런 비율 & 버스 & 계산에서의 역할 \\
\midrule
\nt{GLUT}(글루탐산) & 대다수 & 주소 & 양의 가중치, $w>0$ \\
\nt{GABA}(GABA) & 피질의 20--25\% & 주소 & 음의 가중치, $w<0$; 안정화 \\
\nt{DOPA}(도파민) & $\sim 6\cdot10^{-6}$ & 브로드캐스트 & 예측 오차 $\delta$ \\
\nt{NORA}(노르아드레날린) & $\sim 6\cdot10^{-7}$ & 브로드캐스트 & 이득 $g$(가설) \\
\nt{ACET}(아세틸콜린) & 적음 & 둘 다 & 학습률 $\alpha$; “쓰기/읽기”(가설) \\
\nt{SERO}(세로토닌) & $\sim 3\cdot10^{-6}$ & 브로드캐스트 & 할인 $\gamma$(가설) \\
\bottomrule
\end{tabular}
\caption{계산의 언어로 본 뇌의 주요 신경전달물질. 오른쪽 열은 크게 단순화한 것이다. 글루탐산, GABA, 도파민에 대해서는 믿을 만하고, 나머지는 가설이다.}
\label{tab:transmitters}
\end{table}

\section{부호는 수신자가 정한다}
\label{sec:receiver}

글루탐산은 양의 가중치, GABA는 음의 가중치라고 말했을 때 나는 조금 속였다. 어떤 분자도 부호를 품고 있지 않다. 분자는 열쇠일 뿐이다. 그것이 붙잡혔을 때 무슨 일이 일어날지는 \emph{자물쇠}, 즉 수신 세포의 수용체가 정한다.

가장 좋은 예는 도파민이다. 도파민에는 두 계열의 수용체가 있다~\cite{Beaulieu2011}. 한쪽에서는 세포의 흥분을 쉽게 하고, 다른 쪽에서는 어렵게 한다. 행동을 고르는 영역에서 어떤 뉴런은 첫 번째 유형의 수용체를, 다른 뉴런은 두 번째 유형을 지니고 있어서, 같은 도파민 신호가 한 무리는 강화하고 다른 무리는 약화한다. 네트워크의 노드마다 다르게 해석하는 브로드캐스트 패킷 하나와 같다.

\begin{keyblock}
\begin{proposition}[메시지의 의미는 수신자가 정한다]
\label{prop:receptor}
신경전달물질 작용의 부호와 속도는 물질 자체가 아니라 그것이 결합하는 수용체가 정한다. 수용체는 두 종류이다. \emph{이온성} 수용체는 그 자체가 이온 채널이며 1밀리초도 안 되어 열린다. 트랜지스터의 게이트처럼 전류를 직접 제어한다. \emph{대사성} 수용체는 세포 안에서 화학 반응의 연쇄를 시작하며 수백 밀리초 이상에 걸쳐 작용한다. 이것은 오히려 설정 레지스터에 값을 써서 세포의 이후 동작을 바꾸는 일에 가깝다. 주소 버스는 주로 전자를 통해, 브로드캐스트 버스는 주로 후자를 통해 말한다.
\end{proposition}
\end{keyblock}

그렇다면 왜 “글루탐산은 흥분시킨다”고 말할 수 있는가? 실제로 만나는 글루탐산 수용체는 거의 모두 흥분성이고, GABA 수용체는 거의 모두 억제성이기 때문이다. 이것은 자연법칙이 아니라 매우 믿을 만한 관례이며, 규칙~\eqref{eq:dalesign}은 바로 이 관례에 기대고 있다.

\section{가져갈 것}

압도적 다수의 뉴런은 뉴런마다 한 가지 부호의 가중치를 가진 주소 지정 데이터 버스이다. 극소수는 뇌의 넓은 영역에 몇 가지 공통 제어 신호를 전달하는 브로드캐스트 채널이다. 십만 개 중 두 개 정도의 뉴런이 나머지 네트워크 전체가 어떻게 학습하고 어떤 모드로 작동할지를 좌우한다. 엔지니어에게는 놀라운 일이 아니다. 어느 컴퓨터에서나 제어 레지스터는 메모리 셀보다 비교할 수 없이 적지만, 그것 없이는 기계가 돌아가지 않는다.

다음 장에서는 가장 수가 많은 유형인 \nt{GLUT} 뉴런만으로 이루어진 네트워크가 연속 시간에서 동작하면서 무엇을 계산할 수 있는지 살펴본다.

\section{연습문제}

\begin{exercise}[\lvlA]
\label{ex:01:1}
\emph{무더기와 전선.} 표~\ref{tab:types}에 따라(로그 눈금으로 범위의 가운데 값; \nt{ACET}은 출력 $10^5$개): (a)~뉴런 하나를 사람 한 명으로 치면 \nt{GLUT} 뉴런은 지구 인구의 몇 배이고, 브로드캐스트 뉴런 전체는 얼마만 한 도시인가? (b)~\nt{DOPA}, \nt{SERO}, \nt{NORA}, \nt{ACET}의 말단 비율은 뉴런 비율의 몇 배인가?
\end{exercise}

\begin{exercise}[\lvlA]
\label{ex:01:2}
\emph{0으로 되돌리기.} 틈 속 글루탐산은 $c_0e^{-t/\tau_c}$로 줄어든다. $c_0=1\mM$, $\tau_c=1\ms$이고, 수신자는 $10$~\textmu M보다 높은 농도를 감지한다. (a)~방출 한 번 뒤 선로는 얼마 동안 “통화 중”인가? (b)~펌프가 일부 차단되어 $\tau_c=10\ms$이다. 선로가 제때 비워지는 최대 방출 빈도는 얼마인가?
\end{exercise}

\begin{exercise}[\lvlB]
\label{ex:01:3}
\emph{버스트는 어디로 옮겨 가는가.} 그림~\ref{fig:rpe}의 실험에서 전등 뒤에 간격 $\Delta$로 상태 $s_0\to\dots\to s_{K-1}$이 이어지고 $T=K\Delta$이다. 보상 $r$은 마지막 전이에서 오고, 전등의 순간은 예측할 수 없다. 학습된 $V(s_k)$와 전등에 대한 응답 $\delta$를 구하라. $\gamma=e^{-\Delta/\tau_\gamma}$로 놓고 $T=1$~s, $\Delta=0.1$~s, $\gamma=0.95$일 때 $\tau_\gamma$를 구하라.
\end{exercise}

\begin{exercise}[\lvlB]
\label{ex:01:4}
\emph{모델링: 오차로 배우기.} 연습문제~\ref{ex:01:3}의 사슬을 $K=10$, $\gamma=0.95$, $r=1$과 규칙 $V(s_t)\leftarrow V(s_t)+\alpha\delta_t$로 시뮬레이션하라. $\alpha=0.05$, $0.1$, $0.2$, $0.5$일 때 전등에 대한 응답이 주스에 대한 응답을 처음 앞지르는 시행 $n^*$는 몇 번째인가? $n^*$는 $\alpha$에 어떻게 의존하는가?
\end{exercise}

\begin{exercise}[\lvlB]
\label{ex:01:5}
\emph{설계: 숫자 하나 브로드캐스트하기.} 숫자 $\delta$를 뉴런 $10^{10}$개 모두에 전달해야 한다. 중계기를 포함한 모든 수신자는 앞 층에서 말단 $10$개를 받아야 하고, 한 단계마다 $1\ms$가 더해진다. 뉴런당 출력이 $10^4$개(\nt{GLUT})일 때와 $10^{5.5}$개(\nt{DOPA})일 때, 가장 경제적인 중계 트리에는 뉴런 몇 개와 단계 몇 개가 필요한가?
\end{exercise}

\begin{exercise}[\lvlB]
\label{ex:01:6}
\emph{균형은 왜 민감한가.} 네트워크가 \nt{GLUT} $80\,\%$, \nt{GABA} $20\,\%$로 이루어지고, 모두가 모두와 가중치 $J_E/N$과 $-J_I/N$으로 연결되어 있다. $\tau\dot r=-r+Wr+h$. $J_E=10$일 때 $W$의 유일한 0이 아닌 고윳값은 $0.5$이다. 억제 전류와 흥분 전류의 비를 구하고, 억제를 몇 퍼센트만 약화해도 네트워크가 사태(avalanche)로 치닫는지 구하라.
\end{exercise}

\begin{exercise}[\lvlC]
\label{ex:01:7}
\emph{탐구: \nt{SERO}는 $\gamma$인가?} 연습문제~\ref{ex:01:3}에 따르면 전등에 대한 \nt{DOPA} 뉴런의 응답은 $re^{-T/\tau_\gamma}$이다. 지연 $T=1,\dots,5$~s를 각각 $n$번씩 제시하고, $\ln\delta$는 산포 $0.5$로 측정된다. $\tau_\gamma=2$~s와 $2.4$~s를 구별하려면(유의수준 $0.05$, 검정력 $0.8$) $n$이 얼마나 필요한가? $\gamma$ 말고 어떤 메커니즘이 $T$에 따라 같은 감소를 낳을 수 있는가?
\end{exercise}

\chapter{\nt{GLUT} 뉴런은 무엇을 계산하는가}
\label{sec:glutcomputer}

\section{게임의 규칙}

\ref{sec:neurontypes}장에서 보았듯이 뇌 뉴런의 압도적 다수는 \nt{GLUT}이다. 이들은 흥분만 시키며, 가중치는 양수뿐이다. 이런 뉴런을 원하는 만큼 가져와서 묻자. 살아 있는 몸에서 일어나는 일만 허용한다면, 이런 네트워크는 무엇을 계산할 수 있는가?

살아 있는 몸에는 모든 소자를 한꺼번에 전환하는 클록 발생기가 없다. 뉴런은 저마다 연속 시간에서 독립적으로 작동한다. 스파이크는 제각각의 지연을 두고 아무 때나 오고 간다. 빛, 소리, 압력에 반응하는 감각 뉴런이 \emph{입력}이고, 근육을 제어하는 뉴런이 \emph{출력}이다. 그 사이에 네트워크가 있으며, 누구도 네트워크에 언제 계산을 시작하고 언제 끝낼지 알려 주지 않는다. 전자공학자에게 이것은 소자 지연을 미리 정확히 알 수 없는 \emph{비동기} 회로이다.

뉴런 모델로는 연속 시간에서 제대로 작동하는 모델 가운데 가장 단순한 것을 쓰자. 뉴런 막의 전압 $V$는 휴지 상태에서 0이다. 뉴런 $j$에서 스파이크가 하나 올 때마다 전압이 $w_j\ge0$만큼 뛰어오르고, 그 뒤 시간 상수 $\tau$로 저절로 0까지 흘러내린다.
\begin{empheq}[box=\keybox]{equation}
\tau\,\frac{\dd V}{\dd t} \;=\; -V \;+\; \tau\sum_j w_j \sum_k \delta\bigl(t-t_j^{(k)}-d_j\bigr), \qquad w_j\ge 0 .
\label{eq:glutneuron}
\end{empheq}
여기서 $t_j^{(k)}$는 뉴런 $j$의 스파이크 시각, $d_j$는 그 뉴런에서 오는 경로의 지연, $\delta$는 순간적인 도약을 뜻하는 델타 함수이다. $V$가 문턱값 $\theta$에 이르면 뉴런은 스파이크를 내고, $V$는 0으로 초기화되며, 약 1밀리초 동안 뉴런은 다시 발화하지 못한다. 전형적인 수치는 이렇다. $\tau$는 뉴런에 따라 1밀리초 미만에서 20밀리초까지, 지연은 1밀리초 미만에서 수십 밀리초까지이다. 이것이 \emph{누설 적분 발화 뉴런}, 곧 문턱값이 있는 RC 회로이다. 의도적인 단순화이다. 피질 뉴런에서는 입력 가지가 스스로 국소 스파이크를 내며~\cite{Smith2013}, 뉴런 하나는 오히려 작은 다층 네트워크처럼 행동한다(\ref{sec:synthesis}장). 이 장의 질문에는 RC 회로 하나로 충분하다.

논리 게이트와의 가장 큰 차이는 누설이다. 뉴런은 스파이크를 영원히 기억하지 않고 대략 $\tau$ 동안만 기억하며, 이 창 안에 들어온 것만 더해진다.

\section{OR와 AND}

게이트부터 시작하자. 스파이크 하나가 전압을 문턱값 위로 올리면, 즉 $w\ge\theta$이면 뉴런은 어느 입력에서 오든 모든 스파이크에 발화한다. 이것이 \emph{OR}이다.

AND는 더 흥미롭다. $w<\theta<2w$라고 하자. 스파이크 하나로는 모자라고 둘이 필요하다. 그런데 이제 `둘 다 왔는가'라는 질문은 \emph{어느 시간 안에} 왔는지를 말하기 전에는 의미가 없다. 첫 스파이크가 시각 $0$에, 둘째가 시간 $\Delta$ 뒤에 왔다고 하자. 둘째가 올 때 첫째에서 남은 것은 $we^{-\Delta/\tau}$이고, $we^{-\Delta/\tau}+w\ge\theta$이면 뉴런이 발화한다. 여기서 일치 창이 나온다.
\begin{empheq}[box=\keybox]{equation}
\Delta \;\le\; \Delta_{\max} \;=\; \tau\,\ln\frac{w}{\theta-w} .
\label{eq:window}
\end{empheq}
이것은 논리 AND가 아니라 \emph{동시성 검출기}, 곧 시간 속의 AND이다(그림~\ref{fig:coincidence}). $\theta=1.5w$이면 창은 $\tau\ln2\approx0.7\tau$이다. $\tau=10\ms$인 피질 뉴런에서는 약 7밀리초이다. $\tau$가 1밀리초보다 짧은 청각계 뉴런에서는 창이 1밀리초 미만으로 줄어든다.

\begin{figure}[t]
\centering
\begin{tikzpicture}[>=Stealth,x=0.9cm,y=1.6cm]
\foreach \dx/\lab/\c [count=\i] in {0.5/{일치함}/red!70!black, 2.6/{일치하지 않음}/blue!60!black}{
 \begin{scope}[xshift={(\i-1)*7.2cm}]
  \draw[->,gray!70!black] (0,0) -- (6.3,0) node[below left,black] {\small $t$};
  \draw[->,gray!70!black] (0,0) -- (0,1.75) node[left,black] {\small $V$};
  \draw[densely dashed,gray!60!black] (0,1.5) -- (6.0,1.5) node[right,black] {\small $\theta$};
  \draw[very thick,\c] (0,0) -- (1,0) -- (1,1)
      plot[domain=1:1+\dx,samples=30] (\x,{exp(-(\x-1)/1.5)});
  \pgfmathsetmacro{\v}{exp(-\dx/1.5)+1}
  \draw[very thick,\c] (1+\dx,{exp(-\dx/1.5)}) -- (1+\dx,{min(\v,1.5)});
  \ifdim\v pt<1.5pt
    \draw[very thick,\c] plot[domain=1+\dx:6,samples=40] (\x,{\v*exp(-(\x-1-\dx)/1.5)});
  \else
    \draw[very thick,\c] (1+\dx,1.5) -- (1+\dx,0) -- (6,0);
    \draw[->,thick,\c] (1+\dx,1.55) -- (1+\dx,1.72) node[above,font=\scriptsize] {스파이크};
  \fi
  \draw[->,thick] (1,-0.35) -- (1,-0.08); \draw[->,thick] (1+\dx,-0.35) -- (1+\dx,-0.08);
  \draw[decorate,decoration={brace,amplitude=3pt,mirror}] (1,-0.42) -- (1+\dx,-0.42) node[midway,below=3pt] {\scriptsize $\Delta$};
  \node[\c] at (3.5,-0.95) {\small \lab};
 \end{scope}}
\end{tikzpicture}
\caption{동시성 검출기, 문턱값 $\theta=1.5w$. 왼쪽에서는 둘째 스파이크가 $\Delta<\Delta_{\max}$ 뒤에 와서 합이 문턱값에 이르고 뉴런이 발화했다. 오른쪽은 $\Delta>\Delta_{\max}$이다. 첫 스파이크에서 거의 아무것도 남지 않아 뉴런은 침묵했다.}
\label{fig:coincidence}
\end{figure}

AND를 얻는 다른 방법은 정확한 시각이 아니라 지속 시간을 쓰는 것이다. 빛이 켜져 있는 동안 감각 뉴런은 스파이크를 자주, 고르게 낸다. 그런 입력 하나로는 문턱값에 모자라고 둘이면 충분하다면, 뉴런은 두 입력이 모두 활성인 동안 작동한다. 이렇게 네트워크는 \emph{수준}으로 계산하며, 스파이크가 도착하는 정확한 시각은 중요하지 않다. 아래 내용의 대부분이 바로 이런 방식이다.

\section{할 수 없는 것}

이제 NOT을 만들어 보자. 입력이 조용할 때 작동하는 뉴런이다. 만들 수 없다. 입력 스파이크가 없으면 \eqref{eq:glutneuron}에는 도약이 하나도 없고, 전압은 문턱값 아래인 0에 머문다. 뉴런이 많은 네트워크라면? 역시 안 되며, 이것은 모든 네트워크에 대해 한꺼번에 증명된다.

발화율로 말하는 것이 가장 편하다. $r_i$를 뉴런 $i$의 발화율, $I_i$를 입력에서 오는 유입, $f_i$를 그 뉴런의 전달 특성, 즉 발화율이 총 유입에 어떻게 의존하는지라고 하자. 적분 뉴런의 전달 특성은 감소하지 않는다. 유입이 많을수록 스파이크가 잦다. 평형에 이른 네트워크의 발화율은 다음 방정식을 만족한다.
\begin{equation}
r_i \;=\; f_i\Bigl(\sum_j w_{ij}\,r_j + I_i\Bigr), \qquad w_{ij}\ge 0 .
\label{eq:ratenet}
\end{equation}

\begin{keyblock}
\begin{proposition}[단조성]
\label{prop:monotone}
네트워크~\eqref{eq:ratenet}가 \nt{GLUT} 뉴런만으로 이루어지고 완전한 침묵에서 출발한다고 하자. 입력을 강하게 하면 어떤 뉴런의 정상 상태 발화율도 줄어들지 않는다. 이런 네트워크가 계산하는 것은 모두 입력의 \emph{단조} 함수이다.
\end{proposition}
\end{keyblock}

증명. 침묵 $r^{(0)}=0$에서 출발해 발화율을 \eqref{eq:ratenet}의 우변에 대입해 나가자. $r^{(k+1)}=F\bigl(r^{(k)}\bigr)$. 가중치가 음이 아니고 $f_i$가 감소하지 않으므로 우변 $F$는 어느 인수에 대해서도 감소하지 않는다. $r^{(1)}\ge0=r^{(0)}$이므로 귀납적으로 $r^{(k+1)}\ge r^{(k)}$이다. 발화율은 커지기만 하며 \eqref{eq:ratenet}의 가장 작은 해에 이른다. 입력 $I$를 더 큰 $I'$로 바꾸면 매 단계에서 $r^{(k)}(I)\le r^{(k)}(I')$이고, 극한에서도 그렇다. 실제 네트워크는 단계가 아니라 연속적으로 평형에 이르지만 같은 결론이 성립한다. 연결이 양수이면 두 실행 사이의 부등식이 처음에 성립할 경우 내내 유지된다.

개별 스파이크에 대해서는 이 명제가 글자 그대로 성립하지 않는다. 여분의 입력 스파이크 하나가 뉴런을 조금 일찍 발화시키면, 1밀리초의 불응기 때문에 그다음 스파이크가 사라질 수 있다. 그러나 이 효과는 약하고 믿을 수 없어서 계산에 쓸 수 없다. 발화율에 대해서는 단조성이 엄밀하다.

결과는 인버터가 없는 여느 회로와 같다. NOT이 없다. 배타적 OR(XOR)도 없다. 두 번째 입력을 더해 출력을 끌 수 없기 때문이다. 그리고 가장 곤란한 점은 \emph{활동을 활동으로 멈출 수 없다}는 것이다. 활동을 떠받치는 것을 치울 수 있을 뿐이다.

\section{전선 두 가닥: 켜짐과 꺼짐}

막다른 길에서 나가는 방법은 몸 스스로가 알려 준다. 많은 감각 기관에서 자극은 두 벌의 뉴런으로 전달된다. 시각에서는 어떤 세포는 빛이 강해질 때, 다른 세포는 약해질 때 응답한다~\cite{Schiller1992}. 이들을 \emph{켜짐} 뉴런과 \emph{꺼짐} 뉴런이라 부르자. 네트워크는 전선 $x$ 하나 대신 쌍 $(x,\bar x)$를 받으며, 네트워크가 계산할 수 없는 `없음'은 센서가 직접 알려 준다.

전선 쌍이 있으면 NOT은 공짜이다. 전선 두 가닥을 맞바꾸기만 하면 된다. AND와 OR는 각각 뉴런 둘로 만든다.
\begin{empheq}[box=\keybox]{equation}
\begin{aligned}
x\wedge y \;&\to\; \bigl(\;x\wedge y,\;\; \bar x\vee\bar y\;\bigr),\\
x\vee y \;&\to\; \bigl(\;x\vee y,\;\; \bar x\wedge\bar y\;\bigr).
\end{aligned}
\label{eq:dualrail}
\end{empheq}
오른쪽의 연산은 모두 단조이므로 명제~\ref{prop:monotone}에 어긋나지 않는다.

비동기 회로에서 이중 레일 부호에는 더 중요한 두 번째 장점이 있다. 전선 쌍의 상태는 셋이다. `예', `아니오', 그리고 둘 다 조용할 때의 \emph{`아직 답 없음'}. 수신자는 계산이 끝났다는 것을 타이머가 아니라 신호 자체로 안다. 모든 쌍에서 정확히 한 가닥이 작동하면 끝난 것이다. 자극과 자극 사이에는 센서가 조용하고, 네트워크는 침묵으로 돌아가 다음 차례를 준비한다. 비동기 전자공학에서는 이 기법으로 소자의 지연이 어떻든 올바르게 동작하는 회로를 만든다~\cite{Sparso2001}.

\begin{keyblock}
\begin{proposition}[비동기 계산]
\label{prop:dualrail}
각 입력이 `켜짐--꺼짐' 전선 쌍으로 들어오면, \nt{GLUT} 뉴런 네트워크는 입력의 어떤 논리 함수든 계산할 수 있다. 답도 전선 쌍으로 나오며, 답이 준비되었다는 것은 모든 쌍에서 전선 한 가닥이 작동했다는 것으로 보인다. 이를 위해 클록은 필요 없다. 대가는 뉴런이 약 두 배로 드는 것이다.
\end{proposition}
\end{keyblock}

예로 배타적 OR를 보자. `두 자극 가운데 정확히 하나가 변했다'이다. 답의 정방향 전선은 $(x\wedge\bar y)\vee(\bar x\wedge y)$, 역방향 전선은 $(x\wedge y)\vee(\bar x\wedge\bar y)$이다. 뉴런 여섯 개이며, 그중 어느 것도 억제를 필요로 하지 않는다(그림~\ref{fig:dualxor}).

\begin{figure}[t]
\centering
\begin{tikzpicture}[>=Stealth,x=1cm,y=1cm,
  nrn/.style={circle,draw,very thick,fill=blue!6,minimum size=0.75cm,inner sep=0pt,font=\small},
  inp/.style={font=\small}]
\node[inp] (X)  at (0,3.3) {$x$};
\node[inp] (Xb) at (0,2.2) {$\bar x$};
\node[inp] (Y)  at (0,1.1) {$y$};
\node[inp] (Yb) at (0,0)   {$\bar y$};
\node[nrn] (A1) at (3.2,3.3) {$2$};
\node[nrn] (A2) at (3.2,2.2) {$2$};
\node[nrn] (A3) at (3.2,1.1) {$2$};
\node[nrn] (A4) at (3.2,0)   {$2$};
\node[nrn] (O)  at (7.2,2.75) {$1$};
\node[nrn] (Ob) at (7.2,0.55) {$1$};
\foreach \s/\t in {X/A1,Yb/A1,Xb/A2,Y/A2,X/A3,Y/A3,Xb/A4,Yb/A4}{\draw[->,thick,blue!60!black] (\s) -- (\t);}
\draw[->,thick,blue!60!black] (A1) -- node[pos=0.45,above,sloped,font=\scriptsize,black] {$x\wedge\bar y$} (O);
\draw[->,thick,blue!60!black] (A2) -- node[pos=0.45,below,sloped,font=\scriptsize,black] {$\bar x\wedge y$} (O);
\draw[->,thick,blue!60!black] (A3) -- node[pos=0.45,above,sloped,font=\scriptsize,black] {$x\wedge y$} (Ob);
\draw[->,thick,blue!60!black] (A4) -- node[pos=0.45,below,sloped,font=\scriptsize,black] {$\bar x\wedge\bar y$} (Ob);
\draw[->,very thick,red!70!black] (O) -- (8.6,2.75) node[right,black] {$x\oplus y$: `예'};
\draw[->,very thick,red!70!black] (Ob) -- (8.6,0.55) node[right,black] {$\overline{x\oplus y}$: `아니오'};
\node[below,font=\small,gray!40!black] at (3.2,-0.55) {AND: 문턱값 $2$};
\node[below,font=\small,gray!40!black] at (7.2,-0.55) {OR: 문턱값 $1$};
\node[below,font=\small,gray!40!black] at (0,-0.55) {전선 쌍};
\end{tikzpicture}
\caption{\nt{GLUT} 뉴런만으로 이중 레일 부호 위에 만든 배타적 OR. 원 안의 숫자는 입력 하나의 가중치를 단위로 한 문턱값이다. AND 뉴런 넷이 입력의 네 조합을 알아보고, OR 뉴런 둘이 이들을 모아 답의 전선 쌍을 만든다.}
\label{fig:dualxor}
\end{figure}

\section{시계 대신 지연}

시간을 다루는 계산에는 전선의 지연과 동시성 검출기만으로 충분하다. 가장 좋은 예는 뇌가 소리가 어디서 왔는지 알아내는 방식이다.

옆쪽에 있는 음원의 소리는 먼 귀보다 가까운 귀에 먼저 닿는다. 그 차이는 방향에 따라 달라서, 음원이 바로 앞에 있으면 0이고, 사람에서 음원이 정확히 옆에 있으면 $0.22\,\text{m}/343\,\text{m/s}\approx0.6\ms$이다. 뇌는 약 10\emph{마이크로초}의 차이를 구별한다~\cite{KlumppEady1956,Grothe2010}. 그런데 뉴런 자체는 1밀리초의 몇 분의 1보다 나은 정밀도로 발화하지 못한다. 어떻게 가능한가?

왼쪽 귀에서 오는 전선을 동시성 검출기 열을 따라 왼쪽에서 오른쪽으로, 오른쪽 귀에서 오는 전선을 오른쪽에서 왼쪽으로 깔자(그림~\ref{fig:itd}). 신호는 전선을 따라 속도 $v$로 달린다. 길이 $L$인 열의 왼쪽 끝에서 거리 $x$에 있는 검출기까지 왼쪽 신호는 시간 $x/v$, 오른쪽 신호는 $(L-x)/v$ 걸린다. 오른쪽 귀에 소리가 $\delta t$만큼 늦게 왔다면 두 신호는 다음 지점에서 동시에 만난다.
\begin{empheq}[box=\keybox]{equation}
\frac{x}{v} \;=\; \frac{L-x}{v} + \delta t
\quad\Longrightarrow\quad
x \;=\; \frac{L}{2} + \frac{v\,\delta t}{2} .
\label{eq:itd}
\end{empheq}
두 신호가 창~\eqref{eq:window} 안에 도착한 검출기만 발화한다. 발화한 검출기의 번호가 곧 음원의 방향이다. 시간 차이가 \emph{위치}로 바뀌었다.

\begin{figure}[t]
\centering
\begin{tikzpicture}[>=Stealth,
  nrn/.style={circle,draw,very thick,fill=blue!6,minimum size=0.7cm,inner sep=0pt}]
\foreach \k in {0,...,6}{\node[nrn] (D\k) at (1.4*\k,0) {\scriptsize $2$};}
\node[nrn,fill=red!15] at (D4) {\scriptsize $2$};
\draw[->,thick,blue!60!black] (-1.4,0.9) node[left,black] {\small 왼쪽 귀} -- (9.2,0.9);
\draw[->,thick,orange!85!black] (9.8,-0.9) node[right,black] {\small 오른쪽 귀} -- (-0.8,-0.9);
\foreach \k in {0,...,6}{
  \draw[->,blue!60!black] (1.4*\k,0.9) -- (D\k.north);
  \draw[->,orange!85!black] (1.4*\k,-0.9) -- (D\k.south);}
\draw[->,thick,red!70!black] (D4.north east) ++(0.1,0.05) -- ++(0.5,0.45) node[above right,font=\small,black] {발화};
\node[font=\small,gray!40!black] at (4.2,-1.5) {지연 증가 $\longrightarrow$ \hspace{2.2cm} $\longleftarrow$ 지연 증가};
\pic[scale=1.4] at (-2.3,1.75) {speaker};
\node[font=\scriptsize,align=center] at (-2.3,2.35) {왼쪽 음원};
\end{tikzpicture}
\caption{소리 방향 알아내기. 두 귀에서 온 신호가 서로를 향해 달린다. 원 하나하나는 문턱값이 $2$인 동시성 검출기이다. 오른쪽 귀에 소리가 늦게 왔으면 신호는 식~\eqref{eq:itd}에 따라 가운데보다 오른쪽에서 만난다. 네트워크의 출력은 발화한 뉴런의 번호이다. 여기서는 음원이 왼쪽에 있어 왼쪽 귀에 소리가 먼저 왔다.}
\label{fig:itd}
\end{figure}

여기에는 클록도, 억제도, 이중 레일 부호조차 없다. 네트워크는 `예'나 `아니오'로 답하지 않고 많은 뉴런 중 하나를 가리킨다. 지연선과 동시성 검출기로 된 이런 회로는 새의 청각 뇌줄기에서 실제로 작동하는 것으로 보인다~\cite{Carr1990}. 마이크로초 정밀도는 뉴런 하나하나의 정밀도가 아니라, 검출기가 많고 각자 자기 지연을 본다는 데서 나온다. 포유류에서는 이런 지연의 열이 발견되지 않았다. 포유류에서는 같은 문제를 \nt{GLYC} 뉴런이 시간을 정확히 맞춰 가하는 억제로 조율된 동시성 검출기가 푸는 것으로 보인다~\cite{Brand2002,Grothe2010}. 억제가 일치 창을 어떻게 좁히는지는 \ref{sec:gaba}장에서 \nt{GABA}를 예로 살펴본다. 글리신도 수신자의 염소 채널을 열어 같은 방식으로 억제한다.

\section{메모리}

컴퓨터에는 메모리가 필요하다. 이런 네트워크에서 메모리는 스스로를 유지하는 활동이다. 서로 강하게 연결된 \nt{GLUT} 뉴런 무리를 하나 잡아 발화율 하나 $r$로 기술하자(그림~\ref{fig:recurrentgroup}a). 평형에서 $r=f(wr+I)$이고, 여기서 $w$는 무리가 자기 자신에게 주는 피드백의 세기, $I$는 바깥에서 들어오는 유입이다. 이 방정식을 곡선 $f(wr+I)$와 직선 $r$의 교점으로 그래프를 그려 풀자(그림~\ref{fig:bistable}).

\begin{figure}[t]
\centering
% (b): tau dr/dt = -r + f(w r + I), f as in fig:bistable, w=1, I=0.6; Euler dt=0.001 tau.
\begin{tikzpicture}[>=Stealth,
  nrn/.style={circle,draw,thick,fill=blue!6,minimum size=0.5cm,inner sep=0pt}]
\begin{scope}
\node[font=\small] at (-0.5,1.55) {(a)};
\foreach \k in {1,...,6}{\node[nrn] (n\k) at ({1.4+1.0*cos(60*\k)},{1.0*sin(60*\k)}) {};}
\foreach \a/\b in {1/2,2/3,3/4,4/5,5/6,6/1,1/4,2/5,3/6}{\draw[<->,blue!60!black] (n\a) -- (n\b);}
\foreach \k in {2,3,4}{\draw[->,thick,green!45!black] ($(n\k)+(-0.95,0)$) -- (n\k);}
\node[font=\small,green!45!black] at (-0.35,-0.45) {$I$};
\node[font=\Large] at (3.1,0) {$\approx$};
\node[nrn,minimum size=0.85cm,font=\small] (R) at (4.35,-0.1) {$r$};
\draw[->,thick,blue!60!black] (R) edge[loop above,looseness=6] node[above,font=\small] {$w$} (R);
\draw[->,thick,green!45!black] (3.55,-0.95) node[left,font=\small] {$I$} -- (R);
\end{scope}
\begin{scope}[xshift=6.3cm,y=2.6cm,x=1.05cm]
\node[font=\small] at (-0.6,1.25) {(b)};
\draw[->,gray!70!black] (0,0) -- (7.3,0) node[below left,font=\small,black] {$t/\tau$};
\draw[->,gray!70!black] (0,0) -- (0,1.12) node[left,font=\small,black] {$r$};
\foreach \t in {0,2,4,6}{\draw[gray!60!black] (\t,-0.02) -- (\t,0.02); \node[below,font=\scriptsize] at (\t,0) {\t};}
\draw[densely dotted,gray!60!black] (0,0.5) -- (7,0.5) node[right,font=\scriptsize,black,align=left] {쓰기\\문턱값};
\fill[green!45!black,opacity=0.25] (1,-0.2) rectangle (2,-0.13);
\fill[gray!60!black,opacity=0.35] (1,-0.29) rectangle (1.5,-0.22);
\draw[very thick,blue!60!black] plot coordinates {(0.0,0.002) (0.1,0.002) (0.2,0.002) (0.3,0.002) (0.4,0.002) (0.5,0.002) (0.6,0.002) (0.7,0.002) (0.8,0.002) (0.9,0.002) (1.0,0.002) (1.1,0.083) (1.2,0.165) (1.3,0.242) (1.4,0.313) (1.5,0.378) (1.6,0.437) (1.7,0.491) (1.8,0.539) (1.9,0.583) (2.0,0.623) (2.1,0.643) (2.2,0.665) (2.3,0.688) (2.4,0.71) (2.5,0.732) (2.6,0.753) (2.7,0.773) (2.8,0.792) (2.9,0.81) (3.0,0.826) (3.2,0.855) (3.4,0.88) (3.6,0.9) (3.8,0.917) (4.0,0.931) (4.4,0.953) (4.8,0.967) (5.2,0.977) (5.6,0.984) (6.0,0.988) (6.5,0.992) (7.0,0.994)};
\draw[very thick,gray!60!black,densely dashed] plot coordinates {(0.0,0.002) (0.1,0.002) (0.2,0.002) (0.3,0.002) (0.4,0.002) (0.5,0.002) (0.6,0.002) (0.7,0.002) (0.8,0.002) (0.9,0.002) (1.0,0.002) (1.1,0.083) (1.2,0.165) (1.3,0.242) (1.4,0.313) (1.5,0.378) (1.6,0.358) (1.7,0.336) (1.8,0.313) (1.9,0.291) (2.0,0.269) (2.2,0.228) (2.4,0.191) (2.6,0.159) (2.8,0.133) (3.0,0.11) (3.2,0.091) (3.4,0.076) (3.6,0.063) (3.8,0.052) (4.0,0.043) (4.4,0.03) (4.8,0.021) (5.2,0.015) (5.6,0.011) (6.0,0.008) (6.5,0.006) (7.0,0.004)};
\node[font=\scriptsize,blue!60!black,anchor=west] at (3.3,0.8) {유입 $1\tau$: 기록됨};
\node[font=\scriptsize,gray!50!black,anchor=west] at (2.3,0.3) {유입 $0.5\tau$: 꺼짐};
\end{scope}
\end{tikzpicture}
\caption{자기 자신과 강하게 연결된 무리. (a)~서로를 흥분시키는 많은 \nt{GLUT} 뉴런은 세기 $w$로 자기 자신을 흥분시키는 발화율 $r$의 뉴런 하나처럼 행동한다. 바깥에서 유입 $I$가 들어온다. (b)~$w=1$이고 곡선 $f$가 그림~\ref{fig:bistable}과 같을 때 무리의 발화율을 시간에 따라 그렸다. 유입 $I=0.6$은 축 아래 막대로 표시한 시간 동안 가해진다. 유입이 $1\tau$ 동안 유지되면 무리는 불안정한 점 $r=0.5$를 넘어서 불이 붙고, 유입이 끝난 뒤에도 계속 켜져 있다. $0.5\tau$뿐이면 그 점에 이르지 못하고 다시 꺼진다. 비트를 쓰려면 유입이 약 $0.7\tau$보다 오래 지속되어야 한다.}
\label{fig:recurrentgroup}
\end{figure}

\begin{figure}[t]
\centering
\begin{tikzpicture}[>=Stealth,x=5cm,y=3.2cm]
\draw[->,gray!70!black] (0,0) -- (1.1,0) node[below left,black] {\small $r$};
\draw[->,gray!70!black] (0,0) -- (0,1.1);
\draw[thick,gray!60!black] (0,0) -- (1.05,1.05) node[right,black] {\small $r$};
\draw[very thick,blue!60!black] plot[domain=0:1.05,samples=80] (\x,{1/(1+exp(-(\x-0.5)/0.08))});
\draw[very thick,red!70!black,densely dashed] plot[domain=0:1.05,samples=80] (\x,{1/(1+exp(-(\x+0.3-0.5)/0.08))});
\fill[blue!60!black] (0.002,0.002) circle (0.018);
\draw[blue!60!black,fill=white,thick] (0.5,0.5) circle (0.018);
\fill[blue!60!black] (0.998,0.998) circle (0.018);
\node[below right,font=\small] at (0.02,0.0) {꺼짐};
\node[right,font=\small] at (0.52,0.46) {불안정};
\node[below,font=\small] at (0.99,0.96) {켜짐};
\node[anchor=east,font=\small,red!70!black] at (0.19,0.62) {$f(wr+I)$};
\node[anchor=west,font=\small,blue!60!black] at (0.45,0.1) {$f(wr)$};
\end{tikzpicture}
\caption{피드백이 강한 \nt{GLUT} 뉴런 무리로 만든 메모리. 실선 곡선은 바깥 유입이 없을 때이다. 직선 $r$과 안정한 교점 둘, 즉 `꺼짐'과 `켜짐'이 있고, 그 사이에 불안정한 교점이 있다. 짧은 유입 $I$(파선 곡선)가 가해지면 위쪽 교점만 남는다.}
\label{fig:bistable}
\end{figure}

피드백이 강하면 교점은 셋이다. 아래쪽은 침묵, 위쪽은 무리가 한껏 켜진 상태, 가운데는 불안정하다. 두 안정 상태를 가진 소자, 곧 플립플롭이 생겼다. 여기에 1을 쓰기는 쉽다. 짧은 유입이 곡선을 들어 올리면 아래쪽 교점이 사라지고, 무리에 불이 붙는다. 그리고 유입이 끝난 뒤에도 계속 켜져 있다(그림~\ref{fig:recurrentgroup}b). 너무 짧은 유입은 무리를 불안정한 점까지 끌어올리지 못하고, 무리는 다시 꺼진다. 자극이 사라진 뒤 몇 초 동안 유지되는 이런 자기 유지 활동은 실제로 뇌에서 관찰되며 단기 기억의 바탕으로 여겨진다~\cite{Wang2001}. 그러나 몇 초를 기억하는 방법이 이것만은 아니다. 기록은 시냅스 자체의 느린 변수에도 저장할 수 있다. 스파이크 버스트 뒤의 촉진은 \ref{sec:code}장의 `대시' 시냅스처럼 약 1초 동안 유지되고, 네트워크는 짧은 폭발 사이사이에 거의 침묵할 수 있다. 플립플롭이 아니라 충전된 커패시터이다~\cite{Mongillo2008}. 기억을 유지하는 동안의 이런 `조용한' 구간은 관찰되며~\cite{Stokes2015}, 끊임없는 스파이크 없이 저장하는 편이 에너지가 덜 든다. 피질이 어떤 경우에 두 방법 중 어느 것을 쓰는지는 논의 중이다.

그러면 어떻게 지우는가? 명제~\ref{prop:monotone}에 따르면 활동으로는 지울 수 없다. 무언가를 \emph{치우는} 수밖에 없다. 첫째 방법은 지지를 치우는 것이다. 무리가 다른 무리에서 오는 유입과 함께일 때만 켜져 있다면, 기억은 그 유입과 함께 꺼진다. 둘째는 피로이다. 실제 시냅스는 오래 작동하면 신경전달물질 저장량이 고갈되고~\cite{Tsodyks1997}, 뉴런에서는 흥분성을 떨어뜨리는 느린 전류가 자라난다. 피드백이 약해지고 위쪽 교점이 사라지면, 무리는 몇 초 뒤 저절로 꺼진다. 신호로 켜지고 시간으로만 꺼지는 메모리이다.

\section{순서열}

클록 대신 \emph{사건이 시동하는 타이머}를 둘 수 있다. \nt{GLUT} 뉴런 무리를 사슬로 잇자. 첫째가 둘째를, 둘째가 셋째를 흥분시키는 식이다. 바깥 신호가 첫 무리에 불을 붙이면 활동의 파동이 사슬을 따라 달린다. 각 고리는 앞 고리보다 지연 한두 개만큼 뒤에, 그리고 딱 한 번만 켜진다. 스파이크 직후 뉴런은 불응 상태이고, 파동은 이미 앞으로 나아갔기 때문이다.

이런 사슬은 사건으로부터의 시간을 잰다. $k$번째 고리가 발화했다면 시동 이후 지연 $k$개만큼의 시간이 흘렀다는 뜻이다. 그 출력을 근육에 연결하면 스스로 펼쳐지는 동작 순서열을 얻는다. 명금류가 노래할 때 뇌의 한 영역에 있는 \nt{GLUT} 뉴런들은 하나씩 엄격히 차례대로, 저마다 노래의 자기 순간에 발화한다. 이런 사슬과 매우 비슷하다~\cite{Hahnloser2002}.

클록과 달리 사슬은 시동되기 전까지 침묵하고, 한 번만 작동하며, 자신에게 연결된 것들에 대해서만 시간을 잰다. 네트워크에는 여전히 공통 박자가 없다.

\section{부족한 것}

억제 없이 \nt{GLUT} 뉴런만으로도 많은 것을 만들 수 있다. 그러나 세 가지가 부족하며, 세 가지 모두 명제~\ref{prop:monotone} 때문이다.

\emph{선택.} 네트워크는 방향 하나, 행동 하나를 골라야 한다. 두 자극이 거의 같으면 그림~\ref{fig:itd}의 네트워크에서는 두 출력이 모두 발화하고, 강한 쪽을 위해 약한 쪽을 누를 수단이 없다.

\emph{초기화.} 기억을 신호로 지울 수 없다.

\emph{안정성.} 엔지니어의 설계가 아닌 방식으로 연결이 생기는 네트워크에서는 양의 피드백 루프를 피할 수 없고, 그 안에서 활동이 사태처럼 불어날 수 있다(\ref{sec:neurontypes}장). 유일한 브레이크는 역시 피로인데, 느리고 거칠다.

세 가지 문제 모두 약은 하나이다. 스파이크를 스파이크로 끄는 뉴런이다. 그것이 \nt{GABA}이며, 이 뉴런은 계산하기 위해서가 아니라 고르고, 초기화하고, 계산을 통제하기 위해 필요하다.

\section{연습문제}

\begin{exercise}[\lvlA]
\label{ex:02:1}
\emph{소리를 위한 검출기 열.} 귀에서 온 신호가 그림~\ref{fig:itd}의 전선을 따라 속도 $5$~m/s로 달리고, 가장 큰 도착 시간 차는 $0.64\ms$이다. 검출기는 $\tau=0.5\ms$, $\theta=1.5w$인 뉴런~\eqref{eq:glutneuron}이며, 이웃한 검출기가 $10$~\textmu s를 구별하도록 배치된다. 열에는 검출기가 몇 개 있고, 소리 하나에 몇 개가 발화하는가?
\end{exercise}

\begin{exercise}[\lvlB]
\label{ex:02:2}
\emph{입력이 같지 않으면 창이 밀린다.} $\tau=0.5\ms$, $\theta=1.5$인 동시성 검출기~\eqref{eq:glutneuron}가 가중치 $1$과 $0.8$인 두 입력에서 스파이크를 하나씩 받는다. 일치 창과 그 중심을 구하라. 한쪽 귀의 입력이 다른 쪽보다 $20\,\%$ 약하면 그림~\ref{fig:itd}의 열은 몇 마이크로초 틀리는가?
\end{exercise}

\begin{exercise}[\lvlB]
\label{ex:02:3}
\emph{진동하지 않는 네트워크.} $f_i$가 감소하지 않고 $i\ne j$에서 $w_{ij}\ge0$인 네트워크 $\tau\dot r_i=-r_i+f_i\bigl(\textstyle\sum_jw_{ij}r_j+I_i\bigr)$는 단조이며 지속적으로 진동하지 않는다. 모든 순환에 억제 연결이 짝수 개 있는 네트워크는 치환 $r_i\to\pm r_i$로 이것에 귀착됨을 증명하라. 두 무리의 상호 억제는 진동할 수 있는가? \nt{GLUT}--\nt{GABA} 루프는?
\end{exercise}

\begin{exercise}[\lvlB]
\label{ex:02:4}
\emph{설계: 이중 레일 가산기.} 각 비트는 전선 쌍 $(x,\bar x)$로 전달된다. 합의 쌍 $(S,\bar S)$와 자리올림의 쌍 $(C,\bar C)$를 내는 세 비트 전가산기를 \nt{GLUT} 뉴런으로 만들고 가중치와 문턱값을 밝혀라. 뉴런 여덟 개 안에 맞추어라. 그림~\ref{fig:dualxor}처럼 AND와 OR 소자로 만들면 몇 개가 필요한가?
\end{exercise}

\begin{exercise}[\lvlB]
\label{ex:02:5}
\emph{모델링: 쓰기에 걸리는 시간.} 그림~\ref{fig:recurrentgroup}의 무리: $\tau\dot r=-r+f(r+I)$, $f(u)=1/\bigl(1+e^{-(u-0.5)/0.08}\bigr)$이고, 쓰기 전에는 아래쪽 상태에 있다. 유입 $I$를 시간 $T$ 동안 가한다. $I=0.6$에서 비트를 쓰는 가장 작은 $T$와, 그보다 작으면 어떤 $T$로도 비트가 써지지 않는 문턱 $I_c$를 수치로 구하라.
\end{exercise}

\begin{exercise}[\lvlA]
\label{ex:02:6}
\emph{타이머로서의 사슬.} 사슬의 고리는 \nt{GLUT} 뉴런 $100$개이고, 지연은 $2\ms$에 독립적인 산포 $0.2\ms$가 있다. (a)~$1$~s를 재려면 뉴런이 몇 개 필요하고 상대 오차는 얼마인가? 오차는 구간에 어떻게 의존하는가? (b)~실제 타이머의 오차는 구간과 상관없이 $5\,\%$이다. 어떤 산포 원천이 이런 결과를 주는가?
\end{exercise}

\begin{exercise}[\lvlC]
\label{ex:02:7}
\emph{탐구: 플립플롭인가 커패시터인가.} 뉴런 $100$개가 비트 하나를 $10$~s 동안 유지한다. 플립플롭: 뉴런마다 초당 스파이크 $20$개. 커패시터: 촉진 $u$가 $\tau_F=1$~s로 줄어들고, $u\ge u_{\max}/2$이면 비트가 읽히며, 갱신은 뉴런마다 스파이크 세 개짜리 버스트이다. 커패시터는 몇 배 더 경제적인가? 무리를 $200\ms$ 동안 침묵시키면 비트는 어떻게 되는가?
\end{exercise}

\chapter{\nt{GLUT}와 \nt{GABA}의 협업}
\label{sec:gaba}
\begin{chaptergoals}
\item 금지 $a\wedge\bar b$로 \nt{GLUT}와 \nt{GABA}가 비트당 전선 하나로 완전해지는 방식;
\item 분로가 빼지 않고 나누는 이유와 억제가 일치 창을 좁히는 방식;
\item $\beta>1$인 상호 억제가 승자 하나를 고르는 방식과 신호로 기억을 초기화하는 방식;
\item 빠른 억제가 네트워크를 안정시키는 조건과 느린 억제가 리듬을 만드는 조건.
\end{chaptergoals}

\section{게임의 규칙}

\nt{GLUT} 뉴런만으로 된 네트워크는 고를 줄도, 기억을 초기화할 줄도, 활동이 사태로 번지지 않게 붙잡을 줄도 모른다. 모두 단조성 때문이다(명제~\ref{prop:monotone}). 여기에 수가 두 번째로 많은 유형인 \nt{GABA} 뉴런을 더하자. 규칙은 전과 같다. 살아 있는 몸에서 일어나는 일만 허용하고, 클록 신호는 없다.

뉴런 모델은 \eqref{eq:glutneuron}와 같지만, \nt{GABA} 뉴런에서 온 스파이크는 수신자의 전압을 올리지 않고 내린다. 이제 가중치에는 두 부호가 있고, 부호는 송신자가 정한다. 규칙~\eqref{eq:dalesign}이다.

회로의 파라미터 몇 가지. 피질 뉴런의 5분의 1에서 4분의 1이 \nt{GABA} 뉴런이다~\cite{Hendry1987}. 출력은 짧아서, 먼 영역이 아니라 이웃을 억제한다. 억제는 1밀리초 뒤에 도착해 몇 밀리초 동안 지속된다. 입력이 없으면 \nt{GABA} 뉴런은 침묵한다. 누군가를 억제하려면 자신이 먼저 흥분을 받아야 하며, 보통 같은 네트워크의 \nt{GLUT} 뉴런에게서 받는다.

그래서 \nt{GLUT} 뉴런 $A$에서 뉴런 $B$로 가는 음의 연결은 중개자를 거쳐야 한다. $A$가 \nt{GABA} 뉴런을 흥분시키고, 그 뉴런이 $B$를 억제한다. 부호를 바꾸는 값은 뉴런 하나와 지연 하나, 약 1밀리초이다.

억제에서는 연결이 \emph{누구에게서} 오는지뿐 아니라 \emph{어디에} 붙는지도 중요하다. 붙는 자리가 연결이 하는 일을 상당 부분 결정한다~\cite{Markram2004}. \nt{GLUT} 출력은 주로 수신자의 입력 가지, 즉 \emph{가지돌기}에 붙어 데이터를 나른다. 이런 입력 하나하나는 합의 한 항이다. 뉴런 세포체에 붙는 출력은 합 전체에 작용한다. 많은 빠른 \nt{GABA} 뉴런이 이런 식으로 수신자의 응답을 나누고, 수신자가 언제 발화할지를 정한다. 수신자의 스파이크가 생겨나는 자리에 붙는 출력은 최종 결정권, 즉 출력 전체에 대한 금지이다. 그리고 다른 뉴런의 전선 말단에 붙는 출력은 입력 선로 하나의 방출 확률 $p$를 바꾸고 나머지는 건드리지 않는다~\cite{Rudomin1999}. \ref{sec:pain}장의 통증 관문이 특히 이렇게 작동한다. 뇌 바깥에서는 출력이 근육 섬유(\ref{sec:muscles}장)와 장기(\ref{sec:chemoutput}장)에 붙으며, 어떤 출력은 아무 데도 붙지 않고 신경전달물질을 조직 공간이나 혈액으로 내보낸다(\ref{sec:serocomputer}장과~\ref{sec:chemoutput}장).

\section{NOT과 금지}

이제 NOT을 만들 수 있는가? 거의 그렇다. 입력이 조용할 때 작동하는 뉴런은 어디선가 흥분을 얻어야 한다. 살아 있는 뇌에는 그것이 있다. 모든 뉴런에는 다른 뉴런 수천 개에서 오는 배경 활동이 끊임없이 들어온다. 배경이 뉴런 $Y$를 작동 상태로 유지하고, 입력 $x$가 \nt{GABA} 중개자를 거쳐 그것을 억제한다고 하자. 이것이 NOT이다.

그러나 더 유용한 것은 \emph{금지}, 곧 “$a$, 단 $b$일 때는 아님”이다.
\begin{empheq}[box=\keybox]{equation}
y \;=\; a \wedge \bar b .
\label{eq:veto}
\end{empheq}
뉴런 $Y$는 $a$에서 흥분을, $b$가 흥분시키는 \nt{GABA} 뉴런에서 억제를 받는다. 여기서는 배경이 필요 없다. 흥분은 $a$ 자신이 준다. 금지와 OR로 모든 것을 조립할 수 있다. 예를 들어 배타적 OR는 $x\oplus y=(x\wedge\bar y)\vee(\bar x\wedge y)$로, 금지 뉴런 둘, \nt{GABA} 중개자 둘, OR 뉴런 하나이다.

\begin{keyblock}
\begin{proposition}[\nt{GLUT}+\nt{GABA}의 완전성]
\label{prop:gabacomplete}
\nt{GLUT} 뉴런과 \nt{GABA} 뉴런으로 된 네트워크는 입력의 어떤 논리 함수든 계산할 수 있으며, 각 비트는 전선 한 가닥으로 전달된다. 이를 위해 \ref{sec:glutcomputer}장의 “켜짐--꺼짐” 센서는 더 이상 필요 없다. 모든 입력이 조용할 때 1을 내야 하는 함수라면 배경 활동의 원천이 필요하다.
\end{proposition}
\end{keyblock}

증명: 상수 1이 있으면 금지~\eqref{eq:veto}와 OR는 함수적으로 완전하다. NOT $x$는 “1, 단 $x$일 때는 아님”이라는 금지이기 때문이다. 그리고 모든 입력이 조용할 때 0을 내는 함수는 1 없이도 금지와 OR로 조립된다.

\section{운동 방향}

\ref{sec:glutcomputer}장의 시간 속 AND처럼, 시간 속의 금지는 운동 방향 검출기를 준다.

이웃한 두 센서 $A$와 $B$가 서로 $\Delta x$ 떨어져 있다고 하자. 출력 뉴런 $Y$는 $B$에서 흥분을, $A$에서는 \nt{GABA} 중개자를 거쳐 지연 $d$로 억제를 받는다(그림~\ref{fig:direction}). 속도 $v$로 $A$에서 $B$로 달리는 빛 점은 시각 $0$에 $A$를 지나고, 억제는 시각 $d$에 $Y$에 이른다. 빛 점은 시각 $\Delta x/v$에 $B$를 지나고, 그때 흥분이 도착한다. 이 두 시각이 창~\eqref{eq:window}의 정밀도 안에서 일치하면 억제가 흥분을 끄고 $Y$는 침묵한다. 이는 대략 다음 속도에서 일어난다.
\begin{empheq}[box=\keybox]{equation}
v^* \;=\; \frac{\Delta x}{d} .
\label{eq:vstar}
\end{empheq}
$B$에서 $A$로 움직일 때는 $B$에서 오는 흥분이 시각 $0$에 도착하고, $A$에서 오는 억제는 $Y$가 이미 발화한 뒤인 시각 $\Delta x/v+d$에야 도착한다.

\begin{figure}[t]
\centering
\begin{tikzpicture}[>=Stealth,
  nrn/.style={circle,draw,very thick,fill=blue!6,minimum size=0.9cm,inner sep=0pt},
  gab/.style={nrn,fill=red!10},
  inh/.style={-{Bar[width=7pt]},thick,red!70!black}]
\node[nrn] (A) at (0,2) {$A$};
\node[nrn] (B) at (3,2) {$B$};
\node[gab] (G) at (0.9,0.6) {\small \nt{GABA}};
\node[nrn] (Y) at (3,-0.6) {$Y$};
\draw[->,thick] (A) -- (G);
\draw[inh] (G) -- node[below left=-2pt] {\scriptsize 지연 $d$} (Y);
\draw[->,thick] (B) -- (Y);
\draw[->,very thick,red!70!black] (Y) -- (4.6,-0.6) node[right,black] {\small 출력};
\draw[->,thick,orange!85!black] (-0.6,2.9) -- (3.6,2.9) node[right,black] {\small 침묵};
\draw[<-,thick,orange!85!black] (-0.6,3.4) -- (3.6,3.4) node[right,black] {\small 응답};
\foreach \yy/\dd in {2.9/-1,3.4/1}{
  \foreach \k in {1,2,3}{\draw[orange!70!yellow,line width=0.8pt] (1.5+\dd*0.12+\dd*0.13*\k,\yy+0.1-0.1*\k+0.1) -- ++(\dd*0.25,0);}
  \fill[yellow!70!orange,draw=orange!70!black] (1.5,\yy) circle (0.14);}
\node[font=\scriptsize,align=center] at (1.5,3.95) {빛 점};
\end{tikzpicture}
\caption{운동 방향 검출기. $B$는 $Y$를 흥분시키고, $A$는 \nt{GABA} 중개자를 거쳐 지연 $d$로 $Y$를 억제한다. $A$에서 $B$로 약 $\Delta x/d$의 속도로 움직이면 억제가 흥분을 끈다. 반대로 움직이면 억제가 늦는다. 짧은 줄이 달린 노란 점은 센서 위를 달리는 빛이다.}
\label{fig:direction}
\end{figure}

망막에서 운동 방향 선택성은 바로 이렇게 만들어지는 것으로 보인다. 운동이 응답을 일으켜서는 안 되는 쪽에서 오는 비대칭 억제이다~\cite{Barlow1965}.

\section{뺄셈인가 나눗셈인가}

지금까지 억제는 뺄셈을 했다. 실제로는 더 섬세하다.

시냅스는 채널을 연다. 즉 컨덕턴스를 켠다. 흥분성 시냅스의 평형 전위는 문턱값보다 훨씬 위, 휴지 전위보다 약 $70\mV$ 위에 있다. 열린 채널은 전압을 위로 끌어올린다. 억제성 \nt{GABA} 시냅스의 채널은 염소를 통과시키며, 염소의 평형 전위는 휴지 전위 근처에 있다. 그래서 열린 채널은 전압을 아래로가 아니라 \emph{휴지 전위 쪽으로} 끈다. 전압을 휴지 전위에서부터 재면, 누설 컨덕턴스 $g_L$, 흥분성 컨덕턴스 $g_E$, 억제성 컨덕턴스 $g_I$를 가진 막의 정상 상태 전압은 다음과 같다.
\begin{empheq}[box=\keybox]{equation}
V_\infty \;=\; \frac{g_E\,E_E}{g_L+g_E+g_I} ,
\label{eq:shunt}
\end{empheq}
여기서 $E_E\approx70\mV$는 흥분성 시냅스의 평형 전위이다. 이것은 전압 분배기에 대한 옴의 법칙이다. $g_E$는 $E_E$ 쪽으로, $g_L$과 $g_I$는 0 쪽으로 끈다.

$g_I$는 분자에 음의 부호로 있지 않고 분모에 있다. 억제는 흥분에서 \emph{빼지} 않고 흥분을 \emph{나눈다}(그림~\ref{fig:shunt}). 이런 억제를 분로 억제라고 한다. 열린 염소 채널이 막을 단락시키기 때문이다. 네트워크에게 이것은 이득 조절이다. $g_I$가 이웃의 총활동에 비례하면, 각 뉴런은 자기 입력의 절대 세기가 아니라 전체 활동 가운데 그 입력이 차지하는 몫에 응답한다. 총활동으로 나누는 이런 연산은 뇌의 여러 부위에서 나타나며, 뇌의 표준 연산 가운데 하나로 여겨진다~\cite{Carandini2012}. 같은 네트워크가 약한 입력에서도 강한 입력에서도 작동한다.

단서가 있다. 식~\eqref{eq:shunt}는 스파이크를 내지 않는 뉴런에 관한 것이다. 발화하는 뉴런에서는 전압이 늘 문턱값 근처에 머물러 억제성 컨덕턴스를 지나는 전류가 거의 일정하므로, 분로는 스파이크 \emph{발화율}에 주로 뺄셈으로 작용한다~\cite{HoltKoch1997}. 뉴런에 균형 잡힌 흥분성, 억제성 배경 컨덕턴스를 함께 더해 주면, 즉 살아 있는 피질의 흔들리는 균형 체제와 같으면 발화율이 나누어진다~\cite{Chance2002}. 그러니 뇌는 나눗셈을 분로 하나에서가 아니라 억제와 배경 잡음의 조합에서 얻는다~\cite{Carandini2012}.

\begin{figure}[t]
\centering
\begin{tikzpicture}[>=Stealth,x=1.25cm,y=0.05cm]
\draw[->,gray!70!black] (0,0) -- (5.4,0) node[right,black] {\small $g_E/g_L$};
\draw[->,gray!70!black] (0,0) -- (0,68) node[above,black] {\small $V_\infty$, mV};
\foreach \v in {20,40,60}{\draw[gray!60!black] (-0.05,\v) -- (0.05,\v) node[left=3pt,font=\scriptsize] {\v};}
\foreach \x in {1,2,3,4,5}{\draw[gray!60!black] (\x,-1.5) -- (\x,1.5) node[below=5pt,font=\scriptsize] {\x};}
\draw[very thick,blue!60!black] plot[domain=0:5,samples=60] (\x,{70*\x/(1+\x)});
\draw[very thick,red!70!black] plot[domain=0:5,samples=60] (\x,{70*\x/(3+\x)});
\draw[thick,densely dashed,gray!50!black] plot[domain=0.56:5,samples=60] (\x,{70*\x/(1+\x)-25});
\node[right,font=\small,blue!60!black] at (5.02,58.3) {억제 없음};
\node[right,font=\small,red!70!black] at (5.02,45.5) {나눗셈: $g_I=2g_L$};
\node[right,font=\small,gray!40!black] at (5.02,32.5) {$25\mV$ 뺄셈};
\end{tikzpicture}
\caption{분로 억제는 전압에서 빼지 않고 전압을 나눈다. 파란 곡선은 억제가 없을 때의 정상 상태 전압~\eqref{eq:shunt}, 빨간 곡선은 억제성 컨덕턴스 $g_I=2g_L$이 있을 때이다. 빨간 곡선은 파란 곡선을 가로로 세 배 늘인 것과 같다. 입력을 $1+g_I/g_L=3$으로 나눈 셈이다. 파선은 일정한 $25\mV$를 빼는 억제이다. 곡선 전체가 내려갈 뿐 기울기는 변하지 않는다.}
\label{fig:shunt}
\end{figure}

두 번째 결과도 있다. 막의 시간 상수는 $\tau_{\text{eff}}=C/(g_L+g_E+g_I)$이고, 억제는 이것을 줄인다. 일치 창~\eqref{eq:window}는 $\tau$에 비례하므로, 흥분과 함께 온 억제는 \emph{창을 좁힌다}. 입력이 뉴런을 직접 흥분시키면서 동시에 1밀리초 지연으로 \nt{GABA} 중개자를 거쳐 억제하는 회로는 뇌에서 가장 흔한 회로 가운데 하나이다. 뉴런은 처음 1--2밀리초 안에 온 것에만 응답할 시간이 있다~\cite{Pouille2001}.

\section{선택}

이제 \nt{GLUT} 네트워크에 가장 부족했던 것, 여럿 가운데 하나를 고르는 일이다. 입력이 $I_1$과 $I_2$인 \nt{GLUT} 뉴런 무리 둘을 잡자. 각 무리는 자기 \nt{GABA} 중개자를 거쳐 세기 $\beta$로 상대를 억제한다(그림~\ref{fig:wta}). 발화율 $r_1,r_2$에 대해 다음을 얻는다.
\begin{equation}
\tau\,\frac{\dd r_1}{\dd t} = -r_1 + \bigl[I_1-\beta r_2\bigr]_+ ,\qquad
\tau\,\frac{\dd r_2}{\dd t} = -r_2 + \bigl[I_2-\beta r_1\bigr]_+ ,
\label{eq:wta}
\end{equation}
여기서 $[u]_+=\max(u,0)$이다. 발화율은 음수가 될 수 없다.

\begin{figure}[t]
\centering
\begin{tikzpicture}[>=Stealth,
  nrn/.style={circle,draw,very thick,fill=blue!6,minimum size=0.9cm,inner sep=0pt},
  gab/.style={nrn,fill=red!10,minimum size=0.75cm},
  inh/.style={-{Bar[width=7pt]},thick,red!70!black}]
\node[nrn] (E1) at (0,0) {$r_1$};
\node[nrn] (E2) at (4,0) {$r_2$};
\node[gab] (G1) at (1.4,1.1) {};
\node[gab] (G2) at (2.6,-1.1) {};
\draw[->,thick] (E1) -- (G1); \draw[inh] (G1) -- (E2);
\draw[->,thick] (E2) -- (G2); \draw[inh] (G2) -- (E1);
\draw[->,thick,blue!60!black] (-1.6,0) node[left,black] {$I_1$} -- (E1);
\draw[->,thick,blue!60!black] (5.6,0) node[right,black] {$I_2$} -- (E2);
\end{tikzpicture}
\caption{둘 중 하나 고르기. \nt{GLUT} 뉴런 무리 둘(파란색)이 \nt{GABA} 중개자(빨간색)를 거쳐 서로를 억제한다.}
\label{fig:wta}
\end{figure}

두 뉴런이 모두 작동하는 동안 시스템~\eqref{eq:wta}는 선형이며, 그 거동은 두 고윳값이 정한다. 두 발화율이 같은 방향으로 벗어날 때는 $-(1+\beta)/\tau$, 반대 방향으로 벗어날 때는 $-(1-\beta)/\tau$이다. 억제가 약하면($\beta<1$) 두 고윳값 모두 음수이다. 두 뉴런이 다 작동하고, 억제는 약한 쪽을 조금 누를 뿐이다. 억제가 강하면($\beta>1$) 두 번째 고윳값이 양수가 된다. $r_1$과 $r_2$ 사이의 차이는 무엇이든 한 뉴런이 완전히 침묵할 때까지 커진다.

\begin{keyblock}
\begin{proposition}[승자독식]
\label{prop:wta}
상호 억제가 1보다 강하면($\beta>1$), 두 무리가 모두 작동하는 상태는 불안정하다. 네트워크는 한 무리가 작동하고 다른 무리는 침묵하는 상태에 이른다. 발화율 $r_1=I_1$인 승자는 $I_2<\beta I_1$인 한 승리를 지킨다.
\end{proposition}
\end{keyblock}

이것을 모델로 확인했다. $\beta=0.5$, 입력 $1.0$과 $0.9$에서는 두 무리가 모두 발화율 $0.73$과 $0.53$으로 작동한다. $\beta=2$에서 같은 입력이면 두 번째 무리는 완전히 침묵한다. 이런 선택에는 기억이 있다. 나중에 입력을 맞바꾸어도 이전 승자가 계속 승자로 남는다. $1.0<2\cdot0.9$이기 때문이다.

무리가 백 개여도 마찬가지이다. 각 무리가 나머지를 억제하고, 하나만 살아남는다.

\section{기억 초기화}

\ref{sec:glutcomputer}장의 메모리는 피로로만 꺼졌다. 여기에 “초기화” 신호로 흥분되어 무리를 억제하는 \nt{GABA} 뉴런을 더하자. 억제가 그림~\ref{fig:bistable}의 곡선 $f(wr+I)$를 끌어내리면 위쪽 교점이 사라지고, 무리는 꺼진다. 그리고 초기화 신호가 끝난 뒤에도 아래쪽 상태에 머문다. 이제 메모리는 한 신호로 써지고 다른 신호로 지워진다. 세트 입력과 리셋 입력이 있는 플립플롭과 같다.

더 멋진 방법은 이런 무리 둘을 그림~\ref{fig:wta}처럼 상호 억제로 잇고, 각각에 자기 되먹임을 주는 것이다. 한 무리의 입력에 신호가 오면 셀은 그 무리의 상태로 넘어가고, 셀은 마지막 결정을 기억한다. NOR 소자 둘을 교차 연결한 래치와 똑같은 구조이다.

\section{안정성과 리듬}

\nt{GLUT} 네트워크의 세 번째 문제, 사태가 남았다. 자기 되먹임 $w_{EE}$를 가진 \nt{GLUT} 뉴런 무리와, 첫 무리가 세기 $w_{IE}$로 흥분시키고 첫 무리를 세기 $w_{EI}$로 억제하는 \nt{GABA} 뉴런 무리를 잡자. 발화율 $E$와 $I$에 대해 다음을 얻는다.
\begin{equation}
\tau_E\frac{\dd E}{\dd t} = -E + w_{EE}E - w_{EI}I + h, \qquad
\tau_I\frac{\dd I}{\dd t} = -I + w_{IE}E ,
\label{eq:ei}
\end{equation}
여기서 $h$는 바깥 입력이다~\cite{WilsonCowan1972}. 억제가 없으면 $w_{EE}>1$에서 활동이 한없이 커진다. 스파이크 하나가 다음 스파이크를 하나보다 많이 낳는다. 억제가 있으면 정상 상태 발화율
\begin{equation}
E^* \;=\; \frac{h}{1-w_{EE}+w_{EI}w_{IE}}
\label{eq:eistar}
\end{equation}
은 분모가 양수일 때 유한하다. 그러나 그것으로는 부족하다. 평형이 끌어당기기도 해야 한다. \eqref{eq:ei}의 선형 해석에서 두 조건이 나온다.

\begin{keyblock}
\begin{proposition}[안정 조건]
\label{prop:ei}
평형~\eqref{eq:eistar}은 억제가 다음과 같을 때 안정하다.
\begin{enumerate}
\item[(i)] \emph{충분히 강하다}: $w_{EI}\,w_{IE} > w_{EE}-1$;
\item[(ii)] \emph{충분히 빠르다}: $\tau_I < \tau_E/(w_{EE}-1)$.
\end{enumerate}
(i)이 깨지면 활동이 사태처럼 불어난다. (i)은 성립하지만 (ii)가 깨지면 억제가 늦어 활동이 진동하기 시작한다.
\end{proposition}
\end{keyblock}

조건~(i)은 시스템~\eqref{eq:ei} 행렬의 행렬식이고, 조건~(ii)는 그 대각합이다. 두 번째 조건의 의미는 제어기를 튜닝해 본 사람이면 누구나 안다. 늦은 되먹임은 편차를 가라앉히지 않고 오히려 흔든다.

\begin{figure}[t]
\centering
\begin{tikzpicture}[>=Stealth]
\draw[->,gray!70!black] (0,0) -- (10.9,0) node[right,black] {\small $t$, ms};
\draw[->,gray!70!black] (0,0) -- (0,3.3) node[left,black] {\small $E$};
\foreach \t in {0,50,100,150}{\draw[gray!70!black] (\t*0.07,0.06) -- (\t*0.07,-0.06) node[below,font=\scriptsize] {\t};}
\input{figures/ei_traces}
\node[font=\small,gray!40!black,anchor=west] at (1.3,3.45) {억제 없음: 사태};
\node[font=\small,blue!60!black,anchor=west] at (7.4,1.25) {빠른 억제};
\node[font=\small,red!70!black,anchor=west] at (7.4,2.55) {느린 억제};
\end{tikzpicture}
\caption{입력을 켠 뒤, 되먹임 $w_{EE}=1.5$, $\tau_E=10\ms$인 \nt{GLUT} 무리의 발화율. 억제가 없으면(파선) 활동이 한없이 커진다. 세기 $w_{EI}w_{IE}=2$, $\tau_I=2\ms$인 억제가 있으면(파란 곡선) 수준 $E^*=h/(1-1.5+2)=0.67h$에 자리 잡는다. 억제가 같되 느리면, $\tau_I=30\ms$(빨간 곡선), 활동이 진동한다.}
\label{fig:ei}
\end{figure}

피질은 바로 이런 체제에서 작동한다고 여겨진다. 피질 자신의 되먹임은 억제가 없으면 사태로 치달을 만큼 강하고, 빠른 억제만이 피질을 동작점에 붙잡아 둔다~\cite{Tsodyks1997isn}.

이 체제에는 바깥에서 알아볼 수 있는 역설적인 특징이 있다~\cite{Tsodyks1997isn}. \eqref{eq:ei}에서 \nt{GABA} 무리에 바깥 입력 $h_I$를 직접 더하자. 정상 상태 발화율은 다음과 같다.
\begin{equation}
E^* \;=\; \frac{h-w_{EI}\,h_I}{D},\qquad
I^* \;=\; \frac{w_{IE}\,h+(1-w_{EE})\,h_I}{D},\qquad
D \;=\; 1-w_{EE}+w_{EI}w_{IE} .
\label{eq:isn}
\end{equation}
$w_{EE}>1$이면 두 번째 식에서 $h_I$의 계수는 음수이다. \nt{GABA} 뉴런을 밀어 주면 그 발화율이 \emph{떨어진다}. 원인은 루프에 있다. 여분의 억제가 \nt{GLUT} 무리를 누르고, \nt{GLUT} 무리는 \nt{GABA} 뉴런을 덜 흥분시키며, \nt{GABA} 뉴런은 얻은 것보다 더 많이 잃는다. 그림~\ref{fig:ei}의 수치($w_{EE}=1.5$, $w_{EI}w_{IE}=2$, $D=1.5$)에서는 더한 입력 한 단위마다 $I^*$가 3분의 1 단위씩 줄어든다. 이 특징은 피질에서 발견되었다. 시각 피질 뉴런의 응답이 주변 자극으로 억눌릴 때, 그 뉴런들이 받는 억제 전류는 늘지 않고 흥분 전류와 함께 줄어든다~\cite{Ozeki2009}. 나중에 같은 특징이 피질의 여러 영역과 층에서 발견되었다~\cite{Sanzeni2020}.

조건~(ii)가 깨질 때의 진동도 뇌에서 나타난다. “\nt{GLUT}가 \nt{GABA}를 흥분시키고, \nt{GABA}가 \nt{GLUT}를 억제한다”는 루프는 몇 밀리초의 지연으로 주기 $12$--$50\ms$ 정도(주파수 $20$--$80$\,Hz)의 리듬을 낳으며, 피질의 이런 빠른 리듬은 잘 알려져 있다~\cite{Cardin2009,Buzsaki2012}. 이것은 컴퓨터의 클록이 아니다. 리듬은 국소적이고 네트워크가 활성일 때만 생긴다. 그러나 몇 주기 동안 리듬은 시간을 뉴런이 발화할 수 있는 창들로 나눈다.

\section{요약}

\begin{table}[t]
\centering
\small
\begin{tabular}{>{\raggedright}p{3.3cm}>{\raggedright}p{4.4cm}>{\raggedright\arraybackslash}p{4.4cm}}
\toprule
& \nt{GLUT}만 & \nt{GLUT} + \nt{GABA} \\
\midrule
NOT & “켜짐--꺼짐” 센서를 통해서만 & 금지 $a\wedge\bar b$; 배경 활동이 있으면 NOT \\
비트당 전선 수 & 둘 & 하나 \\
운동 방향 & 없음 & 지연된 금지 \\
이득 조절 & 없음 & 나눗셈: 분로와 배경 잡음의 조합 \\
일치 창 & $\sim\tau$ & 억제로 좁아짐 \\
여럿 중 하나 고르기 & 없음 & 상호 억제, $\beta>1$ \\
기억 초기화 & 피로로만 & \nt{GABA}를 거친 신호로 \\
안정성 & 피로로만 & 빠른 음의 되먹임 \\
리듬 & 필요 없음 & 느린 억제에서 생김 \\
\bottomrule
\end{tabular}
\caption{\nt{GABA}가 \nt{GLUT} 뉴런 네트워크에 더하는 것.}
\label{tab:gaba}
\end{table}

\nt{GABA}는 네트워크에 새로운 \emph{계산}을 거의 더하지 않는다. 그것들은 모두 이중 레일 센서로도 계산할 수 있었다. 대신 \nt{GABA}는 \emph{제어}를 더한다(표~\ref{tab:gaba}). 선택, 초기화, 이득 조절, 안정성이다. 뇌에서 흥분은 데이터를 나르고, 억제는 어떤 데이터가 언제 얼마나 크게 지나갈지를 정한다. 피질 뉴런 네다섯 개 중 하나가 맡는 일로서 매우 합리적인 분업이다.

\section{연습문제}

\begin{exercise}[\lvlA]
\label{ex:03:1}
\emph{숫자로 본 분로.} $E_E=70\mV$. \eqref{eq:shunt}로 $g_E=g_L$과 $4g_L$일 때 억제가 없는 경우와 분로 $g_I=2g_L$이 있는 경우의 $V_\infty$를 구하라. $g_E=g_L$에서 분로와 같은 양을 빼는 뺄셈이라면 $g_E=4g_L$에서 $V_\infty$는 얼마인가? 분로는 일치 창을 몇 배 좁히는가?
\end{exercise}

\begin{exercise}[\lvlB]
\label{ex:03:2}
\emph{안정 조건의 유도.} 선형 시스템~\eqref{eq:ei} 행렬의 대각합과 행렬식을 구하고, 여기서 명제~\ref{prop:ei}의 두 조건을 유도하라. 조건~(ii)의 경계에서 평형은 진동하며 안정성을 잃는다. 그림~\ref{fig:ei}의 수치에서 이 진동의 주기를 구하라.
\end{exercise}

\begin{exercise}[\lvlB]
\label{ex:03:3}
\emph{지연에서 생기는 리듬.} 모델~\eqref{eq:ei}의 느린 억제를 지연 $d$가 있는 빠른 억제로 바꾸자. $\tau_E\dot E=-E(t)-GE(t-d)+h$. $\tau_E=10\ms$일 때 진동을 낳는 가장 작은 지연 $d_c$와, $G=2$와 $G=5$에서의 진동 주기를 구하라. 연습문제~\ref{ex:03:2}와 달리 이 진동은 피질의 빠른 리듬 $12$--$50\ms$에 들어가는가?
\end{exercise}

\begin{exercise}[\lvlB]
\label{ex:03:4}
\emph{모델링: 기억이 있는 선택.} $\beta=2$, $\tau=1$에서 \eqref{eq:wta}를 적분하라. (a)~$I_1=1$에서 $I_2$를 $0$에서 $3$까지 천천히 올렸다가 다시 내려라. 어떤 $I_2$에서 네트워크가 전환되는가? (b)~입력 차 $\varepsilon$이 10분의 1로 줄면, 침묵에서 출발한 결정 시간은 얼마나 늘어나는가?
\end{exercise}

\begin{exercise}[\lvlB]
\label{ex:03:5}
\emph{설계: 속도 대역 검출기.} 그림~\ref{fig:direction}의 검출기는 $A$에서 $B$로 가는 통과 시간이 $5$--$20\ms$인 움직임에서 침묵해야 한다. 지연 $d_k$인 \nt{GABA} 중개자는 $A$의 스파이크 뒤 구간 $[d_k,\,d_k+5\ms]$ 동안 $Y$를 금지한다. $B$에서 오는 흥분은 곧바로 도착한다. 중개자가 몇 개, 어떤 지연으로 필요한가?
\end{exercise}

\begin{exercise}[\lvlB]
\label{ex:03:6}
\emph{금지와 배경.} “입력마다 \nt{GABA} 중개자, $f=1$인 입력 조합마다 금지, 공통 OR”라는 일반 회로로 패리티 $x\oplus y\oplus z$를 만들면 뉴런이 몇 개 필요한가? 입력을 직접 받는 뉴런 두 개, \nt{GABA}와 \nt{GLUT}로 패리티를 만들고 가중치와 문턱값을 밝혀라.
\end{exercise}

\begin{exercise}[\lvlC]
\label{ex:03:7}
\emph{설계: 백 중 하나 고르기.} \nt{GLUT} 무리 $100$개가 저마다 나머지 모두를 똑같이 억제해야 하지만 자기 자신은 억제하지 않아야 한다. 선형 \nt{GABA} 중개자는 무리들에게서 흥분을 받고 무리들을 억제한다. 무리들은 서로를 흥분시킬 수 있지만 자기 자신은 흥분시킬 수 없다. 중개자의 최소 개수는 얼마인가? 직접적인 흥분 연결이 전혀 없다면?
\end{exercise}

\begin{exercise}[\lvlC]
\label{ex:03:8}
\emph{탐구: 억제는 언제 역설적인가.} 그림~\ref{fig:ei}의 모델($w_{EI}=2$, $w_{IE}=1$)에서 \nt{GLUT} 무리는 전달 함수 $f(u)=[u]_+^2$를 얻었고, \nt{GABA} 무리는 선형으로 남았다. 어떤 입력 $h_c$보다 크면 \nt{GABA} 무리에 더한 입력이 그 무리 자신의 발화율을 낮추는가? 시뮬레이션으로 확인하라.
\end{exercise}

\chapter{모스 부호: \nt{GLUT} 뉴런과 \nt{GABA} 뉴런이 말하는 언어에 대해 우리가 아는 것}
\chaptermark{모스 부호}
\label{sec:code}

\section{기호 하나짜리 알파벳}

모스 부호에는 점과 대시라는 두 기호가 있고, 그 사이에 세 가지 길이의 쉼이 있다. \ref{sec:neurontypes}장에서 보았듯이 뉴런의 알파벳은 더 빈약하다. 기호가 하나뿐이다. 스파이크는 약 1밀리초 동안 지속되고, 한 뉴런의 스파이크는 높이와 모양이 모두 같다. 그러므로 메시지 전체는 쉼에, 즉 시각의 열 $t_1,t_2,t_3,\dots$에 담긴다. 대시가 없고 점의 리듬만이 의미를 싣는 부호이다(그림~\ref{fig:morse}).

\begin{figure}[t]
\centering
\begin{tikzpicture}[>=Stealth,x=0.08cm,y=1cm]
% Morse letter: dot, dash, dot
\node[left,font=\small] at (-6,2.55) {모스:};
\fill[black] (0,2.45) rectangle (6,2.65);
\fill[black] (14,2.45) rectangle (32,2.65);
\fill[black] (40,2.45) rectangle (46,2.65);
\node[right,font=\small,gray!40!black] at (52,2.55) {점, 대시, 점: 의미는 길이와 쉼에 있다};
% stimulus onset
\node[left,font=\small] at (-6,0.4) {뉴런:};
\draw[->,thick,green!45!black] (0,-0.55) -- (0,-0.05);
\node[below,font=\scriptsize,green!45!black] at (0,-0.55) {자극};
\draw[gray!70!black] (0,0) -- (152,0) node[right,black] {\small $t$};
% spikes: single ones blue, the burst red
\foreach \s in {14,26,41,88,112,131}{\draw[very thick,blue!60!black] (\s,0) -- (\s,0.8);}
\foreach \s in {60,63,66,69}{\draw[very thick,red!70!black] (\s,0) -- (\s,0.8);}
% latency
\draw[<->,thick] (0,1.1) -- node[above,font=\scriptsize] {$t_1$: 첫 스파이크 시각} (14,1.1);
% burst
\draw[decorate,decoration={brace,amplitude=4pt},red!70!black] (58,0.95) -- (71,0.95) node[midway,above=4pt,font=\scriptsize] {버스트: `대시'};
% counting windows
\foreach \a/\b/\n in {0/50/3,50/100/5,100/150/2}{
  \draw[decorate,decoration={brace,amplitude=4pt,mirror},gray!50!black] (\a+1,-0.2) -- (\b-1,-0.2) node[midway,below=4pt,font=\small] {$n=\n$};}
\node[font=\scriptsize,gray!40!black] at (75,-1.05) {발화율 부호: 창마다 스파이크 수 $n$};
\foreach \x in {0,50,100,150}{\draw[gray!60!black] (\x,-0.06) -- (\x,0.06);}
\end{tikzpicture}
\caption{모스 부호와 뉴런의 부호. 같은 스파이크 열을 여러 방식으로 읽을 수 있다. $50\ms$ 창마다 스파이크를 세거나(\ref{sec:rate}절), 지연 $t_1$을 재거나~\eqref{eq:latency}, `대시'인 버스트를 알아챌 수 있다~\eqref{eq:burst}.}
\label{fig:morse}
\end{figure}

어떤 쉼이 가능한가? 아래로는 불응기가 제한한다. 스파이크 뒤 약 1밀리초 동안 뉴런은 다시 발화할 수 없으므로 초당 수백 개보다 많은 스파이크는 나오지 않는다. 위로는 아무 제한이 없다. 뉴런은 몇 분 동안 침묵할 수 있다. 피질 \nt{GLUT} 뉴런의 발화율은 초당 1개 미만에서 수십 개 사이에 있고, 대부분의 뉴런은 대부분의 시간을 아래쪽 경계 근처에서 작동한다.

\ref{sec:glutcomputer}장과 \ref{sec:gaba}장에서 우리는 스파이크로 수를 적는 방식을 그때그때 하나씩 가정했다. 이제 직접 묻자. \nt{GLUT} 뉴런과 \nt{GABA} 뉴런은 서로 어떤 언어로 말하는가? 완전한 답은 없다. 이 언어는 해독되지 않았다. 그러나 확실히 알려진 것도 있으며, 알려진 것의 거의 전부는 통신 엔지니어처럼 추론해서 얻을 수 있다. 채널 용량, 잡음, 수신기, 전송 비용이다.

\section{얼마나 전달할 수 있는가}

상한부터 시작하자. 수신자가 스파이크 시각을 정밀도 $\Delta t$로 구별한다고 하자. $\Delta t$보다 작은 이동은 수신자에게 구별되지 않는다. 시간을 불응기보다 짧은 길이 $\Delta t$의 칸으로 자르면, 각 칸에는 스파이크가 하나 있거나 없다(그림~\ref{fig:cells}). 평균 발화율이 $r$이면 스파이크가 칸에 들어갈 확률은 $p=r\Delta t$이다. 흐름이 가장 많은 정보를 나르는 것은 칸들이 독립일 때이며, 이때 각 칸은 $p\ll1$에서 엔트로피 $-p\log_2p-(1-p)\log_2(1-p)\approx p\log_2(e/p)$비트를 준다. $\Delta t$로 나누면 전송 속도의 한계와 스파이크 하나당 몫을 얻는다~\cite{MacKay1952}.
\begin{empheq}[box=\keybox]{equation}
R_{\max} \;\approx\; r\,\log_2\frac{e}{r\,\Delta t}\ \ \text{bit/s},
\qquad
\frac{R_{\max}}{r} \;\approx\; \log_2\frac{e}{r\,\Delta t}\ \ \text{스파이크당 bit}.
\label{eq:spikeentropy}
\end{empheq}
$r$이 초당 스파이크 10개이고 $\Delta t=1\ms$이면, 스파이크당 약 8비트, 초당 약 80비트이다.

\begin{figure}[t]
\centering
\begin{tikzpicture}[>=Stealth,x=0.5cm,y=1cm]
% time cut into cells of width Delta t; a cell holds one impulse or none
\foreach \k in {0,...,23}{\draw[gray!55] (\k,0) rectangle (\k+1,0.7);}
\foreach \k in {2,7,8,15,21}{
  \fill[red!10] (\k,0) rectangle (\k+1,0.7);
  \draw[very thick,red!70!black] (\k+0.5,0.05) -- (\k+0.5,0.65);}
\foreach \k in {0,...,23}{
  \pgfmathtruncatemacro{\bit}{(\k==2)||(\k==7)||(\k==8)||(\k==15)||(\k==21)}
  \ifnum\bit=1 \def\bitcol{red!70!black}\else \def\bitcol{gray!60!black}\fi
  \node[font=\scriptsize,text=\bitcol] at (\k+0.5,-0.25) {\bit};}
\draw[<->,gray!60!black] (0,0.95) -- (1,0.95) node[midway,above,font=\scriptsize,black] {$\Delta t$};
\node[font=\small,anchor=east] at (-0.3,0.35) {스파이크};
\node[font=\small,anchor=east] at (-0.3,-0.25) {비트};
\draw[decorate,decoration={brace,amplitude=5pt,mirror},gray!60!black] (0,-0.55) -- (24,-0.55)
  node[midway,below=6pt,font=\scriptsize,black,align=center] {세기만 하는 수신자: `$n=5$';\\ 칸을 보는 수신자: $24$칸 중 정확히 어느 $5$칸인가, 즉 $\log_2\binom{24}{5}\approx15$비트};
\end{tikzpicture}
\caption{용량~\eqref{eq:spikeentropy}은 어디서 오는가. 시간을 길이 $\Delta t$의 칸으로 자르면, 칸마다 스파이크가 있거나($1$) 없다($0$). 스파이크 흐름은 이진 문자열이다. 스파이크가 칸에 들어갈 확률은 $p=r\Delta t$이다. 칸이 잘수록 문자열이 길어지고 비트가 많아진다. 스파이크 계수기는 1이 \emph{어디에} 있는지에 적힌 것을 거의 모두 잃는다.}
\label{fig:cells}
\end{figure}

스파이크당 비트 수는 시간 정밀도에 로그로만 의존한다. 정밀도가 $10\ms$이면 스파이크당 약 5비트가 남는다. 그러나 수신자가 1초 동안의 스파이크를 세기만 한다면, $r=10$에서 1초 전체에 대해 2비트도 얻지 못한다(\ref{sec:rate}절).

\begin{keyblock}
\begin{proposition}[스파이크 채널의 용량]
\label{prop:capacity}
평균 발화율 $r$인 스파이크 흐름을 시간 정밀도 $\Delta t$로 읽으면 \eqref{eq:spikeentropy}보다 많이 나르지 못한다. 이 용량의 거의 전부는 스파이크의 \emph{시각}에 들어 있다. 스파이크를 세기만 하는 수신자는 같은 흐름에서 시각을 1밀리초 정밀도로 구별하는 수신자보다 수십 배 적게 뽑아낸다.
\end{proposition}
\end{keyblock}

이것은 상한이지 측정이 아니다. 자극을 알고 있는 감각 뉴런에서 측정하면 스파이크당 1비트에서 몇 비트가 나온다. 가장 좋은 경우 그 정밀도에 대해 계산한 한계의 절반까지 이른다~\cite{Rieke1997}. 따라서 적어도 신경계의 입구에서는 스파이크 시각이 버려지지 않고 쓰인다.

\section{잡음 있는 채널}

용량~\eqref{eq:spikeentropy}은 잡음 없이 계산했다. 잡음이 어디서 생기는지에 대해서는 세 가지 사실이 확실히 알려져 있다.

\emph{스파이크 발생기 자체는 정확하다.} 조직 조각과 함께 뇌에서 꺼낸 피질 뉴런에 시간에 따라 빠르게 변하는 같은 전류를 여러 번 주입하면, 스파이크는 약 1밀리초의 정밀도로 반복된다~\cite{Mainen1995}. 전류가 일정하면 스파이크 시각은 시행마다 흩어진다. 정확한 시각을 만드는 것은 뉴런 자체가 아니라 입력의 빠른 변화이다.

\emph{시냅스는 믿을 수 없다.} \nt{GLUT} 시냅스에 도착한 스파이크는 어떤 확률 $p$로만 신경전달물질 한 꾸러미를 내보낸다. 해마 시냅스에서는 스파이크의 절반 넘게가 수신자에게 그냥 전달되지 않으며, 원인은 전도 실패가 아니라 바로 방출의 무작위성이다~\cite{AllenStevens1994}. 통신 엔지니어에게 이것은 소거 채널이다. 두 뉴런 사이의 여러 시냅스가 사정을 어느 정도 구해 준다.

\emph{살아 있는 피질에서 스파이크 흐름은 무작위로 보인다.} 스파이크 간격은 무작위 푸아송 흐름과 거의 같게 흩어져 있고, 같은 자극을 되풀이하면 창 안의 스파이크 수가 분산이 평균에 가깝게 변한다~\cite{Softky1993}.

정확한 소자가 어떻게 무작위 흐름을 낳는가? 뉴런은 입력 수천 개를 받고, 하나하나는 믿을 수 없으며, 흥분과 억제는 \ref{sec:gaba}장에서처럼 균형을 이룬다. 막전압은 문턱값 근처에서 흔들리고, 스파이크 시각은 이 흔들림이 정한다. 이것이 잡음인지 우리가 읽지 못하는 메시지인지는 대체로 열린 문제이다. `잡음'의 일부는 이미 메시지로 밝혀졌다. 시각 피질에서 변동의 상당 부분은 동물 자신의 움직임을 따라간다~\cite{Stringer2019}. 여기에는 원리적인 어려움이 있다. 잘 압축된 메시지는 부호를 모르는 사람에게 무작위 흐름과 구별되지 않는다.

\section{발화율: 느리지만 믿을 만하다}
\label{sec:rate}

가장 단순한 부호는 \emph{발화율 부호}이다. 수는 뉴런이 창 $T$ 동안 내는 스파이크의 개수에 적힌다. 소거도 시각의 흔들림도 이런 부호에는 두렵지 않다. 그러나 대가가 있다. 흐름이 푸아송이면 창 안의 스파이크 수 $n$은 평균 $rT$, 산포 $\sqrt{rT}$를 가진다.
\begin{empheq}[box=\keybox]{equation}
\frac{\sigma_n}{\bar n} \;=\; \frac{1}{\sqrt{r\,T}} .
\label{eq:rateerror}
\end{empheq}
발화율을 10퍼센트 정밀도로 알려면 스파이크 백 개가 필요하다. 초당 스파이크 20개라면 5초이다. 이웃한 수준은 산포 두 배쯤 떨어져 있어야 구별되므로, 발화율 $r$까지 시간 $T$ 동안 약 $\sqrt{rT}$개의 수준이 구별된다. $r=10$, $T=1\,\text{s}$이면 서너 수준, 2비트 미만이다.
뇌는 그렇게 느리게 작동하지 않는다. 사람은 한순간 보여 준 그림에서 동물을 알아보며, 뇌 속의 응답은 약 $150\ms$ 뒤에 나타난다~\cite{Thorpe1996}. 거칠게 추정하면 그동안 신호는 연속된 처리 단계 열 개쯤을 한 단계에 $10$--$20\ms$씩 지나간다. 피질의 보통 발화율에서 각 뉴런은 한 단계 동안 스파이크를 0개나 1개 낼 시간밖에 없고, 이런 창에서 그 발화율은 정의되지 않는다. 출구는 두 가지이다. 뉴런을 많이 쓰거나, 시각을 보는 것이다.

\emph{많은 뉴런.} 뉴런 $N$개가 같은 수를 나르고 수신자가 그 스파이크를 더한다고 하자. 잡음이 독립이면 \eqref{eq:rateerror}에서 $rT$ 대신 $NrT$가 들어간다. $r=20$인 뉴런 백 개는 $15\ms$ 동안 스파이크 서른 개를 주어 약 20퍼센트의 정밀도를 낸다. 공간이 시간을 대신한다. 그러나 이웃한 피질 뉴런의 잡음은 독립이 아니다. 이웃한 쌍의 스파이크 수 상관 계수는 보통 약 $0.1$이다~\cite{Zohary1994}. 상관 계수가 $c$이면 뉴런 $N$개에 대한 평균의 분산은 다음과 같다.
\begin{empheq}[box=\keybox]{equation}
\sigma^2_N \;=\; \sigma^2_1\,\frac{1+(N-1)\,c}{N}
\;\xrightarrow[N\to\infty]{}\; c\,\sigma^2_1 .
\label{eq:correlated}
\end{empheq}

\begin{keyblock}
\begin{proposition}[평균의 한계]
\label{prop:pooling}
잡음이 상관되어 있으면 무리에 대한 평균은 오차를 기껏해야 $1/\sqrt{c}$배 줄인다. $c=0.1$이면 세 배이며, 이 이득은 이미 뉴런 수십 개에서 달성된다. 더 보태 봐야 소용없다.
\end{proposition}
\end{keyblock}

모든 상관이 똑같이 해로운 것은 아니다. 정밀도를 제한하는 것은 공통 잡음 가운데 \emph{신호 자체를 닮은} 부분, 즉 측정하는 양이 변했을 때와 같은 방식으로 무리 전체를 미는 부분뿐이다~\cite{MorenoBote2014}. 뇌에 이런 잡음이 실제로 얼마나 있는지는 아직 밝히는 중이다.

많은 뉴런을 쓰는 다른 방법도 있다. 소리 방향 검출기(그림~\ref{fig:itd})처럼 수를 \emph{어느} 뉴런이 발화했는지에 적는 것이다. 이것이 \emph{위치} 부호이며, 알려진 것 중 가장 믿을 만하다. 감각 뉴런에서 전선의 번호는 피부의 어느 점이 닿았는지, 소리의 높이가 얼마인지를 말해 주며, 이는 여러 번 확인되었다. 뉴런 수로 보면 비싸지만 빠르다. 메시지 하나에 지연 하나가 든다.

\section{시각: 빠르지만, 어디서부터 재는가}

두 번째 출구는 수를 스파이크 시각에 적는 것이다. 뉴런~\eqref{eq:glutneuron}에 시각 $t=0$부터 전압을 수준 $U$까지 올릴 일정한 입력을 가하자. 전압은 $V=U\bigl(1-e^{-t/\tau}\bigr)$로 커지고 다음 시각에 문턱값 $\theta$에 이른다.
\begin{empheq}[box=\keybox]{equation}
t_1 \;=\; \tau\,\ln\frac{U}{U-\theta}, \qquad U>\theta .
\label{eq:latency}
\end{empheq}
강한 입력은 이른 스파이크, 약한 입력은 늦은 스파이크, 문턱 아래 입력은 스파이크 없음이다. $\tau=10\ms$에서 입력 $U=4\theta$는 $2.9\ms$ 뒤에, $U=2\theta$는 $6.9\ms$ 뒤에, $U=1.2\theta$는 $17.9\ms$ 뒤에 첫 스파이크를 준다. 스파이크 단 하나가 수를 나른다.

그런데 $t_1$을 어느 순간부터 재는가? 공통 시계는 없으며, 수신자는 자극이 언제 시작되었는지 모른다. 답은 세 가지이다.

\emph{상대 지연.} 두 뉴런이 같은 자극을 보면, 두 지연의 차 $t_1^A-t_1^B$는 시작 시각이 아니라 입력에만 의존한다. 시각은 지연이 있는 동시성 검출기가 비교한다(그림~\ref{fig:itd}). 뉴런 $N$개가 스파이크를 하나씩 냈다면, 그 도착 \emph{순서}만으로 최대 $\log_2N!$비트를 나른다. 뉴런 열 개면 약 22비트인데, `발화했는가 아닌가'라는 답은 10비트를 넘지 못한다. 망막에서는 출력 뉴런들의 첫 스파이크 상대 지연으로 그림을 빠르고 믿을 만하게 복원할 수 있음이 밝혀졌다~\cite{Gollisch2008}.

\emph{낱말의 시작을 알리는 일제 발화.} 자극이 급격히 바뀌면 수많은 뉴런이 거의 동시에 발화한다. 이런 일제 발화는 모스 부호에서 낱말 사이의 긴 쉼처럼 그 자체로 시작 표시가 되고, 수신자는 거기서부터 지연을 잰다.

\emph{국소 리듬.} \ref{sec:gaba}장에서 \nt{GLUT}--\nt{GABA} 루프는 주기 $12$--$50\ms$의 국소 리듬을 낳았다. 이것은 클록이 아니지만, 그 한 주기 안에서 입력이 강한 뉴런은 입력이 약한 뉴런보다 먼저 발화하며, \eqref{eq:latency}는 \emph{위상} 부호가 된다. 수는 스파이크가 주기의 어느 부분에 왔는지에 적힌다. 이런 부호는 관찰되었다. 공간의 특정 장소를 맡은 뉴런은 동물이 `자기' 장소를 지나가면서 느린 국소 리듬의 주기 안에서 점점 더 일찍 발화한다~\cite{OKeefe1993}. 뇌가 위상을 얼마나 널리 쓰는지는 아직 모른다.

\section{부호는 수신자가 고른다}

스파이크 열에는 발화율도, 시각도, 순서도 있다. 그중 무엇을 읽을지는 수신자가 정하며, 파라미터 하나, 즉 자신의 시간 상수 $\tau$로 정한다.

방정식~\eqref{eq:glutneuron}를 시간에 대해 평균하자. 발화율 $r_j$인 독립 흐름들이 뉴런에 들어오면 평균 전압과 그 상대적 흔들림은 다음과 같다.
\begin{empheq}[box=\keybox]{equation}
\bar V \;=\; \tau\sum_j w_j\,r_j ,
\qquad
\frac{\sigma_V}{\bar V} \;\approx\; \frac{1}{\sqrt{2\,r_{\text{in}}\,\tau}} ,
\label{eq:reader}
\end{empheq}
여기서 두 번째 식은 가중치가 같을 때의 것이고, $r_{\text{in}}$은 모든 입력의 발화율 합이다. $\tau$가 크면 전압은 거의 흔들리지 않고 발화율의 가중합을 따라간다. 수신자는 \emph{적분기}이며, 발화율을 읽고 시각은 잊는다. $\tau$가 작으면 전압이 스파이크 사이에 떨어질 시간이 있어서, 여러 스파이크가 창~\eqref{eq:window} 안에 들어왔을 때만 문턱값에 이른다. 수신자는 \emph{동시성 검출기}이며, 시각을 읽는다(그림~\ref{fig:reader}). 초당 스파이크 5개인 입력 천 개는 $\tau=10\ms$에서 약 10퍼센트의 흔들림을 주며, 거의 순수한 적분이다. \ref{sec:gaba}장에서처럼 억제가 $\tau$를 $2\ms$로 줄이면 흔들림은 20퍼센트 남짓으로 커지고, 뉴런은 일치를 알아채기 시작한다.

\begin{figure}[t]
\centering
\begin{tikzpicture}[>=Stealth,x=0.115cm,y=1cm]
% common input: the same spike train for both receivers
\foreach \s in {6,21,22,38,52,55.5,59,62.5,66,69.5,73,76.5,80,83.5,87,90.5,94,97.5}{\draw[thick,gray!40!black] (\s,5.25) -- (\s,5.65);}
\draw[gray!70!black] (0,5.25) -- (105,5.25);
\node[left,font=\small] at (-1,5.45) {입력};
\node[above,font=\scriptsize,gray!40!black] at (13.5,5.65) {드묾};
\node[above,font=\scriptsize,gray!40!black] at (21.5,6.0) {쌍};
\node[above,font=\scriptsize,gray!40!black] at (75,5.65) {잦음};
% axes and thresholds
\foreach \y/\th in {2.7/2.5,0/1.5}{
  \draw[gray!70!black] (0,\y) -- (106,\y) node[right,black] {\small $t$};
  \draw[densely dashed,gray!60!black] (0,\y+0.6*\th) -- (104,\y+0.6*\th) node[right,black,font=\small] {$\theta$};
}
\node[anchor=south west,font=\small] at (0,4.3) {$\tau=10\ms$: 적분기};
\node[anchor=south east,font=\small] at (104,1.0) {$\tau=2\ms$: 동시성 검출기};
% integrator: tau=10 ms, theta=2.5w
\begin{scope}[yshift=2.7cm]
\foreach \a/\b/\v in {0/6/0.000,6/21/1.000,21/22/1.223,22/38/2.107,38/52/1.425,52/55.5/1.351,55.5/59/1.952,59/62.5/2.376,62.5/66/0.000,66/69.5/1.000,69.5/73/1.705,73/76.5/2.201,76.5/80/0.000,80/83.5/1.000,83.5/87/1.705,87/90.5/2.201,90.5/94/0.000,94/97.5/1.000,97.5/104/1.705}{\draw[very thick,blue!60!black] plot[domain=\a:\b,samples=14] (\x,{0.6*\v*exp(-(\x-\a)/10)});}
\foreach \s/\p/\q in {6/0.000/1.000,21/0.223/1.223,22/1.107/2.107,38/0.425/1.425,52/0.351/1.351,55.5/0.952/1.952,59/1.376/2.376,62.5/1.674/2.500,62.5/2.500/0.000,66/0.000/1.000,69.5/0.705/1.705,73/1.201/2.201,76.5/1.551/2.500,76.5/2.500/0.000,80/0.000/1.000,83.5/0.705/1.705,87/1.201/2.201,90.5/1.551/2.500,90.5/2.500/0.000,94/0.000/1.000,97.5/0.705/1.705}{\draw[very thick,blue!60!black] (\s,0.6*\p) -- (\s,0.6*\q);}
\foreach \s in {62.5,76.5,90.5}{\draw[->,thick,red!70!black] (\s,1.58) -- (\s,1.95);}
\end{scope}
% coincidence: tau=2 ms, theta=1.5w
\begin{scope}[yshift=0cm]
\foreach \a/\b/\v in {0/6/0.000,6/21/1.000,21/22/1.001,22/38/0.000,38/52/1.000,52/55.5/1.001,55.5/59/1.174,59/62.5/1.204,62.5/66/1.209,66/69.5/1.210,69.5/73/1.210,73/76.5/1.210,76.5/80/1.210,80/83.5/1.210,83.5/87/1.210,87/90.5/1.210,90.5/94/1.210,94/97.5/1.210,97.5/104/1.210}{\draw[very thick,blue!60!black] plot[domain=\a:\b,samples=14] (\x,{0.6*\v*exp(-(\x-\a)/2)});}
\foreach \s/\p/\q in {6/0.000/1.000,21/0.001/1.001,22/0.607/1.500,22/1.500/0.000,38/0.000/1.000,52/0.001/1.001,55.5/0.174/1.174,59/0.204/1.204,62.5/0.209/1.209,66/0.210/1.210,69.5/0.210/1.210,73/0.210/1.210,76.5/0.210/1.210,80/0.210/1.210,83.5/0.210/1.210,87/0.210/1.210,90.5/0.210/1.210,94/0.210/1.210,97.5/0.210/1.210}{\draw[very thick,blue!60!black] (\s,0.6*\p) -- (\s,0.6*\q);}
\foreach \s in {22}{\draw[->,thick,red!70!black] (\s,0.98) -- (\s,1.35);}
\end{scope}
\foreach \x in {0,50,100}{\draw[gray!60!black] (\x,-0.06) -- (\x,0.06) node[below,font=\scriptsize] {\x\,ms};}
\end{tikzpicture}
\caption{같은 입력, 서로 다른 두 수신자. 스파이크마다 전압이 $w$만큼 올랐다가 뉴런 모델~\eqref{eq:glutneuron}에서처럼 흘러내린다. 빨간 화살표는 수신자의 스파이크이다. 느린 수신자($\tau=10\ms$, $\theta=2.5w$)는 스파이크가 \emph{많은} 곳에서 발화하고 쌍은 알아채지 못한다. 빠른 수신자($\tau=2\ms$, $\theta=1.5w$)는 창~\eqref{eq:window} 안의 쌍에만 발화하고, 잦지만 일치하지 않는 흐름은 흘려보낸다.}
\label{fig:reader}
\end{figure}

측정에 따르면 활성 피질에서 막 컨덕턴스는 휴지 상태보다 몇 배 커진다~\cite{Destexhe2003}. 그러니 거기서는 $\tau$가 더 짧고, 작동 중인 피질 뉴런은 적분기보다 동시성 검출기에 더 가깝다.

\begin{keyblock}
\begin{proposition}[부호는 수신자가 정한다]
\label{prop:reader}
같은 스파이크 열이 느린 수신자에게는 발화율을, 빠른 수신자에게는 시각을, 신경전달물질이 방출되는 조직 공간에게는 몇 초에 걸쳐 평활된 농도 수준을 나른다(\ref{sec:serocomputer}장). 어떤 부호가 쓰이는지는 송신자가 아니라 수신자의 시간 상수가 정하며, 그 시간 상수는 억제가 조정한다.
\end{proposition}
\end{keyblock}
\section{점과 대시: 필터로서의 시냅스}

지금까지 시냅스는 가중치가 달린 전선이었다. 실제로는 시냅스의 응답이 앞선 스파이크에 의존하며, 이것이 뉴런의 언어에 두 번째 기호를 준다.

\emph{고갈: 시냅스는 변화를 전달한다.} 방출할 때마다 방출 준비가 된 신경전달물질 저장량의 몫 $U$가 쓰이고~\cite{Tsodyks1997}, 저장량은 수백 밀리초 정도의 시간 상수 $\tau_{\text{rec}}$로 회복된다. 발화율이 $r$이면 스파이크당 응답 진폭은 대략 $A(r)=A_0/(1+U r\tau_{\text{rec}})$ 수준에 자리 잡으며, 여기서 $A_0$는 쉬고 난 시냅스의 응답이다. 수신자가 단위 시간에 받는 전류는 다음과 같다.
\begin{empheq}[box=\keybox]{equation}
r\,A(r) \;=\; \frac{A_0\,r}{1+U r\,\tau_{\text{rec}}}
\;\xrightarrow[r\to\infty]{}\; \frac{A_0}{U\tau_{\text{rec}}} .
\label{eq:depression}
\end{empheq}
$U=0.5$, $\tau_{\text{rec}}=0.8\,\text{s}$이면 포화는 이미 초당 스파이크 $1/(U\tau_{\text{rec}})=2.5$개에서 시작된다. 발화율이 5에서 20으로 오르면 정상 상태 전류는 겨우 3분의 1 늘어난다. 대신 새 발화율의 첫 스파이크들은 저장량이 아직 예전 그대로일 때 도착하므로 전류가 한순간 네 배로 뛰었다가 $\tau_{\text{rec}}/(1+Ur\tau_{\text{rec}})\approx90\ms$ 만에 줄어든다(그림~\ref{fig:depression}).

이런 시냅스는 발화율이 아니라 그 \emph{변화}, 본질적으로 상대 변화 $\Delta r/r$를 전달한다~\cite{Abbott1997depression}. 전자공학자에게 이것은 전선 안에 바로 넣은 고역 통과 필터이다.

\begin{figure}[t]
\centering
\begin{tikzpicture}[>=Stealth,x=12.5cm,y=1cm]
% input rate
\draw[->,gray!70!black] (0,2.6) -- (0.9,2.6) node[right,black] {\small $t$};
\draw[->,gray!70!black] (0,2.6) -- (0,4.3) node[above,black] {\small $r$};
\draw[very thick,blue!60!black] (0,2.9) -- (0.2,2.9) -- (0.2,3.8) -- (0.8,3.8);
\node[left,font=\scriptsize] at (0,2.9) {$5$};
\node[left,font=\scriptsize] at (0,3.8) {$20$};
\node[right,font=\small,blue!60!black] at (0.4,4.1) {시냅스 입력의 발화율};
% output current
\draw[->,gray!70!black] (0,0) -- (0.9,0) node[right,black] {\small $t$};
\draw[->,gray!70!black] (0,0) -- (0,2.45) node[left,black] {\small $rA$};
\draw[densely dashed,gray!60!black] (0,0.667) -- (0.8,0.667);
\draw[very thick,red!70!black] (0,0.5) -- (0.2,0.5) -- (0.2,2.0)
   plot[domain=0.2:0.8,samples=60] (\x,{0.667+(2.0-0.667)*exp(-(\x-0.2)/0.089)});
\node[left,font=\scriptsize] at (0,0.5) {$1.7$};
\node[left,font=\scriptsize] at (0,2.0) {$6.7$};
\node[right,font=\scriptsize] at (0.8,0.667) {$2.2$};
\node[right,font=\small,red!70!black] at (0.28,1.6) {수신자로 가는 전류: 도약 뒤 $\approx90\ms$ 만에 감소};
\draw[gray!60!black] (0.2,-0.08) -- (0.2,0.08); \node[below,font=\scriptsize] at (0.2,0) {$0.2\,\text{s}$};
\draw[gray!60!black] (0.8,-0.08) -- (0.8,0.08); \node[below,font=\scriptsize] at (0.8,0) {$0.8\,\text{s}$};
\end{tikzpicture}
\caption{$U=0.5$, $\tau_{\text{rec}}=0.8\,\text{s}$일 때 식~\eqref{eq:depression}에 따른 고갈 시냅스. 전류는 초당 $A_0$ 단위이고, 파선은 정상 상태 전류이다. 정상 상태 전류는 $1.7$에서 $2.2$로 늘었을 뿐이지만, 도약하는 순간에는 $6.7$까지 오른다. 시냅스는 수신자에게 발화율이 얼마인지가 아니라 발화율이 바뀌었다는 것을 알린다.}
\label{fig:depression}
\end{figure}

\emph{촉진: 대시로서의 버스트.} 다른 시냅스들은 방출 확률이 낮지만, 스파이크마다 수십 밀리초 동안 그 확률이 커진다. 이런 시냅스를 지나는 단일 스파이크는 대개 사라진다. 그러나 \emph{버스트}, 즉 몇 밀리초 간격으로 연이어 나오는 스파이크 몇 개는 거의 확실히 통과한다. 촉진이 없어도 버스트의 스파이크 $k$개가 모두 사라질 확률은
\begin{empheq}[box=\keybox]{equation}
P_{\text{loss}} \;=\; (1-p)^k ,
\label{eq:burst}
\end{empheq}
이며, $p=0.2$이면 단일 스파이크에서 $0.8$, 네 개짜리 버스트에서 $0.41$이다. 촉진은 이 확률을 더 줄인다. 이렇게 뉴런에게 두 번째 기호가 생긴다. 단일 스파이크는 점, 버스트는 대시이다. 고갈 시냅스는 쉼 뒤의 첫 스파이크를 가장 잘 전달하고, 촉진 시냅스는 버스트를 통과시키고 홀로 된 점은 죽인다. 버스트가 더 믿을 만하게 전달된다는 것은 측정되었다. 뇌가 실제로 버스트를 별도의 기호로 쓴다는 것은 그럴듯한 가설이다~\cite{Lisman1997}.

\section{\nt{GABA} 뉴런은 어떻게 말하는가}

피질에서 가장 수가 많은 것은 빠른 \nt{GABA} 뉴런이다~\cite{Tremblay2016}. 이들의 스파이크는 \nt{GLUT} 뉴런보다 짧고, 초당 수백 개의 발화율로 고르게, 오랫동안, 거의 지치지 않고 스파이크를 낼 수 있다. 이들의 시냅스는 더 믿을 만하다. 각 수신자와의 연결이 많은 방출 자리로 이루어져 있어 스파이크가 사라지는 일이 거의 없다. 이들의 출력은 수신자에서 스파이크가 생겨나는 자리 가까이에 붙으므로, 수신자가 입력을 어떻게 더하는지가 아니라 발화할지, 언제 발화할지를 제어한다.

이들 자신은 이웃한 많은 \nt{GLUT} 뉴런에게서 흥분을 받고 주변의 총활동을 알린다. \ref{sec:gaba}장의 나눗셈에 딱 필요한 것이다. 끝으로 빠른 \nt{GABA} 뉴런은 서로 연결되어 쉽게 함께 발화한다. 앞에서 말한 국소 리듬을 정하는 것이 바로 이들이다~\cite{Cardin2009}.

모스 부호의 언어로 말하면 \nt{GABA} 뉴런은 글자가 아니라 \emph{쉼}을 맡는다. 이들은 \nt{GLUT} 스파이크의 끊임없는 흐름을 수신자가 발화할 수 있는 창들로 자르고, 일치 창을 좁히며, 흐름이 너무 시끄러워지면 전체를 낮춘다.

\begin{keyblock}
\begin{proposition}[역할 분담]
\label{prop:punctuation}
\nt{GLUT} 뉴런은 메시지의 내용, 즉 \emph{무엇}과 \emph{얼마나}를 나른다. 빠른 \nt{GABA} 뉴런은 메시지의 표지, 즉 \emph{언제} 말해도 되는지와 \emph{얼마나 크게}를 나른다. 또 다른 부류는 억제하는 뉴런을 억제하여 관문이 \emph{누구에게} 열리는지를 정한다. 이 분담의 개별 부분은 측정되었다. 일반 규칙으로서는 작업 가설이다.
\end{proposition}
\end{keyblock}

더 느린 \nt{GABA} 뉴런들도 있는데, 그 출력은 수신자의 개별 입력에 붙는다. 이들은 \ref{sec:gaba}장의 금지~\eqref{eq:veto}처럼 작동한다. \nt{GLUT} 스파이크의 흐름 하나를 닫고 나머지는 건드리지 않는다.

\nt{GABA} 뉴런을 빠른 것과 느린 것으로 나누는 것은 오래된 그림이며, 완전하지 않다. 지금은 피질에서 적어도 세 개의 큰 부류를 구별하며~\cite{Tremblay2016}, 세 번째 부류는 뜻밖의 구조를 가진다. 이 부류는 \nt{GLUT} 뉴런이 아니라 다른 \nt{GABA} 뉴런, 무엇보다 입력을 닫는 뉴런들을 억제한다~\cite{Pfeffer2013}. 억제의 억제는 이중 부정이다. 이런 뉴런은 수신자에게 흥분성 스파이크를 하나도 보내지 않고 관문을 \emph{연다}. 이 뉴런을 켜는 것은 피질의 다른 영역에서 오는 신호와 브로드캐스트 채널이므로, 이 뉴런은 네트워크 한 구역 전체의 허가 입력처럼 작동한다. 이것이 활성인 동안 금지가 풀리고 흐름이 통과한다~\cite{Pi2013}. 부록~\ref{sec:othermediators}에서는 이 기법을 탈억제라고 부른다.
\section{절약하는 언어}

뉴런의 언어에 대해 확실히 알려진 마지막 사실은 그것이 비싸다는 것이다. 스파이크 하나와 그것이 일으키는 시냅스의 일에는 ATP 분자가 약 $10^9$개 들고(설치류 피질에 대한 추정~\cite{Attwell2001}), 분자 하나는 약 $10^{-19}$줄을 낸다. 따라서 스파이크 하나의 값은 $E_{\text{spk}}\sim10^{-10}$줄 정도이다. 뇌 전체는 약 $20$와트를 쓰며, 그 절반이 신호 전달에 쓰인다면, 즉 $P\sim10$와트이면 뉴런 하나의 평균 발화율은 다음과 같다.
\begin{empheq}[box=\keybox]{equation}
\bar r \;\approx\; \frac{P}{N\,E_{\text{spk}}}
\;\sim\; \frac{10\ \text{W}}{8.6\cdot10^{10}\cdot10^{-10}\ \text{J}}
\;\approx\; 1\ \text{spike/s}.
\label{eq:energy}
\end{empheq}
피질에 대한 더 자세한 추정은 이보다도 적은 값을 준다~\cite{Lennie2003}. 뇌의 부호는 \emph{희소}할 수밖에 없다. 뉴런 가운데 작은 몫만이 빠르게 작동할 수 있다.

\eqref{eq:spikeentropy}를 보면 이것이 그리 나쁘지 않다. 정밀도가 $1\ms$일 때, 초당 $100$번 발화하는 뉴런의 스파이크는 5비트를 넘지 못하지만, 초당 한 번 발화하는 뉴런의 스파이크는 11비트까지 나른다. 스파이크마다 값을 치러야 한다면 드물게 말하는 편이 이득이다. 피질은 실제로 정밀도를 에너지와 맞바꾼다. 먹이가 부족하면 뉴런은 스파이크 발화율은 유지하되 시냅스 전류에 덜 쓰고, 그 응답은 덜 정밀해진다~\cite{Padamsey2022}.

모스 부호에서 자주 쓰는 글자는 짧다. 영어 글에서 가장 흔한 글자 E는 점 하나이다. 알려진 바로는 뉴런의 부호도 다른 형태로 같은 규칙을 따른다. 예상된 것에는 스파이크를 적게 쓴다~\cite{Barlow1961}. 자극이 일정하면 뉴런은 발화율을 점점 낮추고, \ref{sec:glutcomputer}장의 센서는 수준이 아니라 변화에 응답하며, \ref{sec:neurontypes}장의 \nt{DOPA} 뉴런은 보상 예측 오차만을 알린다.

\begin{keyblock}
\begin{proposition}[절약]
\label{prop:economy}
에너지는 뉴런의 평균 발화율을 초당 스파이크 1개 정도로 제한한다. 그래서 \nt{GLUT} 뉴런의 언어는 희소하며, 스파이크를 새 소식, 즉 변한 것이나 예측되지 않은 것에 쓴다. 에너지 한계는 측정되었다. 뇌의 부호가 줄당 비트 수를 최대로 하도록 맞추어져 있다는 것은 근거가 탄탄한 가설이다.
\end{proposition}
\end{keyblock}

\section{요약}

\begin{table}[t]
\centering
\footnotesize
\setlength{\tabcolsep}{4pt}
\renewcommand{\arraystretch}{1.15}
\begin{tabular}{>{\raggedright}p{2.6cm}>{\raggedright}p{2.3cm}>{\raggedright}p{3.0cm}>{\raggedright}p{2.0cm}>{\raggedright\arraybackslash}p{3.6cm}}
\toprule
적는 방식 & 나르는 것 & 읽는 쪽 & 메시지당 시간 & 알려진 것 \\
\midrule
발화한 뉴런의 번호(위치 부호) & 어떤 값인가 & 모든 수신자; 동시성 검출기 & 지연 하나 & 감각 뉴런과 소리 방향 검출에서 확립됨 \\
뉴런 하나의 발화율 & 크기 & $\tau$가 큰 적분기; 조직 공간(\ref{sec:serocomputer}장) & 수백 ms--수 초 & 확립됨; $10$--$20\ms$의 처리 단계에는 너무 느림 \\
무리의 발화율 & 크기 & 입력이 수천 개인 뉴런 & $10$--$50\ms$ & 확립됨; 이득은 상관~\eqref{eq:correlated}으로 제한됨 \\
첫 스파이크 시각, 순서 & 크기, 순서 & 지연이 있는 동시성 검출기 & 지연 하나 & 신경계 입구에서 확립됨; 피질에서는 논의 중 \\
국소 리듬 속의 위상 & 크기, 위치 & 공통 리듬을 가진 뉴런 & 주기 $12$--$50\ms$ 및 더 느림 & 일부 영역에서 관찰됨; 일반적 역할은 가설 \\
버스트(`대시') & `중요함': 믿을 만한 전달 & 촉진 시냅스 & $\sim10\ms$ & 신뢰성은 측정됨; 별도 기호라는 것은 가설 \\
발화율의 변화 & 새 소식 & 고갈 시냅스 & $\sim0.1\,\text{s}$ & 확립됨 \\
빠른 \nt{GABA} 뉴런의 스파이크 & 언제, 얼마나 크게 & 이웃한 모든 \nt{GLUT} 뉴런 & $1$--$30\ms$ & 리듬과 나눗셈은 측정됨; 규칙으로서의 `구두점'은 가설 \\
\bottomrule
\end{tabular}
\caption{\nt{GLUT} 뉴런과 \nt{GABA} 뉴런이 스파이크로 메시지를 적는 방식. 오른쪽 열은 측정된 것과 추정된 것을 구분한다.}
\label{tab:code}
\end{table}

표~\ref{tab:code}는 우리가 아는 것을 모은 것이며, 우리가 모르는 것도 거기서 보인다. 피질이 무리의 발화율뿐 아니라 정확한 스파이크 시각을 얼마나 쓰는지 우리는 모른다. 뉴런에게 `낱말', 즉 전체로 읽히는 여러 뉴런 스파이크의 안정된 조합이 있는지 모른다. 그리고 뇌의 여러 부분에서 언어가 같은지도 모른다. 모스 부호는 그 표가 알려져 있기에 하룻저녁이면 익힐 수 있다. 뉴런 언어의 표는 아직 만들어야 하며, 그것을 만드는 사람은 아마 통신 엔지니어일 것이다.

\section{연습문제}

\begin{exercise}[\lvlA]
\label{ex:04:1}
\emph{숫자로 본 용량.} $\Delta t=1\ms$. (a)~스파이크 값이 \eqref{eq:energy}와 같을 때, 발화율이 초당 스파이크 $1$개인 뉴런은 줄당 몇 비트를 주는가? (b)~$r=10$에서 1초 동안의 스파이크를 세기만 하는 수신자는 한계~\eqref{eq:spikeentropy}보다 몇 배 적게 뽑아내는가?
\end{exercise}

\begin{exercise}[\lvlB]
\label{ex:04:2}
\emph{최적 발화율.} 뉴런은 전력 $P_0+rE_{\text{spk}}$를 쓴다. 여기서 $P_0$는 $20$~W의 나머지 절반을 모든 뉴런에 나눈 값이고, $E_{\text{spk}}=10^{-10}$~J이다. $\Delta t=1\ms$인 \eqref{eq:spikeentropy}에 따른 줄당 비트 수 $R_{\max}(r)/(P_0+rE_{\text{spk}})$는 어떤 발화율 $r$에서 가장 큰가?
\end{exercise}

\begin{exercise}[\lvlB]
\label{ex:04:3}
\emph{어떤 상관이 해로운가.} 뉴런 $N=1000$개, 분산 $\sigma^2=1$, 쌍별 상관 $c=0.1$. 기울기가 모두 같은 $f'=\mathbf 1$일 때와, $\mathbf 1^{\mathsf T}f'=0$인 기울기 $\pm1$일 때의 피셔 정보 $\mathcal I=f'^{\mathsf T}\Sigma^{-1}f'$($\Sigma$는 공분산 행렬)을 계산하라. 둘은 몇 배 다른가?
\end{exercise}

\begin{exercise}[\lvlB]
\label{ex:04:4}
\emph{모델링: 동시성 검출기.} 뉴런~\eqref{eq:glutneuron}가 초당 스파이크 $5$개인 푸아송 입력 $1000$개를 받고 $w=1$이다. 문턱값 $\theta$는 자유 전압이 1초에 한 번 넘도록 잡는다. $\tau=10$과 $2\ms$에서 동시 입력 몇 개 $m$이 확률 $1/2$로 뉴런을 발화시키는지 모델로 구하라.
\end{exercise}

\begin{exercise}[\lvlB]
\label{ex:04:5}
\emph{설계: 상대 변화 검출기.} 시냅스~\eqref{eq:depression}를 고르라. 초당 스파이크 $40$개에서의 정상 상태 전류가 $10$개에서보다 $10\,\%$ 넘게 높지 않아야 하고, $40$으로 뛰어오른 뒤 전류가 시간 상수 $50\ms$ 이하로 새 수준에 이르러야 한다. 가장 작은 $\tau_{\text{rec}}$는 얼마이고, 그때 $U$는 얼마인가?
\end{exercise}

\begin{exercise}[\lvlA]
\label{ex:04:6}
\emph{첫 스파이크 시각과 버스트.} (a)~\eqref{eq:latency}에서 첫 스파이크가 정확히 시간 상수 하나 뒤에 오게 하는 입력을 구하라. 이 입력은 무엇에 의존하는가? (b)~방출 확률이 $p=0.2$일 때, 모두 잃을 확률이 $5\,\%$ 아래가 되려면 버스트에 스파이크가 몇 개 있어야 하는가?
\end{exercise}

\begin{exercise}[\lvlC]
\label{ex:04:7}
\emph{탐구: 배치에는 몇 비트가 있는가.} 뉴런은 독립 칸들~\eqref{eq:spikeentropy}이고, $r=10$, $\Delta t=1\ms$이다. (a)~칸 $20$개로 된 `낱말'의 엔트로피를 스파이크 수의 엔트로피와 나머지로 나누어라. (b)~기록 $N=30$개와 $10^4$개로부터 낱말 빈도로 엔트로피를 추정하는 것을 시뮬레이션하라. 추정치는 얼마나 낮게 나오는가?
\end{exercise}

\chapter{\nt{SERO} 뉴런은 무엇을 계산하는가}
\label{sec:serocomputer}
\begin{chaptergoals}
\item 억제성 수용체를 가진 페이스메이커로 \nt{SERO} 뉴런 하나가 NOT과 NOR가 되는 방식;
\item 체적 방출에서 독립 전선이 $\ell\sim\sqrt{D\tau}$로 정해져 몇 개뿐인 이유;
\item 농도가 스파이크를 수준으로 평활하는 방식과 음의 되먹임이 이를 붙잡을 수 있는 방식;
\item \nt{SERO}가 계획 지평의 후보인 이유와 지평이 인내심을 정하는 방식.
\end{chaptergoals}

\section{첫 번째 브로드캐스트 채널}

지금까지 이 책의 뉴런은 모두 주소 버스로 통신했다. \nt{GLUT}와 \nt{GABA}는 정해진 수신자에게 스파이크를 보냈고, 우리는 그런 네트워크가 무엇을 계산하며 어떤 언어로 말하는지 알아보았다. 이제 \ref{sec:buses}절에서 본 두 번째 버스, 즉 브로드캐스트 버스로 넘어간다. 그 첫 번째 채널은 \nt{SERO}이다. 이번에도 살아 있는 생물에 실제로 있는 것만 쓰기로 한다.

이 장에서는 이후 모든 브로드캐스트 채널이 쓰게 될 세 가지 장치가 등장한다. 고른 배경 입력이 있으면 누가 멈추기 전까지 스스로 작동하는 뉴런, 실제로는 조직의 부피인 전선, 그리고 개별 스파이크 대신 쓰이는 느린 농도다. \ref{sec:dofa}--\ref{sec:acet}장의 \nt{DOPA}, \nt{NORA}, \nt{ACET}도 같은 부품으로 조립된다.

엔지니어의 눈으로 보면 \nt{SERO} 뉴런에는 중요한 차이가 네 가지 있다.

\emph{첫째, 부호는 수신자가 고른다.} 세로토닌에는 두 부호의 수용체가 모두 있으므로 규칙~\eqref{eq:dalesign}은 여기서 통하지 않는다. 부호를 정하는 것은 송신자가 아니라 수신자의 수용체다(명제~\ref{prop:receptor})~\cite{BarnesSharp1999}.

\emph{둘째, \nt{SERO} 뉴런은 고른 배경 입력만 있으면 스스로 작동한다.} 깨어 있을 때 이 뉴런은 1초에 몇 번씩 고르게 스파이크를 내며, 스파이크마다 따로 밀어 줄 필요가 없다. 즉 페이스메이커다. 다만 일정한 배경 입력은 필요하다. 이 입력이 없는 뇌 절편에서는 이런 세포 대부분이 침묵한다. $\alpha_1$ 수용체를 통해 노르아드레날린을 주면 고른 리듬이 돌아오고, 이 수용체를 차단하면 리듬이 꺼진다~\cite{VandermaelenAghajanian1983}. 생체에서는 이 배경을 \nt{NORA}가 공급한다. 회로로 보면 전원이 필요한 발진기다. 전원이 있는 동안에는 억제가 멈출 때까지 스스로 돈다.

\emph{셋째, 느리다.} 세로토닌 수용체는 거의 모두 대사성이다. 스스로 채널을 여는 대신 세포 안에서 연쇄 반응을 일으킨다. 그래서 응답이 느리다. 주 수용체를 통한 억제 전류는 약 1초 안에 오르고 내린다~\cite{CourtneyFord2016}. \nt{GLUT} 빠른 시냅스의 응답보다 수십 배 길다.

\emph{넷째, 신경전달물질을 전선이 아니라 부피 속으로 내보낸다.} \nt{SERO} 뉴런의 출력이 닿는 영역에서 세로토닌은 좁은 틈이 아니라 세포 사이 공간으로 나와 주위로 퍼지는 경우가 많다~\cite{Bunin1998}. 수신자는 농도를 느낄 뿐, 분자가 어느 송신자에게서 왔는지는 알 수 없다.

이런 시간 척도에서는 개별 스파이크가 더 이상 중요하지 않으므로, 네트워크를 발화율 $r_j$와 농도로 기술하겠다. 뉴런 $j$가 신경전달물질을 내보내는 부피 속 세로토닌 농도 $c$는 방출로 늘고, 재흡수로 신경전달물질이 빠져나가면서 줄어든다.
\begin{empheq}[box=\keybox]{equation}
\tau_c\,\frac{\dd c}{\dd t} \;=\; -c \;+\; \kappa\sum_j r_j ,
\qquad
r_i \;=\; f_i\bigl(b_i + s_i\,g\,c\bigr), \quad s_i=\pm1 .
\label{eq:seroneuron}
\end{empheq}
이것은 명제~\ref{prop:reader}를 보완하는, 스파이크열을 읽는 또 하나의 방법이다. 부피는 개별 스파이크도 그 순서도 보지 못하고, 시간 $\tau_c$ 동안 평균한 발화율만 본다. 첫 번째 식은 막과 똑같은 RC 회로다. $\tau_c$는 부피 속 신경전달물질의 수명으로, 직접 측정된 적은 없고 여기서는 1초 정도로 둔다. $\kappa$는 스파이크 하나가 내놓는 신경전달물질의 양이다. 두 번째 식은 수신자를 기술한다. $b_i$는 수신자 자신의 입력으로, 페이스메이커에서는 양수다(여기에는 \nt{NORA}에서 오는 고른 배경 입력도 들어 있다). $g$는 수용체의 감도, 부호 $s_i$는 수신자가 고르며, $f_i$는 감소하지 않는 전달 특성이다.

\section{뉴런 하나로 만드는 NOT}

억제성 수용체를 가진 페이스메이커, $s=-1$을 생각하자. 주위에 세로토닌이 없는 동안에는 자기 입력 $b$만으로 작동한다. 세로토닌이 나타나면 입력이 $gc$만큼 줄고, $gc>b$이면 뉴런은 침묵한다. 이것이 NOT이다. 입력이 없을 때 출력이 있다. 여기에 든 뉴런은 하나뿐이다(그림~\ref{fig:serogates}). 명제~\ref{prop:monotone}은 여기에 적용되지 않는다. 그 명제는 양의 가중치를 전제로 했다.

같은 페이스메이커에 서로 다른 두 부피의 세로토닌이 작용하면, 뉴런은 둘 중 하나에라도 세로토닌이 있으면 침묵한다. 이것이 NOR이며, NOR로는 어떤 논리 함수든 만들 수 있다. \nt{GLUT} 네트워크의 이중 레일 부호는 여기서 필요 없다. 신호가 없다는 사실은 누가 멈추기 전까지 작동하는 뉴런이 계산한다.

\begin{figure}[t]
\centering
\begin{tikzpicture}[>=Stealth,
  nrn/.style={circle,draw,very thick,fill=blue!6,minimum size=0.95cm,inner sep=0pt},
  pace/.style={nrn,double,double distance=1.2pt},
  inh/.style={-{Bar[width=7pt]},thick,red!70!black}]
\node[pace] (N1) at (0,0) {};
\draw[inh] (-1.6,0) node[left,black] {$c$} -- (N1);
\draw[->,very thick,red!70!black] (N1) -- (1.3,0) node[right,black] {$\bar c$};
\node[below=0.55cm] at (0,0) {\small NOT};
\node[pace] (N2) at (5,0) {};
\draw[inh] (3.4,0.45) node[left,black] {$c_1$} -- (N2);
\draw[inh] (3.4,-0.45) node[left,black] {$c_2$} -- (N2);
\draw[->,very thick,red!70!black] (N2) -- (6.3,0) node[right,black] {$\overline{c_1\vee c_2}$};
\node[below=0.55cm] at (5,0) {\small NOR};
\end{tikzpicture}
\caption{\nt{SERO} 뉴런으로 만든 논리. 이중 원은 페이스메이커로, 고른 배경 입력이 있으면 억제될 때까지 스스로 작동한다. 끝에 가로줄이 있는 선은 억제성 수용체($s=-1$)이고, $c$, $c_1$, $c_2$는 수신자 주변 부피의 세로토닌 농도다. 왼쪽은 NOT, 오른쪽은 NOR이다.}
\label{fig:serogates}
\end{figure}

\begin{keyblock}
\begin{proposition}[\nt{SERO} 논리의 완전성]
\label{prop:serocomplete}
네트워크에 억제성 수용체를 가진 페이스메이커가 있고 서로 다른 부피로의 방출이 섞이지 않는다면, 네트워크~\eqref{eq:seroneuron}은 어떤 논리 함수든 계산할 수 있으며, 각 비트는 전선 하나로 전달된다.
\end{proposition}
\end{keyblock}

논리만 놓고 보면 \nt{SERO} 네트워크는 \nt{GLUT}보다 강하다. 그러나 “서로 다른 부피로의 방출이 섞이지 않는다”는 조건은 뇌에서 성립하지 않는다(\ref{sec:volume}절).

\section{음의 되먹임}

세로토닌 뉴런에는 자기 자신에 붙은 억제성 수용체가 있다~\cite{Sprouse1987}. 이들이 내보낸 세로토닌은 다시 자신에게 돌아와 자신을 억누른다. 엔지니어에게는 익숙한 회로, 즉 음의 되먹임 증폭기다(그림~\ref{fig:serofeedback}).

무엇이 측정되었고 무엇이 모델인지 짚어 두자. \nt{SERO} 뉴런이 모여 있는 봉선핵 자체에서는 세로토닌이 공통 부피가 아니라 시냅스를 통해 점대점으로 이웃을 억제한다. 수송체가 세로토닌을 빠르게 치워 틈 밖에 쌓이지 못하게 하기 때문이다. 한 번의 방출은 발화에 짧은 멈춤만 만들 뿐 평균 발화율은 바꾸지 않는다~\cite{CourtneyFord2016}. 따라서 아래에서 공통 부피를 통해 닫히는 고리는 모델이다. 이런 되먹임이 평균 발화율 수준에서 작동한다면 무엇을 할지 계산해 보는 것이다.

\begin{figure}[t]
\centering
\begin{tikzpicture}[>=Stealth,
  nrn/.style={circle,draw,very thick,fill=blue!6,minimum size=0.95cm,inner sep=0pt},
  pace/.style={nrn,double,double distance=1.2pt},
  blk/.style={draw,very thick,rounded corners=2pt,fill=orange!10,minimum height=0.9cm,minimum width=2.2cm,align=center,font=\small},
  inh/.style={-{Bar[width=7pt]},thick,red!70!black}]
\node[pace] (N) at (0,0) {};
\node[blk] (V) at (3.6,0) {부피\\$\tau_c$};
\draw[->,very thick] (N) -- node[above] {\small $r$} (V);
\draw[->,very thick] (V) -- (6.0,0) node[right] {\small $c$: 모든 수신자에게};
\draw[inh] (V.south) -- ++(0,-0.9) -| node[pos=0.25,below] {\small $-g\,c$} (N.south);
\draw[->,thick,blue!60!black] (-1.7,0) node[left,black] {\small $b$} -- (N);
\end{tikzpicture}
\caption{자기 신경전달물질을 통한 되먹임을 가진 \nt{SERO} 뉴런. 농도 $c$는 모든 수신자에게 가고 뉴런 자신도 억제한다. 모델에서 이것은 수준 조절기다. 뇌에서 이 고리가 그런 역할을 한다는 것은 가설이다.}
\label{fig:serofeedback}
\end{figure}

이 고리가 무엇을 하는지 계산해 보자. 전달 특성이 선형, 즉 $u>0$에서 $f(u)=au$라 하자. 평형에서 $c=\kappa r$이고 $r=a(b-gc)$이다. 하나를 다른 하나에 대입하면
\begin{empheq}[box=\keybox]{equation}
r \;=\; \frac{a\,b}{1+G}, \qquad G \;=\; a\,g\,\kappa .
\label{eq:feedback}
\end{empheq}
$G$는 고리 이득이다. 이것은 모든 음의 되먹임 증폭기에 쓰이는 바로 그 공식이며, 같은 두 가지 이점을 준다.

\emph{편차에 대한 강인성.} 뉴런의 흥분성 $a$가 비율 $\varepsilon$만큼 변했다고 하자. 되먹임이 없다면 발화율도 같은 비율만큼 변할 것이다. 되먹임이 있고 $G\gg1$이면 발화율 $r\approx b/(g\kappa)$는 $a$에 거의 의존하지 않는다. 그 변화가 $1+G$배 작아진다. $G=9$이면 편차가 10분의 1로 줄어든다.

\emph{빠르기.} $r=a(b-gc)$를 식~\eqref{eq:seroneuron}의 첫 번째 식에 대입하면 $\tau_c\,\dd c/\dd t=-(1+G)\,c+\kappa ab$가 된다. 농도는 시간 상수 $\tau_c/(1+G)$로 자기 수준에 다가간다. 되먹임이 없을 때보다 $1+G$배 빠르다.

이것이 \nt{SERO} 시스템이 할 수 있는 첫 번째 계산이다. 온도 조절기가 온도를 유지하듯 뇌 전체의 세로토닌 수준을 정해진 값에 유지하는 것이다. 이는 가설이다. 실험에서 자가 억제는 아직 평균 발화율을 바꾸지 않는 짧은 멈춤으로만 보인다.

\section{스파이크에서 수준으로}

두 번째 계산은 식~\eqref{eq:seroneuron}의 첫 번째 식에 숨어 있다. 뉴런은 개별 스파이크를 내지만 수신자는 농도를 느끼며, 농도는 스파이크를 시간 $\tau_c$에 걸쳐 평활한다. 출처의 발화율이 변하면 농도는 늦게 따라간다.
\begin{equation}
c(t) \;=\; \frac{\kappa}{\tau_c}\int_{-\infty}^{t} e^{-(t-t')/\tau_c}\,\sum_j r_j(t')\,\dd t' .
\label{eq:lowpass}
\end{equation}
이것은 최근 몇 $\tau_c$ 동안 출처 활동의 이동 평균, 즉 저역 통과 필터다(그림~\ref{fig:lowpass}). 가장 단순한 디지털-아날로그 변환기도 펄스를 RC 회로에 통과시켜 그 빈도를 수준으로 바꾼다. \nt{SERO} 시스템은 수십만 개의 뉴런으로 한꺼번에 같은 일을 한다.

이 양은 동시에 메모리이기도 하다. 출처가 침묵한 뒤에도 농도는 $\tau_c$ 정도의 시간 동안 남아 있다. 비트가 아니라 서서히 사라지는 흔적이다.

\begin{figure}[t]
\centering
\begin{tikzpicture}[>=Stealth,x=1.45cm,y=1cm]
\draw[gray!70!black] (0,3.2) -- (8.1,3.2);
\node[left,font=\small] at (-0.05,3.4) {스파이크};
\node[above,font=\scriptsize,gray!40!black] at (1,3.65) {초당 $2$회};
\node[above,font=\scriptsize,gray!40!black] at (3.5,3.65) {초당 $8$회};
\node[above,font=\scriptsize,gray!40!black] at (6.5,3.65) {초당 $2$회};
\draw[->,gray!70!black] (0,0) -- (8.3,0) node[right,black] {\small $t$};
\draw[->,gray!70!black] (0,0) -- (0,2.75) node[left,black] {\small $c$};
\foreach \s in {0.136,0.529,0.760,1.223,1.714,1.748,1.755,2.663,2.701,2.734,3.414,3.493,3.719,3.800,3.928,3.948,4.074,4.327,4.420,4.589,4.728,4.736,4.914,5.025,5.205,5.220,6.224,6.544,7.178}{\draw[thick,blue!60!black] (\s,3.2) -- (\s,3.6);}
\foreach \a/\b/\v in {0.000/0.136/0.000,0.136/0.529/1.000,0.529/0.760/1.675,0.760/1.223/2.330,1.223/1.714/2.466,1.714/1.748/2.509,1.748/1.755/3.425,1.755/2.663/4.402,2.663/2.701/2.775,2.701/2.734/3.672,2.734/3.414/4.553,3.414/3.493/3.306,3.493/3.719/4.055,3.719/3.800/4.235,3.800/3.928/4.905,3.928/3.948/5.316,3.948/4.074/6.211,4.074/4.327/6.476,4.327/4.420/6.028,4.420/4.589/6.493,4.589/4.728/6.483,4.728/4.736/6.642,4.736/4.914/7.589,4.914/5.025/7.351,5.025/5.205/7.579,5.205/5.220/7.331,5.220/6.224/8.221,6.224/6.544/4.012,6.544/7.178/3.914,7.178/8.000/3.076}{\draw[very thick,orange!85!black] plot[domain=\a:\b,samples=10] (\x,{0.25*\v*exp(-(\x-\a))});}
\foreach \s/\p/\q in {0.136/0.000/1.000,0.529/0.675/1.675,0.760/1.330/2.330,1.223/1.466/2.466,1.714/1.509/2.509,1.748/2.425/3.425,1.755/3.402/4.402,2.663/1.775/2.775,2.701/2.672/3.672,2.734/3.553/4.553,3.414/2.306/3.306,3.493/3.055/4.055,3.719/3.235/4.235,3.800/3.905/4.905,3.928/4.316/5.316,3.948/5.211/6.211,4.074/5.476/6.476,4.327/5.028/6.028,4.420/5.493/6.493,4.589/5.483/6.483,4.728/5.642/6.642,4.736/6.589/7.589,4.914/6.351/7.351,5.025/6.579/7.579,5.205/6.331/7.331,5.220/7.221/8.221,6.224/3.012/4.012,6.544/2.914/3.914,7.178/2.076/3.076}{\draw[very thick,orange!85!black] (\s,0.25*\p) -- (\s,0.25*\q);}
\draw[thick,densely dashed,black] plot[domain=0:2,samples=30] (\x,{0.25*2*(1-exp(-\x))})
  plot[domain=2:5,samples=40] (\x,{0.25*(8-6.271*exp(-(\x-2)))})
  plot[domain=5:8,samples=40] (\x,{0.25*(2+5.688*exp(-(\x-5)))});
\foreach \x in {0,2,5,8}{\draw[gray!60!black] (\x,-0.05) -- (\x,0.05) node[below=3pt,font=\scriptsize] {\x\,s};}
\end{tikzpicture}
\caption{스파이크에서 수준으로. 위는 \nt{SERO} 뉴런의 스파이크다. 2초 동안 드물게, 3초 동안 자주, 그다음 다시 드물게 나온다. 아래는 $\tau_c=1\,\text{s}$일 때의 세로토닌 농도~\eqref{eq:lowpass}다. 점선은 스파이크가 완벽하게 고른 간격으로 나올 때의 농도다.}
\label{fig:lowpass}
\end{figure}

\section{전선은 곧 부피다}
\label{sec:volume}

명제~\ref{prop:serocomplete}의 조건으로 돌아가자. 가까운 곳에서 여러 송신자가 내보낸 세로토닌은 하나의 농도로 합쳐진다.

\emph{여기서 전선은 섬유가 아니라 조직의 부피다.} 두 신호가 서로 다르게 남으려면, 신경전달물질이 퍼지는 거리보다 더 멀리 떨어진 영역으로 방출되어야 한다. 확산 계수가 $D$인 분자는 시간 $\tau$ 동안 다음 거리만큼 이동한다.
\begin{empheq}[box=\keybox]{equation}
\ell \;\sim\; \sqrt{D\,\tau}\, .
\label{eq:diffusionlength}
\end{empheq}
세포 사이 틈의 굴곡은 작은 분자의 확산을 약 $2.6$배 늦춘다~\cite{Sykova2008}. 조직 속 세로토닌 자체의 확산 계수는 측정된 적이 없으므로, 분자 크기로 보아 $10^{-6}$~cm$^2$/s의 몇 배 정도로 추정하자. 신호가 $\tau\sim1$~s 동안 산다면(이것도 추정이다) $\ell\sim\sqrt{3\cdot10^{-6}}$~cm $\approx20$~\textmu m이다. 이런 네트워크에서 주소는 \emph{장소}다. 수신자는 자신이 어디에 있는지로만 신호가 누구에게서 왔는지 안다.

실제 \nt{SERO} 뉴런은 이런 개별 주소가 거의 남지 않도록 생겼다. 각 뉴런의 출력은 뇌의 넓은 영역에 흩뿌려진다. 세포 하나의 가지가 서로 멀리 떨어진 여러 영역에 닿는다~\cite{GagnonParent2014}. 세포당 방출 지점은 추정으로 약 $10^5$개이며(표~\ref{tab:types}), 직접 센 적은 없다. 서로 다른 \nt{SERO} 뉴런의 “전선”은 겹치고 섞인다. 그래도 완전히 전선 하나로 합쳐지지는 않는다. \nt{SERO} 뉴런은 몇 개의 하위 시스템으로 나뉘며, 이들은 서로 다른 영역으로 출력을 보내고 같은 사건에 서로 다르게 반응한다~\cite{Ren2018}. 그러나 이런 하위 시스템은 수십만 개가 아니라 몇 개뿐이다.

\begin{keyblock}
\begin{proposition}[독립 전선의 수]
\label{prop:volumewires}
체적 전달을 하는 네트워크에서 독립 신호의 수는 뉴런의 수가 아니라, 뉴런들이 신경전달물질을 내보내는 크기 $\ell\sim\sqrt{D\tau}$의 서로 겹치지 않는 영역의 수로 제한된다. 각 뉴런의 출력이 큰 부피에 흩뿌려져 있으면, 네트워크 전체가 전달하는 것은 공통 변수 몇 개뿐이다.
\end{proposition}
\end{keyblock}

그래서 명제~\ref{prop:serocomplete}의 논리 회로는 뇌에서 만들 재료가 없다. 반면 조절기~\eqref{eq:feedback}와 필터~\eqref{eq:lowpass}에는 개별 전선이 필요 없다. 공통된 양 하나면 충분하다. 필터는 목표 영역에서 측정된 부피 방출~\cite{Bunin1998}에서 바로 나온다. 수준 조절기는 가설이다.

\section{계획 지평}
\label{sec:horizon}

공통된 느린 양 하나는 어디에 쓸모가 있을까? 좋은 후보는 네트워크 전체에서 같아야 하고 몇 초나 몇 분보다 빠르게 변하지 않는 매개변수다. \ref{sec:neurontypes}장에서 이미 그런 매개변수가 나왔다. 오차 식~\eqref{eq:rpe}의 계수 $\gamma$로, 미래의 보상이 즉시 받는 보상보다 얼마나 싼지를 나타낸다. 연속 시간에서는 그 자리를 \emph{지평} $\tau_\gamma$가 차지한다. 시간 $D$를 기다려야 받는 보상 $R$의 현재 가치는 $R\,e^{-D/\tau_\gamma}$이다.

지평은 기다릴 가치가 있는지를 정한다. 작은 보상 $r_0$를 바로 받거나 큰 보상 $R$을 기다릴 수 있다고 하자. 다음 조건이 성립하는 동안은 기다리는 편이 이득이다.
\begin{empheq}[box=\keybox]{equation}
D \;<\; \tau_\gamma\,\ln\frac{R}{r_0} .
\label{eq:patience}
\end{empheq}
$\tau_\gamma=2\,\text{s}$이고 나중 보상이 세 배 크다면 $2.2\,\text{s}$까지 기다릴 만하다. 지평이 두 배 길면 $4.4\,\text{s}$까지다. 인내심은 지평에 비례한다.

식~\eqref{eq:patience}는 지수형 할인에 기댄다. 몇 초 지연에서 측정된 할인은 쌍곡선 $R/(1+D/\tau_\gamma)$에 더 가깝다. 원숭이에서 선택 속 지연 보상의 가치도, 그 보상을 알리는 신호에 대한 \nt{DOPA} 뉴런의 응답도 이렇게 떨어진다~\cite{KobayashiSchultz2008}. 쌍곡선에서는 조건~\eqref{eq:patience}가 $D<\tau_\gamma\,(R/r_0-1)$로 바뀐다. 보상 비에 대한 의존이 로그가 아니라 선형이 되고, 같은 수치에서 기다림의 한계는 거의 두 배로 늘어난다. 지연이 무작위이면 같은 지수 함수에서 쌍곡선이 나오며(연습문제~\ref{ex:05:7}), 인내심이 시간 척도에 비례한다는 성질은 그대로 남는다.

가설에 따르면 지평을 정하는 것이 바로 \nt{SERO}다~\cite{Doya2002}. 이 역할에 필요한 것은 다 갖추었다. 네트워크 전체에 하나뿐이고 몇 초에 걸쳐 변하는 공통 수준이다. 직접적인 실험도 있다. 세로토닌 뉴런을 인위적으로 켜면 동물은 지연 보상을 포기하기 전까지 더 오래 기다리고~\cite{Miyazaki2014}, 보상이 드물어진 곳에서 보상을 얻으려고 더 오래 애쓴다~\cite{Lottem2018}. 네트워크 안에서 세로토닌 농도가 정확히 어떻게 $\tau_\gamma$로 바뀌는지는 알려져 있지 않다. 다음 장에서 지평은 \nt{DOPA}가 방송하는 오차에 들어간다.

\section{정리}

\begin{table}[t]
\centering
\small
\begin{tabular}{>{\raggedright}p{3.6cm}>{\raggedright}p{4.3cm}>{\raggedright\arraybackslash}p{4.3cm}}
\toprule
& \nt{GLUT} & \nt{SERO} \\
\midrule
반응 시간 & 밀리초 & 1초 미만~1초 \\
작용의 부호 & $+$만 & $+$와 $-$, 수신자가 고름 \\
입력이 없을 때 & 침묵 & 고른 배경 입력이 있으면 스스로 작동 \\
주소 지정 & 점대점 & 부피, 주소는 장소 \\
NOT & “꺼짐” 검출기를 통해서만 & 뉴런 하나 \\
메모리 & 자기 유지 활동, 피로로 꺼짐 & 농도, 시간 $\tau_c$에 걸쳐 사라짐 \\
주요 계산 & 동시성, 지연, 논리, 순서열 & 평활화; 수준 조절과 지평 $\tau_\gamma$(가설) \\
\bottomrule
\end{tabular}
\caption{\nt{GLUT} 뉴런만으로 된 네트워크와 \nt{SERO} 뉴런만으로 된 네트워크가 계산하는 것.}
\label{tab:glutvssero}
\end{table}

논리 소자로서(표~\ref{tab:glutvssero}) \nt{SERO} 뉴런은 \nt{GLUT}보다 풍부하지만, 통신 시스템으로서는 느리고 주소가 거의 없는 브로드캐스트다. 그래서 \nt{SERO} 뉴런은 프로세서가 되려 하지 않는다. 대신 \ref{sec:neurontypes}장에서 말한 몇 가지 공통된 양을 평활하고 방송하며, 가설에 따르면 계획 지평도 그중 하나다. 페이스메이커, 부피, 느린 농도는 앞으로 세 번 더 만난다. \nt{DOPA}, \nt{NORA}, \nt{ACET}에서다.

\section{연습문제}

\begin{exercise}[\lvlA]
\label{ex:05:1}
\emph{전선은 곧 부피다.} 세로토닌에 대해 $D=3\cdot10^{-6}$~cm$^2$/s로 두자(\ref{sec:volume}절). 식~\eqref{eq:diffusionlength}로 $1\,\text{s}$ 동안 사는 신호의 $\ell$과, 신경전달물질이 $5$~cm 퍼지는 데 걸리는 시간을 구하라. 그렇다면 \nt{SERO}의 브로드캐스트를 가능하게 하는 것은 확산인가, 방출 지점 $\sim10^5$개로 갈라지는 전선인가?
\end{exercise}

\begin{exercise}[\lvlA]
\label{ex:05:2}
\emph{수준 조절기.} 식~\eqref{eq:seroneuron}에서 $f(u)=au$일 때 상대 감도 $S=(a/r)\,\partial r/\partial a$가 $G\gg1$일 때만이 아니라 정확히 $1/(1+G)$임을 보여라. $G=9$에서 흥분성 $a$가 $20\,\%$ 늘었다. 선형화 없이 $r$은 몇 퍼센트 변하는가?
\end{exercise}

\begin{exercise}[\lvlB]
\label{ex:05:3}
\emph{고리 안의 느린 수용체.} \nt{SERO} 수용체는 지연을 두고 응답한다. $r(t)=a\bigl(b-gc(t-T)\bigr)$이고 부피는 식~\eqref{eq:seroneuron}을 따른다. $G>1$일 때 고리는 $T<T^*=\tau_c\arccos(-1/G)/\sqrt{G^2-1}$일 때만 안정함을 보여라. $\tau_c=1\,\text{s}$에서 $G=2$와 $G=9$일 때 $T^*$를 구하라. 지연이 수백 밀리초이면 편차를 얼마나 억제할 수 있는가?
\end{exercise}

\begin{exercise}[\lvlB]
\label{ex:05:4}
\emph{수준의 산탄 잡음.} 각 스파이크는 식~\eqref{eq:lowpass}에 따라 $c$에 $\tau_c=1\,\text{s}$인 지수 함수를 더한다. \nt{SERO} 뉴런 $3\cdot10^5$개가 모두 초당 $2$회의 푸아송 스파이크를 한 부피에 내보낼 때 $c$의 변동 계수를 구하라. 발화율이 초당 $4$회인 뉴런 하나로 $\mathrm{CV}=10\,\%$를 얻으려면 $\tau_c$가 얼마여야 하는가?
\end{exercise}

\begin{exercise}[\lvlB]
\label{ex:05:5}
\emph{시뮬레이션: 되먹임 대 잡음.} 발화율 $r_i=a_i[b-gc]_+$, $a_i$가 $[2.8,\,5.2]$에서 균등한 푸아송 페이스메이커 $N=50$개와 $\kappa=1/N$, $\tau_c=1\,\text{s}$인 부피~\eqref{eq:seroneuron}을 시뮬레이션하라. 되먹임($b=1$, $g=2.25$)은 $b=0.1$, $g=0$에 비해 수준 $c$의 CV를 몇 배 줄이는가? 뉴런들의 발화율도 고르게 만드는가?
\end{exercise}

\begin{exercise}[\lvlB]
\label{ex:05:6}
\emph{설계: \nt{SERO} 논리로 만드는 XOR.} 게이트는 $b-g\sum_kc_k>0$이면 $r_0$를, 아니면 $0$을 자기 부피 $\tau_c\,\dd c/\dd t=-c+\kappa r$로 내보내는 페이스메이커이며, NOR을 계산한다. 이런 게이트 다섯 개로 XOR을 만들고, $\tau_c=1\,\text{s}$, $G'=g\kappa r_0/b=2$일 때 정착 시간을 구하라.
\end{exercise}

\begin{exercise}[\lvlB]
\label{ex:05:7}
\emph{지연을 모를 때의 인내심.} (a)~동물은 기다림이 $6\,\text{s}$ 이하이면 세 배 큰 보상을 기다린다. 이는 어떤 지평 $\tau_\gamma$를 뜻하는가? (b)~지연 $D$가 평균 $\bar D$인 지수 분포를 따른다고 하자. $\bar D<\tau_\gamma(R/r_0-1)$일 때 기다리는 편이 이득임을 보이고, $\tau_\gamma=2\,\text{s}$, $R=3r_0$에서 고정 지연과 비교하라.
\end{exercise}

\begin{exercise}[\lvlC]
\label{ex:05:8}
\emph{농도는 어떻게 지평을 정하는가.} \nt{DOPA} 뉴런은 $V$를 직접 가중치 $k_1$로 받고, $\tau_d=0.1\,\text{s}$로 평활하는 \nt{GABA} 중개 뉴런을 거쳐 가중치 $-k_2$로 받는다. 느린 $V$에 대해 합은 $(k_1-k_2)V+k_2\tau_d\dot V$이다. $\tau_\gamma=2\,\text{s}$인 식~\eqref{eq:ctd}에 맞게 $k_1$, $k_2$를 정하라. \nt{SERO}가 $k_1$에 $m$을 곱한다면, 지평을 두 배로 늘리려면 $m$을 얼마나 바꿔야 하는가?
\end{exercise}

\chapter{학습할 줄 아는 신경망: \nt{DOPA}를 더하다}
\chaptermark{학습하는 신경망}
\label{sec:dofa}
\begin{chaptergoals}
\item 시냅스가 국소 신호만 쓸 수 있는 이유와 그래서 역전파가 배제되는 이유;
\item 헤브 규칙이 배우는 것: 입력이 가장 크게 변하는 방향;
\item 적격 흔적과 브로드캐스트 $\delta$가 잡음을 보상 기울기를 오르는 걸음으로 바꾸는 방식;
\item 비평자가 연속 시간에서 $\delta$를 내는 방식과 공유 전선 하나로는 학습이 느린 이유.
\end{chaptergoals}

\section{게임의 규칙}

앞 장들의 모든 회로에서 가중치는 엔지니어가 정했다. 생체에는 엔지니어가 없다. 가중치는 작동하는 동안 저절로 정해져야 하고, 그 결과 신경망이 쓸모 있는 일을 해야 한다. 이것을 학습이라고 한다.

기존 규칙에 하나가 더해진다. 시냅스는 자기 가중치를 스스로 바꾸며, 알 수 있는 것은 \emph{자기 주변}에서 일어나는 일뿐이다. 즉 송신자의 활동 $x$, 수신자의 활동 $y$, 그리고 시냅스까지 도달하는 물질이다. 누구도 시냅스에 개별 보정값을 보내 줄 수 없다.

이 규칙은 인공 신경망의 핵심 방법인 오차 역전파를 배제한다. 역전파에서는 출력 오차가 순방향 통과가 끝난 뒤 별도의 단계에서 같은 가중치를 거쳐 신경망을 거꾸로 지나가고, 가중치 하나하나가 각자의 보정값을 받는다. 이를 위해서는 순방향 전선을 정확히 되풀이하는 역방향 전선과 공통 클록이 필요하다. 뇌에서는 둘 다 발견되지 않았다. 적어도 인공 신경망과 같은 형태로는 없다~\cite{Lillicrap2020,WhittingtonBogacz2019}.

가중치의 변화는 느린 연속 과정으로 쓰겠다.
\begin{equation}
\frac{\dd w}{\dd t} \;=\; \eta\,\Phi(\text{시냅스가 아는 것}),
\label{eq:plasticity}
\end{equation}
여기서 $\eta$는 학습률이며, 스파이크 하나가 지나가는 동안 가중치가 거의 변하지 않을 만큼 작다. 이 장 전체는 함수 $\Phi$가 어떤 모양일 수 있는지에 관한 것이다.

\section{가중치는 어디에 저장되는가}
\label{sec:weightstore}

“가중치를 바꾸는” 시냅스란 무엇인가? 연결에는 두 쪽이 있다. 송신자 전선의 말단(\emph{시냅스 전} 쪽)과 수용체가 있는 수신자 막의 한 구역(\emph{시냅스 후} 쪽)이며, 그 사이에 좁은 틈이 있다. \nt{GLUT} 연결에서 수신자 막은 보통 \emph{가시}, 곧 수신자 입력 가지에 난 작은 돌기 위에 있다. 연결의 가중치는 세 인수로 나누면 편리하다~\cite{delCastillo1954}.
\begin{empheq}[box=\keybox]{equation}
w \;=\; N_s\,p\,A_1 .
\label{eq:quantal}
\end{empheq}
여기서 $N_s$는 두 뉴런 사이의 방출 부위 수, $p$는 스파이크 하나가 한 부위에서 신경전달물질 한 꾸러미를 방출할 확률(\ref{sec:code}장의 바로 그 $p$), $A_1$은 한 꾸러미에 대한 수신자의 응답, 즉 그 꾸러미를 붙잡는 수용체의 수다. 인수 $p$는 송신자에, $A_1$은 수신자에 있고, $N_s$에는 양쪽이 모두 필요하다. 새 방출 부위가 자라나고 옛 부위가 사라지며, 이 과정은 평생 계속된다.

\emph{결정은 수신자가 한다.} 전형적인 \nt{GLUT} 시냅스에서 글루탐산 수용체 가운데 하나는 두 조건이 겹칠 때만 전류를 흘린다. 송신자에서 글루탐산이 왔고, 수신자 막이 이미 흥분해 있어야 한다~\cite{Nowak1984}. 이것은 분자 하나에 내장된 헤브 규칙이며, \ref{sec:glutcomputer}장의 동시성 검출기를 시냅스 입구에 그대로 설치한 것이다. 이 수용체가 작동하면 수신자 안으로 칼슘이 들어가고, 칼슘이 가중치 변화를 일으킨다. 입력과 출력이 한꺼번에 보이는 곳은 수신자뿐이므로 결정도 수신자에서 내려진다.

\emph{결정의 실행은 어느 쪽이든 할 수 있다.} 가장 많이 연구된 장기 강화는 수신자에서 일어난다. 막에 새 수용체가 끼워지고~\cite{Malinow2002} 가시가 커진다~\cite{Matsuzaki2004}. 즉 $A_1$이 커진다. 그러나 수신자는 $p$도 바꿀 수 있다. 수신자가 틈을 거쳐 송신자 쪽으로 거슬러 가는 물질을 내보내면(부록~\ref{sec:othermediators}의 역방향 채널) 송신자는 신경전달물질을 덜 방출하기 시작한다~\cite{WilsonNicoll2001}. 한편 \ref{sec:code}장의 단기 가중치 변화인 고갈과 촉진은 송신자가 자신의 최근 활동에 따라 스스로 조절하는 $p$의 변화다.

엔지니어의 눈으로 보면 가중치 레지스터가 두 칩에 나뉘어 있다. 학습 규칙은 수신자가 실행하고, 그 결과는 자기 쪽에도, 역방향 채널을 통해 송신자 쪽에도 기록할 수 있다. 그래서 이 장의 규칙에서 “국소적”이라는 말은 연결의 한쪽이 아니라 연결 전체를 뜻한다.

\section{교사 없이 배우기}

시냅스가 할 수 있는 가장 단순한 일은 송신자와 수신자가 함께 활동할 때 강해지는 것이다.
\begin{equation}
\frac{\dd w}{\dd t} \;=\; \eta\,x\,y .
\label{eq:hebb}
\end{equation}
이것이 헤브 규칙이다~\cite{Hebb1949}. 이 규칙은 국소적이고 클록이 필요 없다. 그러나 \ref{sec:glutcomputer}장의 \nt{GLUT} 네트워크와 같은 문제, 곧 양의 되먹임이 있다. 강한 가중치는 $y$를 키우고, 큰 $y$는 가중치를 더 키우며, 가중치는 한없이 자란다. 처방은 가중치에 대한 음의 되먹임이다. 가중치가 $y^2$에 비례해 줄어들게도 하자. 입력이 $\mathbf x$이고 출력이 선형 $y=\mathbf w\cdot\mathbf x$인 뉴런에서는 다음을 얻는다.
\begin{empheq}[box=\keybox]{equation}
\frac{\dd\mathbf w}{\dd t} \;=\; \eta\,y\,\bigl(\mathbf x - y\,\mathbf w\bigr) .
\label{eq:oja}
\end{empheq}
이 규칙도 여전히 국소적이다. 시냅스는 자기 입력 $x_i$, 출력 $y$, 가중치 $w_i$를 안다.

\begin{keyblock}
\begin{proposition}[헤브 규칙이 찾는 것]
\label{prop:oja}
규칙~\eqref{eq:oja}에 따르면 가중치 벡터는 입력이 가장 크게 변하는 방향, 곧 입력 상관행렬 $C=\langle\mathbf x\mathbf x^{\mathsf T}\rangle$의 주 고유벡터 방향으로 향하며 길이가 1로 수렴한다~\cite{Oja1982}. 뉴런은 자기 입력을 가장 잘 대표하는 하나의 값으로 압축해 내보내는 법을 배운다.
\end{proposition}
\end{keyblock}

증명. 식~\eqref{eq:oja}을 입력에 대해 평균하면 $\dd\mathbf w/\dd t=\eta\bigl(C\mathbf w-(\mathbf w^{\mathsf T}C\mathbf w)\,\mathbf w\bigr)$이다. 고정점은 단위 길이의 $C$ 고유벡터들이다. $\mathbf w$가 고윳값 $\lambda_k$를 갖는 고유벡터 $\mathbf e_k$ 근처에 있으면, 다른 벡터 $\mathbf e_j$ 방향의 작은 섭동은 $e^{\eta(\lambda_j-\lambda_k)t}$처럼 자란다. 모든 $\lambda_j<\lambda_k$인 점, 즉 주 벡터만 안정하다.

스파이크의 시각을 보는 헤브 규칙도 있다. 송신자의 스파이크가 수신자의 스파이크보다 $s$밀리초 \emph{먼저} 오면 가중치는 $A_+e^{-s/\tau_+}$만큼 커지고, $s$밀리초 \emph{뒤에} 오면 $A_-e^{-s/\tau_-}$만큼 작아진다. $\tau_\pm$는 20밀리초 정도다. 이 규칙은 \nt{GLUT} 뉴런의 시냅스에서 측정되었다~\cite{BiPoo1998}. 이 규칙은 응답의 \emph{원인이었을 수 있는} 연결을 강화하고, 늦게 온 연결을 약화한다(그림~\ref{fig:stdp}). 같은 흥분 순서를 신경망에 여러 번 흘려 보내면 각 고리에서 다음 고리로 가는 연결이 저절로 자란다. 엔지니어 없이 조립된 \ref{sec:glutcomputer}장의 사슬이다.

\begin{figure}[t]
\centering
\begin{tikzpicture}[>=Stealth,x=0.07cm,y=1.3cm]
\draw[->,gray!70!black] (-66,0) -- (68,0) node[right,black] {\small $s$};
\draw[->,gray!70!black] (0,-1.2) -- (0,1.35) node[above,black] {\small 가중치 변화};
\draw[very thick,blue!60!black] plot[domain=0.5:60,samples=50] (\x,{exp(-\x/20)});
\draw[very thick,red!70!black] plot[domain=-60:-0.5,samples=50] (\x,{-0.6*exp(\x/20)});
\draw[blue!60!black,thick] (0.5,0) -- (0.5,1);
\draw[red!70!black,thick] (-0.5,0) -- (-0.5,-0.6);
\foreach \x in {-40,-20,20,40}{\draw[gray!60!black] (\x,-0.04) -- (\x,0.04);}
\node[below,font=\scriptsize] at (20,-0.04) {$20\ms$};
\node[above,font=\scriptsize] at (-20,0.04) {$-20\ms$};
\node[align=left,font=\small,blue!60!black,anchor=west] at (22,0.95) {송신자가 먼저:\\ 연결 강화};
\node[align=right,font=\small,red!70!black,anchor=east] at (-12,-0.95) {송신자가 나중:\\ 연결 약화};
\end{tikzpicture}
\caption{시각에 따른 헤브 규칙. 가로축은 $s=t_{\text{post}}-t_{\text{pre}}$, 즉 수신자의 스파이크가 송신자의 스파이크보다 얼마나 늦었는지다. $\tau_\pm=20\ms$, $A_-=0.6A_+$. 폭이 수십 밀리초인 이 창은 동시성 창~\eqref{eq:window}과 같은 규모다.}
\label{fig:stdp}
\end{figure}

그러나 이 절의 규칙들은 모두 \emph{자주} 일어나는 것을 가르칠 뿐, \emph{쓸모 있는} 것을 가르치지는 않는다. $x$와 $y$만 아는 시냅스는 그 행동이 주스를 가져왔는지 통증을 가져왔는지 모른다. 결과로부터 배우려면 시냅스에 세 번째 신호가 필요하다.

\section{세 번째 인수}

세 번째 신호는 \nt{DOPA}가 가져온다. 도파민 뉴런은 수 하나, 곧 보상 예측 오차 $\delta$~\eqref{eq:rpe}를 브로드캐스트한다(\ref{sec:neurontypes}장). 가중치 변화가 세 인수, 즉 송신자 활동, 수신자 활동, $\delta$의 곱에 비례한다고 하자.

어려움은 보상이 신경망이 행동을 고른 순간이 아니라 수백 밀리초나 몇 초 뒤에 온다는 데 있다. $\delta$가 도착할 무렵이면 곱 $xy$는 이미 오래전에 0이고, 시냅스에는 자기가 참여했다는 기억이 필요하다. 가장 단순한 기억은 익숙한 RC 회로다.
\begin{empheq}[box=\keybox]{equation}
\tau_e\,\frac{\dd e}{\dd t} \;=\; -e \;+\; x\,y ,
\qquad
\frac{\dd w}{\dd t} \;=\; \eta\,\delta(t)\,e(t) .
\label{eq:threefactor}
\end{empheq}
값 $e$를 \emph{적격 흔적}이라고 부른다~\cite{Fremaux2016}. 함께 일어난 활동은 시냅스에 표지를 남기고, 이 표지는 시간 $\tau_e$ 동안 사라진다. 그 사이에 도파민이 오면 가중치는 $\delta$의 부호대로 바뀌고, 오지 않으면 표지는 흔적 없이 사라진다(그림~\ref{fig:trace}). \nt{GLUT} 시냅스에서 측정된 결과는 이렇다. 함께 일어난 활동 뒤 1--2초 안에 온 도파민은 그 활동을 연결 강화로 바꾸고, 더 늦게 온 도파민은 그러지 못한다~\cite{Yagishita2014}. 세 번째 신호가 도파민이 아니라 수신자 자신의 강한 흥분인 경우에도 몇 초 길이의 가소성 창이 발견되었다~\cite{Bittner2017}.

\begin{figure}[t]
\centering
\begin{tikzpicture}[>=Stealth,x=7cm,y=1cm]
\fill[orange!12] (0.7,-0.2) rectangle (0.8,5.2);
% rows: x*y, delta, traces, weights
\foreach \y/\lab in {4.4/{$x\,y$},3.1/{$\delta$},1.4/{흔적 $e$},0/{가중치 $w$}}{
  \draw[gray!70!black] (0,\y) -- (1.42,\y);
  \node[left,font=\small] at (-0.02,\y+0.3) {\lab};}
\draw[very thick,blue!60!black] (0,4.4) -- (0,4.9) -- (0.2,4.9) -- (0.2,4.4);
\node[right,font=\scriptsize,blue!60!black] at (0.21,4.8) {송신자와 수신자가 함께 활동};
\draw[very thick,orange!85!black] (0,3.1) -- (0.7,3.1) -- (0.7,3.7) -- (0.8,3.7) -- (0.8,3.1) -- (1.4,3.1);
\node[right,font=\scriptsize,orange!85!black] at (0.81,3.6) {도파민 도착};
% traces, each in units of its own maximum
\draw[very thick,blue!60!black] plot[domain=0:0.2,samples=20] (\x,{1.4+1.1*(1-exp(-\x))/(1-exp(-0.2))})
  plot[domain=0.2:1.4,samples=60] (\x,{1.4+1.1*exp(-(\x-0.2))});
\draw[thick,densely dashed,gray!50!black] plot[domain=0:0.2,samples=30] (\x,{1.4+1.1*(1-exp(-\x/0.05))/(1-exp(-4))})
  plot[domain=0.2:1.4,samples=80] (\x,{1.4+1.1*exp(-(\x-0.2)/0.05)});
\node[right,font=\scriptsize,blue!60!black] at (1.0,2.15) {$\tau_e=1\,\text{s}$};
\node[right,font=\scriptsize,gray!40!black] at (0.3,1.65) {$\tau_e=50\ms$};
% weights
\draw[very thick,blue!60!black] (0,0.25) -- (0.7,0.25) -- (0.8,0.8) -- (1.4,0.8) node[right,font=\scriptsize] {$\tau_e=1\,\text{s}$: 가중치 증가};
\draw[thick,densely dashed,gray!50!black] (0,0.2) -- (1.4,0.2) node[right,font=\scriptsize,gray!40!black] {$\tau_e=50\ms$: 변화 없음};
% time axis marks
\foreach \x/\l in {0/0,0.2/0.2,0.7/0.7,1.4/1.4}{\draw[gray!60!black] (\x,-0.05) -- (\x,0.05) node[below=3pt,font=\scriptsize] {\l\,s};}
\draw[<->,thick,gray!50!black] (0.2,5.35) -- node[above,font=\scriptsize,black] {보상 지연 $0.5\,\text{s}$} (0.7,5.35);
\end{tikzpicture}
\caption{규칙~\eqref{eq:threefactor}에 따른 적격 흔적. 함께 일어난 활동은 $0.2\,\text{s}$ 동안 이어지고, 도파민은 활동이 끝나고 0.5초 뒤에 온다. 흔적은 각자의 최댓값에 대한 비율로 나타냈다. 느린 흔적($\tau_e=1\,\text{s}$)은 이 순간에도 아직 살아 있지만, 빠른 흔적($\tau_e=50\ms$)은 이미 오래전에 사라졌다.}
\label{fig:trace}
\end{figure}

\begin{keyblock}
\begin{proposition}[브로드캐스트 신호에 의한 학습]
\label{prop:threefactor}
규칙~\eqref{eq:threefactor}은 최근 $\tau_e$ 동안 신경망의 작동에 참여한 시냅스만을, 공통 신호 $\delta$가 정한 방향으로 바꾼다. 브로드캐스트 신호와 국소적인 적격 흔적이 함께 주소 지정된 보정을 만든다. 행동과 결과 사이의 지연은 $\tau_e$를 넘지 않는 한 허용된다.
\end{proposition}
\end{keyblock}

\section{왜 작동하는가: 탐색으로서의 잡음}
\label{sec:dofagradient}

규칙~\eqref{eq:threefactor}은 왜 더 나은 쪽으로 이끄는가? 답은 \ref{sec:code}장의 잡음에 있다. 뉴런의 출력을 $y=f(u)+\xi$라 하자. 여기서 $u=\sum_jw_jx_j$이고, $\xi$는 평균이 0이고 분산이 $\sigma^2$인 무작위 흔들림이다. 보상 $R$은 $y$에 따라 정해진다. 적격 흔적으로 $x_jy$ 전체가 아니라 그 무작위 부분 $x_j\xi$를 쓰고, 세 번째 인수로 오차 $\delta=R-\bar R$를 쓰자. $\bar R$는 기대 보상이다. 잡음이 작으면 $R\approx R\bigl(f(u)\bigr)+R'\,\xi$이므로 다음이 성립한다.
\begin{empheq}[box=\keybox]{equation}
\bigl\langle \delta\,\xi\,x_j \bigr\rangle \;=\; \sigma^2\,R'\,x_j \;=\; \frac{\sigma^2}{f'(u)}\,\frac{\partial\langle R\rangle}{\partial w_j} .
\label{eq:perturbation}
\end{empheq}
평균 걸음은 기대 보상의 기울기 방향으로 간다~\cite{Williams1992}. 누구도 도함수를 계산하지 않았다. 신경망이 무작위로 벗어났고, 도파민이 결과가 나아졌는지 알렸으며, 참여한 시냅스들이 유리한 쪽으로 움직였다(그림~\ref{fig:perturbation}). \ref{sec:code}장에서 메시지 전달을 방해하던 잡음이 여기서는 탐색이 된다.

\begin{figure}[t]
\centering
\begin{tikzpicture}[>=Stealth,x=2cm,y=3cm]
\draw[->,gray!70!black] (0,0) -- (4.2,0) node[below left,font=\small,black] {출력 $y$};
\draw[->,gray!70!black] (0,0) -- (0,1.2) node[left,font=\small,black] {보상 $R$};
% reward hill
\draw[very thick,blue!60!black] plot[domain=0:4,samples=60] (\x,{max(0,1-((\x-3)/2.2)^2)});
\node[font=\scriptsize,blue!60!black,anchor=south] at (3,1.02) {최선의 출력};
% noise around f(u)
\draw[thick,gray!60!black] plot[domain=0.6:2.4,samples=50] (\x,{0.28*exp(-((\x-1.5)/0.3)^2/2)});
\draw[densely dashed,gray!60!black] (1.5,0) -- (1.5,0.535);
\node[below,font=\small] at (1.5,0) {$f(u)$};
\node[font=\scriptsize,gray!50!black] at (1.5,-0.2) {흔들림 $\xi$};
% expected reward
\draw[densely dotted,gray!60!black] (0.5,0.535) node[left,font=\scriptsize,black] {$\bar R$} -- (2.1,0.535);
% two random deviations
\fill[green!45!black] (2.1,0.833) circle (0.035);
\draw[very thick,green!45!black] (2.1,0.535) -- (2.1,0.833);
\node[font=\scriptsize,green!45!black,anchor=west,align=left] at (2.18,0.6) {$\xi>0$: 더 좋음,\\ $\delta>0$};
\fill[red!70!black] (0.9,0.089) circle (0.035);
\draw[very thick,red!70!black] (0.9,0.535) -- (0.9,0.089);
\node[font=\scriptsize,red!70!black,anchor=east,align=right] at (0.83,0.27) {$\xi<0$: 더 나쁨,\\ $\delta<0$};
% resulting step
\draw[->,very thick,orange!85!black] (1.55,0.4) -- (1.95,0.4) node[right,font=\scriptsize,black] {걸음};
\end{tikzpicture}
\caption{탐색으로서의 잡음~\eqref{eq:perturbation}. 뉴런의 출력은 $f(u)$ 주위에서 흔들린다. 오른쪽으로의 무작위 이탈(초록)은 기대보다 많은 보상을 주어 $\delta>0$이고, 왼쪽으로의 이탈(빨강)은 기대보다 적어 $\delta<0$이다. 두 경우 모두 곱 $\delta\,\xi$는 양수이고, 가중치는 보상이 커지는 오른쪽으로 움직인다. 평균적으로 걸음은 기울기 $R'$에 비례한다. 언덕 꼭대기에서는 양쪽 이탈이 똑같이 나쁘고, 걸음들이 서로 상쇄된다.}
\label{fig:perturbation}
\end{figure}

\begin{keyblock}
\begin{proposition}[역방향 전선 없는 경사 하강]
\label{prop:gradient}
신경망 활동에 무작위 이탈이 있고, 브로드캐스트 신호가 보상 예측 오차와 같으며 각 상황에서 평균이 0이면, 규칙~\eqref{eq:threefactor}은 평균적으로 가중치를 기대 보상의 기울기 방향으로 바꾼다. 이를 위해 역방향 전선도, 별도의 학습 단계도 필요 없다.
\end{proposition}
\end{keyblock}

이제 도파민이 왜 보상 자체가 아니라 보상 예측의 \emph{오차}를 알리는지 분명해진다. 세 번째 인수가 보상 자체 $R\ge0$이었다면 식~\eqref{eq:perturbation}의 평균은 변하지 않겠지만 두 가지 문제가 생긴다. 첫째, 쓸모 있는 부분에 항 $\bar R\,\xi x_j$가 더해진다. 평균은 0이지만 강한 잡음이다. 둘째, 더 나쁜 문제로, 실제 시냅스는 적격 흔적 $x_jy$에서 무작위 부분과 무작위가 아닌 부분을 가려낼 수 없다. 입력이 주어지면 $x_jf(u)$는 시도마다 같은 수다. 그 상황에서 $\delta$의 평균이 0이면 흔적의 이 부분은 아무것도 바꾸지 않지만, $\delta\ge0$이면 이 부분이 신경망이 한 모든 일을, 옳은 것이든 그른 것이든, 거듭 강화한다. 그래서 $\bar R$는 평생의 평균이 아니라 \emph{주어진 상황에서의} 기대값이어야 한다. 이것을 계산하는 것은 아래에서 다룰 비평자의 몫이다.

이를 모델로 확인했다. 두 \nt{GLUT} 뉴런 무리가 그림~\ref{fig:wta}처럼 상호 억제를 통해 두 행동 가운데 하나를 고른다. “시작” 신호에서 각 무리로 가는 가중치는 처음에 같고, 누가 이길지는 잡음이 정한다. 행동 $A$는 확률 $0.8$로, 행동 $B$는 확률 $0.2$로 주스를 가져오며, 주스는 선택하고 0.5초 뒤에 온다. $\tau_e=1\,\text{s}$, $\delta=R-\bar R$일 때 500회 시도 동안 $A$를 고른 비율은 처음 100회의 $0.65$--$0.8$에서 마지막 100회의 $0.96$--$1.0$으로 올랐고, 가중치는 $0.85$--$1.35$ 범위에 머물렀다. 보상 지연보다 짧은 $\tau_e=50\ms$에서는 신경망이 아무것도 배우지 못했다. $A$를 고른 비율은 절반 근처에 머물렀다. 기대값을 빼지 않은 $\delta=R$에서도 신경망은 $A$를 고르게 되었지만, 같은 500회 동안 승자의 가중치가 천 배로 커졌고 계속 자랐다. 그리고 같은 $\delta=R$에 학습률을 열 배로 높이자, 세 번의 실행 가운데 한 번은 신경망이 처음의 무작위 선택, 그것도 더 나쁜 쪽을 영영 굳혀 버렸다.

\section{전선 하나의 대가}

신호 $\delta$는 모두에게 하나뿐인데, 그것이 평가하는 잡음은 결과에 영향을 주는 뉴런 $N$개 모두의 흔들림의 합이다. 각 시냅스는 $\delta$ 안에서 자기의 쓸모 있는 몫과 이웃 $N-1$개의 잡음을 함께 본다. 자기 몫을 가려내려면 많은 시도에 걸쳐 평균해야 하고, 학습 시간은 $N$에 비례해 늘어난다~\cite{Werfel2005}. 역전파는 이런 대가를 치르지 않는다.

\begin{keyblock}
\begin{proposition}[브로드캐스트의 대가]
\label{prop:broadcastcost}
하나의 공통 신호로 학습하는 시간은 잡음이 결과에 영향을 주는 뉴런의 수에 비례해 늘어난다. 이런 학습은 몇 가지 선택지 중에서 고르는 작은 회로에는 알맞지만, 수십억 개의 가중치를 조율하는 데에는 맞지 않는다.
\end{proposition}
\end{keyblock}

뇌는 실제로 이렇게 일을 나누는 것으로 보인다. \nt{DOPA} 뉴런의 출력은 무엇보다 행동을 고르는 영역으로 가고(\ref{sec:neurontypes}장), 가중치의 압도적 다수가 있는 피질의 거대한 \nt{GLUT} 네트워크는 주로 외부 교사 없이 국소 규칙으로 배운다. 예전에는 이 규칙들을 헤브 규칙~\eqref{eq:oja}으로 환원했다. 지금은 피질에 고유한 목표, 곧 자기의 다음 입력을 \emph{예측하는} 목표가 있다는 증거가 늘고 있다. 그렇다면 오차는 예측한 것과 실제로 온 것의 차이로서 그 자리에서 계산되고~\cite{Keller2018}, 교사는 신경망 자체에 내장되어 있다. 세 번째 신호가 수신자 자신의 강한 흥분인 몇 초 길이의 가소성 창~\cite{Bittner2017}은 시냅스에 국소적인 오차 신호가 실제로 있음을 보여 준다. 이것이 피질의 일반 규칙인지는 가설이다.

\section{오차는 어디서 오는가: 연속 시간의 비평자}

남은 질문은 \nt{DOPA} 뉴런이 $\delta$를 어떻게 계산하느냐다. 강화 학습 프로그램에서는 예측 오차를 $t,t+1,\dots$ 단계로 센다. 생체에는 단계가 없으므로 오차를 연속 시간으로 쓰자~\cite{Doya2000}. $r(t)$를 보상의 흐름, $V(t)$를 앞으로 올 모든 보상의 기대값이라 하되, 더 먼 보상일수록 낮게 친다.
\begin{equation}
V(t) \;=\; \int_t^{\infty} e^{-(s-t)/\tau_\gamma}\,r(s)\,\dd s .
\end{equation}
미분하면 $\dd V/\dd t=V/\tau_\gamma-r$을 얻는다. 이 등식의 잔차가 바로 예측 오차다.
\begin{empheq}[box=\keybox]{equation}
\delta(t) \;=\; r(t) \;+\; \frac{\dd V}{\dd t} \;-\; \frac{V}{\tau_\gamma} .
\label{eq:ctd}
\end{empheq}
상수 $\tau_\gamma$는 계획 지평선으로, 할인율 $\gamma$의 연속 버전이다. 가설에 따르면 이것을 정하는 것은 \nt{SERO}이며(\ref{sec:horizon}절), 그렇다면 식~\eqref{eq:patience}은 얼마나 기다릴 가치가 있는지를 정하는 규칙이다. 세 경우를 확인해 보자(그림~\ref{fig:ctdcases}). 주스를 약속하는 전등이 켜졌다. $V$가 계단처럼 뛰어오르고 $\dd V/\dd t$가 커서 급증이 생긴다. 약속된 주스가 왔다. $r>0$이지만 같은 순간 기대가 소진되어 $\dd V/\dd t\approx-r$이고, 급증은 없다. 주스가 오지 않았다. 기대가 보상 없이 사라져 $\dd V/\dd t<0$이고, 급감이 생긴다.

\begin{figure}[t]
\centering
\begin{tikzpicture}[>=Stealth,x=1.15cm,y=1cm]
\foreach \col/\ttl/\lampon/\juice in {0/{예상 못 한 주스}/0/1, 1/{약속된 주스}/1/1, 2/{주스가 오지 않음}/1/0}{
 \begin{scope}[xshift=\col*4.6cm]
  \node[font=\small] at (1.5,3.55) {\ttl};
  \foreach \yy in {2.4,1.2,0}{\draw[gray!60!black] (0,\yy) -- (3.1,\yy);}
  % reward r
  \ifnum\juice=1
    \fill[green!45!black] (1.95,2.4) rectangle (2.05,2.85);
    \pic[scale=0.55] at (2.0,3.1) {juice};
  \else
    \pic[scale=0.55] at (2.0,3.1) {nojuice};
  \fi
  % expectation V and error delta
  \ifnum\lampon=1
    \pic[scale=0.55] at (1.0,3.1) {lamp};
    \draw[very thick,blue!60!black] (0,1.2) -- (1,1.2) -- (1,1.45)
      plot[domain=1:2,samples=20] (\x,{1.2+0.25*exp((\x-1)/1.5)}) -- (2,{1.2+0.25*exp(1/1.5)}) -- (2,1.2) -- (3.1,1.2);
    \draw[very thick,red!70!black] (0,0) -- (0.95,0) -- (0.95,0.5) -- (1.05,0.5) -- (1.05,0) -- (3.1,0);
    \ifnum\juice=0
      \draw[very thick,red!70!black] (1.95,0) -- (1.95,-0.45) -- (2.05,-0.45) -- (2.05,0);
      \node[font=\scriptsize,red!70!black,anchor=west] at (2.1,-0.35) {급감};
    \else
      \node[font=\scriptsize,gray!50!black,anchor=north,align=center] at (2.0,-0.08) {$r$과 $V$의 하락이\\ 서로 상쇄};
    \fi
  \else
    \draw[very thick,blue!60!black] (0,1.2) -- (3.1,1.2);
    \draw[very thick,red!70!black] (0,0) -- (1.95,0) -- (1.95,0.5) -- (2.05,0.5) -- (2.05,0) -- (3.1,0);
  \fi
  \ifnum\col=0
    \node[font=\small,anchor=east] at (-0.1,2.55) {$r$};
    \node[font=\small,anchor=east] at (-0.1,1.35) {$V$};
    \node[font=\small,anchor=east] at (-0.1,0.1) {$\delta$};
  \fi
 \end{scope}}
\end{tikzpicture}
\caption{오차~\eqref{eq:ctd}의 세 경우. 위에서 아래로 보상 $r$, 기대 $V$, 오차 $\delta$다. 왼쪽에는 기대가 없어서 주스 전체가 뜻밖의 일이다. 가운데에서는 전등이 $V$를 계단처럼 끌어올린다. 이것이 $\dd V/\dd t$의 도약이고, $\delta$의 급증은 전등에서 일어난다. 그 뒤 $V$는 지평선이 요구하는 만큼 정확히 자라고 $\delta=0$이다. 주스가 오는 순간 보상 $r$과 $V$의 하락이 서로 상쇄된다. 오른쪽에서는 $V$가 보상 없이 떨어져 급감이 생긴다. 이것은 그림~\ref{fig:rpe}의 급증, 이동, 급감과 같지만, 이제는 그것들이 어디서 오는지 보인다.}
\label{fig:ctdcases}
\end{figure}

우리가 이미 가진 회로로 식~\eqref{eq:ctd}을 계산하려면 무엇이 필요한가? 세 가지다.

\emph{추정값 $V$ 자체.} 이것은 \nt{GLUT} 뉴런 무리, 곧 비평자가 내보낸다. 비평자의 가중치는 같은 규칙~\eqref{eq:threefactor}으로 같은 $\delta$를 써서 배운다. $\delta>0$이면 기대가 너무 낮았던 것이고, 그 직전에 활동한 연결이 강화된다.

\emph{도함수 $\dd V/\dd t$.} 이것은 수준이 아니라 변화에 반응하는 회로라면 무엇이든 계산할 수 있다. \ref{sec:gaba}장의 방향 검출기처럼 \nt{GABA} 중개자를 거친 지연 억제와 직접 흥분을 함께 쓰는 회로, 또는 \ref{sec:code}장의 고갈되는 시냅스~\eqref{eq:depression}가 그렇다.

\emph{전선 하나에 부호 싣기.} 오차는 양수일 수도 음수일 수도 있지만, 스파이크 빈도는 양수뿐이다. 해법은 \ref{sec:serocomputer}장의 \nt{SERO} 뉴런과 같다. \nt{DOPA} 뉴런은 페이스메이커다. 초당 몇 번의 고른 스파이크는 $\delta=0$을, 버스트는 $\delta>0$을, 멈춤은 $\delta<0$을 뜻한다. 음의 오차에는 여유가 더 적다. 빈도를 0 아래로 내릴 수 없기 때문이다. 실제 \nt{DOPA} 뉴런에서도 급감은 급증보다 얕다~\cite{Bayer2005}.

도파민이 $\delta$처럼 행동한다는 것은 측정되었다. 뇌가 $\dd V/\dd t$를 정확히 어떻게 계산하는지는 가설이다.

\section{행위자와 비평자}

이제 모든 것을 합쳐 보자(그림~\ref{fig:actorcritic}). 상황 뉴런들은 \nt{GABA}를 통해 승자를 고르는 행동 무리(명제~\ref{prop:wta}), 곧 \emph{행위자}를 흥분시키고, $V$를 내보내는 비평자도 흥분시킨다~\cite{SuttonBarto2018}. \nt{DOPA} 뉴런은 보상 $r$과 $V$를 받아 식~\eqref{eq:ctd}을 계산하고, $\delta$를 행위자와 비평자의 모든 시냅스에 브로드캐스트한다.

\begin{figure}[t]
\centering
\begin{tikzpicture}[>=Stealth,
  nrn/.style={circle,draw,very thick,fill=blue!6,minimum size=0.9cm,inner sep=0pt,font=\small},
  gab/.style={circle,draw,very thick,fill=red!10,minimum size=0.6cm,inner sep=0pt},
  dofa/.style={circle,draw,very thick,double,double distance=1.2pt,fill=orange!15,minimum size=1.35cm,inner sep=0pt,font=\scriptsize},
  inh/.style={-{Bar[width=7pt]},thick,red!70!black},
  bc/.style={->,thick,densely dashed,orange!85!black}]
\node[nrn] (S1) at (0,2.4) {$x_1$};
\node[nrn] (S2) at (0,1.2) {$x_2$};
\node[nrn] (S3) at (0,0) {$x_3$};
\node[left=0.15cm,align=right,font=\small] at (-0.5,1.2) {상황};
\node[nrn] (A) at (4,2.6) {$A$};
\node[nrn] (B) at (4,1.0) {$B$};
\node[gab] (G) at (5.2,1.8) {};
\draw[->,thick] (A) -- (G); \draw[->,thick] (B) -- (G);
\draw[inh] (G) to[bend left=35] (A);
\draw[inh] (G) to[bend right=35] (B);
\node[above,font=\small] at (4.6,3.05) {행위자: 행동 선택};
\node[nrn] (V) at (4,-1.1) {$V$};
\node[below,font=\small] at (4,-1.6) {비평자};
\foreach \s in {S1,S2,S3}{\foreach \t in {A,B,V}{\draw[->,gray!60!black] (\s) -- (\t);}}
\node[dofa] (D) at (8.0,-1.1) {\nt{DOPA}};
\draw[->,thick] (V) -- node[above,font=\scriptsize] {$V,\ \dd V/\dd t$} (D);
\draw[->,thick,green!45!black] (10.0,-1.1) node[right,black,font=\small] {보상 $r$} -- (D);
\pic[scale=1.1] at (10.6,-0.5) {juice};
\draw[bc] (D.north) -- (8.0,4.0) -- node[above,font=\small,black] {$\delta$: 행위자와 비평자의 모든 시냅스로} (2.2,4.0) -- (2.2,3.0);
\draw[bc] (D.south) -- (8.0,-2.4) -- (2.2,-2.4) -- (2.2,-0.6);
\end{tikzpicture}
\caption{행동 선택을 배우는 신경망. 회색 화살표는 가중치가 바뀌는 \nt{GLUT} 시냅스, 빨간 원은 \nt{GABA} 뉴런, 이중 원은 식~\eqref{eq:ctd}으로 $\delta$를 계산하는 \nt{DOPA} 페이스메이커, 점선 화살표는 $\delta$의 브로드캐스트다.}
\label{fig:actorcritic}
\end{figure}

신경전달물질 작용의 부호는 수신자가 고르며(명제~\ref{prop:receptor}), 행동 선택 영역에서는 뉴런 일부가 한 계열의 도파민 수용체를, 다른 일부가 다른 계열의 수용체를 가진다. 뇌 절편 실험에서 도파민은 함께 일어난 활동과 맞물려 앞의 뉴런에서는 시냅스를 강화하고, 뒤의 뉴런에서는 약화한다~\cite{Shen2008}. 모델로 말하면 뒤의 뉴런에서는 식~\eqref{eq:threefactor}의 $\delta$ 부호가 뒤집혀 있다. 앞의 뉴런은 \emph{할 만한 일}을, 뒤의 뉴런은 \emph{하지 말아야 할 일}을 배운다.

\section{정리}

\begin{table}[t]
\centering
\footnotesize
\setlength{\tabcolsep}{4pt}
\renewcommand{\arraystretch}{1.15}
\begin{tabular}{>{\raggedright}p{2.6cm}>{\raggedright}p{2.7cm}>{\raggedright}p{3.1cm}>{\raggedright\arraybackslash}p{5.0cm}}
\toprule
규칙 & 시냅스가 아는 것 & 신경망이 배우는 것 & 알려진 사실 \\
\midrule
빈도에 따른 헤브 규칙~\eqref{eq:oja} & $x$, $y$, 자기 가중치 & 입력 압축: 주 방향 & 함께 활동할 때의 강화는 측정됨. 가중치 제한 메커니즘은 연구 중 \\
시각에 따른 헤브 규칙 & $\sim20\ms$ 안에서 $x$와 $y$의 스파이크 순서 & 인과 관계, 순서열 & \nt{GLUT} 시냅스에서 측정됨 \\
3요소 규칙~\eqref{eq:threefactor} & $x$, $y$, 적격 흔적 $e$, 공통 신호 $\delta$ & 보상을 가져오는 행동 & 활동 뒤 몇 초 안의 도파민 효과는 측정됨. 선택 학습의 주된 메커니즘이라는 것은 근거가 탄탄한 가설 \\
비평자~\eqref{eq:ctd} & 위와 같음 & 보상 기대 $V$ & 도파민과 $\delta$의 유사성은 측정됨. $\dd V/\dd t$를 계산하는 회로는 가설 \\
오차 역전파 & 가중치마다의 개별 보정값 & 무엇이든 배우지만 역방향 전선과 클록이 필요 & 뇌에서 이 형태로는 발견되지 않음. 그 근사를 찾는 중 \\
\bottomrule
\end{tabular}
\caption{\nt{GLUT}, \nt{GABA}, \nt{DOPA} 뉴런으로 이루어진 신경망의 학습 규칙.}
\label{tab:learning}
\end{table}

학습 규칙은 표~\ref{tab:learning}에 정리했다. \nt{GLUT}는 데이터를 나르고, \nt{GABA}는 흐름을 제어하며, \nt{DOPA}는 신경망의 변화 가운데 어떤 것을 남길지 결정한다.

\section{연습문제}

\begin{exercise}[\lvlA]
\label{ex:06:1}
\emph{적격 흔적의 수명.} 함께 일어난 활동 $xy=1$이 $e=0$에서 시작해 $T_a=0.2\,\text{s}$ 동안 이어지고, 도파민은 활동이 끝나고 $d=0.5\,\text{s}$ 뒤에 온다. (a)~이 순간 흔적~\eqref{eq:threefactor}은 $\tau_e=1\,\text{s}$일 때가 $\tau_e=50\ms$일 때보다 몇 배 큰가? (b)~흔적이 가장 큰 $\tau_e$는 얼마인가?
\end{exercise}

\begin{exercise}[\lvlA]
\label{ex:06:2}
\emph{헤브 규칙의 표류.} 송신자와 수신자가 각각 초당 $10$회의 독립 푸아송 스파이크열을 내고, 스파이크 쌍마다 그림~\ref{fig:stdp}의 창($\tau_\pm=20\ms$, $A_-=0.6A_+$)대로 가중치가 바뀐다. 완전히 무작위인 활동 한 시간 동안 가중치는 $A_+$의 몇 배만큼 커지는가? 표류가 사라지는 $\tau_-$는 얼마인가?
\end{exercise}

\begin{exercise}[\lvlB]
\label{ex:06:3}
\emph{공분산이 아니라 상관.} 명제~\ref{prop:oja}의 규칙은 $C=\langle\mathbf x\mathbf x^{\mathsf T}\rangle$의 주 고유벡터를 찾는다. 빈도는 양수다. $\mathbf x=\mathbf m+\mathbf n$, $\mathbf m=(\mu,0)$, 잡음의 공분산은 $\operatorname{diag}(s_1^2,s_2^2)$이다. 뉴런이 $\mathbf e_1$을 배우는 조건은? $\mu=1$, $s_1=0.1$, $s_2=0.5$이면 무엇을 배우는가?
\end{exercise}

\begin{exercise}[\lvlA]
\label{ex:06:4}
\emph{급증 속의 지평선.} 전등이 예측할 수 없는 순간에 켜지고, 크기 $R$의 주스가 $L=1\,\text{s}$ 뒤에 온다. 식~\eqref{eq:ctd}으로 학습된 기대에서는 전등과 주스 사이에 $\delta\equiv0$임을 보이고, $\tau_\gamma=2\,\text{s}$와 $\tau_\gamma=4\,\text{s}$일 때 전등에서 생기는 급증의 넓이를 구하라.
\end{exercise}

\begin{exercise}[\lvlB]
\label{ex:06:5}
\emph{브로드캐스트의 대가는 어디서 오는가.} $N$개의 뉴런이 독립적으로 흔들리고($\xi_i\sim\mathcal N(0,\sigma^2)$), 오차는 $\delta=a\sum_i\xi_i$, 시냅스 $j$는 기울기를 $\delta\,\xi_jx_j$로 추정한다. 시도 한 번의 신호 대 잡음비를 구하라. 뉴런 하나가 $5\,\text{s}$짜리 시도 $100$회로 배운다면 $N=10^4$과 $N=10^{10}$에서는 학습에 얼마나 걸리는가?
\end{exercise}

\begin{exercise}[\lvlB]
\label{ex:06:6}
\emph{설계: $\delta$를 계산하는 회로.} \nt{DOPA} 뉴런은 $r$, 가중치 $k_1$인 $V$, 그리고 $\tau_d$로 평활한 $V$의 복사본을 가중치 $-k_2$로 받는다. (a)~느린 $V$에 대해 출력이 식~\eqref{eq:ctd}이 되도록 $k_1$, $k_2$를 정하라. (b)~$V=V_0e^{t/\tau_\gamma}$이면 참 $\delta=0$이다. $\tau_\gamma=2\,\text{s}$일 때 회로의 거짓 오차가 $V/\tau_\gamma$의 $5\,\%$ 이하가 되는 $\tau_d$는?
\end{exercise}

\begin{exercise}[\lvlB]
\label{ex:06:7}
\emph{모델링: 보상 지연의 한계.} \texttt{models/\allowbreak dofa\_learning.py}를 복사해 \texttt{run()}에 보상 지연 $d$를 넣어라. $d=0.5\,\text{s}$에서 $\tau_e=0.05$, $0.2$, $1$, $4\,\text{s}$일 때, 그리고 $\tau_e=1\,\text{s}$, $d=3\,\text{s}$일 때 $500$회 시도 중 마지막 100회에서 $A$를 고른 비율을 구하라. 보상까지 살아남는 흔적의 비율(연습문제~\ref{ex:06:1})이 얼마일 때 학습이 사라지는가?
\end{exercise}

\begin{exercise}[\lvlC]
\label{ex:06:8}
\emph{브로드캐스트의 대가를 피할 수 있는가.} 뉴런 $N$개: $y_i=w_i+\xi_i$, $\xi_i\sim\mathcal N(0,\sigma^2)$, 보상 $R=-\sum_i(y_i-w_i^*)^2$, 규칙 $\Delta w_i=\eta\,(R-\langle R\rangle)\,\xi_i$. 최적의 $\eta$에서 오차 $\sum_i(w_i-w_i^*)^2$가 $e$분의 1로 줄어드는 데 몇 번의 시도가 필요한가? $N$개 중 무작위로 고른 $m$개만 흔들리면 어떤가? 희소 탐색이 도움이 되는가?
\end{exercise}

\chapter{도파민에 기반한 현대의 대안 학습 이론}
\chaptermark{\nt{DOPA}의 다른 이론들}
\label{sec:dofaalt}
\begin{chaptergoals}
\item 오차에 비대칭으로 반응하는 \nt{DOPA} 뉴런들이 보상 분포 전체를 기록하는 방식;
\item 일부 \nt{DOPA} 뉴런이 오차의 부호 대신 중요도나 위험을 알리는 방식;
\item 전선 하나가 빠른 학습 오차와 행동 속도를 정하는 느린 수준을 함께 나르는 방식;
\item \nt{DOPA} 신호가 수 하나가 아니라 다발인 이유와 오차 자체에 도전하는 이론.
\end{chaptergoals}

\section{전선에는 실제로 무엇이 흐르는가}

\ref{sec:dofa}장에서 \nt{DOPA} 시스템은 전선 하나였고, 그 전선으로는 수 하나, 곧 보상 예측 오차 $\delta$~\eqref{eq:ctd}가 흘렀다. 이 그림은 많은 것을 설명한다. 그러나 도파민을 정밀하게 측정할수록 그림에 들어맞지 않는 것이 더 많이 발견된다. 어떤 뉴런은 불쾌한 자극에 보상처럼 반응한다. 동물이 목표를 향해 달리는 동안에는 예상 밖의 일이 전혀 없는데도 도파민 농도가 오른다. \nt{DOPA} 뉴런들은 저마다 다르게 행동한다.

엔지니어에게 이 발견들은 모두 한 가지 질문으로 모인다. \emph{브로드캐스트 버스 위의 신호는 어떤 형식인가?} 수 하나인가, 벡터인가? 전선 위의 신호는 하나인가, 주파수로 나뉜 여러 개인가? $\delta$처럼 미래를 말하는가, 과거를 말하는가? 아래에서 현대의 답 여섯 가지를 살펴본다. 각각에 대해 측정된 것과 추정된 것을 구분한다.

\section{수 하나가 아니라 분포}

오차~\eqref{eq:ctd}는 비평자에게 보상의 \emph{평균} 기대값을 가르친다. 그런데 확실한 주스 반 방울과 확률 1/2의 주스 한 방울은 평균이 같다. 둘을 구분하려면 보상의 분포 전체가 필요하다.

가장 단순한 문제를 보자. 예고 신호 뒤에 크기가 무작위인 보상 $r$가 온다. \nt{DOPA} 뉴런이 많고, $i$번째 뉴런은 저마다 자기 추정값 $V_i$를 맡아서 보상 순간의 오차가 $\delta_i=r-V_i$라고 하자. 또 각 뉴런은 양의 오차를 가중치 $\alpha_i^+$로, 음의 오차를 가중치 $\alpha_i^-$로 반영한다고 하자. 보상이 올 때마다 추정값은 다음만큼 이동한다.
\begin{equation}
\Delta V_i \;=\; \eta\,\bigl(\alpha_i^+\,[\delta_i]_+ \;-\; \alpha_i^-\,[-\delta_i]_+\bigr).
\label{eq:asymmetric}
\end{equation}
평균적으로 양의 보정과 음의 보정이 균형을 이루면 추정값은 더 변하지 않는다.
\begin{empheq}[box=\keybox]{equation}
\alpha_i^+\,\bigl\langle[r-V_i]_+\bigr\rangle \;=\; \alpha_i^-\,\bigl\langle[V_i-r]_+\bigr\rangle ,
\qquad
\kappa_i \;=\; \frac{\alpha_i^+}{\alpha_i^++\alpha_i^-} .
\label{eq:expectile}
\end{empheq}
$\kappa_i=1/2$이면 보통의 평균이다. $\kappa_i>1/2$이면 뉴런은 “낙관론자”다. 그 추정값은 분포의 위쪽 끝으로 치우친다. $\kappa_i<1/2$이면 “비관론자”다.

예를 들어 크기 $1$인 주스 방울이 확률 $1/2$로 온다고 하자. 그러면~\eqref{eq:expectile}는 $\kappa_i\cdot\tfrac12(1-V_i)=(1-\kappa_i)\cdot\tfrac12V_i$, 곧 $V_i=\kappa_i$를 준다(그림~\ref{fig:expectiles}). $\kappa_i$가 서로 다른 뉴런 집합은 통계학에서 분위수 집합이 분포를 기록하듯 보상 분포를 기록한다.

\begin{figure}[t]
\centering
\begin{tikzpicture}[>=Stealth,x=7cm,y=3cm]
\draw[->,gray!70!black] (-0.08,0) -- (1.15,0) node[right,black] {\small 보상};
\draw[->,gray!70!black] (-0.08,0) -- (-0.08,0.75) node[left,black] {\small 확률};
\fill[blue!25] (-0.03,0) rectangle (0.03,0.5);
\fill[blue!25] (0.97,0) rectangle (1.03,0.5);
\node[below,font=\small] at (0,0) {$0$};
\node[below,font=\small] at (1,0) {$1$};
\node[left,font=\scriptsize] at (-0.08,0.5) {$1/2$};
\foreach \k/\c/\lab/\h in {0.2/blue!60!black/{비관론자, $\kappa=0.2$}/0.62, 0.5/gray!50!black/{평균, $\kappa=0.5$}/0.42, 0.8/red!70!black/{낙관론자, $\kappa=0.8$}/0.22}{
  \draw[very thick,\c] (\k,0) -- (\k,\h);
  \node[above,font=\scriptsize,\c] at (\k,\h) {\lab};}
\end{tikzpicture}
\caption{\nt{DOPA} 뉴런 집합이 기록한 보상 분포. 보상은 확률 $1/2$로 $0$ 또는 $1$이다(파란 막대). 비율 $\kappa$를 가진 뉴런은 추정값 $V=\kappa$에 도달한다~\eqref{eq:expectile}.}
\label{fig:expectiles}
\end{figure}

측정된 것: \nt{DOPA} 뉴런마다 “반전점”, 곧 반응이 급증에서 급감으로 바뀌는 보상 크기가 다르다. 그리고 양의 오차와 음의 오차에 대한 반응 비대칭이 큰 뉴런일수록 더 큰 보상에서 반전한다. 이는~\eqref{eq:expectile}가 요구하는 그대로다~\cite{Dabney2020}. 뉴런 무리의 반응으로부터 분포의 모양을 복원할 수 있다. 이 편차가 부수 효과가 아니라 주된 기능이라는 것은 가설이다.

\begin{keyblock}
\begin{proposition}[비대칭에서 나오는 분포]
\label{prop:expectile}
\nt{DOPA} 뉴런들이 양의 오차를 반영하는 비율 $\kappa_i$에서 서로 다르면, 그 추정값들은 보상 분포의 서로 다른 점~\eqref{eq:expectile}으로 수렴한다. 뉴런 무리는 평균이 아니라 분포를 전달하며, 따라서 위험도 전달한다.
\end{proposition}
\end{keyblock}

\section{오차만이 아니라 중요도}

오차 식~\eqref{eq:ctd}에 따르면 전기 충격이나 쓴맛 같은 불쾌한 사건은 급감을 일으켜야 한다. 기대보다 나빴기 때문이다. 일부 \nt{DOPA} 뉴런은 실제로 그렇게 반응하지만, 다른 뉴런들은 불쾌한 사건에 보상처럼 \emph{급증}으로 반응한다~\cite{Matsumoto2009}. 이 뉴런들이 알리는 것은 오차의 부호가 아니라 무언가 \emph{중요한} 일, 곧 예상 밖이고 강하며 주의가 필요한 일이 일어났다는 사실이다~\cite{BrombergMartin2010}.

엔지니어는 이것을 두 신호로 생각하면 편하다. 첫째는 부호 있는 오차 $\delta$로, 가중치를 어느 쪽으로 바꿀지 알려 준다. 둘째는 그 크기 $|\delta|$ 또는 사건의 새로움으로, 가중치를 \emph{얼마나 크게} 바꿀지, 곧 학습률을 알려 준다. 예상 밖의 사건 뒤에는 더 빨리 배워야 한다. 의미로 보면 이는 \ref{sec:neurontypes}장의 노르아드레날린 이득에 가깝다.

세 번째 방식도 있다. 별도의 축을 위한 별도의 전선이다. 행동 선택 영역 가운데 하나로 가는 \nt{DOPA} 뉴런들은 보상이 아니라 자극의 새로움과 세기에 반응하며, 동물이 위험을 피하도록 배우는 데 필요하다~\cite{Menegas2018}. 이 전선으로 흐르는 것은 “더 좋다, 더 나쁘다”가 아니라 “위험하다, 아니다”이다.

측정된 것: \nt{DOPA} 뉴런 무리마다 불쾌한 자극과 새로운 자극에 다르게 반응하고, 출력마다 가르치는 내용이 다르다. 뇌가 중요도 신호를 정확히 어떻게 쓰는지, 곧 학습률로 쓰는지, 주의 신호로 쓰는지, 위험 축의 별도 오차로 쓰는지는 논쟁 중이다.

\section{가르치는 것만이 아니라 재촉하는 것}

도파민에는 즉각적인 작용도 있다. 도파민이 많을수록 동물은 더 기꺼이, 더 빨리 행동한다. 이것을 학습과 어떻게 양립시킬까?

엔지니어의 답은 \emph{주파수 분할}이다. 수백 밀리초 길이의 빠른 급증과 급감은 $\delta$를 싣고 학습시킨다. 몇 분에 걸쳐 변하는 느린 수준은 다른 양, 곧 단위 시간당 평균 보상 $\bar R$를 싣는다~\cite{Niv2007}. 시간 $\tau_a$ 동안 수행한 행동이 노력 $C/\tau_a$를 요구한다고 하자. 빠를수록 비싸다. 대신 지체하는 매초는 $\bar R$, 곧 그 시간에 얻을 수 있었던 보상만큼의 비용이 든다. 총비용 $C/\tau_a+\bar R\,\tau_a$는 다음에서 최소가 된다.
\begin{empheq}[box=\keybox]{equation}
\tau_a^{*} \;=\; \sqrt{\frac{C}{\bar R}} .
\label{eq:vigor}
\end{empheq}
환경이 풍요로울수록 시간은 비싸지고 빨리 행동하는 편이 이득이다. $\bar R$가 두 배가 되면 행동 시간은 $\sqrt2\approx1.4$배 줄어든다. 여기서 $\bar R$는 \ref{sec:dofa}장에서 빼 주었던 바로 그 기대 보상이다.

목적지에서의 도파민 방출은 스파이크 발화율이 변하지 않아도 달라질 수 있다. 출력 말단의 국소 신호가 방출을 조절하기 때문이다~\cite{Mohebi2019}. 그러나 느린 성분은 스파이크 자체에도 있다. 다음 절의 실험에서는 보상에 다가갈 때의 완만한 상승이 \nt{DOPA} 뉴런의 발화율에서도 보인다~\cite{Kim2020}. 따라서 느린 수준은 스파이크 발화율과 방출의 국소 조절이라는 두 원천에서 합쳐지며, 어느 쪽이 주된지는 출력과 과제에 따라 다른 것으로 보인다.

\begin{keyblock}
\begin{proposition}[한 전선 위의 두 신호]
\label{prop:multiplex}
\nt{DOPA} 시스템은 시간 척도로 나뉜 적어도 두 가지 양을 전달한다. 학습시키는 빠른 오차 $\delta$와, 시간의 가격~\eqref{eq:vigor}을 정하고 그로써 행동 속도를 정하는 느린 수준이다. 빠른 성분과 느린 성분이 다르게 행동한다는 것은 측정되었다. 느린 성분이 정확히 $\bar R$라는 것은 가설이다.
\end{proposition}
\end{keyblock}

\section{다가갈수록 오르는 신호}

동물이 복도를 따라 먼 보상을 향해 달리면, 목표가 가까워지는 몇 초 동안 행동 선택 영역의 도파민 농도가 완만하게 오른다~\cite{Hamid2016}. \ref{sec:dofa}장에 따르면 이런 일은 없어야 한다. 모든 것이 계획대로 진행되므로 $\delta$는 0이어야 한다.

첫째 설명: 이 상승은 $\delta$가 아니라 기대 $V$ 자체, 곧 앞 절의 “목표가 가까우니 힘을 낼 만하다”라는 신호다~\cite{Hamid2016}. 둘째 설명은 $\delta$를 유지한다~\cite{Kim2020}. 기대가 정확하다면 보상 $R$로 가는 길에서 지평 $\tau_\gamma$일 때 기대는 $V=Re^{-(T-t)/\tau_\gamma}$로 자라고, 오차 식~\eqref{eq:ctd}에 따라 $\dd V/\dd t-V/\tau_\gamma=0$이다. 그런데 멀리서는 기대가 과소평가되어 있다가 다가갈수록 추정이 정확해진다고 하자. 그러면 기대는 더 가파르게, 예를 들어 $\tau_v<\tau_\gamma$인 $V=Re^{-(T-t)/\tau_v}$로 자라고,
\begin{empheq}[box=\keybox]{equation}
\delta(t) \;=\; \frac{\dd V}{\dd t}-\frac{V}{\tau_\gamma} \;=\; V\Bigl(\frac{1}{\tau_v}-\frac{1}{\tau_\gamma}\Bigr) \;>\; 0 .
\label{eq:ramp}
\end{empheq}
오차는 양수이고 $V$와 함께 커진다. 바로 그 완만한 상승이다(그림~\ref{fig:ramp}). $\tau_\gamma=4\,\text{s}$, $\tau_v=1\,\text{s}$이면 초당 $\delta=0.75\,V$가 된다. 이 가설은 가상 복도에서 동물을 목표 쪽으로 순간 이동시켜 검증되었다. 도파민은 완만한 수준이 아니라 도함수에 걸맞게 급증으로 반응했다~\cite{Kim2020}.

\begin{figure}[t]
\centering
\begin{tikzpicture}[>=Stealth,x=1.6cm,y=2.6cm]
\draw[->,gray!70!black] (-4.2,0) -- (0.5,0) node[below left,black] {\small $t$};
\draw[->,gray!70!black] (-4.2,0) -- (-4.2,1.2);
\foreach \x/\l in {-4/{$T-4$},-2/{$T-2$},0/{$T$}}{\draw[gray!60!black] (\x,-0.03) -- (\x,0.03); \node[below,font=\scriptsize] at (\x,0) {\l\,s};}
\draw[thick,densely dashed,gray!60!black] plot[domain=-4:0,samples=40] (\x,{exp(\x/4)});
\draw[very thick,blue!60!black] plot[domain=-4:0,samples=60] (\x,{exp(\x)});
\draw[very thick,red!70!black] plot[domain=-4:0,samples=60] (\x,{0.75*exp(\x)});
\node[anchor=south east,font=\small,gray!50!black] at (-1.3,0.77) {$\tau_v=\tau_\gamma$일 때 $V$: $\delta=0$};
\node[right,font=\small,blue!60!black] at (0.05,1.0) {$\tau_v=1\,\text{s}$일 때 $V$};
\node[right,font=\small,red!70!black] at (0.05,0.72) {\eqref{eq:ramp}의 $\delta$};
\end{tikzpicture}
\caption{보상에 다가갈 때의 도파민 상승을 예측 오차로 본 그림, $\tau_\gamma=4\,\text{s}$. 점선은 지평이 요구하는 그대로 자라는 기대($\delta=0$), 파란 선은 더 가파르게 자라는 기대, 빨간 선은 오차~\eqref{eq:ramp}이다.}
\label{fig:ramp}
\end{figure}

측정된 것: 완만한 상승은 도파민 농도에도, \nt{DOPA} 뉴런의 스파이크 발화율에도 있다~\cite{Kim2020}. 그리고 환경의 순간적 도약은 도파민의 도약을 일으킨다. 두 설명 중 어느 것이 옳은지, 또는 출력마다 둘 다 옳은지는 논쟁 중이다.

\section{뒤돌아보기}

위의 이론은 모두 \emph{앞}을 본다. 현대 이론 가운데 하나는 방향을 뒤집는다~\cite{Jeong2022}. 뇌가 예고 신호로 보상을 예측하도록 배우는 것이 아니라, 보상을 받은 뒤 \emph{뒤돌아본다}고 하자. 어떤 사건이 다른 사건보다 더 자주 보상에 앞섰는가? 이 연관의 척도는 두 확률의 차이다.
\begin{empheq}[box=\keybox]{equation}
M \;=\; P(\text{보상 전의 예고 신호}) \;-\; P(\text{무작위 구간의 예고 신호}) .
\label{eq:retro}
\end{empheq}
예고 신호가 보상 없이도 자주 나타나면 둘째 항이 커지고 연관은 약해진다. 이 이론에서 도파민이 알리는 것은 예측 오차가 아니라, 주어진 사건이 보상의 \emph{원인}일 가능성이 얼마나 되는가이다.

뒤돌아보기 메커니즘에는 \ref{sec:dofa}장의 3요소 규칙과 같은 것이 필요하다. 보상 순간에 그 앞에 무엇이 있었는지 읽어 낼 수 있도록, 모든 사건은 사라져 가는 흔적을 남겨야 한다. 차이는 시냅스가 아니라 \nt{DOPA}가 전달하는 내용에 있다.

이론~\eqref{eq:retro}과 오차~\eqref{eq:ctd}가 서로 다른 것을 예측하는 실험에서, 도파민 방출은 여러 경우 앞의 이론을 따랐다~\cite{Jeong2022}. 그러나 다른 실험에서는 학습 중 도파민 반응이 보상 순간에서 예고 신호 순간으로 점차 옮겨 갔다. 오차~\eqref{eq:ctd}에 의한 학습이 요구하는 그대로다~\cite{Amo2022}. 어느 이론이 뇌를 기술하는지는 아직 모른다.

\section{보상에 대한 오차만이 아니다}

마지막 확장은 오차가 \emph{무엇에 대한} 것인가에 관한 것이다. 기대한 주스를 가치는 같지만 맛이 다른 주스로 바꾸면, 가치는 변하지 않았으므로 오차~\eqref{eq:ctd}는 0이다. 그런데 \nt{DOPA} 뉴런은 이런 교체에 반응한다~\cite{Takahashi2017}. 또 인위적으로 일으킨 도파민 급증은 보상이 전혀 없어도 두 중립 자극 사이의 연관을 동물에게 가르칠 수 있다~\cite{Sharpe2017}. \nt{DOPA}는 보상뿐 아니라 \emph{무엇이} 일어날지에 대한 예측 오차도 알리는 것으로 보인다.

형식적으로 이는~\eqref{eq:ctd}의 단순한 일반화다. 보상 흐름 $r$ 하나 대신 맛, 장소, 소리 같은 사건의 특징 집합 $\varphi_k$를 두고, 각각에 대해 고유한 기대 $M_k$와 고유한 오차를 둔다~\cite{Gardner2018}.
\begin{empheq}[box=\keybox]{equation}
\delta_k(t) \;=\; \varphi_k(t) \;+\; \frac{\dd M_k}{\dd t} \;-\; \frac{M_k}{\tau_\gamma} ,
\qquad k=1,\dots,K .
\label{eq:vectortd}
\end{empheq}
보상은 특징 가운데 하나일 뿐이다. 오차 벡터는 무엇이 좋은지뿐 아니라 무엇 다음에 무엇이 오는지, 곧 세계 모델도 가르친다.

\nt{DOPA} 뉴런의 이질성도 이와 부합한다. 한 과제에서 뉴런마다 서로 다른 양을 싣는다. 어떤 뉴런은 보상과, 어떤 뉴런은 움직임과, 또 어떤 뉴런은 주의가 향하는 방향과 연관된다~\cite{Engelhard2019}.

\begin{keyblock}
\begin{proposition}[전선 대신 다발]
\label{prop:bundle}
\nt{DOPA} 시스템의 신호는 수 하나가 아니다. 그것은 뉴런에 걸쳐(분포~\eqref{eq:expectile}, 서로 다른 특징~\eqref{eq:vectortd}, 별도의 위험 축) 그리고 시간에 걸쳐(빠른 오차와 느린 수준~\eqref{eq:vigor}) 나뉘어 있다. 스칼라 오차~\eqref{eq:ctd}는 이 신호의 1차 근사다. 문턱값이 있는 가산기가 뉴런의 1차 근사인 것과 같다.
\end{proposition}
\end{keyblock}

\section{요약}

\begin{table}[t]
\centering
\footnotesize
\setlength{\tabcolsep}{4pt}
\renewcommand{\arraystretch}{1.15}
\begin{tabular}{>{\raggedright}p{2.5cm}>{\raggedright}p{3.2cm}>{\raggedright}p{3.6cm}>{\raggedright\arraybackslash}p{4.2cm}}
\toprule
이론 & 전선에 흐르는 것 & 핵심 측정 & 알려진 것 \\
\midrule
예측 오차(\ref{sec:dofa}장) & 스칼라 $\delta$~\eqref{eq:ctd} & 급증, 예고 신호로의 이동, 급감 & 측정됨; 1차 근사 \\
분포 & 비대칭이 서로 다른 $\delta_i$의 집합 & 뉴런마다 다른 반전점 & 실험과 부합; 편차의 역할은 가설 \\
중요도 & $|\delta|$, 새로움; 위험 축 & 일부 뉴런의 불쾌 자극에 대한 급증 & 측정됨; 쓰임새는 논쟁 중 \\
시간의 가격 & 느린 수준 $\bar R$ & 도파민이 행동을 빠르게 함; 느린 성분은 출력 말단의 방출과 스파이크 발화율 모두에 있음 & 빠른 성분과 느린 성분의 분리는 측정됨; $\bar R$는 가설 \\
다가갈수록 상승 & 추정이 정확해질 때의 $V$ 또는 $\delta$~\eqref{eq:ramp} & 완만한 상승, 복도에서 순간 이동 시 도약 & 상승은 측정됨; 설명은 논쟁 중 \\
뒤돌아보기 & 인과 연관의 척도~\eqref{eq:retro} & 두 이론의 예측이 갈리는 실험 & 결과가 서로 모순됨 \\
오차 벡터 & 특징별 $\delta_k$~\eqref{eq:vectortd} & 맛 교체에 대한 반응; 보상 없는 학습 & 측정됨; 벡터 형태는 가설 \\
\bottomrule
\end{tabular}
\caption{\nt{DOPA} 시스템이 전달하는 것: \ref{sec:dofa}장의 표준 그림과 그 현대적 확장.}
\label{tab:dofaalt}
\end{table}

표~\ref{tab:dofaalt}의 이론들은 서로를 배제하지 않는다. 거의 모두가 예측 오차를 기반으로 유지하면서 그 형식을 넓힌다. 예측 오차와 진지하게 경쟁하는 것은 뒤돌아보기뿐이며, 이 논쟁은 실험이 결정할 것이다. 엔지니어에게 결론은 간단하다. 전선이 실제로는 수신처가 서로 다른 여러 전선이라면, 명제~\ref{prop:broadcastcost}의 브로드캐스트 비용은 줄어든다.

\section{연습문제}

\begin{exercise}[\lvlA]
\label{ex:07:1}
\emph{주스 한 방울의 분포 점.} 크기 $1$인 주스가 확률 $p=0.2$로 오고, 그렇지 않으면 보상은 $0$이다. \eqref{eq:expectile}로 $\kappa_i=0.2$, $0.5$, $0.8$일 때의 $V_i$를 구하라. 이 보상의 중앙값은 얼마이며, 점~\eqref{eq:expectile}는 분위수와 어떻게 다른가?
\end{exercise}

\begin{exercise}[\lvlB]
\label{ex:07:2}
\emph{두 뉴런으로 분포 복원.} 보상은 $0$ 또는 $x$이고, $x$의 확률은 $p$이다. 규칙~\eqref{eq:asymmetric}을 따르는 두 뉴런이 $V(1/2)=0.25$와 $V(0.8)=0.5$를 알렸다. $p$, $x$와 보상의 표준편차를 복원하라. $\kappa=0.2$인 뉴런은 무엇을 알리겠는가?
\end{exercise}

\begin{exercise}[\lvlB]
\label{ex:07:3}
\emph{시뮬레이션: 비대칭에서 나오는 분포.} $\kappa_i=0.1,\dots,0.9$이고 규칙~\eqref{eq:asymmetric}, $\alpha_i^+=2\kappa_i$, $\alpha_i^-=2(1-\kappa_i)$를 따르는 \nt{DOPA} 뉴런 아홉 개를, $[0,1]$에서 균등한 보상으로 시뮬레이션하라. $V_i$의 흩어짐은 $\eta$에 어떻게 의존하며, 이 분포를 같은 확률의 두 방울 $0.25$와 $0.75$와 구별하려면 어떤 $\eta$가 필요한가?
\end{exercise}

\begin{exercise}[\lvlA]
\label{ex:07:4}
\emph{시간의 가격.} (a)~\eqref{eq:vigor}를 유도하고 최적점에서의 총비용을 구하라. (b)~뇌가 $\bar R$를 $k$배 틀리게 추정했다. $k=2$와 $k=4$일 때 상대 초과 비용을 구하라. 느린 도파민 수준은 $\bar R$를 얼마나 정확히 알려야 하는가?
\end{exercise}

\begin{exercise}[\lvlB]
\label{ex:07:5}
\emph{순간 이동 시의 급증.} 기대가 $\tau_v=1\,\text{s}$, $\tau_\gamma=4\,\text{s}$로~\eqref{eq:ramp}를 따라 자란다. 동물을 $T-2\,\text{s}$ 지점에서 $T-1\,\text{s}$ 지점으로 순간 이동시킨다. $\delta$ 급증의 면적을 $R$의 비율로 구하고, 그것이 이동 전 램프의 몇 초 분량을 대신하는지 구하라. 도파민이 $V$ 자체를 알린다면 순간 이동은 무엇을 보여 주겠는가?
\end{exercise}

\begin{exercise}[\lvlB]
\label{ex:07:6}
\emph{설계: 전선 위의 두 신호.} \nt{DOPA} 전선에 초당 $s=1/60$ Hz씩 오르는 수준과 스파이크 $A=3$개짜리 사건이 흐른다. 흔적 $\tau_e=1\,\text{s}$인 시냅스는 발화율에서 $\tau_f$로 평활한 그 복사본을 뺀다. 사건의 손실이 $10\,\%$ 이하이고, 흔적 시간 동안 수준의 기여가 사건 기여의 $10\,\%$ 미만이 되는 $\tau_f$는?
\end{exercise}

\begin{exercise}[\lvlB]
\label{ex:07:7}
\emph{실험 설계: 뒤돌아보기 대 오차.} 한 구간에 예고 신호가 확률 $0.1$로 나타나고, 그 뒤 보상이 확률 $0.8$로 온다. (a)~예고 신호 없는 보상을 구간당 $0.08$ 추가할 때, (b)~보상 없이 예고 신호 빈도를 두 배로 할 때 $\Delta P=P(\text{보상}\mid\text{예고})-P(\text{보상}\mid\text{없음})$와~\eqref{eq:retro}의 $M$을 비교하라.
\end{exercise}

\begin{exercise}[\lvlC]
\label{ex:07:8}
\emph{다발과 브로드캐스트 비용.} 연습문제~\ref{ex:06:8}의 모델에서 $N$개 뉴런이 저마다 전선을 가진 $K$개 그룹으로 나뉘고, 각 전선에는 다른 그룹 오차의 비율 $\rho$가 새어 든다. 학습 상수가 $N/K+2+\rho^2N(1-1/K)$임을 보여라. $K\to\infty$, $N=10^{10}$, $\rho=0.01$일 때 이는 몇 번의 시도를 주는가?
\end{exercise}

\chapter{언제 탐색하고 언제 결정할까: \nt{NORA}를 더하다}
\chaptermark{\nt{NORA}를 더하다}
\label{sec:nora}

\section{무엇이 빠져 있는가}

\ref{sec:dofa}장에서 네트워크는 잡음 덕분에 학습했다. 활동의 무작위한 편차가 탐색이었고, \nt{DOPA}는 그 편차로 결과가 나아졌는지를 알려 주었다. 그러나 거기서 정하지 않은 매개변수가 하나 남았다. \emph{얼마나} 탐색할 것인가이다. 잡음이 너무 강하면 네트워크는 갈팡질팡하며 이미 아는 것을 활용하지 못한다. 너무 약하면 매번 같은 것만 고르고, 가까운 어딘가가 더 나아졌다는 것을 알아채지 못한다. 강화 학습에서는 이를 활용과 탐색 사이의 선택이라 하며, 모든 알고리즘에는 이를 위한 노브가 있다. \ref{sec:neurontypes}장에서 우리는 뇌에서 이 노브를 돌리는 것이 \nt{NORA}라고 가정했다.

사람의 노르아드레날린 뉴런은 수만 개이고(표~\ref{tab:types}), 거의 모두 작은 무리 하나에 모여 있으며, 그 출력은 사실상 뇌 전체로 퍼진다. 다만 이 무리의 부분마다 서로 다른 영역을 어느 정도 나누어 맡는다는 사실이 밝혀졌다~\cite{Poe2020}. 이것은 \nt{DOPA}보다도 좁은 브로드캐스트 채널이다. 작동 방식은 두 가지다~\cite{AstonJones2005}. \emph{배경} 모드에서 뉴런은 초당 1회 미만에서 몇 회 사이의 발화율로 고르게 스파이크를 내며, 이 발화율은 각성 수준과 함께 천천히 변한다. \emph{돌발} 모드에서는 현재 과제에 중요한 사건에, 대략 결정을 내리는 순간에 짧은 버스트로 응답한다. 편리하지만 정확하지 않은 측정점은 동공 지름이다. 동공은 이 뉴런들의 활동과 함께 변하지만, 다른 영역의 활동~\cite{Joshi2016}과 피질로 가는 콜린성 출력~\cite{Reimer2016}과도 함께 변한다. 동공이 보여 주는 것은 \nt{NORA} 하나가 아니라 전반적인 각성 수준이다.

\section{이득}

측정된 사실은 이것이다. 노르아드레날린은 뉴런의 배경 활동을 자극에 대한 응답보다 더 강하게 억제하여 신호 대 배경 비를 높인다~\cite{Foote1975NE}. 이때 개별 뉴런 특성 곡선의 기울기에 문자 그대로 어떤 계수가 곱해진다는 것까지 실험에서 나오지는 않는다. \ref{sec:neurontypes}장에서처럼 이 효과를 전달하는 단순한 모델을 쓰자. 노르아드레날린은 수신자 전달 특성의 기울기를 바꾼다~\cite{ServanSchreiber1990}. 그러면 약한 문턱값 아래 입력(배경)은 더 작게, 강한 입력은 더 크게 나온다. 뉴런이 다음 발화율로 응답한다고 하자.
\begin{empheq}[box=\keybox]{equation}
r \;=\; F\bigl(g\,(u-\theta)\bigr), \qquad F(z)=\frac{r_{\max}}{1+e^{-z}} ,
\label{eq:gain}
\end{empheq}
여기서 $u$는 입력 전류, $\theta$는 문턱값, $g$는 모델에서 \nt{NORA}가 조절하는 이득이다. $g$가 작으면 발화율은 넓은 범위에서 입력을 매끄럽게 따라간다. 뉴런은 \emph{아날로그 증폭기}다. $g$가 크면 문턱값보다 조금만 높은 입력도 최대 발화율을, 낮은 입력은 0을 낸다. 뉴런은 \emph{비교기}, 거의 디지털 소자다(그림~\ref{fig:gaincurves}). 노브 하나가 같은 뉴런을, 전자공학에서라면 서로 다른 회로가 맡을 두 모드 사이에서 전환한다.

\begin{figure}[t]
\centering
\begin{tikzpicture}[>=Stealth,x=1.1cm,y=2.4cm]
\draw[->,gray!70!black] (-3.3,0) -- (3.4,0) node[below left,black] {\small $u$};
\draw[->,gray!70!black] (-3.3,0) -- (-3.3,1.2) node[left,black] {\small $r$};
\draw[densely dashed,gray!60!black] (0,0) -- (0,1.1);
\node[below,font=\small] at (0,0) {$\theta$};
\draw[densely dotted,gray!60!black] (-3.3,1) -- (3.2,1) node[right,black,font=\small] {$r_{\max}$};
\draw[very thick,blue!60!black] plot[domain=-3.2:3.2,samples=60] (\x,{1/(1+exp(-0.8*\x))});
\draw[very thick,orange!85!black] plot[domain=-3.2:3.2,samples=60] (\x,{1/(1+exp(-2*\x))});
\draw[very thick,red!70!black] plot[domain=-3.2:3.2,samples=120] (\x,{1/(1+exp(min(-8*\x,9)))});
\node[blue!60!black,font=\small,anchor=west] at (0.5,0.12) {작은 $g$: 증폭기};
\node[red!70!black,font=\small,anchor=east] at (-0.3,0.85) {큰 $g$: 비교기};
\end{tikzpicture}
\caption{이득 $g$의 세 값에서의 전달 특성~\eqref{eq:gain}.}
\label{fig:gaincurves}
\end{figure}

큰 이득은 잡음에 도움이 될까? 항상 그렇지는 않다. 두 입력이 $\Delta u$만큼 다르고, 어느 쪽이 큰지 가려야 한다고 하자. 잡음 $\sigma_{\text{pre}}$가 증폭기 \emph{앞}에서 입력에 더해지면 신호와 잡음이 함께 증폭되어 그 비는 변하지 않는다. 반면 잡음 $\sigma_{\text{post}}$가 증폭기 \emph{뒤}에서, 즉 \ref{sec:code}장에서처럼 믿을 수 없는 시냅스와 네트워크 자체의 무작위 스파이크에서 더해지면, 신호 대 잡음비
\begin{empheq}[box=\keybox]{equation}
d' \;=\; \frac{g\,\Delta u}{\sqrt{g^2\sigma_{\text{pre}}^2+\sigma_{\text{post}}^2}}
\label{eq:gainsnr}
\end{empheq}
는 $g\sigma_{\text{pre}}$가 $\sigma_{\text{post}}$와 비슷해질 때까지 $g$와 함께 커진다~\cite{ServanSchreiber1990}.

\begin{keyblock}
\begin{proposition}[이득이 도움이 될 때]
\label{prop:gainnoise}
이득을 높임으로써 \nt{NORA}는 증폭기 \emph{뒤}, 즉 네트워크 안에서 생기는 잡음에 대해서만 입력 구별을 개선한다. 입력과 함께 들어온 잡음은 이득으로 없앨 수 없다.
\end{proposition}
\end{keyblock}

\section{양자택일에서의 이득}

\ref{sec:gaba}장의 선택으로 돌아가자. 입력이 $I_1$, $I_2$인 두 \nt{GLUT} 뉴런 무리가 세기 $\beta$로 서로 억제한다. 이제 각 무리에 이득 $g$가 있다. 선택 방정식~\eqref{eq:wta}에서 괄호 안의 식에 $g$가 곱해진다. 두 무리가 모두 작동하는 동안 정상 상태의 발화율은 $r_1=g(I_1-\beta r_2)$, $r_2=g(I_2-\beta r_1)$에서 구한다. 두 식을 더하고 빼면 합 $r_1+r_2=g(I_1+I_2)/(1+g\beta)$와 차 $r_1-r_2=g(I_1-I_2)/(1-g\beta)$를 얻는다. 둘의 비는 네트워크가 한 입력을 다른 입력보다 얼마나 뚜렷하게 선호하는지를 나타낸다.
\begin{empheq}[box=\keybox]{equation}
\frac{r_1-r_2}{r_1+r_2} \;=\; \varepsilon\,\frac{1+g\beta}{1-g\beta}, \qquad \varepsilon=\frac{I_1-I_2}{I_1+I_2}, \quad g\beta<1 .
\label{eq:wtagain}
\end{empheq}
작은 입력 차 $\varepsilon$은 $(1+g\beta)/(1-g\beta)$배로 증폭되며, 이 인수는 $g\beta$가 1에 다가가면 한없이 커진다. $g\beta>1$이면 명제~\ref{prop:wta}에서처럼 두 무리가 모두 작동하는 상태는 불안정하다. 한쪽이 이기고 다른 쪽은 침묵한다(그림~\ref{fig:pitchfork}).

\begin{figure}[t]
\centering
\begin{tikzpicture}[>=Stealth,x=3.2cm,y=1.5cm]
\draw[->,gray!70!black] (0,0) -- (2.15,0) node[below left,black] {\small $g\beta$};
\draw[->,gray!70!black] (0,-1.25) -- (0,1.35) node[above,black] {\small $(r_1-r_2)/(r_1+r_2)$};
\draw[gray!60!black] (1,-0.04) -- (1,0.04); \node[below right,font=\small] at (1,0) {$1$};
\node[left,font=\scriptsize] at (0,1) {$1$}; \node[left,font=\scriptsize] at (0,-1) {$-1$};
\draw[very thick,blue!60!black] plot[domain=0:0.905,samples=60] (\x,{0.05*(1+\x)/(1-\x)}) -- (1,1) -- (2.05,1);
\draw[very thick,blue!60!black] (1.105,-1) -- (2.05,-1);
\draw[thick,densely dashed,gray!60!black] (1,0) -- (2.05,0);
\node[font=\small,align=center] at (0.45,-0.62) {숙고 중:\\두 무리 모두\\작동};
\node[font=\small,align=center] at (1.55,0.45) {결정함:\\한 무리만 작동};
\node[font=\small,blue!60!black,anchor=south west] at (1.05,-0.97) {약한 입력이 먼저 이겼다: 유지된다};
\end{tikzpicture}
\caption{이득에 따른 양자택일. 입력 차는 $\varepsilon=0.05$이다. $g\beta>1$이면 두 결정이 모두 안정하며(실선; 약한 입력의 승리는 $g\beta>(1+\varepsilon)/(1-\varepsilon)\approx1.1$에서), 네트워크는 이미 내린 결정을 유지한다. 대칭 상태(점선)는 불안정하다.}
\label{fig:pitchfork}
\end{figure}

\begin{keyblock}
\begin{proposition}[이득이 선택 모드를 전환한다]
\label{prop:gainwta}
상호 억제 $\beta$를 가진 선택 네트워크에서 이득 $g$는 네트워크를 경계 $g\beta=1$ 너머로 옮긴다. 경계 아래에서 네트워크는 “숙고”한다. 두 무리가 모두 작동하고 입력 차는 식~\eqref{eq:wtagain}에 따라 증폭된다. 경계 위에서 네트워크는 “결정”했다. 한 무리만 작동하고 결정이 유지된다. \nt{NORA}는 $g$를 바꿈으로써 가중치를 하나도 건드리지 않고 네트워크를 이 두 모드 사이에서 전환할 수 있다.
\end{proposition}
\end{keyblock}

\section{온도로서의 이득}

이제 선택에 네트워크 자체에서, 즉 증폭기 뒤에서 생기는 잡음을 더하자. 어느 무리가 이길지는 $g(u_A-u_B)+\eta$의 부호가 정한다. 여기서 $u_A-u_B$는 입력 차, 즉 \ref{sec:dofa}장에서 학습한 가중치의 차이고, $\eta$는 척도 $\sigma$인 잡음이다. 잡음이 로지스틱 분포를 따른다면(가우스 분포여도 곡선은 거의 같다) $A$를 고를 확률은
\begin{empheq}[box=\keybox]{equation}
P(A) \;=\; \frac{1}{1+e^{-g\,(u_A-u_B)/\sigma}} .
\label{eq:softmax}
\end{empheq}
프로그래머라면 여기서 온도 $T=\sigma/g$인 softmax 선택을 알아볼 것이다. 이득은 역온도다. $g$가 작으면 눈에 띄게 나은 선택지도 못한 선택지보다 그리 자주 뽑히지 않는다. 네트워크는 많이 탐색한다. $g$가 크면 알려진 최선의 선택지가 거의 항상 뽑힌다. 네트워크는 아는 것을 활용한다.

식~\eqref{eq:wtagain}와 \eqref{eq:softmax}의 $g$가 정확히 무엇인지 짚어 두자. 이것은 선택하는 순간 현재 과제 입력 차에 걸리는 \emph{유효} 이득이다. \nt{NORA}의 두 모드는 이 이득에 반대로 작용한다. 결정 순간의 돌발 버스트는 이득을 높여 네트워크가 선택을 굳히게 한다. 반대로 높은 배경 활동은 탐색과 함께 간다. 사람에서는 무작위 선택 직전에 기저 동공이 더 넓고~\cite{JepmaNieuwenhuis2011}, 쥐에서는 앞띠피질로 들어가는 \nt{NORA} 입력이 선택을 무작위 모드로 바꾼다~\cite{Tervo2014stochastic}. 따라서 높은 배경 \nt{NORA}에는 \emph{작은} 유효 $g$가, 버스트에는 큰 유효 $g$가 대응한다.

\section{얼마나 탐색할 것인가}

어떤 이득이 더 나을까? 모델로 확인했다. 네트워크는 식~\eqref{eq:softmax}에 따라 두 행동 가운데 하나를 고른다. 각 행동은 어떤 확률로 보상을 주며, 네트워크는 \ref{sec:dofa}장에서처럼 선택한 행동의 보상으로 그 확률을 추정한다. 때때로 세계가 바뀐다. 두 확률이 새로 정해진다. 바뀌는 것은 선택한 행동일 수도 있고, 네트워크가 오랫동안 시도하지 않은 행동일 수도 있다. 후자는 탐색으로만 알 수 있다. 그림~\ref{fig:invertedU}는 여러 고정 $g$에서 네트워크가 가능한 최선의 보상 중 얼마를 얻는지 보여 준다.

\begin{figure}[t]
\centering
\begin{tikzpicture}[>=Stealth,x=1.2cm,y=1cm]
\draw[->,gray!70!black] (0,0) -- (7.6,0) node[below left,black] {\small $g$};
\draw[->,gray!70!black] (0,0) -- (0,4.4) node[above,black,font=\small] {최선 보상 대비 비율};
\foreach \x/\l in {0/0.5,1/1,2/2,3/4,4/8,5/16,6/32,7/64}{\draw[gray!60!black] (\x,-0.06) -- (\x,0.06); \node[below,font=\scriptsize] at (\x,0) {\l};}
\foreach \v/\y in {0.8/0.8,0.9/2.4,1.0/4.0}{\draw[gray!60!black] (-0.06,\y) -- (0.06,\y); \node[left,font=\scriptsize] at (0,\y) {\v};}
\draw[very thick,blue!60!black] plot[mark=*,mark size=1.3pt] coordinates {(0,0.368) (1,0.880) (2,1.728) (3,2.784) (4,3.456) (5,3.616) (6,3.488) (7,3.264)};
\draw[very thick,red!70!black] plot[mark=*,mark size=1.3pt] coordinates {(0,0.368) (1,0.672) (2,1.184) (3,1.840) (4,2.176) (5,1.920) (6,1.456) (7,1.104)};
\node[blue!60!black,font=\small,anchor=south] at (5,3.7) {조용한 세계};
\node[red!70!black,font=\small,anchor=north] at (5.1,1.25) {빠르게 변하는 세계};
\draw[->,thick,blue!60!black] (5,3.25) -- (5,3.55);
\draw[->,thick,red!70!black] (4,1.8) -- (4,2.1);
\node[font=\small\itshape,gray!35!black,anchor=west] at (0.15,2.3) {맹목적으로 탐색};
\node[font=\small\itshape,gray!35!black] at (6.45,2.55) {변화를 놓침};
\end{tikzpicture}
\caption{고정 이득 $g$에서 가능한 최선의 보상 대비 비율($g$축은 로그 눈금). 파란 곡선은 세계가 평균 1000번 시도에 한 번, 빨간 곡선은 20번에 한 번 바뀌는 경우다. 화살표는 최선의 $g$를 표시한다. 빠르게 변하는 세계에서는 그 값이 절반이다. 4000번 시도의 실행 60회 평균.}
\label{fig:invertedU}
\end{figure}

$g=0.5$에서 네트워크는 지식을 거의 활용하지 못해 가능한 보상의 약 4분의 3을 얻고, $g$가 매우 크면 변화를 놓친다. 최선의 값은 세계가 1000번 시도에 한 번 바뀔 때 $g=16$(비율 $0.976$), 20번에 한 번 바뀔 때 $g=8$($0.886$)이다.

\begin{keyblock}
\begin{proposition}[뒤집힌 U]
\label{prop:explore}
네트워크가 모으는 보상은 이득에 대해 뒤집힌 U 모양으로 의존한다. $g$가 너무 작으면 시도를 탐색에 낭비하고, 너무 크면 변화를 알아채지 못한다. 세계가 빨리 변할수록 최선의 $g$는 작다.
\end{proposition}
\end{keyblock}

뒤집힌 U는 실험에서도 관찰된다. 동물은 \nt{NORA} 뉴런의 배경 활동이 적당하고 돌발 응답이 뚜렷할 때 과제를 가장 잘 풀고, 배경 활동이 높으면 산만해져 과제 사이를 오간다~\cite{AstonJones2005}. 이는 앞 절에서 본 두 모드의 차이와 들어맞는다. 결정 순간의 돌발 버스트는 바로 선택해야 할 때 큰 유효 $g$를 준다. 버스트 없는 높은 배경 활동은 무관한 입력까지 모든 것을 증폭하므로 현재 과제에 대한 유효 $g$가 떨어진다. 이것이 탐색이다.

\section{누가 노브를 돌리는가}

최선의 이득이 세계가 얼마나 빨리 변하는지에 달려 있다면 이득을 맞추어 조정하는 편이 좋다. 그런데 \nt{NORA} 뉴런은 세계가 바뀌었다는 것을 어떻게 알까? 자연스러운 표지는 \emph{뜻밖의 일}이다. 예측 오차가 평소보다 커지는 것이다. 여기서 두 종류의 불확실성을 구별해야 한다. 보상이 그냥 무작위라면 오차는 늘 크고, 그것은 예상된 일이다. 반면 오차가 평소 크기보다 갑자기 커졌다면 무언가가 바뀐 것이다. 한 가설에 따르면 \nt{NORA}가 알리는 것은 바로 이 예상치 못한 불확실성이고, 예상된 불확실성은 \nt{ACET}가 알린다~\cite{YuDayan2005}.

이 신호로 무엇을 할까? 이득을 낮추어 더 탐색할 수 있다. 학습률을 높여 낡은 추정을 더 빨리 잊을 수도 있다. 실험에 따르면 급격한 변화 뒤에 사람은 더 빨리 학습하고, 이때 동공이 커진다~\cite{Nassar2012}. 다시 말하지만 동공은 \nt{NORA}만이 아니라 \nt{ACET}의 지표이기도 하다~\cite{Reimer2016}. 둘을 한꺼번에 할 수도 있다. 또 다른 가설에 따르면 \nt{NORA}의 돌발 버스트는 \emph{인터럽트}로 작동한다. 네트워크가 현재 상태를 초기화하고 새 상황에 맞게 재구성하게 하는 신호로~\cite{BouretSara2005}, 프로세서의 공통 인터럽트 선과 같다.

찾지 못한 것도 솔직히 말해 두자. 모델에서 우리는 두 가지 단순한 조정 방식을 모두 시험했다. 평활한 오차가 장기 평균을 넘으면 $g$를 낮추거나 학습률을 높이는 것이다. 한동안 오래 가만히 있다가 한동안 자주 바뀌는 세계에서, 두 방식의 최적 설정은 최선 보상의 $0.926$--$0.927$을 얻었다. 최선의 고정 $g$($g=8$에서 $0.925$)나 고정이지만 충분히 큰 학습률($0.928$)과 거의 같고, 설정이 나쁘면 그보다 적었다. 그림~\ref{fig:invertedU}의 곡선은 꼭대기 근처에서 완만하므로, 이런 세계에서는 조정할 것이 거의 없다. 조정이 값을 하려면 모드마다 전혀 다른 설정이 필요해야 한다. \nt{NORA}가 어떤 제어 법칙을 구현하는지는 아직 밝혀지지 않았다.

\section{정리}

\begin{table}[t]
\centering
\footnotesize
\setlength{\tabcolsep}{4pt}
\renewcommand{\arraystretch}{1.15}
\begin{tabular}{>{\raggedright}p{3.0cm}>{\raggedright}p{4.4cm}>{\raggedright\arraybackslash}p{6.0cm}}
\toprule
\nt{NORA}가 하는 일 & 방법 & 알려진 것 \\
\midrule
뉴런 응답의 기울기를 바꾼다 & 식~\eqref{eq:gain}의 이득 $g$ & 신호 대 잡음비 증가는 측정됨; 곱셈 모델은 단순화 \\
선택을 전환한다 & 네트워크를 $g\beta=1$ 너머로 옮긴다 & 모델~\eqref{eq:wtagain}에서 나옴; 문턱값의 직접 측정은 없음 \\
얼마나 탐색할지 정한다 & $g$는 역온도~\eqref{eq:softmax}; 배경은 작은 유효 $g$(탐색), 버스트는 큰 $g$(결정) & 뒤집힌 U와 두 작동 모드는 측정됨; 탐색과의 연관은 근거가 탄탄한 가설 \\
변화를 알린다 & 예상치 못한 불확실성 & 변화 뒤 동공과 학습률 증가는 측정됨; 그것이 \nt{NORA} 신호라는 것은 가설 \\
끼어들어 초기화한다 & 뜻밖의 사건에 대한 돌발 버스트 & 중요한 사건에 대한 버스트는 측정됨; “인터럽트”는 가설 \\
\bottomrule
\end{tabular}
\caption{\nt{GLUT}, \nt{GABA}, \nt{DOPA}로 된 네트워크에 \nt{NORA}가 더하는 것.}
\label{tab:nora}
\end{table}

\nt{DOPA}는 네트워크에 가중치를 \emph{어느 쪽으로} 바꿀지 알려 준다. \nt{NORA}는 가중치를 하나도 바꾸지 않지만, 네트워크가 이미 아는 것을 \emph{어떻게} 쓸지 정한다(표~\ref{tab:nora}). 노브 하나, 즉 이득이 뉴런을 증폭기에서 비교기로, 선택 네트워크를 “숙고 중”에서 “결정함”으로, 선택을 탐색에서 활용으로 바꾼다. \nt{NORA}가 세계의 변화 속도를 어떻게 아는지는 열린 문제다.

\section{연습문제}

\begin{exercise}[\lvlA]
\label{ex:08:1}
\emph{뇌 전체에 노브 하나.} (a)~표~\ref{tab:types}로 \nt{NORA} 뉴런 하나당 뇌 뉴런이 몇 개인지 추정하라. (b)~증폭기 앞 잡음은 $\sigma_{\text{pre}}=0.5$, 뒤 잡음은 $\sigma_{\text{post}}=4$이다. 식~\eqref{eq:gainsnr}의 $d'$이 극한의 $90\,\%$에 이르는 $g$는 얼마인가?
\end{exercise}

\begin{exercise}[\lvlB]
\label{ex:08:2}
\emph{이득이 있는 선택.} $\tau\,\dd r_1/\dd t=-r_1+g[I_1-\beta r_2]_+$, $\tau\,\dd r_2/\dd t=-r_2+g[I_2-\beta r_1]_+$, $\tau=20\ms$. (a)~$g\beta=0.5$와 $0.9$에서 $r_1-r_2$가 감쇠하는 시간은? (b)~$\varepsilon=0.05$에서 두 무리가 모두 작동하는 $g\beta$의 상한과, 약한 입력의 승리가 안정해지는 $g\beta$의 하한을 구하라.
\end{exercise}

\begin{exercise}[\lvlA]
\label{ex:08:3}
\emph{잡음에서 나오는 softmax.} $K$개 무리가 저마다 굼벨 잡음 $P(\eta_k<z)=\exp(-e^{-z/\sigma})$를 받고, $gu_k+\eta_k$가 가장 큰 무리가 이긴다. $P(k)$를 구하라. $u_A-u_B=0.1$, $\sigma=1$인 두 선택지 중 나은 쪽이 $90\,\%$ 선택되는 $g$는?
\end{exercise}

\begin{exercise}[\lvlB]
\label{ex:08:4}
\emph{최선 이득의 추정.} 그림~\ref{fig:invertedU}의 모델에서 변화 뒤 보상 확률은 $[0,1]$에서 균등하다. 네트워크는 $e^{g\Delta}$번에 한 번 못한 행동을 시도한다. 이를 $1/h$와 같게 두면 $g^*\approx\ln(1/h)/\bar\Delta$이다. $\bar\Delta=\langle|p_A-p_B|\rangle$와 $h=0.001$, $0.01$, $0.05$에서의 $g^*$를 구하라.
\end{exercise}

\begin{exercise}[\lvlB]
\label{ex:08:5}
\emph{시뮬레이션: 이득과 변화 속도.} \texttt{models/\allowbreak nora\_explore.py}의 사본(현재 폴더에 파일을 쓴다)으로 $h=0.001$부터 $0.2$까지 $g^*(h)$를 구하라. 연습문제~\ref{ex:08:4}의 추정은 어디서 맞는가? 변화가 빠를 때 $g^*$가 더 줄지 않는 이유는?
\end{exercise}

\begin{exercise}[\lvlB]
\label{ex:08:6}
\emph{설계: 선택 네트워크를 위한 인터럽트.} 연습문제~\ref{ex:08:2}의 네트워크($\beta=2$, $\tau=20\ms$, $g=1$)는 입력이 이제 $I_1=0.9$, $I_2=1.0$인데도 $r_1=0.9$, $r_2=0$을 유지한다. \nt{NORA}가 시간 $T_p$ 동안 $g$를 $g_{\text{lo}}$로 낮춘다. 최소 $T_p$를 $g_{\text{lo}}=0$에서는 해석적으로, $g_{\text{lo}}=0.25$에서는 시뮬레이션으로 구하라.
\end{exercise}

\begin{exercise}[\lvlC]
\label{ex:08:7}
\emph{조정이 값을 할 때.} 오라클은 지금이 조용한 시기인지 변덕스러운 시기인지 알고 시기마다 $g$를 따로 정한다. (a)~이 장의 세계(시기당 $1000$번 시도, $h=0.001$과 $0.05$)에서 최선의 고정 $g$ 대비 이득을 시뮬레이션으로 구하라. (b)~이득이 더 큰 세계를 만들고, 뜻밖의 일에 따른 조정이 그중 얼마를 가져가는지 구하라.
\end{exercise}

\chapter{언제 쓰고 언제 읽을까: \nt{ACET}를 더하다}
\chaptermark{\nt{ACET}를 더하다}
\label{sec:acet}
\begin{chaptergoals}
\item 같은 연결에 저장하고 학습하는 메모리에 쓰기 허용 선이 필요한 이유;
\item \nt{ACET}가 입력을 강하게, 내부 연결을 약하게 하고 가중치 변화를 쉽게 하는 방식;
\item 쓰기와 읽기에 모두 알맞은 \nt{ACET} 수준이 없어 스위치가 필요한 이유;
\item 최적 필터가 입력 가중치와 학습률을 수 하나 $P/R$로 정하는 이유.
\end{chaptergoals}

\section{무엇이 빠져 있는가}

메모리 칩에는 주소선과 데이터선 말고도 선이 하나 더 있다. 쓰기 허용 신호다. 이 신호가 꺼져 있으면 메모리는 읽히기만 하고, 켜지면 데이터선에 실린 값이 기록된다. 이런 선이 없으면 메모리는 작동하지 않는다. 읽을 때마다 무언가를 함께 쓴다면, 읽은 내용이 잘못 읽은 것까지 포함해 곧바로 다시 기록되기 때문이다.

뇌에서는 이 문제가 칩보다 더 절실하다. \ref{sec:glutcomputer}장에서 메모리는 자기 연결로 켜진 상태를 유지하는 \nt{GLUT} 뉴런 무리였다(그림~\ref{fig:bistable}). \ref{sec:dofa}장에서는 이 연결이 작동하는 도중에 헤브 규칙으로 학습한다는 것을 보았다. 즉 같은 시냅스가 옛것을 저장하면서 새것도 받아들이며, 따로 떨어진 쓰기 단계는 클록과 마찬가지로 없다. 네트워크는 보이는 것을 기억해야 할 순간과 아는 것을 떠올려야 할 순간을 무엇으로 구별할까? \ref{sec:neurontypes}장에서 우리는 이 선을 \nt{ACET}가 맡는다고 가정했다. 이 장에서는 그런 선이 어떻게 작동해야 하는지, 그리고 왜 그것 없이는 안 되는지 확인한다.

\section{전선 하나의 세 가지 작용}

뇌 안에서 작동하는 아세틸콜린 뉴런은 많지 않으며, 몇 개의 작은 무리에 모여 있고, 그 출력은 피질 전체로 퍼진다. \nt{DOPA}, \nt{NORA}와 같은 브로드캐스트 채널이다. 피질의 아세틸콜린 수준은 주의 깊게 깨어 있을 때 높고 서파 수면에서 낮다~\cite{Hasselmo1999}. 도파민과 마찬가지로 신호의 의미는 수신자의 수용체가 정한다(명제~\ref{prop:receptor}). 아세틸콜린에는 채널을 여는 빠른 수용체도, 세포 안의 반응을 일으키는 느린 수용체도 있다. 여기서 중요한 작용은 세 가지다.

\emph{첫째, 내부 연결을 약하게 한다.} 후각 피질 절편 실험에서 아세틸콜린은 한 영역 안 뉴런 사이의 연결을 통한 전달을 강하게 억제하고, 바깥에서 그 영역으로 들어오는 입력은 거의 건드리지 않았다~\cite{HasselmoBower1992}.

\emph{둘째, 입력을 강하게 한다.} 아세틸콜린은 뉴런의 흥분성을 높이고 일부 입력 경로의 전달을 쉽게 한다. 감각 기관에서 온 신호가 네트워크 자체의 활동보다 커진다~\cite{Hasselmo2006}.

\emph{셋째, 가중치 변화를 쉽게 한다.} 아세틸콜린 수준이 높으면 연결이 더 쉽게 바뀐다~\cite{Hasselmo2006}.

엔지니어에게 이 세 작용은 하나의 동작이다. “읽기” 모드에서 “쓰기” 모드로의 전환이다. 네트워크는 자기 자신의 말을 듣지 않고 입력을 듣기 시작하며, 들은 것을 기억한다.

\section{완성하는 네트워크}

서로 연결된 \nt{GLUT} 뉴런 $N$개를 생각하자. 이들의 연결에 헤브 규칙으로 몇 개의 패턴, 즉 동시에 활성인 뉴런 $pN$개의 집합을 저장한다~\cite{Hebb1949}. 이런 네트워크에 패턴의 절반을 주면 연결이 나머지 절반을 흥분시킨다. 네트워크는 그림~\ref{fig:bistable}의 무리가 스스로를 유지했듯이 일부로부터 패턴을 \emph{완성}한다. 이것이 연상 메모리다. 내용의 일부가 주소 역할을 한다. \ref{sec:gaba}장에서처럼 \nt{GABA}는 활성 뉴런의 총수를 $pN$ 근처로 유지해 두 패턴이 함께 켜지지 못하게 한다.

아세틸콜린 수준을 $0$에서 $1$ 사이의 $a$로 나타내고, 그 세 작용을 모델에 곧바로 적자.
\begin{empheq}[box=\keybox]{equation}
\tau\frac{\dd r_i}{\dd t} = -r_i + F\Bigl(\tfrac{1+a}{2}\,x_i + (1-a)\sum_j W_{ij}\,r_j - g\Bigr),
\qquad
\frac{\dd W_{ij}}{\dd t} = \eta\,a\,(r_i-p)(r_j-p) .
\label{eq:acetnet}
\end{empheq}
여기서 $r_i$는 뉴런 $i$의 발화율, $x_i$는 감각 기관에서 오는 입력, $F$는 문턱값이 있는 전달 특성, $g$는 활성 뉴런이 $pN$보다 많아지면 커지는 \nt{GABA}의 전체 억제다. 인수 $\frac{1+a}{2}$는 입력 증폭, $1-a$는 내부 연결의 약화, $\eta a$는 학습률이다. 두 번째 식은 드문 패턴에 알맞은 형태로 쓴 헤브 규칙~\eqref{eq:hebb}이다~\cite{Tsodyks1988}. 두 뉴런이 모두 평소보다 활발하면 가중치가 커지고, 한쪽만 활발하면 작아진다. 가중치의 음수 부분은 언제나처럼 \nt{GABA} 중개 뉴런이 맡는다. 클록은 없다. 뉴런도 가중치도 연속적으로 변한다.

\begin{figure}[t]
\centering
\begin{tikzpicture}[>=Stealth,
  nrn/.style={circle,draw,very thick,fill=blue!6,minimum size=0.8cm,inner sep=0pt,font=\small},
  acet/.style={circle,draw,very thick,double,double distance=1.2pt,fill=orange!15,minimum size=1.35cm,inner sep=0pt,font=\scriptsize},
  bc/.style={->,thick,densely dashed,orange!85!black}]
\foreach \i/\y in {1/2.4,2/1.2,3/0}{
  \node[nrn] (N\i) at (3,\y) {$r_\i$};
  \draw[->,thick,blue!60!black] (0,\y) node[left,black] {\small $x_\i$} -- (N\i);}
\node[left,align=right,font=\small] at (-0.5,1.2) {감각\\기관에서};
\draw[->,thick,gray!60!black] (N1) to[bend left=30] (N2);
\draw[->,thick,gray!60!black] (N2) to[bend left=30] (N1);
\draw[->,thick,gray!60!black] (N2) to[bend left=30] (N3);
\draw[->,thick,gray!60!black] (N3) to[bend left=30] (N2);
\draw[->,thick,gray!60!black] (N1.east) .. controls (4.6,2.0) and (4.6,0.4) .. (N3.east);
\node[right,font=\small,gray!40!black,align=left] at (4.7,1.2) {내부\\연결 $W$};
\node[acet] (A) at (1.2,-1.9) {\nt{ACET}};
\draw[bc] (A) -- (1.6,-0.15);
\node[font=\small,orange!60!black,anchor=east] at (0.95,-0.9) {$+$ 입력};
\draw[bc] (A) -- (4.3,0.6);
\node[font=\small,orange!60!black,anchor=west] at (4.3,0.1) {$-$ 내부 연결, $+$ 학습률};
\node[right,font=\small,align=left] at (2.2,-1.9) {높은 수준: “쓰기”\\낮은 수준: “읽기”};
\end{tikzpicture}
\caption{쓰기 허용 선으로서의 \nt{ACET}. 뉴런 $r_i$는 감각 기관에서 입력 $x_i$를 받고(파란 화살표), 패턴이 저장된 내부 연결 $W$(회색)를 통해 서로를 흥분시킨다. 아세틸콜린(점선 화살표)은 식~\eqref{eq:acetnet}에서처럼 입력을 강하게 하고, 내부 연결을 약하게 하며, 가중치 변화를 쉽게 한다.}
\label{fig:acetnet}
\end{figure}

\section{쓰기와 읽기는 서로 방해한다}

쓰기와 읽기가 실제로 충돌하는 과제로 모델~\eqref{eq:acetnet}을 시험해 보자. 뉴런 $200$개로 된 네트워크가 이미 활성 뉴런 $20$개짜리 패턴 다섯 개를 저장하고 있다. 이제 여섯 번째 패턴을 기억해야 하는데, 이 패턴은 옛 패턴 하나와 절반이 겹친다. 뉴런 열 개를 공유한다. 이런 일은 늘 일어난다. 새 방은 익숙한 방과 닮았고, 새 얼굴은 옛 얼굴과 닮았다.

\emph{낮은 $a$에서의 쓰기.} 먼저 아세틸콜린의 처음 두 작용을 세 번째 작용과 떼어 놓자. 학습률은 최대로 두고, 입력 증폭과 내부 연결 약화만 바뀌게 한다. $a=0$이면 입력이 공유 뉴런 열 개를 켜고, 이들이 내부 연결을 통해 \emph{옛} 패턴의 나머지 절반을 켠다. \nt{GABA}는 둘이 동시에 켜져 있지 못하게 하므로, 자기 연결의 지원을 받는 옛 패턴이 입력에만 의지하는 새 패턴을 밀어낸다. 네트워크는 눈앞에 있는 것이 아니라 기억하는 것을 본다. 그리고 헤브 규칙은 본 것을 성실하게 기록한다.

결국 새 패턴은 전혀 기억되지 않는다. 새 패턴의 고유한 절반을 주어도 네트워크는 아무것도 떠올리지 못한다. 옛 패턴은 망가진다. 고유한 절반으로 불러내면 평균적으로 남의 뉴런 열 개 중 여섯 개를 끌고 나온다. $a=0.1$에서는 새 패턴이 기억되긴 하지만 옛 패턴과 합쳐지고, 손상은 오히려 더 심하다. 둘 다 자기 절반으로 불러내면 남의 뉴런을 약 여덟 개씩 끌고 나온다. $a=0.2$부터는 쓰기가 깨끗하다. 떠올릴 때 끼어드는 뉴런이 하나 미만이다(실행 10회 평균).

\emph{실제 학습률로 쓰기.} 이제 아세틸콜린이 식~\eqref{eq:acetnet}에서처럼 학습률도 정한다고 하자. 한 번 제시로 새 패턴이 완전히 기억되는 것은 $a\ge0.9$일 때뿐이다. $a=0.75$에서는 10분의 8이 기억되고, $a\le0.5$에서는 전혀 기억되지 않는다.

\emph{읽기.} 패턴 다섯 개가 깨끗하게 기록된 네트워크에 패턴의 절반을 주었다가 그것마저 치운다. $a\le0.2$에서는 네트워크가 패턴을 완전히 완성하고 스스로 유지한다. $a=0.3$에서는 거의 완전히 완성한다. $a\ge0.4$에서는 내부 연결이 너무 약해져 힌트가 사라진 뒤 활동이 꺼진다(그림~\ref{fig:acetsim}).

\begin{figure}[t]
\centering
\begin{tikzpicture}[>=Stealth,x=9cm,y=3.2cm]
\draw[->,gray!70!black] (0,0) -- (1.1,0) node[right,black] {\small $a$};
\draw[->,gray!70!black] (0,0) -- (0,1.2);
\foreach \x in {0,0.2,0.4,0.6,0.8,1.0}{\draw[gray!60!black] (\x,-0.02) -- (\x,0.02); \node[below,font=\scriptsize] at (\x,0) {$\x$};}
\foreach \y in {0,0.5,1}{\draw[gray!60!black] (-0.01,\y) -- (0.01,\y); \node[left,font=\scriptsize] at (0,\y) {$\y$};}
\node[rotate=90,font=\small] at (-0.1,0.6) {복원된 패턴의 비율};
\draw[very thick,blue!60!black] plot[mark=*,mark size=1.6pt] coordinates {(0,1.00) (0.1,1.00) (0.2,1.00) (0.3,0.96) (0.4,0.04) (0.5,0.00) (0.75,0.00) (1.0,0.00)};
\draw[very thick,red!70!black] plot[mark=*,mark size=1.6pt] coordinates {(0.1,0.00) (0.2,0.00) (0.3,0.00) (0.5,0.00) (0.75,0.80) (0.9,1.00) (1.0,1.00)};
\node[font=\small,blue!60!black,anchor=west,align=left] at (0.03,0.5) {읽기:\\절반으로\\패턴 복원};
\node[font=\small,red!70!black,anchor=east,align=right] at (1.0,0.35) {쓰기:\\새 패턴을\\한 번에};
\end{tikzpicture}
\caption{모델~\eqref{eq:acetnet}: 뉴런 $200$개, 패턴당 활성 뉴런 $20$개, 실행 10회 평균. 파란 점은 읽기다. 힌트를 치운 뒤 절반으로부터 패턴의 몇 분의 몇이 복원되는지를 나타낸다. 빨간 점은 쓰기다. 아세틸콜린이 학습률도 정할 때, 옛 패턴과 절반이 닮은 새 패턴이 한 번 제시 뒤 얼마나 떠올려지는지를 나타낸다. 읽기는 $a\le0.3$에서, 쓰기는 $a\ge0.9$에서 작동한다.}
\label{fig:acetsim}
\end{figure}

\begin{keyblock}
\begin{proposition}[쓰기와 읽기는 서로 다른 모드를 요구한다]
\label{prop:writeread}
같은 연결이 옛것을 저장하면서 새것을 학습하고, 아세틸콜린이 학습률도 정하는 모델~\eqref{eq:acetnet}에서는 쓰기와 읽기가 모두 똑같이 잘 되는 수준 $a$가 없다. 쓰려면 내부 연결을 약하게 해야 한다. 그러지 않으면 네트워크는 입력이 아니라 떠올린 것을 기록한다. 읽으려면 내부 연결을 강하게 해야 한다. 그러지 않으면 일부로부터 패턴을 완성하지 못한다. 스위치가 필요하며, 브로드캐스트 전선 하나가 그 역할을 할 수 있다.
\end{proposition}
\end{keyblock}

아세틸콜린이 학습률을 바꾸지 않는다면 두 일 모두에 좁은 띠 $0.2\le a\le0.3$이 쓸 만할 것이다. 그러나 그러면 네트워크는 읽을 때마다 학습하게 되고, 읽은 것을 그 오류까지 포함해 다시 기록하게 된다. 학습하는 동안 내부 연결을 억제한다는 발상 자체는 절편 측정을 바탕으로 만든 모델에서 가져왔다~\cite{HasselmoAndersonBower1992}. 위의 수치는 우리의 단순화된 모델에 관한 것이다. 뇌에서 모드의 경계가 정확히 어디에 있는지는 알려져 있지 않다.

\section{눈을 얼마나 믿을 것인가}

“쓰기/읽기” 스위치는 더 일반적인 질문의 극단적인 경우다. 입력을 얼마나 믿고 기억을 얼마나 믿을 것인가? 이 질문에는 엔지니어들이 가진 정확한 답이 있으며, 그 답은 왜 전선 하나가 모드와 학습률을 동시에 조절하는지를 밝혀 준다.

네트워크가 천천히 떠도는 양 $s(t)$를 추적한다고 하자. 1초마다 그 분산은 $q$만큼 커진다. 감각 기관은 $y(t)=s(t)+$(세기 $R$인 잡음)을 알려 준다. 연속 시간에서의 최선 추정 $\hat s$는 칼만--뷰시 필터가 준다~\cite{KalmanBucy1961}.
\begin{empheq}[box=\keybox]{equation}
\frac{\dd\hat s}{\dd t} \;=\; \frac{P}{R}\,\bigl(y-\hat s\bigr),
\qquad
\frac{\dd P}{\dd t} \;=\; q-\frac{P^2}{R} ,
\label{eq:kalman}
\end{empheq}
여기서 $P$는 추정 오차의 분산, 즉 네트워크가 자기가 아는 것을 얼마나 확신하지 못하는지다. 첫 번째 식은 뉴런 모델~\eqref{eq:glutneuron}의 막과 같은 RC 회로이지만, 시간 상수 $R/P$가 변한다. 네트워크가 확신하지 못하면 $P$가 크고 시간 상수가 작아서 추정이 입력을 빠르게 따라간다. 이것이 “쓰기”다. 확신하면 $P$가 작고, 추정은 입력을 거의 듣지 않고 기억을 붙든다. 이것이 “읽기”다.

정상 상태에서 $P=\sqrt{qR}$이고 기억의 시간 상수는 $\sqrt{R/q}$이다. 세계가 100초에 1단위씩 움직여 $q=0.01$이고 감각 잡음이 $R=1$이면, 기억은 약 10초를 유지해야 한다.

여기서 핵심은 수 하나 $P/R$가 두 가지를 한꺼번에 정한다는 점이다. 기억에 대한 입력의 가중치와, 저장된 추정이 바뀌는 속도다. 이것이 바로 식~\eqref{eq:acetnet}에서 아세틸콜린이 하는 일이다. 가설에 따르면~\cite{YuDayan2005} 아세틸콜린이 네트워크에 알리는 것이 $P$와 비슷한 양, 즉 네트워크가 미리 알고 있는 \emph{예상된} 불확실성이다. 이 채널이 늘 느린 것은 아니다. \nt{ACET} 뉴런의 일부는 보상과 처벌에 20밀리초 남짓 만에 응답하며, 강화가 뜻밖일수록 더 강하게 응답한다~\cite{Hangya2015}.

\begin{keyblock}
\begin{proposition}[입력 가중치와 학습률을 위한 노브 하나]
\label{prop:kalman}
최적 필터~\eqref{eq:kalman}에서 입력이 기억과 섞이는 가중치와 학습률은 같은 양 $P/R$이다. 그러므로 둘을 신호 하나로 조절하는 것이 합리적이다. 세계의 성질이 변하지 않으면 이 신호는 일정한 수준 $\sqrt{q/R}$에 이르며, 세계의 성질이 바뀔 때만 조정하면 된다.
\end{proposition}
\end{keyblock}

명제의 두 번째 문장은 \ref{sec:nora}장의 결과를 설명한다. 거기서 뜻밖의 일에 따라 학습률을 조정하는 방식은 최선의 고정값을 이기지 못했다. 변화의 성질이 변하지 않는 세계에서는 최선의 학습률이 곧 고정값이다. 이득은 세계가 갑자기 달라질 때 생긴다. 그러면 추정은 갑자기 입력에서 멀어지는데, $P$는 그 사실을 모른다. 올바른 행동은 $P$를 급격히 올려 쓰기 모드로 넘어가는 것이다. \ref{sec:nora}장에서 논의한 가설에 따르면 이런 \emph{예상치 못한} 변화를 알리는 것은 \nt{NORA}다~\cite{YuDayan2005}. 분업은 자연스럽다. \nt{NORA}는 세계가 바뀌었음을 알아채고, \nt{ACET}는 새 상황을 학습할 때까지 네트워크를 쓰기 모드로 유지한다.

\section{읽기 모드로서의 수면}

아세틸콜린 수준이 높은 것이 쓰기 모드라면, 낮을 때 네트워크는 무엇을 할까? 서파 수면에서는 아세틸콜린 수준이 떨어지고, 내부 연결이 억제에서 풀려나며, 감각 기관에서 오는 입력은 거의 없다. 읽기 모드의 네트워크가 힌트 없이 저장된 것을 재생한다. 활동 기록에 따르면 낮에 함께 작동한 뉴런들은 서파 수면에서 다시 함께 발화하며~\cite{WilsonMcNaughton1994}, 순서까지 같다~\cite{Skaggs1996}. 가설에 따르면~\cite{Hasselmo1999} 이 재생이 낮에 빠르게 학습한 것을 장기 기억으로 옮겨 쓴다(짧은 재생 폭발을 인위적으로 늘리면 기억이 좋아진다~\cite{FernandezRuiz2019}). 낮에는 입력에서 쓰고, 밤에는 안에서 다시 쓴다. 엔지니어에게는 낮에 빠른 버퍼에 쓰고 밤에 그것을 느린 영구 메모리로 내보내는 것과 비슷하다.

\section{정리}

\begin{table}[t]
\centering
\footnotesize
\setlength{\tabcolsep}{4pt}
\renewcommand{\arraystretch}{1.15}
\begin{tabular}{>{\raggedright}p{3.2cm}>{\raggedright}p{4.2cm}>{\raggedright\arraybackslash}p{6.0cm}}
\toprule
\nt{ACET}가 하는 일 & 방법 & 알려진 것 \\
\midrule
내부 연결을 약하게 한다 & 식~\eqref{eq:acetnet}의 인수 $1-a$ & 절편에서 측정됨 \\
입력을 강하게 한다 & 인수 $\frac{1+a}{2}$ & 흥분성과 입력 경로 전달의 증가는 측정됨 \\
학습을 쉽게 한다 & 학습률 $\eta a$ & 측정됨 \\
“쓰기/읽기”를 전환한다 & 명제~\ref{prop:writeread} & 모델에서 나옴; 모드의 직접 측정은 없음 \\
예상된 불확실성을 알린다 & 식~\eqref{eq:kalman}의 $P$ & 가설 \\
수면 중 재생을 위해 네트워크를 풀어 준다 & 서파 수면에서 낮은 $a$ & 수면 중 낮은 수준과 재생은 측정됨; 장기 기억으로의 재기록은 가설 \\
\bottomrule
\end{tabular}
\caption{\nt{GLUT}, \nt{GABA}, \nt{DOPA}, \nt{NORA}로 된 네트워크에 \nt{ACET}가 더하는 것.}
\label{tab:acet}
\end{table}

\nt{DOPA}는 네트워크에 가중치를 어느 쪽으로 바꿀지를, \nt{NORA}는 아는 것을 얼마나 과감하게 쓸지를, \nt{ACET}는 세계를 들을지 자신을 들을지를 알려 준다(표~\ref{tab:acet}). 이것이 표~\ref{tab:types}의 마지막 브로드캐스트 제어선이다. 쓰기 허용 신호다. 이것이 없다면, 패턴을 읽는 바로 그 연결에 패턴을 저장하는 메모리는 읽을 때마다 스스로를 망가뜨릴 것이다.

\section{연습문제}

\begin{exercise}[\lvlA]
\label{ex:09:1}
\emph{얼마나 기억할까.} $q=0.01$, $R=1$인 필터~\eqref{eq:kalman}는 약 $10$~s를 기억한다. (a)~센서 잡음의 진폭이 두 배가 되었을 때, (b)~세계가 백 배 빨리 떠돌게 되었을 때 $P_\infty$와 기억 $\sqrt{R/q}$를 구하라. (c)~네트워크가 1분을 기억하려면 잡음이 몇 배 커져야 하는가?
\end{exercise}

\begin{exercise}[\lvlB]
\label{ex:09:2}
\emph{변화 뒤의 쓰기 모드.} 세계가 바뀌어 $q=0.01$, $R=1$인 필터~\eqref{eq:kalman}의 불확실성이 $P_0\to\infty$로 뛰었다. (a)~$P$는 어떤 시간 상수로 $P_\infty$로 돌아가는가? (b)~몇 초 뒤에 학습률이 정상 상태 값보다 $10\,\%$ 이하로만 크게 되는가?
\end{exercise}

\begin{exercise}[\lvlB]
\label{ex:09:3}
\emph{고정 학습률.} 네트워크는 $\dd\hat s/\dd t=(y-\hat s)/T$로 학습한다. 세계는 $\dd s/\dd t=w$로 떠돌고 $y=s+n$이며, $w$와 $n$은 세기 $q$와 $R$인 백색 잡음이다. 최선의 $T$와 그때의 오차 분산을 구하라. $T$가 $2$배, $10$배 틀리면 분산은 몇 배가 되는가?
\end{exercise}

\begin{exercise}[\lvlB]
\label{ex:09:4}
\emph{쓰기의 경계는 어디인가.} \texttt{models/\allowbreak acet\_memory.py}에서 문턱값은 $\theta=0.35$, 패턴 내부 가중치는 $w=0.041$(이상적으로는 $0.045$)이다. 새 패턴은 옛 패턴과 뉴런 $m=10$개를 공유한다. 식~\eqref{eq:acetnet}에서 $a$가 얼마일 때 옛 패턴의 나머지 뉴런이 이 $m$개로부터 켜지지 않는가(문턱값은 날카롭다)?
\end{exercise}

\begin{exercise}[\lvlB]
\label{ex:09:5}
\emph{시뮬레이션: 메모리 용량.} \texttt{models/\allowbreak acet\_memory.py}의 함수 \texttt{read\_sweep}를 고쳐, $a=1$에서 패턴 $M$개($M$은 $5$부터 $100$까지)를 쓰고 $a=0$에서 처음 열 개를 절반으로부터 읽어라. 어떤 $M$에서 오류가 나타나며, 한계 $M_{\max}$는 $N$과 $p$에 어떻게 의존하는가?
\end{exercise}

\begin{exercise}[\lvlB]
\label{ex:09:6}
\emph{설계: 새지 않는 쓰기.} 학습률이 $\eta a$이면 네트워크는 읽을 때도 학습한다. $a\le0.3$에서 0이고 $a\ge0.9$에서 $\eta a$ 못지않은 단조 함수 $\eta(a)\le\eta$를 제안하라. 확인해 보라. 남의 뉴런 세 개가 섞인 힌트로 $a=0.2$에서 $20$번 읽은 뒤, 그중 몇 개가 패턴에 들어갔는가?
\end{exercise}

\begin{exercise}[\lvlC]
\label{ex:09:7}
\emph{밤에 버퍼를 어떻게 옮겨 쓸까.} (a)~수면 중 $a=0.05$이고 재생은 $0.1\,\text{s}$ 지속되며, 낮에는 $a=1$에서 $0.2\,\text{s}$이다. 학습률이 $\eta a$일 때와 연습문제~\ref{ex:09:6}의 $\eta(a)$일 때, 낮의 한 번 제시만큼 가중치를 바꾸려면 재생이 몇 번 필요한가? (b)~저장소는 버퍼보다 $100$배 느리게 학습하고 밤마다 패턴을 $0.1\,\text{s}$씩 $20$번 듣는다. 몇 밤이 필요한가?
\end{exercise}

\chapter{기계 전체}
\chaptermark{기계 전체}
\label{sec:synthesis}
\begin{chaptergoals}
\item \nt{GLUT}, \nt{GABA}, 브로드캐스트 채널 넷이 세 층의 기계 하나를 이루는 방식;
\item 식 하나가 모든 노브 $\delta$, $\tau_\gamma$, $g$, $a$를 모으는 방식;
\item 세계가 바뀔 때 노브들이 함께 동작할 수 있는 방식과 그중 가설인 부분;
\item 이 기계가 보통 컴퓨터와 다른 점과 아직 모르는 것.
\end{chaptergoals}

\section{지금까지 조립한 것}

우리는 표~\ref{tab:types}에서 뉴런 유형을 하나씩 골라, 그것이 계산 기계에 무엇을 더하는지 물었다. 규칙은 매번 같았다. 살아 있는 생물에 실제로 있는 것만 쓰고, 공통 클럭은 쓰지 않는다. 이제 부품을 하나의 기계로 조립하고 함께 어떻게 동작하는지 살펴보자.

\ref{sec:neurontypes}장의 그림은 확인되었고 더 정확해졌다. 거대한 주소 네트워크인 \nt{GLUT} 뉴런이 데이터를 나른다. \nt{GABA} 뉴런은 이 데이터의 흐름을 제어한다. 그리고 네트워크 위에는 네 개의 브로드캐스트 채널이 걸려 있고, 각 채널에는 자기 노브가 있다. \nt{DOPA}는 가중치 변화 중 무엇을 남길지 알려 주고, \nt{SERO}는 전체 수준과, 가설에 따르면 계획 지평을 유지하며, \nt{NORA}는 탐색과 결정 사이에서, \nt{ACET}는 쓰기와 읽기 사이에서 고른다(그림~\ref{fig:machine}).

\begin{figure}[t]
\centering
\begin{tikzpicture}[>=Stealth,
  blk/.style={draw,very thick,rounded corners=3pt,align=center,font=\small},
  reg/.style={draw,very thick,double,double distance=1.2pt,circle,minimum size=1.25cm,inner sep=0pt,font=\scriptsize,fill=orange!15},
  bc/.style={->,thick,densely dashed,orange!85!black}]
% sensors, bus, actions
\node[blk,fill=green!8,text width=1.7cm] (S) at (0.9,0) {센서\\\scriptsize 켜짐/꺼짐};
\draw[very thick,rounded corners=4pt,fill=blue!5] (2.6,-1.35) rectangle (9.0,1.35);
\node[font=\small] at (5.8,0.95) {\nt{GLUT} 버스: 데이터};
\node[font=\scriptsize,align=center] at (4.3,0.25) {동시성, 지연,\\논리(\ref{sec:glutcomputer}장)};
\node[font=\scriptsize,align=center] at (7.4,0.25) {기억, 부호\\(\ref{sec:code}, \ref{sec:acet}장)};
\draw[thick,red!70!black,fill=red!8,rounded corners=2pt] (2.8,-1.15) rectangle (8.8,-0.55);
\node[font=\scriptsize,red!60!black] at (5.8,-0.85) {\nt{GABA}(\ref{sec:gaba}장): 금지, 선택, 나눗셈, 리듬};
\node[blk,fill=blue!5,text width=2.0cm] (A) at (10.9,0) {행동\\선택};
\draw[->,very thick] (S) -- (2.6,0);
\draw[->,very thick] (9.0,0) -- (A);
\draw[->,very thick] (A) -- (12.6,0) node[right,font=\small] {근육};
\pic[scale=1.1] at (13.2,-0.5) {spring};
\pic[scale=1.15] at (0.4,0.85) {eyepic};
\pic[scale=1.05] at (1.35,0.85) {speaker};
% registers
\node[reg] (D) at (3.4,3.2) {\nt{DOPA}};
\node[reg] (C) at (5.6,3.2) {\nt{SERO}};
\node[reg] (N) at (7.8,3.2) {\nt{NORA}};
\node[reg] (H) at (10.0,3.2) {\nt{ACET}};
\node[font=\scriptsize,align=center,above=1pt] at (D.north) {오차 $\delta$};
\node[font=\scriptsize,align=center,above=1pt] at (C.north) {수준, $\tau_\gamma$};
\node[font=\scriptsize,align=center,above=1pt] at (N.north) {이득 $g$};
\node[font=\scriptsize,align=center,above=1pt] at (H.north) {쓰기 $a$};
\draw[bc] (D.south) -- (3.4,1.35);
\draw[bc] (C.south) -- (5.6,1.35);
\draw[bc] (N.south) -- (7.8,1.35);
\draw[bc] (H.south) -- (10.0,2.1) -- (8.8,1.35);
\draw[bc] (H.south east) to[bend left=20] (A.north east);
\draw[->,thick] (2.85,1.35) -- (2.85,2.75) -- (D.west) node[pos=0.25,left,font=\scriptsize] {$V$};
\draw[->,thick,green!45!black] (0.4,3.2) node[left,black,font=\small] {보상 $r$} -- (2.2,3.2) -- (D.west);
\pic[scale=0.9] at (-0.45,3.7) {juice};
\draw[->,thick,densely dotted,gray!50!black] ([yshift=0.45cm]D.north) to[bend left=18] node[above,font=\scriptsize,black] {놀람} ([yshift=0.45cm]N.north);
\end{tikzpicture}
\caption{기계 전체. \nt{GLUT} 버스는 데이터를 센서에서 행동 선택으로 나르고, \nt{GABA}는 그 안의 흐름을 제어한다. 이중 원은 브로드캐스트 채널이고, 파선 화살표는 그 신호를 네트워크 전체에 보내는 것을 나타낸다. 채널마다 자기 노브가 있다. \nt{DOPA}는 보상 $r$과, 버스 안의 비평자가 내놓는 기대 $V$를 받는다(\ref{sec:dofa}장). 점선 화살표는 \nt{NORA}가 비정상적으로 큰 오차를 보고 놀람을 안다는 가설이다(\ref{sec:nora}장).}
\label{fig:machine}
\end{figure}

\section{모든 노브를 담은 식 하나}

노브들을 하나의 도식적인 식으로 모을 수 있다. 이것은 새 모델이 아니라 \ref{sec:glutcomputer}--\ref{sec:acet}장 모델의 요약이다. 각 기호는 해당 장에서 가져왔다.
\begin{empheq}[box=\keybox]{equation}
\begin{aligned}
\tau\,\frac{\dd r_i}{\dd t} \;&=\; -r_i + F\Bigl(g\,\Bigl[\tfrac{1+a}{2}\,x_i + (1-a)\sum_j W_{ij}\,r_j - \sum_k M_{ik}\,q_k - \theta\Bigr]\Bigr),\\
\frac{\dd W_{ij}}{\dd t} \;&=\; \eta\,a\,\delta(t)\,e_{ij}(t),
\qquad \tau_e\,\frac{\dd e_{ij}}{\dd t} = -e_{ij} + r_i\,r_j,
\qquad \delta = r + \frac{\dd V}{\dd t} - \frac{V}{\tau_\gamma} .
\end{aligned}
\label{eq:machine}
\end{empheq}
여기서 $r_i$는 \nt{GLUT} 뉴런의 발화율, $x_i$는 센서 입력, $W_{ij}\ge0$는 이들 사이의 가중치이다(\ref{sec:glutcomputer}장). $q_k$는 \nt{GABA} 뉴런의 발화율, $M_{ik}\ge0$는 그 억제의 세기이다(\ref{sec:gaba}장). $e_{ij}$는 적격 흔적, $\delta$는 \nt{DOPA}의 오차이다(\ref{sec:dofa}장). $\tau_\gamma$는 가설상 \nt{SERO}가 정하는 지평, $g$는 \nt{NORA}의 이득(\ref{sec:nora}장), $a$는 \nt{ACET}의 수준이다(\ref{sec:acet}장).

\eqref{eq:machine}에 무엇이 없는지 보자. 클럭이 없다. 모든 것이 미분방정식으로 쓰여 있고, 뉴런은 저마다 스스로 변한다. 별도의 메모리가 없다. 계산에 쓰이는 바로 그 $W_{ij}$와 바로 그 활동 $r_i$가 기억한다. 대다수 시냅스에는 개별 보정이 없다. 그 학습은 국소적 적격 흔적과 하나의 공통된 수의 곱이다. 출력마다 따로 오차 전선이 가는 것은 운동을 가르치는 회로뿐이다(\ref{sec:motorlearning}장). 그리고 제어 신호는 $\delta$, $\tau_\gamma$, $g$, $a$의 네 종류뿐인데, 이것으로 뉴런 800억 개를 다룬다. 사실 각 신호는 하나의 수가 아니라 병렬 선로의 작은 다발이지만(명제~\ref{prop:bundle}), 다발 속 선로는 몇 개에 불과하다.

\begin{keyblock}
\begin{proposition}[아키텍처]
\label{prop:architecture}
지금 우리가 이해하는 한, 뇌라는 기계는 세 층으로 기술할 수 있다. 데이터는 \nt{GLUT} 주소 버스를 따라 흐른다. 데이터의 흐름은 주소 지정된 \nt{GABA} 뉴런이 이끌고, 제한하고, 정규화한다. 동작 모드는 몇 개의 브로드캐스트 채널이 정하는데, 각 채널은 네트워크의 넓은 영역이 공유하는 병렬 선로의 작은 다발이다. 앞의 두 층은 확고하게 확립되어 있다. 브로드캐스트 채널 사이의 역할 분담은 대부분 가설이다.
\end{proposition}
\end{keyblock}

\section{모든 유형을 표 하나에}

표~\ref{tab:synthesis}는 이 책의 모든 뉴런 유형을 정리한다. 자기 장을 갖지 못한 네 유형은 표에서 한 줄씩을 차지한다. 그중 둘은 기계의 출력을 다루는 장에서 역할을 찾았다. \nt{GLYC}은 척수의 루프에서(\ref{sec:muscles}장), \nt{ADRE}는 화학적 출력에서(\ref{sec:chemoutput}장) 역할을 한다. 나머지 둘의 계산상 역할은 단순하거나 알려져 있지 않다. 뉴런 유형을 정하지 않는 물질은 부록~\ref{sec:othermediators}에 모았다.

\begin{table}[t]
\centering
\footnotesize
\setlength{\tabcolsep}{3pt}
\renewcommand{\arraystretch}{1.15}
\begin{tabular}{>{\raggedright}p{1.3cm}>{\raggedright}p{1.0cm}>{\raggedright}p{3.9cm}>{\raggedright}p{1.5cm}>{\raggedright}p{1.8cm}>{\raggedright\arraybackslash}p{3.8cm}}
\toprule
유형 & 장 & 무엇을 계산하는가 & 시간 & 버스 & 알려진 것 \\
\midrule
\nt{GLUT} & \ref{sec:glutcomputer}, \ref{sec:code}, \ref{sec:sensors}, \ref{sec:motorlearning}, \ref{sec:dendrites} & 데이터: 동시성, 지연, 논리, 기억, 순서열 & ms & 주소 & 확립됨 \\
\nt{GABA} & \ref{sec:gaba}, \ref{sec:code}, \ref{sec:pain}, \ref{sec:motorlearning} & 흐름 제어: 금지, 선택, 나눗셈, 안정성, 리듬; 통증의 관문 & ms & 주소 & 확립됨; 발화율 나눗셈은 배경 잡음과 함께 \\
\nt{SERO} & \ref{sec:serocomputer}, \ref{sec:pain}, \ref{sec:sleep} & 수준 조절기, 평활화; 지평 $\tau_\gamma$; 통증의 음량; 수면 모드 & 초 & 체적 & 자기 억제는 측정됨; 지평은 가설 \\
\nt{DOPA} & \ref{sec:dofa}, \ref{sec:dofaalt} & 예측 오차 $\delta$; 느린 수준은 시간의 가격 & $0.1$\,s와 분 & 브로드캐스트 & $\delta$는 측정됨; 확장은 \ref{sec:dofaalt}장 \\
\nt{NORA} & \ref{sec:nora}, \ref{sec:pain}, \ref{sec:chemoutput}, \ref{sec:sleep} & 이득 $g$: 탐색이냐 결정이냐; 놀람; 장기의 “가속 페달” & 초 & 브로드캐스트 & 신호 대 잡음비 증가는 측정됨; 탐색에서의 역할은 가설 \\
\nt{ACET} & \ref{sec:acet}, \ref{sec:muscles}, \ref{sec:chemoutput}, \ref{sec:sleep} & 쓰기냐 읽기냐; 학습 속도; 근육으로의 출력; 장기의 “브레이크” & 초 & 둘 다 & 세 효과는 측정됨; 모드 전환은 가설 \\
\nt{GLYC} & \ref{sec:muscles}, \ref{sec:pain} & 척수와 뇌간의 음의 가중치: 길항근의 금지 & ms & 주소 & 확립됨 \\
\nt{HIST} & --- & 전역 신호 “깨어 있어라” & 느림 & 브로드캐스트 & 계산상 역할은 연구되지 않음 \\
\nt{ADRE} & \ref{sec:chemoutput} & 혈액을 통한 온몸의 “가속 페달”; 뇌 안에서는 미상 & 느림 & 브로드캐스트 & 뇌 밖에서는 확립됨; 뇌 안의 역할은 불분명 \\
\nt{ASPA} & --- & 약한 양의 가중치 & ms & 주소 & 역할은 논쟁 중 \\
\bottomrule
\end{tabular}
\caption{이 책의 모든 뉴런 유형: 각각이 무엇을 계산하며 그것이 얼마나 알려져 있는가.}
\label{tab:synthesis}
\end{table}

\section{노브들은 어떻게 함께 동작하는가}

지금까지 각 장은 노브 하나만 돌렸다. 이제 사건 하나를 기계 전체에 걸쳐 따라가 보자(그림~\ref{fig:timeline}). 동물은 오래전에 행동 $A$가 주스를 가져온다는 것을 배웠다. 어느 순간 세상이 바뀐다. $A$는 더 이상 주스를 가져오지 않고, 동물이 거의 시도하지 않는 행동 $B$가 주스를 가져온다.

\emph{오차.} $A$ 뒤에 주스가 오지 않자 \nt{DOPA} 뉴런은 발화 급감으로 응답한다. 오차 식~\eqref{eq:ctd}에 따라 $\delta<0$이다. 방금 $A$를 고른 시냅스들은 적격 흔적을 지니고 있으므로 약해진다.

\emph{놀람.} 급감 한 번은 흔한 우연이다. 그러나 급감이 되풀이되고 오차가 평소보다 커진다. \ref{sec:nora}장의 가설에 따르면 이것이 바로 \nt{NORA}의 신호, 곧 세상이 바뀌었다는 신호이다. 이득 $g$가 떨어지고 선택이 부드러워지며, 잡음 때문에 $B$를 시도하는 일이 잦아진다.

\emph{쓰기.} \ref{sec:acet}장의 가설에 따르면 변화 뒤에 \nt{ACET}가 올라가 네트워크를 쓰기 모드로 바꾼다. 옛 습관을 저장한 내부 연결은 약해지고, 센서 입력은 강해지며, 가중치는 쉽게 변한다.

\emph{새 습관.} $B$를 시도하자 아무도 기대하지 않은 주스가 온다. 발화 급증, $\delta>0$이 일어나고 $B$를 고른 시냅스들이 강해진다. 이런 시도가 몇 번 이어지면 기대 $V$가 새 세상을 따라잡고 오차는 다시 작아진다.

\emph{복귀.} 놀람이 지나가면 \nt{NORA}는 높은 이득을 되돌리고 선택은 다시 단호해진다. \nt{ACET}가 내려가고, 읽기 모드의 네트워크는 배운 것을 활용한다. 이 모든 동안 \nt{SERO}는, 가설에 따르면, 지평 $\tau_\gamma$를 정한다. 얼마나 먼 주스까지 기다릴 가치가 있는지를 정하는 것이다.

\begin{figure}[t]
\centering
\begin{tikzpicture}[>=Stealth,x=1.1cm,y=0.95cm]
\foreach \y/\lab in {4/{$\delta$ (\nt{DOPA})},3/{$g$ (\nt{NORA})},2/{$a$ (\nt{ACET})},1/{$B$ 선택}}{
  \draw[gray!60!black] (0,\y-0.35) -- (10,\y-0.35);
  \node[left,font=\scriptsize] at (0,\y) {\lab};}
\draw[densely dashed,gray!60!black] (2,0.4) -- (2,4.55) node[above,font=\scriptsize,black] {세상이 바뀜};
\draw[densely dashed,gray!60!black] (5,0.4) -- (5,4.55) node[above,font=\scriptsize,black] {$B$에 대한 첫 주스};
\pic[scale=0.85] at (2,5.25) {nojuice};
\pic[scale=0.85] at (5,5.25) {juice};
% delta: dips after the change, burst at the first juice for B, then back to zero
\draw[thick,red!70!black] (0,3.65) -- (2.3,3.65) -- (2.3,3.3) -- (2.5,3.3) -- (2.5,3.65) -- (3.2,3.65) -- (3.2,3.3) -- (3.4,3.3) -- (3.4,3.65) -- (4.1,3.65) -- (4.1,3.35) -- (4.3,3.35) -- (4.3,3.65) -- (5.0,3.65) -- (5.0,4.35) -- (5.2,4.35) -- (5.2,3.65) -- (5.8,3.65) -- (5.8,4.1) -- (6.0,4.1) -- (6.0,3.65) -- (6.6,3.65) -- (6.6,3.85) -- (6.8,3.85) -- (6.8,3.65) -- (10,3.65);
% gain: high, drops after surprise, recovers
\draw[thick,blue!60!black] (0,3.2) -- (3.0,3.2) -- (3.4,2.75) -- (5.8,2.75) -- (7.2,3.2) -- (10,3.2);
% ACh: low, rises after the change, falls after learning
\draw[thick,green!45!black] (0,1.75) -- (2.8,1.75) -- (3.3,2.25) -- (6.8,2.25) -- (7.8,1.75) -- (10,1.75);
% choice of B
\draw[thick,black] (0,0.7) -- (3.0,0.7) plot[domain=3:10,samples=40] (\x,{0.7+0.6/(1+exp(-(\x-5.3)*1.8))});
\draw[->,gray!70!black] (0,0.2) -- (10.2,0.2) node[below left,font=\scriptsize,black] {시간};
\end{tikzpicture}
\caption{기계 전체를 거쳐 추적한 사건 하나. 계산 결과가 아니라 \ref{sec:dofa}--\ref{sec:acet}장의 모델과 가설에 따른 정성적 도식이다. 변화 뒤에 \nt{DOPA}는 발화 급감으로 응답한다. 가설에 따르면 급감이 되풀이되면 놀람 신호가 올라가고, \nt{NORA}는 이득을 낮추며, \nt{ACET}는 쓰기를 켠다. $B$에 대한 첫 주스는 급증을 일으키고, 선택은 $B$로 옮겨 가며, 노브들은 제자리로 돌아온다.}
\label{fig:timeline}
\end{figure}

어느 노브도 다른 노브가 무엇을 하는지 모른다는 점에 주목하자. 각 노브는 자기 특징, 곧 오차, 오차의 비정상적인 크기, 불확실성에 반응할 뿐이다. 그런데도 동작의 순서는 저절로 짜인다. 입력이 준비되면 각 블록이 동작하는 비동기 회로와 같다. 우리가 아는 한, 이런 협조된 동작 전체는 아직 검증된 적이 없다. 노브들은 하나하나 측정되었지만, 그 공동 동작은 가설이다.

\section{이 기계가 컴퓨터와 다른 점}

\begin{table}[t]
\centering
\footnotesize
\setlength{\tabcolsep}{4pt}
\renewcommand{\arraystretch}{1.15}
\begin{tabular}{>{\raggedright}p{2.2cm}>{\raggedright}p{4.2cm}>{\raggedright\arraybackslash}p{6.6cm}}
\toprule
& 일반 컴퓨터 & 뇌 \\
\midrule
시간 & 공통 클럭, 약 1기가헤르츠 & 공통 클럭이 없다; 뉴런은 저마다 스스로 변하고, 창은 사건과 국소 리듬이 정한다(\ref{sec:glutcomputer}, \ref{sec:gaba}, \ref{sec:code}장) \\
소자 & 신뢰할 수 있는 이진 게이트 & 아날로그 문턱 소자; 시냅스는 스파이크 대부분을 잃는다(\ref{sec:code}장) \\
연결의 부호 & 무엇이든 & 송신 뉴런이 정한다(\ref{sec:neurontypes}장) \\
기억 & 프로세서와 분리 & 계산과 같은 연결 안에, 그리고 자기 유지 활동 안에(\ref{sec:glutcomputer}, \ref{sec:acet}장) \\
학습 & 오차 역전파 & 국소 규칙, 하나의 브로드캐스트 수 $\delta$(\ref{sec:dofa}장), 운동 회로 출력으로 가는 오차 전선(\ref{sec:motorlearning}장) \\
제어 & 레지스터와 명령 버스 & 네 개의 브로드캐스트 채널, 각각 몇 개 선로의 다발(\ref{sec:serocomputer}, \ref{sec:dofa}--\ref{sec:acet}장) \\
에너지 & 프로세서 하나에 수십--수백 와트 & 뇌 전체에 약 $20$와트, 뉴런당 평균 초당 약 한 번의 스파이크(\ref{sec:code}장) \\
\bottomrule
\end{tabular}
\caption{뇌와 일반 컴퓨터.}
\label{tab:computer}
\end{table}

표~\ref{tab:computer}는 주요 차이를 정리한다. 각 차이는 자연이 미처 고치지 못한 결함이 아니라 하나의 설계를 이루는 부분이다. 공통 클럭이 없으니 “켜짐---꺼짐” 센서와 동시성 검출기가 필요하다. 신뢰할 수 없는 소자에는 많은 뉴런의 중복이 필요하다. 계산과 합쳐진 기억에는 쓰기가 읽기를 망치지 않도록 \nt{ACET}가 필요하다. 역방향 전선 없는 학습에는 \nt{DOPA}와 잡음이 필요하다. 그리고 초당 스파이크 한 번이라는 에너지 한계 때문에 언어 전체가 인색해져서 스파이크를 새 소식에만 쓴다.

\section{우리가 모르는 것}

이 요약은 우리가 모르는 것의 목록으로 끝내는 것이 가장 정직하다. 항목 하나하나가 엔지니어를 위한 과제이다.

\emph{부호.} 우리는 스파이크 채널의 용량을 알고, 수신자가 메시지의 어느 부분을 읽을지 고른다는 것도 안다(\ref{sec:code}장). 그러나 피질이 스파이크의 정확한 시각을 얼마나 활용하는지, 뉴런에게 “단어”가 있는지는 알려져 있지 않다.

\emph{여러 층을 거치는 학습.} 브로드캐스트 신호는 느리게 가르치고, 결과에 영향을 주는 뉴런이 하나 늘 때마다 더 느려진다(명제~\ref{prop:broadcastcost}). 그렇다면 피질이 다층 회로의 수십억 가중치를 어떻게 조정하는지는 알려져 있지 않다. 층 사이의 오차를 스파이크 버스트가 나르는 모델도 있다~\cite{Payeur2021}. 답의 일부는 뉴런 자신의 가지 안에서 이루어지는 계산에 있을지도 모른다. 거기서는 뉴런 하나가 작은 다층 네트워크처럼 행동한다. 피질 뉴런 하나의 모델 응답을 재현하려면 인공 네트워크에 다섯--여덟 층이 필요하고~\cite{Beniaguev2021}, 사람 뉴런의 가지는 배타적 논리합(XOR)까지 계산할 수 있다~\cite{Gidon2020}. 또 다른 후보는 자기 지도 학습이다. 피질이 자기 입력을 예측하도록 배운다면 오차는 각 층에서 그 자리에서 생기고, 공통 신호는 필요 없다(\ref{sec:dofa}장)~\cite{Keller2018}.

\emph{브로드캐스트 채널이 실제로 전하는 것.} \nt{DOPA}에는 표준 그림과 그 여섯 가지 확장이 있다(\ref{sec:dofaalt}장). \nt{NORA}와 \nt{ACET}에는 그럴듯한 가설이 있다(\ref{sec:nora}, \ref{sec:acet}장). \nt{SERO}에는 대부분 추측뿐이다.

\emph{시간 맞추기.} 우리는 함수를 계산하고, 사건으로부터 시간을 재고, 국소 리듬으로 흐름을 창으로 자르는 법을 보았다. 그러나 공통 클럭이 없는 거대한 네트워크가 멀리 떨어진 영역의 작업을 어떻게 맞추는지는 일부만 이해된다.

\emph{에너지.} 모든 것에 20와트. 부호가 왜 희소해야 하는지는 이해하지만(명제~\ref{prop:economy}), 같은 일을 같은 전력으로 해내는 기계는 아직 아무도 만들지 못한다~\cite{MehonicKenyon2022}.

뇌는 엔지니어가 조립하지는 않았지만, 어떤 엔지니어라도 알아볼 법칙에 따라 조립된 계산 기계이다. RC 회로, 문턱값, 되먹임, 필터, 버스, 브로드캐스트 레지스터. 그 회로도를 우리는 일부만 읽었다. 나머지는 회로도를 읽을 줄 아는 사람들이 읽을 것이다.

\section{이 책의 다음 내용}

이 장은 기계의 계산 핵심을 조립했다. 나머지 장은 그 바깥으로 나간다. \ref{sec:sensors}장과 \ref{sec:recognition}장은 입력을 다룬다. 센서가 빛, 소리, 분자를 어떻게 스파이크로 바꾸고, 기계가 그 속에서 사물을 어떻게 알아보는지 다룬다. \ref{sec:muscles}--\ref{sec:chemoutput}장은 출력과 그것을 떠받치는 것을 다룬다. 근육, 경보 신호로서의 통증, 개별 오차 전선에 의한 운동 학습, 몸으로 가는 느린 화학적 출력이다. \ref{sec:debugging}장은 뇌-컴퓨터 인터페이스를 포함해 기계를 밖에서 디버깅하는 법을 다룬다. \ref{sec:eetypes}장과 \ref{sec:dendrites}장은 뉴런을 다시 분류하는데, 이번에는 전기적 기능과 입력의 수에 따라 나눈다. \ref{sec:sleep}장은 세상과 연결이 끊겼을 때 기계가 하는 일을, \ref{sec:livingmachine}장과 \ref{sec:dishbrain}장은 칩 위의 살아 있는 뉴런으로 조립한 기계를 다룬다.

\section{연습문제}

이 장의 연습문제는 여러 장의 결과를 잇는다. 각 문제는 적어도 두 장에 기댄다.

\begin{exercise}[\lvlA]
\label{ex:10:1}
\emph{버스와 레지스터.} 네 브로드캐스트 채널은 각각 약 다섯 선로의 다발이다(명제~\ref{prop:bundle}). 각 선로는 초당 한 번 변하고 네 단계로 구별된다. \nt{GLUT} 버스(발화율~\eqref{eq:energy}로 발화하는 뉴런 $8.6\cdot10^{10}$개, \eqref{eq:spikeentropy}에 따른 정밀도 $1\ms$)가 1초에 전송하는 것을 제어 채널로 보내려면 몇 년이 걸리는가?
\end{exercise}

\begin{exercise}[\lvlB]
\label{ex:10:2}
\emph{한 루프 위의 두 노브.} 루프~\eqref{eq:ei}에서 \eqref{eq:machine}의 노브는 흥분 쪽에 있다: $\tau_E\dot E=-E+g[(1-a)w_{EE}E-w_{EI}I+h]$, \nt{GABA}는 노브의 제어를 받지 않는다. 그림~\ref{fig:ei}의 수치에서 \nt{NORA} 버스트가 $g$를 $6$까지 올린다. 루프가 안정한 가장 작은 $a$는? \nt{NORA}와 \nt{ACET} 노브를 독립적으로 돌릴 수 있는가?
\end{exercise}

\begin{exercise}[\lvlB]
\label{ex:10:3}
\emph{브로드캐스트의 대가는 어디서 오는가.} $N$개 뉴런의 출력 $y_k=f(u_k)+\xi_k$에 분산 $\sigma^2$의 독립 가우스 잡음이 있고, $R\approx R_0+\sum_kR'_k\xi_k$이며 $R'_k$는 모두 같고, \nt{DOPA}는 $\delta=R-R_0$를 보낸다. 한 번의 시도에서 보정 $\delta\,\xi_jx_j$의 신호 대 잡음비를 구하라. 시도 횟수는 $N$에 따라 어떻게 늘어나는가?
\end{exercise}

\begin{exercise}[\lvlB]
\label{ex:10:4}
\emph{소리와 섬광 묶기.} 소리는 $5\ms$ 뒤에, 섬광은 $30\ms$에 첫 스파이크 지연~\eqref{eq:latency}($\tau=10\ms$, $U$는 $1.2\theta$에서 $4\theta$까지)을 더한 뒤에 버스에 도착한다. 각 입력의 배경 발화는 초당 스파이크 $5$개이다. 모든 대비에 대해 소리 선로의 지연과 가장 작은 동시성 창을 고르라. 배경은 초당 몇 번의 거짓 묶기를 일으키는가?
\end{exercise}

\begin{exercise}[\lvlB]
\label{ex:10:5}
\emph{시뮬레이션: 누가 변화를 알아채는가.} 비평자는 $y=s+\text{잡음}$($R=1$)을 본다. 확률 $1/200$로 $s$가 분산 $J^2=4$인 새 값으로 도약한다. $q=0$이고, $E\leftarrow0.9E+\delta^2/(P+R)$가 $20$을 넘으면 $P$를 $J^2$로 올리는 칼만 필터를 시뮬레이션하라. 그 오차는 가장 좋은 고정 $\alpha$일 때보다 얼마나 작은가?
\end{exercise}

\begin{exercise}[\lvlA]
\label{ex:10:6}
\emph{습관을 바꾸는 데 몇 번의 시도가 드는가.} 네트워크는 $Q_A=0.8$, $Q_B=0.2$를 배웠고 \eqref{eq:softmax}에 따라 $\sigma=1$, $g=8$로 고른다. 이제 $A$는 확률 $0.2$로 주스를 준다. $Q_A\leftarrow Q_A+\alpha(r-Q_A)$이고 $Q_B$는 변하지 않는다. $\alpha=0.1$과 $0.5$일 때 $A$를 몇 번 고르고 나면 $A$를 고를 확률이 $0.6$까지 떨어지는가? \nt{NORA}가 $g$를 $2$로 낮추면?
\end{exercise}

\begin{exercise}[\lvlC]
\label{ex:10:7}
\emph{전선 대신 다발.} 출력 $y_k=w_k+\xi_k$, 보상 $R=-\sum_ky_k^2$이다. 다발의 한 선로가 $G$개 뉴런의 그룹에 그 출력의 오차를 보내고, $w_k\leftarrow w_k+\eta\,\delta_l\xi_k$이다. 가장 좋은 $\eta$에서 $\sum w_k^2=\varepsilon\sum w_k^2(0)$에 이르려면 $\approx(G+2)\ln(1/\varepsilon)$번의 시도가 필요함을 보여라. $\varepsilon=0.01$일 때, 선로 다섯 개를 가진 뉴런 $10^6$개의 회로가 3초에 한 번 시도한다면 학습에 얼마나 걸리는가?
\end{exercise}

\chapter{기계의 입력: 눈, 귀, 그 밖의 센서는 어떻게 연결되는가}
\chaptermark{기계의 입력}
\label{sec:sensors}
\begin{chaptergoals}
\item 변환기로서의 감각 기관: 수용체 게이트, 이득 조절, 스파이크 인코더;
\item 광자의 산탄 잡음이 어둠 속 시각을 제한하는 이유와 어스름에 더 느리게 보는 이유;
\item 자동 이득 조절이 거대한 입력 범위를 스파이크에 담아 센서가 변화를 알리게 하는 방식;
\item 눈이 영상을 백 배로 압축하는 방식과 귀가 필터 뱅크로 일하는 방식.
\end{chaptergoals}

\section{변환기로서의 센서}

그림~\ref{fig:machine}에서 데이터는 왼쪽의 “센서” 블록으로부터 기계에 들어왔다. 이제 그 블록을 열어 보자. 엔지니어에게 모든 감각 기관은 아날로그 입력단과 끝에 아날로그-디지털 변환기를 갖춘 \emph{변환기}이다. 다만 출력되는 디지털 값은 숫자가 아니라 스파이크이다.

감각 기관의 회로는 모두 같다(그림~\ref{fig:transducer}). 빛, 압력, 진동, 분자 같은 물리량이 감각 세포 막의 수용체 단백질에 작용한다. 이 단백질은 게이트처럼 동작한다. 스스로, 또는 짧은 반응 사슬을 거쳐 이온 채널을 연다. \emph{수용체 전류}가 생기고, 그와 함께 막에 아날로그 전압이 생긴다. 이 전압은 이득 조절을 거쳐 스파이크 인코더에 들어간다. 인코더는 이미 익숙한 적분 뉴런~\eqref{eq:glutneuron}으로, 수준을 발화율과 스파이크 시각으로 바꾼다. 출력은 거의 언제나 같다. 눈, 귀, 후각, 피부의 감각 뉴런은 글루탐산을 방출하므로, 입력은 곧바로 \nt{GLUT} 버스에 연결된다.

\begin{figure}[t]
\centering
\begin{tikzpicture}[>=Stealth,
  blk/.style={draw,very thick,rounded corners=2pt,align=center,font=\footnotesize,minimum height=1.0cm,text width=1.9cm,fill=blue!5}]
\node[align=center,font=\small] (P) at (0.1,0) {빛, 소리,\\압력,\\분자};
\node[blk] (R) at (3.0,0) {수용체\\게이트};
\node[blk] (G) at (6.5,0) {이득\\조절};
\node[blk] (E) at (10.0,0) {스파이크\\인코더};
\node[align=center,font=\small] (B) at (13.4,0) {\nt{GLUT}\\버스};
\draw[->,very thick] (P) -- (R);
\foreach \ic/\xx in {sun/-1.1,speaker/-0.3,press/0.5,molecule/1.3}{\pic[scale=0.85] at (\xx,1.05) {\ic};}
\draw[->,very thick] (R) -- node[above,font=\scriptsize] {전류} (G);
\draw[->,very thick] (G) -- node[above,font=\scriptsize] {수준} (E);
\draw[->,very thick,red!70!black] (E) -- node[above,font=\scriptsize,black] {스파이크} (B);
\draw[decorate,decoration={brace,amplitude=5pt,mirror},gray!60!black] (1.8,-0.8) -- (7.7,-0.8) node[midway,below=5pt,font=\scriptsize,black] {아날로그 부분};
\draw[decorate,decoration={brace,amplitude=5pt,mirror},gray!60!black] (8.8,-0.8) -- (11.2,-0.8) node[midway,below=5pt,font=\scriptsize,black] {스파이크};
\end{tikzpicture}
\caption{변환의 사슬로서의 감각 기관. 수용체 단백질이 채널을 열어 전류를 만들고, 이득 조절이 그것을 환경에 맞추며, 인코더~\eqref{eq:glutneuron}가 수준을 스파이크로 바꿔 \nt{GLUT} 버스로 보낸다. 눈과 귀의 첫 단들은 스파이크를 전혀 내지 않는다. 신호를 멀리 보내야 할 때까지 신호는 아날로그로 남는다.}
\label{fig:transducer}
\end{figure}

전자 엔지니어에게 익숙한 세부가 하나 있다. 눈과 귀에서 맨 처음 세포들은 스파이크를 전혀 내지 않는다. 보드 위의 아날로그 단처럼 이웃에게 매끄러운 전압을 넘긴다. 스파이크는 신호를 멀리 날라야 하는 곳에서만 나타난다. 아날로그 신호는 수 밀리미터만 가도 감쇠하고 잡음이 끼지만, 스파이크는 \ref{sec:neurontypes}장에서 보았듯이 같은 높이로 전선 끝까지 간다. 디지털화는 긴 선로가 시작되는 곳에서 일어난다.

\section{물리의 한계에서}

뇌의 센서는 감도의 물리적 한계 근처에서 동작한다. 어둠 속에서 빛 센서는 \emph{광자 하나}에 응답한다. 빛 입자 하나를 흡수하면 약 1피코암페어의 전류가 생기는데, 이는 센서 자체의 잡음 위로 드러난다~\cite{Baylor1979}. 소리 센서는 자기 감지 다발이 1나노미터보다 작게 움직여도 알아챈다~\cite{Hudspeth1989}.

빛이 적을 때 정밀도를 제한하는 것은 센서가 아니라 빛 자체이다. 광자는 무작위로, 푸아송 흐름으로 도착한다. 적분 시간 $T$ 동안 센서에 평균 $N=\Phi T$개의 광자가 들어오면, 그 수는 $\sqrt N$만큼 요동한다. 증가분 $\Delta N$이 눈에 띄려면 이 요동보다 몇 배 커야 하므로, 구별할 수 있는 밝기 변화의 비율은 다음과 같다.
\begin{empheq}[box=\keybox]{equation}
\frac{\Delta I}{I} \;\approx\; \frac{k}{\sqrt{\Phi\,T}} .
\label{eq:photonnoise}
\end{empheq}
이것은 스파이크 수에 대한 식~\eqref{eq:rateerror}와 같은 식이다. 광자의 산탄 잡음과 스파이크의 산탄 잡음은 같은 법칙을 따른다. 어둠 속에서 눈이 정밀도를 높이는 방법은 광자를 더 오래 모으는 것 하나뿐이고, 실제로 적분 시간은 수십분의 1초까지 늘어난다. 그래서 어스름 속에서는 우리가 더 느리게 본다.

\section{동적 범위: 이득 조절}

센서의 가장 어려운 공학 문제는 범위이다. 조도는 별이 뜬 밤에서 햇빛 비치는 눈밭까지 약 열 자릿수 변하고, 소리의 세기는 가청 문턱에서 통증 문턱까지 열두 자릿수 변한다. 반면 스파이크 발화율은 0에서 초당 수백까지, 세 자릿수가 안 되게 변한다. 이런 입력을 이런 출력에 선형으로 담을 수는 없다.

해법은 여느 수신기와 같다. \emph{자동 이득 조절}이다. 눈의 감각 세포 응답은 다음 식으로 잘 기술된다.
\begin{empheq}[box=\keybox]{equation}
r \;=\; r_{\max}\,\frac{I}{I+\sigma} ,
\label{eq:nakarushton}
\end{empheq}
여기서 $\sigma$는 응답이 절반에 이르는 밝기이다~\cite{NakaRushton1966}. \eqref{eq:nakarushton} 자체는 두세 자릿수만 담는다. 그러나 $\sigma$는 일정하지 않다. 밝은 배경에서 오래 동작하면 평균 밝기를 따라 커진다. 응답 곡선은 밝기의 로그 축을 따라 이동하며, 매번 가파른 구간을 지금 입력이 있는 곳에 놓는다(그림~\ref{fig:adaptation}). 이것은 \ref{sec:gaba}장의 분로 억제와 같은, 전체 활동에 의한 나눗셈이다~\cite{Carandini2012}. 다만 여기서 분모는 지난 몇 초 동안의 평균 밝기이다.

\begin{figure}[t]
\centering
\begin{tikzpicture}[>=Stealth,x=1.25cm,y=2.8cm]
\draw[->,gray!70!black] (-2,0) -- (6.4,0) node[below left,font=\small,black] {$\log_{10} I$};
\draw[->,gray!70!black] (-2,0) -- (-2,1.15) node[left,font=\small,black] {$r/r_{\max}$};
\foreach \u in {-2,0,2,4,6}{\draw[gray!60!black] (\u,-0.02) -- (\u,0.02); \node[below,font=\scriptsize] at (\u,0) {$\u$};}
\draw[densely dashed,gray!60!black] (-2,0.5) -- (6.2,0.5);
\foreach \s/\c/\lab in {0/blue!60!black/{어두운 배경},2/green!45!black/{중간},4/red!70!black/{밝은 배경}}{
  \pgfmathsetmacro{\lo}{max(-2,\s-3)}
  \draw[very thick,\c] (-2,0) -- (\lo,0) plot[domain=\lo:6.2,samples=80] (\x,{1/(1+10^(\s-\x))});
  \node[font=\scriptsize,\c,anchor=west] at (\s+0.15,0.42) {\lab};}
\foreach \s/\ic in {0/moon,2/cloud,4/sun}{\pic at (\s+0.6,0.64) {\ic};}
\end{tikzpicture}
\caption{\eqref{eq:nakarushton}에 따른 빛 센서의 이득 조절. 각 곡선은 밝기의 두세 자릿수를 담는다. 더 밝은 배경으로 옮겨 가면 $\sigma$가 커지고 곡선은 로그 축을 따라 이동한다. 이동하는 곡선들이 함께 전체 범위를 덮는다.}
\label{fig:adaptation}
\end{figure}

이득 조절에는 \ref{sec:code}장에서 본 대가가 따른다. 센서는 밝기가 아니라 \emph{배경에 대한} 밝기를 알린다. 일정한 입력은 점차 응답을 일으키지 못하게 되고, 변화는 매번 응답을 일으킨다. 기계의 입력은 처음부터 수준이 아니라 변화를 말하며, 여기에 명제~\ref{prop:economy}와 같은 스파이크 절약이 있다~\cite{Barlow1961}. \ref{sec:glutcomputer}장의 켜짐 센서와 꺼짐 센서도 여기서 나온다~\cite{Schiller1992}. 한쪽은 더 밝아졌다고, 다른 쪽은 더 어두워졌다고 알린다.

엔지니어들은 이 기법을 \emph{이벤트 카메라}에서 재현했다. 이런 카메라는 고정된 주기로 프레임을 찍지 않는다. 각 픽셀이 공통 클럭 없이 스스로, 밝기의 로그가 정해진 비율만큼 변하면 “더 밝음” 또는 “더 어두움” 이벤트를 낸다~\cite{Lichtsteiner2008}. 이것은 \ref{sec:glutcomputer}장의 켜짐---꺼짐 이중 레일 부호를 실리콘으로 구현한 것 그대로이며, 일반 이미지 센서로는 얻을 수 없는 범위와 속도를 카메라에 준다~\cite{Gallego2022}.

\section{눈: 전선 하나로의 압축}

눈에는 약 $10^8$개의 빛 센서가 있지만~\cite{Curcio1990}, 영상을 뇌로 나르는 출력 뉴런은 약 백만 개이다. 신호가 눈을 떠나기 전에 약 백 배로 압축된다. 주요 압축 기법은 우리에게 익숙하다. 각 출력 뉴런은 작은 점의 밝기와 그 주변 밝기의 차이에 응답한다. \ref{sec:gaba}장에서처럼 억제를 통한 이웃 빼기이다. 영상의 고른 부분은 거의 비용이 들지 않고, 가장자리와 변화가 전송된다. 출력 뉴런은 전문화되어 있다. 어떤 것은 밝아짐을, 어떤 것은 어두워짐을, 또 어떤 것은 특정 방향의 움직임을 알린다(그림~\ref{fig:direction}). 생쥐에서는 이런 출력 채널이 서른 개 넘게 세어졌다~\cite{Baden2016}.

이렇게 만들어진 케이블의 대역폭은 얼마일까? 기니피그의 출력 뉴런 기록으로부터 $\sim10^5$개의 뉴런에서 초당 약 $875$천 비트가 측정되었다. 이를 사람의 출력 뉴런 백만 개로 환산하면 초당 $10^7$비트 정도가 된다~\cite{Koch2006}. 옛날 10메가비트 네트워크 케이블과 같다. 뉴런당 평균 발화율이 초당 스파이크 몇 개라면, 이는 스파이크당 몇 비트로, \ref{sec:code}장에서 얻은 것과 같다.

\section{귀: 필터 뱅크}

귀는 다른 문제를 푼다. 소리를 주파수별로 분해해야 한다. 엔지니어라면 대역 통과 필터 뱅크를 달 것이고, 귀는 바로 그 일을 기계적으로 한다. 소리 진동은 탄성 칸막이를 따라 달리는데, 칸막이의 성질은 길이를 따라 매끄럽게 변하고, 각 구간은 자기 주파수에 공진한다. 그 구간에 앉은 센서는 자기 대역에만 응답한다(그림~\ref{fig:filterbank}). 주파수는 위치에 따라 대략 로그로 배치된다. 옥타브마다 비슷한 몫의 센서가 할당된다. 동작한 센서의 번호가 곧 주파수이며, 이것이 \ref{sec:code}장의 위치 부호이다.

\begin{figure}[t]
\centering
\begin{tikzpicture}[>=Stealth,x=1.35cm,y=2.2cm]
\draw[->,gray!70!black] (-0.3,0) -- (7.6,0) node[above left,font=\small,black] {주파수, Hz};
\draw[->,gray!70!black] (-0.3,0) -- (-0.3,1.15) node[left,font=\small,black] {응답};
\foreach \k/\f in {0/125,1/250,2/500,3/1000,4/2000,5/4000,6/8000}{
  \draw[gray!60!black] (\k,-0.02) -- (\k,0.02); \node[below,font=\scriptsize] at (\k,0) {\f};
  \draw[thick,blue!60!black] plot[domain={max(\k-1.2,-0.3)}:{\k+0.7},samples=60] (\x,{1/(1+((\x-\k)/(\x<\k?0.45:0.2))^4)});}
% a piano keyboard: seven white keys per octave, like the octaves on the axis
\foreach \i in {0,...,48}{\draw[gray!50!black,fill=white,line width=0.3pt] ({\i/7},-0.44) rectangle ({(\i+1)/7},-0.26);}
\foreach \o in {0,...,6}{\foreach \j in {0,1,3,4,5}{\fill[black] ({\o+(\j+1)/7-0.035},-0.35) rectangle ({\o+(\j+1)/7+0.035},-0.26);}}
\end{tikzpicture}
\caption{대역 통과 필터 뱅크로서의 귀, 도식. 각 곡선은 한 구간의 센서가 여러 주파수의 소리에 보이는 응답이다. 중심 주파수는 로그 눈금 위의 건반처럼 한 옥타브 간격이다. 실제 필터는 고주파 쪽 경사가 가파르고 저주파 쪽 경사가 완만하다. 아래는 건반이다. 귀에서처럼 건반에서도 옥타브마다 같은 자리가 할당된다.}
\label{fig:filterbank}
\end{figure}

귀의 공진기는 수동적이지 않다. 센서 일부는 소리에 맞춰 스스로 칸막이를 흔들어 약한 진동을 진폭으로 수천 배($66$--$76$\,dB) 증폭하고, 소리가 커지면 증폭은 줄어든다~\cite{Ruggero1997,Robles2001}. 엔지니어에게 이것은 자기 발진 경계 바로 옆에 둔 양의 되먹임 증폭기이다. \ref{sec:gaba}장의 네트워크와 같은 경계 위의 삶이다. 경계 근처에서 시스템은 작은 신호에 특히 민감하다. 음량 압축도 같은 증폭기가 한다. 작은 소리는 크게, 큰 소리는 작게 증폭한다.

귀에는 두 번째 부호도 있다. 낮은 주파수에서 뉴런의 스파이크는 소리와 위상을 맞춰 나간다. 뉴런은 모든 파동마다는 아니지만 각 파동의 특정 구간에서 발화한다. 기니피그에서 위상 고정은 약 $600$\,Hz에서 약해지기 시작해 약 $3.5$\,kHz까지 여전히 드러난다~\cite{PalmerRussell1986}. 그래서 스파이크 시각에는 소리의 정확한 시간이 남고, 거기서 두 귀에 소리가 도착하는 시간 차가 나온다. \ref{sec:glutcomputer}장의 네트워크는 이 시간 차로 방향을 찾는다~\cite{Grothe2010}. 높은 주파수에서는 스파이크가 파동을 따라가지 못하고 위치 부호만 남는다. 귀 하나가 \ref{sec:code}장의 두 부호를 모두 쓴다. 뉴런이 따라갈 수 있는 곳에서는 시간을, 그렇지 못한 곳에서는 위치를 쓴다.

\section{나머지 센서}

\begin{table}[t]
\centering
\footnotesize
\setlength{\tabcolsep}{3pt}
\renewcommand{\arraystretch}{1.15}
\begin{tabular}{>{\raggedright}p{1.9cm}>{\raggedright}p{2.6cm}>{\raggedright}p{4.0cm}>{\raggedright\arraybackslash}p{4.6cm}}
\toprule
감각 & 무엇을 재는가 & 입력의 구조 & 이 책의 기법 \\
\midrule
시각 & 빛 & 센서 $\sim10^8$개, $\sim100$배 압축, $\sim10^7$ bit/s & 이득 조절, 이웃 빼기, 두 전선, 움직임 방향 \\
청각 & 공기 압력 & 기계적 필터 뱅크, 능동 증폭기 & 위치 부호와 타이밍 부호, 지연 \\
촉각 & 압력, 진동 & 빨리 적응하는 센서와 느리게 적응하는 센서 & “고역 통과 필터 --- 저역 통과 필터” 쌍 \\
온도 & 열 & 따뜻함과 차가움의 별도 센서 & 이중 레일 부호 \\
후각 & 분자 & 수용체 유형 $\sim400$개, 센서마다 한 유형 & 조합 부호: 냄새는 동작한 유형의 집합 \\
미각 & 음식 속 물질 & 다섯 가지 맛: 단맛, 쓴맛, 짠맛, 신맛, 감칠맛 & 전용 선로; 일부 세포는 글루탐산이 아니라 ATP나 세로토닌을 방출 \\
몸의 위치 & 근육의 신장 & 신장의 길이와 속도 센서 & 운동 조절기를 위한 되먹임 \\
\bottomrule
\end{tabular}
\caption{센서는 어떻게 연결되어 있는가. 셋째 열의 구조는 모두 측정되었다. 오른쪽 열은 이를 앞 장들의 기법과 연결한다.}
\label{tab:sensors}
\end{table}

나머지 감각은 표~\ref{tab:sensors}에 모았다. 거의 모든 행에서 앞 장들의 기법을 알아볼 수 있다. 촉각은 필터 쌍을 쓴다. 어떤 압력 센서는 닿는 순간 응답하고 곧 조용해지며, 다른 센서는 압력이 있는 동안 응답을 유지한다. 온도는 따뜻함과 차가움을 따로 맡는 두 전선이 전한다. 가장 특이한 입력은 후각이다. 사람의 후각 수용체 유형은 약 사백 가지이고, 후각 뉴런 하나는 한 유형만 지닌다~\cite{Mombaerts2004}. 냄새는 한 유형이 아니라 여러 유형의 집합을 활성화하며, 메시지는 \emph{어떤} 유형이 동작했는가이다. 엔지니어에게 이것은 길이 약 사백의 이진 벡터이고, 가능한 값의 수는 천문학적이다.

\begin{keyblock}
\begin{proposition}[기계의 입력]
\label{prop:sensors}
모든 감각은 감도의 물리적 한계 근처에서 동작하는 변환기이다. 거대한 동적 범위를 자동 이득 조절로 압축하고, \nt{GLUT} 버스로 무엇보다 \emph{변화}를 보낸다. 이를 위해 센서는 기계 자신과 같은 기법을 쓴다. 이중 레일 부호, 위치 부호, 타이밍 부호, 배경에 의한 나눗셈이다. 센서의 구조는 측정되었다. 그 부호가 스파이크당 최대 정보에 맞춰져 있다는 것은 근거가 탄탄한 가설이다.
\end{proposition}
\end{keyblock}

입력도 다른 모든 것처럼 비동기로 동작한다. 각 센서는 알릴 것이 있을 때 스파이크를 내고, 감각마다 도착 지연이 다르다. 소리는 수 밀리초, 빛은 수십 밀리초 뒤에 온다. 공통 클럭이 없는데 기계가 이것들을 어떻게 하나의 세계상으로 합치는가는 \ref{sec:synthesis}장의 시간 맞추기 문제 가운데 하나이다.

\section{연습문제}

\begin{exercise}[\lvlA]
\label{ex:11:1}
\emph{광자를 얼마나 모아야 하는가.} 어스름 속에서 센서에 초당 광자 $100$개가 들어온다. 증가분은 요동~\eqref{eq:photonnoise}의 $3$배보다 크면 눈에 띈다. (a)~센서 하나가 밝기의 $10\,\%$ 변화를 알아채는 데 얼마나 걸리는가? (b)~$0.2\,\text{s}$ 안에 해내려면 센서 몇 개를 합해야 하며, 축 방향 해상도는 몇 배 떨어지는가?
\end{exercise}

\begin{exercise}[\lvlA]
\label{ex:11:2}
\emph{스파이크 하나에 몇 비트.} (a)~눈은 초당 스파이크 $3$--$5$개로 발화하는 출력 뉴런 $10^6$개로 $\sim10^7$ bit/s를 보낸다. $\Delta t=1\ms$일 때 용량~\eqref{eq:spikeentropy}의 몇 분의 몇을 쓰는가? (b)~냄새는 $\approx400$개 수용체 유형 중 $20$개를 켠다. 이 집합은 몇 비트를 나르며, 무작위 냄새 두 개는 평균 몇 개의 유형을 공유하는가?
\end{exercise}

\begin{exercise}[\lvlB]
\label{ex:11:3}
\emph{곡선~\eqref{eq:nakarushton} 하나로 할 수 있는 것.} (a)~응답이 $r_{\max}$의 $10\,\%$에서 $90\,\%$로 커지는 밝기 범위는? 몇 자릿수인가? (b)~밝기 열 자릿수를 덮으려면 이동하는 곡선의 위치가 몇 개 필요한가? 음량 열두 자릿수라면?
\end{exercise}

\begin{exercise}[\lvlB]
\label{ex:11:4}
\emph{귀에서 타이밍 부호의 한계.} 두 귀의 도착 시간 차는 최대 $0.22\,\text{m}/343\,\text{m/s}$이다. (a)~스파이크가 순음과 위상을 맞춰 나갈 때, 어느 주파수까지 방향이 한 가지로 정해지는가? (b)~스파이크는 $0.5\ms$만큼 흔들리는데 $20\,\mu\text{s}$를 구별해야 한다. 초당 스파이크 $100$개인 뉴런 몇 개를 $50\ms$ 동안 평균해야 하는가?
\end{exercise}

\begin{exercise}[\lvlB]
\label{ex:11:5}
\emph{눈의 케이블에 이벤트 카메라를.} 출력이 $10^7$ bit/s인 $10^6$픽셀 카메라는 픽셀의 $\ln I$가 $\ln1.2$만큼 변하면 이벤트(주소와 부호)를 보낸다. 밝기 차가 $4$배이고 길이가 $500$픽셀인 가장자리는 최대 얼마나 빨리 움직일 수 있는가? 픽셀당 $8$비트인 일반 카메라는 같은 출력으로 초당 몇 프레임을 보낼 수 있는가?
\end{exercise}

\begin{exercise}[\lvlB]
\label{ex:11:6}
\emph{시뮬레이션: 이득 조절.} 센서~\eqref{eq:nakarushton}는 $\tau=1\,\text{s}$로 $\sigma$를 로그로, $\tau\,\dd\ln\sigma/\dd t=\ln I-\ln\sigma$, 또는 선형으로, $\tau\,\dd\sigma/\dd t=I-\sigma$ 조정한다. 배경은 $1\to100\to1$로 도약한다. 각 법칙에서 위로, 아래로 도약한 뒤 $+10\,\%$ 시험 자극에 대한 응답이 적응된 값의 절반까지 회복되는 데 몇 초가 걸리는가?
\end{exercise}

\begin{exercise}[\lvlC]
\label{ex:11:7}
\emph{곡선은 어떤 세계에 맞춰져 있는가.} 출력 잡음이 작고 일정할 때, 밝기에 대한 정보는 $r(I)=r_{\max}\Phi_I(I)$일 때 최대이다. $\Phi_I$는 밝기의 분포 함수이다. (a)~곡선~\eqref{eq:nakarushton}이 최적인 분포는 무엇이며, 그 $\log_{10}I$의 퍼짐은? (b)~출력 잡음이 푸아송일 때 최적인 곡선은?
\end{exercise}

\chapter{이미지와 영상의 인식}
\chaptermark{인식}
\label{sec:recognition}

\section{과제: 픽셀이 달라도 같은 것}

\ref{sec:sensors}장은 영상을 기계의 입력까지 가져왔다. 눈의 출력 뉴런 약 백만 개, 주로 가장자리와 변화를 전하는 압축된 흐름이다. 그러나 센서와 행동 사이에는 지금까지 말하지 않은 단계가 있다. 그림 속 고양이는 픽셀의 집합이 아니다. 고양이를 반 프레임 옮기고, 돌리고, 멀리 떨어뜨리고, 조명을 절반 끄면 거의 모든 픽셀이 달라지지만, 답 “고양이”는 그대로여야 한다. 엔지니어는 이것을 \emph{불변성}이라고 부른다. 분류기의 답은 이동, 크기, 회전, 조명에 따라 바뀌면 안 된다.

어려움은 간단한 실험에서 잘 드러난다. $24$개 점으로 된 1차원 “망막”과, 어디에나 놓일 수 있는 다섯 점짜리 도형 두 개를 생각하자. 입력을 한 위치에 기억된 견본과 비교하는 분류기는 $49$--$52\%$의 경우에 맞힌다. 동전 던지기와 같다. 이동은 픽셀 단위 비교를 완전히 무너뜨린다. 다른 회로가 필요하다.

\section{단계들의 파이프라인}

뇌가 이 일을 얼마나 빨리 하는지는 \ref{sec:code}장에서 이미 안다. 순간적으로 보여 준 그림 속 동물은 약 $150\ms$ 만에 인식되고, 신호는 약 열 개의 단계를 단계당 $10$--$20\ms$씩 거친다~\cite{Thorpe1996}. 각 단계에서 뉴런은 스파이크를 0개 또는 1개 낼 시간밖에 없다. 따라서 빠른 인식은 오래 숙고하거나 되풀이하지 않는, 파이프라인을 한 번 통과하는 것이다. 문제는 각 단계가 무엇을 하느냐이다.

시각의 첫 단계들에서 얻은 측정이 시사하는 답은~\cite{HubelWiesel1962} 번갈아 나오는 두 연산으로 이루어진다(그림~\ref{fig:conveyor}).

\emph{선택성.} $S$ 단계의 뉴런은 입력의 작은 영역을 보고, 거기에 특정한 부분들의 조합이 있을 때만 발화한다. 이 가장자리가 있고 그 옆에 저 가장자리가 있을 때이다. 이것은 부분들의 AND이다. 어느 한 부분만으로는 넘을 수 없는 문턱값을 가진, \ref{sec:glutcomputer}장의 문턱 뉴런이다.

\emph{불변성.} $C$ 단계의 뉴런은 같은 조합을 서로 다른 위치에서 찾는 많은 $S$ 뉴런의 출력을 모으고, 그 조합이 어디에서든 발견되면 발화한다. 이것은 위치들에 대한 OR, 더 정확히는 최댓값의 선택이다.

1차원 모델에서 두 연산은 다음과 같이 쓴다.
\begin{empheq}[box=\keybox]{equation}
s_{f,p} \;=\; \Bigl[\sum_i t_{f,i}\,x_{p+i} - \theta\Bigr]_+ ,
\qquad
c_f \;=\; \max_p\, s_{f,p} ,
\label{eq:scstage}
\end{empheq}
여기서 $x$는 입력, $t_{f,i}$는 특징 $f$의 견본, $p$는 위치, $\theta$는 문턱값이다. 첫 식은 응답을 선택적으로, 둘째 식은 위치와 무관하게 만든다. 많은 입력 중 최댓값은 \ref{sec:gaba}장의 수단으로 얻는다. 상호 억제는 가장 강한 입력만 남기고(명제~\ref{prop:wta}), 전체 활동에 의한 나눗셈은 밝기가 달라도 응답을 고르게 한다~\cite{Carandini2012}.

\begin{figure}[t]
\centering
\begin{tikzpicture}[>=Stealth,
  px/.style={draw,minimum size=0.36cm,inner sep=0pt},
  on/.style={px,fill=blue!55!black},
  snode/.style={circle,draw,very thick,minimum size=0.62cm,inner sep=0pt,font=\scriptsize,fill=blue!6},
  cnode/.style={circle,draw,very thick,minimum size=0.8cm,inner sep=0pt,font=\scriptsize,fill=red!10}]
% two inputs: the same shape at two positions
\foreach \row/\sh/\lab in {0/1/{프레임 1},1/7/{프레임 2}}{
  \begin{scope}[yshift=-\row*1.0cm]
  \node[left,font=\scriptsize] at (-0.25,0) {\lab};
  \foreach \i in {0,...,13}{\node[px] at (\i*0.4,0) {};}
  \foreach \d in {0,1,3,4}{\node[on] at ({(\sh+\d)*0.4},0) {};}
  \end{scope}}
% S stage: detectors of the shape at several positions
\foreach \k in {0,...,5}{
  \node[snode] (S\k) at ({0.8+0.8*\k},1.7) {$S$};}
\node[font=\scriptsize,left] at (-0.25,1.7) {$S$ 단계: 부분들의 AND};
\draw[->,thick,blue!60!black] (1.2,0.25) -- (S1.south);
\draw[->,thick,orange!85!black] (3.6,0.25) -- (S4.south);
\node[font=\scriptsize,blue!60!black,anchor=east] at (1.25,0.85) {프레임 1};
\node[font=\scriptsize,orange!85!black,anchor=west] at (3.9,0.85) {프레임 2};
% C stage
\node[cnode] (C) at (2.8,3.25) {$C$};
\node[font=\scriptsize,left] at (-0.25,3.25) {$C$ 단계: 위치들에 대한 OR};
\foreach \k in {0,...,5}{\draw[->,gray!60!black] (S\k.north) -- (C);}
\draw[->,very thick,red!70!black] (C.north) -- ++(0,0.6) node[above,font=\scriptsize,black] {두 프레임 모두에서 “도형이 있다”};
% next stages
\draw[decorate,decoration={brace,amplitude=5pt},gray!60!black] (6.3,1.4) -- (6.3,-1.3);
\node[right,font=\scriptsize,align=left] at (6.5,0.05) {같은 물체,\\다른 픽셀};
\draw[->,thick,densely dashed,gray!60!black] (3.6,3.25) -- (5.9,3.25) node[right,font=\scriptsize,black,align=left] {다음 $S$, $C$ 쌍:\\더 크고, 더 복잡하고,\\더 불변적};
\end{tikzpicture}
\caption{파이프라인의 단계 한 쌍. $S$ 뉴런들은 같은 부분 조합을 서로 다른 위치에서 찾는다. $C$ 뉴런은 그 조합이 어디에서든 발견되면 응답한다~\eqref{eq:scstage}. 프레임 1과 프레임 2의 도형은 서로 다른 $S$ 뉴런을 흥분시키지만 같은 $C$ 뉴런을 흥분시킨다. 다음 단계 쌍들은 점점 더 크고 복잡한 조합에 대해 같은 일을 되풀이한다.}
\label{fig:conveyor}
\end{figure}

같은 1차원 모델에서 단계 쌍~\eqref{eq:scstage}는 도형을 어느 위치에서나 오류 없이 인식하고, 입력의 각 점을 확률 $0.1$로 무작위로 뒤집으면 $86\%$, 확률 $0.2$로 뒤집으면 $70\%$의 경우에 인식한다. 동전은 여전히 절반이다. 이런 쌍 여러 개를 위로 쌓아 보자. 첫 쌍은 가장자리를, 다음 쌍은 가장자리의 조합을, 더 위는 물체의 부분을, 맨 위는 물체 전체를 찾는다. 쌍을 하나 지날 때마다 조합은 커지고, 답은 물체가 어디에 어떻게 놓였는지에 점점 덜 좌우된다~\cite{Riesenhuber1999}.

\begin{keyblock}
\begin{proposition}[선택성과 불변성의 파이프라인]
\label{prop:conveyor}
빠른 인식은 열 개쯤의 단계를 한 번 통과하는 것이다. 각 단계 쌍에서 부분들의 AND~\eqref{eq:scstage}는 응답을 선택적으로 만들고, 위치들에 대한 OR은 이동과 무관하게 만든다. 이 두 연산을 번갈아 하면 물체를 구별하면서도 물체가 어디에 어떻게 보이는지에 거의 좌우되지 않는 답이 나온다. 첫 단계들의 구조와 통과 속도는 측정되었다. 사슬 전체가 이 두 연산으로 환원된다는 것은 단순화이다.
\end{proposition}
\end{keyblock}

이 구조가 낯익어 보인다면 그럴 만하다. 바로 이 교대, 곧 모든 위치에서 특징을 찾고 이웃 위치들을 합치는 방식을 엔지니어들이 인공 \emph{합성곱} 신경망으로 옮겼다~\cite{Fukushima1980}. 그리고 최근에는 반대 방향으로 갔다. 물체 인식을 학습한 합성곱 신경망이 시각의 윗단계 뉴런 응답을 가장 잘 설명하는 모델로 드러났다~\cite{Yamins2014}. 맨 위의 결과는 간단히 읽힌다. 마지막 단계 뉴런 약 백 개의 발화율을 합친 것으로부터, 선형 판정 블록이 제시 후 0.1초 만에 물체를 알아본다~\cite{Hung2005}. 이것은 \ref{sec:code}장의 집단 발화율 부호이다.

\section{없는 교사: 영상으로 배우기}

합성곱 신경망은 이름표가 붙은 수백만 장의 그림으로 학습한다. 뇌에게는 거의 아무도 이름표를 주지 않는다. 아이는 물체를 몇 시간씩 보지만 “컵”이라는 말은 가끔 듣는다. 그렇다면 $C$ 단계는 어떤 $S$ 뉴런들을 합쳐야 할지 어떻게 알까? “여기 있는 이 도형”과 “저기 있는 이 도형”을 합치려면 그것이 같은 도형이라는 것을 이미 알아야 한다.

실마리는 영상 자체가 준다. 세상은 연속적으로 변한다. 시야 속 물체는 움직이고, 돌고, 조명이 바뀌는 동안에도 같은 물체로 남고, 시선이 다른 물체로 옮겨 가는 일은 가끔뿐이다. \emph{느리게} 변하는 것은 대개 물체에 속하고, 빠르게 변하는 것은 물체의 위치와 모습에 속한다~\cite{WiskottSejnowski2002}. 따라서 $C$ 뉴런에게는 \ref{sec:dofa}장의 흔적을 가진 헵 규칙이면 충분하다. 동시에 온 것만이 아니라 잇달아 온 것을 연결하면 된다~\cite{Foldiak1991}.
\begin{empheq}[box=\keybox]{equation}
\tau_{\text{tr}}\,\frac{\dd\bar y_k}{\dd t} \;=\; -\bar y_k + y_k ,
\qquad
\frac{\dd\mathbf w_k}{\dd t} \;=\; \eta\,\bar y_k\,\bigl(\mathbf x - \mathbf w_k\bigr) .
\label{eq:traceinvariance}
\end{empheq}
여기서 $y_k$는 억제를 통한 경쟁에서 이긴 $C$ 뉴런 $k$의 활동, $\bar y_k$는 그 흔적으로 규칙~\eqref{eq:threefactor}에서와 같은 RC 회로이며, $\mathbf x$는 $S$ 뉴런들의 출력이다. 물체의 첫 모습에서 이긴 뉴런은 둘째 모습이 올 때에도 그것을 기억하고 있어서, 그 모습도 자기 쪽으로 끌어당긴다.

\begin{figure}[t]
\centering
\begin{tikzpicture}[>=Stealth,x=1.0cm,y=1.0cm]
\foreach \i/\lab in {0/1,1/2,2/3,3/4}{
  \node[draw,thick,fill=blue!8,minimum width=0.9cm,minimum height=0.55cm,font=\scriptsize] at (\i+0.5,2.2) {$A$(\lab)};}
\foreach \i/\lab in {0/4,1/3,2/2,3/1}{
  \node[draw,thick,fill=orange!12,minimum width=0.9cm,minimum height=0.55cm,font=\scriptsize] at (\i+6.5,2.2) {$B$(\lab)};}
\node[font=\scriptsize,gray!40!black] at (5.0,2.2) {휴지};
\draw[->,gray!70!black] (0,0) -- (10.4,0) node[below left,font=\scriptsize,black] {시간};
% traces
\draw[thick,blue!60!black] (0,0.05) .. controls (0.4,1.2) and (0.8,1.3) .. (1.2,1.35) -- (4,1.4) .. controls (4.6,0.5) and (5.2,0.15) .. (6,0.08) -- (10,0.05);
\draw[thick,orange!85!black] (0,0.05) -- (6,0.05) .. controls (6.4,1.2) and (6.8,1.3) .. (7.2,1.35) -- (10,1.4);
\node[font=\scriptsize,blue!60!black,anchor=west] at (1.4,1.65) {흔적 $\bar y_1$: $A$가 지나가는 내내 유지};
\node[font=\scriptsize,orange!85!black,anchor=west] at (7.2,1.65) {흔적 $\bar y_2$};
\draw[decorate,decoration={brace,amplitude=4pt},gray!60!black] (4.0,2.6) -- (0,2.6);
\draw[decorate,decoration={brace,amplitude=4pt,mirror},gray!60!black] (0,-0.35) -- (4,-0.35) node[midway,below=4pt,font=\scriptsize,black] {$A$의 모든 모습이 뉴런 1과 연결된다};
\end{tikzpicture}
\caption{영상이 불변성을 가르치는 방식, 도식. 물체 $A$가 네 위치를 지나가고, 휴지가 있은 뒤 물체 $B$가 온다. $A$의 첫 모습에서 이긴 뉴런의 흔적~\eqref{eq:traceinvariance}은 지나가는 내내 유지되어, $A$의 모든 모습이 한 뉴런과 연결된다. 휴지가 흔적을 지우므로 $B$는 다른 뉴런에게 간다.}
\label{fig:tracevideo}
\end{figure}

우리는 이것을 모델로 확인했다(그림~\ref{fig:tracevideo}). $C$ 뉴런 두 개가 $S$ 뉴런 $16$개의 입력을 놓고 경쟁한다. 두 도형이 각각 여덟 위치에 있다. 도형은 모든 위치를 위치당 $50\ms$씩 지나가고, $0.3\,\text{s}$ 휴지 뒤 다음 무작위 도형이 온다. 이름표는 전혀 없다. 흔적이 짧으면($\tau_{\text{tr}}=5\ms$) 뉴런들은 입력을 \emph{위치}에 따라 나누고, 답은 동전보다 나을 것 없이 도형과 일치한다. 열 번 실행한 평균이 $52\%$이다. 흔적이 한 위치의 제시 시간과 비슷하면, 곧 $\tau_{\text{tr}}\approx20\ms$부터는 각 뉴런이 자기 도형의 \emph{모든} 위치를 모은다($20$, $50$, $200\ms$에서 열 번 실행 모두). 불변성이 영상만으로 학습된 것이다. 물체 사이의 휴지가 중요하다. 휴지가 없으면 흔적이 다음 물체로 흘러 들어가 둘을 뒤섞는다.

같은 현상이 실험에서도 보인다. 눈이 한 곳에서 다른 곳으로 도약하는 동안 실험에서 몰래 물체를 바꿔치기하면, 윗단계 뉴런들은 한두 시간 만에 바꿔치기한 물체에 원래 물체처럼 응답하기 시작한다. 잇달아 온 것을 연결하는 것이다~\cite{LiDiCarlo2008}. 뇌의 불변성은 실제로 시간의 연속성에서 학습된다. 이 규칙이 시각 학습 전체를 얼마나 설명하는지는 가설이다.

\section{영상은 곧 움직임이다}

영상은 학습을 위한 프레임의 흐름일 뿐 아니라 그 자체가 과제이다. 무엇이, 어디로, 얼마나 빨리 움직이는가. 첫걸음은 익숙하다. \ref{sec:gaba}장의 방향 검출기이다(그림~\ref{fig:direction}). 여기서는 지연된 억제가 한 방향의 움직임을 지운다. 이런 검출기는 시야 전체에 퍼져 있고, 그다음은 형태 특징과 똑같이 다룬다. 같은 방향의 검출기 여러 개를 서로 다른 위치에서 모으는 단계는 큰 물체의 움직임에, 그 물체가 어디에 있든 응답한다. 이것은 같은 AND와 OR의 쌍~\eqref{eq:scstage}이며, 다만 특징이 “가장자리”가 아니라 “시간 $d$ 동안 옮겨 간 가장자리”일 뿐이다.

영상은 또 예측 가능하다. 다음 프레임은 이전 프레임과 거의 같다. \emph{예측 부호화} 가설에 따르면 윗단계는 아랫단계에 예측을 보내고, 위로는 주로 오차, 곧 예측이 고려하지 못한 것이 올라간다~\cite{RaoBallard1999}. 이것은 명제~\ref{prop:economy}와 같은 스파이크 절약이다. 이미 아는 것이 아니라 새 소식에 값을 치른다. 개별 관찰은 이 가설과 부합하지만, 파이프라인의 일반 구조로서는 증명되지 않았다.

\section{뇌와 합성곱 신경망}

유사성은 어디까지 가는가? 물체 하나를 빠르게 인식하는 데서는 정말 깊다~\cite{Yamins2014}. 그러나 뇌에는 단순한 합성곱 신경망에 없는 것도 있다.

\emph{되먹임.} 뇌의 파이프라인은 순방향만이 아니다. 단계마다 뒤로, 옆으로 가는 연결이 있다. 가려지거나 조명이 나쁜 어려운 그림에서는 윗단계의 응답이 더 늦게 형성되고, 이 늦은 응답은 순방향 네트워크보다 되먹임 네트워크로 더 잘 기술된다~\cite{Kar2019}. 한 번의 통과로 쉬운 경우를 풀고, 어려운 경우에는 여러 바퀴가 필요하다.

\emph{강건성.} 인공 네트워크는 눈에 띄지 않는 픽셀 위조로 속일 수 있다~\cite{Szegedy2014}. 사람이 보지 못하는 변화가 “판다”를 “긴팔원숭이”로 바꾼다~\cite{Goodfellow2015}. 사람의 시각은 이런 위조에 훨씬 강하지만 무적은 아니다. 여러 네트워크를 한꺼번에 속이는 그림은 짧게 보여 주면 사람의 선택도 옮겨 놓는다~\cite{Elsayed2018}. 강건성이 왜 이렇게 다른지는 열린 문제이다. 후보는 잡음, 전체 활동에 의한 나눗셈, 되먹임이다.

\emph{교사.} 네트워크에는 이름표 붙은 예가 수백만 개 필요하고, 뇌에는 주로 이름표 없는 영상~\eqref{eq:traceinvariance}과 무엇이 중요한지에 대한 \nt{DOPA}의 약간의 힌트면 된다.

\emph{에너지.} 뇌 전체가 약 $20$와트이고(명제~\ref{prop:economy}) 인식은 그중 일부만 쓴다. 학습된 합성곱 신경망은 그래픽 가속기에서 수백 와트를 쓴다.

\section{착시: 기계는 어디서, 왜 틀리는가}
\label{sec:illusions}

남의 회로를 이해하려는 엔지니어는 회로에 특이한 신호를 넣고 어디서 틀리는지 본다. 시각에 대해서는 이런 신호가 오래전에 고안되었다. 바로 착시이다. 착시는 고장이 아니다. 착시 하나하나는 기계의 특정 메커니즘을 가리킨다. 그 메커니즘은 보통 환경에서는 올바로 동작하지만, 특별히 고른 그림 앞에서는 정체를 드러낸다.

\emph{합리적 기대.} 입력에는 잡음이 있고, 기계는 세상에서 흔히 일어나는 것으로 그 입력을 보충한다. 느린 움직임은 빠른 움직임보다 흔하다. 입력을 믿기 어려울 때, 예를 들어 그림의 대비가 낮을 때 속도 추정은 느린 쪽으로 치우치고, 흐린 물체는 선명한 물체보다 느리게 움직이는 것처럼 보인다~\cite{Weiss2002}. 많은 밝기 착시도 같은 식으로 설명된다. 우리는 픽셀의 밝기가 아니라, 그 밝기가 과거에 대개 어떤 표면에 해당했는지를 본다~\cite{Purves2011}. 어떤 사람에게는 흰색과 금색으로, 다른 사람에게는 파란색과 검은색으로 보이는 드레스 사진도 같은 메커니즘이다. 그림이 모호하고, 사람들은 무의식적으로 가정하는 조명에서 서로 갈린다~\cite{LaferSousa2015,Wallisch2017}. 사람 사이의 차이는 측정되었다. 뇌가 여기서 확률론의 규칙에 따라 입력과 기대를 결합한다는 것은 근거가 탄탄한 가설이다.

\emph{이득 조절.} 폭포를 1분 동안 바라본 뒤 움직이지 않는 바위를 보라. 바위가 위로 떠오른다. “아래” 채널과 “위” 채널은 \eqref{eq:dualrail}에서처럼 전선 한 쌍이다. 이득 조절~\eqref{eq:nakarushton}이 지친 채널을 약하게 만들어 쌍의 균형이 옮겨 간 것이다. 주변이 중심의 모습을 바꾸는 기울기 착시와 대비 착시는 \ref{sec:gaba}장에서처럼 이웃 활동에 의한 나눗셈으로 설명된다~\cite{Schwartz2009}. 조심해야 한다는 교훈적인 예도 있다. 격자의 흰 줄이 교차하는 곳의 어두운 점은 오랫동안 단순한 이웃 억제로 설명되었지만, 변형한 격자로 한 실험은 이 설명이 통하지 않음을 보였다~\cite{SchillerCarvey2005}.

\emph{지연 보상.} 눈에서 윗단계까지 가는 데는 수십 밀리초가 걸리고, 그 사이에 움직이는 물체는 이미 옮겨 간다. 물체 옆에 정지한 점이 번쩍이면, 둘이 일치했는데도 물체가 섬광보다 앞서 있는 것처럼 보인다~\cite{Nijhawan1994}. 한 설명에 따르면 기계는 지연을 보상하려고 움직임을 앞으로 연장한다. \ref{sec:muscles}장에서와 같다. 되먹임이 늦으면 예측해야 한다. 이 착시에는 설명이 여럿 있고 논쟁은 끝나지 않았다.

\emph{시간 부호의 오류.} 색 전환이 날카로운 고리로 된 정지 그림이 돌아가는 것처럼 보인다. 대비가 강한 부분은 흐린 부분보다 빨리 처리되고, 지연의 차이~\eqref{eq:latency}를 방향 검출기가 움직임으로 읽는다~\cite{Conway2005}. 착시를 일으키는 것은 작은 불수의적 눈 도약과 눈 깜박임이다. 도약할 때마다 입력이 새로 고쳐지고, 검출기는 다시 “움직임”을 본다~\cite{OteroMillan2012}. 이것은 \ref{sec:code}장의 시간 부호를 잘못 읽은 것이다.

\emph{한 그림 속 두 그림.} 한 면으로 우리를 보다가 다른 면으로 보는 정육면체 골격, 또는 두 눈에 보여 준 서로 다른 그림. 지각은 몇 초마다 뒤바뀐다. 가장 좋은 모델은 두 해석 사이의 상호 억제, 느린 피로, 잡음이다~\cite{MorenoBote2007}. 곧 \ref{sec:muscles}장에서 걸음의 리듬을 만든 바로 그 회로~\eqref{eq:halfcentre}이다. 전환 시간은 잡음과 피로가 정한다. 모델~\cite{MorenoBote2007}에 따르면 주된 역할은 잡음이 하고 피로는 거들 뿐이다.

인공 네트워크는 어떨까? 영상의 다음 프레임을 예측하도록 학습한 네트워크는 같은 정지 고리에서 회전을 보고, 대조 그림에서는 보지 않는다~\cite{Watanabe2018}. 이미지의 잡음을 없애고 선명도를 높이도록 학습한 네트워크는 사람과 같은 색 착시와 밝기 착시에 넘어간다~\cite{GomezVilla2020}. 그러므로 착시의 일부는 신경 조직의 특성이 아니라 기계가 배우는 과제 자체의 결과이다. 어떤 착시가 뇌에만 있는지는 어떤 시각 모델에게나 좋은 검사이다.

\begin{keyblock}
\begin{proposition}[메커니즘의 지문으로서의 착시]
\label{prop:illusions}
착시는 보통 세상에서 쓸모 있는 메커니즘이 특이한 입력을 받을 때 생긴다. 기대가 믿기 어려운 입력을 이기고, 이득 조절이 채널 쌍을 옮기고, 지연 보상이 물체를 앞으로 끌어가고, 지연 차이가 움직임으로 읽히고, 잡음과 피로를 동반한 상호 억제가 뒤바뀐다. 착시 하나하나는 자기 메커니즘을 가리키는 시험 신호이다. 착시 자체는 측정되었다. 그 설명은 근거가 탄탄한 것부터 논쟁 중인 것까지 있다.
\end{proposition}
\end{keyblock}

\section{정리}

\begin{table}[t]
\centering
\footnotesize
\setlength{\tabcolsep}{4pt}
\renewcommand{\arraystretch}{1.15}
\begin{tabular}{>{\raggedright}p{3.0cm}>{\raggedright}p{4.6cm}>{\raggedright\arraybackslash}p{5.2cm}}
\toprule
기계가 하는 일 & 방법 & 알려진 것 \\
\midrule
$150\ms$ 만에 인식 & 열 개쯤의 단계를 한 번 통과 & 속도는 측정됨 \\
응답을 선택적으로 & 부분들의 AND, $S$ 단계~\eqref{eq:scstage} & 첫 단계들에서 측정됨 \\
응답을 불변으로 & 위치들에 대한 OR, $C$ 단계; 억제와 나눗셈 & 첫 단계들에서 측정됨; 윗단계에서는 단순화 \\
답을 냄 & 마지막 단계 뉴런 약 백 개의 발화율, 선형 판독 & 측정됨 \\
이름표 없이 학습 & 연속 영상에 흔적을 가진 헵 규칙~\eqref{eq:traceinvariance} & 물체 바꿔치기가 응답을 바꿈은 측정됨; 주된 규칙으로서는 가설 \\
움직임을 봄 & 방향 검출기, 이어서 같은 AND/OR 쌍 & 측정됨 \\
예측 가능한 것을 아낌 & 위로 예측 오차가 올라감 & 가설 \\
어려운 경우를 품 & 되먹임, 여러 바퀴 & 늦은 응답과 되먹임의 역할은 측정됨 \\
예측 가능하게 틀림 & 착시: 기대, 이득 조절, 지연, 잡음과 피로를 동반한 상호 억제 & 착시는 측정됨; 설명은 근거 있는 것부터 논쟁 중인 것까지 \\
\bottomrule
\end{tabular}
\caption{기계는 이미지와 영상을 어떻게 인식하는가.}
\label{tab:recognition}
\end{table}

인식은 우리가 이미 가진 부품으로 조립된다(표~\ref{tab:recognition}). 문턱값은 부분들의 AND를, 상호 억제는 위치들에 대한 OR을, 나눗셈은 밝기와의 무관함을, 흔적을 가진 헵 규칙은 없는 교사를 대신한다. 엔지니어들은 시각에서 선택성과 불변성의 교대를 가져와 그 위에 합성곱 신경망을 세웠다. 반대로 뇌는 가장 중요한 것을 아직 내주지 않았다. 이름표 없는 영상으로, 거의 에너지를 쓰지 않고 배우는 능력이다.

\section{연습문제}

\begin{exercise}[\lvlA]
\label{ex:12:1}
\emph{한 단계에서의 발화율.} 파이프라인의 한 단계는 $15\ms$ 걸리고, 뉴런은 초당 스파이크 $20$개로 발화하며, 잡음은 $c=0.1$로 상관되어 있다~\eqref{eq:correlated}. 뉴런 백 개 집단의 발화율 상대 오차는 얼마이며, 집단이 한없이 클 때의 극한은? 한 단계에서 발화율의 몇 단계를 구별할 수 있는가?
\end{exercise}

\begin{exercise}[\lvlA]
\label{ex:12:2}
\emph{인식 한 번의 에너지.} $150\ms$ 동안의 한 번 통과에서 뉴런 $10^7$개씩인 열 단계가 비용 $10^{-10}$~J인 스파이크를 최대 하나씩 낸다. (가)~인식은 이 $150\ms$ 동안 뇌 전체 에너지의 몇 분의 몇을 쓰는가? (나)~그림을 $10\ms$에 인식하는 $300$~W 가속기는 몇 배를 쓰는가?
\end{exercise}

\begin{exercise}[\lvlB]
\label{ex:12:3}
\emph{$S$ 단계의 문턱값.} $24$개 점의 “망막”, 도형 $A=11011$과 $B=10101$, 견본~\eqref{eq:scstage}은 도형의 1에서 $+1$, 0에서 $-1$이다. (가)~$S$ 뉴런이 뒤집힌 점 하나는 봐주되 다른 도형에는 침묵하는 문턱값은? (나)~$100\times100$ 그림에서 $5\times5$ 특징 열 개에 $S$ 뉴런이 몇 개 필요한가?
\end{exercise}

\begin{exercise}[\lvlB]
\label{ex:12:4}
\emph{두 그림 중 하나는 얼마나 오래 가는가.} 전환 회로~\eqref{eq:halfcentre}에서 $\tau\ll\tau_a$이고, 집단~1이 이기는 동안 $r_2=0$, $r_1=I-a_1$이다. (가)~$I=1$, $\beta=2$, $b=2$, $\tau_a=0.4\,\text{s}$일 때 승리 지속 시간을 구하라. (나)~어떤 $\tau_a$라야 그림이 3초마다 바뀌는가?
\end{exercise}

\begin{exercise}[\lvlB]
\label{ex:12:5}
\emph{시뮬레이션: 불변성의 대가.} \texttt{models/\allowbreak recognition\_models.py}의 함수 \texttt{two\_stage}($4000$번 시행)로, $N=24$에서 $S$, $C$ 단계 쌍이 네 번에 한 번 틀리는 점 뒤집기 확률 $p$를 구하라. $p=0.1$에서 $N=8$부터 $N=192$까지 정답 비율은 어떻게 변하는가?
\end{exercise}

\begin{exercise}[\lvlB]
\label{ex:12:6}
\emph{시뮬레이션: 흔적의 창.} \texttt{models/\allowbreak recognition\_models.py}의 \texttt{trace\_learning}을 휴지 $0.3\,\text{s}$로 열 번 실행해, $5\ms$에서 $2\,\text{s}$까지의 $\tau_{\text{tr}}$ 중 어느 값에서 모든 실행이 불변성을 오류 없이 학습하는지 구하라. 이 창의 상한은 어디서 오는가?
\end{exercise}

\begin{exercise}[\lvlB]
\label{ex:12:7}
\emph{어디서든 움직임.} $0.1^\circ$ 간격으로 늘어선 이웃 점들에 \ref{sec:gaba}장의 방향 검출기가 있다. 지연 $d$를 가진 $A$로부터의 억제는, 둘의 시간 차가 $5\ms$ 이내이면 $B$로부터의 흥분을 지운다. 그 위에 $C$ 단계가 있다. 초당 $5$에서 $20^\circ$ 속도의 역방향 움직임에서 $C$가 침묵하려면 어떤 지연이 필요한가?
\end{exercise}

\begin{exercise}[\lvlC]
\label{ex:12:8}
\emph{픽셀 위조.} 깨끗한 도형에서 모델 \texttt{two\_stage}($N=24$)의 답을 바꾸려면 점 몇 개를 뒤집어야 하는가? 역시 이동 불변인 분류기 “입력의 1이 $3.5$개보다 많으면 $A$”라면? $p=0.1$에서 이 분류기의 정답 비율은, 단계 쌍의 $0.86$과 비교해 얼마인가?
\end{exercise}

\chapter{기계의 출력: 뇌는 근육을 어떻게 제어하는가}
\chaptermark{기계의 출력}
\label{sec:muscles}

\section{빠른 출력}

기계가 세상에서 빠르게 하는 일은 모두 근육으로 한다. 말, 시선, 걸음, 미소는 모두 근육의 수축이다. 두 번째 출력, 곧 느린 출력인 몸의 화학은 \ref{sec:chemoutput}장에서 다룬다. 그림~\ref{fig:machine}에서 출력은 “근육”이라는 화살표였다. 이제 그 화살표 뒤에 무엇이 있는지 보자.

사슬의 마지막 고리는 \nt{ACET} 뉴런이다. 표~\ref{tab:types}에 “근육으로의 출력”이라고 적힌 바로 그 뉴런이다. 근육 섬유 하나는 정확히 그런 뉴런 하나에게서 입력을 받고, 뉴런 하나는 수백에서 이천 개 남짓의 섬유 무리를 제어한다~\cite{Feinstein1955mu}. 섬세한 조정이 필요한 근육에서는 단위가 더 작고, 다리의 큰 근육에서는 더 크다. 뉴런과 그 섬유를 합쳐 \emph{운동 단위}라고 부른다. 뇌가 따로 켤 수 있는 근육의 가장 작은 조각이다.

근육 위 뉴런의 시냅스는 \ref{sec:code}장의 피질 시냅스와 전혀 다르게 만들어져 있다. 거기서는 스파이크 대부분이 사라졌다. 여기서는 스파이크 하나가 섬유를 동작시키는 데 필요한 양의 몇 배나 되는 아세틸콜린을 방출한다~\cite{WoodSlater2001}. 이것은 \emph{안전 여유}를 가진 시냅스이다. 뉴런의 스파이크 하나가 섬유의 수축 정확히 한 번이다. 기계 내부에서는 신뢰할 수 없는 연결을 허용하고 오류를 중복으로 바로잡을 수 있다. 출력에서는 오류가 곧 움직임을 대가로 치르므로, 신뢰성을 하드웨어로 곧장 사 두었다.

\section{힘: 발화율과 단위의 수}

스파이크 하나는 단일 수축, 곧 연축을 일으킨다. 힘은 수십 밀리초에 걸쳐 커졌다가 줄어든다. 좋은 근사식은 다음과 같다~\cite{Fuglevand1993}.
\begin{equation}
h(t) \;=\; F_1\,\frac{t}{T_c}\,e^{\,1-t/T_c},
\label{eq:twitch}
\end{equation}
여기서 $T_c$는 상승 시간으로 $50\ms$ 정도이고, $F_1$은 연축 한 번의 힘이다. 스파이크가 더 자주 오면 연축이 합쳐진다.
\begin{empheq}[box=\keybox]{equation}
F(t) \;=\; \sum_k h\bigl(t-t_k\bigr) .
\label{eq:forcesum}
\end{empheq}
이것은 \ref{sec:serocomputer}장에서 스파이크를 농도로 바꾸던 것과 같은 저역 통과 필터이다. 다만 이제 출력이 농도가 아니라 힘이다. 근육은 스파이크를 힘으로 바꾸는 디지털-아날로그 변환기이다(그림~\ref{fig:tetanus}). 우리 모델에서 초당 스파이크 다섯 개일 때 연축은 거의 겹치지 않고 힘은 $0.2F_1$에서 $1.1F_1$까지 오르내린다. 열다섯 개일 때는 이미 $1.8F_1$과 $2.2F_1$ 사이에 머물고, 마흔 개일 때는 거의 평탄해져 약 $5.4F_1$이 된다. 실제 근육은 스파이크가 잦으면 포화하므로, 선형 합~\eqref{eq:forcesum}은 초당 스파이크 수십 개까지만 맞다.

\begin{figure}[t]
\centering
\begin{tikzpicture}[>=Stealth,x=20cm,y=0.55cm]
\draw[->,gray!70!black] (0,0) -- (0.66,0) node[right,font=\small,black] {$t$, s};
\draw[->,gray!70!black] (0,0) -- (0,6.4) node[left,font=\small,black] {$F/F_1$};
\foreach \t in {0,0.2,0.4,0.6}{\draw[gray!60!black] (\t,-0.1) -- (\t,0.1); \node[below,font=\scriptsize] at (\t,0) {\t};}
\foreach \f in {1,3,5}{\draw[gray!60!black] (-0.004,\f) -- (0.004,\f); \node[left,font=\scriptsize] at (0,\f) {\f};}
\draw[thick,blue!60!black] plot file {figures/muscle_twitch_5.dat};
\draw[thick,green!45!black] plot file {figures/muscle_twitch_15.dat};
\draw[thick,red!70!black] plot file {figures/muscle_twitch_40.dat};
\node[font=\scriptsize,blue!60!black,anchor=west] at (0.605,0.9) {5 Hz};
\node[font=\scriptsize,green!45!black,anchor=west] at (0.605,2.1) {15 Hz};
\node[font=\scriptsize,red!70!black,anchor=west] at (0.605,5.4) {40 Hz};
\end{tikzpicture}
\caption{디지털-아날로그 변환기로서의 근육: 운동 단위 하나의 규칙적인 스파이크에 대한 힘~\eqref{eq:forcesum}, $T_c=50\ms$. 드문 스파이크는 따로따로 연축을 일으키고, 잦은 스파이크는 고른 힘으로 합쳐진다. 그림~\ref{fig:lowpass}에서 잦은 스파이크가 고른 농도로 합쳐지는 것과 같다.}
\label{fig:tetanus}
\end{figure}

뇌에게는 힘의 노브가 두 개 있다. 각 단위의 스파이크 발화율과, 켜진 단위의 수이다. 그리고 둘째 노브는 매우 영리하게 만들어져 있다. 근육 속 단위들은 서로 다르다. 가장 약한 단위는 가장 강한 단위보다 백 배 약하다. 점점 더 큰 힘이 필요해지면 단위들은 \emph{크기 순서대로}, 작은 것부터 큰 것으로 켜진다~\cite{Henneman1957}. 이 순서를 계산하는 주체는 없다. 작은 뉴런은 같은 공통 입력에 그냥 더 쉽게 흥분하므로, 켜지는 순서는 하드웨어 자체가 정한다.

이것이 무엇을 주는가? 순서상 다음 단위가 앞 단위보다 $q$배 강하다고 하자. $q$는 1보다 조금 크다. 켜진 모든 단위의 힘은 등비수열의 합이고, 그 거의 전부가 마지막 단위들에서 온다. 그러면 새로 켜진 단위가 더하는 힘은
\begin{empheq}[box=\keybox]{equation}
\frac{\Delta F}{F} \;\approx\; q-1 \;=\; \text{const} ,
\label{eq:sizeprinciple}
\end{empheq}
곧 힘의 단계가 힘 자체에 비례한다. 약한 힘은 작은 몫으로, 강한 힘은 큰 몫으로 조절한다. 이것은 \emph{부동소수점 수}이다. 지수는 켜진 단위의 수가 정하고, 가수는 그 스파이크 발화율이 정한다. \ref{sec:sensors}장의 센서와 같은 로그 눈금이다. 기계의 입력과 출력은 세상을 상대적인 비율로 잰다.

부동소수점에는 이면이 있다. 힘의 퍼짐도 힘 자체에 비례해 커진다. 강한 움직임은 약한 움직임보다 필연적으로 덜 정밀하다. 영향력 있는 한 가설에 따르면, 신호에 의존하는 바로 이 잡음이 우리의 움직임이 매끄러운 이유를 설명한다. 매끄러운 명령이 끝점에서 가장 작은 퍼짐을 준다~\cite{HarrisWolpert1998}.

\section{당기되 밀지 못한다: 관절 하나에 두 전선}

근육은 당길 줄만 안다. \nt{GLUT} 뉴런이 음의 신호를 낼 수 없듯이 근육은 밀 수 없다. 해법은 \ref{sec:glutcomputer}장에서 익숙한 \emph{두 전선}이다. 관절마다 반대 방향으로 당기는 근육 한 쌍이 붙어 있다(그림~\ref{fig:antagonists}). 두 힘이 $F_1$, $F_2$이고 지레팔이 $\ell$이면, 관절의 회전력과 강성은 다음과 같다.
\begin{empheq}[box=\keybox]{equation}
M \;=\; \ell\,(F_1-F_2), \qquad K \;\approx\; \kappa_m\,(F_1+F_2),
\label{eq:antagonists}
\end{empheq}
여기서 $\kappa_m$은 근육의 힘과 강성을 잇는 계수이다. 긴장한 근육은 이완한 근육보다 더 세게 튕긴다. 힘의 차는 회전력을, 합은 강성을 정한다~\cite{Hogan1984}. 두 전선으로 뇌는 독립된 두 양을 제어한다. 관절을 어디로 돌릴지와 얼마나 단단히 붙잡을지이다. 찻잔을 쏟지 않고 들고 있어야 할 때 우리는 두 근육을 동시에 긴장시킨다. 회전력은 같고 강성은 높아져서, 우연한 충격이 팔을 덜 흔든다.

\begin{figure}[t]
\centering
\begin{tikzpicture}[>=Stealth,
  nrn/.style={circle,draw,very thick,minimum size=0.95cm,inner sep=0pt,font=\scriptsize},
  inh/.style={-{Bar[width=7pt]},thick,red!70!black}]
% the joint: a bar on a pivot, two springs pulling down on either side
\fill[gray!30] (4.6,-0.25) -- (5.4,-0.25) -- (5,0.15) -- cycle;
\draw[line width=3pt,gray!50!black] (2.4,0.2) -- (7.6,0.2);
\fill[black] (5,0.2) circle (0.08);
\draw[decorate,decoration={coil,segment length=5pt,amplitude=4pt},very thick,blue!60!black] (3.0,0.2) -- (3.0,-1.6);
\draw[decorate,decoration={coil,segment length=5pt,amplitude=4pt},very thick,orange!85!black] (7.0,0.2) -- (7.0,-1.6);
\draw[very thick] (2.5,-1.6) -- (3.5,-1.6); \draw[very thick] (6.5,-1.6) -- (7.5,-1.6);
\node[font=\small,blue!60!black] at (2.35,-0.75) {$F_1$};
\node[font=\small,orange!85!black] at (7.65,-0.75) {$F_2$};
\draw[->,thick] (5.9,0.75) arc (60:120:1.8) node[above,font=\scriptsize,pos=0.5] {$M=\ell(F_1-F_2)$};
% motor neurons and reflex circuit
\node[nrn,fill=green!12] (M1) at (1.6,-3.0) {\nt{ACET}};
\node[nrn,fill=green!12] (M2) at (8.4,-3.0) {\nt{ACET}};
\node[nrn,fill=red!10] (G) at (5.0,-3.0) {\nt{GLYC}};
\node[nrn,fill=blue!6] (S) at (1.6,-5.0) {\nt{GLUT}};
\draw[->,very thick] (M1) -- (3.0,-1.7);
\draw[->,very thick] (M2) -- (7.0,-1.7);
\draw[->,thick] (S) -- (M1);
\draw[->,thick] (S) -- (G);
\draw[inh] (G) -- (M2);
\draw[->,thick,densely dashed,gray!60!black] (3.15,-1.2) to[bend left=25] node[pos=0.35,right,font=\scriptsize,black] {신장} (S.north east);
\draw[->,thick,blue!60!black] (-0.3,-3.0) node[left,font=\scriptsize,black,align=right] {뇌의\\명령} -- (M1);
\draw[->,thick,blue!60!black] (10.3,-3.0) node[right,font=\scriptsize,black,align=left] {뇌의\\명령} -- (M2);
\end{tikzpicture}
\caption{관절 하나에 근육 둘, 그리고 가장 빠른 되먹임 루프. \nt{ACET} 뉴런은 뇌의 명령을 받아 관절을 서로 다른 방향으로 당긴다. 힘의 차는 관절을 돌리고, 합은 더 단단하게 만든다. 왼쪽 근육의 신장 센서(\nt{GLUT})는 그 근육 자신의 \nt{ACET} 뉴런을 흥분시키고, \nt{GLYC} 중개 뉴런을 거쳐 길항근의 뉴런을 억제한다. \ref{sec:gaba}장의 금지이다.}
\label{fig:antagonists}
\end{figure}

\section{지연된 되먹임}

근육은 눈을 감고 일하지 않는다. 근육마다 표~\ref{tab:sensors}에서 말한 신장 센서가 있고, 가장 짧은 루프는 척수 안에서 곧바로 닫힌다(그림~\ref{fig:antagonists}). 근육이 갑자기 늘어나면 신장 센서가 그 근육 자신의 \nt{ACET} 뉴런을 흥분시키고, 근육이 수축해 관절을 되돌린다. 동시에 억제성 중개 뉴런을 거쳐 길항근이 꺼져서 방해하지 않는다. 이 중개 뉴런이 바로 \nt{GLYC}이다. 표~\ref{tab:types}에서 자기 장을 갖지 못한 그 유형이며, 그 자리는 바로 여기, 척수의 빠른 루프이다.

루프마다 지연 $d$가 있다. 팔의 척수 루프는 약 $25$--$30\ms$, 피질을 거치는 루프는 그 1.5--3배, 눈을 거치는 루프는 $100$--$200\ms$이다. 그리고 명제~\ref{prop:ei}에서 보았듯이 지연된 되먹임은 시스템을 흔든다. $d$초 전의 정보로 편차 $x$를 속도 $G$로 바로잡는 가장 단순한 조절기 $\dd x/\dd t=-G\,x(t-d)$의 안정 조건은 다음과 같다~\cite{Hayes1950}.
\begin{empheq}[box=\keybox]{equation}
G\,d \;<\; \frac{\pi}{2} ,
\label{eq:delaystable}
\end{empheq}
경계에서 시스템은 주파수 $1/(4d)$로 진동한다. 우리는 이것을 모델로 확인했다. $d=30\ms$일 때 편차는 초당 $G=40$에서 감쇠하고, 경계 $52$를 조금 넘는 초당 $G=55$에서 커지며 진동한다.

여기서 두 가지 결론이 나온다. 첫째, 루프가 길수록 더 느리게 바로잡아야 한다. 지연이 150밀리초인 눈을 거쳐서는 팔을 약 0.1초보다 빨리 바로잡을 수 없다. 그렇지 않으면 팔이 떨린다. 둘째, 척수 루프의 안정 경계 $1/(4\cdot30\ms)\approx8$ Hz는 건강한 사람 누구에게나 있는 미세 떨림의 주파수, 곧 초당 8--12회에 가깝다. 신장 루프는 이 떨림의 유력한 원천 가운데 하나이다~\cite{McAuleyMarsden2000}.

그렇다면 되먹임이 늦는데 뇌는 어떻게 빠르고 정확하게 움직일까? 널리 퍼진 가설에 따르면 \emph{예측한다}. 몸의 내부 모델을 유지하며, 센서를 기다리지 않고 명령의 결과가 어떨지 계산한다~\cite{WolpertGhahramani2000}. 센서는 매 움직임이 아니라 모델을 바로잡는 데 필요하다. 엔지니어라면 이것을 상태 관측기를 갖춘 피드포워드 제어라고 부를 것이다.

\section{운동 리듬 발생기}

걷기, 숨쉬기, 씹기는 리듬 운동이다. 여기에 클럭이 필요할까? 아니다. 척수와 뇌간에는 리듬적인 입력이 전혀 없어도 스스로 리듬을 만들어 내는 작은 네트워크가 있다~\cite{MarderBucher2001}. 가장 단순한 그런 네트워크는 우리가 이미 가진 부품으로 조립된다. 그림~\ref{fig:wta}에서처럼 \nt{GABA} 또는 \nt{GLYC} 중개 뉴런을 거쳐 서로 억제하고, \ref{sec:glutcomputer}장의 기억처럼 천천히 지치는 \nt{GLUT} 뉴런 무리 두 개이다.
\begin{empheq}[box=\keybox]{equation}
\tau\,\frac{\dd r_i}{\dd t} = -r_i + \bigl[I - \beta\,r_j - a_i\bigr]_+ ,
\qquad
\tau_a\,\frac{\dd a_i}{\dd t} = -a_i + b\,r_i ,
\label{eq:halfcentre}
\end{empheq}
여기서 $a_i$는 무리 $i$의 피로, $j$는 다른 무리이다. $\beta>1$인 상호 억제가 승자를 고른다(명제~\ref{prop:wta}). 승자는 일하면서 지친다. 피로가 입력을 잡아먹으면, 이미 쉬고 난 패자가 앞으로 치고 나와 이번에는 자기가 지치기 시작한다. 두 무리는 번갈아 일하고, 각 무리는 쌍의 자기 근육을 제어한다(그림~\ref{fig:halfcentre}).

\begin{figure}[t]
\centering
\begin{tikzpicture}[>=Stealth,x=3.2cm,y=2.4cm]
\draw[->,gray!70!black] (0,0) -- (4.2,0) node[right,font=\small,black] {$t$, s};
\draw[->,gray!70!black] (0,0) -- (0,1.2) node[left,font=\small,black] {$r$};
\foreach \t in {0,1,2,3,4}{\draw[gray!60!black] (\t,-0.03) -- (\t,0.03); \node[below,font=\scriptsize] at (\t,0) {\t};}
\draw[thick,blue!60!black] plot file {figures/halfcentre1.dat};
\draw[thick,orange!85!black] plot file {figures/halfcentre2.dat};
\node[font=\scriptsize,blue!60!black,anchor=west] at (1.6,1.28) {무리 1: “앞으로”};
\node[font=\scriptsize,orange!85!black,anchor=west] at (2.7,1.28) {무리 2: “뒤로”};
\end{tikzpicture}
\caption{모델~\eqref{eq:halfcentre}에 따른 리듬 발생기: $I=1$, $\beta=2$, $b=2$, $\tau=20\ms$, $\tau_a=0.4\,\text{s}$. 상호 억제와 피로를 가진 두 무리가, 입력은 일정한데도 주기 $0.84\,\text{s}$로 번갈아 일한다.}
\label{fig:halfcentre}
\end{figure}

우리 모델에서 입력이 일정할 때 두 무리는 주기 $0.84\,\text{s}$로 교대한다. 피로 시간 $\tau_a$의 약 두 배이다. 위에서 오는 일정한 명령 “걸어라”가 그 자리에서 리듬으로 바뀐다. 이것은 \ref{sec:gaba}장의 \nt{GLUT}--\nt{GABA} 루프 리듬과 같은 또 하나의 국소 리듬이다. 네트워크가 일할 때만 생기고, 자기가 제어하는 근육에만 박자를 준다.

\begin{keyblock}
\begin{proposition}[기계의 출력]
\label{prop:muscles}
뇌는 조절기의 계층으로 근육을 제어한다. 맨 위에서 행동을 고른다(\ref{sec:dofa}장). 그 아래에서 국소 발생기가 일정한 명령을 리듬으로 바꾸고, 척수의 빠른 루프가 관절을 붙잡는다. 힘은 부동소수점으로 정해진다. 지수는 크기 순서대로 켜지는 단위의 수로, 가수는 스파이크 발화율로 정한다. 각 관절은 당기는 근육 한 쌍으로 제어되며, 두 힘의 차가 회전력을, 합이 강성을 정한다. 이 장치들은 측정되었다. 뇌가 몸의 내부 모델에 기댄다는 것은 근거가 탄탄한 가설이다.
\end{proposition}
\end{keyblock}

\section{정리}

\begin{table}[t]
\centering
\footnotesize
\setlength{\tabcolsep}{4pt}
\renewcommand{\arraystretch}{1.15}
\begin{tabular}{>{\raggedright}p{3.0cm}>{\raggedright}p{4.6cm}>{\raggedright\arraybackslash}p{5.2cm}}
\toprule
출력이 하는 일 & 방법 & 알려진 것 \\
\midrule
스파이크를 힘으로 바꿈 & 연축의 합~\eqref{eq:forcesum}, 저역 통과 필터 & 측정됨; 선형 합은 단순화 \\
힘을 조절함 & 크기 순서대로 단위를 켬, 부동소수점~\eqref{eq:sizeprinciple} & 켜지는 순서는 측정됨 \\
당길 줄만 알면서 밂 & 근육 쌍: 회전력은 차, 강성은 합~\eqref{eq:antagonists} & 측정됨 \\
관절을 붙잡음 & 신장 루프, \nt{GLYC}를 통한 길항근 금지 & 측정됨 \\
떨지 않음 & $Gd<\pi/2$~\eqref{eq:delaystable} & 제어 이론에서 나옴; 떨림에 대한 기여는 유력한 원인 \\
빠르게 움직임 & 내부 모델에 의한 예측 & 근거가 탄탄한 가설 \\
리듬 있게 걷고 숨쉼 & 리듬 발생기~\eqref{eq:halfcentre} & 발생기는 측정됨; 가장 단순한 회로는 모델 \\
\bottomrule
\end{tabular}
\caption{기계는 근육을 어떻게 제어하는가.}
\label{tab:muscles}
\end{table}

기계의 출력은 나머지 모든 것과 같은 기법으로 만들어져 있다(표~\ref{tab:muscles}). 부호 대신 두 전선, DAC 대신 저역 통과 필터, 선형 눈금 대신 로그 눈금, 클럭 발생기 대신 상호 억제와 피로이다. 입력(\ref{sec:sensors}장)과 출력은 세상을 상대적인 비율로 재고, 그 사이에서 이 책이 다루는 기계가 일한다.

\section{연습문제}

\begin{exercise}[\lvlA]
\label{ex:13:1}
\emph{숫자로 보는 부동소수점.} 근육에 힘이 $f_n=f_1q^{\,n-1}$인 운동 단위 $N=100$개가 있다. 가장 강한 단위는 가장 약한 단위보다 $100$배 강하다. (a)~$q$, 상대 단계~\eqref{eq:sizeprinciple}, 데시벨 단위의 동적 범위를 구하라. (b)~전체 힘이 같은 동일한 단위 $100$개로 된 “선형 DAC”에서 힘 $10f_1$일 때의 단계는?
\end{exercise}

\begin{exercise}[\lvlA]
\label{ex:13:2}
\emph{안전 여유의 값.} 근육 섬유 위의 시냅스는 스파이크마다 평균 $m$인 푸아송 개수의 소포를 방출하고, 섬유는 소포 $n_{\text{th}}=20$개 이상에서 동작한다. 안전 여유는 $S=m/n_{\text{th}}$이다. 초당 스파이크 $20$개로 일하는 단위는 $S=2$와 $S=4$에서 하루에 몇 번 실패하는가?
\end{exercise}

\begin{exercise}[\lvlB]
\label{ex:13:3}
\emph{필터로서의 근육.} $T_c=50\ms$인 연축~\eqref{eq:twitch}은 선형 시스템의 임펄스 응답이다. (a)~전달 함수 $H(s)$와 차단 주파수를 구하라. (b)~규칙적인 스파이크의 발화율이 얼마일 때 힘의 맥동 폭이 평균 힘의 $10\%$가 되는가?
\end{exercise}

\begin{exercise}[\lvlB]
\label{ex:13:4}
\emph{지연보다 빨리 바로잡을 수는 없다.} 조절기 $\dd x/\dd t=-G\,x(t-d)$. (a)~편차가 진동 없이 가장 빨리 감쇠하는 것은 $Gd=1/e$일 때이며 그 속도가 $1/d$임을 보여라. (b)~척수 루프 $d=30\ms$와 눈을 거치는 루프 $d=150\ms$에 대해 이 $G$와 감쇠 시간 상수를 구하라.
\end{exercise}

\begin{exercise}[\lvlB]
\label{ex:13:5}
\emph{리듬 발생기의 주기.} $\tau\ll\tau_a$일 때 모델~\eqref{eq:halfcentre}의 주기 $T$는 방정식 $(\beta-1)(1-E^{2+b})/(\beta-E)=b\,(1-E^{1+b})/(1+b)$로 정해진다. 여기서 $E=e^{-T/(2\tau_a)}$이다. (a)~$\beta=b=2$, $\tau_a=0.4$~s일 때 $T$를 구하라. (b)~$T$가 입력 $I$와 무관하며, 리듬은 $b>\beta-1$일 때만 가능함을 보여라.
\end{exercise}

\begin{exercise}[\lvlB]
\label{ex:13:6}
\emph{발생기 시뮬레이션.} \texttt{models/\allowbreak muscle\_models.py}에서 함수 \texttt{halfcentre}를 가져와(파일 자체는 실행하지 말 것), 처음 두 주기를 버리고 $b=0.8$, $1.05$, $1.2$, $1.5$, $2$, $4$와 $I=0.5$, $1$, $2$에서 주기를 재라. 명령 “더 빨리 걸어라”가 $I$를 통해 템포를 바꿀 수 있는가?
\end{exercise}

\begin{exercise}[\lvlB]
\label{ex:13:7}
\emph{강성 대 반사.} 관절이 신장 루프로 감싸여 있다: $\dd\theta/\dd t=-a\theta(t)-G\theta(t-d)$, $a=K/B$, $d=30\ms$. 이를 연습문제~\ref{ex:13:4}로 환원해, 편차가 시간 상수 $15\ms$로 진동 없이 감쇠하는 가장 작은 $a$와 필요한 $G$를 구하라. 근육의 공동 긴장이 커지면 $G$를 어떻게 바꿔야 하는가?
\end{exercise}

\begin{exercise}[\lvlC]
\label{ex:13:8}
\emph{템포 노브.} 연습문제~\ref{ex:13:5}와~\ref{ex:13:6}에 따르면 발생기~\eqref{eq:halfcentre}의 입력 $I$는 템포가 아니라 진폭만 바꾼다. 이 책의 부품으로 모델을 최소한으로 바꿔, 어떤 $I_{\min}$ 아래에서는 리듬이 없고 그 위에서는 $I$가 커질수록 주기가 줄어들되 $I$가 세 배가 되면 적어도 절반이 되게 하라. $\tau\ll\tau_a$에서 $I_{\min}$을 구하고 시뮬레이션으로 확인하라.
\end{exercise}

\chapter{통증: 경보 신호}
\chaptermark{통증}
\label{sec:pain}

\section{측정이 아니라 경보}

\ref{sec:sensors}장의 센서는 모두 세상을 쟀다. 빛이 얼마나 되는지, 소리의 높이가 어떤지, 압력이 얼마인지를 쟀다. 통증은 다르게 만들어져 있다. 통증은 “저기 무엇이 있다”가 아니라 “위험하다, 모든 것을 멈춰라”를 알린다. 엔지니어에게 이것은 측정 채널이 아니라 \emph{경보 신호}이다. 기계가 지금 무슨 일을 하든 그것을 뚫고 들어가야 하는 최우선 순위의 선로이다.

경보 장치는 측정 채널과 세 가지 성질에서 다르고, 통증은 세 가지를 모두 갖고 있다. 첫째, 적응으로 약해지기 어렵다. 통증 센서는 촉각 센서보다 훨씬 약하게 적응하고, 가벼운 손상 뒤에는 오히려 더 민감해진다~\cite{LaMotte1982hyper}. 강한 가열 뒤 응답이 약해지는 현상은 한 유형에서만 측정되었다~\cite{Treede1995heat}. 빛 센서는 일정한 배경에서 조용해지지만(그림~\ref{fig:adaptation}), 1분 뒤 조용해지는 경보 장치는 쓸모가 없다. 둘째, 문턱값이 높다. 통증 센서는 자극이 위험에 가까워야만 동작한다. 예를 들어 가열에 대한 문턱값은 유형에 따라 센서마다 $41^\circ$에서 $46$--$48^\circ$까지이다~\cite{Treede1995heat, Weidner1999cfib}. 셋째, 이 장에서 가장 중요한 점으로, 경보의 음량은 센서만이 아니라 기계 자신이 정한다. 같은 상처라도 상황에 따라 다르게 아프다. 이것이 어떤 익숙한 부품으로 조립되어 있는지 보자.

통증 센서는 \ref{sec:sensors}장의 다른 감각 뉴런과 마찬가지로 \nt{GLUT} 뉴런이다. 그중 일부는 글루탐산과 함께 부록~\ref{sec:othermediators}의 P 물질을 방출하는데, 이 물질은 전달을 천천히 강화한다.

\section{두 전선: 빠른 통증과 느린 통증}

손가락을 부딪치면 통증을 두 번 느낀다. 먼저 짧고 날카로운 통증, 1초 뒤에 둔하고 오래가는 통증이다. 이것은 서로 다른 두 전선이다. 빠른 전선은 절연되어 있고 스파이크를 초속 $5$에서 $30$미터의 속도 $v$로 전도한다. 느린 전선은 가늘고 절연이 없으며 초속 $0.5$--$2$미터이다. 발에서 오는 신호는 약 1미터를 지나고, 도착 시간
\begin{empheq}[box=\keybox]{equation}
t \;=\; \frac{L}{v}
\label{eq:paintime}
\end{empheq}
은 빠른 전선에서 수십 밀리초, 느린 전선에서 약 1초이다. 두 전선은 두 메시지를 나른다. 빠른 전선은 \emph{어디}와 \emph{언제}를 말한다. 그 스파이크는 \ref{sec:code}장의 시간 부호에서처럼 더 일찍, 더 정확하게 도착한다. 느린 전선은 \emph{얼마나 나쁜지}를 말하며, 그 길고 둔한 통증은 손상이 가실 때까지 기계를 경보 모드에 붙잡아 둔다.

\section{척수의 관문}

통증 신호가 처음 도착하는 곳은 척수이고, 거기서 신호는 \emph{관문}을 지난다~\cite{MelzackWall1965}. 이 관문의 구체적인 뉴런은 비교적 최근에야 발견되었다~\cite{Duan2014}. 통증을 뇌로 전달하는 출력 \nt{GLUT} 뉴런에는 통증 전선과 촉각 전선이 모인다. 그리고 그 옆에는 촉각이 흥분시키는 억제성 중개 뉴런, 곧 \nt{GABA} 또는 \nt{GLYC}가 있다(그림~\ref{fig:gate}). 촉각은 관문을 닫는다. 원래의 관문 이론~\cite{MelzackWall1965}에서는 반대로 통증이 이 중개 뉴런을 꺼서, 강한 통증은 관문을 연다. 이것은 이론의 가정이며, 발견된 관문 뉴런에서 직접 보이지는 않았다.

\begin{figure}[t]
\centering
\begin{tikzpicture}[>=Stealth,
  nrn/.style={circle,draw,very thick,minimum size=1.0cm,inner sep=0pt,font=\scriptsize},
  inh/.style={-{Bar[width=7pt]},thick,red!70!black},
  bc/.style={->,thick,densely dashed,orange!85!black}]
\node[nrn,fill=blue!6] (Z) at (6,0) {\nt{GLUT}};
\node[nrn,fill=red!10] (Q) at (3.6,-1.6) {\nt{GABA}};
\draw[->,very thick,red!70!black] (0,0.6) node[left,font=\small,black,align=right] {통증 $\nu$} -- (5.45,0.15);
\draw[->,very thick,blue!60!black] (0,-2.6) node[left,font=\small,black,align=right] {촉각 $\mu$} -- (2.9,-2.0);
\draw[->,thick,blue!60!black] (1.8,-2.35) -- (5.5,-0.35) node[pos=0.8,below right=-2pt,font=\scriptsize] {$w_{z\mu}$};
\draw[inh] (1.6,0.5) -- (3.35,-1.1) node[pos=0.5,left=1pt,font=\scriptsize] {$w_{q\nu}$};
\draw[inh] (Q) -- (Z) node[pos=0.5,above left=-1pt,font=\scriptsize] {$\beta$};
\draw[->,very thick] (Z) -- (8.2,0) node[right,font=\small,align=left] {뇌로: $z$};
\draw[bc] (6,2.1) node[above,font=\scriptsize,black,align=center] {위에서: \nt{SERO}, \nt{NORA}, 오피오이드} -- (Z.north) node[pos=0.45,right,font=\scriptsize,black] {이득 $g$};
\end{tikzpicture}
\caption{모델~\eqref{eq:gate}에 따른 척수의 통증 관문. 출력 \nt{GLUT} 뉴런은 통증 $\nu$와 촉각 $\mu$로부터의 약한 흥분을 받는다. 억제성 중개 뉴런(\nt{GABA} 또는 \nt{GLYC})은 촉각에 흥분하고, 관문 이론의 가정에 따르면 통증에 꺼진다. 위에서는 브로드캐스트 채널을 통해 이득 $g$가 온다.}
\label{fig:gate}
\end{figure}

\ref{sec:gaba}장에서처럼 이것을 발화율로 쓰자. 중개 뉴런의 발화율 $q$와 출력 발화율 $z$는 다음과 같다.
\begin{empheq}[box=\keybox]{equation}
q \;=\; \bigl[w_{q\mu}\,\mu + q_0 - w_{q\nu}\,\nu\bigr]_+ ,
\qquad
z \;=\; \bigl[\,g\,(\nu + w_{z\mu}\,\mu) - \beta\,q\,\bigr]_+ ,
\label{eq:gate}
\end{empheq}
여기서 $\nu$는 통증 입력, $\mu$는 촉각 입력, $w$는 연결 가중치(첫 첨자는 받는 쪽, 둘째 첨자는 보내는 쪽), $\beta$는 억제의 세기, $q_0$는 중개 뉴런 자신의 배경 활동, $g$는 위에서 정하는 이득이다(아래에서 다룬다). 이것은 \ref{sec:gaba}장의 금지~\eqref{eq:veto}, 곧 “통증, 단 촉각이 없을 때”이다. 여기에 관문 이론에서 가정으로 가져온 수정이 하나 있다. 강한 통증은 금지 뉴런을 스스로 끈다.

우리는 $w_{q\mu}=1$, $q_0=0.2$, $w_{q\nu}=0.5$, $w_{z\mu}=0.3$, $\beta=1.5$로 모델을 계산했다(그림~\ref{fig:gatecurves}). 촉각이 없으면 통증은 $\nu\approx0.18$부터 통과하기 시작한다. 촉각이 있으면 문턱값이 $0.86$으로 오르고, $\nu=1$에서 출력은 $1$에서 $0.25$로 네 배 줄어든다. 그러나 강한 통증 $\nu=2$에서는 촉각이 더 이상 아무것도 바꾸지 못한다. 관문이 활짝 열렸다. 실생활에서도 그렇다. 멍든 곳을 문지르면 도움이 되지만 심한 부상에는 소용이 없다. 통증 없는 촉각은 관문을 지나지 못한다. 만지는 것만으로는 경보가 울리지 않는다.

\begin{figure}[t]
\centering
\begin{tikzpicture}[>=Stealth,x=5.2cm,y=1.35cm]
\draw[->,gray!70!black] (0,0) -- (2.12,0) node[right,font=\small,black] {통증 $\nu$};
\draw[->,gray!70!black] (0,0) -- (0,4.3) node[left,font=\small,black] {$z$};
\foreach \t in {0,0.5,1,1.5,2}{\draw[gray!60!black] (\t,-0.05) -- (\t,0.05); \node[below,font=\scriptsize] at (\t,0) {\t};}
\foreach \f in {1,2,3,4}{\draw[gray!60!black] (-0.01,\f) -- (0.01,\f); \node[left,font=\scriptsize] at (0,\f) {\f};}
\draw[very thick,red!70!black] plot file {figures/pain_gate_notouch.dat};
\draw[very thick,blue!60!black] plot file {figures/pain_gate_touch.dat};
\draw[very thick,orange!85!black,densely dashed] plot file {figures/pain_gate_sens.dat};
\draw[very thick,red!70!black] (0.08,3.9) -- (0.2,3.9) node[right,font=\scriptsize,black] {촉각 없음};
\draw[very thick,blue!60!black] (0.08,3.45) -- (0.2,3.45) node[right,font=\scriptsize,black] {촉각 있음};
\draw[very thick,orange!85!black,densely dashed] (0.08,3.0) -- (0.2,3.0) node[right,font=\scriptsize,black] {감작 후, $g=2$};
\end{tikzpicture}
\caption{통증 세기에 따른 관문~\eqref{eq:gate}의 출력. 촉각은 문턱값을 올리고 약한 통증과 중간 통증을 누그러뜨리지만, 강한 통증은 누그러뜨리지 못한다. 감작(파선)은 이득을 두 배로 올린다. 같은 통증이 두 배 크게 전달되고 문턱값은 내려간다.}
\label{fig:gatecurves}
\end{figure}

\section{위에서 오는 음량 노브}

관문은 아래에서만 제어되지 않는다. 뇌간에서 관문으로 하행 경로가 내려와 \eqref{eq:gate}의 이득 $g$를 바꾼다. \nt{SERO}, \nt{NORA}, 그리고 부록~\ref{sec:othermediators}의 오피오이드 펩타이드이다~\cite{Fields2004}. 이들은 \ref{sec:serocomputer}--\ref{sec:acet}장과 같은 브로드캐스트 채널이며, 다만 여기서 그 노브는 경보의 음량이다. 노브는 양쪽으로 돌릴 수 있다. 같은 하행 회로가 통증을 누그러뜨릴 수도, 키울 수도 있다.

기계는 왜 경보를 줄이는가? 경보는 인터럽트이고, 인터럽트가 늘 적절하지는 않기 때문이다. 포식자에게서 달아나는 동물은 다친 발 때문에 멈추면 안 된다. 먼저 달리고, 상처는 나중에 핥는다. 통증이 줄어들 것이라는 기대도 통증을 누그러뜨리며, 오피오이드 수용체를 차단하면 이 효과의 일부가 사라진다~\cite{Levine1978}. 뇌에서 계산된 기대가 오피오이드 채널을 타고 관문으로 내려가는 것이다. 이 그림 자체는 측정되었다. 뇌가 음량을 정확히 어떻게 정하는지는 가설이다.

\section{감작: 학습하는 경보}

손상 뒤에는 관문의 출력 뉴런이 더 민감해진다. 같은 입력이 더 큰 출력을 일으키고 문턱값은 내려간다~\cite{Woolf1983}. 모델~\eqref{eq:gate}에서 이것은 이득 $g$의 증가이며, 가장 간단한 기술은 통증 자체가 충전하는 느린 RC 회로이다.
\begin{empheq}[box=\keybox]{equation}
\tau_s\,\frac{\dd g}{\dd t} \;=\; -(g-1) + k_s\,z ,
\label{eq:sensitization}
\end{empheq}
여기서 $\tau_s$는 민감도가 정상으로 돌아오는 시간이다. 손상 자극이 멈춘 뒤에도 감작이 유지되므로~\cite{Woolf1983}, 이 시간은 통증 자체의 지속 시간보다 길다. $k_s$는 충전의 세기이다. 조직이 손상되어 있는 동안 관문의 출력은 $g$를 1보다 높게 유지하고, 상처 근처에서 일어나는 모든 일이 더 크게 전달된다(그림~\ref{fig:gatecurves}의 파선: $g=2$에서 문턱값은 $0.11$로 내려가고 출력은 두 배가 된다). 이것은 경보 장치의 합리적인 설정이다. 손상이 낫기 전까지는 안전하게 가는 편이 낫다.

엔지니어에게 이것은 느린 누설을 가진 양의 되먹임이다. 통증은 이득을 올리고, 이득은 통증을 올린다. 루프가 약하면 상처가 나았을 때 $g$는 정상으로 돌아온다. 그러나 루프가 강하면 경보 장치는 손상이 없어진 뒤에도 켜진 상태에 걸려 있을 수 있다. 그림~\ref{fig:bistable}의 강한 되먹임을 가진 집단과 같다. 모델~\eqref{eq:sensitization}은 단순화이다. 중추 감작이 존재한다는 것은 측정되었다.

\section{교사이자 인터럽트로서의 통증}

기계 안에서 통증은 경보 말고도 두 가지 일을 한다.

\emph{교사.} 통증은 음의 보상이다. 오차 식~\eqref{eq:ctd}에 $r<0$으로 들어간다. \ref{sec:dofa}장에서 말한 모든 것이 여기서도 통한다. 통증의 전조는 미리 오차를 일으키기 시작하고, 네트워크는 통증으로 이어지는 것을 피하도록 배운다. 사람 실험은 통증으로 학습할 때의 뇌 활동이 바로 시간차 오차를 따른다는 것을 보여 주며~\cite{Seymour2004}, \nt{DOPA} 뉴런 일부는 불쾌한 사건에도 응답한다(\ref{sec:dofaalt}장).

\emph{인터럽트.} 통증은 주의를 현재 과제에서 떼어 낸다. 이것은 부작용이 아니라 통증의 목적이다~\cite{EcclestonCrombez1999}. \ref{sec:nora}장에서는 같은 “인터럽트” 역할을 \nt{NORA}의 짧고 날카로운 버스트에 대해 논의했다. 그리고 가장 빠른 응답은 뇌에 닿지도 않는다. 뜨거운 것에서 손을 떼는 반응은 그림~\ref{fig:antagonists}와 같은 근육 쌍과 같은 길항근 금지로 척수 안에서 곧바로 닫힌다.

\begin{keyblock}
\begin{proposition}[경보 신호]
\label{prop:pain}
통증은 측정이 아니라 우선순위가 높은 경보 신호이다. 측정 채널보다 훨씬 약하게 적응하고, 두 전선(빠른 전선은 “어디”, 느린 전선은 “얼마나 나쁜지”)을 따라 가며, 촉각이 닫는(그리고 관문 이론의 가정에 따르면 강한 통증이 여는) 관문을 지나고, 그 음량은 위에서 브로드캐스트 채널이 정한다. 통증에 의한 관문 열림을 빼면 이 장치들은 측정되었다. 간단한 모델~\eqref{eq:gate}와~\eqref{eq:sensitization}은 단순화이다.
\end{proposition}
\end{keyblock}

\section{정리}

\begin{table}[t]
\centering
\footnotesize
\setlength{\tabcolsep}{4pt}
\renewcommand{\arraystretch}{1.15}
\begin{tabular}{>{\raggedright}p{3.0cm}>{\raggedright}p{4.9cm}>{\raggedright\arraybackslash}p{4.9cm}}
\toprule
통증이 하는 일 & 방법 & 알려진 것 \\
\midrule
경보를 울림 & 높은 문턱값, 촉각보다 약한 적응 & 문턱값은 측정됨; 적응 비교는 추정 \\
“어디”와 “얼마나 나쁜지”를 알림 & 속도가 다른 두 전선~\eqref{eq:paintime} & 측정됨 \\
관문을 지남 & \nt{GABA}/\nt{GLYC}를 통한 금지~\eqref{eq:gate} & 관문은 측정됨; 통증에 의한 열림은 이론의 가정; 모델은 단순화 \\
음량을 바꿈 & 하행 \nt{SERO}, \nt{NORA}, 오피오이드 & 측정됨; 음량 선택 규칙은 가설 \\
손상 뒤에 학습함 & 이득의 증가~\eqref{eq:sensitization} & 감작은 측정됨; 모델은 단순화 \\
회피를 가르침 & \eqref{eq:ctd}의 음의 보상 & 시간차 오차와의 유사성은 측정됨 \\
작업을 중단시킴 & 주의 포획; 회피 반사 & 측정됨 \\
\bottomrule
\end{tabular}
\caption{경보 신호로서의 통증.}
\label{tab:pain}
\end{table}

통증은 우리가 이미 본 부품으로 조립되어 있다(표~\ref{tab:pain}). \nt{GLUT} 센서, 두 전선, 억제성 중개 뉴런을 통한 금지, 브로드캐스트 음량 노브, 느린 RC 회로, 예측 오차, 인터럽트이다. 새로운 것은 하나뿐이다. 통증은 기계가 음량을 스스로 정하는 유일한 입력이다. 주인이 누그러뜨릴 수도 있고 다시 설정할 수도 있는 경보 장치이다.

\section{연습문제}

\begin{exercise}[\lvlA]
\label{ex:14:1}
\emph{두 전선.} 신호는 발에서 온다, $L=1$~m. (가)~일제 사격이 같은 유형의 전선 다발을 따라 가고, 전선의 속도는 범위~\eqref{eq:paintime}를 채운다. 빠른 유형과 느린 유형에서 척수에 도착할 때 일제 사격은 얼마나 늘어나는가? (나)~$15$~m/s와 $1$~m/s 전선을 따라 오는 통증의 두 파는 차이가 $0.1$~s 이상이면 구별된다. 어떤 거리부터 합쳐지는가?
\end{exercise}

\begin{exercise}[\lvlB]
\label{ex:14:2}
\emph{촉각이 아플 때.} 본문 매개변수의 관문~\eqref{eq:gate}에서 감작에 따라 이득 $g$가 커진다. $g$가 얼마가 될 때까지 통증 없는 촉각은 어떤 세기 $\mu$에서도 관문을 지나지 못하는가? $g$가 얼마일 때 촉각 $\mu=1$이 통과하기 시작하는가?
\end{exercise}

\begin{exercise}[\lvlB]
\label{ex:14:3}
\emph{감작의 안정성.} 입력은 일정하고, 관문 출력은 $g$에 선형이다: $z=gA-C$, $A=\nu+w_{z\mu}\mu$, $C=\beta q\ge0$. (가)~모델~\eqref{eq:sensitization}의 평형 $g^*$와 그 안정 조건을 구하라. (나)~$k_s=0.5$일 때, 촉각이 없을 때와 촉각 $\mu=1$이 있을 때 모델이 폭주하기 시작하는 통증 세기를 구하라.
\end{exercise}

\begin{exercise}[\lvlA]
\label{ex:14:4}
\emph{통증의 전조.} 시각 $0$의 신호가 시각 $T=2$~s의 통증을 예측한다. \eqref{eq:ctd}에서 $r(t)=-P\,\delta(t-T)$, $\tau_\gamma=2$~s이다. (가)~신호 시각의 $\delta$ 급변의 넓이를 구하라. (나)~시각 $t_e=1$~s의 행동이 통증을 취소한다. 이때 $\delta$ 급변은 얼마이며, 네트워크에게 무엇을 가르치는가?
\end{exercise}

\begin{exercise}[\lvlB]
\label{ex:14:5}
\emph{걸려 버리는 경보.} \texttt{models/\allowbreak pain\_gate.py}의 관문에 동역학~\eqref{eq:sensitization}과 포화 $z\le4$를 더하라. 촉각 $\mu=1$은 늘 있고, 손상 $\nu=2$는 시간 $T$ 동안 이어진 뒤 $\nu=0$이 된다. $k_s=4$, $4.6$, $5$, $6$, $8$일 때, 그 뒤에도 경보가 켜진 채로 남는 가장 작은 $T$($\tau_s$ 단위)를 시뮬레이션으로 구하라.
\end{exercise}

\begin{exercise}[\lvlB]
\label{ex:14:6}
\emph{관문 설계.} $\beta=1.5$, $w_{z\mu}=0.3$, $g=1$에서, 통증 문턱값이 촉각 없이 $0.2$, 촉각 $\mu=1$과 함께 $1.0$이고, $\nu_\times=2.5$에서 촉각이 통증을 더 이상 누그러뜨리지 않도록 \eqref{eq:gate}의 $q_0$, $w_{q\nu}$, $w_{q\mu}$를 고르라. 이득 $g$가 얼마가 될 때까지 통증 없는 촉각이 관문을 지나지 못하는가?
\end{exercise}

\begin{exercise}[\lvlC]
\label{ex:14:7}
\emph{음량 노브.} 작업은 단위 시간당 가치 $V$를 가져오고, 손상 $\nu$를 안고 일하는 비용은 $c\nu$이므로, $\nu>V/c$이면 중단하는 편이 이득이다. 통증은 $z>\theta=0.5$일 때 작업을 중단시킨다. (가)~$V/c=0.25$, $0.5$, $2$일 때 촉각 없는 관문~\eqref{eq:gate}에서 이 문턱값을 주는 이득 $g^*$는? (나)~기계는 이 책의 어떤 신호로 $V$와 $c$를 추정할 수 있을까?
\end{exercise}

\chapter{기계는 움직임을 어떻게 배우는가}
\chaptermark{기계는 움직임을 어떻게 배우는가}
\label{sec:motorlearning}

\section{세 교사}

\ref{sec:muscles}장에서 우리는 기계가 근육을 어떻게 움직이는지 보았고, 한 가설에서 멈췄다. 빠르고 정확한 움직임에는 명령의 결과를 예측하는 몸의 내부 모델이 필요하다는 가설이다~\cite{WolpertGhahramani2000}. 이 모델은 어디서 오는가? 배워야 하고, 지금까지 우리가 배운 방식과는 다르게 배워야 한다.

이 책에는 이미 두 교사가 있었다. 첫째는 \emph{교사 없음}이다. 헵 규칙(\ref{sec:dofa}장)은 자주 함께 일어나는 것을 강화하고 입력을 압축한다. 둘째는 \emph{보상}이다. \nt{DOPA}는 모두에게 하나의 수, 곧 “기대보다 좋은가 나쁜가”를 보낸다. 움직임에는 셋째 교사, \emph{오차}가 필요하다. 손이 컵을 왼쪽으로 2센티미터 빗나갔다. 이것은 “나쁘다”가 아니라 각 출력에게 어느 쪽으로 얼마나 틀렸는지 알려 주는 벡터이다. 엔지니어라면 이것을 지도 학습이라고 부를 것이다.

둘째 교사와 셋째 교사의 차이는 모두에게 전선 하나와 출력마다 따로 전선 하나의 차이이다. 명제~\ref{prop:broadcastcost}에서 하나의 공통 수에 의한 학습은 잡음이 결과에 영향을 주는 뉴런이 하나 늘 때마다 느려졌다. 그러나 각 출력 $i$가 자기 오차 $e_i$를 받으면, 가중치는 가장 단순한 규칙으로 바꿀 수 있다.
\begin{empheq}[box=\keybox]{equation}
\frac{\dd w_{ij}}{\dd t} \;=\; -\eta\,e_i(t)\,x_j(t) .
\label{eq:lms}
\end{empheq}
이것은 적응 필터 기술에서 오래전부터 알려진 최소 제곱 규칙이다~\cite{WidrowHoff1960}. 가중치는 자기 입력과 자기 출력의 오차의 곱에 비례해 변하고, 평균적으로 오차 제곱의 기울기를 따라 내려간다. \ref{sec:dofa}장에서처럼 탐색용 잡음은 필요 없다. 보정의 방향이 오차 전선을 따라 곧바로 온다. 그리고 다른 출력의 오차는 시냅스에 오지 않으므로, 속도는 출력 수가 늘어도 떨어지지 않는다.

\begin{keyblock}
\begin{proposition}[세 교사]
\label{prop:teachers}
헵 규칙은 교사 없이 자주 일어나는 것을 가르친다. \nt{DOPA} 신호는 네트워크 전체에 하나의 수로 이로운 것을 가르치며, 뉴런이 늘 때마다 느려진다. 출력마다 가는 오차 전선은 규칙~\eqref{eq:lms}에 따라 정확한 것을 가르치며, 속도는 출력 수와 무관하다. 셋째 방식의 대가는 출력마다 따로 가는 전선과, 정답을 아는 누군가이다.
\end{proposition}
\end{keyblock}

\section{오차 전선}

움직임의 정답을 누가 알 수 있을까? 센서 자신이다. 손이 목표에서 왼쪽으로 벗어나면 눈과 \ref{sec:sensors}장의 신장 센서가 그것을 알린다. 이 오차를 알맞은 출력으로 배달하는 회로가 필요하다.

뇌에는 바로 이 일에 맞춰 만들어진 것으로 보이는 영역이 있다~\cite{Marr1969,Albus1971}. 그 회로는 이례적으로 규칙적이어서 거의 결정 같고, 그 안에서 규칙~\eqref{eq:lms}를 쉽게 알아볼 수 있다(그림~\ref{fig:errorwire}).

\emph{입력 확장층.} 명령과 몸 상태에 대한 신호가 작은 \nt{GLUT} 뉴런으로 된 거대한 층에 들어온다. 이 뉴런은 약 700억 개로, 뇌의 나머지 전부를 합친 것보다 많다~\cite{Azevedo2009}. 각 뉴런은 몇 개의 입력을 섞고, 신호는 매우 긴 벡터로 펼쳐진다. 엔지니어에게 익숙한 기법이다. 입력의 차원을 높이면 새 공간에서는 거의 어떤 원하는 관계라도 단순한 가중합으로 얻을 수 있다~\cite{Cover1965}.

\emph{출력 뉴런.} 가중합은 약 3000만 개의 큰 출력 뉴런이 계산하며~\cite{Andersen1992cereb}, 각각 확장층으로부터 10만 개 정도의 입력을 받는다. 그리고 이들은 \nt{GABA} 뉴런이다. 학습 결과는 흥분이 아니라 억제로, \ref{sec:gaba}장의 금지처럼 전달된다.

\emph{오차 전선.} 각 출력 뉴런은 특별한 입력을 하나 더 받는다. 정확히 하나의 \nt{GLUT} 전선인데, 초당 약 한 번으로 드물게 동작하지만 동작할 때마다 출력 뉴런에 강력한 폭발을 일으킨다~\cite{Ito2001}. 이것이 바로 $e_i$이다. 폭발의 확률은 움직임 오차의 방향에 따라 달라진다~\cite{Herzfeld2018}. 오차 폭발이 입력 활동과 겹치면 그 입력이 약해진다. \eqref{eq:lms}의 항 $-\eta\,e_i x_j$ 그대로이다.

\begin{figure}[t]
\centering
\begin{tikzpicture}[>=Stealth,
  gran/.style={circle,draw,thick,fill=blue!6,minimum size=0.32cm,inner sep=0pt},
  outn/.style={circle,draw,very thick,fill=red!10,minimum size=1.1cm,inner sep=0pt,font=\scriptsize},
  inh/.style={-{Bar[width=7pt]},very thick,red!70!black}]
% inputs
\foreach \y/\lab in {1.6/{명령},0.4/{몸 상태}}{
  \draw[->,thick,blue!60!black] (-0.6,\y) node[left,font=\scriptsize,black] {\lab} -- (0.35,\y);}
% expansion layer
\foreach \i in {0,...,11}{\node[gran] (g\i) at ({0.8+0.28*mod(\i,2)},{-0.3+0.23*\i}) {};}
\foreach \i in {0,...,11}{\pgfmathsetmacro{\yy}{mod(\i,3)>0 ? 1.6 : 0.4}\draw[gray!60!black] (0.35,\yy) -- (g\i);}
\node[font=\scriptsize,align=center] at (0.95,-1.05) {확장층\\$\sim7\cdot10^{10}$ \nt{GLUT}};
% parallel wires to the output neuron
\foreach \i in {0,...,11}{\draw[->,gray!70!black] (g\i.east) -- (4.1,{1.0+0.07*(\i-5.5)});}
\node[outn] (P) at (4.6,1.0) {\nt{GABA}};
\node[font=\scriptsize,align=center,above=2pt] at (P.north) {출력 뉴런\\입력 $x_j$ $\sim10^5$개, 가중치 $w_{ij}$};
\draw[inh] (P) -- (6.8,1.0) node[right,font=\scriptsize,align=left] {움직임\\보정};
% error wire
\draw[->,very thick,red!70!black,densely dashed] (4.6,-1.4) node[below,font=\scriptsize,black,align=center] {오차 전선 하나 $e_i$\\\nt{GLUT}, 초당 $\sim1$회} -- (P.south);
\end{tikzpicture}
\caption{뇌가 구현하는 것으로 보이는 지도 학습 회로. 명령과 몸 상태는 \nt{GLUT} 뉴런의 거대한 층에서 긴 벡터 $x_j$로 펼쳐진다. 출력 \nt{GABA} 뉴런은 이를 가중치 $w_{ij}$로 더한다. 유일한 오차 전선이 $e_i$를 가져온다. 오차가 활성 입력과 겹치면 규칙~\eqref{eq:lms}에서처럼 그 입력의 가중치가 약해진다.}
\label{fig:errorwire}
\end{figure}

한 가지 어려움은 \ref{sec:dofa}장에서 익숙하다. 오차는 그것을 일으킨 입력보다 늦게 온다. 빗나감은 명령 뒤 수십--수백 밀리초가 지나서야 보인다. 따라서 시냅스는 자기가 활성이었음을 기억해야 한다. 규칙~\eqref{eq:threefactor}에서처럼 적격 흔적이 필요하다. 그리고 실험으로 보면 이 흔적의 길이는 과제마다 되먹임 지연에 맞춰져 있다. 오차가 100밀리초 뒤에 오는 곳에서는 그보다 100밀리초 앞서 활성이었던 입력이 가장 강하게 약해진다~\cite{Suvrathan2016}.

\begin{keyblock}
\begin{proposition}[오차 전선]
\label{prop:errorwire}
그림~\ref{fig:errorwire}의 회로에서 각 출력 뉴런은 정확히 하나의 오차 전선을 받고, 오차와 입력이 겹치면 그 입력이 약해진다. 이것은 거대한 차원으로 확장된 입력에 적용한 최소 제곱 규칙~\eqref{eq:lms}이다. 회로의 구조와 겹칠 때의 약화는 측정되었다. 회로 전체가 지도 학습으로 동작한다는 것은 유력한 가설이다.
\end{proposition}
\end{keyblock}

\section{적응 필터로서의 내부 모델}

이런 회로는 무엇을 배우는가? 가장 자연스러운 답은 예측이다. 확장층에서처럼 시간 상수가 서로 다른 필터 묶음을 거친 명령 $u(t)$가 입력으로 들어온다고 하자: $x_j(t)=(g_j*u)(t)$. 출력은 그 가중합, 곧 센서가 알려 줄 것에 대한 예측이다.
\begin{empheq}[box=\keybox]{equation}
\hat y(t) \;=\; \sum_j w_j\,(g_j*u)(t), \qquad e(t) \;=\; \hat y(t) - y(t) ,
\label{eq:adaptivefilter}
\end{empheq}
가중치는 \eqref{eq:lms}에 따라 학습한다. 이것이 \emph{적응 필터}이다. 미지의 시스템을 재현하도록 자기 전달 함수를 스스로 맞추는 회로이다~\cite{Fujita1982}. 여기서 미지의 시스템은 질량, 탄성, 지연을 가진 자기 몸이다. 학습된 필터가 바로 \ref{sec:muscles}장의 내부 모델이다. 센서를 기다리지 않고 센서가 무엇을 알려 줄지 말해 준다.

이런 모델은 두 번 쓸모 있다. 첫째, 그 예측을 늦게 오는 되먹임 대신 넣을 수 있고, 그러면 안정 조건~\eqref{eq:delaystable}이 더 이상 손발을 묶지 않는다. 모델에 따라 빠르게 바로잡을 수 있다. 둘째, 예측한 것과 받은 것의 차이는 \emph{예상 밖}의 것에 대한 신호이다. 기계가 스스로 한 일은 모두 모델이 예측해 빼냈고, 남는 것은 세상이 한 일이다. 그래서 한 가설에 따르면 우리는 스스로를 간지럽힐 수 없다~\cite{Blakemore1998}.

\section{세상이 바뀔 때: 빠른 것과 느린 것}

쉽게 해 볼 수 있는 실험으로 회로를 확인하자. 사람이 목표를 향해 손을 뻗는데 세상이 갑자기 바뀐다. 예를 들어 팔에 옆으로 힘이 가해진다. 처음 움직임들은 빗나가고, 그다음 빗나감이 줄어든다. 힘을 없애면 손은 처음에 \emph{반대쪽}으로 빗나간다. 모델이 힘을 배웠고 계속 그것을 보상하는 것이다. 이것은 그냥 “잘되기 시작했다”가 아니라 모델이 학습되었다는 직접 증거이다.

실험은 더 미묘한 것도 보여 준다. 두 과정이 동시에 학습한다. 빨리 배우고 빨리 잊는 빠른 과정과, 천천히 배우고 오래 기억하는 느린 과정이다~\cite{Smith2006}. 보정은 매 움직임 뒤에, 사건 단위로 갱신된다.
\begin{empheq}[box=\keybox]{equation}
e = p - (x_f + x_s), \qquad x_f \leftarrow A_f\,x_f + B_f\,e, \qquad x_s \leftarrow A_s\,x_s + B_s\,e ,
\label{eq:tworate}
\end{empheq}
여기서 $p$는 섭동, $x_f$와 $x_s$는 두 과정이 학습한 보정, $A$는 보정이 다음 움직임까지 얼마나 유지되는지, $B$는 오차로부터 얼마나 강하게 배우는지이다. 빠른 과정은 $B$가 크고 $A$는 1보다 눈에 띄게 작다. 느린 과정은 반대이다.

\begin{figure}[t]
\centering
\begin{tikzpicture}[>=Stealth,x=0.034cm,y=1.9cm]
\draw[->,gray!70!black] (0,0) -- (355,0) node[right,font=\small,black] {움직임};
\draw[->,gray!70!black] (0,-1.15) -- (0,1.25);
\foreach \v in {-1,1}{\draw[gray!60!black] (-3,\v) -- (3,\v); \node[left,font=\scriptsize] at (0,\v) {$\v$};}
\foreach \n in {100,200,300}{\draw[gray!60!black] (\n,-0.04) -- (\n,0.04); \node[below,font=\scriptsize] at (\n,0) {\n};}
\draw[thin,gray!70!black] plot file {figures/tworate_p.dat};
\draw[thick,orange!85!black] plot file {figures/tworate_fast.dat};
\draw[thick,blue!60!black] plot file {figures/tworate_slow.dat};
\draw[very thick,black] plot file {figures/tworate_net.dat};
\node[font=\scriptsize,gray!50!black] at (120,1.12) {섭동 $p$};
\node[font=\scriptsize,black,anchor=west] at (238,0.52) {합: 복귀};
\node[font=\scriptsize,blue!60!black,anchor=west] at (150,0.39) {느린 $x_s$};
\node[font=\scriptsize,orange!85!black,anchor=west] at (60,0.08) {빠른 $x_f$};
\draw[densely dashed,gray!60!black] (226,-1.15) -- (226,1.25);
\node[font=\scriptsize,align=center,anchor=south west] at (228,-1.1) {오차가\\더 없음};
\end{tikzpicture}
\caption{모델~\eqref{eq:tworate}에 따른 두 학습 과정, $A_f=0.92$, $B_f=0.1$, $A_s=0.996$, $B_s=0.02$. 섭동 $p=1$로 이백 번 움직인 뒤, 합이 0으로 돌아올 때까지 반대 섭동 $p=-1$로 여섯 번 움직인다. 파선 뒤로는 오차가 더 이상 들어오지 않고, 합은 스스로 옛 보정 쪽으로 돌아간다. 빠른 과정은 반대 섭동을 잊었고, 느린 과정은 원래 섭동을 기억한다.}
\label{fig:tworate}
\end{figure}

모델~\eqref{eq:tworate}는 이것 없이는 이해하기 어려운 실험을 설명한다(그림~\ref{fig:tworate}). 우리 계산에서 섭동이 있는 이백 번의 움직임 동안 보정의 합은 $0.84$에 이르고, 그 대부분인 $0.63$을 느린 과정이 쥐고 있다. 이어서 반대 섭동으로 여섯 번 움직이면 합이 0으로 돌아온다. 빠른 과정은 음수 $-0.53$으로 갔고, 느린 과정은 아직 양수 $0.46$이다. 이제 오차를 알려 주지 않으면 빠른 과정은 수십 번의 움직임 동안 자기 보정을 잊고, 느린 과정은 훨씬 오래 걸리며, 합은 마흔 번의 움직임 동안 \emph{저절로} $0.37$까지 올라간다. 옛 기술이 아무 훈련 없이 돌아오는 것이다. 이런 자발적 회복은 사람에서 측정되었다. 과정이 하나인 모델은 이것을 줄 수 없다. 오차가 없으면 그냥 0으로 감쇠한다.

왜 과정이 둘일까? \ref{sec:acet}장의 최적 필터가 그럴듯한 공학적 답을 준다. 몸은 여러 원인으로 여러 속도로 변한다. 근육은 몇 분 만에 지치고, 몇 달에 걸쳐 자라거나 약해진다. 교란에 빠른 것과 느린 것이 있다면, 가장 좋은 추정은 시간 상수가 다른 두 필터로 이루어지고, 각 필터가 자기 교란을 잡는다~\cite{Kording2007}. 빠른 과정은 일시적 보정, 곧 “팔이 지쳤다”이고, 느린 과정은 “팔이 달라졌다”이다. 아직 가설이지만, 하루에 배운 기술이 다음 날까지 일부 사라지는 이유와 두 번째에는 더 빨리 배워지는 이유를 설명한다.

\section{정리}

\begin{table}[t]
\centering
\footnotesize
\setlength{\tabcolsep}{4pt}
\renewcommand{\arraystretch}{1.15}
\begin{tabular}{>{\raggedright}p{3.0cm}>{\raggedright}p{4.6cm}>{\raggedright\arraybackslash}p{5.2cm}}
\toprule
회로가 하는 일 & 방법 & 알려진 것 \\
\midrule
교사와 함께 배움 & 출력마다 가는 오차 전선, 규칙~\eqref{eq:lms} & 회로 구조와 겹칠 때의 약화는 측정됨; 지도 학습은 유력한 가설 \\
입력을 확장함 & \nt{GLUT} 뉴런 700억 개 & 뉴런 수는 측정됨; 확장의 역할은 이론 \\
오차의 지연을 고려함 & 지연에 맞춘 적격 흔적 & 측정됨 \\
내부 모델을 만듦 & 적응 필터~\eqref{eq:adaptivefilter} & 근거가 탄탄한 가설 \\
두 속도로 배움 & 빠른 과정과 느린 과정~\eqref{eq:tworate} & 후효과와 자발적 복귀는 측정됨; 두 과정은 모델 \\
\bottomrule
\end{tabular}
\caption{기계는 움직임을 어떻게 배우는가.}
\label{tab:motorlearning}
\end{table}

이제 기계에는 세 교사가 있다(표~\ref{tab:motorlearning}). 헵은 자주 일어나는 것을 압축하고, \nt{DOPA}는 하나의 수로 이로운 것을 고르도록 가르치며, 오차 전선은 정확성을 가르친다. 그리고 각 출력에 자기 오차가 가므로 빨리 가르친다. 뇌는 셋을 모두 쓰는 것으로 보이며, 각각을 가장 싼 곳에서 쓴다. 브로드캐스트 신호는 몇 안 되는 결정에, 개별 전선은 정답을 센서가 직접 알려 주는 몸의 정밀 조정에 쓴다.

\section{연습문제}

\begin{exercise}[\lvlA]
\label{ex:15:1}
\emph{오차 전선 회로의 용량.} 회로에 입력이 $10^5$개씩인 출력 뉴런 $3\cdot10^7$개가 있고, 각각 자기 오차 전선(초당 $1$회)을 갖는다. 입력이 $n$개인 뉴런은 벡터 $\approx2n$개의 임의의 표지를 학습한다~\cite{Cover1965}. (a)~회로는 연합을 몇 개 저장하는가? (b)~$\Delta t=10\ms$에서 \eqref{eq:spikeentropy}로 뉴런의 용량을 채우는 최소 시간을 추정하라.
\end{exercise}

\begin{exercise}[\lvlB]
\label{ex:15:2}
\emph{최소 제곱 규칙의 속도.} 규칙~\eqref{eq:lms}의 모드는 속도 $\eta\lambda_k(C)$로 감쇠한다. 백색 잡음 위의 필터~\eqref{eq:adaptivefilter} 다섯 개에 대해 $C_{jk}=2\sqrt{\tau_j\tau_k}/(\tau_j+\tau_k)$이다. $\tau_j$가 $2$배씩($10$--$160\ms$), 그리고 $4$배씩 커질 때 가장 빠른 속도와 가장 느린 속도의 비를 구하라.
\end{exercise}

\begin{exercise}[\lvlB]
\label{ex:15:3}
\emph{두 과정의 분석.} 그림~\ref{fig:tworate}의 매개변수로 모델~\eqref{eq:tworate}를 $\mathbf x_{n+1}=M\mathbf x_n+\mathbf b\,p$로 쓰라. (a)~$M$의 고윳값으로부터 움직임 단위의 시간 상수를 구하라. (b)~일정한 $p=1$에서 정상 상태의 $x_f$, $x_s$와 그 합을 구하라.
\end{exercise}

\begin{exercise}[\lvlB]
\label{ex:15:4}
\emph{시뮬레이션: 두 번째는 더 빠르다.} \texttt{models/\allowbreak motor\_learning.py}의 매개변수로 모델~\eqref{eq:tworate}에서 $p=1$일 때 $x_f+x_s\ge0.8$이 될 때까지의 움직임 수를 구하라. (a)~학습하지 않은 기계; (b)~그림~\ref{fig:tworate}에서처럼 $p=1$로 $200$번 움직이고 합이 0이 될 때까지 반대 섭동을 준 뒤. 과정이 하나인 모델도 이렇게 빨라질 수 있는가?
\end{exercise}

\begin{exercise}[\lvlB]
\label{ex:15:5}
\emph{내부 모델을 가진 조절기.} 대상 $\dd x/\dd t=ku$는 지연 $d$를 두고 관측된다. 조절기 $u=-G\hat x$는 $\hat x(t)=x(t-d)+\hat k\int_{t-d}^{t}u\,\dd s$로 예측한다. (a)~$\hat k=k=1$, $d=150\ms$에서 시간 상수 $20\ms$를 주는 $G$를 구하고 한계~\eqref{eq:delaystable}와 비교하라. (b)~$\hat k=0.4k$일 때 회로는 안정한가?
\end{exercise}

\begin{exercise}[\lvlB]
\label{ex:15:6}
\emph{적응 필터가 근육을 배운다.} 대상은 넓이가 1인 연축~\eqref{eq:twitch}, $T_c=50\ms$이다. 필터~\eqref{eq:adaptivefilter}는 $\tau_j=12.5$, $25$, $50$, $100$, $200\ms$인 RC 회로이고, 입력은 $10\ms$마다 새로운 정규 난수이다. $\eta=50$과 $200$에서 규칙~\eqref{eq:lms}가 오차를 $0.5\%$로 낮추는 데 몇 초가 걸리는가? 왜 $300$~s 뒤에도 가중치는 최적과 거리가 먼가?
\end{exercise}

\begin{exercise}[\lvlC]
\label{ex:15:7}
\emph{최적 필터로서의 두 과정.} 섭동은 잡음 분산이 $q_f$, $q_s$인 과정 $p^{f,s}_{n+1}=A_{f,s}\,p^{f,s}_n+w^{f,s}_n$의 합이고, 오차는 분산 $R$의 잡음과 함께 관측된다. 칼만 필터는~\eqref{eq:tworate}를 준다. $A_f=0.92$, $A_s=0.996$에서 $B_f=0.1$, $B_s=0.02$가 나오는 $q_f/R$, $q_s/R$은? $q_s$를 네 배로 하면 $B_f$, $B_s$는 어떻게 바뀌는가?
\end{exercise}

\chapter{두 번째 출력: 뇌는 어떻게 몸의 화학을 조절하는가}
\chaptermark{화학적 출력}
\label{sec:chemoutput}

\section{빠른 출력과 느린 출력}

\ref{sec:muscles}장에서 우리는 기계의 빠른 출력인 근육을 보았다. 그러나 기계에는 느린 두 번째 출력도 있다. 뇌는 체온, 심박수, 몸속 물과 당의 양을 필요한 범위 안에 유지하고, 몸을 달리기나 잠에 대비시킨다. 이를 위해 뇌는 관절을 움직이는 대신 화학 신호를 내보낸다. 경로는 두 가지다. 첫째는 내장 기관, 즉 심장, 혈관, 장, 샘으로 곧장 가는 신경이다. 둘째는 혈액이다. 일부 뉴런과 샘은 물질을 이웃 세포와의 틈이 아니라 곧바로 혈류로 내보낸다.

혈액은 이 책에 등장한 버스 가운데 가장 넓게 방송하는 버스다. \ref{sec:serocomputer}--\ref{sec:acet}장의 브로드캐스트 채널은 신호를 뇌의 넓은 영역에 보냈다. 혈액은 신호를 몸의 모든 세포에 보낸다. 혈액은 약 1분 만에 몸을 한 바퀴 돌므로, 메시지는 수십 초 안에 모든 곳에 닿고 물질이 분해될 때까지 몇 분에서 몇 시간 동안 유지된다. 이런 방송에는 주소가 전혀 없다. 명제~\ref{prop:receptor}에서처럼 메시지의 의미는 수신자가 정한다. 그 물질에 대한 수용체를 가진 세포만 반응한다.

\section{가속 페달과 브레이크: 기관으로 가는 두 가닥의 전선}

가장 좋은 예는 심장이다. 심장에는 \ref{sec:serocomputer}장의 \nt{SERO} 뉴런처럼 자체 페이스메이커가 있다. 페이스메이커는 스스로 리듬을 만들며, 신경을 모두 끊어도 심장은 1분에 약 100번 뛴다. 뇌는 리듬을 만들지 않고 조정만 한다. 그것도 부호가 반대인 두 가닥의 전선으로 한다(그림~\ref{fig:gasbrake}). 한 전선으로는 \nt{NORA}가 온다. 이것이 리듬을 빠르게 하는 \emph{가속 페달}이다. 다른 전선으로는 \nt{ACET}이 온다. 이것이 리듬을 늦추는 \emph{브레이크}다. 대략
\begin{empheq}[box=\keybox]{equation}
f \;\approx\; f_0 \;+\; a_S\,S \;-\; a_P\,P ,
\label{eq:gasbrake}
\end{empheq}
이다. 여기서 $f_0\approx100$회/분은 페이스메이커의 고유 리듬이고, $S$와 $P$는 가속 페달과 브레이크의 활동이다. 쉬고 있을 때는 브레이크가 밟혀 있어서 심장은 보통 1분에 60--70번쯤 뛴다. 운동할 때는 브레이크를 놓고 가속 페달을 밟는다.

\begin{figure}[t]
\centering
\begin{tikzpicture}[>=Stealth,
  blk/.style={draw,very thick,rounded corners=2pt,align=center,font=\small},
  pace/.style={circle,draw,very thick,double,double distance=1.2pt,minimum size=1.8cm,inner sep=0pt,font=\scriptsize,fill=red!8,align=center},
  inh/.style={-{Bar[width=7pt]},thick,red!70!black}]
\node[blk,fill=blue!5,text width=1.6cm] (B) at (0,0) {뇌의\\명령};
\node[pace] (H) at (6.2,0) {페이스\\메이커};
\draw[->,very thick,orange!85!black] (B.north east) to[bend left=20] node[above,font=\scriptsize,black,align=center] {\nt{NORA}: 가속 페달, 수 초} (H.north west);
\draw[inh,very thick,blue!60!black] (B.south east) to[bend right=20] node[below,font=\scriptsize,black,align=center] {\nt{ACET}: 브레이크, 1초 미만} (H.south west);
\draw[->,very thick] (H) -- (8.4,0) node[right,font=\small] {박동수 $f$};
\node[blk,fill=orange!10,text width=1.6cm] (A) at (0,-2.6) {\nt{ADRE}\\혈액으로};
\draw[->,thick,orange!85!black] (B) -- (A);
\foreach \y/\lab in {-2.0/{심장},-2.6/{혈관},-3.2/{간, 근육}}{
  \draw[->,thick,densely dashed,orange!85!black] (A.east) -- (4.3,\y) node[right,font=\scriptsize,black] {\lab};}
\end{tikzpicture}
\caption{심장 페이스메이커로 가는 부호가 반대인 두 가닥의 전선: 가속 페달 \nt{NORA}는 느리고 브레이크 \nt{ACET}은 빠르다. 아래는 같은 가속 페달을 브로드캐스트 버스로 보내는 경로다. \nt{ADRE} 세포가 아드레날린을 혈액으로 내보내면, 알맞은 수용체를 가진 모든 기관이 신호를 받는다(점선 화살표).}
\label{fig:gasbrake}
\end{figure}

이것은 \ref{sec:glutcomputer}장의 두 전선 부호이자 \ref{sec:muscles}장의 근육 쌍이다. 다만 이제 전선이 관절이 아니라 페이스메이커의 템포를 조절한다. 그리고 두 전선은 속도가 다르다. 둘 다 저역 통과 필터처럼 작동하지만, 브레이크는 가속 페달보다 몇 배 빠르게 변화를 통과시킨다~\cite{Berger1989}. \nt{ACET}의 경로는 짧다. 수용체와 결합한 단백질이 세포 안의 2차 전달자 없이 칼륨 채널을 직접 연다~\cite{Logothetis1987gbg}. 반면 \nt{NORA}는 세포 안의 연쇄 반응을 거쳐 작용한다. 엔지니어라면 속도가 다른 두 액추에이터를 쓰는 기법을 알아볼 것이다. 빠른 쪽은 즉각적인 미세 조정을 맡고, 느린 쪽은 크고 오래가는 변화를 맡는다.

여기서 \nt{ADRE}도 제 역할을 찾는다. 표~\ref{tab:types}에서는 뇌 안에서 \nt{ADRE}의 역할이 불분명하다고 적었다. 뇌 밖에서는 분명하다. 가속 페달 경로의 세포 일부는 기관으로 전선을 보내지 않고 아드레날린을 곧바로 혈액으로 내보낸다. 이것은 같은 가속 페달이지만 브로드캐스트 버스를 탄다. 수십 초 안에 심장, 혈관, 간, 근육에 한꺼번에 닿고 몇 분 동안 유지된다. 주소가 있는 가속 페달 \nt{NORA}는 기관 하나를 조정하고, 브로드캐스트 가속 페달 \nt{ADRE}는 몸 전체를 달리기 모드로 전환한다.

\section{되먹임이 걸린 캐스케이드}

많은 호르몬은 샘 하나가 아니라 캐스케이드가 분비한다. 뇌 바닥에 있는 작은 뉴런 무리가 첫 번째 신호를 혈액으로 내보낸다. 이 신호는 중간 샘이 두 번째 신호를 내보내게 하고, 두 번째 신호는 최종 샘이 호르몬을 내보내게 한다. 몸에 작용하는 것은 이 호르몬이다. 그리고 최종 호르몬은 앞 단계로 돌아가 그 단계들을 억제한다. 그 결과 음의 되먹임으로 감싼 3단 증폭기, 즉 식~\eqref{eq:feedback}의 \nt{SERO} 시스템과 같은 레벨 조절기가 된다.

각 단을 시상수 $\tau$와 이득 $g$를 가진 RC 회로로, 되먹임을 계수 $k$로 나타내자.
\begin{equation}
\tau\,\frac{\dd x_1}{\dd t} = -x_1 + g\bigl[u - k\,x_3\bigr]_+ ,\qquad
\tau\,\frac{\dd x_2}{\dd t} = -x_2 + g\,x_1 ,\qquad
\tau\,\frac{\dd x_3}{\dd t} = -x_3 + g\,x_2 ,
\label{eq:cascade}
\end{equation}
여기서 $u$는 뇌의 명령, $x_3$은 최종 호르몬의 농도다. 루프 이득은 $K=g^3k$이다. 평형에서 $x_3=g^3u/(1+K)$이며, 되먹임이 있는 모든 조절기처럼 루프 이득이 크면 농도가 각 단 이득의 편차에 둔감해진다. 그러나 각 단은 위상을 늦춘다. 주파수 $\omega=\sqrt3/\tau$에서 같은 단 세 개를 합치면 위상이 반 바퀴 늦어지고 진폭은 8분의 1로 줄어든다. 여기서 제어 이론의 고전적인 결과가 나온다.
\begin{empheq}[box=\keybox]{equation}
K \;<\; 8 ,
\label{eq:cascadestable}
\end{empheq}
그렇지 않으면 조절기는 주기 약 $2\pi\tau/\sqrt3\approx3.6\,\tau$로 진동하기 시작한다.

\begin{keyblock}
\begin{proposition}[정밀도 대 안정성]
\label{prop:cascade}
비례 되먹임이 걸린 3단 캐스케이드에서 농도 오차는 $1/(1+K)$이고, 캐스케이드는 $K<8$일 때만 안정하다. 따라서 이런 조절기는 농도를 약 10퍼센트보다 정확하게 유지하지 못한다. 이득을 높이려 하면 조절기는 발진기로 바뀐다.
\end{proposition}
\end{keyblock}

우리는 이를 모델~\eqref{eq:cascade}로 확인했다(그림~\ref{fig:cascade}). $K=4$에서는 켜진 뒤 농도가 조금 흔들리다가 자리를 잡는다. $K=7$에서도 자리를 잡지만 느리다. $K=12$에서는 진동이 잦아들지도 않고 끝없이 커지지도 않는다. 각 단은 음의 신호를 낼 수 없으므로, 캐스케이드는 평균 농도의 $0.83$배와 $1.24$배 사이를 오가는 일정한 맥동에 들어가며 주기는 $3.7$--$3.9\,\tau$이다.

\begin{figure}[t]
\centering
\begin{tikzpicture}[>=Stealth,x=0.3cm,y=1.2cm]
\draw[->,gray!70!black] (0,0) -- (41.5,0) node[right,font=\small,black] {$t/\tau$};
\draw[->,gray!70!black] (0,0) -- (0,2.8) node[left,font=\small,black] {$x_3/x_3^*$};
\foreach \t in {0,10,20,30,40}{\draw[gray!60!black] (\t,-0.05) -- (\t,0.05); \node[below,font=\scriptsize] at (\t,0) {\t};}
\foreach \f in {1,2}{\draw[gray!60!black] (-0.3,\f) -- (0.3,\f); \node[left,font=\scriptsize] at (0,\f) {\f};}
\draw[densely dashed,gray!60!black] (0,1) -- (40,1);
\draw[thick,blue!60!black] plot file {figures/cascade_K4.dat};
\draw[thick,red!70!black] plot file {figures/cascade_K12.dat};
\node[font=\scriptsize,blue!60!black,anchor=west] at (16,0.55) {$K=4$: 자리를 잡는다};
\node[font=\scriptsize,red!70!black,anchor=west] at (16,1.75) {$K=12$: 맥동한다};
\end{tikzpicture}
\caption{식~\eqref{eq:cascade}에 따른 되먹임이 걸린 3단 캐스케이드. 최종 호르몬의 농도를 평형값 $x_3^*$에 대한 비로 나타냈다. 루프 이득이 8보다 작으면 농도가 자리를 잡고, 8보다 크면 캐스케이드가 스스로 일정한 맥동을 만든다.}
\label{fig:cascade}
\end{figure}

실제 캐스케이드 호르몬은 정말로 맥동하며 분비된다. 예를 들어 스트레스 호르몬의 농도는 약 1시간 주기로 오르내린다. 한 가설에 따르면 이 맥동은 단계 사이의 지연된 되먹임 자체가 만들어 내며, 따로 리듬 발생기가 필요하지 않다~\cite{Walker2010}. 이는 명제~\ref{prop:ei}의 \nt{GLUT}--\nt{GABA} 루프와 조건~\eqref{eq:delaystable}의 근육 신장 루프에서 진동이 생긴 이유와 같다. 지연되는 음의 되먹임이다.

\section{호르몬의 펄스 부호}

맥동은 부작용일 뿐 아니라 부호가 되기도 한다. 생식을 조절하는 뉴런은 자신의 신호를 대략 1시간에 한 번씩 짧은 분량으로 나누어 혈액에 내보낸다. 수신 샘에 같은 신호를 분량으로 나누지 않고 연속으로 주면, 샘은 펄스를 받을 때처럼 제대로 일하지 않고 아예 멈춘다. 다시 1시간에 한 번씩 분량으로 주면 작동이 회복된다~\cite{Belchetz1978}. 즉 수신자는 농도가 아니라 분량의 \emph{빈도}를 읽는다.

엔지니어에게 여기에는 수수께끼가 없다. 물질이 오래 작용하면 지쳐서 반응을 멈추는 수용체는 \ref{sec:code}장의 고갈되는 시냅스~\eqref{eq:depression}처럼 고역 통과 필터다. 연속 신호는 막고 변화는 통과시킨다. 이런 수신자에게는 펄스로만 무언가를 전달할 수 있으므로, 정보는 농도에서 분량의 빈도와 모양으로 옮겨 간다. 게다가 분량의 빈도가 다르면 같은 샘 안의 서로 다른 세포에 다르게 작용하므로, 화학 회선 하나로 둘 이상의 값을 보낼 수 있다. \ref{sec:dofaalt}장에서 \nt{DOPA} 전선 하나로 빠른 성분과 느린 성분이 함께 전달된 것과 같다.

\section{일주기 리듬: 모드를 바꾸는 느린 발진기}

몸에서 가장 느린 리듬은 하루 주기의 리듬이다. 이 리듬은 세포 안의 지연된 음의 되먹임이 만든다. 유전자가 단백질을 만들고, 그 단백질이 몇 시간 늦게 자기 생산을 억제한다~\cite{TakahashiJS2017}. 다시 지연되는 음의 되먹임이 나온다. 이 책에서 네 번째다. 그리고 다시 리듬이 생기는데, 이번에는 주기가 약 하루다. 주 발진기는 뇌 바닥에 있는 수만 개의 뉴런으로 이루어진 작은 무리이며, 이들의 리듬이 몸의 나머지 세포를 동기화한다.

사람에서 이 발진기의 고유 주기는 정확히 하루가 아니라 평균 $24.18$시간이다~\cite{Czeisler1999}. 날마다 눈의 센서에서 오는 빛이 발진기를 맞춰 준다(\ref{sec:sensors}장). 전자 공학자에게 이것은 \emph{위상 고정 루프}다. 고유 주파수 $\omega_0$인 발진기가 주파수 $\omega$인 외부 신호에 맞춰지며, 위상차 $\varphi$는 다음 식을 따른다.
\begin{empheq}[box=\keybox]{equation}
\frac{\dd\varphi}{\dd t} \;=\; (\omega-\omega_0) \;-\; K_\varphi\,\sin\varphi .
\label{eq:pll}
\end{empheq}
$|\omega-\omega_0|<K_\varphi$이면 고정이 유지된다. 주기 $24.18$시간인 발진기는 날마다 약 11분씩 옮겨져야 하며, 보통은 아침 햇빛만으로 충분하다. 극야나 빛이 없는 동굴에서는 맞춤이 없으므로 리듬은 고유 주기로 흘러간다.

일주기 발진기는 컴퓨터의 클록 발진기가 아니다. 계산 단계를 세지도 않고 뉴런의 스파이크를 동기화하지도 않는다. 이 발진기는 \emph{모드}를 바꾼다. 잠과 깨어 있음, 호르몬 농도, 체온, 먹을 준비를 바꾼다. 이것은 \ref{sec:serocomputer}--\ref{sec:acet}장의 레지스터와 같은 종류의 느린 브로드캐스트 레지스터다. 다만 그 값이 태양에 맞춘 일정표에 따라 바뀐다.

\begin{keyblock}
\begin{proposition}[화학적 출력]
\label{prop:chemoutput}
기계의 느린 출력은 빠른 출력과 같은 부품으로 만들어진다. 기관은 부호가 반대이고 속도가 다른 두 가닥의 전선, 즉 가속 페달 \nt{NORA}와 브레이크 \nt{ACET}을 받고, 몸 전체는 아드레날린 \nt{ADRE}를 비롯한 혈액의 브로드캐스트 신호를 한꺼번에 받는다. 호르몬 농도는 음의 되먹임이 걸린 캐스케이드가 유지하며, 그 정밀도는 안정성에 제한된다. 일부 호르몬은 분량의 빈도로 부호화된다. 일주기 리듬은 빛에 위상 고정되는 느린 발진기다. 경로의 구조, 분량 실험, 주기 실험은 측정되었다. 캐스케이드의 맥동을 되먹임 자체로 설명하는 것은 가설이다.
\end{proposition}
\end{keyblock}

\section{요약}

\begin{table}[t]
\centering
\footnotesize
\setlength{\tabcolsep}{4pt}
\renewcommand{\arraystretch}{1.15}
\begin{tabular}{>{\raggedright}p{3.1cm}>{\raggedright}p{4.8cm}>{\raggedright\arraybackslash}p{4.9cm}}
\toprule
화학적 출력이 하는 일 & 방법 & 알려진 것 \\
\midrule
기관을 조정한다 & 가속 페달 \nt{NORA}와 브레이크 \nt{ACET}~\eqref{eq:gasbrake} & 측정됨; 브레이크가 가속 페달보다 빠르다 \\
몸 전체를 달리기 모드로 바꾼다 & 혈액 속 아드레날린 \nt{ADRE} & 측정됨 \\
호르몬 농도를 유지한다 & 되먹임이 걸린 캐스케이드, $K<8$~\eqref{eq:cascadestable} & 캐스케이드와 되먹임은 측정됨; 루프 자체가 만드는 맥동은 가설 \\
분량의 빈도로 전달한다 & 고역 통과 필터로서의 지친 수용체 & 분량 의존성은 측정됨 \\
하루 동안 모드를 바꾼다 & 빛에 위상 고정되는 발진기~\eqref{eq:pll} & 주기와 빛에 의한 맞춤은 측정됨 \\
\bottomrule
\end{tabular}
\caption{기계는 몸의 화학을 어떻게 조절하는가.}
\label{tab:chemoutput}
\end{table}

기계에는 출력이 두 개 있다(표~\ref{tab:chemoutput}). 빠른 출력은 근육이다. 밀리초 단위이고, 주소가 있으며, 관절 하나하나로 간다. 느린 출력은 몸의 화학이다. 초, 분, 하루 단위이고, 두 가닥의 전선으로 기관에, 혈액으로 모든 곳에 한꺼번에 간다. 두 출력의 부품은 같다. 부호가 반대인 두 전선, 페이스메이커, 되먹임과 그 지연이다. 기계 내부를 이루는 부품도 바로 이것들이다.

\section{연습문제}

\begin{exercise}[\lvlA]
\label{ex:16:1}
\emph{필터로서의 혈액.} 호르몬은 반감기 $t_{1/2}=10$~min으로 혈액에서 제거되고, $T=60$~min마다 짧은 분량으로 분비된다. 농도의 최댓값은 최솟값의 몇 배이고, 분량의 기본 고조파는 몇 배로 감쇠하는가? $t_{1/2}$가 얼마일 때 최댓값이 최솟값의 2배에 불과한가?
\end{exercise}

\begin{exercise}[\lvlA]
\label{ex:16:2}
\emph{일주기 위상 고정.} 발진기~\eqref{eq:pll}의 고유 주기는 $24.18$~h이고, 빛은 위상을 하루에 최대 1시간 옮긴다: $K_\varphi=2\pi/24$~rad/day. 정상 상태 위상차 $\varphi^*$를 분 단위로, 그리고 이완 시간을 구하라. 외부 신호가 $\pm6$~h 도약한 뒤 불일치가 1시간 아래로 줄어드는 데 며칠이 걸리는가?
\end{exercise}

\begin{exercise}[\lvlB]
\label{ex:16:3}
\emph{정밀도 대 안정성.} 선형화한 캐스케이드~\eqref{eq:cascade}를 같은 단 $n$개로 늘렸다. 임계 이득 $K_c(n)$과, $n=3$과 $n=6$에서 달성할 수 있는 최소 오차 $1/(1+K_c)$를 구하라. $n\to\infty$일 때 이 오차는 어디로 수렴하는가?
\end{exercise}

\begin{exercise}[\lvlB]
\label{ex:16:4}
\emph{가속 페달과 브레이크.} 식~\eqref{eq:gasbrake}에서 가속 페달과 브레이크의 기여는 $\tau_S=2$~s와 $\tau_P=0.4$~s인 RC 회로의 출력이고, 명령은 각각 분당 $80$회와 $60$회를 넘지 않는다. $f_0=100$이고, 안정 시 브레이크 $35$, 가속 페달 $0$, $f=65$이다. $f=120$에 도달하는 데 얼마나 걸리는가? 브레이크를 놓지 않고 가속 페달만 쓰면 어떤가?
\end{exercise}

\begin{exercise}[\lvlB]
\label{ex:16:5}
\emph{모델링: 지연이 있는 캐스케이드.} \texttt{models/\allowbreak chem\_models.py}의 함수 \texttt{cascade}를 복사해(파일은 실행하지 말 것) 되먹임에 지연 $d$를 넣어라. 첫 단이 $x_3(t-d)$를 본다. $d/\tau=0$, $0.5$, $1$, $2$에서 $K_c$와 경계에서의 주기를 구하고, $0.9K_c$와 $1.1K_c$에서 시뮬레이션으로 확인하라.
\end{exercise}

\begin{exercise}[\lvlB]
\label{ex:16:6}
\emph{분량을 읽는 수신자.} 수용체가 지친다: $y=[c-a]_+$, $\tau_a\,\dd a/\dd t=c-a$, $\tau_a=30$~min. 호르몬은 평균 용량 $\bar c$는 같게, 한 번에 $6$~min씩 주기 $T$로 들어온다. $T=15$, $60$, $120$~min과 $T\to\infty$에서 평균 반응을 구하라. 농도가 $\bar c$로 일정하면 반응은 얼마인가?
\end{exercise}

\begin{exercise}[\lvlC]
\label{ex:16:7}
\emph{발진기인가, 되먹임인가?} 캐스케이드~\eqref{eq:cascade}의 지연된 되먹임이 만드는 맥동과 별도의 페이스메이커가 만드는 맥동을 구별하는 실험을 제안하라. $K$를 1.5배밖에 바꿀 수 없고 주기를 $5\%$ 정확도로 측정한다면, 주기의 이동으로 두 가설을 구별할 수 있는가?
\end{exercise}

\chapter{기계 디버깅}
\chaptermark{기계 디버깅}
\label{sec:debugging}

\section{회로도 없이 기계를 디버깅하는 법}

회로도 없이 작동하는 보드를 받은 엔지니어는 네 가지 기법으로 디버깅한다. \emph{프로브}를 대고 전선에 무엇이 흐르는지 본다. \emph{시험 신호를 넣고} 무엇이 바뀌는지 본다. 회로의 일부를 \emph{끊고} 무엇이 멈추는지 본다. 그리고 \emph{제어 레지스터를 읽어} 보드가 어떤 모드로 동작하는지 알아낸다. 이 책 전체가 뇌의 회로도를 읽으려는 시도다. 이 장에서는 네 기법 가운데 엔지니어가 이미 쓸 줄 아는 것이 무엇이고, 앞 장들이 무엇을 기대하라고 말해 주는지 살펴본다.

오늘날 이런 디버깅의 가장 눈에 띄는 예는 뇌-컴퓨터 인터페이스다. 뉴런의 스파이크를 읽어 외부 기계를 위한 명령으로 바꾸거나, 반대로 바깥에서 뇌로 신호를 넣는 장치다. 이 장의 대부분은 이 인터페이스, 특히 Neuralink가 하는 일을 다룬다.

\section{프로브: 전극은 무엇을 듣는가}

뇌의 프로브는 조직에 꽂은 가느다란 전극이다. 전극은 반경 약 100마이크로미터 안에 있는 뉴런, 보통 한 개에서 몇 개의 스파이크를 듣는다. 전극에 잡힌 스파이크는 수십에서 수백 마이크로볼트, 길이 약 1밀리초의 파형이며, 그 모양으로 이웃 뉴런을 서로 구별할 수 있다. 가장 잘 들리는 것은 큰 \nt{GLUT} 뉴런이다. 피질에서 대다수를 차지하고(표~\ref{tab:types}) 전류도 더 세다.

주된 어려움은 바로 보인다. 오늘날 좋은 프로브는 전극 수백에서 수천 개이지만, 뇌의 뉴런은 $8.6\cdot10^{10}$개다. 우리는 기계의 대략 1억분의 1을 엿듣는다. 그렇게 보잘것없는 표본에서 어떻게 무언가를 알 수 있을까? 답은 \ref{sec:code}장에 있다. 이웃한 뉴런들의 잡음은 서로 상관되어 있어서, 무리에 대한 평균은 한계에 부딪힌다(명제~\ref{prop:pooling}). 뉴런 100개는 1000개와 거의 같은 양을 전달한다. 게다가 운동 피질이 푸는 문제는 저차원이다. 팔의 관절은 많지 않고(\ref{sec:muscles}장), 팔을 제어하는 수백 개 뉴런의 활동은 공통 변수 열 개 남짓으로 이루어진 공간에 들어간다~\cite{Gallego2017}. 뉴런 100개짜리 프로브는 이 변수를 거의 다 듣는다. 만약 뇌가 뉴런마다 독립된 메시지를 저장한다면 이런 엿듣기는 쓸모가 없을 것이다.

\section{데이터 흐름: 프로브에서 바로 압축하기}

프로브는 엄청난 흐름을 낸다. 스파이크의 모양을 보려면 각 전극을 초당 약 $2$만 번, 약 10비트 정밀도로 디지털화해야 한다. 전극 1000개는
\begin{empheq}[box=\keybox]{equation}
\begin{aligned}
R_{\text{원시}} \;&=\; N\,f_s\,b \;\approx\; 10^3\cdot2\cdot10^4\cdot10 \;\approx\; 2\cdot10^8\ \text{bit/s},\\
R_{\text{스파이크}} \;&\approx\; N\,r\,\log_2\frac{e}{r\,\Delta t} \;\approx\; 8\cdot10^4\ \text{bit/s}
\end{aligned}
\label{eq:probedata}
\end{empheq}
를 낸다. 두 번째 추정은 $r=10$~spikes/s, 정밀도 $\Delta t=1\ms$에서의 스파이크 흐름의 용량~\eqref{eq:spikeentropy}이다. 흐름에 실제로 들어 있는 것은 2500분의 1 크기다. 이식 장치에는 이것이 결정적인 차이다. 머리 안에서 조직을 데우지 않고 쓸 수 있는 전력으로는 원시 흐름을 피부를 통해 무선으로 보낼 수 없다. 그래서 이식 장치는 전극 옆 칩에서 바로 스파이크를 골라내고, 이벤트, 즉 채널 번호와 스파이크 시각만 밖으로 보내야 한다. Neuralink 시스템을 처음 소개한 논문은 아직 거기까지 가지 못했다. 칩이 신호를 증폭하고 디지털화하면 원시 흐름 전체가 케이블로 나가고, 스파이크는 외부 프로그램이 골라낸다. 칩 위의 압축과 무선 통신은 계획으로 언급되어 있다~\cite{Musk2019}.

익숙한 기법이다. 눈은 긴 케이블로 보내기 전에 신호를 100분의 1로 압축한다(\ref{sec:sensors}장). 이벤트 카메라는 프레임이 아니라 변화를 보낸다. 프로브도 전선에 맞추려면 뇌 자신이 하는 일을 해야 한다. 스파이크로, 그리고 새 소식만 말하는 것이다.

\section{Neuralink가 하는 일}

Neuralink 장치는 이런 제약에 맞춰 만든 프로브다~\cite{Musk2019}. 딱딱한 바늘 배열 대신, 머리카락보다 가는 유연한 실 여러 가닥에 전극이 줄지어 달려 있다. 시스템을 처음 소개한 논문에서는 실 $96$가닥에 전극 최대 $3072$개, 실마다 $32$개였다. 사람에게 이식한 장치는 회사 발표에 따르면 실 $64$가닥에 약 1000채널이다. 유연한 실은 조직을 덜 손상시키고, 두개골 안에서 뇌가 늘 조금씩 움직이는 것도 더 잘 견딘다. 실은 로봇이 혈관을 피해 꽂는다. 앞서 말했듯이 첫 논문에서는 원시 흐름~\eqref{eq:probedata}이 케이블로 나갔다. 회사 발표에 따르면 사람용 장치는 다르다. 본체는 완전히 피부로 덮이고, 스파이크는 안에서 골라내며, 데이터는 무선으로 나가고, 전력은 선 없이 들어온다.

첫 장치의 과제는 읽기다. 실을 팔 움직임을 계획하는 피질 부위에 꽂으면, 디코더가 스파이크를 화면의 커서 움직임으로 바꾼다. 팔이 마비된 사람이 움직임을 상상해서 컴퓨터를 조작한다. 발상 자체는 새롭지 않다. 전극 100개 남짓의 딱딱한 배열을 쓴 인터페이스는 오래전부터 커서를 조작하게 해 주었고~\cite{Hochberg2006}, 손으로 글씨 쓰는 움직임을 상상해 글을 쓰게 했으며~\cite{Willett2021}, 말하려는 시도를 분당 약 60단어의 속도로 화면의 글로 바꾸기까지 했다~\cite{Willett2023}. Neuralink의 새로움은 공학에 있다. 채널 수, 무선, 밀폐된 본체, 로봇 삽입, 즉 실험실 밖에서 몇 년 동안 지닐 수 있는 프로브를 만들려는 시도다. 우리가 아는 한 사람 대상 시험 결과는 주로 동료 심사 논문이 아니라 회사 자체가 발표했다. 회사는 특히 삽입 뒤 실 일부가 밀려나 작동하는 채널 수가 줄었다고 밝혔다. 디버깅에서는 흔한 골칫거리다. 보드에 대해 움직이는 프로브는 이미 다른 전선을 듣는다.

회사의 두 번째 방향은 쓰기다. 시력을 잃은 사람의 시각 피질에 신호를 넣는 시각 장치다. 이에 대해서는 아래 쓰기 절에서 다룬다.

\section{읽기: 수신자로서의 디코더}

인터페이스의 디코더는 명제~\ref{prop:reader}의 의미에서 수신자다. 메시지의 어느 부분을 읽을지 스스로 정한다. 사실상 모든 인터페이스는 \emph{무리의 발화율}을 읽는다. 각 채널의 스파이크를 수십 밀리초 창마다 세고, 의도한 움직임이 나오도록 고른 가중치로 더한다. 이것은 \ref{sec:rate}절의 무리 발화율 부호이며, 우연히 고른 것이 아니다. 창마다 스파이크가 몇 개뿐이면 뉴런 하나의 발화율은 정해지지 않지만, 100개를 합치면 이미 쓸 만하다.

정보를 얼마나 꺼낼 수 있을까? “생각 속의 손”으로 글쓰기는 분당 약 90자, 즉 초당 1.5자의 속도였다~\cite{Willett2021}. 글자당 5비트로 잡아도 초당 8비트가 안 된다. Neuralink 발표에 따르면 커서 조작은 초당 약 10비트다. 한편 상한~\eqref{eq:spikeentropy}은 뉴런 \emph{하나}에 초당 수십 비트, 100개면 수천 비트다. 인터페이스는 엿들은 뉴런이 전달할 수 있는 양의 작은 일부만 꺼낸다. 병목은 전선이 아니다. 무리가 독립 변수를 몇 개나 전달하는지(명제~\ref{prop:pooling}), 그리고 \ref{sec:code}장에서 보았듯이 아직 완전히 해독되지 않은 부호를 디코더가 얼마나 잘 이해하는지가 병목이다.

\ref{sec:motorlearning}장에서 익숙한 두 번째 특징도 있다. 디코더와 뇌는 서로에게서 배운다. 뇌는 커서가 어디로 갔는지 보고 오차를 받아 명령을 조정한다. 오차 전선을 통한 운동 학습과 같다. 엔지니어는 때때로 디코더의 가중치를 다시 맞춘다. 실 아래의 뉴런이 바뀌고 네트워크의 연결이 다시 학습되기 때문이다. 결국 서로에게 맞춰 가는 두 학습자가 생긴다. 이 쌍이 폭주하지 않게 하려면 학습 속도를 떼어 놓는 것이 편하다. 한쪽은 빠르게, 다른 쪽은 느리게 배우게 한다.

\begin{keyblock}
\begin{proposition}[뉴런 1000개짜리 프로브가 작동하는 이유]
\label{prop:probe}
보잘것없는 비율의 뉴런을 엿듣는 인터페이스가 움직임을 제어할 수 있는 것은, 무리의 활동이 저차원이고 그 잡음이 상관되어 있기 때문이다. 뉴런 수백 개가 공통 변수를 거의 다 전달한다. 같은 이유로 인터페이스는 초당 몇 비트밖에 꺼내지 못한다. 스파이크 흐름이 담을 수 있는 양이 아니라 과제의 독립 변수 개수만큼이다. 인터페이스의 성능은 측정되었다. 부호를 이해하면 디코더가 얼마나 좋아질지는 열린 문제다.
\end{proposition}
\end{keyblock}

\section{쓰기: 읽기보다 어렵다}

기계에 신호를 넣는 것은 읽는 것보다 어렵다. 전류를 흘리는 전극은 뉴런 하나가 아니라 주변의 모든 뉴런과 전선을 유형과 부호를 가리지 않고 흥분시킨다. \emph{어느} 뉴런이 \emph{언제} 발화했는지가 중요한 부호(\ref{sec:code}장)를 이런 전극으로 쓸 수는 없다. \emph{어디에}, \emph{얼마나 세게}를 대략 정할 수 있을 뿐이다.

시각에는 이것으로 어느 정도 충분하다. 시각 피질의 전극에 전류를 흘리면 사람은 시야의 특정 위치에서 빛나는 점을 본다. 점을 동시에, 최대 16개까지 켜자 시력을 잃은 사람이 몇몇 글자를 알아보고 물체의 경계를 구별했다~\cite{Fernandez2021}. 이것은 위치 부호이며, 쓸 수 있다. 그러나 \ref{sec:sensors}장을 떠올려 보자. 눈은 뇌로 픽셀이 아니라 압축된 설명, 즉 가장자리, 변화, 밝아짐과 어두워짐을 따로따로의 전선으로 보낸다. 피질에 “빛의 점”을 쓰는 것은 입력에서 전혀 다른 부호를 기다리는 기계에 픽셀을 쓰는 것이다. 그래서 인공 시각은 아직 전극 수로 기대할 만한 수준보다 거칠다. 전극이 필요 없는 시험 신호도 있다. 착시는 눈을 통해 특이한 입력을 넣어 기계를 평소 모드에서 벗어나게 하고, 오류 하나하나가 각자의 메커니즘을 가리킨다(\ref{sec:illusions}절).

\section{프로브가 보지 못하는 것}

전극은 주소 버스, 즉 \nt{GLUT} 뉴런의 스파이크를 듣고, 작은 \nt{GABA} 뉴런은 덜 잘 듣는다. 그러나 \ref{sec:serocomputer}--\ref{sec:acet}장에서 보았듯이 기계 전체의 동작 모드는 브로드캐스트 레지스터, 즉 \nt{SERO}, \nt{DOPA}, \nt{NORA}, \nt{ACET}의 농도가 정한다. 전극은 이것을 듣지 못한다. 원천 뉴런이 극히 적고, 그 신호는 프로브 옆의 스파이크가 아니라 부피 속의 농도이기 때문이다. 데이터 버스는 보지만 제어 레지스터를 보지 못하는 디버거는 늘 놀라게 된다. 같은 입력이 다른 응답을 내는데, 어딘가에서 $g$, $a$, $\tau_\gamma$가 바뀌었기 때문이다. 농도는 조직 안의 화학 센서로 잴 수 있지만~\cite{Bunin1998}, 인터페이스에서는 아직 이런 센서를 쓰지 않는다. 기계의 느린 두 번째 출력인 혈액 속 호르몬(\ref{sec:chemoutput}장)도 프로브의 시야 밖에 있다. 호르몬은 전극이 아니라 혈액 시료로 잰다.

프로브는 가중치도 보지 못한다. 그런데 기억과 학습한 것의 거의 전부가 바로 가중치에 저장된다. 가중치는 응답이 어떻게 바뀌는지로 짐작할 수 있을 뿐이다. 또 프로브는 사람이 느끼는 방식의 통증도 보지 못한다. \ref{sec:pain}장에서 보았듯이 경보 신호의 세기는 기계 자신, 즉 척수의 관문과 위에서 오는 조절기가 정하므로, 센서로는 얼마나 아픈지 읽을 수 없다.

\section{요약: 디버깅과 열린 문제}

\begin{table}[t]
\centering
\footnotesize
\setlength{\tabcolsep}{4pt}
\renewcommand{\arraystretch}{1.15}
\begin{tabular}{>{\raggedright}p{2.2cm}>{\raggedright}p{3.6cm}>{\raggedright}p{3.4cm}>{\raggedright\arraybackslash}p{4.0cm}}
\toprule
디버깅 기법 & 인터페이스가 하는 일 & 책이 설명하는 것 & 알려진 것 \\
\midrule
프로브 & 뉴런 수백--수천 개를 듣는다 & 저차원성과 상관된 잡음(\ref{sec:code}장) & 측정됨 \\
프로브에서 압축 & 원시 신호 대신 이벤트를 보낸다 & 스파이크 흐름의 용량~\eqref{eq:spikeentropy} & 공학적으로 해결됨 \\
읽기 & 무리 발화율 $\to$ 움직임, 글, 말 & 디코더는 수신자, 무리 발화율 부호 & 작동함; 초당 몇 비트 \\
공동 학습 & 뇌와 디코더가 서로에게 맞춘다 & 오차 전선(\ref{sec:motorlearning}장) & 관찰됨; 맞춤의 규칙은 연구 대상 \\
쓰기 & 전류가 시야에 점을 켠다 & 위치 부호는 쓸 수 있고 타이밍 부호는 못 쓴다(\ref{sec:sensors}장) & 점은 측정됨; 쓸모 있는 시각은 아직 없음 \\
레지스터 읽기 & 거의 하지 않는다 & 브로드캐스트 채널(\ref{sec:serocomputer}--\ref{sec:acet}장) & 화학 센서는 실험에 있음 \\
\bottomrule
\end{tabular}
\caption{기계 디버깅: 뇌-컴퓨터 인터페이스가 할 수 있는 일과 이 책의 각 장이 그에 대해 말하는 것.}
\label{tab:debugging}
\end{table}

표~\ref{tab:debugging}이 이 장을 요약한다. 뇌-컴퓨터 인터페이스는 이미 주소 버스에 프로브를 대고 초당 몇 비트를 읽을 수 있다. 이것이 애초에 작동한다는 사실은 \ref{sec:code}장의 저차원성과 상관된 잡음으로 설명된다. 기계에 쓰는 일은 거칠게만 할 수 있다. 기계의 입력 부호가 압축되어 있고 개별 전선에 주소가 매겨져 있기 때문이다. 그리고 기계를 학습 가능하고 조정 가능하게 만드는 제어 레지스터와 가중치는 아직 프로브에 거의 보이지 않는다.

\ref{sec:synthesis}장의 열린 문제 하나하나는 디버거의 과제이기도 하다. 어떤 부호를 읽을지는 프로브가 더 조밀해지고 디코더가 발화율만이 아니라 스파이크 타이밍을 쓸 줄 알게 되면 정해질 것이다. 다층 네트워크가 어떻게 학습하는지는 여러 층에 동시에 프로브를 대고 학습하면서 응답이 어떻게 바뀌는지 보면 확인할 수 있다. 브로드캐스트 채널이 무엇을 전달하는지는 전기 프로브에 화학 프로브가 더해지면 보일 것이다. 이 책은 기계의 동작과 모든 엔지니어가 아는 법칙으로 기계의 회로도를 읽으려는 시도다. 지금 만들어지고 있는 디버거는 이 회로도를 프로브로 확인할 수 있게 해 줄 것이다.

\section{연습문제}

\begin{exercise}[\lvlA]
\label{ex:17:1}
\emph{무선 채널 예산.} 피부를 통한 전송은 비트당 $1$~nJ이 들고, 송신기는 최대 $10$~mW를 쓸 수 있다. $N=1000$ 전극에서 원시 흐름~\eqref{eq:probedata}과 스파이크 흐름에 필요한 전력은 얼마인가? 원시 흐름이라면 비트당 에너지가 얼마여야 하는가?
\end{exercise}

\begin{exercise}[\lvlA]
\label{ex:17:2}
\emph{뉴런 100개에 든 비트.} 디코더는 $r=20$~spikes/s인 뉴런 $N$개의 스파이크를 $50\ms$ 창마다 더한다. 잡음의 쌍별 상관은 $c=0.1$이고, 신호는 무리 발화율을 (제곱평균으로) $30\%$ 바꾼다. $N=10$, $100$, $1000$, $N\to\infty$에서 창당 속도 $\tfrac12\log_2(1+\text{SNR})$을 추정하라. $c=0$, $N=1000$이면?
\end{exercise}

\begin{exercise}[\lvlB]
\label{ex:17:3}
\emph{키보드 인터페이스의 속도.} 인터페이스가 초당 1.5번 $N=32$개 글자 가운데 하나를 고른다. 의도한 글자는 확률 $P$로 나오고 오류는 똑같이 확률이 나뉜다. $P=0.99$, $0.94$, $0.8$에서 초당 몇 비트를 전달하는가? $P=0.94$에서 $10$~bit/s를 내려면 초당 몇 글자가 필요한가?
\end{exercise}

\begin{exercise}[\lvlB]
\label{ex:17:4}
\emph{모델링: 무리 디코더.} 뉴런 $N$개, $\lambda_i=[b+m\,\mathbf v\cdot\mathbf p_i]_+$, $b=20$, $m=15$~spikes/s, 창 $50\ms$, $\mathbf v$는 성분마다 표준편차 $0.5$. 모든 뉴런이 공통 잡음이 섞인 $\mathbf v+\boldsymbol\xi$, $\sigma_\xi=0.2$를 본다. 비편향 무리 벡터 디코더를 시뮬레이션하라. $N$이 얼마일 때 두 잡음이 같아지며, $N\to\infty$에서 창당 성분마다 몇 비트인가?
\end{exercise}

\begin{exercise}[\lvlB]
\label{ex:17:5}
\emph{두 학습자.} 뇌는 명령 $m=b\,v^*$를, 디코더는 $v=d\,m$을 낸다. $\langle v^{*2}\rangle=1$. 둘 다 시도마다 $\tfrac12(v-v^*)^2$에 대해 속도 $\eta$로 경사 하강을 한 걸음 한다. $b=1.2$, $d=1$에서 시작해 $\eta=0.8$, $1.1$, $1.5$, $2$로 시뮬레이션하라. 각각은 $\eta<2$에서 안정하다. 쌍은 $\eta$가 얼마일 때 안정한가?
\end{exercise}

\begin{exercise}[\lvlB]
\label{ex:17:6}
\emph{설계: 프로브 파이프라인.} 채널 $1024$개, 채널당 초당 $2$만 샘플, 뉴런 $10$~spikes/s, 무선 $1$~nJ/bit. 샘플의 백색 잡음에 대해 어떤 문턱값이면 채널당 거짓 스파이크가 $0.1$~Hz 이하인가? $20\ms$마다 스파이크 수를 $2$비트로 보낸다면 전력은 얼마이고, 원시 $10$비트 샘플보다 몇 배 작은가?
\end{exercise}

\begin{exercise}[\lvlC]
\label{ex:17:7}
\emph{인터페이스 속도의 한계.} 대역폭 $W=5$~Hz인 변수 $D=10$개에 대해 $\text{SNR}\approx0.9$(신호를 닮은 잡음)와 $90$(뉴런 $1000$개의 독립 잡음)일 때 $R\le DW\log_2(1+\text{SNR})$을 추정하라. 측정값은 약 $10$~bit/s다. 격차의 두 설명, 즉 신호를 닮은 잡음과 부호를 모른다는 것을 어떤 실험으로 구별하겠는가?
\end{exercise}

\chapter{전기적 기능에 따른 뉴런의 유형}
\chaptermark{계측기로서의 뉴런}
\label{sec:eetypes}
\begin{chaptergoals}
\item 뉴런이 동작하는 소자들: 적분기, 필터, 변환기, 발진기, 래치;
\item 세포 하나가 대역 통과 공진기가 되거나 \nt{SERO}가 켜는 래치가 되는 방식;
\item 되먹임이 누설 뉴런으로 무누설 적분기를 만드는 방식과 정밀한 조정이 필요한 이유;
\item 세포 유형이 아니라 이온 채널과 브로드캐스트 채널이 정하는 한 세포의 모드인 이유.
\end{chaptergoals}

\ref{sec:neurontypes}장에서 우리는 뉴런을 그 출력이 내보내는 신경전달물질에 따라 분류했다. 전자 공학자라면 다르게 분류할 것이다. 뉴런이 자기 신호에 무엇을 하는지, 즉 어떤 소자로 동작하는지에 따라서다. 이 책의 주요 소자는 \ref{sec:glutcomputer}장부터 알고 있다. 누설 적분 뉴런~\eqref{eq:glutneuron}, 곧 비교기가 달린 RC 회로다. 그러나 뇌에는 다른 소자도 있다. 표~\ref{tab:eetypes}는 이들을 네 무리로 모았으며, 거의 모두 이 책에 이미 나왔다.

\begin{center}
\footnotesize
\setlength{\tabcolsep}{3pt}
\renewcommand{\arraystretch}{1.15}
\begin{longtable}{>{\raggedright}p{3.1cm}>{\raggedright}p{3.3cm}>{\raggedright}p{4.7cm}>{\raggedright\arraybackslash}p{2.4cm}}
\caption{소자로서의 뉴런: 이름, 전자 공학의 대응물, 소자가 하는 일, 책에서 나오는 곳.}\label{tab:eetypes}\\
\toprule
뉴런 & 소자 & 하는 일 & 책에서 나오는 곳 \\
\midrule
\endfirsthead
\toprule
뉴런 & 소자 & 하는 일 & 책에서 나오는 곳 \\
\midrule
\endhead
\bottomrule
\endfoot
\multicolumn{4}{l}{\emph{필터와 적분기}} \\
누설 적분 뉴런 & 비교기가 달린 RC 적분기 & 창 $\tau$ 안의 입력을 더하고 문턱값에서 발화한다 & \ref{sec:glutcomputer}장, \eqref{eq:glutneuron} \\
동시성 검출기 & 시간 창이 있는 AND & 입력이 거의 함께 올 때만 발화한다; 매우 작은 $\tau$ & \ref{sec:glutcomputer}, \ref{sec:code}장, \eqref{eq:window} \\
위상성 뉴런 & 고역 통과 필터, 에지 검출기 & 입력의 변화에 반응하고 곧 잠잠해진다 & \ref{sec:sensors}장 \\
적응 뉴런 & 자동 이득 조절이 달린 고역 통과 필터 & 입력이 일정하면 발화율이 떨어진다 & \ref{sec:sensors}장, \eqref{eq:nakarushton} \\
공진기 & 공진 회로, 대역 통과 필터 & 자기 주파수의 입력에 가장 잘 반응한다 & \ref{sec:resonator}절 \\
무누설 적분기 & 연산 증폭기 적분기 & 속도를 위치로 바꾸고 그 위치를 유지한다 & \ref{sec:perfectint}절 \\
\midrule
\multicolumn{4}{l}{\emph{변환기}} \\
긴장성 뉴런 & 전압---주파수 변환기 & 발화율이 입력을 선형으로 따르고 거의 지치지 않는다 & \ref{sec:code}장 \\
등급 전위 뉴런(스파이크 없음) & 아날로그 증폭기, 버퍼 & 매끄러운 전압을 짧은 거리로 전달한다 & \ref{sec:sensors}장 \\
정류 뉴런 & 반파 정류기 & 발화율은 음수가 될 수 없다, $[u]_+$ & \ref{sec:glutcomputer}, \ref{sec:gaba}장 \\
“켜짐---꺼짐” 쌍 & 정류기 쌍, 이중 레일 부호 & 하나는 “더 크다”에, 다른 하나는 “더 작다”에 반응한다 & \ref{sec:glutcomputer}장, \eqref{eq:dualrail} \\
분로 억제 뉴런 & 아날로그 나눗셈기, 이득 조절 & 입력에서 빼지 않고 입력을 나눈다 & \ref{sec:gaba}장, \eqref{eq:shunt} \\
\midrule
\multicolumn{4}{l}{\emph{발진기와 시간}} \\
페이스메이커 & 이완 발진기, 멀티바이브레이터 & 스스로 일정한 주파수로 스파이크를 낸다 & \ref{sec:serocomputer}, \ref{sec:dofa}장 \\
입력이 있는 페이스메이커 & 전압 제어 발진기 & 자기 리듬이 있고 입력이 이를 빠르게 하거나 늦춘다 & \ref{sec:chemoutput}장 \\
버스트 뉴런 & 게이트형 버스트 발생기 & 잦은 스파이크의 버스트, 즉 “대시”를 낸다 & \ref{sec:code}장, \eqref{eq:burst} \\
위상 고정 뉴런 & 위상 검출기, 위상 고정 루프 & 입력 파동의 특정 위상에서 발화한다 & \ref{sec:sensors}, \ref{sec:chemoutput}장 \\
지연 축삭 & 지연선 & 신호를 정해진 시간만큼 늦춘다 & \ref{sec:glutcomputer}장, 그림~\ref{fig:itd} \\
\midrule
\multicolumn{4}{l}{\emph{기억과 전환}} \\
쌍안정 뉴런 & 슈미트 트리거, 래치 & 짧은 입력이 오래가는 상태로 전환한다 & \ref{sec:bistable}절 \\
가지돌기 가지가 있는 뉴런 & 멀티플렉서, 작은 다층 네트워크 & 허용 입력이 있을 때만 가지가 입력을 통과시킨다 & \ref{sec:synthesis}장 \\
\end{longtable}
\end{center}

표 전체보다 중요한 단서가 하나 있다. 이것들은 서로 다른 세포가 아니라 서로 다른 \emph{모드}다. 같은 뉴런이 느린 입력에서는 적분기가 되고, 억제가 $\tau$를 짧게 하면 동시성 검출기가 된다(\ref{sec:code}장). 깨어 있을 때는 페이스메이커이고, 잘 때는 조용한 수신기다. 어떤 소자가 될지는 막의 이온 채널과, 그 채널을 조정하는 브로드캐스트 채널이 정한다.

\section{이 책에 아직 나오지 않은 네 소자}

\subsection{공진기}
\label{sec:resonator}

어떤 뉴런에는 전압의 모든 변화에 맞서는 느린 전류가 있다. 전압이 오르면 끌어내리고 내리면 끌어올리지만, 지연이 있다. 전자 공학자에게 이런 전류는 인덕턴스처럼 행동하며, 막의 커패시턴스와 함께 공진 회로를 이룬다. 뉴런은 특정 주파수, 보통 몇 Hz에서 수십 Hz의 입력에 가장 강하게 반응하고, 그보다 느리거나 빠른 입력에는 약하게 반응한다~\cite{HutcheonYarom2000}. 세포 하나로 된 대역 통과 필터다. 잡음 섞인 입력에서 자기 주파수의 리듬, 예를 들어 \ref{sec:gaba}장의 국소 리듬을 골라낸다.

\subsection{무누설 적분기}
\label{sec:perfectint}

누설 적분기는 입력을 시간 $\tau$, 즉 수십 밀리초 동안만 기억한다. 그러나 눈이 움직이지 않고 머물려면 자기 위치를 몇 초 동안 기억해야 한다. “돌려라”라는 명령은 짧은 속도로 오지만, 눈은 새 위치에 붙들어 두어야 한다. 네트워크는 이 문제를 \ref{sec:glutcomputer}장의 기억과 같은 되먹임으로 푼다. 시상수 $\tau$인 뉴런 무리가 가중치 $w$로 자신을 흥분시키면 $\tau\,\dd r/\dd t=-r+w\,r+I$이고, 무리의 시상수는
\begin{empheq}[box=\keybox]{equation}
\tau_{\text{eff}} \;=\; \frac{\tau}{1-w}
\label{eq:perfectint}
\end{empheq}
이다. $\tau=0.1\,\text{s}$, $w=0.995$이면 $20\,\text{s}$가 된다. 하나하나는 10분의 1초 만에 잊는 뉴런들로 만든 거의 이상적인 적분기다~\cite{Seung1996}. 대가는 정밀한 조정이다. $w$가 1보다 조금만 커도 적분기는 포화로 치닫고, 조금만 작아도 빨리 잊는다. 그래서 이런 적분기에는 \ref{sec:motorlearning}장에서 설명한 것 같은 학습에 의한 끊임없는 조정이 필요하다.

\subsection{쌍안정 뉴런}
\label{sec:bistable}

\ref{sec:glutcomputer}장과 \ref{sec:gaba}장에서는 플립플롭을 뉴런 무리로 만들었다. 그러나 세포 하나로 된 플립플롭도 있다. 어떤 뉴런에는 한번 켜지면 스스로 흥분을 유지하는 전류가 있다. 짧은 입력이 뉴런을 오래가는 발화 상태, 즉 \emph{플래토}로 옮기고, 억제가 다시 휴지 상태로 되돌린다. 이것은 슈미트 트리거다. 뉴런에는 안정한 상태가 둘 있고 그 사이에 히스테리시스가 있다. 예를 들어 근육을 제어하는 \nt{ACET} 뉴런이 이렇게 행동하며, 이 성질은 브로드캐스트 채널이 켠다. 세로토닌이 뉴런에 닿으면 플래토가 나타난다~\cite{Hounsgaard1988}. 같은 세포가 \nt{SERO}가 있으면 래치이고, 없으면 평범한 적분기다.

\subsection{적분기와 공진기}

이론은 흥분성 세포를 두 부류로 나눈다~\cite{Izhikevich2007}. \emph{적분기}는 0에서 시작하는 전압 제어 발진기를 닮았다. 문턱값을 조금 넘으면 뉴런은 얼마든지 느리게 발화할 수 있고, 발화율은 입력에 따라 매끄럽게 커진다. 이런 뉴런은 들어온 것을 모두 더하며, 입력의 크기를 발화율로 잘 전달한다. \emph{공진기}는 최소 주파수가 있는 발진기를 닮았다. 느리게 발화하지 못하고, 문턱값 수준의 입력에서 곧바로 유한한 주파수로 스파이크를 내며, 자기 주파수의 입력을 선호한다. 앞의 것은 \ref{sec:code}장의 발화율 부호에 자연스러운 소자이고, 뒤의 것은 타이밍 부호에 자연스러운 소자다.

\begin{keyblock}
\begin{proposition}[뉴런 하나, 여러 소자]
\label{prop:eetypes}
전기적 기능으로 보면 뉴런은 누설 적분기와 무누설 적분기, 동시성 검출기, 고역 통과 필터와 대역 통과 필터, 전압---주파수 변환기, 정류기, 나눗셈기, 발진기, 지연선, 플립플롭으로 동작한다. 주어진 세포가 어떤 소자가 될지는 그 세포의 이온 채널과, 그 채널을 조정하는 브로드캐스트 채널이 정한다. 모드 자체는 측정되었다. 뇌가 이런 재구성을 계산에 얼마나 쓰는지는 대부분 가설이다.
\end{proposition}
\end{keyblock}

엔지니어에게 이것은 프로그래머블 로직 어레이와 비슷하다. 하드웨어는 하나이고 소자는 구성이 정한다. 다만 여기서는 구성을 펌웨어를 쓸 때가 아니라 동작 중에 바꾼다. \ref{sec:serocomputer}--\ref{sec:acet}장에서 다룬 느린 브로드캐스트 채널이 바꾼다.

\section{연습문제}

\begin{exercise}[\lvlA]
\label{ex:18:1}
\emph{무누설 적분기의 허용 오차.} $\tau=0.1$~s인 뉴런들이 식~\eqref{eq:perfectint}에 따라 눈의 위치를 $T=1$~s 동안 양쪽으로 $5\%$ 이내의 오차로 유지한다. (a)~허용되는 $w$의 범위를 구하라. (b)~$w$는 시냅스 $K$개의 가중치 합이고 각각 독립적인 $10\%$ 오차가 있다. $K$가 얼마일 때 합의 편차가 허용 범위에 들어가는가?
\end{exercise}

\begin{exercise}[\lvlB]
\label{ex:18:2}
\emph{회로로서의 공진기.} \ref{sec:resonator}절의 느린 전류를 변수 $W$로 나타내자: $C\,\dd V/\dd t=-g_LV-g_wW+I$, $\tau_w\,\dd W/\dd t=V-W$. 이것이 $R_w=1/g_w$, $L_w=\tau_w/g_w$인 병렬 가지를 가진 $RC$ 회로임을 보이고, $C/g_L=10\ms$, $\tau_w=100\ms$, $\alpha=g_w/g_L=4$에서 공진 주파수와 그곳에서 $|Z|$가 $|Z(0)|$보다 몇 배 큰지 구하라.
\end{exercise}

\begin{exercise}[\lvlB]
\label{ex:18:3}
\emph{적분기는 어떻게 발화를 시작하는가.} (a)~뉴런 $\tau\,\dd V/\dd t=-V+\mu$, $\tau=10\ms$, 문턱값 $\theta$, 0으로 재설정. $\mu/\theta-1$이 얼마일 때 발화율이 $1$~Hz인가? (b)~모델 $\tau\,\dd v/\dd t=v^2+\mu$(스파이크는 $v$가 $+\infty$로 가는 것, 재설정은 $-\infty$로)는 $\mu=0.01$에서 발화율이 얼마인가?
\end{exercise}

\begin{exercise}[\lvlB]
\label{ex:18:4}
\emph{모델링: 적분기와 공진기.} $g_L=1$, $\tau_m=10\ms$, $\tau_w=100\ms$, 문턱값 $1$, 재설정 $V\to0$($W$는 그대로)인 연습문제~\ref{ex:18:2}의 모델에 $\mu=1.5\sin2\pi ft$, $f=1$--$40$~Hz를 넣는다. $\alpha=0$과 $\alpha=4$에서 뉴런은 어떤 주파수 대역에서 스파이크를 내는가? 조건 $1.5\,|Z(f)|>1$과 비교하라.
\end{exercise}

\begin{exercise}[\lvlB]
\label{ex:18:5}
\emph{세포 하나로 된 래치.} 플래토 전류가 있는 뉴런: $\tau\,\dd r/\dd t=-r+F(wr+I)$, $F(x)=\min(1,\max(0,x))$, $\tau=10\ms$, $w=1.5$, 배경 $I=-0.2$. (a)~$I=0.3$ 펄스가 뉴런을 $r=0$에서 플래토로 옮기는 최소 길이는? (b)~긴 펄스 $I=-0.2-B$가 뉴런을 휴지 상태로 되돌리는 최소 진폭 $B$는?
\end{exercise}

\begin{exercise}[\lvlB]
\label{ex:18:6}
\emph{적분기의 자기 조정.} 시선을 고정할 때 입력은 $I=0$이고, 미끄러짐 센서가 표류 속도 $e=\dd r/\dd t$를 주며, $\dd w/\dd t=-\eta\,e\,r$이다. $w=0.95$, $\tau=0.1$~s, $\langle r^2\rangle=1$에서 고정 $60$~s 뒤 $|1-w|$가 연습문제~\ref{ex:18:1}의 허용 범위에 드는 $\eta$를 구하라. 눈을 돌리는 동안에도 학습하면 무엇이 망가지는가?
\end{exercise}

\begin{exercise}[\lvlC]
\label{ex:18:7}
\emph{소자를 왜 재구성하는가?} 브로드캐스트 채널이 연습문제~\ref{ex:18:2} 모델의 $\alpha$를 $0$과 $4$ 사이에서 바꾼다. 전류의 백색 잡음 속 약한 사인파에 대해, $11$~Hz와 $1$~Hz에서 공진기는 적분기보다 신호 대 잡음비가 몇 배 좋은가? 바로 모드 전환이 유리한 과제를 고안하라.
\end{exercise}

\chapter{가지돌기 개수에 따른 뉴런}
\chaptermark{가지돌기 개수}
\label{sec:dendrites}

\section{회로 매개변수로서의 입력 개수}

\ref{sec:neurontypes}장에서는 뉴런을 출력이 무엇을 내보내는지에 따라, \ref{sec:eetypes}장에서는 어떤 소자로 동작하는지에 따라 분류했다. 엔지니어라면 곧바로 알아챌 세 번째 특징도 있다. 뉴런에 \emph{입력이 몇 개 있는가}이다. 입력은 가지돌기가 모은다. 가지돌기는 다른 세포의 축삭이 달라붙는 가지다(\ref{sec:dofa}장, \ref{sec:weightstore}절). 어떤 뉴런에는 가지돌기가 아예 없고, 어떤 뉴런에는 한두 개가 있으며, 또 어떤 뉴런에는 입력 10만 개를 모으는 나무가 있다. 전자 공학자에게 이것은 입력 부하 능력, 즉 팬인(fan-in)이며, 거의 언제나 과제에 맞춰져 있다.

가지돌기의 개수와 크기가 왜 그렇게 중요한지는 간단한 추정으로 알 수 있다. 막은 단위 면적당 커패시턴스가 $c_m\approx1$~\textmu F/cm$^2$인 커패시터이며, 가지돌기를 포함한 뉴런의 면적 $A$가 클수록 전압을 $\Delta V$만큼 올리는 데 더 많은 전하가 필요하다. 누설을 무시하면 입력 전류 $I$가 이 일을 하는 데 걸리는 시간은 다음과 같다.
\begin{empheq}[box=\keybox]{equation}
t \;\approx\; \frac{c_m\,A\,\Delta V}{I} .
\label{eq:dendriteload}
\end{empheq}
짧은 가지돌기가 몇 개 달린 작은 뉴런은 면적이 약 $2\cdot10^{-5}$~cm$^2$이고 커패시턴스는 약 $20$~pF이다. 전류가 몇 나노암페어인 큰 시냅스 하나가 이 뉴런을 10분의 1밀리초도 안 되어 $20\mV$ 올린다. 큰 나무를 가진 뉴런은 면적이 열다섯 배쯤 크고 커패시턴스는 약 $300$~pF이다. 전류가 약 $0.1$~nA인 보통 시냅스라면 이 뉴런을 충전하는 데 $60\ms$가 걸릴 것이다. 누설 시상수보다 훨씬 길므로, 결국 끝내 충전하지 못한다. 이런 뉴런에는 동시에 오는 입력 수십 개가 필요하다. 여기서 수치는 크기 수준일 뿐이지만 결론은 수치에 좌우되지 않는다. \emph{입력이 적으면 빠르고 정확하며, 입력이 많으면 느리고 합으로 동작한다}.

\begin{figure}[t]
\centering
\begin{tikzpicture}[>=Stealth,
  soma/.style={circle,draw,very thick,fill=blue!6,minimum size=0.55cm,inner sep=0pt},
  ax/.style={->,very thick,red!70!black},
  den/.style={very thick,blue!60!black}]
% 0 dendrites: sensor on a line
\begin{scope}[xshift=0cm]
\draw[den,decorate,decoration={snake,amplitude=1pt,segment length=4pt}] (-0.9,0) -- (0,0);
\node[font=\scriptsize] at (-0.9,0.3) {센서};
\draw[ax] (0,0) -- (1.0,0);
\draw[thick] (0.1,0) -- (0.1,0.45);
\node[soma,minimum size=0.4cm] at (0.1,0.65) {};
\node[font=\scriptsize,align=center] at (0.05,-0.75) {$0$\\센서 달린 선로};
\end{scope}
% 1 dendrite
\begin{scope}[xshift=3.3cm]
\draw[den] (-0.9,0) -- (-0.28,0);
\node[soma] at (0,0) {};
\draw[ax] (0.28,0) -- (0.9,0);
\node[font=\scriptsize,align=center] at (0,-0.75) {$1$\\버퍼, 리피터};
\end{scope}
% 2 dendrites: one per ear
\begin{scope}[xshift=6.2cm]
\draw[den] (-0.28,0) -- (-0.9,0); \draw[den] (0.28,0) -- (0.9,0);
\node[soma] at (0,0) {};
\draw[ax] (0,-0.28) -- (0,-0.6);
\node[font=\scriptsize] at (-0.9,0.25) {왼쪽}; \node[font=\scriptsize] at (0.9,0.25) {오른쪽};
\node[font=\scriptsize,align=center] at (0,-1.0) {$2$\\시간 AND};
\end{scope}
% ~4 short dendrites: mixer
\begin{scope}[xshift=9.0cm]
\foreach \a in {45,135,225,315}{\draw[den] (\a:0.28) -- (\a:0.7); \fill[blue!60!black] (\a:0.72) circle (0.06);}
\node[soma] at (0,0) {};
\draw[ax] (0.28,0) -- (1.0,0);
\node[font=\scriptsize,align=center] at (0,-1.0) {$\sim4$\\믹서};
\end{scope}
% many: summing tree
\begin{scope}[xshift=12.0cm]
\foreach \a in {60,90,120}{\draw[den] (0,0.28) -- ++(\a:0.55) coordinate (b); \foreach \d in {-20,20}{\draw[den] (b) -- ++(\a+\d:0.35);}}
\foreach \a in {-120,-60}{\draw[den] (0,-0.28) -- ++(\a:0.45);}
\node[soma] at (0,0) {};
\draw[ax] (0.28,0) -- (1.0,0);
\node[font=\scriptsize,align=center] at (0,-1.0) {$10^4$--$10^5$\\가산기};
\end{scope}
\end{tikzpicture}
\caption{입력 개수에 따른 다섯 부류. 파란색은 가지돌기(입력), 빨간색은 축삭(출력)이다. 입력이 적을수록 뉴런은 개별 신호를 더 빠르고 정확하게 전달하고, 많을수록 더 잘 더하고 분류한다. 세포 모양을 그린 것이 아니라 도식이다.}
\label{fig:dendriteclasses}
\end{figure}

\section{가지돌기 0개: 선로 끝의 센서}

촉각, 통증, 근육 신장을 감지하는 감각 뉴런(\ref{sec:sensors}, \ref{sec:pain}, \ref{sec:muscles}장)은 세포체에 가지돌기가 하나도 없다. 축삭은 세포체에서 가지 하나로 나와 곧 둘로 갈라진다. 한 가지는 피부나 근육으로 가고 그 끝이 곧 센서다. 다른 가지는 척수로 간다. 스파이크는 세포체를 거치지 않고 센서에서 곧장 척수로 달린다. 엔지니어에게 이것은 먼 끝에 센서가 달린 통신 선로이고, 세포체는 전원 장치처럼 옆에 매달려 있다. 세포체는 선로에 전력을 대지만 전송에는 참여하지 않는다. 이득은 속도와 신뢰성이다. 경로 어디에도 신호를 합산하고 전달할지 다시 결정해야 하는 곳이 없다.

빛과 소리의 센서는 더 극단적인 경우다. 이들은 입력을 모으는 뉴런이 아니라 변환기 자체이며, 입력은 시냅스가 아니라 수용체 단백질이다(그림~\ref{fig:transducer}).

\section{가지돌기 1개: 버퍼}

가지돌기 하나와 축삭 하나를 가진 뉴런은 입력 선로 하나를 받아 그대로 넘긴다. 후각 뉴런이 그렇다. 하나뿐인 가지돌기가 코의 표면으로 나가며 한 종류의 수용체만 지닌다~\cite{Mombaerts2004}. 망막에서 빛 센서와 눈의 출력 뉴런 사이에 있는 전달 세포, 귀의 센서를 뇌와 잇는 뉴런도 그렇다. 엔지니어라면 이들을 버퍼나 리피터라고 부를 것이다. 이들은 계산하지 않고 신호 하나를 배달하며, 때로는 \ref{sec:sensors}장에서처럼 센서의 매끄러운 전압을 스파이크로 바꾸는 식으로 신호의 형태를 바꾼다.

\section{가지돌기 2개: 시간 AND}

소리의 방향을 정하는 동시성 검출기(그림~\ref{fig:itd})는 포유류에서 흔히 서로 반대쪽을 향한 가지돌기를 정확히 두 개 가진다. 한쪽에는 왼쪽 귀의 입력이, 다른 쪽에는 오른쪽 귀의 입력이 온다~\cite{Grothe2010}. 왜 입력을 떼어 놓을까? 식~\eqref{eq:shunt}을 떠올려 보자. 막의 한 부위에서 흥분성 채널이 많이 열리면 그곳의 전압은 평형 전위에 가까워지고, 다음 입력은 점점 적게 보탠다. 그래서 한쪽 귀에서 온 입력이 한 가지돌기에 많이 모이면 그 가지돌기는 포화되어, 그것만으로는 뉴런을 문턱값까지 끌어올리지 못한다. 반면 각 가지돌기에서 하나씩 온 입력은 거의 선형으로 더해진다. 가지돌기 두 개는 뉴런을 더 엄격한 AND로 만든다. 신호가 양쪽에서 함께 왔을 때만 발화하라. 모델에서 이런 분리가 실제로 동시성 검출을 개선한다는 것이 밝혀졌다~\cite{AgmonSnir1998}.

\section{짧은 가지돌기 몇 개: 믹서와 중계기}

\emph{믹서.} \ref{sec:motorlearning}장의 운동 학습 회로에서는 약 700억 개의 작은 \nt{GLUT} 뉴런이 입력을 아주 긴 벡터로 확장한다. 이들 각각은 짧은 가지돌기가 약 네 개뿐이고, 가지돌기마다 입력 하나의 말단을 감싸는 “발톱”으로 끝난다. 따라서 믹서 하나는 많은 입력 가운데 약 네 개만 보고, 자기 네 입력에 대한 작은 AND 소자로 동작한다. 입력 선로 $M$개에서 서로 다른 네 개 묶음을 $\binom{M}{4}$가지 고를 수 있다. $M=100$이면 거의 400만 가지다. 왜 이렇게 많이 필요할까? 입력이 $n$개인 문턱 소자는 무작위 패턴을 최대
\begin{empheq}[box=\keybox]{equation}
P_{\max} \;\approx\; 2n
\label{eq:cover}
\end{empheq}
개까지 두 부류로 올바르게 나눌 수 있다~\cite{Cover1965}. 같은 회로의 출력 뉴런은 믹서에서 약 $10^5$개의 입력을 받으므로, \eqref{eq:cover}에 따라 상한은 서로 다른 대응 약 $2\cdot10^5$개다. 이것은 이상적인 소자의 한계다. 흥분성 시냅스처럼 모든 가중치가 양수이고 잡음에 대한 여유가 필요하면, 같은 뉴런에 대한 추정은 약 $4\cdot10^4$개의 대응이다~\cite{Brunel2004}. 입력이 적은 믹서는 바로 큰 가산기가 서로 다르고 거의 독립인 입력을 많이 갖게 하려고 필요하다~\cite{Marr1969,Albus1971}.

\emph{중계기.} 귀에서 방향 검출기로 가는 경로에는 짧은 가지돌기가 몇 개뿐인 뉴런이 있다. 이들은 거대한 시냅스를 몇 개만 받고, 어떤 뉴런은 사실상 하나만 받는다. \eqref{eq:dendriteload}에 따르면 이것이 딱 필요한 구성이다. 커패시턴스는 작고 전류는 크므로, 뉴런은 입력 스파이크마다 1밀리초도 안 되어 발화하며 시간 정밀도를 거의 잃지 않는다. 이런 중계기 가운데 하나는 \nt{GLUT}를 받아 \nt{GLYC}를 낸다. 수십 마이크로초 정밀도의 인버터로, 방향 검출기에 빠른 억제를 공급한다~\cite{BorstSoria2012}.

\section{입력 수천 개: 가산기와 분류기}

반대쪽 끝에는 입력이 1만 개쯤인 피질의 \nt{GLUT} 뉴런과, 입력이 10만 개인 운동 학습 회로의 출력 \nt{GABA} 뉴런이 있다. 이런 뉴런은 \eqref{eq:dendriteload}에 따라 느리고 개별 입력에 반응할 수 없다. 이들은 더한다. \ref{sec:glutcomputer}장의 적분 뉴런이며, 분류기를 만드는 바로 그 가산기다. 큰 나무는 새로운 것도 보탠다. 나무의 가지는 각자 문턱값을 가진 별도의 하위 회로로 동작하므로, 뉴런 하나가 작은 다층 네트워크처럼 행동한다~\cite{Beniaguev2021,Gidon2020}.

\section{출력이기도 한 가지돌기}

축삭이 아예 없고 가지돌기가 입력이자 출력인 뉴런도 있다. 이 뉴런은 신호를 받고 스스로 신경전달물질을 내보낸다. 가장 많이 연구된 예는 운동 방향 판별에 참여하는 망막의 \nt{GABA} 뉴런이다(\ref{sec:gaba}장, 그림~\ref{fig:direction}). 이런 뉴런에서는 가지 하나하나가 스스로 세포체에서 자기 끝 쪽으로 가는 움직임에 반응하고, 그 끝에서 억제를 낸다~\cite{Euler2002}. 뉴런 하나가 가지마다 하나씩, 독립된 방향 검출기 여러 개다. 엔지니어에게 이것은 소자 하나가 아니라 패키지 하나에 채널이 여럿 든 칩이다.

\section{요약}

\begin{table}[t]
\centering
\footnotesize
\setlength{\tabcolsep}{3pt}
\renewcommand{\arraystretch}{1.15}
\begin{tabular}{>{\raggedright}p{1.9cm}>{\raggedright}p{3.4cm}>{\raggedright}p{2.6cm}>{\raggedright}p{1.8cm}>{\raggedright\arraybackslash}p{3.8cm}}
\toprule
가지돌기 & 예 & 소자 & 입력 & 알려진 것 \\
\midrule
$0$ & 촉각, 통증, 신장 센서 & 끝에 센서가 달린 선로 & --- & 확립됨 \\
$1$ & 후각 뉴런, 망막의 전달 세포, 귀의 뉴런 & 버퍼, 리피터 & $1$--소수 & 확립됨 \\
$2$ & 소리 방향 검출기 & 시간 AND & 두 다발 & 구조는 확립됨; 입력 분리의 이득은 모델에서 입증 \\
$\sim4$ & 운동 학습 회로의 믹서 & 작은 AND 소자 & $\sim4$ & 확립됨; 입력 확장의 역할은 근거가 탄탄한 가설 \\
소수, 짧음 & 귀에서 오는 경로의 중계기 & 중계기, 인버터 & $1$--몇 개의 거대 시냅스 & 확립됨 \\
수천 & 피질의 \nt{GLUT} 뉴런, 운동 학습의 출력 뉴런 & 가산기, 분류기 & $10^4$--$10^5$ & 확립됨; 하위 회로로서의 가지는 개별 사례에서 측정 \\
출력 가지돌기 & 망막의 방향 \nt{GABA} 뉴런 & 다채널 칩 & 많음 & 측정됨 \\
\bottomrule
\end{tabular}
\caption{가지돌기 개수에 따른 뉴런.}
\label{tab:dendrites}
\end{table}

\begin{keyblock}
\begin{proposition}[입력 개수는 과제를 따른다]
\label{prop:dendrites}
센서, 중계기, 동시성 검출기처럼 개별 신호의 속도와 정밀도가 필요한 곳에서는 뉴런의 가지돌기가 적고 입력이 적으며 시냅스가 크다. 작은 커패시턴스~\eqref{eq:dendriteload}는 빨리 충전된다. 더하고 분류해야 하는 곳에서는 뉴런이 입력 수천 개를 가진 큰 나무를 갖는다. 부류의 구조는 측정되었다. 계산적 설명은 근거가 탄탄하지만 많은 부류에서는 아직 가설이다.
\end{proposition}
\end{keyblock}

표~\ref{tab:dendrites}는 앞의 두 분류, 즉 신경전달물질에 따른 분류(표~\ref{tab:types})와 소자에 따른 분류(표~\ref{tab:eetypes})를 보완한다. 세 분류 모두 같은 소자를 다른 쪽에서 본다. 무엇을 내보내는지, 무엇을 계산하는지, 얼마나 많이 듣는지다.

\section{연습문제}

\begin{exercise}[\lvlA]
\label{ex:19:1}
\emph{큰 뉴런에는 입력이 몇 개 필요한가.} $C=300$~pF, $\tau_m=20\ms$인 뉴런이 각각 $0.1$~nA의 전류원인 시냅스를 받고, 문턱값은 휴지 상태보다 $20\mV$ 높다. (a)~문턱값에 아예 도달하려면, $10\ms$ 안에, $5\ms$ 안에 도달하려면 시냅스 몇 개가 동시에 작동해야 하는가? (b)~$C=20$~pF인 뉴런은 $4$~nA 시냅스 하나로 얼마 만에 문턱값에 도달하는가?
\end{exercise}

\begin{exercise}[\lvlA]
\label{ex:19:2}
\emph{왜 네 개인가.} 믹서는 $M=100$개 가운데 $k$개 선로에 대한 AND이고, 패턴에서 선로의 절반이 활성이며, 믹서는 $7\cdot10^{10}$개다. (a)~$k=2$, $4$, $40$일 때 패턴에 반응하는 믹서는 몇 개인가? (b)~믹서는 직접 선로 $M=100$개에 비해 입력 $10^5$개인 출력 뉴런의 대응 수~\eqref{eq:cover}를 몇 배 늘리는가?
\end{exercise}

\begin{exercise}[\lvlB]
\label{ex:19:3}
\emph{$2n$은 어디서 오는가.} 원점을 지나는 초평면은 $n$차원 공간에서 일반 위치에 있는 점 $P$개를 $C(P,n)=2\sum_{k=0}^{n-1}\binom{P-1}{k}$가지 방식으로 두 부류로 나눈다. (a)~$P=2n$이면 $2^P$개 분할 가운데 정확히 절반이 분리 가능함을 증명하라. (b)~$n=100$, $P=180$과 $220$에서 분리 가능한 비율은 얼마인가?
\end{exercise}

\begin{exercise}[\lvlB]
\label{ex:19:4}
\emph{엄격한 AND로서의 가지돌기 두 개.} 가지돌기는 누설 $g_L$인 구획이고, 입력 하나가 그 위에 전위 $E_E$인 컨덕턴스 $g_0=0.5\,g_L$을 연다. 세포체는 $\kappa(V_1+V_2)$를 보고, 문턱값은 $\theta=1.1\,\kappa E_E$이다. 왜 한 가지돌기에 입력이 아무리 많아도 뉴런이 발화하지 않는가? 각 가지돌기에 입력이 몇 개 필요한가?
\end{exercise}

\begin{exercise}[\lvlB]
\label{ex:19:5}
\emph{중계기 설계.} 스파이크 지터는 $\sigma_t=\sigma_V/\dot V$이고, $\dot V$는 문턱값에서의 전압 기울기다. 잡음 $\sigma_V=1\mV$에서 $\sigma_t\le20$~\textmu s가 필요하다. 누설은 무시하라. $C=300$~pF인 뉴런에는 $0.1$~nA의 동기 시냅스가 몇 개 필요한가? $C=20$~pF이고 $4$~nA 시냅스가 하나인 뉴런의 지터는?
\end{exercise}

\begin{exercise}[\lvlB]
\label{ex:19:6}
\emph{모델링: 긴 가지돌기의 충전.} 끝이 닫힌 길이 $L=\ell/\lambda$의 수동 케이블 $\tau_m\partial_tV=\lambda^2\partial_x^2V-V+i/g$를 시뮬레이션하라($40$개 구획, 오일러법). 먼 끝에 짧은 전류 펄스를 넣는다. $L=0.2$, $1$, $2$에서 가까운 끝은 언제, 같은 전하를 가지돌기 전체에 고르게 퍼뜨렸을 때 전압의 몇 분의 몇까지 오르는가?
\end{exercise}

\begin{exercise}[\lvlC]
\label{ex:19:7}
\emph{탐구: 최적의 입력 개수.} 커패시턴스는 $C=C_0+nc_s$이고, 전류 $I_s$인 입력 $m$개가 동시에 작동하며, 뉴런은 대응 $P\approx2n$개를 기억한다~\eqref{eq:cover}. $m$이 일정할 때와 비율 $m=\varphi n$이 일정할 때 응답 시간은 $P$에 어떻게 의존하는가? 비용 기준을 제안하고 중계기와 분류기의 최적 $n$을 구하라.
\end{exercise}

\chapter{잠자는 동안 뉴런은 무엇을 하는가}
\chaptermark{수면}
\label{sec:sleep}

\section{수면은 꺼짐이 아니라 다른 모드다}

꺼진 컴퓨터를 본 엔지니어는 아무것도 하지 않는다고 말할 것이다. 잠든 뇌는 그렇지 않다. 잠자는 동안에도 뉴런은 조용하지 않다. 깊은 잠에서 피질은 스파이크를 내고, 렘수면에서는 낮과 거의 같은 만큼 낸다. 바뀌는 것은 뉴런이 \emph{일하는지}가 아니라 \emph{어떻게} 일하는지다. 수면은 기계의 모드 묶음이며, 이 모드를 바꾸는 것은 낮과 같은 브로드캐스트 채널이다.

모드마다 “레지스터 상태”를 적을 수 있다(표~\ref{tab:sleepmodes}). 낮에는 \nt{ACET}, \nt{NORA}, \nt{SERO}가 일하고, 센서가 연결되어 있으며, 근육이 명령을 따른다. 느리고 깊은 서파 수면에서는 세 채널이 모두 잠잠해지고 감각 기관의 입력이 약해진다~\cite{Hasselmo1999}. 피질의 아세틸콜린 방출이 떨어지고, \nt{NORA} 뉴런과 \nt{SERO} 뉴런은 낮보다 드물게 스파이크를 낸다~\cite{Marrosu1995,AstonJonesBloom1981,McGintyHarper1976}. 렘수면의 그림은 기묘하다. \nt{ACET}은 다시 낮 수준으로 오르지만~\cite{Marrosu1995}, \nt{NORA}와 \nt{SERO}는 거의 완전히 입을 다문다~\cite{AstonJonesBloom1981,McGintyHarper1976,JacobsAzmitia1992}. 렘수면에서 몸의 근육은 꺼져 있다. 근육을 제어하는 \nt{ACET} 뉴런이 \nt{GABA}와 \nt{GLYC}를 통해 강하게 억제된다~\cite{BrooksPeever2012}. \ref{sec:gaba}장의 금지와 같은 것이지만, 출력 전체에 한꺼번에 걸린다.

\begin{table}[t]
\centering
\footnotesize
\setlength{\tabcolsep}{4pt}
\renewcommand{\arraystretch}{1.15}
\begin{tabular}{>{\raggedright}p{3.1cm}>{\raggedright}p{2.9cm}>{\raggedright}p{3.2cm}>{\raggedright\arraybackslash}p{3.4cm}}
\toprule
& 각성 & 서파 수면 & 렘수면 \\
\midrule
\nt{ACET} (쓰기, \ref{sec:acet}장) & 높음 & 낮음 & 높음 \\
\nt{NORA} (이득, \ref{sec:nora}장) & 중간, 폭발적 발화 & 낮음 & 거의 조용 \\
\nt{SERO} (\ref{sec:serocomputer}장) & 중간 & 낮음 & 거의 조용 \\
센서 입력 & 연결됨 & 약해짐 & 약해짐 \\
근육 출력 & 연결됨 & 약해짐 & 억제로 꺼짐 \\
피질 활동 & 고름 & 느린 시소: 활동---침묵 & 낮처럼 고름 \\
\bottomrule
\end{tabular}
\caption{세 모드에서 기계의 레지스터 상태. 모든 행은 측정되었다. \nt{ACET}, \nt{NORA}, \nt{SERO}의 수준은 수면 단계별 직접 기록으로 측정되었다~\cite{Marrosu1995,AstonJonesBloom1981,McGintyHarper1976}.}
\label{tab:sleepmodes}
\end{table}

\section{언제 잘까: 히스테리시스가 있는 릴레이}

기계가 언제 수면으로 전환할지는 무엇이 정할까? 두 부분으로 된 간단한 회로가 잘 들어맞는다~\cite{Borbely1982}. 첫 부분은 \emph{수면 압력} $S$다. 깨어 있는 동안 쌓이고 자는 동안 풀린다. 낮에 충전되고 밤에 방전되는 커패시터다.
\begin{empheq}[box=\keybox]{equation}
\frac{\dd S}{\dd t} \;=\; \frac{1-S}{\tau_r}\ \ \text{(각성)},\qquad
\frac{\dd S}{\dd t} \;=\; -\frac{S}{\tau_d}\ \ \text{(수면)} .
\label{eq:sleeppressure}
\end{empheq}
둘째 부분은 문턱값 두 개다. 기계는 $S$가 위 문턱값 $H(t)$에 이르면 잠들고, 아래 문턱값 $L(t)$까지 내려가면 깬다. 두 문턱값은 \ref{sec:chemoutput}장의 일주기 리듬에 맞춰 부드럽게 오르내린다. 그 결과는 히스테리시스가 있는 릴레이, 즉 \ref{sec:bistable}절의 슈미트 트리거이며, 커패시터와 함께 일주기 리듬~\eqref{eq:pll}에 위상이 맞춰지는 이완 발진기가 된다.

\begin{figure}[t]
\centering
\begin{tikzpicture}[>=Stealth,x=0.16cm,y=4.2cm]
\foreach \a/\b in {0/4.3,21.2/29.3,45.4/53.4,69.5/72}{\fill[blue!8] (\a,0) rectangle (\b,1);}
\draw[->,gray!70!black] (0,0) -- (75,0) node[right,font=\small,black] {$t$, h};
\draw[->,gray!70!black] (0,0) -- (0,1.08) node[left,font=\small,black] {$S$};
\foreach \t in {0,24,48,72}{\draw[gray!60!black] (\t,-0.02) -- (\t,0.02); \node[below,font=\scriptsize] at (\t,0) {\t};}
\draw[thick,densely dashed,red!70!black] plot file {figures/sleep_H.dat};
\draw[thick,densely dashed,green!45!black] plot file {figures/sleep_L.dat};
\draw[very thick,blue!60!black] plot file {figures/sleep_S.dat};
\node[font=\scriptsize,red!70!black,anchor=west] at (73,0.72) {$H$: 잠들기};
\node[font=\scriptsize,green!45!black,anchor=west] at (73,0.24) {$L$: 깨기};
\node[font=\scriptsize,blue!60!black] at (25.25,0.93) {수면};
\node[font=\scriptsize,blue!60!black] at (49.4,0.93) {수면};
\end{tikzpicture}
\caption{$\tau_r=18.2$~h, $\tau_d=4.2$~h에서의 수면 압력 $S$~\eqref{eq:sleeppressure}. 낮에는 $S$가 위 문턱값 $H$까지 충전되고, 밤에는(색칠한 부분) 아래 문턱값 $L$까지 방전된다. 두 문턱값은 일주기 리듬과 함께 오르내린다. 기계는 스스로 약 $16$시간 각성, $8$시간 수면의 체제에 이른다.}
\label{fig:twoprocess}
\end{figure}

보통 값 $\tau_r=18.2$~h, $\tau_d=4.2$~h를 넣은 우리 모델에서 기계는 스스로 $16.1$시간 각성과 $8.0$시간 수면에 이른다(그림~\ref{fig:twoprocess}). 모델은 이상해 보이는 현상도 설명한다. 40시간 동안 자지 않으면 압력은 $0.90$까지 오르지만, 모델에서 회복 수면은 평소보다 하루 더 긴 것이 아니라 $8.8$시간뿐이다. 방전은 지수적이어서 높은 충전은 빨리 빠진다. “빚”은 수면의 길이가 아니라 깊이로 갚는다.

물리적으로 $S$ 역할을 하는 것은 무엇일까? 유력한 후보는 부록~\ref{sec:othermediators}의 “부하 계수기” 아데노신이다. 아데노신 농도는 깨어 있는 동안 오르고 자는 동안 떨어진다~\cite{PorkkaHeiskanen1997,Dunwiddie2001}. 수면 압력이 바로 아데노신이고 다른 무엇도 아니라는 것은 가설이다.

\section{서파 수면: 이완 발진기가 된 네트워크 전체}

깊은 잠에서 피질 뉴런은 무리 전체가 번갈아 일한다. 1초에 한 번보다 드물게, 수백 밀리초 동안 함께 켜지고(\emph{활동 상태}) 함께 조용해진다(\emph{침묵 상태}). 이 시소의 주파수는 $1$~Hz 아래이고, 마취 상태에서는 대개 $0.3$--$0.4$~Hz다~\cite{Steriade1993}. 이 느린 시소는 이미 가진 부품으로 만들 수 있다. 자신에게 강한 되먹임을 거는 \nt{GLUT} 뉴런 무리, 즉 안정한 상태가 둘인 \ref{sec:glutcomputer}장의 기억을 가져와서, \ref{sec:muscles}장의 리듬 발생기에서처럼 느린 피로를 더하자.
\begin{empheq}[box=\keybox]{equation}
\tau\,\frac{\dd r}{\dd t} = -r + F\bigl(w\,r + I - a\bigr),
\qquad
\tau_a\,\frac{\dd a}{\dd t} = -a + b\,r .
\label{eq:updown}
\end{empheq}
무리는 켜지고, 일하고, 지친다. 피로 $a$가 위 상태를 잠식하면 무리는 침묵으로 떨어진다. 침묵하는 동안 피로가 풀리면 아래 상태가 사라지고 무리는 다시 켜진다. 수면 압력과 같은 이완 발진기이지만 약 8만 배 빠르다.

\begin{figure}[t]
\centering
\begin{tikzpicture}[>=Stealth,x=2.9cm,y=1.0cm]
\begin{scope}[yshift=1.9cm]
\draw[->,gray!70!black] (0,0) -- (4.1,0);
\draw[->,gray!70!black] (0,0) -- (0,1.25) node[left,font=\small,black] {$r$};
\draw[thick,blue!60!black] plot file {figures/updown_sleep.dat};
\node[font=\scriptsize,anchor=west] at (4.15,0.5) {수면: $b=5$};
\end{scope}
\draw[->,gray!70!black] (0,0) -- (4.1,0) node[right,font=\small,black] {$t$, s};
\draw[->,gray!70!black] (0,0) -- (0,1.25) node[left,font=\small,black] {$r$};
\draw[thick,red!70!black] plot file {figures/updown_wake.dat};
\node[font=\scriptsize,anchor=west] at (4.15,0.5) {낮: $b=1$};
\foreach \t in {0,1,2,3,4}{\draw[gray!60!black] (\t,-0.05) -- (\t,0.05); \node[below,font=\scriptsize] at (\t,0) {\t};}
\end{tikzpicture}
\caption{두 모드에 있는 같은 무리~\eqref{eq:updown}: $w=5$, $I=2$, $\theta=2.5$, $\tau=10\ms$, $\tau_a=0.3$~s. 피로가 강하면($b=5$, 서파 수면처럼) 무리는 주기 $1.1$~s로 활동과 침묵 사이를 오간다(모델의 값; 피질에서 주기는 1초보다 길고, 마취 상태에서는 약 $3$~s). 피로를 약하게 하면($b=1$, 아세틸콜린이 이렇게 작용한다) 같은 무리가 고르게 일한다.}
\label{fig:updown}
\end{figure}

그러면 무엇이 네트워크를 시소에서 낮의 고른 작동으로 옮길까? 아세틸콜린은 뉴런에서 피로를 만드는 느린 전류를 억누른다~\cite{McCormick1992}. 우리 모델에서 피로가 강한 $b=5$일 때 무리는 주기 $1.1$~s로 오간다. 측정된 시소보다 빠르지만 크기 수준은 같다. 정수 개의 주기에 걸쳐 평균하면 시간의 약 $35\%$ 동안 일한다. 피로가 약한 $b=1$에서는 같은 무리가 고르게 일한다(그림~\ref{fig:updown}). 레지스터 하나, 즉 \nt{ACET} 수준이 발진기를 증폭기로 바꾼다. 수면 중에는 느린 시소 위에 초당 $7$--$15$번 진동하는 더 빠른 리듬의 폭발도 얹힌다. 이것은 피질과 중간 중계 핵 사이의 \nt{GLUT}--\nt{GABA} 루프가 만들며~\cite{SteriadeMcCormickSejnowski1993}, \ref{sec:gaba}장의 리듬과 친척이다.

\section{재생: 하루를 빨리 감기}

\ref{sec:acet}장에서 우리는 서파 수면에서 낮에 함께 일한 뉴런들이 같은 순서로 다시 함께 발화하는 것을 보았다. 공학적인 세부 하나를 더하자. 재생은 \emph{빨리 감기}로 진행된다. 낮에 몇 초 걸린 순서가 잠자는 동안 10분의 몇 초 만에 재생된다~\cite{Nadasdy1999}. 해마에서는 약 20배~\cite{LeeWilson2002}, 피질에서는 6--7배 빠르다~\cite{Euston2007}. 게다가 아무 때나 재생되는 것이 아니다. 한 영역의 짧은 재생 폭발은 피질의 활동-침묵 시소와 빠른 리듬 폭발에 맞춰진다. 이 맞물림을 인위적으로 강화하면 기억이 좋아진다~\cite{Maingret2016}. 이는 이미 우연의 일치가 아니라 인과 관계다.

기계는 왜 하루를 빨리 감을까? 엔지니어는 \emph{파국적 망각}이라는 어려움을 안다. 새것을 하나씩 차례로 배우는 네트워크는 옛것을 망가뜨린다. 해결책은 \emph{섞기}다. 새것을 옛것과 섞어, 조금씩, 여러 번 배우는 것이다. 두 학습 시스템 가설~\cite{McClelland1995}에 따르면 빠른 기억은 하루를 통째로 기록하고, 잠자는 동안 그것을 느린 피질에 조금씩, 섞어서, 여러 번 재생해 준다. 이것은 \ref{sec:acet}장의 “빠른 버퍼 --- 느린 저장소” 쌍과 같고, \ref{sec:motorlearning}장의 빠른 학습자와 느린 학습자 쌍과도 같다. 재생과 그 맞물림은 측정되었다. 그 목적이 느린 기억의 섞어 배우기라는 것은 근거가 탄탄한 가설이다.

\section{가중치의 재정규화}

낮에는 \ref{sec:dofa}장의 헤브 규칙이 주로 연결을 강화한다. 강화만 하면 가중치가 커지고, 에너지 소비도 함께 커지며 민감도는 떨어진다. 모든 입력이 강한 뉴런은 입력을 구별하지 못하게 된다. \emph{시냅스 항상성} 가설~\cite{TononiCirelli2014}에 따르면 수면은 무엇보다 가중치를 작동 수준으로 되돌리는 데 필요하다. 모든 연결이 같은 비율로 약해지므로 연결 사이의 \emph{비}, 즉 배운 것은 보존된다. 엔지니어에게 이것은 규칙~\eqref{eq:oja}과 같은 정규화이지만, 동작 중이 아니라 따로 떼어 둔 시간에 한다.

직접 측정은 이와 모순되지 않는다. 생쥐에서 피질 시냅스의 접촉 면적은 깨어 있은 뒤보다 잔 뒤에 약 $18\%$ 작고, 감소는 크기에 비례하며, 가장 크고 안정된 접촉은 거의 바뀌지 않는다~\cite{deVivo2017}. 가중치의 스케일링은 측정되었다. 그것이 수면의 주된 과제라는 것은 가설이다.

\section{렘수면: 입력과 출력에서 분리된 기계}

렘수면에서 피질은 낮과 거의 같이 일하지만, 감각 기관의 입력은 약해지고 근육 출력은 억제로 꺼진다. 엔지니어에게 익숙한 모드다. \emph{벤치 시험}이다. 회로는 최대 출력으로 동작하지만, 입력은 내부 발생기에 물려 있고 출력은 아무것도 망가뜨리지 않도록 액추에이터에서 분리되어 있다. 한 가설에 따르면 꿈은 바로 낮 동안 모은 세계 모델을 이렇게 내부에서 돌려 보는 것이다. 이때 네트워크가 정확히 무엇을 하는지, 학습을 하는지는 알려지지 않았다. 여기서 측정된 것은 표~\ref{tab:sleepmodes}에 적은 것뿐이다. 어느 채널이 켜져 있고 어느 채널이 조용한지, 그리고 출력이 막혀 있다는 것이다.

\section{유지 보수와 국소 수면}

오랫동안 잠자는 동안 세포 사이 틈이 넓어져 뇌가 대사 산물을 더 빨리 씻어 낸다고 여겨졌다~\cite{Xie2013}. 다른 방법으로 세척을 잰 최근 실험은 반대 결과를 얻었다. 잠자는 동안 세척이 더 느리다~\cite{Miao2024}. 이 문제는 지금 열려 있으며, 우리도 열어 둔다.

더 확실한 사실이 있다. 수면은 기계 전체가 아니라 국소 네트워크의 성질이다. 동물을 오래 재우지 않으면, 동물이 깨어 움직이는 동안에도 피질의 일부 영역이 서파 수면에서처럼 잠깐 침묵으로 빠지고, 그 순간 동물은 더 자주 실수한다~\cite{Vyazovskiy2011}. 네트워크 하나하나가 스스로 지치고 스스로 전환한다. 전체 모드는 브로드캐스트 채널이 모든 네트워크를 한꺼번에 전환할 때 생긴다.

\begin{keyblock}
\begin{proposition}[모드 묶음으로서의 수면]
\label{prop:sleep}
잠자는 동안 뉴런은 꺼지지 않는다. 브로드캐스트 채널이 기계를 다른 모드로 옮긴다. 언제 잘지는 일주기 리듬에 맞춰진 히스테리시스 릴레이~\eqref{eq:sleeppressure}가 정한다. 서파 수면에서 네트워크는 이완 발진기~\eqref{eq:updown}가 되어 하루를 빨리 감아 재생한다. 렘수면에서는 입력과 출력에서 분리된 채 일한다. 모드, 재생, 가중치 스케일링은 측정되었다. 그 목적, 즉 섞어 배우기와 재정규화는 근거가 탄탄한 가설이다.
\end{proposition}
\end{keyblock}

\section{요약}

\begin{table}[t]
\centering
\footnotesize
\setlength{\tabcolsep}{4pt}
\renewcommand{\arraystretch}{1.15}
\begin{tabular}{>{\raggedright}p{3.2cm}>{\raggedright}p{4.4cm}>{\raggedright\arraybackslash}p{5.2cm}}
\toprule
일어나는 일 & 방법 & 알려진 것 \\
\midrule
모드 전환 & 레지스터 \nt{ACET}, \nt{NORA}, \nt{SERO}(표~\ref{tab:sleepmodes}) & 측정됨 \\
언제 잘까 & 커패시터 $S$와 히스테리시스 릴레이~\eqref{eq:sleeppressure} & 모델이 실험을 잘 기술함; $S$로서의 아데노신은 가설 \\
느린 시소 & 되먹임과 피로가 있는 무리~\eqref{eq:updown} & 시소는 측정됨; 모델은 가장 단순한 회로 \\
하루 빨리 감기 & $6$--$20$배 빠르고 시소에 맞물린 재생 & 측정됨; 맞물림을 강화하면 기억이 좋아짐 \\
재생은 왜 & 느린 기억의 섞어 배우기 & 근거가 탄탄한 가설 \\
가중치 재정규화 & 수면 중 연결의 스케일링 & 접촉 감소는 측정됨; 수면의 주된 역할이라는 것은 가설 \\
렘수면 & 입력과 출력을 분리한 채 동작 & 모드는 측정됨; 목적은 알려지지 않음 \\
세척 & 세포 사이 틈의 확장 & 상반된 결과 \\
국소 수면 & 개별 영역이 침묵으로 빠진다 & 측정됨 \\
\bottomrule
\end{tabular}
\caption{잠자는 동안 뉴런이 하는 일.}
\label{tab:sleep}
\end{table}

표~\ref{tab:sleep}이 이 장을 요약한다. 수면은 기계 작동의 휴지가 아니라 동작 중의 유지 보수로 드러났다. 같은 브로드캐스트 채널에 의한 모드 전환, 걸음 리듬과 같은 부품으로 만든 발진기, 느린 기억을 위한 하루 빨리 감기, 가중치 재정규화다. 낮에 기계는 기록하고, 밤에는 기록한 것을 다시 쓰고 정리한다.

\section{연습문제}

\begin{exercise}[\lvlA]
\label{ex:20:1}
\emph{일주기 리듬이 없는 릴레이.} 문턱값이 일정하다고 하자: $H=0.67$, $L=0.17$(그림~\ref{fig:twoprocess} 모델 문턱값의 평균). $\tau_r=18.2$~h, $\tau_d=4.2$~h에서 \eqref{eq:sleeppressure}로부터 각성과 수면의 길이, 발진기의 주기를 유도하라. 문턱값이 오르내리는 모델은 왜 $16.1+8.0=24$~h에 이르는가? 이 효과는 \ref{sec:chemoutput}장의 어느 소자에 해당하는가?
\end{exercise}

\begin{exercise}[\lvlA]
\label{ex:20:2}
\emph{수면 빚.} 압력 $S_0$에서 시작한 수면은 $t_s=\tau_d\ln(S_0/L)$ 동안 이어진다. $\tau_d=4.2$~h, $L$은 깰 때의 아래 문턱값이다. (a)~$40$~h 동안 자지 않은 뒤 $S_0=0.90$이고 수면은 $8.8$~h였다. $L$은 얼마였는가? 평소의 $L\approx0.092$이면 수면은 얼마나 이어졌겠는가? (b)~하룻밤 사이 시냅스 접촉 면적이 $18\%$ 줄어든다. 낮 동안 가중치는 얼마나 커져야 하는가?
\end{exercise}

\begin{exercise}[\lvlB]
\label{ex:20:3}
\emph{문턱값 설계.} $\tau_r=18.2$~h, $\tau_d=4.2$~h인 릴레이~\eqref{eq:sleeppressure}가 정확히 $16$~h 각성과 $8$~h 수면을 내도록 일정한 문턱값 $H$와 $L$을 고르라. 어느 날 밤 $4$시간 만에 깨웠다면, 이 기계는 문턱값이 오르내리는 기계보다 무엇이 나쁜가?
\end{exercise}

\begin{exercise}[\lvlB]
\label{ex:20:4}
\emph{느린 시소의 주기.} 모델~\eqref{eq:updown}에서 $F(x)=1/\bigl(1+e^{-(x-\theta)/k}\bigr)$, $w=5$, $I=2$, $\theta=2.5$, $k=0.25$, $b=5$, $\tau\ll\tau_a=0.3$~s. (a)~곡선 $r=F(wr+I-a)$에서 활동과 침묵이 사라지는 전환점 $a_{\text{hi}}$와 $a_{\text{lo}}$를 구하라. (b)~활동에서 $r\approx1$, 침묵에서 $r\approx0$으로 보고 주기와 활동 비율을 구하라.
\end{exercise}

\begin{exercise}[\lvlB]
\label{ex:20:5}
\emph{시소는 언제 꺼지는가.} 연습문제~\ref{ex:20:4}의 모델에서 고정점은 곡선 $a=wr+I-F^{-1}(r)$과 직선 $a=br$의 교점이다. 교점이 전환점 사이의 가운데 가지에 있는 동안 진동이 일어난다. $b$가 얼마일 때 그런가? \nt{ACET}은 $b$를 바꿔 어떻게 발진기를 증폭기로 바꾸는가?
\end{exercise}

\begin{exercise}[\lvlB]
\label{ex:20:6}
\emph{모델링: 빚인가, 위상인가.} \texttt{sleep\_models.py}에서 $48$~h 뒤의 첫 깸을 잡고, 기계를 $D=24$, $32$, $40$, $48$, $64$~h 동안 깨어 있게 하라(\texttt{stay\_awake}, \texttt{days=8}). 각 $D$에서 회복 수면은 얼마나 이어지는가? 그 길이는 무엇이 정하는가?
\end{exercise}

\begin{exercise}[\lvlB]
\label{ex:20:7}
\emph{모델링: 망각과 섞기.} 선형 뉴런 $y=\mathbf w\cdot\mathbf x$, $n=40$이 규칙 $\mathbf w\leftarrow\mathbf w-0.5\,(y-y^*)\,\mathbf x$로 배운다. 과제 A와 B는 각각 패턴 $10$개, $x_j\sim\mathcal N(0,1/n)$, $y^*=\pm1$. A를 $200$ 에포크, 이어서 B를 $200$ 에포크 학습한 뒤 A에 대한 평균제곱오차는? A와 B를 섞어 $200$ 에포크 학습한 뒤에는?
\end{exercise}

\begin{exercise}[\lvlC]
\label{ex:20:8}
\emph{탐구: 섞기인가, 재정규화인가?} 연습문제~\ref{ex:20:7}의 뉴런에 문턱값을 두고, 밤의 두 절차를 생각하자: (i)~모든 가중치에 $0.82$를 곱한다; (ii)~옛 패턴을 새 패턴과 섞어 반복한다. 각각은 가중치의 비와 뉴런의 응답에 무엇을 하는가? 잠든 뒤 시냅스 크기로 두 가설을 어떻게 구별하겠는가?
\end{exercise}

\chapter{살아 있는 뉴런으로 만든 기계}
\chaptermark{살아 있는 뉴런으로 만든 기계}
\label{sec:livingmachine}

\section{회로도가 보이는 테스트 벤치}

\ref{sec:debugging}장에서 우리는 회로도가 없는 기계를 디버깅했다. 프로브는 800억 개 뉴런 가운데 천 개만 들을 수 있고, 나머지는 모두 추측해야 한다. 이런 처지의 엔지니어라면 간단한 일을 할 것이다. 회로의 작은 조각을 브레드보드 위에 조립해서 부품 하나하나가 보이게 하는 것이다. 뉴런으로도 이것을 할 수 있다.

피질 뉴런을 떼어 내어 전극 배열이 있는 칩 위에서 키운다. 전극은 수십 개에서 수만 개까지다. 몇 주가 지나면 세포는 돌기를 뻗어 수천에서 수십만 개 뉴런의 네트워크로 연결된다. 대부분은 \nt{GLUT}이고 \nt{GABA}는 더 적은데, 그 비율은 표본에 따라 다르다. 해마 배양에서는 뉴런의 약 $6\%$이다~\cite{Benson1994}. 각 전극은 \ref{sec:debugging}장의 프로브처럼 들을 수도 있고, 테스트 신호처럼 전류를 넣을 수도 있다. 두 번째 종류의 테스트 벤치는 \emph{오가노이드}다. 인간 줄기세포에서 키운 지름 1밀리미터 정도의 3차원 신경 조직 덩어리이다~\cite{Lancaster2013}. 이런 테스트 벤치에서 살아 있는 네트워크는 몇 달 동안 작동한다. 오가노이드 실험을 네트워크로 원격 수행하면서 100일 넘게 연속으로 돌릴 수 있는 시설도 있다~\cite{Jordan2024}. 이 분야를 \emph{생물학적 컴퓨팅} 또는 \emph{오가노이드 지능}이라고 부른다~\cite{Smirnova2023}.

테스트 벤치는 뇌와 두 가지 중요한 점에서 다르다. 우선 아주 작다. 뉴런 수가 10만 배 적다. 그리고 브로드캐스트 채널이 없다. 피질 배양에는 \nt{DOPA}, \nt{SERO}, \nt{NORA}, \nt{ACET} 뉴런이 없다. 데이터 버스와 흐름 제어는 있지만 제어 레지스터는 없는 기계이다. 이 기계가 무엇을 할 수 있는지 보자.

\section{네트워크가 스스로 하는 일}

가만히 두어도 배양은 조용하지 않고, 고르게 잡음을 내지도 않는다. 이따금 네트워크 거의 전체가 한꺼번에 발화하는 \emph{일제 발화}가 일어나고, 다시 잠잠해진다. 일제 발화 사이의 간격은 배양에 따라 1초에서 몇 분이다~\cite{Wagenaar2006}. 엔지니어에게는 익숙한 그림이다. 강한 양의 되먹임을 걸고 부하를 떼어 낸 증폭기는 발진기가 된다.

메커니즘은 이미 알고 있다. \nt{GLUT} 뉴런 사이의 강한 상호 연결은 사태를 일으킨다. \ref{sec:gaba}장에서 명제~\ref{prop:ei}의 조건이 깨졌을 때와 같다. 사태는 시냅스의 신경전달물질 저장량을 빠르게 고갈시키고~\eqref{eq:depression}, 네트워크는 조용해진다. 저장량은 시간 상수 $\tau_{\text{rec}}$로 회복되고, 새 사태에 충분한 비율 $x^*$까지 회복되면 다음 일제 발화가 온다. 일제 발화 간격은
\begin{empheq}[box=\keybox]{equation}
T_{\text{burst}} \;\approx\; \tau_{\text{rec}}\,\ln\frac{1}{1-x^*} .
\label{eq:burstinterval}
\end{empheq}
\ref{sec:code}장의 $\tau_{\text{rec}}=0.8\,\text{s}$와 $x^*=0.9$를 넣으면 약 2초, $x^*=0.99$이면 약 4초가 나온다. 이것은 이 책의 $\tau_{\text{rec}}$를 쓴 모델의 추정값이며, 측정된 간격의 짧은 쪽 끝에 들어맞는다. 미세 배양에서 직접 측정해 보면 일제 발화 뒤 신경전달물질 저장량이 회복되는 데 수십 초가 걸린다~\cite{CohenSegal2011}. 즉 거기서는 $\tau_{\text{rec}}$가 더 크다. 그런 $\tau_{\text{rec}}$를 쓰면 식~\eqref{eq:burstinterval}은 느린 배양에서처럼 수십 초에서 몇 분을 준다. 이것은 이완형 발진기이며, \ref{sec:muscles}장의 리듬 발생기와 친척이다. 거기서는 뉴런 무리의 피로가 발진기를 다시 충전했고, 여기서는 신경전달물질 저장량이 그 일을 한다.

일제 발화는 네트워크가 할 일이 없을 때 하는 일이다. 살아 있는 뇌에서도 비슷한 그림이 깊은 잠(\ref{sec:sleep}장)과, 진짜 입력이 들어오기 전의 발달 중인 뇌에서 나타난다. 이 유사성은 시사하는 바가 있지만, 메커니즘이 같다는 것은 가설이다.

\section{도파민 교사 없이 네트워크를 가르치는 법}

배양에는 \nt{DOPA}가 없으므로 “더 좋다, 더 나쁘다”라는 신호는 우리가 직접 전극으로 넣어야 한다. 실험을 보면 테스트 벤치 위의 네트워크는 실제로 응답을 바꾼다. 그리고 가장 흥미로운 것은 \emph{무엇으로} 가르치느냐이다.

\emph{입을 다무는 교사.} 한 전극을 통해 1--3초마다 한 번씩 네트워크를 자극하다가, 다른 전극에서 자극 후 $50\pm10\ms$에 원하는 응답이 나타나면 자극을 5분 동안 끈다. 이런 주기를 몇 번 거치면 원하는 응답이 자극에 곧바로 나타난다~\cite{Shahaf2001}. 여기서 보상은 침묵이다. 자극은 네트워크를 흔들고, 자극을 멈추면 네트워크는 올바른 응답을 낸 상태에 머문다. 이것은 \ref{sec:dofa}장의 무작위 섭동을 통한 경사 하강(명제~\ref{prop:gradient})이며, 여기서는 흔드는 것 자체가 오차 역할을 한다. 맞을 때까지 흔든다.

\emph{테니스를 치는 네트워크.} 전극이 촘촘히 배열된 테스트 벤치 실험에서, 약 100만 개의 뉴런을 라켓이 하나인 컴퓨터 테니스 게임에 연결했다~\cite{Kagan2022}. 공의 위치는 “감각” 전극에 전류로 넣었다. 공이 어디 있는지는 장소로, 얼마나 먼지는 빈도로 부호화했다. 두 “운동” 영역의 활동이 라켓을 위아래로 움직였다. 학습 규칙은 특이했다. 라켓이 공을 받아 치면 네트워크는 짧고 \emph{예측 가능한} 자극을 받았고, 놓치면 몇 초 동안 \emph{예측 불가능한} 무작위 전류를 받았다. 20분 동안 게임하는 사이에 랠리가 길어졌다. 세션 안에서 유의미한 향상은 완전한 되먹임일 때만 있었다. 자극 대신 침묵을 준 경우에도 배양은 세션 끝에 되먹임이 없을 때보다 더 잘 쳤지만 효과가 작았고, 며칠에 걸친 세션 간의 안정적인 학습은 없었다. 이 실험의 자세한 분석은 \ref{sec:dishbrain}장에 있다.

저자들은 네트워크가 자기 입력이 예측 가능해지도록 행동한다는 가설로 결과를 설명했다~\cite{Friston2010}. \ref{sec:code}장의 스파이크 절약, \ref{sec:recognition}장의 프레임 예측과 같은 생각이다. 입력이 예상대로일 때 기계는 덜 소비한다. 그러나 실험 자체도, 특히 배양의 “지각력”에 관한 저자들의 표현도 날카로운 비판을 받았다. 효과가 작고 더 단순하게 설명할 수도 있다는 것이다~\cite{Balci2023}. 정직하게 평가하면 이렇다. 테스트 벤치의 배양이 되먹임에 따라 행동을 바꾼다는 것은 측정된 사실이고, 그것이 바로 자기 입력을 예측하도록 학습한다는 것은 가설이다.

\emph{몸을 움직이는 네트워크.} 비슷한 방식으로 배양을 단순한 가상 “몸”에 연결하고, 훈련 자극의 시공간 패턴을 골라 몸이 목표로 움직이도록 가르쳤다~\cite{Bakkum2008}. 이 실험들 어디에도 도파민은 없다. 교사는 엔지니어가 고른 전류 패턴이다. 교사가 프로그램이나 도파민 자체가 되면 무엇이 달라지는지는 \ref{sec:dishdopamine}--\ref{sec:cartpole}절에서 다룬다.

\begin{figure}[t]
\centering
\begin{tikzpicture}[>=Stealth,
  blk/.style={draw,very thick,rounded corners=2pt,align=center,font=\footnotesize,minimum height=0.8cm,fill=blue!5}]
% (a) closed loop
\node[font=\small] at (1.5,4.3) {(a) 닫힌 루프};
\node[blk,text width=2.6cm] (G) at (1.5,3.3) {게임: 공과 라켓};
\draw[very thick,fill=orange!10,rounded corners=3pt] (0.5,0.2) rectangle (2.5,1.6);
\foreach \x in {0.7,1.0,...,2.35}{\foreach \y in {0.4,0.7,1.0,1.3}{\fill[gray!60] (\x,\y) circle (0.04);}}
\draw[->,thick,blue!60!black] (G.west) -| (-0.5,0.9) -- (0.5,0.9);
\node[font=\scriptsize,blue!60!black,anchor=east] at (-0.6,2.1) {공 $\to$ 전류};
\draw[->,thick,red!70!black] (2.5,0.9) -- (3.5,0.9) |- (G.east);
\node[font=\scriptsize,red!70!black,anchor=west] at (3.6,2.1) {스파이크 $\to$ 라켓};
\node[font=\scriptsize] at (1.5,-0.2) {전극 배열 위의 뉴런};
\node[font=\scriptsize,align=center] at (1.5,-0.85) {히트: 예측 가능한 자극\\미스: 무작위 전류};
% (b) reservoir
\begin{scope}[xshift=7.9cm]
\node[font=\small] at (1.6,4.3) {(b) 저수지};
\node[blk,text width=1.1cm] (X) at (-0.4,1.0) {입력\\$u(t)$};
\draw[very thick,fill=orange!10] (1.6,1.0) circle (0.8);
\foreach \a/\r in {20/0.35,80/0.5,150/0.3,210/0.55,280/0.4,330/0.2,110/0.15}{\fill[red!60!black] ({1.6+\r*cos(\a)},{1.0+\r*sin(\a)}) circle (0.05);}
\node[font=\scriptsize] at (1.6,-0.2) {오가노이드: 교사 없음};
\node[blk,text width=1.8cm,fill=green!8] (W) at (4.0,1.0) {판독\\$W_{\text{out}}$};
\draw[->,thick] (X) -- (0.8,1.0);
\draw[->,thick] (2.4,1.0) -- node[above,font=\scriptsize] {$\mathbf r(t)$} (W);
\draw[->,thick] (W) -- (4.0,2.3) node[above,font=\scriptsize] {$y(t)$};
\node[font=\scriptsize,green!40!black] at (4.0,0.3) {교사 있는 학습};
\end{scope}
\end{tikzpicture}
\caption{살아 있는 네트워크가 계산하게 만드는 두 가지 방법. (a) 닫힌 루프: 공의 위치를 전류로 넣고, 운동 영역의 스파이크가 라켓을 움직이며, 타격 뒤의 예측 가능한 자극 또는 무작위 자극이 교사 역할을 한다. (b) 저수지~\eqref{eq:reservoir}: 오가노이드에는 학습 신호를 주지 않는다. 오가노이드는 입력을 기억을 가진 수많은 비선형 응답으로 쪼개고, 오차로 학습하는 것은 바깥의 선형 판독 하나뿐이다. 이때 오가노이드 자체의 연결도 교사 없이 바뀐다.}
\label{fig:livingmachine}
\end{figure}

\section{살아 있는 저수지}

살아 있는 네트워크를 쓰는 다른 방법도 있다. 아예 가르치지 않는 것이다. 공학에서는 이것을 \emph{저수지 컴퓨팅}(reservoir computing)이라고 부른다~\cite{Maass2002,JaegerHaas2004}. 기억을 가진 기성 비선형 시스템을 가져온다. 입력에 대한 응답이 풍부하고 최근 과거에 의존하기만 하면 무엇이든 된다. 입력 $u(t)$가 시스템을 흔들면 응답 벡터 $\mathbf r(t)$가 나오고, 학습하는 것은 선형 판독뿐이다.
\begin{empheq}[box=\keybox]{equation}
y(t) \;=\; W_{\text{out}}\,\mathbf r(t) ,
\label{eq:reservoir}
\end{empheq}
그 가중치는 \ref{sec:motorlearning}장의 최소제곱 규칙~\eqref{eq:lms}으로 맞춘다. 저수지는 \ref{sec:motorlearning}장의 작은 \nt{GLUT} 뉴런들이 하던 일을 한다. 입력을 긴 벡터로 쪼개어, 필요한 클래스들이 평면으로 분리될 수 있게 만든다(\ref{sec:dendrites}장).

전극 배열 위의 오가노이드는 이런 저수지 역할에 알맞았다. 전류 패턴으로 넣은 말소리를 선형 판독 하나만 학습시킨 뒤 화자별로 구분했더니 정확도가 약 $78\%$였다~\cite{Cai2023}. $78\%$라는 정확도는 저자들이 직접 보고하고 언론이 되풀이한 숫자이다. 쓸모 있는 결과지만, 그 안에서 “지능”이 어디에 있는지 이해하는 것이 중요하다. 오가노이드는 비선형성과 감쇠하는 기억을 제공하고, 오차에 의한 학습은 바깥의 실리콘에서 이루어진다(그림~\ref{fig:livingmachine}). 그렇다고 오가노이드가 그대로인 것은 아니다. 저자들에 따르면 훈련 동안 오가노이드 자체의 기능적 연결성이 바뀌었고, 이를 오가노이드 안에서 일어나는 교사 없는 학습이라고 부른다~\cite{Cai2023}. 정확도 가운데 얼마가 이 변화에서 오고 얼마가 판독에서 오는지는 분리되지 않았다.

\section{테스트 벤치가 기계에 대해 말해 주는 것}

\emph{에너지.} 추정~\eqref{eq:energy}에 따르면 스파이크 하나는 약 $10^{-10}$~J이 든다. 뉴런 100만 개의 배양이 초당 스파이크 하나의 빈도로 발화하면 신호 전달에 쓰는 전력은
\begin{empheq}[box=\keybox]{equation}
P \;\approx\; 10^{6}\times1\,\text{s}^{-1}\times10^{-10}\,\text{J} \;=\; 10^{-4}\,\text{W},
\label{eq:dishpower}
\end{empheq}
즉 0.1밀리와트이다. 이 스파이크를 자극하고 기록하고 처리하는 전자 장치는 몇 자릿수 더 많이 쓴다. \ref{sec:debugging}장에서처럼 병목은 계산이 아니라 계산과의 통신이다.

\emph{빠진 부품.} 테스트 벤치는 제어 레지스터 없이 \nt{GLUT} 버스와 \nt{GABA} 흐름 제어가 얼마나 많은 일을 할 수 있는지 보여 준다. 스스로 조직되고, 일제 발화를 만들고, 되먹임에 따라 응답을 바꾸고, 좋은 저수지 역할을 한다. 제어 레지스터 없이 할 수 없는 것은 무엇을 성공으로 볼지 스스로 정하는 일이다. 테스트 벤치에서는 엔지니어가 이 신호를 넣고, 뇌에서는 \ref{sec:dofa}--\ref{sec:acet}장에서 본 것처럼 브로드캐스트 채널이 넣는다.

\emph{윤리.} 오가노이드는 인간 세포에서 키우며, 복잡해질수록 그런 조직이 무언가를 느낄 수 있느냐는 질문이 무거워진다. 공학적 관점은 부분적인 답을 준다. \ref{sec:pain}장의 통증은 센서 하나의 신호가 아니라 관문, 음량 조절기, 브로드캐스트 채널을 갖춘 경보 시스템 전체이며, 테스트 벤치에는 그런 것이 없다. 그러나 이것은 추론이지 증명이 아니다. 이 분야 스스로도 처음부터 실험과 함께 윤리적 평가를 진행하자고 제안한다~\cite{Smirnova2023}.

\begin{keyblock}
\begin{proposition}[테스트 벤치 위의 살아 있는 네트워크]
\label{prop:livingmachine}
전극 배열 위의 \nt{GLUT} 뉴런과 \nt{GABA} 뉴런 네트워크는 스스로 일제 발화를 만들고~\eqref{eq:burstinterval}, 되먹임에 따라 응답을 바꾸며, 바깥에서 학습된 판독을 위한 저수지~\eqref{eq:reservoir} 역할을 할 수 있다. 이것은 모두 측정되었다. 이런 네트워크가 어떤 일반적 규칙, 예를 들어 자기 입력을 예측하는 규칙에 따라 학습한다는 것은 가설이며, 그 “지각력”에 관한 과장된 표현은 대다수 전문가가 받아들이지 않는다.
\end{proposition}
\end{keyblock}

\section{빛의 섬광으로 주는 도파민}
\label{sec:dishdopamine}

배양에 도파민을 교사로 준 직접 실험으로 우리가 아는 것은 하나뿐이다. 연구자들이 원격으로 접속하는 오가노이드 플랫폼에서 수행되었다~\cite{Jordan2024}. 오가노이드를 적시는 액체에 빛에 반응하는 “케이지”에 갇힌 도파민을 넣었다. 묶여 있는 분자는 수용체에 작용하지 않는다. 케이지 도파민의 농도는 $1$~mM이다. 네트워크에 보상을 주어야 할 때 자외선을 비추면 케이지가 분해되고 도파민이 부피 속으로 풀려난다.

루프는 이렇게 짜여 있다. 같은 자극을 되풀이해서 넣고, 그 자극이 이전보다 많은 스파이크를 일으키면 섬광을 켜서 도파민을 방출한다. $240$번 반복하는 동안 유발 스파이크 수는 전체적으로 늘었다. 저자들 스스로 이런 실험 하나만으로는 증가가 도파민 때문이라고 결론 내릴 수 없다고 쓴다~\cite{Jordan2024}.

엔지니어에게 이것은 3요소 규칙~\eqref{eq:threefactor}을 접시 안에 조립하려는 첫 시도이다. 송신자와 수신자의 활동이 흔적을 남기고, 공통 신호가 변화를 굳힐지 결정한다. 그러나 여기서 신호는 양수뿐이다. 보상이 있거나 없거나이며, 장치는 음의 오차 $\delta<0$를 내지 못한다. 게다가 도파민은 \eqref{eq:lowpass}의 농도처럼 퍼지고 천천히 제거되므로, 섬광이 잦으면 그저 전체 수준만 올라갈 수 있다. 학습을 이것과 구별하려면 대조 실험 두 가지가 필요하다. 네트워크의 응답과 무관한 무작위 시점에 같은 수의 섬광을 주는 것, 그리고 케이지 도파민 없이 같은 섬광을 주는 것이다. 또 휴식 뒤의 확인도 필요하다. 도파민이 사라진 뒤에도 변화가 남았는가?

\section{도파민이 실리콘을 가르친다}
\label{sec:biohybrid}

반대 방향의 실험에서는 살아 있는 세포가 전자 소자를 가르친다~\cite{Keene2020}. 도파민을 분비하는 세포를 유기 뉴로모픽 소자 위의 미세 채널에서 키운다. 이 소자는 트랜지스터이며, 채널의 전도도가 시냅스 가중치 역할을 한다. 세포에서 나온 도파민이 소자의 전극에서 산화되면서 전도도를 바꾼다. 이 장치는 틈새에서 신경전달물질을 치우는 메커니즘도 재현하므로, 가중치를 오래 바꿀 수 있을 뿐 아니라 원래대로 되돌릴 수도 있다.

회로로 보면 이것은 화학 입력을 받는 누설 적분기이다. 가중치는 도파민 흐름을 누적하고, “제거”가 얼마나 빨리 잊는지를 정한다. \eqref{eq:lowpass}와 같은 RC 회로이다. 중요한 것은 연결 방향이다. \ref{sec:debugging}장의 프로브는 브로드캐스트 레지스터를 보지 못한다. 그것이 화학적이기 때문이다. 이 소자는 바로 화학 레지스터를 읽어서 전기적 가중치로 바꾼다. 뇌-컴퓨터 인터페이스에게 이것은 데이터 버스뿐 아니라 제어 레지스터까지 듣는 센서로 가는 길이다.

\section{자기 도파민 뉴런을 가진 오가노이드}
\label{sec:midbrainorg}

인간 줄기세포로 \nt{DOPA} 뉴런이 있는 뇌 영역을 닮은 오가노이드를 키울 수 있다. 그 안에는 도파민을 분비하는, 작동하는 도파민 뉴런이 있다~\cite{Jo2016}. 이런 오가노이드는 주로 질병 연구를 위해 키운다. 그 도파민이 다른 네트워크의 교사 역할을 한 실험은 찾지 못했다.

살아 있는 뉴런으로 만든 기계에게 이것은 빠진 부품이다. \ref{sec:livingmachine}장의 테스트 벤치는 제어 레지스터 없는 데이터 버스와 흐름 제어이다. 여기에는 브로드캐스트 신호의 원천이 이미 있다. 모자란 것은 비평가, 즉 예측 오차~\eqref{eq:ctd}를 계산하고 이 뉴런들을 조종하는 네트워크이다. 피질 오가노이드를 이런 오가노이드와 연결해서 오차에 따라 도파민이 분비되게 하는 것은 열린 문제이며, 아직은 실험이 아니라 가설이다.

\section{도파민 없이 막대의 균형을 잡는 오가노이드}
\label{sec:cartpole}

최근 실험 가운데 가장 완전한 것은 도파민을 전혀 쓰지 않는다~\cite{Robbins2026}. 생쥐 피질 오가노이드를 “카트 위의 막대” 과제에 닫힌 루프로 연결했다. 오가노이드의 활동이 카트를 움직이고, 막대의 위치는 자극으로 되돌아 들어간다. 시도와 시도 사이에 오가노이드는 짧은 고주파 학습 자극을 받는다. 어떤 자극을 줄지는 강화 학습 프로그램이 고른다.

측정된 결과는 다음과 같다.
\begin{itemize}
\item 대부분의 오가노이드에서 프로그램이 고른 학습 신호가 무작위로 고른 신호나 신호가 없을 때보다 좋은 결과를 냈다.
\item $45$분 휴식 뒤에는 향상이 남지 않았다.
\item 두 종류의 글루탐산 수용체 AMPA와 NMDA를 모두 차단하면 향상이 사라졌다.
\end{itemize}

마지막 항목은 학습된 것이 어디에 사는지 말해 준다. \nt{GLUT} 시냅스 안이며, \ref{sec:weightstore}절에서 입력과 출력의 동시성 검출기로 작동하는 수용체를 통해서이다. 두 번째 항목은 무엇이 모자란지 보여 준다. 빠른 변화는 있지만 굳히기는 없다. \ref{sec:motorlearning}장에서 느린 학습자 없이 빠른 학습자만 있을 때와 같다. 교사의 역할도 바뀌었다. 전류 패턴을 고르는 엔지니어 대신 프로그램이 들어섰지만, 교사는 여전히 접시 바깥에 있다.

배양을 전극 배열에서 학습시킨 더 이른 실험들 중에서도 도파민을 쓴 것은 하나도 찾지 못했다. 학습 신호는 모두 전기 자극이었다.

\section{누가 배양을 가르치는가}

\begin{table}[t]
\centering
\footnotesize
\setlength{\tabcolsep}{4pt}
\renewcommand{\arraystretch}{1.15}
\begin{tabular}{>{\raggedright}p{2.7cm}>{\raggedright}p{2.5cm}>{\raggedright}p{3.2cm}>{\raggedright\arraybackslash}p{4.4cm}}
\toprule
실험 & 가르치는 주체 & 학습 신호 & 알려진 것 \\
\midrule
원하는 응답에서 자극 멈추기 & 엔지니어 & 자극의 종료 & 응답이 굳어진다 --- 측정됨 \\
DishBrain (\ref{sec:dishbrain}장) & 엔지니어 & 예측 가능한 전류 또는 무작위 전류 & 세션 안의 유의미한 랠리 증가는 완전한 되먹임일 때만, 침묵일 때는 효과가 더 작다 --- 측정됨. 세션 간에는 불안정. 원인은 논쟁 중 \\
카트 위의 막대 & 강화 학습 프로그램 & 프로그램이 고른 고주파 자극 & 무작위보다 낫지만 휴식 뒤에는 남지 않는다 --- 측정됨 \\
섬광으로 주는 도파민 & “스파이크가 늘면 보상” 규칙 & 빛으로 풀려난 도파민 & 스파이크 증가 --- 관찰. 도파민의 역할은 증명되지 않음 \\
바이오하이브리드 시냅스 & 살아 있는 세포 & 세포에서 나온 도파민 & 소자 가중치의 장기 변화와 복귀 --- 측정됨 \\
\nt{DOPA} 뉴런을 가진 오가노이드 & --- & --- & 도파민 원천은 있다. 교사로 쓰인 적은 없다 \\
\bottomrule
\end{tabular}
\caption{누가 무엇으로 칩 위의 살아 있는 뉴런을 가르치는가.}
\label{tab:dishteachers}
\end{table}

배양 자체가 학습하는 모든 실험에서 교사는 아직 바깥에 있다(표~\ref{tab:dishteachers}). 엔지니어, 프로그램, 또는 엔지니어가 정한 규칙이다. 도파민이 접시 안에 들어온 것은 단 한 번뿐이고, 그 역할은 아직 증명해야 한다. \ref{sec:dofa}장에서처럼 비평가, 오차, 흔적을 갖춘 진짜 브로드캐스트 채널을 접시 안에 조립하는 것은 아직 아무도 내딛지 않은 다음 걸음이다.


\section{정리}

\begin{table}[t]
\centering
\footnotesize
\setlength{\tabcolsep}{4pt}
\renewcommand{\arraystretch}{1.15}
\begin{tabular}{>{\raggedright}p{2.8cm}>{\raggedright}p{4.2cm}>{\raggedright\arraybackslash}p{5.8cm}}
\toprule
실험 & 살아 있는 네트워크가 하는 일 & 알려진 것 \\
\midrule
입력 없는 네트워크 & 1초에서 몇 분 간격의 일제 발화~\eqref{eq:burstinterval} & 측정됨. 시냅스 고갈을 통한 메커니즘은 널리 받아들여진 모델 \\
입을 다무는 교사 & 자극을 끄면 원하는 응답이 굳어진다 & 측정됨 \\
테니스 게임 & 랠리가 길어진다. 세션 안의 유의미한 향상은 완전한 되먹임일 때만 & 행동 변화는 측정됨. 세션 간 학습은 불안정. 효과의 크기와 입력 예측이라는 설명은 논쟁 중 \\
가상의 몸 & 알맞게 고른 자극 패턴으로 목표를 향해 움직인다 & 측정됨 \\
저수지 & 판독~\eqref{eq:reservoir}을 위한 비선형성과 기억 & 측정됨. 오차에 의한 학습은 바깥의 실리콘에서. 오가노이드의 연결도 바뀐다 \\
에너지 & 뉴런 100만 개당 약 $10^{-4}$~W~\eqref{eq:dishpower} & \eqref{eq:energy}에 따른 추정 \\
\bottomrule
\end{tabular}
\caption{살아 있는 뉴런으로 만든 기계가 보여 준 것.}
\label{tab:livingmachine}
\end{table}

살아 있는 뉴런으로 만든 기계(표~\ref{tab:livingmachine})는 제어 레지스터 없이 데이터 버스와 흐름 제어가 무엇을 할 수 있는지 보여 주는 테스트 벤치이다. 이들은 많은 것을 할 수 있지만, 무엇이 좋은지를 스스로 정하지는 못한다. 테스트 벤치에서는 자기 전류 패턴을 가진 엔지니어가 이 역할을 하고, 뇌에서는 이 책의 중간 부분에서 다룬 몇 안 되는 브로드캐스트 전선이 한다.

\section{연습문제}

\begin{exercise}[\lvlA]
\label{ex:21:1}
\emph{일제 발화 간격.} 일제 발화가 저장량을 0까지 고갈시키고, 그 뒤 $\tau_{\text{rec}}\,\dd x/\dd t=1-x$이며, $x=x^*$에서 새 일제 발화가 시작된다. $\tau_{\text{rec}}=0.8$~s. (a)~$x^*=0.9$와 $0.99$일 때의 간격을 구하라. (b)~문턱값 $x^*$가 $0.99$ 근처에서 $\pm0.005$만큼 무작위로 변한다. 간격은 어느 범위에 있는가?
\end{exercise}

\begin{exercise}[\lvlA]
\label{ex:21:2}
\emph{테스트 벤치의 에너지.} (a)~추정~\eqref{eq:dishpower}은 뇌 전체($8.6\cdot10^{10}$개 뉴런)에 대해 얼마의 전력을 주는가? (b)~테스트 벤치가 전극 $1024$개를 $20$~kHz로, 표본당 $10$~비트로 기록한다. 비트당 $1$~nJ일 때 이 데이터 흐름의 전송은 배양 자체의 전력보다 몇 배 비싼가?
\end{exercise}

\begin{exercise}[\lvlB]
\label{ex:21:3}
\emph{설계: 일제 발화 억제.} 여러 전극을 통한 드문 자극 때문에 모든 시냅스가 빈도 $r_b$로 신경전달물질을 방출한다. 저장량은 $\tau_{\text{rec}}\,\dd x/\dd t=1-x-Ur_b\tau_{\text{rec}}x$, $U=0.5$, $\tau_{\text{rec}}=0.8$~s. 저장량이 $x^*=0.9$; $0.99$에 이르지 못하게 하는 가장 작은 $r_b$는 얼마이며, 그때 시냅스는 얼마나 약해지는가?
\end{exercise}

\begin{exercise}[\lvlB]
\label{ex:21:4}
\emph{저수지의 기억은 뉴런 수를 넘지 않는다.} 출력이 $N$개인 저수지 $\mathbf r(t)$가 평균 0, 분산 1인 백색 입력 $u(t)$를 받는다. $m_k$는 판독 $\mathbf w_k\cdot\mathbf r(t)$와 $u(t-k)$의 상관계수 제곱의 최댓값이다. 기억 용량 $\mathrm{MC}=\sum_{k\ge0}m_k\le N$임을 증명하라.
\end{exercise}

\begin{exercise}[\lvlB]
\label{ex:21:5}
\emph{모의실험: 실리콘 저수지.} 저수지 $\mathbf r(t+1)=\tanh(W\mathbf r(t)+\mathbf v\,u(t)+\mathbf c)$, $N=30$, $W$는 스펙트럼 반지름 $0.9$인 무작위 행렬, $v_i\sim U(-0.5,0.5)$, $c_i\sim U(-0.2,0.2)$, $u\sim U(-1,1)$, 판독은 릿지 회귀이다. $\tanh$가 있을 때와 없을 때 연습문제~\ref{ex:21:4}의 기억 용량을 구하라. 어느 저수지가 더 많이 기억하는가?
\end{exercise}

\begin{exercise}[\lvlB]
\label{ex:21:6}
\emph{살아 있는 저수지의 판독.} 오가노이드가 채널 $1024$개와 두 클래스의 학습 기록 $P=240$개를 준다. (a)~연습문제~\ref{ex:19:3}의 공식으로 무작위 라벨이 분리될 확률이 $1\%$ 이하가 되려면 특징을 몇 개($n$)까지 남길 수 있는가? (b)~검증 기록 $100$개에서 정확도 $78\%$는 찍기보다 몇 시그마 높은가?
\end{exercise}

\begin{exercise}[\lvlC]
\label{ex:21:7}
\emph{연구: 브로드캐스트 채널 없는 교사.} 전극 $8$--$1024$개로 테스트 벤치에서 \ref{sec:dofa}장의 3요소 규칙을 구현하는 자극 프로토콜과, 가중치 학습을 네트워크 상태의 이동과 구별하는 판독 고정 대조 실험을 제안하라. 어떤 결과가 학습 가설을 반증하겠는가?
\end{exercise}

\chapter{DishBrain}
\chaptermark{DishBrain}
\label{sec:dishbrain}

\section{접시 속 네트워크가 테니스를 친다}

DishBrain은 전극이 있는 칩 위에서 키운 살아 있는 뉴런 배양이 라켓 하나짜리 단순화된 테니스를 치는 테스트 벤치에 저자들이 붙인 이름이다~\cite{Kagan2022}. 살아 있는 뉴런으로 만든 기계를 다룬 실험 가운데 가장 많이 논의된 것이며(이런 기계의 개관은 \ref{sec:livingmachine}장), 이 책에는 특히 알맞다. 여기서는 작은 네트워크 전체에 손이 닿는다. 입력은 모두 실험자가 정하고, 출력은 모두 기록된다. 앞 장들의 질문, 즉 신호가 어떻게 부호화되는지, 누가 가르치는지, 프로브에 무엇이 보이는지를 눈도 근육도 \nt{DOPA} 뉴런도 없는 네트워크에 던질 수 있다. 엔지니어가 남의 테스트 벤치를 분석하듯 이 실험을 뜯어보자. 회로, 입력, 출력, 교사, 결과, 논쟁 순이다.

\section{테스트 벤치}

뉴런은 생쥐 태아의 피질에서 칩당 약 $800\,000$개를 얻거나, 인간 줄기세포에서 칩당 약 100만 개를 키운다. 뉴런은 몇 주 동안 칩 위에서 자라며 서로 얽혀 들어가고, 스스로 스파이크와 버스트를 내기 시작한다. 그 아래 $8$~mm$^2$ 면적에 백금 전극 $26\,000$개가 깔려 있다. 동시에 기록할 수 있는 것은 최대 $1024$개이고, 자극은 실제로 여덟 개를 통해 준다.

배양 둘레에 루프가 닫혀 있다(그림~\ref{fig:dishloop}). 컴퓨터가 게임을 모의한다. 공이 날아가 벽에 튕기고, 라켓이 위아래로 움직인다. 공의 위치는 여덟 개의 “감각” 전극에 흐르는 전류로 바뀐다. 두 “운동” 영역의 스파이크를 $10\ms$ 창마다 세어 라켓을 움직인다. 매번 공을 치거나 놓친 뒤 배양은 자극을 하나 더 받는다. 이것이 되먹임이다. $10\ms$ 창은 배양이 아니라 테스트 벤치 컴퓨터의 클록이다. 네트워크 자체는 이 책의 모든 네트워크처럼 여전히 비동기로 작동한다.

\begin{figure}[t]
\centering
\begin{tikzpicture}[>=Stealth,
  blk/.style={draw,very thick,rounded corners=3pt,align=center,font=\small,minimum height=1.2cm},
  lab/.style={font=\scriptsize,align=center}]
\node[blk,fill=blue!6,text width=3.4cm] (C) at (0,0) {칩 위의 배양\\\scriptsize 뉴런 $\sim10^6$개, 전극 $26\,000$개};
\node[blk,fill=orange!10,text width=3.0cm] (G) at (7.2,0) {컴퓨터:\\테니스 게임};
\draw[->,very thick,blue!60!black] (G.north west) to[bend right=25] node[lab,above] {공: \emph{장소} --- 8개 전극 중 어느 것,\\\emph{빈도} --- 초당 4--40회} (C.north east);
\draw[->,very thick,red!70!black] (C.south east) to[bend right=25] node[lab,below] {라켓: “위”와 “아래” 영역의\\$10\ms$당 스파이크 수} (G.south west);
\draw[->,thick,densely dashed,green!45!black] (G.east) -- ++(0.9,0) |- ++(0,-2.6) -| node[lab,pos=0.25,below] {히트: 예측 가능한 자극, 미스: 예측 불가능한 자극} (C.south);
\end{tikzpicture}
\caption{DishBrain 루프. 파란 화살표는 입력이다. 공의 위치를 위치 부호와 발화율 부호로 적는다(\ref{sec:code}장). 빨간 화살표는 출력이다. 두 영역의 스파이크 수 차이가 라켓을 움직인다. 초록 점선은 랠리마다 주는 되먹임으로, 테스트 벤치의 유일한 “교사”이다.}
\label{fig:dishloop}
\end{figure}

\section{입력과 출력}

\emph{입력}은 \ref{sec:code}장의 부호 두 가지로 동시에 적는다. 여덟 전극 가운데 어느 것을 자극하느냐가 공의 높이를 알려 준다. 이것은 \ref{sec:sensors}장의 센서처럼 위치 부호이다. 자극 펄스가 얼마나 자주 오느냐가 공이 얼마나 먼지를 알려 준다. 공이 먼 벽에 있으면 초당 $4$회이고, 라켓 쪽 벽에 가까워질수록 선형으로 늘어 초당 $40$회가 된다. 이것은 발화율 부호이며, 공 신호 펄스의 진폭은 $75$~mV이다. 배양은 이 부호들 가운데 어느 것도 미리 “알지” 못한다. 부호는 실험자가 정했고, 네트워크는 그것을 읽는 법을 배워야 한다.

\emph{출력}은 \ref{sec:debugging}장의 인터페이스와 같은 방식으로 읽는다. 한 운동 영역의 스파이크는 라켓을 위로, 다른 영역의 스파이크는 아래로 민다. 가장 단순하게 쓰면 창마다
\begin{empheq}[box=\keybox]{equation}
\Delta p \;=\; c\,\bigl(n_\uparrow - n_\downarrow\bigr),
\label{eq:dishreadout}
\end{empheq}
여기서 $n_\uparrow$, $n_\downarrow$는 $10\ms$ 동안 두 영역의 스파이크 수이고, $c$는 테스트 벤치의 계수이다. 이것은 \ref{sec:glutcomputer}장의 이중 레일 부호이다. 움직임의 부호는 신호의 부호가 아니라 두 전선 가운데 어느 쪽이 더 큰 소리를 내느냐가 전한다.

이런 출력의 잡음을 보자. 한 영역이 모두 합쳐 초당 $R$개의 스파이크를 낸다면, 창 $T=10\ms$ 동안 평균 $RT$개이고, \eqref{eq:rateerror}에 따라 상대 산포는 $1/\sqrt{RT}$이다. $R=1000$이면 창마다 약 $30$퍼센트, $R=100$이면 100퍼센트가 넘는다. 라켓은 떨릴 수밖에 없고, 네트워크가 배울 수 있는 방법은 하나뿐이다. 스파이크 수 차이의 \emph{평균}을 원하는 쪽으로 옮기는 것이다. \ref{sec:debugging}장의 모든 인터페이스를 제한하는 것과 같은 법칙이 이제는 루프의 양쪽에 걸린다.

\section{도파민 없는 교사}

이 테스트 벤치에서 가장 흥미로운 것은 교사이다. 피질 뉴런 배양에는 \nt{DOPA} 뉴런이 없고, 테스트 벤치는 “예상보다 좋다”는 신호를 전혀 보내지 않는다. 되먹임은 다른 것이다. 네트워크가 공을 받아 치면 여덟 감각 전극 모두에 짧고 \emph{예측 가능한} 자극이 한꺼번에 들어간다. $75$~mV 펄스를 초당 $100$회, $0.1\,\text{s}$ 동안이다. 놓치면 저자들이 \emph{예측 불가능하다}고 부르는 자극이 4초 동안 이어진다. $150$~mV 펄스를 초당 $5$회, 무작위 시점에 무작위로 고른 전극에 준다. 그 뒤 자극 없는 4초의 휴지가 있고, 공이 다시 출발한다~\cite{Kagan2022}.

저자들은 네트워크가 자기 입력이 예측 가능해지도록 행동한다는 가설에서 출발했다~\cite{Friston2010}. 엔지니어에게 의미는 간단하다. 배양이 4초의 무작위 잡음에서 벗어나는 방법은 정확히 하나, 공을 받아 치는 것이다. 같은 생각을 우리는 \ref{sec:code}장의 스파이크 절약과 \ref{sec:recognition}장의 프레임 예측에서 만났다.

그런데 여기서 꼭 예측 가능성이 필요할까? 더 거친 메커니즘도 가능하다. 미스 뒤의 예측 불가능한 자극이 네트워크의 최근 변화를 그저 \emph{뒤섞고}, 히트 뒤의 잔잔한 자극은 그것을 건드리지 않는다고 하자. 네트워크의 설정을 벡터 $\boldsymbol\psi$ 하나로 나타내고, 설정 $\boldsymbol\psi$에서 네트워크가 확률 $P_{\text{miss}}(\boldsymbol\psi)$로 공을 놓친다고 하자. 흔들기가 대칭이면, 즉 $\boldsymbol\psi$에서 $\boldsymbol\psi'$로 갈 확률이 되돌아올 확률과 같으면, 정상 상태에서 각 설정을 떠나는 흐름은 들어오는 흐름과 균형을 이루고, 네트워크가 설정 $\boldsymbol\psi$에 머무는 시간의 비율은
\begin{empheq}[box=\keybox]{equation}
\rho(\boldsymbol\psi) \;\propto\; \frac{1}{P_{\text{miss}}(\boldsymbol\psi)} .
\label{eq:dishdensity}
\end{empheq}
열 배 덜 놓치는 설정은 열 배 오래 유지된다. 가중치를 어느 쪽으로 바꿀지 네트워크에 알려 주는 것은 아무것도 없다. 좋은 설정이 덜 자주 부서질 뿐이다.

우리는 배양을 두 숫자로 줄인 모델로 이것을 확인했다. 공이 높이 $y$에 오면 라켓은 점 $a\,y+b$에 오고, 빗나간 거리가 $0.1$보다 작으면 받아 친다. 미스 뒤에는 두 숫자 모두 무작위로 흔들리고, 히트 뒤에는 그대로 남는다. 랠리 $2000$번 동안 받아 친 비율은 처음 200번의 $0.21$에서 마지막 500번의 $0.38$로 올랐다(“배양” 40개, 그림~\ref{fig:dishmodel}). 흔들기를 \emph{모든} 랠리 뒤에 주면 되먹임에 정보가 없고, 비율은 약 $0.12$에 머문다. 흔들기가 아예 없으면 $0.19$에 멈춰 있다. 테스트 벤치의 실험도 이와 맞는다. 세션 안의 유의미한 향상은 완전한 되먹임일 때만 있었다. 자극 대신 침묵을 주었을 때도 배양은 세션 끝에 되먹임이 없을 때보다 잘 쳤지만 효과가 작았다~\cite{Kagan2022}. 이 조건에서도 미스 뒤에는 긴 휴지가, 히트 뒤에는 짧은 휴지가 오므로 결과 사이의 차이가 완전히 사라지지는 않는다는 점에 주의하자.

\begin{figure}[t]
\centering
\begin{tikzpicture}[x=1cm,y=5cm]
\draw[->,gray!70!black] (0,0) -- (10.6,0);
\draw[->,gray!70!black] (0,0) -- (0,0.5) node[above,font=\small,black] {받아 친 비율};
\foreach \v in {0.1,0.2,0.3,0.4}{\draw[gray!60!black] (-0.06,\v) -- (0.06,\v); \node[left,font=\scriptsize] at (0,\v) {\v};}
\foreach \x/\f/\l/\lab in {1.8/0.21/0.38/{미스 뒤\\잡음},5.2/0.13/0.11/{모든 랠리 뒤\\잡음},8.6/0.19/0.19/{잡음 없음}}{
  \fill[blue!35] (\x-0.8,0) rectangle (\x-0.05,\f);
  \fill[red!60!black] (\x+0.05,0) rectangle (\x+0.8,\l);
  \node[below,font=\scriptsize,align=center] at (\x,-0.01) {\lab};
  \node[above,font=\scriptsize] at (\x-0.425,\f) {\f};
  \node[above,font=\scriptsize] at (\x+0.425,\l) {\l};}
\fill[blue!35] (7.4,0.47) rectangle (7.7,0.495); \node[right,font=\scriptsize] at (7.75,0.482) {처음 200};
\fill[red!60!black] (7.4,0.41) rectangle (7.7,0.435); \node[right,font=\scriptsize] at (7.75,0.422) {마지막 500};
\end{tikzpicture}
\caption{“미스 때 뒤섞기” 모델(\texttt{dishbrain\_model.py}): 랠리 $2000$번의 처음과 끝에서 받아 친 비율, 모델 배양 40개의 평균. 흔들기가 히트가 아니라 미스 뒤에 올 때만 학습이 일어난다. 세션 안의 유의미한 향상이 완전한 되먹임일 때만 있었던 실험과 같다.}
\label{fig:dishmodel}
\end{figure}

\begin{keyblock}
\begin{proposition}[뒤섞는 교사]
\label{prop:dishscramble}
실패가 네트워크를 흔들고 성공이 네트워크를 가만두면, 네트워크는 스스로 덜 틀리는 설정 쪽으로 표류한다. 한 설정에 머무는 시간은 실패 확률에 반비례한다~\eqref{eq:dishdensity}. 이런 학습에는 보상 신호도, 그 부호도, 예측 가능성 자체도 필요 없다. 성공 뒤의 되먹임과 실패 뒤의 되먹임이 다르기만 하면 된다. 배양의 행동 변화는 측정되었다. 접시 안에서 작동하는 메커니즘이 입력 예측인지 뒤섞기인지는 밝혀지지 않았다.
\end{proposition}
\end{keyblock}

\ref{sec:dofa}장과 비교해 보자. 거기서도 잡음은 탐색 역할을 했지만, 도파민에는 부호가 있었다. 오차 $\delta$가 가중치를 어느 쪽으로 옮길지 알려 주었고, 걸음은 경사를 따랐다~\eqref{eq:perturbation}. 여기에는 부호가 없고 “두기”와 “흔들기”만 있다. 이런 탐색은 더 거칠고 느리지만 네트워크에 요구하는 것이 적다. \nt{DOPA} 뉴런도, 비평가도, 몇 초에 걸친 흔적도 필요 없다.

\section{무엇을 측정했고 무엇을 두고 다투는가}

게임 한 판은 $20$분이었고, 처음 $5$분과 마지막 $15$분을 비교했다. 완전한 되먹임을 받은 배양에서는 랠리가 길어지고 즉시 미스가 줄었다. 세포 없는 대조, 게임 없는 대조, 컴퓨터 “라켓” 대조에서는 그런 일이 없었다. 인간 세포 배양은 평균적으로 생쥐 배양보다 오래 받아 쳤다. 향상은 통계적으로 확실하지만(생쥐 배양 9개와 인간 배양 11개, 각 그룹 100판 남짓) 작다. 네트워크는 좋은 선수가 되지 않았고, 무작위보다 조금 나아졌을 뿐이다. 향상이 확실한 것은 세션 안에서이고, 며칠에 걸친 세션 간의 안정적인 학습은 저자들도 얻지 못했다~\cite{Kagan2022}. 같은 그룹의 후속 연구는 게임 중 배양의 활동이 \emph{임계} 상태, 즉 감쇠와 사태 사이의 경계에 다가간다는 것을 발견했다. \ref{sec:gaba}장의 바로 그 경계 위의 삶이다. 가까울수록 게임을 더 잘했다~\cite{Habibollahi2023}.

가장 큰 논쟁을 부른 것은 숫자가 아니라 표현이었다. 논문 제목은 접시 속 뉴런이 “지각력”(sentience)을 보인다고 주장했고, 한 무리의 연구자들이 반박했다. 닫힌 루프 안의 행동 변화는 아직 지능도 감각도 아니며, 결론은 더 겸손하게 써야 한다는 것이다~\cite{Balci2023}. 공학적으로는 두 질문이 열려 있고, 둘 다 메커니즘에 관한 것이다. 첫째, 네트워크가 학습하는가, 즉 가중치가 오래 바뀌는가, 아니면 되먹임은 현재 상태를 옮길 뿐인가? 둘째, 꼭 예측 가능성이 필요한가, 아니면 명제~\ref{prop:dishscramble}에서처럼 성공과 실패에 대한 응답이 다르기만 하면 충분한가? 모든 전극에 손이 닿는 테스트 벤치는 두 질문 모두에 답할 수 있게 해 준다. 아마 이것이 이 테스트 벤치의 가장 큰 가치일 것이다.

\section{테스트 벤치 사양: 재현하는 법}
\label{sec:dishspec}

아래에는 이런 테스트 벤치를 다시 조립하는 데 필요한 모든 것을 모았다. 장비, 루프, 입력 부호화, 출력 판독, 되먹임, 실험 프로토콜이다. 매개변수마다 출처를 표시했다. “논문”은 값을 논문~\cite{Kagan2022}의 본문에서 가져왔다는 뜻이다. “선택”은 논문에 값이 없어서 재현하려면 직접 정해야 한다는 뜻이다. 이때는 기본값을 제안하고 그것을 확인하는 방법을 적는다.

\subsection*{장비와 배양}

\begin{center}
\footnotesize
\setlength{\tabcolsep}{4pt}
\renewcommand{\arraystretch}{1.15}
\begin{tabular}{>{\raggedright}p{4.0cm}>{\raggedright}p{6.9cm}>{\raggedright\arraybackslash}p{1.6cm}}
\toprule
매개변수 & 값 & 출처 \\
\midrule
칩 & MaxOne 배열: $8$~mm$^2$ 면적에 백금 전극 $26\,000$개 & 논문 \\
기록 & 동시에 최대 $1024$채널, 초당 $20\,000$ 표본 & 논문 \\
자극 & 기술적으로는 최대 $32$개 전극, 실험에서는 독립 전극 $8$개 & 논문 \\
세포 & 생쥐 태아 피질. 줄기세포에서 얻은 인간 뉴런(두 가지 배양법). 대조용 HEK293T 세포(뉴런 아님) & 논문 \\
파종 & 칩당 약 $10^6$개 세포 & 논문 \\
실험 시점의 배양 나이 & 생쥐: $\sim10$일부터. 인간: 배양법에 따라 $14$--$24$일 또는 $\sim80$일 & 논문 \\
\bottomrule
\end{tabular}
\end{center}

\subsection*{루프와 시간}

루프는 $10\ms$ 창($200$ 표본) 단위로 돈다. 창마다 운동 영역 전극의 스파이크를 세고, 게임을 갱신하고, 다음 자극을 내보낸다. 배양 속 스파이크에서 응답 자극까지의 지연은 약 $5\ms$이다. 게임 중 각 전극의 신호는 차단 주파수 $100$~Hz인 2차 베셀 고역 통과 필터를 지난다. 신호의 절댓값은 $1$~Hz의 1차 베셀 저역 통과 필터로 평활하며, 스파이크 문턱값은 이 평활된 절댓값에 비례한다. 배경의 표준편차 여섯 배라는 문턱값은 논문이 게임이 아니라 개별 활동 측정에 대해 밝힌 것이다. 자극이 배양의 응답으로 세어지지 않도록, $75$~mV를 넘는 큰 스파이크가 동시에 $15$개보다 많이 오면 모든 신호를 차단한다. 모두 논문에 따른 것이다.

\subsection*{공의 부호화}

\begin{center}
\footnotesize
\setlength{\tabcolsep}{4pt}
\renewcommand{\arraystretch}{1.15}
\begin{tabular}{>{\raggedright}p{3.4cm}>{\raggedright}p{7.5cm}>{\raggedright\arraybackslash}p{1.6cm}}
\toprule
매개변수 & 값 & 출처 \\
\midrule
감각 영역 & 전극 $8$개. 이웃한 전극이 이웃한 공 위치를 뜻하도록 배치 & 논문 \\
위치 부호 & 전극 번호 $k=1,\dots,8$이 라켓 중심에 대한 공의 위치를 정한다 & 논문 \\
어느 좌표인가 & 공의 수직 위치. 재현에서는 라켓 기준으로, $y-p\in[-\tfrac12,\tfrac12]$일 때 $k=\lceil 8(y-p+\tfrac12)\rceil$ & 논문; 공식은 선택 \\
발화율 부호 & 초당 $4$--$40$회의 자극 빈도가 라켓까지의 거리를 전한다 & 논문 \\
눈금 방향 & 공이 반대쪽 벽에 있으면 초당 $4$회, 라켓 쪽 벽에 가까워질수록 선형으로 초당 $40$회까지: $f=4+36\,(1-x/L)$, $x$는 라켓 쪽 벽까지의 거리, $L$은 코트 너비 & 논문 \\
펄스 모양 & 짧은 직사각형 2상 펄스: 먼저 양의 위상, 다음 음의 위상 & 논문 \\
진폭 & $75$~mV. 억제된 세포를 억지로 발화시키지 않도록 고른 값 & 논문 \\
위상 길이 & 밝히지 않음. 자극 시스템의 표준 설정을 쓰고 실험 중에는 바꾸지 않는다 & 선택 \\
\bottomrule
\end{tabular}
\end{center}

\subsection*{라켓 판독}

논문에는 운동 영역이 두 개 있다. 첫째 영역의 스파이크는 라켓을 위로, 둘째 영역의 스파이크는 아래로 움직인다. 세포 밀도와 활동의 차이를 반영하도록 두 영역의 기여에 가중치를 준다~\cite{Kagan2022}. 정확한 공식은 나와 있지 않다. 재현에서는 $10\ms$ 창마다 규칙~\eqref{eq:dishreadout} $\Delta p=c\,(n_\uparrow-n_\downarrow)$를 쓰고, 휴지 상태의 평균 활동으로 두 영역을 맞추는 가중치를 준다(선택). 즉 $n_\uparrow$와 $n_\downarrow$를 게임 전에 잰 각 영역의 창당 평균 스파이크 수로 나눈다. 계수 $c$는 평균 스파이크 수 하나만큼의 차이가 창마다 라켓을 코트 높이의 $1\%$만큼 옮기도록 고른다(선택). 그러면 한 영역이 계속 우세할 때 라켓은 $1$~s 동안 코트 전체를 가로지른다.

칩 위의 영역 배치는 논문 본문에 좌표로 주어지지 않았고 모식도만 있다. 재현에는 겹치지 않는 전극 그룹 세 개면 충분하다. 가운데에 감각 전극 여덟 개, 양옆에 기록 전극 그룹 하나씩, 하나는 “위”, 다른 하나는 “아래”이다(선택).

\subsection*{게임}

코트 크기, 라켓 길이, 공의 속도는 논문에 없다. 공이 일정한 속도로 날아가 벽에 튕긴다는 말뿐이다. 재현용(선택): 코트는 $1\times1$, 라켓은 오른쪽 벽에 길이 $0.2$(모델~\texttt{models/dishbrain\_model.py}에서처럼 $|y-p|<0.1$이면 히트), 공은 $1$~s 만에 코트를 가로지른다. 한 번도 받아 치지 못하고 놓친 랠리는 “에이스”, 연속 세 번보다 많이 받아 친 랠리는 “긴 랠리”로 센다(논문).

\subsection*{되먹임}

\begin{center}
\footnotesize
\setlength{\tabcolsep}{4pt}
\renewcommand{\arraystretch}{1.15}
\begin{tabular}{>{\raggedright}p{2.6cm}>{\raggedright}p{8.3cm}>{\raggedright\arraybackslash}p{1.6cm}}
\toprule
사건 & 주는 것 & 출처 \\
\midrule
히트 & 예측 가능한 자극: 감각 전극 $8$개 모두에 동시에, $75$~mV, 초당 $100$회, $100\ms$. 그동안 평소의 공 자극은 중단된다 & 논문 \\
미스 & 예측 불가능한 자극: $150$~mV, 초당 $5$회, $4$~s, 무작위 시점에 $8$개 가운데 무작위로 고른 전극에. 그 뒤 $4$~s 동안 자극 없음, 공은 무작위 방향으로 출발 & 논문 \\
무작위성의 법칙 & 밝히지 않음. 재현에서는 펄스 시점을 평균 빈도 초당 $5$회의 푸아송 과정으로 & 선택 \\
\bottomrule
\end{tabular}
\end{center}

\subsection*{프로토콜과 대조 실험}

세션 하나는 $20$~분이며, 처음 $5$~분과 마지막 $15$~분을 비교한다(논문). 결과 척도는 세 가지이다. 평균 랠리 길이, 에이스 비율, 긴 랠리 비율. 실험 조건(논문):
\begin{itemize}
\item \emph{자극}: 위에서 설명한 완전한 루프.
\item \emph{침묵}: 히트나 미스의 자극 대신 같은 시간 동안 자극을 전혀 주지 않고, 그 뒤 공이 무작위 방향으로 출발한다.
\item \emph{되먹임 없음}: 미스 뒤에 공은 그냥 튕겨 계속 날아가고, 공 위치 자극은 이어진다.
\item \emph{휴지}: 게임은 진행되고 라켓은 활동에 따라 움직이지만 자극은 전혀 없다.
\item \emph{대조}: 배지만 있는 칩. HEK293T 세포. 세포 대신 모의실험을 쓴 “컴퓨터 속 배양”.
\end{itemize}
본 실험의 표본 크기: 생쥐 배양 $9$개(세션 $101$번), 인간 배양 $11$개($138$), 배지만 있는 칩 $6$개($80$), 휴지 배양 $20$개($42$), 모의실험 $3$개($38$), 모두 $399$세션. 되먹임 실험에서는 $3$일 동안 하루 $3$세션씩 $486$세션, 조건은 “자극”, “침묵”, “되먹임 없음”($n=27$, $17$, $15$). 처음 $5$분에서 마지막 $15$분으로 평균 랠리 길이가 늘어난 것은 살아 있는 배양에서만이다. 되먹임 실험에서 시간에 따른 유의미한 향상은 “자극” 조건에서만 있었다. 세션 끝에 “침묵” 조건은 “되먹임 없음”보다 여전히 나았지만 효과가 작았고, 3일 동안 세션 간의 안정적인 학습은 없었다~\cite{Kagan2022}.

\subsection*{테스트 벤치의 주기, 단계별로}

테스트 벤치 컴퓨터는 $10\ms$마다 다음을 한다.
\begin{enumerate}
\item 지난 창 동안 두 운동 영역의 스파이크 수 $n_\uparrow$, $n_\downarrow$를 읽고, 각 영역의 휴지 상태 평균 스파이크 수로 정규화한다.
\item 라켓을 $\Delta p=c\,(n_\uparrow-n_\downarrow)$만큼 옮기고 코트 가장자리 안으로 제한한다.
\item 공을 한 걸음 옮기고, 위쪽, 아래쪽, 왼쪽 벽에서 튕긴다.
\item 공이 오른쪽 벽에 닿았으면: $|y-p|<0.1$일 때 히트, 랠리 수가 하나 늘고 히트 자극($100\ms$)을 준다. 아니면 미스, 랠리를 기록하고 미스 자극($4$~s)을 준 뒤 $4$~s 쉬고 공이 다시 출발한다.
\item 그렇지 않으면 공의 수직 위치로 전극 $k$를, 라켓 쪽 벽까지의 거리로 빈도 $f$를 골라 다음 창까지 전극 $k$에 빈도 $f$로 $75$~mV 2상 펄스를 준다.
\end{enumerate}
이것만으로 발표된 것과 호환되는 테스트 벤치를 조립하고 이 장의 핵심 질문을 확인할 수 있다. 배양을 가르치는 것이 예측 가능성인가, 아니면 히트와 미스에 대한 응답의 단순한 차이인가? 그러려면 논문의 세 조건에 논문에 없는 네 번째 조건을 더하는 것이 좋다. 미스 뒤에 예측 불가능한 자극과 같은 길이, 같은 에너지의 \emph{예측 가능한} 자극을 주는 것이다. 그래도 배양이 학습한다면 중요한 것은 예측 가능성이 아니라 차이이다.


\section{정리}

\begin{table}[t]
\centering
\footnotesize
\setlength{\tabcolsep}{4pt}
\renewcommand{\arraystretch}{1.15}
\begin{tabular}{>{\raggedright}p{2.1cm}>{\raggedright}p{4.2cm}>{\raggedright}p{1.9cm}>{\raggedright\arraybackslash}p{4.6cm}}
\toprule
테스트 벤치의 부분 & 구조 & 이 책의 장 & 알려진 것 \\
\midrule
입력 & 장소(8개 전극 중 어느 것)와 빈도(초당 4--40회) & \ref{sec:code}, \ref{sec:sensors} & 실험자가 정함 \\
출력 & $10\ms$ 동안 두 영역의 스파이크 수 차이~\eqref{eq:dishreadout} & \ref{sec:glutcomputer}, \ref{sec:debugging} & 실험자가 정함. 잡음은 \eqref{eq:rateerror}에 따름 \\
교사 & 히트 뒤 예측 가능한 자극, 미스 뒤 예측 불가능한 자극. 도파민 없음 & \ref{sec:dofa} & 세션 안의 유의미한 향상은 완전한 되먹임일 때만. 침묵일 때는 효과가 더 작음. 세션 간에는 불안정 \\
메커니즘 & 입력 예측 또는 미스 때 뒤섞기~\eqref{eq:dishdensity} & \ref{sec:dofa}, \ref{sec:recognition} & 밝혀지지 않음. 둘 다 결과를 설명함 \\
네트워크 상태 & 게임 중 임계 상태에 접근 & \ref{sec:gaba} & 게임 성적과의 관련은 측정됨 \\
“지각력” & 저자들의 주장 & --- & 보이지 않음. 반박됨 \\
\bottomrule
\end{tabular}
\caption{엔지니어의 눈으로 본 DishBrain.}
\label{tab:dishbrain}
\end{table}

DishBrain은 가장 단순한 형태의 살아 있는 뉴런 기계이다. 입력은 부호로 주어지고, 출력은 스파이크 수로 읽히며, 교사는 성공과 실패에 대한 응답의 차이이다(표~\ref{tab:dishbrain}). 눈도 근육도 도파민도 없는 네트워크가 아주 조금 게임을 배웠고, 그것만으로도 이 네트워크는 이 책의 생각들을 시험하는 테스트 벤치가 된다. 메커니즘이 예측이든, 뒤섞기든, 제3의 무엇이든, 답은 \ref{sec:debugging}장이 권하는 방식으로 찾게 될 것이다. 마침내 전체가 보이는 기계에서 프로브와 테스트 신호와 부분 끄기로 찾는 것이다.

\section{연습문제}

\begin{exercise}[\lvlA]
\label{ex:22:1}
\emph{스파이크 수를 얼마나 옮겨야 하는가.} 두 영역이 합쳐서 초당 $R_\uparrow+R_\downarrow=2000$개의 스파이크를 내고, 흐름은 푸아송이다. (a)~$R=1000$과 $R=100$일 때 한 영역의 $10\ms$당 스파이크 수의 상대 산포는 얼마인가? (b)~공이 $1$~s 동안 날아가는 사이에 변위~\eqref{eq:dishreadout}가 더해진다. 평균 변위가 산포의 두 배가 되는 가장 작은 $R_\uparrow-R_\downarrow$는 얼마인가?
\end{exercise}

\begin{exercise}[\lvlA]
\label{ex:22:2}
\emph{학습을 보려면 랠리가 몇 번 필요한가.} 랠리는 서로 독립이고 받아 친 비율은 이항 분포를 따른다. 받아 친 비율의 증가가 차이의 표준편차 두 배를 넘으려면 두 기간 각각에 랠리가 몇 번 필요한가? (a)~$0.20$에서 $0.25$로; (b)~모델에서처럼 $0.21$에서 $0.38$로.
\end{exercise}

\begin{exercise}[\lvlB]
\label{ex:22:3}
\emph{\eqref{eq:dishdensity}의 유도.} 설정 $\boldsymbol\psi$는 유한 집합 위를 움직인다. 미스(확률 $P(\boldsymbol\psi)>0$) 때 네트워크는 대칭 확률 $K(\boldsymbol\psi,\boldsymbol\psi')$로 $\boldsymbol\psi'$로 옮겨 가고, 히트 때는 그대로 남는다. (a)~$\rho\propto1/P$가 정상 분포임을 증명하라. (b)~$P=0.9$와 $P=0.1$인 두 설정이 있을 때 미스 비율은 얼마인가?
\end{exercise}

\begin{exercise}[\lvlB]
\label{ex:22:4}
\emph{모델의 흡수 영역.} 라켓은 $p=\min\bigl(1,\max(0,ay+b)\bigr)$에 오고, 공은 높이 $y\sim U(0,1)$에 오며, $|p-y|<h=0.1$이면 히트이다. (a)~$|b|<h$이고 $|a-1+b|<h$이면 네트워크가 모든 공을 받아 친다는 것을 증명하고, 이 영역의 넓이를 구하라. (b)~$a\sim U(-1,1)$, $b\sim U(0,1)$일 때 평균 히트 비율을 수치로 구하라.
\end{exercise}

\begin{exercise}[\lvlB]
\label{ex:22:5}
\emph{모의실험: 설정 주위의 울타리.} \texttt{dishbrain\_model.py}의 사본에서 배양 $200$개에 각각 랠리 $20\,000$번을 주라. 마지막 $500$번의 히트 비율은 얼마인가? 흔들기마다 설정을 직사각형 $a\in[-1,2]$, $b\in[-1,1]$ 안으로 잘라 넣으면 어떻게 되는가? 연습문제~\ref{ex:22:3}을 써서 차이를 설명하라.
\end{exercise}

\begin{exercise}[\lvlB]
\label{ex:22:6}
\emph{입력 부호 설계.} 높이 오차가 $h=0.1$보다 작으면 공을 받아 친다. 네트워크는 $T=0.25$~s 동안 입력을 읽고, 빈도 $r$까지 약 $\sqrt{rT}$개의 수준을 구별하며, 자극은 초당 $40$회를 넘지 않는다. (a)~전극 여덟 개의 위치 부호로 충분한가? (b)~전극 하나의 빈도로 높이를 전할 수 있는가?
\end{exercise}

\begin{exercise}[\lvlC]
\label{ex:22:7}
\emph{연구: 예측인가 뒤섞기인가?} 모델을 조정 가능한 숫자 $d$개($p=\sum_{i<d}a_iy^i$)로 일반화하라. 히트 비율 $0.5$에 이르기까지 필요한 랠리 수는 뒤섞기일 때와 오차 부호 규칙~\eqref{eq:perturbation}일 때 $d$에 따라 어떻게 늘어나는가? 어떤 테스트 벤치 실험이 두 가설을 가려내겠는가?
\end{exercise}

\appendix
\renewcommand{\chaptername}{\appendixname}
\chapter{뉴런의 유형을 정하지 않는 물질}
\label{sec:othermediators}

\ref{sec:neurontypes}장에서 우리는 뉴런을 주 신경전달물질에 따라 분류해 열 가지 유형을 얻었다(표~\ref{tab:types}). 이제 복잡한 사정을 더할 때이다. 뉴런은 주 신경전달물질만 내보내지 않으며, 뇌에서는 주 신경전달물질 외에 다른 물질도 일한다. 이 물질들은 네트워크가 신호를 전달하고 저장하는 방식을 바꾼다.

주소 버스와 브로드캐스트 버스(\ref{sec:buses}절) 외에 세 번째 버스가 있다. \emph{역방향 채널}이다. 몇몇 물질은 송신자가 아니라 \emph{수신자}가 내보내며, 이 물질은 틈새를 거슬러 송신자의 출력 쪽으로 가서 송신자가 다음번에 신경전달물질을 얼마나 내보낼지를 바꾼다. 통신망으로 치면 송신자가 자기 전송을 조정하는 데 쓰는 수신 확인과 비슷하다. 뉴런이 연결 가중치를 바꿀 때 쓰는 것이 바로 이 채널이다. 이 채널을 통해 송신자의 출력은 수신자의 활동, 즉 \ref{sec:dofa}장 학습 규칙의 두 번째 인자를 알게 된다.

알려진 신경전달물질의 전체 목록은 길다. 100가지가 넘는 물질이 있고, 대부분은 짧은 단백질 사슬인 신경펩타이드이다. 그러나 이들은 모두 몇 개의 부류로 묶이며, 부류 안의 물질은 비슷하게 행동한다. 주 신경전달물질이 되는 일이 없는 물질을 표~\ref{tab:othermediators}에 모았다. 엔지니어가 물을 세 가지 질문, 즉 신호가 어느 버스로 가는지, 빨리 작용하는지, 보통 부호가 무엇인지를 함께 적었다. 이 물질들은 뉴런의 유형을 정하지 않는다. 주 신경전달물질과 함께 방출되거나, 뉴런이 아닌 세포가 방출하거나, 반대 방향으로 가기 때문이다.

\begin{table}[t]
\centering
\footnotesize
\setlength{\tabcolsep}{4pt}
\renewcommand{\arraystretch}{1.12}
\begin{tabular}{>{\raggedright}p{2.25cm}>{\raggedright}p{2.8cm}>{\raggedright}p{1.8cm}>{\raggedright}p{1.5cm}>{\centering}p{0.8cm}>{\raggedright\arraybackslash}p{4.3cm}}
\toprule
부류 & 물질 & 버스 & 속도 & 부호 & 계산에서의 역할 \\
\midrule
아미노산 & D-세린 & 국소 & 느림 & AND & 글루탐산 수용체의 두 번째 열쇠: “AND” 조건 \\
\midrule
모노아민 & 미량 아민 & 브로드캐스트 & 느림 & $\pm$ & 모노아민 채널의 미세 조정 \\
\midrule
\multirow{2}{2.25cm}{퓨린}
 & ATP & 주소 & 빠름과 느림 & $+$ & 공동 신경전달물질. 뉴런과 교세포 사이의 신호 \\
 & 아데노신 & 체적 & 느림 & $-$ & 활동하면 쌓인다: 부하 계수기~\cite{Dunwiddie2001} \\
\midrule
\multirow{8}{2.25cm}{신경펩타이드 (100가지 이상)}
 & 엔케팔린, 엔도르핀, 다이노르핀 & 체적 & 느림 & $-$ & 전달 억제 \\
 & P 물질 & 체적 & 느림 & $+$ & 느린 증폭 \\
 & 신경펩타이드 Y & 체적 & 느림 & $-$ & 느린 억제 \\
 & 소마토스타틴 & 체적 & 느림 & $-$ & 느린 억제, GABA와 함께 나옴 \\
 & 혈관활성 장펩타이드(VIP) & 체적 & 느림 & $+$ & 탈억제: 억제성 뉴런의 억제 \\
 & 콜레시스토키닌 & 체적 & 느림 & $\pm$ & 일부 억제성 뉴런에서 GABA와 함께 나옴 \\
 & 오렉신 & 브로드캐스트 & 느림 & $+$ & 각성 상태의 안정기 \\
 & 옥시토신, 바소프레신 & 브로드캐스트 & 느림 & $\pm$ & 행동의 전역 설정 \\
\midrule
\multirow{2}{2.25cm}{기체}
 & 일산화질소 NO & 역방향, 체적 & 느림 & $\pm$ & 막을 그대로 통과한다. 가중치 변화 때의 역방향 신호~\cite{Garthwaite2008} \\
 & CO, H$_2$S & 체적 & 느림 & $\pm$ & 역할이 거의 연구되지 않음 \\
\midrule
지질 & 엔도카나비노이드 (아난다마이드, 2-AG) & 역방향 & 느림 & $-$ & 수신자가 송신자의 방출을 줄인다: 가중치에 대한 되먹임~\cite{WilsonNicoll2001} \\
\bottomrule
\end{tabular}
\caption{뉴런의 유형을 정하지 않는 물질. \emph{버스}, \emph{속도}, \emph{부호}는 표~\ref{tab:types}와 같다. 여기에 더해, 체적 버스는 신경전달물질이 방출 지점 둘레의 세포 사이 공간으로 퍼지는 것, 역방향은 수신자에서 송신자 쪽으로 거슬러 가는 것, 국소는 방출 지점 가까이에서 작용하는 것이다. “AND”는 허가 신호이다. 그 자체로는 전류를 일으키지 않지만, 이것이 없으면 다른 입력이 작동하지 않는다. 논리 AND 게이트의 두 번째 입력과 같다. 신경펩타이드는 거의 언제나 주 신경전달물질과 함께, 주로 스파이크 빈도가 높을 때 방출된다~\cite{vandenPol2012}.}
\label{tab:othermediators}
\end{table}

\chapter{기호}
\label{sec:notation}

뉴런의 유형은 주 신경전달물질 영어 이름의 첫 네 글자로 표기한다. \nt{GLUT}는 글루탐산, \nt{GABA}는 감마-아미노뷰티르산, \nt{GLYC}는 글리신, \nt{ACET}는 아세틸콜린, \nt{DOPA}는 도파민, \nt{SERO}는 세로토닌, \nt{HIST}는 히스타민, \nt{NORA}는 노르아드레날린, \nt{ADRE}는 아드레날린, \nt{ASPA}는 아스파르트산이다(표~\ref{tab:types}).

문자는 양보다 적어서, 장에 따라 같은 문자가 다른 것을 뜻하기도 한다. 아래 표에서는 뜻마다 행을 따로 두었다. 오른쪽 열은 그 기호를 도입한 식이다.

\begin{center}
\footnotesize
\setlength{\tabcolsep}{4pt}
\renewcommand{\arraystretch}{1.12}
\begin{longtable}{>{\raggedright}p{2.0cm}>{\raggedright}p{9.6cm}>{\raggedright\arraybackslash}p{2.4cm}}
\toprule
기호 & 뜻 & 도입한 곳 \\
\midrule
\endfirsthead
\toprule
기호 & 뜻 & 도입한 곳 \\
\midrule
\endhead
\bottomrule
\endfoot
$a$ & 선형 전달 특성의 기울기, $f(u)=au$ & \eqref{eq:feedback} \\
$a$ & \nt{ACET} 수준: $0$은 읽기, $1$은 쓰기 & \eqref{eq:acetnet} \\
$a_i$, $a$ & 집단의 피로: 리듬 발생기에서; 서파 수면의 시소에서 & \eqref{eq:halfcentre}, \eqref{eq:updown} \\
$a_S$, $a_P$ & 가속과 제동에 대한 심장 박동의 민감도 & \eqref{eq:gasbrake} \\
$A(r)$, $A_0$ & 빈도 $r$일 때 스파이크에 대한 시냅스 응답; 쉬고 난 시냅스의 응답 & \eqref{eq:depression} \\
$A_1$ & 전달물질 한 꾸러미에 대한 수신자의 응답 & \eqref{eq:quantal} \\
$A$ & 수상돌기를 포함한 뉴런의 막 면적 & \eqref{eq:dendriteload} \\
$A_f$, $A_s$ & 다음 움직임까지 보존되는 보정의 비율, 빠른 과정과 느린 과정 & \eqref{eq:tworate} \\
$b_i$ & 페이스메이커 자체의 유입 & \eqref{eq:seroneuron} \\
$b$ & 리듬 발생기나 수면 시소에서 활동이 집단을 얼마나 빨리 지치게 하는가 & \eqref{eq:halfcentre}, \eqref{eq:updown} \\
$b$ & 디지털화 표본 하나당 비트 수 & \eqref{eq:probedata} \\
$B_f$, $B_s$ & 보정이 오차로부터 배우는 세기, 빠른 과정과 느린 과정 & \eqref{eq:tworate} \\
$c$ & 부피 속 전달물질의 농도 & \eqref{eq:seroneuron} \\
$c$ & 이웃 뉴런 잡음의 상관계수 & \eqref{eq:correlated} \\
$c$ & 스파이크 수 차이가 라켓을 움직이는 계수 & \eqref{eq:dishreadout} \\
$c_f$ & $C$ 뉴런의 응답: 모든 위치에 걸친 $s_{f,p}$의 최댓값 & \eqref{eq:scstage} \\
$c_m$ & 막의 비정전용량 & \eqref{eq:dendriteload} \\
$C$ & 노력의 대가: 시간 $\tau_a$ 동안의 행동은 $C/\tau_a$가 든다 & \eqref{eq:vigor} \\
$d'$ & 두 입력의 변별도: 차이를 잡음으로 나눈 것 & \eqref{eq:gainsnr} \\
$d$, $d_j$ & 신호 경로의 지연 & \eqref{eq:glutneuron}, \eqref{eq:vstar} \\
$d$ & 근육과의 되먹임 루프의 지연 & \eqref{eq:delaystable} \\
$D$ & 확산 계수 & \eqref{eq:diffusionlength} \\
$D$ & \nt{GLUT}--\nt{GABA} 네트워크 발화율 식의 분모 & \eqref{eq:isn} \\
$D$ & 지연된 보상을 기다리는 시간 & \eqref{eq:patience} \\
$e$, $e_{ij}$ & 시냅스의 흔적 & \eqref{eq:threefactor} \\
$e$, $e_i$ & 오차 전선으로 오는 출력 또는 움직임의 오차 & \eqref{eq:lms}, \eqref{eq:adaptivefilter}, \eqref{eq:tworate} \\
$E_{\text{spk}}$ & 시냅스의 일을 포함한 스파이크 하나의 에너지 & \eqref{eq:energy} \\
$E$ & \nt{GLUT} 집단의 발화율 & \eqref{eq:ei} \\
$E_E$ & 흥분성 시냅스의 평형 전위 & \eqref{eq:shunt} \\
$f$, $F$ & 전달 특성: 입력의 함수로서의 발화율 & \eqref{eq:ratenet}, \eqref{eq:gain} \\
$f$, $f_0$ & 심박수; 페이스메이커 고유의 리듬 & \eqref{eq:gasbrake} \\
$f_s$ & 전극 하나의 디지털화 빈도 & \eqref{eq:probedata} \\
$F(t)$, $F_1$ & 근육의 힘; 연축 한 번의 힘 & \eqref{eq:forcesum} \\
$F_1$, $F_2$ & 관절을 반대 방향으로 당기는 두 근육의 힘 & \eqref{eq:antagonists} \\
$g$ & 농도에 대한 수용체의 민감도 & \eqref{eq:seroneuron} \\
$g_L$, $g_E$, $g_I$ & 누설, 흥분, 억제의 전도도 & \eqref{eq:shunt} \\
$g$ & \nt{NORA}가 정하는 이득 & \eqref{eq:gain}, \eqref{eq:machine} \\
$g$ & 기억 네트워크에서 \nt{GABA}에 의한 전체 억제 & \eqref{eq:acetnet} \\
$g$ & 위에서 정하는 통증 관문의 이득 & \eqref{eq:gate} \\
$g$ & 호르몬 연쇄 한 단계의 이득 & \eqref{eq:cascade} \\
$g_j$ & 적응 필터 입력에 있는, 고유 시간 상수를 가진 필터 $j$ & \eqref{eq:adaptivefilter} \\
$G$ & 되먹임 루프의 이득 & \eqref{eq:feedback} \\
$G$ & 지연 있는 루프의 보정 속도 & \eqref{eq:delaystable} \\
$h$, $h_I$ & \nt{GLUT} 집단과 \nt{GABA} 집단의 외부 입력 & \eqref{eq:ei}, \eqref{eq:isn} \\
$h(t)$ & 근육의 단일 연축 & \eqref{eq:twitch} \\
$H(t)$ & 수면 압력의 위쪽 문턱값: 여기에 이르면 기계가 잠든다 & \eqref{eq:sleeppressure} \\
$I$, $I_i$ & 뉴런이나 집단으로의 외부 유입 & \eqref{eq:ratenet}, \eqref{eq:wta}, \eqref{eq:updown} \\
$I$ & \nt{GABA} 집단의 발화율 & \eqref{eq:ei} \\
$I$ & 자극의 밝기 또는 세기 & \eqref{eq:photonnoise}, \eqref{eq:nakarushton} \\
$I$ & 막을 충전하는 입력 전류 & \eqref{eq:dendriteload} \\
$k$ & 버스트 안의 스파이크 수 & \eqref{eq:burst} \\
$k$ & 증가분이 잡음의 몇 배를 넘어야 하는가 & \eqref{eq:photonnoise} \\
$k$ & 호르몬 연쇄의 되먹임 계수 & \eqref{eq:cascade} \\
$k_s$ & 통증이 민감화를 끌어올리는 세기 & \eqref{eq:sensitization} \\
$K$ & 결과의 특징 수 & \eqref{eq:vectortd} \\
$K$ & 관절의 강성 & \eqref{eq:antagonists} \\
$K$ & 호르몬 연쇄의 루프 이득, $K=g^3k$ & \eqref{eq:cascade}, \eqref{eq:cascadestable} \\
$K_\varphi$ & 외부 신호에 대한 일주기 발생기의 동조 세기 & \eqref{eq:pll} \\
$L$ & 지연선의 길이 & \eqref{eq:itd} \\
$L$ & 통증 센서에서 오는 전선의 길이 & \eqref{eq:paintime} \\
$L(t)$ & 수면 압력의 아래쪽 문턱값: 여기에 이르면 기계가 깬다 & \eqref{eq:sleeppressure} \\
$M$ & 예고 신호와 보상 사이 인과 관계의 척도 & \eqref{eq:retro} \\
$M_k$ & 특징 $k$의 기댓값 & \eqref{eq:vectortd} \\
$M_{ik}$ & \nt{GABA} 뉴런 $k$로부터의 억제 세기 & \eqref{eq:machine} \\
$M$ & 관절의 회전력 & \eqref{eq:antagonists} \\
$M$ & 혼합기의 입력선 수 & \eqref{eq:cover} \\
$n$, $\bar n$ & 창 안의 스파이크 수; 그 평균 & \eqref{eq:rateerror} \\
$n$ & 문턱 소자의 입력 수 & \eqref{eq:cover} \\
$n_\uparrow$, $n_\downarrow$ & 창 안에서 `위' 영역과 `아래' 영역의 스파이크 수 & \eqref{eq:dishreadout} \\
$N$ & 뉴런 수 & \eqref{eq:correlated}, \eqref{eq:energy} \\
$N$ & 프로브의 전극 수 & \eqref{eq:probedata} \\
$N_s$ & 두 뉴런 사이의 방출 지점 수 & \eqref{eq:quantal} \\
$p$ & 스파이크당 전달물질 방출 확률 & \eqref{eq:burst} \\
$p$ & 기억 네트워크에서 활성 뉴런의 비율 & \eqref{eq:acetnet} \\
$p$ & 보상하도록 학습하는 섭동 & \eqref{eq:tworate} \\
$P$ & 신호에 쓰는 전력 & \eqref{eq:energy}, \eqref{eq:dishpower} \\
$P$ & 저장된 추정값의 불확실성 & \eqref{eq:kalman} \\
$P$ & `제동'의 활동: 심장으로 가는 \nt{ACET} 전선 & \eqref{eq:gasbrake} \\
$P_{\max}$ & 문턱 소자가 두 클래스로 나눌 수 있는 무작위 패턴의 수 & \eqref{eq:cover} \\
$P_{\text{miss}}$ & 미스 확률 & \eqref{eq:dishdensity} \\
$q$ & 세상이 변하는 속도 & \eqref{eq:kalman} \\
$q_k$ & \nt{GABA} 뉴런 $k$의 발화율 & \eqref{eq:machine} \\
$q_0$ & 억제성 중개자의 배경 활동 & \eqref{eq:gate} \\
$q$ & 통증 관문에서 억제성 중개자의 발화율 & \eqref{eq:gate} \\
$q$ & 다음 운동 단위가 앞의 것보다 몇 배 센가 & \eqref{eq:sizeprinciple} \\
$r$, $r_i$ & 뉴런의 스파이크 발화율 & \eqref{eq:ratenet} \\
$r$, $r_t$ & 보상(보상의 흐름) & \eqref{eq:rpe}, \eqref{eq:ctd} \\
$\mathbf r(t)$ & 저수지의 응답 벡터 & \eqref{eq:reservoir} \\
$R$, $\bar R$ & 받은 보상; 그 기댓값 또는 단위 시간당 평균 & \eqref{eq:perturbation}, \eqref{eq:vigor} \\
$R$ & 필터 입력 잡음의 분산 & \eqref{eq:kalman} \\
$R$, $r_0$ & 지연된 보상과 즉각적인 보상 & \eqref{eq:patience} \\
$R_{\max}$ & 최대 전송 속도, bit/s & \eqref{eq:spikeentropy} \\
$r_{\max}$ & 최대 스파이크 발화율 & \eqref{eq:gain}, \eqref{eq:nakarushton} \\
$R_{\text{raw}}$, $R_{\text{spk}}$ & 프로브의 데이터 흐름: 원시 데이터와 스파이크만, bit/s & \eqref{eq:probedata} \\
$s_j$ & 뉴런 $j$ 출력의 부호 & \eqref{eq:dalesign} \\
$s_i$ & 수신자 수용체의 부호 & \eqref{eq:seroneuron} \\
$s$, $\hat s$ & 세상의 상태; 그 추정값 & \eqref{eq:rpe}, \eqref{eq:kalman} \\
$s_{f,p}$ & 위치 $p$에서 특징 $f$에 대한 $S$ 뉴런의 응답 & \eqref{eq:scstage} \\
$S$ & `가속'의 활동: 심장으로 가는 \nt{NORA} 전선 & \eqref{eq:gasbrake} \\
$S$ & 수면 압력 & \eqref{eq:sleeppressure} \\
$t_1$ & 첫 스파이크의 시각 & \eqref{eq:latency} \\
$t_k$ & $k$번째 스파이크의 시각 & \eqref{eq:forcesum} \\
$t_{f,i}$ & 특징 $f$의 견본 & \eqref{eq:scstage} \\
$T$ & 스파이크를 세는 창; 광자를 모으는 시간 & \eqref{eq:rateerror}, \eqref{eq:photonnoise} \\
$T_c$ & 근육 연축의 상승 시간 & \eqref{eq:twitch} \\
$T_{\text{burst}}$ & 배양의 일제 발화 간격 & \eqref{eq:burstinterval} \\
$u$ & 뉴런의 총입력 & \eqref{eq:gain} \\
$u$, $u(t)$ & 뇌의 명령: 적응 필터의 입력에서; 호르몬 연쇄로 & \eqref{eq:adaptivefilter}, \eqref{eq:cascade} \\
$u(t)$ & 저수지의 입력 & \eqref{eq:reservoir} \\
$U$ & 일정한 입력의 수준 & \eqref{eq:latency} \\
$U$ & 스파이크마다 방출되는 전달물질 저장량의 비율 & \eqref{eq:depression} \\
$v$ & 전선 속 신호의 속도; 점이 움직이는 속도 & \eqref{eq:itd}, \eqref{eq:vstar} \\
$V$ & 막 전압 & \eqref{eq:glutneuron}, \eqref{eq:shunt} \\
$V$ & 미래 보상의 기댓값 & \eqref{eq:rpe}, \eqref{eq:ctd} \\
$w$, $w_{ij}$, $W_{ij}$ & 연결 가중치 & \eqref{eq:glutneuron}, \eqref{eq:ratenet}, \eqref{eq:acetnet}, \eqref{eq:lms}, \eqref{eq:adaptivefilter} \\
$\mathbf w_k$ & $C$ 뉴런 $k$의 입력 가중치 & \eqref{eq:traceinvariance} \\
$w_{q\mu}$, $w_{q\nu}$, $w_{z\mu}$ & 통증 관문의 연결 가중치 & \eqref{eq:gate} \\
$W_{\text{out}}$ & 저수지 선형 판독의 가중치 & \eqref{eq:reservoir} \\
$x$, $y$ & 논리 입력과 출력 & \eqref{eq:dualrail}, \eqref{eq:veto} \\
$x$, $y$ & 송신자와 수신자의 활동 & \eqref{eq:hebb} \\
$x$ & 지연선 위 검출기의 위치 & \eqref{eq:itd} \\
$x_i$ & 감각 기관에서 오는 입력 & \eqref{eq:acetnet} \\
$x_p$ & 위치 $p$의 입력 & \eqref{eq:scstage} \\
$\mathbf x$ & $S$ 뉴런의 출력, $C$ 뉴런의 입력 & \eqref{eq:traceinvariance} \\
$x_j$ & 적응 필터의 입력 $j$: 필터 $g_j$를 거친 명령 & \eqref{eq:lms}, \eqref{eq:adaptivefilter} \\
$x_f$, $x_s$ & 빠른 과정과 느린 과정의 보정 & \eqref{eq:tworate} \\
$x_1$, $x_2$, $x_3$ & 호르몬 연쇄 세 단계의 수준; $x_3$은 최종 호르몬 & \eqref{eq:cascade} \\
$x^*$ & 새 사태가 시작되는 전달물질 저장량의 회복 비율 & \eqref{eq:burstinterval} \\
$y$ & 필터의 잡음 섞인 입력 & \eqref{eq:kalman} \\
$y$, $\hat y$ & 센서 신호; 적응 필터의 예측 & \eqref{eq:adaptivefilter} \\
$y_k$, $\bar y_k$ & $C$ 뉴런 $k$의 활동; 그 흔적 & \eqref{eq:traceinvariance} \\
$y(t)$ & 저수지 판독의 출력 & \eqref{eq:reservoir} \\
$z$ & 통증 관문 출력의 발화율 & \eqref{eq:gate} \\
\midrule
$\alpha_i^\pm$ & 양의 오차와 음의 오차의 가중치 & \eqref{eq:asymmetric} \\
$\beta$ & 상호 억제의 세기 & \eqref{eq:wta} \\
$\beta$ & 통증 관문 출력에 대한 억제의 세기 & \eqref{eq:gate} \\
$\gamma$ & 할인 계수 & \eqref{eq:rpe} \\
$\delta(\cdot)$ & 델타 함수: 순간적인 도약 & \eqref{eq:glutneuron} \\
$\delta$, $\delta_t$ & 보상 예측 오차 & \eqref{eq:rpe}, \eqref{eq:ctd} \\
$\delta t$ & 소리가 두 귀에 도착하는 시간 차 & \eqref{eq:itd} \\
$\Delta$, $\Delta_{\max}$ & 스파이크 간격; 동시성 창 & \eqref{eq:window} \\
$\Delta t$ & 수신자가 시간을 구별하는 정밀도 & \eqref{eq:spikeentropy} \\
$\Delta F$ & 다음 운동 단위가 더하는 힘 & \eqref{eq:sizeprinciple} \\
$\Delta I$ & 구별할 수 있는 밝기 증가분 & \eqref{eq:photonnoise} \\
$\Delta p$ & 창 하나 동안의 라켓 이동 & \eqref{eq:dishreadout} \\
$\Delta u$ & 두 집단의 입력 차이 & \eqref{eq:gainsnr} \\
$\Delta V$ & 막 전압을 얼마나 올려야 하는가 & \eqref{eq:dendriteload} \\
$\Delta x$ & 이웃 센서 사이의 거리 & \eqref{eq:vstar} \\
$\varepsilon$ & 입력의 상대적 차이 & \eqref{eq:wtagain} \\
$\eta$ & 학습률 & \eqref{eq:plasticity}, \eqref{eq:lms}, \eqref{eq:traceinvariance} \\
$\eta$ & 증폭기 뒤 네트워크의 잡음 & \eqref{eq:softmax} \\
$\mu$ & 통증 관문으로 가는 촉각 입력 & \eqref{eq:gate} \\
$\nu$ & 통증 입력 & \eqref{eq:gate} \\
$\theta$ & 발화 문턱값 & \eqref{eq:window}, \eqref{eq:scstage} \\
$\kappa$ & 스파이크 하나당 전달물질의 양 & \eqref{eq:seroneuron} \\
$\kappa_i$ & 양의 오차와 음의 오차 가중치의 비 & \eqref{eq:expectile} \\
$\kappa_m$ & 근육의 힘과 강성 사이의 관계 & \eqref{eq:antagonists} \\
$\ell$ & 전달물질이 퍼지는 거리 & \eqref{eq:diffusionlength} \\
$\ell$ & 근육이 관절을 당기는 팔의 길이 & \eqref{eq:antagonists} \\
$\xi$ & 뉴런 출력의 무작위 흔들림 & \eqref{eq:perturbation} \\
$\rho(\boldsymbol\psi)$ & 설정 $\boldsymbol\psi$에 머문 시간의 비율 & \eqref{eq:dishdensity} \\
$\sigma$ & 산포, 잡음의 크기 & \eqref{eq:rateerror}, \eqref{eq:softmax} \\
$\sigma$ & 응답이 최댓값의 절반이 되는 밝기 & \eqref{eq:nakarushton} \\
$\sigma_{\text{pre}}$, $\sigma_{\text{post}}$ & 증폭기 앞과 뒤의 잡음 & \eqref{eq:gainsnr} \\
$\tau$ & 막의 시간 상수 & \eqref{eq:glutneuron} \\
$\tau$ & 호르몬 연쇄 한 단계의 시간 상수 & \eqref{eq:cascade} \\
$\tau_{\text{eff}}$ & 스스로를 흥분시키는 집단의 시간 상수 & \eqref{eq:perfectint} \\
$\tau_c$ & 부피 속 전달물질의 수명 & \eqref{eq:seroneuron} \\
$\tau_E$, $\tau_I$ & \nt{GLUT} 집단과 \nt{GABA} 집단의 시간 상수 & \eqref{eq:ei} \\
$\tau_e$ & 흔적의 수명 & \eqref{eq:threefactor} \\
$\tau_{\text{tr}}$ & $C$ 뉴런 활동 흔적의 수명 & \eqref{eq:traceinvariance} \\
$\tau_s$ & 손상 뒤 민감도가 돌아오는 시간 & \eqref{eq:sensitization} \\
$\tau_r$, $\tau_d$ & 깨어 있을 때 수면 압력이 쌓이는 시간과 잠잘 때 풀리는 시간 & \eqref{eq:sleeppressure} \\
$\tau_{\text{rec}}$ & 전달물질 저장량의 회복 시간 & \eqref{eq:depression} \\
$\tau_\gamma$ & 계획의 시간 지평 & \eqref{eq:ctd} \\
$\tau_v$ & 기댓값이 갱신되는 속도 & \eqref{eq:ramp} \\
$\tau_a$ & 리듬 발생기나 수면 시소의 피로 시간 & \eqref{eq:halfcentre}, \eqref{eq:updown} \\
$\tau_a^*$ & 행동을 수행하는 최적 시간 & \eqref{eq:vigor} \\
$\Phi$ & 학습 규칙의 우변 & \eqref{eq:plasticity} \\
$\Phi$ & 광자 흐름 & \eqref{eq:photonnoise} \\
$\varphi_k$ & 받은 결과의 특징 $k$ & \eqref{eq:vectortd} \\
$\varphi$ & 일주기 발생기와 외부 신호의 위상 차 & \eqref{eq:pll} \\
$\boldsymbol\psi$ & DishBrain 모델에서 네트워크의 설정 & \eqref{eq:dishdensity} \\
$\omega$, $\omega_0$ & 외부 신호(빛)의 주파수; 일주기 발생기의 고유 주파수 & \eqref{eq:pll} \\
\end{longtable}
\end{center}

\chapter{해답}
\label{sec:answers}

각 장 끝 연습문제의 간단한 답이다. 전체 풀이는 교수자용 별도 해설서에 있다.

\answer{ex:01:1}{(a)~\nt{GLUT} $8\cdot10^{10}$개는 지구 인구의 $10$배이다. 브로드캐스트 뉴런 $\approx1.9\cdot10^6$개는 빈(Wien) 크기의 도시이다. (b)~말단은 $\approx8\cdot10^{14}$개이고 그중 브로드캐스트 말단이 $\approx3\cdot10^{11}$개, 비율 $\approx4\cdot10^{-4}$로, 이 뉴런들의 비율($2.2\cdot10^{-5}$)보다 $\approx16$배 크다.}
\answer{ex:01:2}{(a)~$\tau_c\ln(c_0/c_*)\approx4.6\ms$. (b)~$T\ge\tau_c\ln101$이면 선로가 비워진다. 초당 $\approx220$번 대신 $\approx22$번까지이다.}
\answer{ex:01:3}{$V(s_k)=\gamma^{K-1-k}r$; $\delta_{\text{lamp}}=\gamma^Kr=re^{-T/\tau_\gamma}\approx0.60\,r$(어떤 $\Delta$에서도), $\tau_\gamma\approx1.95\,\text{s}$.}
\answer{ex:01:4}{$n^*=93$, $47$, $25$, $12$. $\alpha$가 작으면 $n^*\approx5/\alpha$이다.}
\answer{ex:01:5}{\nt{GLUT}: $1\to10\to10^4\to10^7$, 뉴런 $\approx10^7$개, $4$단계, $4\ms$. \nt{DOPA}: $1\to10\to3.2\cdot10^5$, 뉴런 $\approx3\cdot10^5$개, $3$단계로 $30$배 적다. \nt{DOPA} 뉴런 수와 같은 자릿수이다.}
\answer{ex:01:6}{$\lambda=f_EJ_E-f_IJ_I$에서 $J_I=37.5$; 전류의 비 $f_IJ_I/f_EJ_E=0.94$; 억제를 겨우 $6.7\,\%$만 약화해도 사태($\lambda\ge1$)가 일어난다.}
\answer{ex:01:7}{각 상태에서 지연마다 $n\ge57$번 제시해야 한다. 같은 감소는 예컨대 $T$와 함께 커지는 내부 시간 추정의 부정확성으로도 생길 수 있다.}

\answer{ex:02:1}{길이 $3.2$~mm, 간격 $25\um$인 열: 검출기 $\approx129$개. 폭 $v\,\tau\ln2=1.7$~mm인 띠, 즉 검출기 $\approx69$개가 발화하는데, 열의 절반이 넘는다. 방향은 띠의 중심이다.}
\answer{ex:02:2}{창은 $-\tau\ln\frac{w_2}{\theta-w_1}\le\delta\le\tau\ln\frac{w_1}{\theta-w_2}$, 즉 $[-0.235,\,0.178]\ms$; 중심 $c=\frac\tau2\ln\frac{w_1(\theta-w_1)}{w_2(\theta-w_2)}=-28$~\textmu s. 방향이 강한 귀 쪽으로 $28$~\textmu s 밀린다. 분해능 간격의 거의 세 배이다.}
\answer{ex:02:3}{$s_iw_{ij}s_j\ge0$인 $s_i$가 필요하며, 이런 $s_i$는 순환이 짝수일 때 그리고 그때에만 존재한다. 상호 억제는 귀착된다(지속 진동 없음). \nt{GLUT}--\nt{GABA} 루프는 귀착되지 않는다(진동할 수 있음).}
\answer{ex:02:4}{뉴런 여섯 개 $g_k=[x+y+z\ge k]$와 $\bar g_k=[\bar x+\bar y+\bar z\ge4-k]$, $k=1,2,3$; $C=g_2$, $\bar C=\bar g_2$, $S=[g_1+\bar g_2+g_3\ge2]$, $\bar S=[\bar g_1+g_2+\bar g_3\ge2]$: 뉴런 여덟 개. AND와 OR로는 $12$개.}
\answer{ex:02:5}{$T_{\min}(0.6)=0.718\,\tau$; $I_c=0.225$, 문턱 근처에서 $T_{\min}\propto(I-I_c)^{-1/2}$.}
\answer{ex:02:6}{(a)~고리 $500$개, 뉴런 $5\cdot10^4$개; 산포 $4.5\ms$($0.45\,\%$), 상대 오차 $\propto T^{-1/2}$. (b)~모든 고리에 공통인, 산포 $5\,\%$의 무작위 지연 배율.}
\answer{ex:02:7}{$\tau_F\ln2=0.69\,\text{s}$마다 갱신: 뉴런당 초당 스파이크 $4.3$개 대 $20$개로 $4.6$배 경제적이다(스파이크당 $10^{-10}$~J일 때 $4\cdot10^{-7}$ 대 $2\cdot10^{-6}$~J). 플립플롭은 비트를 잃는다. 커패시터는 침묵이 갱신 뒤 $0.49\,\text{s}$ 이내에 시작되면 비트를 지킨다(확률 $0.71$).}

\answer{ex:03:1}{$g_E=g_L$에서 $35\to17.5\mV$; $g_E=4g_L$에서 $56\to40\mV$이고, 뺄셈이라면 $38.5\mV$일 것이다. 분로는 $g_E$를 $3$으로 나눈다. 창은 절반으로 좁아진다($\tau_{\text{eff}}$가 $2$배 줄어든다).}
\answer{ex:03:2}{$\operatorname{tr}=\frac{w_{EE}-1}{\tau_E}-\frac1{\tau_I}$, $\det=\frac{1-w_{EE}+w_{EI}w_{IE}}{\tau_E\tau_I}$; 경계에서 $\omega=\sqrt{\det}$, $\tau_I=20\ms$, 주기 $73\ms$.}
\answer{ex:03:3}{$d_c=\tau_E\arccos(-1/G)/\sqrt{G^2-1}$, 주기 $2\pi\tau_E/\sqrt{G^2-1}\in(2d_c,4d_c)$. $G=2$: $d_c=12.1\ms$, 주기 $36\ms$; $G=5$: $d_c=3.6\ms$, 주기 $12.8\ms$. 딱 빠른 리듬의 대역이다.}
\answer{ex:03:4}{(a)~올릴 때는 $I_2=\beta I_1=2$에서, 내릴 때는 $I_2=I_1/\beta=0.5$에서 전환된다. 폭 $1.5$의 고리이다. (b)~$\varepsilon$이 한 자릿수 줄 때마다 $\tau\ln10/(\beta-1)=2.3\tau$씩 늘어난다.}
\answer{ex:03:5}{중개자 세 개, $d_k=5$, $10$, $15\ms$: 금지가 $15\ms$를 덮어야 하고, 하나가 $5\ms$씩 덮는다.}
\answer{ex:03:6}{뉴런 $n+M+1=3+4+1=8$개. 두 개로는: \nt{GABA} $=[x+y+z\ge2]$, 출력 $=[x+y+z-2\,\nt{GABA}\ge1]$.}
\answer{ex:03:7}{$\binom84=70<100\le\binom94=126$: 중개자 $9$개(슈페르너 정리). 직접 연결이 없으면 $100$개: $\beta(J-I)$의 계수가 $n$이다.}
\answer{ex:03:8}{$h_c=\frac1{2w_{EE}}+\frac{w_{EI}w_{IE}-w_{EE}}{4w_{EE}^2}=\frac7{18}\approx0.39$; 시뮬레이션에서는 $h=0.38$과 $0.40$ 사이에서 부호가 바뀐다.}

\answer{ex:04:1}{(a)~$\approx1.1\cdot10^{11}$~bit/J. (b)~$\log_2\sqrt{10}\approx1.7$비트 대 $81$비트: $\approx50$배.}
\answer{ex:04:2}{$r^*=(P_0/E_{\text{spk}})\ln\bigl(1/(r^*\Delta t)\bigr)$, $P_0/E_{\text{spk}}=1.16\,\text{s}^{-1}$: $r^*\approx6$~spike/s, $7.4\cdot10^{10}$~bit/J.}
\answer{ex:04:3}{$\mathcal I=\frac{N}{1+(N-1)c}=9.9$(한계 $10$) 대 $\mathcal I=\frac{N}{1-c}=1111$: $\approx110$배. 상관은 신호를 닮았을 때만 해롭다.}
\answer{ex:04:4}{$\theta\approx68.7$과 $19.9$($w$ 단위), $m=19$와 $m=10$: 빠른 수신자는 평균 입력이 다섯 배 작은데도 일치 입력이 거의 절반만 필요하다.}
\answer{ex:04:5}{$U\tau_{\text{rec}}\ge0.725\,\text{s}$이고 $\tau_{\text{rec}}(1-2U)\le0.05\,\text{s}$이므로 $U\ge0.483$이어야 한다. 가장 작은 $\tau_{\text{rec}}=0.725\,\text{s}$는 $U=1$일 때이다($U=0.5$이면 $1.45\,\text{s}$).}
\answer{ex:04:6}{(a)~$\theta\,e/(e-1)\approx1.58\,\theta$, $\tau$와 $\theta$에 상관없다. (b)~스파이크 $14$개.}
\answer{ex:04:7}{(a)~낱말당 $1.62$비트: $0.77$은 스파이크 수, $0.84$는 그 배치. (b)~$1.13$비트와 $1.59$비트: 낱말 $30$개로는 추정치가 거의 3분의 1을 잃는다.}

\answer{ex:05:1}{$\ell\approx17\um$; $\approx100$일. 뇌 전체로의 배포는 전선의 분지가 맡고, 확산은 각 방출 지점을 $\sim20\um$ 정도로 번지게 할 뿐이다.}
\answer{ex:05:2}{$r=ab/(1+G)$, $S=1/(1+G)$는 정확하다. $+20\,\%$가 아니라 $+1.7\,\%$.}
\answer{ex:05:3}{$G=2$에서 $T^*=1.21\,\text{s}$, $G=9$에서 $0.19\,\text{s}$. 현실적으로는 $G\sim2$--$4$, 억제는 $3$--$5$배.}
\answer{ex:05:4}{전체 발화율 $\lambda$에 대해 $\mathrm{CV}=1/\sqrt{2\lambda\tau_c}\approx9\cdot10^{-4}$. 뉴런 하나라면 $\tau_c=12.5\,\text{s}$가 필요하다.}
\answer{ex:05:5}{두 경우 모두 $c\approx0.40$; $\mathrm{CV}=0.050$ 대 $0.158$, $\sqrt{10}$배. 개별 뉴런의 발화율은 여전히 $a_i$에 비례한다. 조절되는 것은 합뿐이다.}
\answer{ex:05:6}{$N_1=\overline{A\vee B}$, $N_2=\overline{A\vee N_1}$, $N_3=\overline{B\vee N_1}$, $N_4=\overline{N_2\vee N_3}$, XOR $=N_5=\overline{N_4}$. 전환 한 번에 $\tau_c\ln2$가 걸리므로 게이트 $4$개를 거치는 경로는 XOR 하나에 $2.8\,\text{s}$.}
\answer{ex:05:7}{(가)~$\tau_\gamma=6/\ln3\approx5.5\,\text{s}$. (나)~$\bar D<4\,\text{s}$ 대 $D<2.2\,\text{s}$: 지연을 예측할 수 없으면 더 참을성이 생긴다.}
\answer{ex:05:8}{$k_2=1/\tau_d=10\,\text{s}^{-1}$, $k_1=1/\tau_d-1/\tau_\gamma=9.5\,\text{s}^{-1}$; $\tau_\gamma=1/(k_2-mk_1)$. $+2.6\,\%$에서 두 배가 된다($+5.3\,\%$에서는 지평이 무한대).}

\answer{ex:06:1}{$e=(1-e^{-T_a/\tau_e})e^{-d/\tau_e}$. (가)~$0.110$ 대 $4.5\cdot10^{-5}$, 약 $2.5\cdot10^3$배. (나)~$\tau_e^*=T_a/\ln(1+T_a/d)=0.59\,\text{s}$.}
\answer{ex:06:2}{속도는 초당 $\nu_x\nu_y(A_+\tau_+-A_-\tau_-)=0.8A_+$, 한 시간에 약 $2900A_+$; $\tau_-=\tau_+/0.6=33\ms$.}
\answer{ex:06:3}{$\mu^2+s_1^2>s_2^2$이면 $\mathbf e_1$. 여기서는 $1.01>0.25$이므로, $\mathbf e_2$ 방향의 산포가 $25$배 큰데도 뉴런은 $\mathbf e_1$을 배운다. 이 규칙은 공분산행렬이 아니라 상관행렬의 주 벡터를 찾는다.}
\answer{ex:06:4}{전등과 주스 사이에서 $V=Re^{-(t_c+L-t)/\tau_\gamma}$이므로 $\dot V=V/\tau_\gamma$. 급증의 넓이는 $Re^{-L/\tau_\gamma}$: $0.61R$과 $0.78R$.}
\answer{ex:06:5}{신호 대 잡음비는 $1/(N+1)$, 시도 횟수는 $\propto N+1$: $5\cdot10^5$회 $\approx1$개월, $5\cdot10^{11}$회 $\approx8\cdot10^4$년.}
\answer{ex:06:6}{(가)~$k_2=1/\tau_d$, $k_1=1/\tau_d-1/\tau_\gamma$. (나)~거짓 오차는 $-(V/\tau_\gamma)\,\varepsilon/(1+\varepsilon)$, $\varepsilon=\tau_d/\tau_\gamma$이므로 $\tau_d\le0.105\,\text{s}$.}
\answer{ex:06:7}{$0.46$, $0.90$, $0.99$, $0.95$, 그리고 $d=3\,\text{s}$에서 $0.70$. 점들은 흔적 비율 $F=(1-e^{-T_{\text{cue}}/\tau_e})e^{-d/\tau_e}$의 한 곡선 위에 놓인다. $F\gtrsim0.1$이면 학습되고, $F<10^{-2}$이면 무작위 선택이다.}
\answer{ex:06:8}{$\eta^*=1/\bigl(2\sigma^2(N+2)\bigr)$, $N+2$회 시도. $N$개 중 $m$개만 흔들리면 $N(m+2)/m\ge N+2$회: 희소성은 도움이 되지 않는다.}

\answer{ex:07:1}{$V=\kappa p/\bigl(\kappa p+(1-\kappa)(1-p)\bigr)$: $0.059$, $0.20$, $0.50$. 중앙값은 $0$이지만 점~\eqref{eq:expectile}는 $\kappa$에 따라 매끄럽게 변한다.}
\answer{ex:07:2}{$p=1/3$, $x=0.75$, $\sigma_r=0.35$, $V(0.2)=0.083$.}
\answer{ex:07:3}{$V=\sqrt\kappa/(\sqrt\kappa+\sqrt{1-\kappa})$, $0.25$에서 $0.75$까지. 흩어짐 $\propto\sqrt\eta$. 두 방울이면 $V=0.25+0.5\kappa$이고, 두 분포는 양 끝 뉴런에서 $0.05$ 차이 나므로 $\eta\lesssim0.01$이 필요하다.}
\answer{ex:07:4}{(a)~$J^*=2\sqrt{C\bar R}$. (b)~초과 비용 $\bigl(k^{1/2}+k^{-1/2}\bigr)/2-1$: $k=2$에서 $6\,\%$, $k=4$에서 $25\,\%$. “두 배 이내”의 정확도면 충분하다.}
\answer{ex:07:5}{급증의 면적은 $R(e^{-1}-e^{-2})\approx0.23R$로, 램프 $2.3\,\text{s}$ 분량이다. 도파민이 $V$를 알린다면 급증은 없고, 수준이 계단처럼 $e$배로 뛴다.}
\answer{ex:07:6}{사건은 가중치를 $\propto A\tau_f/(\tau_e+\tau_f)$만큼 옮기고, 수준은 오차 $s\tau_f$를 준다: $9\,\text{s}\le\tau_f\le17\,\text{s}$.}
\answer{ex:07:7}{처음에는 $\Delta P=0.80$, $M=0.90$. (a)~$\Delta P=0.74$($-8\,\%$), $M=0.43$($-52\,\%$). (b)~$\Delta P=0.40$($-50\,\%$), $M=0.80$($-11\,\%$). 이중 해리다.}
\answer{ex:07:8}{$N_{\text{eff}}\to2+\rho^2N\approx10^6$번의 시도. 전선 하나일 때보다 $10^4$배 적지만, $1\,\%$의 누설만으로도 여전히 백만 번의 시도가 든다.}

\answer{ex:08:1}{(가)~$\approx1.7\cdot10^6$. (나)~$g=\frac{0.9}{\sqrt{0.19}}\,\sigma_{\text{후}}/\sigma_{\text{전}}\approx16.5$: 답은 잡음의 비에만 달려 있다.}
\answer{ex:08:2}{(가)~고유값은 $-(1\pm g\beta)/\tau$이고, 차는 $\tau/(1-g\beta)$에 걸쳐 감쇠한다. $g\beta=0.5$에서 $40\ms$(입력 차가 세 배로 증폭), $0.9$에서 $200\ms$($19$배). (나)~$0.905$와 $1.105$. 두 경계 사이에서는 강한 입력만 이긴다.}
\answer{ex:08:3}{$P(k)=e^{gu_k/\sigma}/\sum_le^{gu_l/\sigma}$, 즉 식~\eqref{eq:softmax}; $g=10\ln9\approx22$.}
\answer{ex:08:4}{$\bar\Delta=1/3$, $g^*\approx3\ln(1/h)$: $21$, $14$, $9$. 모델에서는 $16$--$23$, $11$, $8$.}
\answer{ex:08:5}{$g^*$는 $h=0.001$에서 $16$--$23$, $h=0.2$에서 $6$--$8$. 추정 $3\ln(1/h)$는 $h\le0.05$에서 1.5배 이내로 맞는다. $h\gtrsim\eta/10$이면 규칙 $Q\leftarrow Q+\eta(R-Q)$가 탐색으로 얻은 것을 미처 흡수하지 못해 $g^*$가 $8$ 근처에 머문다.}
\answer{ex:08:6}{$g_{\text{lo}}=0$에서 $T_p=\tau\ln9=44\ms$(모델 $44\ms$), $g_{\text{lo}}=0.25$에서 $70\ms$. 이득을 0 가까이로 낮춘 $2$--$3\tau$ 길이의 버스트면 된다.}
\answer{ex:08:7}{(가)~최선의 고정값 $0.925$($g=8$), 오라클 $0.929$($16/8$). 얻을 것이 거의 없다. (나)~예를 들어 보상이 드문($p\in[0,0.3]$) 조용한 시기와 $h=0.1$인 변덕스러운 시기: $0.850$ 대 $0.872$. 조정은 $\eta=0.2$에서 약 4분의 1, $\eta=0.05$에서 거의 전부를 가져간다.}

\answer{ex:09:1}{(a)~$P_\infty=0.2$, 기억 $20\,\text{s}$. (b)~$P_\infty=1$, 기억 $1\,\text{s}$. (c)~기억은 잡음 진폭에 비례한다. 여섯 배.}
\answer{ex:09:2}{$P(t)=P_\infty\coth(t\sqrt{q/R})$. (a)~$\frac12\sqrt{R/q}=5\,\text{s}$, 기억의 절반이다. (b)~$t=\frac12\ln21\,\sqrt{R/q}\approx15\,\text{s}$. 높아진 \nt{ACET}가 이만큼 유지되어야 한다.}
\answer{ex:09:3}{분산 $qT/2+R/(2T)$; $T^*=\sqrt{R/q}$, 분산 $\sqrt{qR}=P_\infty$로 필터~\eqref{eq:kalman}와 같다. $(\kappa+1/\kappa)/2$배 증가: $1.25$와 $5.05$.}
\answer{ex:09:4}{$a>a_w=1-\theta/(mw)$: $w=0.041$에서 $0.15$, $w=0.045$에서 $0.22$.}
\answer{ex:09:5}{오류는 $M\approx40$--$60$에서 나타나며, $M=80$에서 끼어드는 뉴런 $1.9$개, $M=100$에서 패턴 비율 $0.86$. 다른 패턴에서 오는 간섭의 분산은 $\approx(M-1)p(1-p)/N$이므로 $M_{\max}\approx80$, $M_{\max}\propto N/p$.}
\answer{ex:09:6}{예를 들어 $\eta(a)=\eta\min\bigl(1,\max(0,(a-0.3)/0.6)\bigr)$. $\eta a$에서는 남의 뉴런 세 개가 모두 패턴에 들어갔고, $\eta(a)$에서는 하나도 들어가지 않았다. $a\ge0.9$에서의 쓰기는 같다(비율 $1.00$).}
\answer{ex:09:7}{(a)~재생 $40$번; 연습문제~\ref{ex:09:6}의 $\eta(a)$로는 영원히 불가능. (b)~재생 $200$번, $10$밤.}

\answer{ex:10:1}{버스 $\approx10^{12}$ bit/s(스파이크당 $11.4$비트), 제어 $\approx40$ bit/s; $2.5\cdot10^{10}\,\text{s}$, 약 800년.}
\answer{ex:10:2}{$a>1/3$($a=0$이면 커지는 진동 $\approx67$~Hz). 아니다: 안정성은 곱 $g(1-a)$가 정한다.}
\answer{ex:10:3}{평균 $\sigma^2R'x_j$, 분산 $(N+1)\,x_j^2\sigma^4R'^2$; 비는 $1/(N+1)$, 필요한 시도는 $n\approx z^2(N+1)\propto N$.}
\answer{ex:10:4}{$L_v\in[32.9,\,47.9]\ms$: 지연 $35.4\ms$, 창 $\pm7.5\ms$; 거짓 묶기는 초당 $2\Delta_{\max}r_ar_v\approx0.38$번.}
\answer{ex:10:5}{$0.098$ 대 가장 좋은 $\alpha=0.2$일 때의 $0.180$: $45\,\%$ 작다.}
\answer{ex:10:6}{$n=\ln\bigl(\ln1.5/(0.6g)\bigr)/\ln(1-\alpha)$: $24$와 $4$; $g=2$이면 $11$과 $2$.}
\answer{ex:10:7}{$S'=S\bigl(1-4\eta\sigma^2+4\eta^2\sigma^4(G+2)\bigr)$, 가장 좋은 $\eta\sigma^2=1/(2(G+2))$; 시도 $\approx10^6$번, 약 한 달.}

\answer{ex:11:1}{(a)~$T=k^2/\bigl(\Phi(\Delta I/I)^2\bigr)=9\,\text{s}$. (b)~센서 $45$개, 점 $\approx7\times7$, 해상도는 $\approx7$배 나빠진다.}
\answer{ex:11:2}{(a)~스파이크당 $2$--$3.3$비트, 한계($9.1$--$9.8$비트)의 $22$--$34\,\%$. (b)~$\log_2\binom{400}{20}\approx111$비트; 공유 유형은 평균 $1$개.}
\answer{ex:11:3}{(a)~$\sigma/9\le I\le9\sigma$, $1.9$자릿수. (b)~$6$과 $7$.}
\answer{ex:11:4}{(a)~$f<343/0.44\approx780$~Hz. (b)~스파이크 $625$개, 뉴런 $125$개.}
\answer{ex:11:5}{이벤트당 $21$비트, 초당 이벤트 $4.8\cdot10^5$개, 가장자리 픽셀당 이벤트 $7$개: 초당 $136$픽셀. 일반 카메라는 초당 $1.25$프레임.}
\answer{ex:11:6}{로그 법칙은 양방향 모두 $\approx1.0\,\text{s}$(이론 $0.96\tau$); 선형 법칙은 위로 $0.2\,\text{s}$, 아래로 $3.0\,\text{s}$(이론 $0.18\tau$와 $3.0\tau$).}
\answer{ex:11:7}{(a)~$\ln I$는 중앙값 $\ln\sigma$, 척도 $1$인 로지스틱 분포: 퍼짐은 자연로그로 $\pi/\sqrt3$, $0.79$자릿수. (b)~$r=r_{\max}\Phi_I(I)^2$.}

\answer{ex:12:1}{$0.60$(상관이 없으면 $0.18$), 극한 $\sqrt{c/(rT)}\approx0.58$: 뉴런 백 개로는 “있다/없다”만 구별한다.}
\answer{ex:12:2}{(a)~스파이크 $10^8$개, $\approx10\,\text{mJ}$: 뇌 전체 $3$~J의 $0.3\,\%$. (b)~$3$~J, $\approx300$배.}
\answer{ex:12:3}{(a)~$\theta_A\in[2,3)$, $\theta_B\in[1,2)$; 다른 도형은 최대 $2$와 $1$을 준다. (b)~$10\cdot96^2\approx9.2\cdot10^4$.}
\answer{ex:12:4}{(a)~$T_d=0.85\,\tau_a=0.34\,\text{s}$. (b)~$T_d\propto\tau_a$이고 $I$와 무관하다: $\tau_a\approx3.5\,\text{s}$(시뮬레이션: $3.1\,\text{s}$).}
\answer{ex:12:5}{$p\approx0.17$. 비율은 천천히 떨어진다: $N=8$, $12$, $24$, $48$, $96$, $192$에서 $0.90$, $0.88$, $0.86$, $0.84$, $0.82$, $0.79$.}
\answer{ex:12:6}{$\tau_{\text{tr}}=20$, $50$, $200\ms$에서 열 번 모두 오류 없음($30\ms$에서는 열 번 중 한 번 실패); $5$--$10\ms$에서는 약 $0.5$, $0.5$--$2\,\text{s}$에서는 품질이 $0.86$까지 떨어진다. 상한: 흔적이 휴지 동안 사라져야 한다, $e^{-0.3\,\text{s}/\tau_{\text{tr}}}\ll1$.}
\answer{ex:12:7}{지연 하나로는 부족하다: 창 $\pm5\ms$로는 $5$에서 $20\ms$까지의 $\Delta x/v$를 덮지 못한다. $5<d_1\le10\ms$와 $15\le d_2\le d_1+10\ms$인 억제 가지 두 개.}
\answer{ex:12:8}{\texttt{two\_stage}: $\Delta=2$, 한 번 뒤집으면 동점, 두 번부터 답이 바뀐다. 1의 계수기: 한 번 뒤집기; $p=0.1$에서 $0.57$.}

\answer{ex:13:1}{(가)~$q=100^{1/99}\approx1.048$, 단계 $\approx4.8\%$; 범위 $\approx2.2\cdot10^3$($67$~dB). (나)~단계 $22f_1$, 곧 $10f_1$에서 $220\%$.}
\answer{ex:13:2}{$P_{\text{fail}}\approx1.8\cdot10^{-4}$와 $2.8\cdot10^{-16}$: $S=2$에서 하루 약 $300$번, $S=4$에서 하루 $5\cdot10^{-10}$번 실패.}
\answer{ex:13:3}{(가)~$H(s)=eF_1T_c/(1+sT_c)^2$, 차단 주파수 $1/(2\pi T_c)\approx3.2$~Hz인 2차 필터. (나)~$r\approx22$~Hz(첫째 고조파로는 $\approx20$).}
\answer{ex:13:4}{(가)~$Gd=1/e$이면 근 $s=-1/d$가 중근이 되고, 다른 어떤 $G$에서도 가장 오른쪽 근이 이보다 왼쪽에 있지 않다. (나)~$d=30\ms$: $G\approx12$~s$^{-1}$, $\tau=30\ms$; $d=150\ms$: $G\approx2.5$~s$^{-1}$, $\tau=150\ms$. 시간 상수는 지연과 같다.}
\answer{ex:13:5}{(가)~$E\approx0.427$, $T\approx1.70\,\tau_a=0.68$~s($\tau=20\ms$인 모델은 $0.84$~s). (나)~방정식에 $I$가 없다; 리듬은 $b>\beta-1$일 때 존재한다.}
\answer{ex:13:6}{$b=0.8$에서는 리듬이 없다; $b=1.05$, $1.2$, $1.5$, $2$, $4$에서 $T\approx2.73$, $1.74$, $1.19$, $0.84$, $0.45$~s; 어떤 $I$에서도 $T=0.840$~s: 입력은 템포가 아니라 진폭만 바꾼다.}
\answer{ex:13:7}{치환 $s=u-a$로 $G=e^{-1-ad}/d$에서 최대 감쇠 속도 $a+1/d$를 얻는다; $a\ge33.3$~s$^{-1}$, $G=e^{-2}/d\approx4.5$~s$^{-1}$; 공동 긴장이 강할수록($a$가 클수록) $G$를 더 작게 잡아야 한다.}
\answer{ex:13:8}{열린 문제. 예를 들어 피로 입력에 문턱값 $b[r_i-r_0]_+$를 두면 $b=4$, $r_0=0.2$에서 $I_{\min}=\beta b r_0/(b-\beta+1)\approx0.53$이고, 주기는 $I=0.6$에서 $1.8$~s부터 $I=3$에서 $0.52$~s까지이다.}

\answer{ex:14:1}{(a)~$0.17$~s와 $1.5$~s. (b)~$11$~cm보다 짧은 거리부터.}
\answer{ex:14:2}{$g\le\beta w_{q\mu}/w_{z\mu}=5$인 동안 통증 없는 촉각은 어떤 $\mu$에서도 통과하지 못한다; 촉각 $\mu=1$은 $g>6$에서 통과한다: 강한 감작은 보통의 접촉을 통증으로 만든다.}
\answer{ex:14:3}{(a)~$g^*=(1-k_sC)/(1-k_sA)$, $k_sA<1$이면 안정. (b)~촉각 없이는 $\nu\ge2$에서, 촉각이 있으면 이미 $\nu\ge1.7$에서.}
\answer{ex:14:4}{$V(t)=-Pe^{-(T-t)/\tau_\gamma}$. (a)~$-Pe^{-T/\tau_\gamma}\approx-0.37P$. (b)~$+Pe^{-(T-t_e)/\tau_\gamma}\approx+0.61P$; $t_e$와 함께 $+P$까지 커지며, 통증을 피하도록 가르친다.}
\answer{ex:14:5}{$k_s=4$에서는 걸림이 없다(둘째 안정 상태는 $k_s\ge4.58$에서 존재); $k_s=4.6$, $5$, $6$, $8$에서 $T_{\min}\approx4.35$, $1.40$, $0.65$, $0.32\,\tau_s$.}
\answer{ex:14:6}{$w_{q\nu}=4/9\approx0.444$, $q_0=2/9\approx0.222$, $w_{q\mu}=49/45\approx1.089$; 통증 없는 촉각은 $g=49/9\approx5.4$까지 안전하다.}
\answer{ex:14:7}{(a)~문턱값은 $\nu_c\ge q_0/w_{q\nu}$이면 $\nu_c(g)=\theta/g$, 아니면 $(\theta+\beta q_0)/(g+\beta w_{q\nu})$; $g^*\approx2.45$, $1.0$, $0.25$. (b)~열린 문제.}

\answer{ex:15:1}{(a)~뉴런당 $2\cdot10^5$, 회로 전체 $6\cdot10^{12}$. (b)~전선당 $\approx8$~bit/s, 최소 $2.5\cdot10^4$~s $\approx7$시간.}
\answer{ex:15:2}{$2$배 간격: 이웃 상관 $0.943$, $\lambda$는 $4.18$에서 $7.3\cdot10^{-4}$까지, 비 $\approx5.7\cdot10^3$; $4$배 간격: 상관 $0.8$, 비 $\approx94$.}
\answer{ex:15:3}{(a)~$\lambda=0.988$과 $0.808$: 움직임 $82$번과 $4.7$번. (b)~$x_s^*=0.690$, $x_f^*=0.172$, 합 $0.862$.}
\answer{ex:15:4}{(a)~움직임 $116$번. (b)~$19$번. 과정이 하나인 모델은 합이 0이면 절약이 없다.}
\answer{ex:15:5}{(a)~모델 없는 한계 $10.5$~s$^{-1}$에 비해 $G=50$~s$^{-1}$. (b)~아니다: $G=50$~s$^{-1}$에서 임계 지연은 $d_c\approx103\ms$; $d=150\ms$에서는 $\hat k>0.445k$가 필요하다(어떤 $G$와 $d$에서든 $\hat k>k/2$).}
\answer{ex:15:6}{오차는 $6$--$30$~s 만에 $0.5\%$ 아래로 떨어지고, 바닥은 $\approx(6$--$9)\cdot10^{-4}$, 최적 가중치에서는 $7.5\cdot10^{-4}$; $\eta=50$의 가중치는 $C$의 조건이 나쁜 방향을 따라 아직 최적과 최대 $0.2$만큼 다르다(연습문제~\ref{ex:15:2}).}
\answer{ex:15:7}{$\mathbf B=A\mathbf K$; $q_f/R\approx0.035$, $q_s/R\approx8.6\cdot10^{-4}$; $4q_s$에서 $B_f\approx0.088$, $B_s\approx0.047$: 느린 과정은 두 배 넘게 빨리 배우고, 빠른 과정은 더 느려진다.}

\answer{ex:16:1}{혈액은 $\tau=t_{1/2}/\ln2=14.4$~min인 RC 회로다. 최댓값 대 최솟값은 $2^{T/t_{1/2}}=64$; 기본 고조파는 $0.55$로 감쇠한다. 비가 $2$가 되는 것은 $t_{1/2}=T=60$~min일 때다.}
\answer{ex:16:2}{$\varphi^*=\arcsin\bigl((\omega-\omega_0)/K_\varphi\bigr)\approx41$~min ($\omega-\omega_0\approx0.047$~rad/day), 이완 시간 $1/(K_\varphi\cos\varphi^*)\approx3.9$일; $+6$~h와 $-6$~h 이동 뒤에는 각각 $7.7$일과 $7.9$일.}
\answer{ex:16:3}{$K_c(n)=\sec^n(\pi/n)$: $8$과 $2.37$, 오차 $11\%$와 $30\%$; $n\to\infty$에서 $K_c\to1$, 오차 $\to50\%$.}
\answer{ex:16:4}{극단 명령(가속 페달 $80$, 브레이크 $0$)으로 $f=120$에 $0.76$~s 만에 도달한다; “가속 페달만”이면 $2.33$~s로 세 배 걸린다.}
\answer{ex:16:5}{$3\arctan(\omega\tau)+\omega d=\pi$, $K_c=(1+\omega^2\tau^2)^{3/2}$: $K_c=8$, $3.54$, $2.50$, $1.76$, 주기 $3.63$, $5.46$, $6.86$, $9.27\,\tau$; $0.9K_c$에서는 진동이 감쇠하고 $1.1K_c$에서는 일정한 진동이 자리 잡는다.}
\answer{ex:16:6}{$0.60\bar c$, $0.87\bar c$, $0.90\bar c$, 극한 $0.91\bar c$; 농도가 일정하면 반응은 0이다. 이런 수신자는 분량만 읽는다.}
\answer{ex:16:7}{열린 문제. 되먹임이라면 리듬은 $K>8$에서만 있고, 주기는 단의 $\tau$에 비례하며 $K$에 거의 의존하지 않는다($K=8.5$에서 $3.64\,\tau$, $K=40$에서 $4.23\,\tau$). $K$를 1.5배 바꾸면 주기는 $2$--$4\%$ 옮겨 가므로 측정 오차보다 작다. 결정적인 것은 루프를 여는 실험(첫 단 입력에 최종 호르몬을 일정하게 주면 이런 리듬은 사라진다)이나, 주기의 여러 위상에서 호르몬을 짧게 추가했을 때의 반응이다.}

\answer{ex:17:1}{$0.2$~W 대 $81$~\textmu W; 원시 흐름에는 $50$~pJ/bit 이하가 필요하다.}
\answer{ex:17:2}{$c=0.1$에서 $5.6$, $8.7$, $9.2$~bit/s, 극한 $9.3$~bit/s; $c=0$, $N=1000$에서 $65$~bit/s.}
\answer{ex:17:3}{$B=\log_2N+P\log_2P+(1-P)\log_2\frac{1-P}{N-1}=4.87$, $4.37$, $3.29$비트/글자, 즉 $7.3$, $6.6$, $4.9$~bit/s; $10$~bit/s에는 초당 $2.3$글자가 필요하다.}
\answer{ex:17:4}{성분당 오차 분산은 $2b/(Nm^2T)+\sigma_\xi^2$; 두 기여는 $N\approx90$에서 같다; 극한은 창당 성분마다 $\tfrac12\log_2(1+0.25/0.04)\approx1.4$비트.}
\answer{ex:17:5}{쌍은 $\eta_bd^2+\eta_db^2<2$, 즉 대략 $\eta<1$에서 수렴한다. $0.8$에서 수렴, $1.1$에서 주기 $2$의 지속 진동, $1.5$에서 비주기 진동, $2$에서 폭주한다. 두 학습자의 속도는 떼어 놓아야 한다.}
\answer{ex:17:6}{문턱값 $\approx4.4\sigma$; $102$~kbit/s, $0.10$~mW 대 원시 흐름의 $205$~mW로 $2000$분의 1; 포화되는(스파이크 $3$개 초과) 창은 $5.7\cdot10^{-5}$뿐이다.}
\answer{ex:17:7}{열린 문제. $\text{SNR}\approx0.9$에서 약 $46$~bit/s, 독립 잡음에서 약 $330$~bit/s. 신호를 닮은 잡음이 원인이라면 꺼낸 정보는 채널 수가 늘어도 포화된다. 부호를 모르는 것이 원인이라면 더 나은 디코더(예: 스파이크 타이밍을 쓰는 디코더)로 정보가 늘어난다.}

\answer{ex:18:1}{(a)~$0.9949\le w\le1.0049$, 즉 $|1-w|\lesssim0.005$. (b)~시냅스 $K\ge400$개.}
\answer{ex:18:2}{$(\omega_R\tau_w)^2=\sqrt{\alpha(\alpha+2+2\varrho)}/\varrho-1$, $\varrho=\tau_m/\tau_w$; $f_R\approx11.1$~Hz; $|Z(\omega_R)|/|Z(0)|=4.6$.}
\answer{ex:18:3}{(a)~$f=1/\bigl(\tau\ln\frac{\mu}{\mu-\theta}\bigr)$; $1$~Hz에는 $\mu/\theta-1\approx3.7\cdot10^{-44}$이 필요하다. (b)~$T=\pi\tau/\sqrt\mu$, $f\approx3.2$~Hz: $f\propto\sqrt\mu$.}
\answer{ex:18:4}{적분기는 $f\lesssim17.7$~Hz에서 발화하고(저역 통과 필터), 공진기는 $5.5$--$22$~Hz 대역에서만 발화한다(대역 통과 필터).}
\answer{ex:18:5}{(a)~$T_{\min}=\frac{\tau}{w-1}\ln\frac{0.5}{0.3}\approx10.2\ms$. (b)~$B>w-1-0.2=0.3$.}
\answer{ex:18:6}{조정의 시상수는 $\tau/(\eta\langle r^2\rangle)$, $\eta=\tau\ln10/60\,\text{s}\approx3.8\cdot10^{-3}$; $I\neq0$이면 가중치가 치우치므로 $e$에서 명령의 사본 $I/\tau$를 빼야 한다.}
\answer{ex:18:7}{열린 문제. 공진기는 $11$~Hz에서 $1.35$배만 유리하고 $1$~Hz에서는 $17$배 불리하다. 재구성의 주된 이득은 문턱값에 의한 대역 선택에 있다(연습문제~\ref{ex:18:4}).}

\answer{ex:19:1}{(a)~시냅스 하나는 $6.7\mV$를 준다($R=66.7$~M$\Omega$). 따라서 문턱값에 아예 도달하려면 최소 $4$개가 필요하다(세 개로는 점근적으로만 도달한다); $10\ms$ 안에는 $8$개; $5\ms$ 안에는 $14$개. (b)~$0.1\ms\ll\tau_m$, 누설에 의한 보정은 $\sim0.25\%$.}
\answer{ex:19:2}{(a)~$1.75\cdot10^{10}$(믹서의 4분의 1이 무엇에나 반응한다), $4.4\cdot10^9$, $0.06$($k=40$에서는 거의 반응하지 않는다). (b)~$10^3$배: $2\cdot10^5$ 대 $200$.}
\answer{ex:19:3}{(a)~이항 계수의 대칭성에 의해 $C(2n,n)=2^{2n-1}$. (b)~$0.93$과 $0.09$; 전이 폭은 $P$에 대해 $\sim\sqrt{2n}$, 상대 폭은 $\sim1/\sqrt n$.}
\answer{ex:19:4}{가지돌기 하나에서는 입력이 몇 개든 $V_{\text{세포체}}<\kappa E_E<\theta$; $g_0=0.5g_L$이면 가지돌기마다 $n\ge3$이 필요하다.}
\answer{ex:19:5}{$I\ge C\sigma_V/\sigma_t=15$~nA, 즉 동기 시냅스 $150$개로 비현실적이다. 작은 뉴런은 $1$~nA면 충분하고, $4$~nA 시냅스에서 지터는 $5$~\textmu s.}
\answer{ex:19:6}{가까운 끝의 최댓값: $0.03\,\tau_m$와 $0.97$($L=0.2$); $0.32\,\tau_m$와 $0.66$($L=1$); $0.79\,\tau_m$와 $0.33$($L=2$). 전체 커패시턴스에 의한 추정~\eqref{eq:dendriteload}은 $L\ll1$에서만 맞다.}
\answer{ex:19:7}{열린 문제. $m$이 고정이면 응답 시간은 기억하는 $P$에 선형으로 늘어난다. 비율 $m=\varphi n$이 고정이면 $n$에 의존하지 않게 되며, 큰 뉴런이 느린 것은 크기 자체가 아니라 많은 입력을 합산하기 때문이다.}

\answer{ex:20:1}{$t_{\text{각성}}=\tau_r\ln\frac{1-L}{1-H}=16.8$~h, $t_{\text{수면}}=\tau_d\ln\frac HL=5.8$~h, 주기 $22.5$~h; 발진기를 $24$~h로 끌어오는 것은 문턱값의 일주기 흔들림, 즉 위상 고정 루프~\eqref{eq:pll}다.}
\answer{ex:20:2}{(a)~$L=0.111$; $L=0.092$이면 수면은 $9.6$~h 이어졌을 것이다. (b)~$18\%$가 아니라 $22\%$($1/0.82=1.22$).}
\answer{ex:20:3}{$H=\frac{p-1}{p-1/q}=0.623$, $L=H/q=0.093$, 여기서 $p=e^{16/\tau_r}$, $q=e^{8/\tau_d}$. $4$~h 만에 깨우면 위상이 영구히 옮겨 간다(잠드는 시각이 $7.2$~h 빨라진다). 문턱값이 오르내리면 위상은 2--3일 안에 돌아온다.}
\answer{ex:20:4}{(a)~$r_\pm=\frac12\bigl(1\pm\sqrt{1-4k/w}\bigr)=0.947,\ 0.053$; $a_{\text{상}}=3.51$, $a_{\text{하}}=0.49$. (b)~활동 $\approx0.33$~s, 침묵 $\approx0.59$~s, 주기 $\approx0.93$~s, 활동 비율 $0.36$.}
\answer{ex:20:5}{$a_{\text{상}}/r_+<b<a_{\text{하}}/r_-$: $3.71<b<9.20$. \nt{ACET}이 $b$를 $3.71$ 아래로 낮추면 교점이 위 가지로 옮겨 가고, 활동 상태($r\approx1$)가 안정해진다.}
\answer{ex:20:6}{$D=24$, $32$, $40$, $48$, $64$에서 $4.9$, $6.8$, $8.8$, $5.5$, $8.9$~h: 길이를 정하는 것은 빚이 아니라 잠드는 시각의 일주기 위상이다.}
\answer{ex:20:7}{B 뒤 A에 대한 오차는 평균 $0.51$($20$개 세트에서 $0.19$부터 $1.08$까지; 응답이 0이면 $1$). 섞어서 배우면 두 오차 모두 $10^{-4}$보다 작다.}
\answer{ex:20:8}{열린 문제. 균일한 스케일링은 가중치의 비를 보존하지만, 문턱 뉴런의 응답은 문턱값도 같은 비율로 스케일링될 때만 보존된다. 섞은 재생은 비를 바꾼다. 큰 접촉이 바뀌지 않는다는 사실은 순수하게 균일한 스케일링과 모순된다.}

\answer{ex:21:1}{(a)~$T=\tau_{\text{rec}}\ln\frac{1}{1-x^*}$: $1.84$~s와 $3.68$~s. (b)~$3.36$에서 $4.24$~s까지. 민감도 $\dd T/\dd x^*=\tau_{\text{rec}}/(1-x^*)=80$~s이므로 $x^*\to1$에서 일제 발화는 불규칙해진다.}
\answer{ex:21:2}{(a)~\eqref{eq:energy}에서처럼 $8.6$~W. (b)~$205$~Mbit/s, $0.2$~W로 배양($10^{-4}$~W)의 $2\cdot10^3$배이다.}
\answer{ex:21:3}{$x_{\text{ss}}=1/(1+Ur_b\tau_{\text{rec}})<x^*$이려면 $r_b>(1/x^*-1)/(U\tau_{\text{rec}})$, 즉 $x^*=0.9$에서 $0.28$~Hz, $x^*=0.99$에서 $0.025$~Hz. 시냅스 세기는 $x_{\text{ss}}$로, 즉 $10\%$와 $1\%$만큼 떨어진다.}
\answer{ex:21:4}{$u(t-k)$는 정규직교이고 $m_k=\|P_Vu(t-k)\|^2$이다. 여기서 $P_V$는 출력이 생성하는 $N$차원 공간으로의 사영이다. 베셀 부등식에 의해 $\sum_k m_k\le N$.}
\answer{ex:21:5}{$\tanh$가 있으면 $\mathrm{MC}\approx9$--$11$, 선형 저수지는 $15$--$18$(시드 세 개), 둘 다 $\le30$이다. 비선형성의 대가는 기억이다.}
\answer{ex:21:6}{(a)~$n\le102$. (b)~$5.6$시그마.}
\answer{ex:21:7}{열린 문제. 핵심 대조 실험: 판독을 고정하고, 자극 없는 휴지 뒤에도 향상이 남는지 확인한다. 남지 않으면 바뀐 것은 가중치가 아니라 상태이다.}

\answer{ex:22:1}{(a)~$32\%$와 $100\%$. (b)~$R_\uparrow-R_\downarrow\ge2\sqrt{2000\,\text{Hz}/1\,\text{s}}\approx89$~Hz, 전체 빈도의 $4.5\%$에 불과하다.}
\answer{ex:22:2}{(a)~$n\ge4(p_1q_1+p_2q_2)/\Delta^2\approx560$. (b)~$\approx56$.}
\answer{ex:22:3}{(a)~세부 균형: $\rho PK$가 대칭이다. (b)~미스 비율은 조화평균 $1/\langle1/P\rangle=0.18$로, 되먹임이 없을 때의 $0.5$와 대비된다.}
\answer{ex:22:4}{(a)~넓이 $4h^2=0.04$인 평행사변형($p$의 제한을 고려하면 수치로 $0.045$). (b)~$\approx0.18$.}
\answer{ex:22:5}{$0.49$: 일부 배양은 거의 언제나 미스인 영역으로 빠져나가 거기서 떠돈다. 제한을 두면 $1.00$: 유계 집합 위에서는 밀도 $\propto1/P$가 흡수 영역에 모인다.}
\answer{ex:22:6}{(a)~수준이 $\ge5$개 필요하다. 전극 $8$개면 충분하다. (b)~아니다. $r_{\max}\ge25/T=100$~Hz가 필요하다.}
\answer{ex:22:7}{열린 문제. 무작위 탐색 시간은 $d$에 따라 대략 부피 비 $(L/h)^d$처럼 늘고, 오차 부호를 쓰는 탐색은 $d$에 선형으로 는다. 두 가설을 가르는 것은 뒤바꾸기 실험이다. 미스 뒤에는 예측 가능하지만 강한 자극을, 히트 뒤에는 예측 불가능하지만 약한 자극을 준다. 그룹마다 배양이 수십 개 필요하다.}


\chapter{실험 부록}
\label{sec:supplement}

여기에는 각 장마다 핵심 주장과 그 근거를 모았다. 그것을 측정한 실험, 정확히 무엇을 측정했는지, 어디에 발표되었는지이다. 오른쪽 열의 상태는 주장이 얼마나 단단한지를 말해 준다. \emph{측정됨}은 직접 실험이 있다는 뜻이다. \emph{이론}은 본문이나 고전적 논문에서 유도한 수학적 귀결이다. \emph{모델}은 이 책의 모델로 계산한 결과이다(스크립트는 \texttt{models/} 폴더에 있다). \emph{추정}은 직접 측정 없는 크기 정도의 수치이다. \emph{가설}은 그럴듯하지만 아직 실험으로 증명되지 않은 설명이다. 실험이 없는 곳에는 그렇다고 적었다.

\section{\ref{sec:neurontypes}장. 뉴런의 유형}

이 장의 수치는 사람의 뇌에서 세포를 직접 센 자료가 있는 곳에서는 그것에 기댄다. “뉴런 하나, 신경전달물질 하나” 규칙은 세포 아틀라스와 예외까지 찾아낸 실험에, \nt{DOPA}의 역할은 스파이크 기록에, 나머지 브로드캐스트 채널의 역할은 가설에 기댄다.

{\footnotesize
\setlength{\tabcolsep}{3pt}\renewcommand{\arraystretch}{1.15}
\begin{longtable}{>{\raggedright}p{0.5cm}>{\raggedright}p{4.0cm}>{\raggedright}p{5.6cm}>{\raggedright}p{2.3cm}>{\raggedright\arraybackslash}p{1.7cm}}
\toprule
No. & 명제 & 실험과 측정한 것 & 출처 & 상태 \\
\midrule\endfirsthead
\toprule
No. & 명제 & 실험과 측정한 것 & 출처 & 상태 \\
\midrule\endhead
1 & 사람의 뇌에는 뉴런이 약 $8.6\cdot10^{10}$개 있고, \nt{GLUT} 뉴런의 거의 전부는 소뇌의 작은 뉴런이다(표~\ref{tab:types}) & 등방성 분획법: 조직을 녹여 세포핵의 균일한 현탁액으로 만들고, 뉴런 표지 NeuN을 가진 핵을 센다. 성인 남성의 뉴런은 $86.1\pm8.1$\,십억(billion) 개이다. 그중 대뇌 피질에 있는 것은 $19\,\%$뿐이고 대부분은 소뇌에 있다 & \cite{Azevedo2009} & 측정됨 \\
2 & 거의 모든 뉴런에는 주된 신경전달물질이 하나 있지만(명제~\ref{prop:dale}) 예외도 있다 & 생쥐 뇌의 전사체 아틀라스: 세포 약 $4\cdot10^6$개, 클러스터 $5322$개. 합성과 포장 유전자에 따라 클러스터에 신경전달물질을 배정했다. 예외는 직접 측정되었다. 절편에서 \nt{DOPA} 뉴런은 같은 소포 수송체를 통해 GABA도 내보내 선조체 뉴런을 억제한다. 일부 \nt{DOPA} 축삭에서는 글루탐산과 도파민이 같은 전선의 서로 다른 말단에서 나온다(전자현미경과 광유전학) & \cite{Strata1999,Yao2023,Tritsch2012,Zhang2015} & 측정됨 \\
3 & 가중치 행렬의 열 전체가 같은 부호이다~\eqref{eq:dalesign} & 명제~\ref{prop:dale}와, 글루탐산 수용체는 거의 모두 흥분성이고 GABA 수용체는 억제성이라는 사실에서 장에서 유도한 것이다. “한 뉴런의 모든 수신자가 같은 부호의 가중치를 받는다”를 일반 규칙으로 직접 검증한 연구는 없다. 반례는 \nt{DOPA} 뉴런의 GABA 공동 방출(2행)과 부호가 다른 수용체를 가진 수신자들(11행)이다 & 이 장의 유도 & 이론 \\
4 & 피질 뉴런의 4분의 3에서 5분의 4는 \nt{GLUT}, 5분의 1에서 4분의 1은 \nt{GABA}이다 & 원숭이 피질 $10$개 영역에서 피질 전체 두께에 걸친 폭 $50$\,\textmu m의 기둥을 GABA로 면역염색했다. $8$개 영역에서 GABA 뉴런은 전체 뉴런의 약 $25\,\%$, 일차 시각 피질에서는 $20\,\%$ 미만이다 & \cite{Hendry1987} & 측정됨 \\
5 & \nt{GLUT} 뉴런의 출력은 약 $10^4$개이다 & 부검 입체해석학: 젊은 남성의 신피질에 시냅스 $1.64\cdot10^{14}$개(전자현미경, 디섹터법)와 뉴런 약 $2.3\cdot10^{10}$개. 비를 구하면 뉴런당 시냅스 $\sim7\cdot10^3$개이다. 평균적으로 입력 수는 출력 수와 같다 & \cite{Tang2001synapses,Pakkenberg1997} & 측정됨 \\
6 & \nt{DOPA} 뉴런의 출력은 말단 $10^5$--$10^6$개로 가지를 친다. 이것은 표적을 덮는 범위에서 얻은 추정치이지 센 값이 아니다 & 쥐의 \nt{DOPA} 뉴런 하나하나를 바이러스로 막 표지하고 축삭을 완전히 재구성했다. 뉴런 하나가 선조체 부피의 $0.45$에서 $5.7\,\%$(평균 $2.7\,\%$)를 덮는다. 측정한 것은 덮는 범위이지 말단 수가 아니다. 말단 수는 덮는 범위에서 자릿수 수준으로 유도했고, 말단을 직접 센 자료는 없다. \nt{SERO}와 \nt{NORA}의 말단 수는 거친 추정치이다 & \cite{Matsuda2009} & 측정됨 \\
7 & 브로드캐스트 뉴런은 적다. \nt{DOPA} $\sim5\cdot10^5$(두 무리 중 큰 쪽, 하한), \nt{SERO} $\sim3\cdot10^5$(두 부분 집계의 합), \nt{HIST} $\sim6\cdot10^4$(표~\ref{tab:types}, 그림~\ref{fig:typepie}) & 사람에서의 입체해석학적 집계: 흑색질의 색소 뉴런 $550\,000$개(배쪽 피개 영역 제외); 등쪽 솔기핵의 뉴런 $235\,000$개(핵의 모든 뉴런)와 다리뇌 피개의 세로토닌 뉴런 $125\,000$개; 시상하부의 히스타민 뉴런 약 $64\,000$개. \nt{NORA}, \nt{GLYC}, \nt{ACET}, \nt{ADRE}는 직접 센 자료가 없어 표에는 거친 자릿수 추정치를 넣었다 & \cite{Pakkenberg1991,Baker1990,Baker1991,Airaksinen1991} & 측정됨 \\
8 & 틈 속 글루탐산은 약 1밀리몰까지 치솟았다가 $\sim1\ms$ 만에 줄어든다 & 시냅스 전달 중 빠르게 떨어지는 경쟁적 차단제가 NMDA 수용체에서 밀려나는 정도로 측정했다(해마 뉴런 배양). 정점 $1.1$\,mM, 감소 시간 상수 $1.2\ms$ & \cite{Clements1992} & 측정됨 \\
9 & 작동 중인 피질에서는 흥분 전류와 억제 전류가 거의 균형을 이루며, 뉴런은 문턱값 근처에 머문다 & 이론: 연결이 강한 네트워크는 스스로 균형 체제에 이른다. 실험: 흰담비 전전두 피질 in vivo에서 “업” 상태 동안 시냅스 전류의 역전 전위는 수백 밀리초 동안 약 $-37$\,mV로 유지되는데, 컨덕턴스는 $\sim21$\,nS만큼 변한다. 흥분과 억제가 함께 오르내린다 & \cite{vanVreeswijk1996,Haider2006} & 측정됨 \\
10 & \nt{DOPA} 뉴런은 보상 예측 오차~\eqref{eq:rpe}를 알린다(그림~\ref{fig:rpe}) & 원숭이 \nt{DOPA} 뉴런의 스파이크: 예고 없는 주스에 버스트; 학습 뒤에는 신호에 버스트, 예측된 주스에는 무응답; 약속한 주스를 주지 않으면 발화율이 움푹 꺼진다. 정량 실험에서 발화율은 현재 보상과 과거 보상의 가중 평균의 차에 비례하지만, 양의 오차에서만 그렇다. 식~\eqref{eq:rpe} 자체는 시간차 알고리즘이다 & \cite{Schultz1997,Bayer2005,SuttonBarto2018} & 측정됨 \\
11 & 부호와 속도는 수용체가 정한다. 이온성은 1밀리초 미만, 대사성은 훨씬 느리다(명제~\ref{prop:receptor}) & 해마 뉴런의 막 조각에 글루탐산을 짧게 가했다. 이온성 수용체를 통한 전류는 $0.2$--$0.6\ms$ 만에 ($20$에서 $80\,\%$까지) 오르고 $2$--$3\ms$ 만에 줄어든다. 피라미드 뉴런의 느린 억제 전위는 대사성 GABA 수용체의 선택적 차단제로 사라진다. 도파민 수용체에는 작용이 반대인 두 계열이 있다. 선조체에서 D1 수용체를 가진 뉴런은 시냅스 강화에, D2를 가진 뉴런은 시냅스 약화에 도파민이 필요하다 & \cite{Colquhoun1992,DutarNicoll1988,Beaulieu2011,Shen2008} & 측정됨 \\
12 & 브로드캐스트 채널은 전선 한 가닥이 아니라 적은 수의 선로로 된 다발이다(\ref{sec:buses}절) & 생쥐 등쪽 솔기핵 세로토닌 뉴런의 바이러스-유전적 표지: 편도체로 출력을 보내는 뉴런과 이마엽 피질로 보내는 뉴런은 핵 안의 다른 자리에 있고, 서로 보완하는 영역으로 가지를 치며, 불쾌한 자극에 반대로 응답한다. 총설: 청반(\nt{NORA}) 뉴런의 하위 무리들은 서로 다른 과정을 부호화한다 & \cite{Ren2018,Poe2020} & 측정됨 \\
13 & \nt{NORA}는 이득 $g$를 정한다 & 적응적 이득 가설; 카테콜아민이 전달 특성의 기울기를 높이는 네트워크 모델. “노르아드레날린과 함께 $g$가 커진다”를 네트워크 전체에서 직접 측정한 자료는 없다. 원숭이에서 동공 지름이 청반 뉴런의 스파이크를 따라간다는 것은 측정되었다 & \cite{AstonJones2005,ServanSchreiber1990,Joshi2016} & 가설 \\
14 & \nt{ACET}은 “쓰기/읽기”를 전환하고 학습률을 정한다 & 가설. 근거: 쥐 후각 피질 절편에서 아세틸콜린 작용제($100$\,\textmu M)는 내부의 “회상” 시냅스를 $60\,\%$ 넘게 약화하지만 입력 시냅스는 $15\,\%$ 미만으로 약화한다 & \cite{Hasselmo2006,HasselmoBower1992} & 가설 \\
15 & \nt{SERO}는 할인율 $\gamma$를 정한다; 세로토닌 뉴런은 초당 스파이크 몇 개로 고르게 작동하고 수면의 한 단계에서는 거의 잠잠하다 & “세로토닌 $=\gamma$” 가설. 가장 강력한 근거: 생쥐 등쪽 솔기핵 세로토닌 뉴런을 광유전학으로 활성화하면 지연된 보상을 기다리다 일찍 포기하는 횟수가 줄어든다. 발화율과 급속 안구 운동 수면 중의 침묵은 측정되었다(기록 총설) & \cite{Doya2002,Miyazaki2014,JacobsAzmitia1992} & 가설 \\
\bottomrule
\end{longtable}}

\section{\ref{sec:glutcomputer}장. \nt{GLUT} 뉴런은 무엇을 계산하는가}

이 장의 논리적 명제는 음이 아닌 가중치로 된 회로에 관한 정리로, 장 안에서 유도했다. 실험은 뉴런 모델을 뒷받침하고, 이 정리들에 따라 조립한 회로(동시성 검출기, 지연선, 자기 유지 무리, 사슬)가 뇌에 실제로 있다는 것을 보여 준다.

{\footnotesize
\setlength{\tabcolsep}{3pt}\renewcommand{\arraystretch}{1.15}
\begin{longtable}{>{\raggedright}p{0.5cm}>{\raggedright}p{4.0cm}>{\raggedright}p{5.6cm}>{\raggedright}p{2.3cm}>{\raggedright\arraybackslash}p{1.7cm}}
\toprule
No. & 명제 & 실험과 측정한 것 & 출처 & 상태 \\
\midrule\endfirsthead
\toprule
No. & 명제 & 실험과 측정한 것 & 출처 & 상태 \\
\midrule\endhead
1 & 뉴런은 문턱값과 초기화가 있는 RC 회로이다~\eqref{eq:glutneuron} & 쥐 피질 피라미드 뉴런(절편)의 세포체에 살아 있는 뇌의 시냅스 흐름과 비슷한 잡음 전류를 주입했다. 평균 발화율이 전류의 평균과 산포에 어떻게 의존하는지는 느린 발화율 적응을 더한 누설 적분 발화 뉴런으로 기술된다. $\tau$와 지연의 범위는 장에서 실험 인용 없이 제시했다 & \cite{Rauch2003} & 측정됨 \\
2 & 모델~\eqref{eq:glutneuron}은 단순화이다. 입력 가지가 스스로 국소 스파이크를 낸다 & 생쥐 시각 피질 피라미드 뉴런의 가지돌기에서 in vivo로 직접 기록: 시각 자극이 NMDA 수용체에 의존하는 국소 가지돌기 스파이크를 일으키며, 이를 억누르면 뉴런의 방위 조율이 넓어진다 & \cite{Smith2013} & 측정됨 \\
3 & 일치 창 $\Delta_{\max}=\tau\ln\frac{w}{\theta-w}$~\eqref{eq:window}; $\theta=1.5w$이면 $0.7\tau$ & 모델~\eqref{eq:glutneuron}의 결과이다. 좁은 합산 창을 직접 측정한 자료는 빠른 억제가 있는 뉴런에 대해 있다(\ref{sec:gaba}장, 그 장 표의 3행) & 이 장의 유도 & 이론 \\
4 & \nt{GLUT} 뉴런만으로 된 네트워크는 입력의 단조 함수만 계산한다(명제~\ref{prop:monotone}) & 장에서의 증명: 우변이 감소하지 않는 침묵에서의 반복. 이것은 모델에 관한 명제이다. 억제 없는 살아 있는 네트워크에서 이를 검증한 실험은 없다 & 이 장의 유도 & 이론 \\
5 & 감각 기관은 전선 쌍을 준다. 어떤 세포는 자극이 강해질 때, 다른 세포는 약해질 때 응답한다 & 망막: 이미 쌍극 세포에서 원뿔 세포의 신호가 켜짐 채널과 꺼짐 채널로 갈라진다. 원숭이에서 켜짐 쌍극 세포를 약리학적으로 차단: 채널은 피질까지 따로 가며, 차단 뒤에는 빛이 강해지는 것은 잘 검출하지 못하지만 약해지는 것은 그렇지 않다 & \cite{Schiller1992} & 측정됨 \\
6 & 이중 레일 부호가 있으면 \nt{GLUT} 뉴런 네트워크는 공통 박자 없이 어떤 논리 함수든 비동기적으로 계산하며, 대가는 두 배의 뉴런 수이다(명제~\ref{prop:dualrail}, \eqref{eq:dualrail}, 그림~\ref{fig:dualxor}) & 이중 레일 부호와 신호 자체로 준비 상태를 판정하는 것은 지연에 둔감한 비동기 회로의 표준 기법이다. 뉴런 여섯 개짜리 배타적 OR 회로는 장에서 조립했다. 뇌가 논리 함수를 바로 이중 레일 부호로 계산한다는 것을 보여 주는 실험은 없다 & \cite{Sparso2001} & 이론 \\
7 & 사람에서 두 귀에 소리가 도착하는 시간의 최대 차는 $\approx0.6\ms$이고, 뇌는 약 10마이크로초를 구별한다 & 두 선택지 실험의 청취자: 시간 차 변별 문턱(정답 75\,\%)은 $150$--$1700$\,Hz 대역 잡음에서 $9$\,\textmu s, $1$\,kHz 순음에서 $11$\,\textmu s, 딸깍 소리에서 $28$\,\textmu s. 상한 $0.6\ms$는 장에서 $0.22\,\text{m}/343\,\text{m/s}$로 계산했다 & \cite{KlumppEady1956} & 측정됨 \\
8 & 새에서는 지연선과 동시성 검출기가 시간 차를 위치로 바꾼다~\eqref{eq:itd}(그림~\ref{fig:itd}) & 원숭이올빼미: 두 귀에서 온 축삭이 동시성 검출기 핵에 반대쪽에서 들어온다. 축삭을 따라 스파이크 도착 시간이 깊이에 따라 규칙적으로 변하고, 핵을 가로지르는 지연의 폭($700$\,\textmu m에 걸쳐 $160$\,\textmu s)은 두 귀 사이 지연의 범위와 일치한다 & \cite{Carr1990} & 측정됨 \\
9 & 포유류에서는 지연의 열이 발견되지 않았다. 같은 문제를 \nt{GLYC} 뉴런의 정확한 시간의 억제가 있는 동시성 검출기가 푼다 & 게르빌루스쥐 내측 상올리브에서 in vivo 기록: 응답이 방향 지도를 이루지 않는다. 미세 피펫으로 글리신과 그 차단제를 가하면, 시간을 정확히 맞춘 \emph{글리신} 억제가 지연의 작동 범위를 정한다는 것이 드러난다. 여기서 억제는 \nt{GABA}가 아니라 \nt{GLYC}에서 온다 & \cite{Brand2002,Grothe2010} & 측정됨 \\
10 & 되먹임이 강한 무리는 플립플롭이다. 자기 유지 활동은 자극 뒤 몇 초 동안 지속된다(그림~\ref{fig:bistable}) & 기억한 점으로 눈을 늦게 움직이는 과제(지연 $1$--$6$\,s)를 하는 원숭이의 전전두 피질: 과제와 관련된 뉴런 $170$개 중 $87$개가 지연 동안 발화율을 바꾸며, 그중 $79\,\%$에서 이 변화는 기억한 점의 방향에 의존한다. 이 활동을 두 안정 상태를 가진 되먹임이 유지한다는 것은 이론이다 & \cite{Funahashi1989,Wang2001} & 측정됨 \\
11 & 유입이 $\approx0.7\tau$보다 길면 비트가 써진다(그림~\ref{fig:recurrentgroup}b) & 방정식 $\tau\,\dd r/\dd t=-r+f(wr+I)$를 $w=1$, $I=0.6$에서 오일러 방법으로 계산했다. \texttt{models/}에 스크립트는 없다. 실험은 없다 & 이 장의 계산 & 모델 \\
12 & 기억은 스파이크 없이 거의 시냅스 촉진만으로 저장할 수 있다 & 이론: 말단의 잔류 칼슘이 기록을 약 1초 유지한다. 근거: 흰담비 내측 전전두 피질(여러 뉴런 동시 기록)에 여운이 긴 촉진성 시냅스로 연결된 피라미드 뉴런 하위 네트워크가 있다. 기억 유지 중 “조용한” 구간 관찰의 총설. 피질이 언제 어떤 방법을 쓰는지는 열린 문제이다 & \cite{Mongillo2008,WangMarkram2006,Stokes2015} & 가설 \\
13 & 기억은 피로로 지워진다. 신경전달물질 저장량이 고갈된다 & 피질 피라미드 뉴런의 쌍 기록: 스파이크가 잦으면 시냅스 응답이 줄고, 감소 속도는 방출 확률로 정해진다. 바로 이것이 자기 유지 무리를 끈다는 것은 직접 보이지 않았다 & \cite{Tsodyks1997} & 측정됨 \\
14 & 무리의 사슬은 사건이 시동하는 타이머이다. 명금류에서는 뉴런들이 엄격히 차례대로 발화한다 & 잠잘 때와 노래할 때 금화조의 고위 발성 중추에서 식별된 뉴런을 기록: 운동핵으로 출력을 보내는 뉴런은 저마다 노래의 정확한 한 순간에 짧은 버스트 하나를 내고, 뉴런들은 차례로 발화한다 & \cite{Hahnloser2002} & 측정됨 \\
\bottomrule
\end{longtable}}

\section{\ref{sec:gaba}장. \nt{GLUT}와 \nt{GABA}의 협업}

이 장의 논리적, 동역학적 명제는 선형 해석과 옴의 법칙으로 장 안에서 유도했다. 실험은 이 명제들이 기대는 회로, 즉 지연된 금지, 흥분 뒤에 따라오는 직접 억제, 억제에 의한 안정화, \nt{GLUT}--\nt{GABA} 루프의 리듬이 뇌에서 실제로 작동한다는 것을 보여 준다.

{\footnotesize
\setlength{\tabcolsep}{3pt}\renewcommand{\arraystretch}{1.15}
\begin{longtable}{>{\raggedright}p{0.5cm}>{\raggedright}p{4.0cm}>{\raggedright}p{5.6cm}>{\raggedright}p{2.3cm}>{\raggedright\arraybackslash}p{1.7cm}}
\toprule
No. & 명제 & 실험과 측정한 것 & 출처 & 상태 \\
\midrule\endfirsthead
\toprule
No. & 명제 & 실험과 측정한 것 & 출처 & 상태 \\
\midrule\endhead
1 & 피질 뉴런의 5분의 1에서 4분의 1이 \nt{GABA} 뉴런이다 & 원숭이 피질 $10$개 영역의 GABA 면역염색: $8$개 영역에서 뉴런의 약 $25\,\%$, 일차 시각 피질에서는 $20\,\%$ 미만. 피질 억제성 뉴런의 다양성에 관한 총설 & \cite{Hendry1987,Markram2004} & 측정됨 \\
2 & 억제가 하는 일은 붙는 자리가 크게 좌우한다. 가지돌기, 세포체, 스파이크가 생기는 자리, 다른 뉴런 전선의 말단 & 총설: 피질 \nt{GABA} 뉴런의 부류는 수신자 위의 표적 자리로 구별된다. 피라미드 뉴런의 세포체와 가지돌기에서 동시 기록: 직접 억제는 가지돌기보다 세포체에서 훨씬 강하다. 척수에서 다른 뉴런 전선의 말단에 가해지는 억제는 입력 선로 하나의 신경전달물질 방출을 줄인다(측정 총설) & \cite{Markram2004,Pouille2001,Rudomin1999} & 측정됨 \\
3 & 중개자를 거친 부호 전환의 값은 지연 하나, 약 1밀리초이다. 흥분 뒤에 따라오는 억제는 일치 창을 좁힌다 & 쥐 해마 절편의 피라미드 뉴런: 중개자 하나를 거친 억제가 직접 흥분 뒤에 짧은 지연으로 따라오며, 두 흥분성 입력이 문턱값까지 합산되는 창은 $2\ms$보다 짧다. 이 억제가 약한 가지돌기에서는 창이 넓다 & \cite{Pouille2001} & 측정됨 \\
4 & 금지 $a\wedge\bar b$~\eqref{eq:veto}는 OR, 상수 1과 함께 함수적으로 완전하다(명제~\ref{prop:gabacomplete}) & 장에서의 증명. 고전적 결과: 억제성 입력이 있는 문턱 소자의 네트워크는 어떤 논리식이든 구현한다. 모델에 관한 명제이며, 실험은 없다 & \cite{McCullochPitts1943} & 이론 \\
5 & 지연된 금지는 최적 속도 $v^*=\Delta x/d$~\eqref{eq:vstar}의 운동 방향 검출기를 준다 & 토끼 망막: 방향 선택성은 “무효” 방향으로 움직일 때 흥분을 끄는 억제로 만들어진다. 선택성 세포의 흥분성 입력과 억제성 입력을 따로 기록한 결과, 별 모양 아마크린 세포(\nt{GABA})가 “무효” 쪽을 향한 돌기로만 비대칭적으로 억제한다는 것이 드러났다. $v^*$ 공식은 장에서 유도했다 & \cite{Barlow1965,Fried2002} & 측정됨 \\
6 & 분로 억제는 전압에서 빼지 않고 전압을 나눈다~\eqref{eq:shunt}(그림~\ref{fig:shunt}) & 염소 평형 전위가 휴지 전위 근처일 때 세 컨덕턴스로 된 분배기에 대한 옴의 법칙. 스파이크를 내지 않는 뉴런에 관한 것이다 & 이 장의 유도 & 이론 \\
7 & 분로는 스파이크 발화율에 뺄셈으로 작용하고, 나눗셈은 억제와 균형 잡힌 배경 잡음의 조합에서 나온다 & 모델: 발화하는 뉴런에서는 분로를 지나는 전류가 거의 일정하여 발화율에 대한 효과가 뺄셈이다. 실험: 쥐 체성감각 피질(절편)의 피라미드 뉴런에 동적 클램프로 흥분성, 억제성 배경 컨덕턴스를 주입했다. 둘을 동시에 균형 있게 늘리면 발화율의 뚜렷한 이동 없이 응답 기울기가 나누어진다. 총설: 총활동으로 나누는 연산은 뇌의 여러 부위에서 나타난다 & \cite{HoltKoch1997,Chance2002,Carandini2012} & 측정됨 \\
8 & $\beta>1$이면 상호 억제가 승자 하나를 고른다(명제~\ref{prop:wta}, \eqref{eq:wta}) & 고윳값 $-(1\pm\beta)/\tau$는 장에서 유도했다. $\beta=0.5$에서 발화율 $0.73$과 $0.53$은 고정점 $r_{1,2}=(I_{1,2}-\beta I_{2,1})/(1-\beta^2)$이다. \texttt{models/}에 스크립트는 없다. 장에서 직접 실험을 제시하지 않는다 & 이 장의 유도 & 이론 \\
9 & \nt{GABA}를 거친 신호로 기억을 초기화; 상호 억제하는 두 무리는 래치이다 & 그림~\ref{fig:bistable}과 명제~\ref{prop:wta}에서 나온다. 실험은 없다 & 이 장의 유도 & 이론 \\
10 & \nt{GLUT}--\nt{GABA} 루프의 평형은 억제가 충분히 강하고 충분히 빠르면 안정하다(명제~\ref{prop:ei}) & 두 집단 모델~\eqref{eq:ei}. 조건 (i)과 (ii)는 그 행렬의 행렬식과 대각합이며 장에서 유도했다 & \cite{WilsonCowan1972} & 이론 \\
11 & $w_{EE}=1.5$, $w_{EI}w_{IE}=2$에서 빠른 억제($\tau_I=2\ms$)는 $E^*=0.67h$를 유지하고, 느린 억제($30\ms$)는 진동을 부추긴다(그림~\ref{fig:ei}) & \eqref{eq:ei}의 수치해, 스크립트 \texttt{figures/make\_ei.py} & 계산 & 모델 \\
12 & 피질은 억제 안정화 체제에서 작동한다. 특징은 \nt{GABA} 뉴런을 밀어 주면 그 발화율이 떨어진다는 것이다~\eqref{eq:isn} & 이론이 역설적 응답을 예측했다. 실험: 고양이 일차 시각 피질의 세포내 기록에서 주변 자극으로 응답이 억눌릴 때 억제성 컨덕턴스는 늘지 않고 흥분성 컨덕턴스와 함께 줄어든다. 생쥐의 시각, 체성감각, 운동 피질에서 PV 뉴런을 광유전학으로 흥분시키면 자극이 없을 때와 가벼운 마취에서도 억제성 뉴런의 발화율이 낮아진다. 그림~\ref{fig:ei}의 수치에서 계수 $-1/3$은 장에서 유도했다 & \cite{Tsodyks1997isn,Ozeki2009,Sanzeni2020} & 측정됨 \\
13 & “\nt{GLUT}가 \nt{GABA}를 흥분시키고, \nt{GABA}가 \nt{GLUT}를 억제한다”는 루프는 주기 $12$--$50\ms$의 빠른 리듬을 낳는다 & 생쥐 체성감각 피질에서 빠른 \nt{GABA} 뉴런을 $8$--$200$\,Hz로 광유전학 자극하면 감마 리듬($20$--$80$\,Hz, 주기 $12$--$50\ms$)이 선택적으로 강해진다. \nt{GLUT} 뉴런 자극은 더 느린 리듬만 강화한다 & \cite{Cardin2009,Buzsaki2012} & 측정됨 \\
\bottomrule
\end{longtable}}

\section{\ref{sec:code}장. 모스 부호}

이 장의 한계 추정은 정보 이론과 계산이다. 채널의 어디서 잡음이 생기는지, 뇌가 얼마나 빨리 작동하는지, 스파이크 하나가 얼마인지는 측정되었다. 피질이 실제로 어떤 부호를 쓰는지는 열린 문제로 남아 있다.

{\footnotesize
\setlength{\tabcolsep}{3pt}\renewcommand{\arraystretch}{1.15}
\begin{longtable}{>{\raggedright}p{0.5cm}>{\raggedright}p{4.0cm}>{\raggedright}p{5.6cm}>{\raggedright}p{2.3cm}>{\raggedright\arraybackslash}p{1.7cm}}
\toprule
No. & 명제 & 실험과 측정한 것 & 출처 & 상태 \\
\midrule\endfirsthead
\toprule
No. & 명제 & 실험과 측정한 것 & 출처 & 상태 \\
\midrule\endhead
1 & 피질 \nt{GLUT} 뉴런의 발화율은 초당 1개 미만에서 수십 개이며, 뉴런은 대부분의 시간을 아래쪽 경계 근처에서 작동한다 & 깨어 있는 쥐의 청각 피질에서 세포에 붙인 피펫으로 기록(뉴런 선택이 활동과 무관): 어느 순간이든 초당 $20$개가 넘는 발화율을 내는 뉴런은 $5\,\%$ 미만이다. 일부 뉴런의 배경 발화율은 초당 $0.01$보다 낮다. 발화율 분포는 로그 정규 분포이다 & \cite{Hromadka2008} & 측정됨 \\
2 & 스파이크 채널의 용량은 $R_{\max}\approx r\log_2\frac{e}{r\Delta t}$~\eqref{eq:spikeentropy}이고, 거의 전부가 스파이크 시각에 있다(명제~\ref{prop:capacity}) & 길이 $\Delta t$ 칸으로 된 이진 문자열의 엔트로피. 장의 수치: 초당 $r=10$, $\Delta t=1\ms$에서 스파이크당 $8$비트, $\Delta t=10\ms$에서 약 $5$비트, 스파이크 계수기는 초당 $2$비트 미만 & \cite{MacKay1952} & 이론 \\
3 & 감각 뉴런은 스파이크당 1비트에서 몇 비트를 전달하며, 가장 좋은 경우 한계의 절반까지 이른다 & 자극을 알고 있는 감각 뉴런의 스파이크로 자극을 복원; 책에 측정 요약이 있다 & \cite{Rieke1997} & 측정됨 \\
4 & 스파이크 발생기는 정확하다. 정확한 시각은 입력의 빠른 변화가 정한다 & 쥐 피질 절편의 뉴런: 빠르게 요동하는 반복 전류에는 스파이크가 $1\ms$보다 나은 정밀도로 반복되고, 일정한 전류에는 시각이 흩어진다 & \cite{Mainen1995} & 측정됨 \\
5 & 시냅스는 믿을 수 없다. 방출의 무작위성 때문에 스파이크의 절반 넘게가 수신자에 이르지 못한다 & 해마 피라미드 뉴런 입력의 최소 자극(절편): 스파이크의 절반 넘게가 응답을 일으키지 않는다. 자극 문턱의 요동, 축삭 가지점에서의 전도 실패, 낮은 실험 온도는 배제했다. 남는 것은 확률적 방출이다 & \cite{AllenStevens1994} & 측정됨 \\
6 & 살아 있는 피질의 스파이크 흐름은 무작위로 보인다. 간격은 푸아송 흐름처럼 흩어지고, 스파이크 수의 분산은 평균에 가깝다. “잡음”의 일부는 동물의 움직임에 관한 메시지이다 & 깨어 있는 원숭이의 시각 영역 V1과 MT 기록: 간격의 변동 계수는 약 $1$이고, 스파이크 수의 분산 대 평균 비는 무작위 과정과 같다. 독립적인 작은 입력을 합하는 모델은 훨씬 작은 산포를 준다. 생쥐 시각 피질에서는 변동의 상당 부분이 동물 자신의 움직임을 찍은 영상으로 예측된다 & \cite{Softky1993,Stringer2019} & 측정됨 \\
7 & 발화율 부호는 느리다. 오차는 $1/\sqrt{rT}$~\eqref{eq:rateerror}인데, 뇌는 $\sim150\ms$ 만에 그림을 알아본다 & 공식은 푸아송 통계로, 장에서 유도했다. 실험: 사람이 $20\ms$ 동안 보여 준 낯선 사진에 동물이 있는지 판단했다. “예”와 “아니오”에 대한 뇌 전위는 약 $150\ms$ 뒤에 갈라진다. 이 시간을 $10$--$20\ms$씩 열 개쯤의 단계로 나눈 것은 장의 추정이지 측정이 아니다 & \cite{Thorpe1996} & 측정됨 \\
8 & 무리에 대한 평균은 상관으로 제한된다. 오차는 기껏해야 $1/\sqrt c$배 줄어든다~\eqref{eq:correlated}(명제~\ref{prop:pooling}) & 원숭이 MT 영역에서 동시에 기록한 뉴런 쌍: 스파이크 수의 상관 계수는 평균 $0.12$이며, 이론적 분석은 이 정도 상관만으로도 무리에서 얻는 이득이 제한됨을 보여 준다. 이론: 한계를 정하는 것은 상관 가운데 신호를 닮은 부분뿐이다. 반대 입장의 모델: 뉴런 $50$--$100$개 무리의 발화율은 스파이크 간격 하나, $10$--$50\ms$ 안에 읽히며, 피질에서 개별 스파이크의 시각은 쓰이지 않는다 & \cite{Zohary1994,MorenoBote2014,ShadlenNewsome1998} & 측정됨 \\
9 & 수는 첫 스파이크 시각~\eqref{eq:latency}, 무리의 스파이크 순서, 국소 리듬의 위상에 적을 수 있다 & 도롱뇽 망막: 짧게 보여 준 그림의 공간 구조가 출력 뉴런 첫 스파이크의 상대 지연에 적히며, 대비와 거의 무관하다. 쥐 해마의 장소 세포: “자기” 장소를 지날 때 스파이크가 세타 리듬($7$--$12$\,Hz) 주기 안에서 점점 일찍 옮겨 가며, 이동 폭은 $100^\circ$에서 $355^\circ$이다. 피질이 스파이크 시각을 얼마나 쓰는지는 열린 문제이다(8행) & \cite{Gollisch2008,OKeefe1993} & 측정됨 \\
10 & 어떤 부호가 읽힐지는 수신자의 시간 상수가 정한다~\eqref{eq:reader}(명제~\ref{prop:reader}). 활성 피질의 뉴런은 동시성 검출기에 더 가깝다 & 식~\eqref{eq:reader}는 장에서 유도했다. in vivo 기록 총설: 활성 피질에서 막 컨덕턴스는 휴지 상태보다 몇 배 크고, 시간 상수는 그만큼 짧다. 피질이 바로 이런 식으로 부호를 전환한다는 것은 직접 보이지 않았다 & \cite{Destexhe2003} & 가설 \\
11 & 고갈 시냅스는 발화율의 변화, 본질적으로 $\Delta r/r$를 전달한다~\eqref{eq:depression}(그림~\ref{fig:depression}) & 피질 피라미드 뉴런의 쌍 기록: 고갈 속도는 방출 확률이 정한다. 고갈이 느리면 응답이 발화율을 반영하고, 빠르면 일치를 반영한다. 쥐 피질에서 측정한 고갈에 기초한 모델: 빠른 입력과 느린 입력에서 같은 상대 발화율 변화가 같은 응답을 준다. 즉 시냅스는 $\Delta r/r$를 전달한다. 수치 $1.7\to2.2$, $6.7$, $90\ms$는 장의 계산이다 & \cite{Tsodyks1997,Abbott1997depression} & 측정됨 \\
12 & 버스트는 “대시”이다. 버스트가 사라질 확률 $(1-p)^k$~\eqref{eq:burst}는 작고, 촉진이 이를 더 줄인다 & 총설: 믿을 수 없는 시냅스에서 단일 스파이크는 자주 사라지지만, 버스트는 촉진 덕분에 믿을 만하게 전달된다. 버스트는 시냅스의 장기 변화를 일으킨다. 뇌가 버스트를 별도의 기호로 읽는다는 것은 가설이다 & \cite{Lisman1997,AllenStevens1994} & 가설 \\
13 & 빠른 \nt{GABA} 뉴런은 리듬과 “구두점”을 정한다. 또 다른 부류는 억제하는 뉴런을 억제하여 관문을 연다(명제~\ref{prop:punctuation}) & 피질 \nt{GABA} 뉴런 부류들의 성질에 관한 총설. 광유전학: 빠른 \nt{GABA} 뉴런을 리듬 있게 흥분시키면 감마 리듬이 생기고, 자극에 대한 응답은 주기의 위상에 의존한다. 생쥐 시각 피질에서 PV 뉴런은 서로를, SST 뉴런은 나머지 모든 부류를, VIP 뉴런은 주로 SST를 억제한다. 깨어 있는 생쥐에서 VIP 뉴런을 흥분시키면 \nt{GLUT} 뉴런의 억제가 풀린다. VIP 뉴런은 강화 신호로 켜진다. 일반 규칙으로서의 역할 분담은 가설이다 & \cite{Tremblay2016,Cardin2009,Pfeffer2013,Pi2013} & 가설 \\
14 & 에너지는 평균 발화율을 초당 스파이크 1개 정도로 제한한다~\eqref{eq:energy} & 해부학적, 생리학적 자료로 본 설치류 회색질의 에너지 예산: 스파이크와 시냅스 전류가 신호 전달 에너지의 $47\,\%$와 $34\,\%$이다. 동시에 활성일 수 있는 뉴런은 $\sim15\,\%$를 넘지 못한다. 사람 피질에서는 동시에 강하게 활성일 수 있는 뉴런이 $\sim1\,\%$를 넘지 못한다. 스파이크당 ATP $\sim10^9$개와 신호 전달 전력 $\sim10$\,W는 이 계산에 기초한 장의 추정이다 & \cite{Attwell2001,Lennie2003} & 이론 \\
15 & 부호는 희소하며 줄당 비트 수가 최대가 되도록 맞추어져 있다. 피질은 정밀도를 에너지와 맞바꾼다(명제~\ref{prop:economy}) & 절약 부호 가설. 근거: 먹이를 제한한 생쥐의 시각 피질에서 AMPA 수용체 컨덕턴스와 시냅스의 ATP 소비가 줄고($-29\,\%$), 스파이크 발화율은 유지되며, 방위 조율은 $32\,\%$ 넓어진다 & \cite{Barlow1961,Padamsey2022} & 가설 \\
\bottomrule
\end{longtable}}

\section{\ref{sec:serocomputer}장. \nt{SERO} 뉴런}

\nt{SERO} 뉴런 자체의 성질(리듬, 자가 억제, 수용체에 따른 부호, 하위 시스템)은 직접 측정되었다. 이 장의 계산, 즉 조절기, 필터, 논리는 이 장의 식에서 나온다. $\tau_c$, 세로토닌의 확산 계수, 방출 지점의 수는 추정이다. 고리가 뇌에서 수준 조절기로 작동한다는 것은 가설이고(봉선핵에서 자가 억제는 점대점이며 짧은 멈춤만 만든다), \nt{SERO}가 지평을 정한다는 것은 이를 지지하는 행동 실험이 있는 가설이다.

{\footnotesize
\setlength{\tabcolsep}{3pt}\renewcommand{\arraystretch}{1.15}
\begin{longtable}{>{\raggedright}p{0.5cm}>{\raggedright}p{4.0cm}>{\raggedright}p{5.6cm}>{\raggedright}p{2.3cm}>{\raggedright\arraybackslash}p{1.7cm}}
\toprule
No. & 명제 & 실험과 측정 내용 & 출처 & 상태 \\
\midrule\endfirsthead
\toprule
No. & 명제 & 실험과 측정 내용 & 출처 & 상태 \\
\midrule\endhead
1 & 세로토닌 작용의 부호는 수신자의 수용체가 고른다(명제~\ref{prop:receptor}) & 세로토닌 수용체는 약리와 기전에 따라 여러 계열로 나뉜다. 5-HT$_{1A}$는 G 단백질을 통해 칼륨 채널을 열어 세포를 억제하고, 5-HT$_{2A}$와 이온성 5-HT$_3$는 흥분시킨다. 한 신경전달물질이 세포에 따라 반대 응답을 낸다. & \cite{BarnesSharp1999} & 측정됨 \\
2 & 깨어 있을 때 \nt{SERO} 뉴런은 초당 몇 개의 스파이크를 고르게 낸다 & 자유롭게 움직이는 고양이의 등쪽 봉선핵 뉴런 기록: 조용히 깨어 있을 때 $2.8\pm0.2$ Hz의 느리고 규칙적인 발화. 서파 수면에서는 발화율이 $34$--$68\,\%$, 렘수면에서는 $98\,\%$ 떨어진다. 마취한 쥐에서는 $1.7\pm0.5$ Hz로 매우 규칙적이다. & \cite{TrulsonJacobs1979, GagnonParent2014, JacobsAzmitia1992} & 측정됨 \\
3 & \nt{SERO} 뉴런은 페이스메이커다. 스파이크마다 밀어 줄 필요는 없지만 고른 배경 입력은 필요하다(절편에서는 $\alpha_1$ 수용체를 통한 노르아드레날린, 생체에서는 \nt{NORA}에서 온다) & 뇌 절편에서 세포는 느리고 고르게 발화하며, 스파이크마다 $200$--$800\ms$의 후과분극이 뒤따른다. 절편에서 자발적으로 발화하는 세포는 생체보다 적다. 고른 리듬에는 $\alpha_1$ 수용체를 통한 노르아드레날린의 지속적 입력이 필요하고, 이를 차단하면 리듬이 꺼진다. & \cite{VandermaelenAghajanian1983} & 측정됨 \\
4 & 수용체는 거의 모두 대사성이다. 응답은 느려서 약 1초 안에 오르고 내린다 & 5-HT$_3$을 뺀 모든 수용체 계열은 G 단백질과 결합한다. 절편에서 유발된 세로토닌 방출은 5-HT$_{1A}$를 통해 1초 안에 오르고 내리는 억제 전류를 만든다. & \cite{BarnesSharp1999, CourtneyFord2016} & 측정됨 \\
5 & 목표 영역에서 세로토닌은 틈뿐 아니라 부피로도 나온다. 봉선핵 자체에서 이웃 억제는 점대점이다 & 절편에서의 고속 주사 순환 전압전류법: 스파이크 하나당 방출량이 시냅스가 있는 영역과 시냅스가 거의 없는 봉선핵에서 같다. 재흡수와 수용체를 차단해도 달라지지 않는다. 말단당 분자 수가 수송체와 수용체보다 훨씬 적고, 세포 밖 농도는 주 수용체의 친화도와 일치한다. 봉선핵에서 5-HT$_{1A}$를 통한 억제는 점대점이다. 수송체가 세로토닌이 틈 밖에 쌓이지 못하게 한다. & \cite{Bunin1998, CourtneyFord2016} & 측정됨 \\
6 & 식~\eqref{eq:seroneuron}의 농도는 재흡수로 줄어든다. $\tau_c$가 1초 정도라는 것은 추정이다 & 생체 뇌에서의 전압전류법: 세포 밖 세로토닌을 제한하는 것은 합성이 아니라 주로 재흡수와 분해다. 인용한 연구에 $\tau_c$의 직접 측정은 없다. 크기 정도는 이 장의 추정이다. & \cite{Hashemi2012serotonin} & 측정됨; $\tau_c$는 추정 \\
7 & 페이스메이커 하나로 만드는 NOT과 NOR; \nt{SERO} 논리의 완전성(명제~\ref{prop:serocomplete}, 그림~\ref{fig:serogates}) & 식~\eqref{eq:seroneuron}에서 나온다. 억제되는 페이스메이커는 NOR이고, NOR은 함수적으로 완전하다. \nt{SERO} 뉴런으로 논리를 만든 실험은 없다. 부피가 분리되어 있어야 한다는 조건은 뇌에서 성립하지 않는다. & 이 장의 유도 & 이론 \\
8 & 세로토닌은 \nt{SERO} 뉴런 자신에 있는 수용체를 통해 그 뉴런을 억제한다 & 마취한 쥐에서 5-HT$_{1A}$ 작용제를 정맥 주사하거나 이온영동으로 주면, 세로토닌 자체와 같은 용량--반응 관계로 등쪽 봉선 뉴런의 발화가 억제된다. 5-HT$_{1B}$ 작용제는 거의 효과가 없다. 절편에서 이웃 세포의 내인성 세로토닌은 5-HT$_{1A}$를 통해 전류와 발화의 짧은 멈춤을 만들지만 평균 발화율은 바꾸지 않는다. & \cite{Sprouse1987, CourtneyFord2016} & 측정됨 \\
9 & 모델에서 공통 부피를 통한 고리는 수준 조절기다. 편차와 시간 상수가 $1+G$배 줄어든다~\eqref{eq:feedback}. 고리가 뇌에서 이렇게 작동한다는 것은 가설이다 & 음의 되먹임 증폭기의 고전적 공식이며, 이 장에서 새로 유도했다. \nt{SERO} 시스템의 고리 이득 $G$는 측정되지 않았다. 측정된 것은 절편에서 자가 억제가 점대점이고, 짧은 멈춤을 만들며, 평균 발화율을 바꾸지 않는다는 점이다. & \cite{Black1934feedback, CourtneyFord2016} & 이론; 뇌에서는 가설 \\
10 & 농도는 발화율의 이동 평균, 즉 저역 통과 필터다~\eqref{eq:lowpass}, 그림~\ref{fig:lowpass} & 선형 방정식~\eqref{eq:seroneuron}의 해. 그림은 $\tau_c=1\,\text{s}$에서 이 공식으로 계산한 것이다. \texttt{models/}에 별도 모델은 없다. & 이 장의 유도 & 이론 \\
11 & 틈의 굴곡은 작은 분자의 확산을 약 $2.6$배 늦춘다 & 점원 방법: 테트라메틸암모늄을 이온영동으로 주입하고 이온 선택성 전극으로 그 농도를 절편과 생체 뇌에서 기록한다. 굴곡도 $\lambda\approx1.6$, 즉 유효 $D$는 자유 확산보다 $\lambda^2\approx2.6$배 작다. 세포 사이 부피의 비율은 약 $0.2$. 이 연구들에서 조직 속 세로토닌 자체의 확산 계수는 측정하지 않았다. 이 장에서 그 값은 추정이다. & \cite{Sykova2008} & 측정됨 \\
12 & 독립 “전선”의 길이 $\ell\sim\sqrt{D\tau}\approx20$~\textmu m~\eqref{eq:diffusionlength}; 전선의 수는 이런 영역의 수로 제한된다(명제~\ref{prop:volumewires}) & 확산 법칙에 $D$와 $\tau$의 추정값(6행과 11행)을 대입한 것이며, 명제~\ref{prop:volumewires}는 그 따름정리다. 체적 전달의 독립 채널 수를 센 직접 실험은 없다. & 이 장의 유도 & 이론 \\
13 & \nt{SERO} 뉴런 하나의 출력은 뇌의 넓은 영역에 퍼져 있다. 방출 지점은 추정으로 $\sim10^5$개 & 쥐 등쪽 봉선 뉴런 하나하나의 완전 재구성: 모든 상행 축삭이 많이 갈라지며, 세포 하나의 축삭 총길이는 최대 $18.7$~cm, 한 세포의 가지가 선조체, 전전두피질, 편도체에 닿는다. 세포당 방출 지점 수는 직접 세지 않았다. & \cite{GagnonParent2014} & 측정됨; $10^5$은 추정 \\
14 & \nt{SERO} 뉴런은 출력과 응답이 다른 몇 개의 하위 시스템으로 나뉜다 & 생쥐에서의 바이러스 유전 표지: 편도체로 가는 뉴런과 전두피질로 가는 뉴런은 분지가 거의 겹치지 않고, 서로 다른 입력을 받으며, 불쾌한 자극에 반대로 반응한다. 각 집단을 켜고 끄면 행동이 서로 다르게 바뀐다. & \cite{Ren2018} & 측정됨 \\
15 & 지수형 할인에서 인내 조건은 $D<\tau_\gamma\ln(R/r_0)$~\eqref{eq:patience}; 측정된 쌍곡선형 할인에서는 $D<\tau_\gamma(R/r_0-1)$ & 할인 $e^{-D/\tau_\gamma}$ 또는 $1/(1+D/\tau_\gamma)$에서 나오며, 이 장에서 유도했다. 몇 초 지연에서 선택하는 원숭이: 할인은 지수보다 쌍곡선에 가깝다. 지연 보상을 알리는 신호에 대한 \nt{DOPA} 뉴런의 응답도 지연에 따라 쌍곡선으로 떨어진다. & 이 장의 유도; \cite{KobayashiSchultz2008} & 이론 \\
16 & \nt{SERO}가 지평 $\tau_\gamma$를 정한다(\ref{sec:horizon}절) & 생쥐에서 등쪽 봉선 \nt{SERO} 뉴런을 광유전학으로 켜면 지연 보상을 더 오래 기다리고, 드물어진 보상을 더 끈질기게 찾는다. 사람에서 세로토닌을 낮추면(트립토판을 뺀 식단) 지연 보상의 할인이 가팔라진다. 농도가 어떻게 $\tau_\gamma$로 바뀌는지는 알려져 있지 않다. & \cite{Doya2002, Miyazaki2014, Lottem2018, Schweighofer2008} & 가설 \\
\bottomrule
\end{longtable}}

\section{\ref{sec:dofa}장. \nt{DOPA}를 이용한 학습}

시냅스의 메커니즘(꾸러미, 동시성 검출기, 칼슘, 가소성 창)과 예측 오차로서의 도파민의 행동은 직접 측정되었다. 규칙들이 기울기로 이어진다는 것과 브로드캐스트의 대가는 정리다. 선택 학습의 결과는 이 책의 모델로 계산한 것이다. 오차를 계산하는 회로는 가설이다.

{\footnotesize
\setlength{\tabcolsep}{3pt}\renewcommand{\arraystretch}{1.15}
\begin{longtable}{>{\raggedright}p{0.5cm}>{\raggedright}p{4.0cm}>{\raggedright}p{5.6cm}>{\raggedright}p{2.3cm}>{\raggedright\arraybackslash}p{1.7cm}}
\toprule
No. & 명제 & 실험과 측정한 것 & 출처 & 상태 \\
\midrule\endfirsthead
\toprule
No. & 명제 & 실험과 측정한 것 & 출처 & 상태 \\
\midrule\endhead
1 & 인공 신경망과 같은 형태의 오차 역전파는 뇌에서 발견되지 않았다 & 직접적인 실험은 없다. 이것은 부재에 관한 명제다. 리뷰들은 이 알고리즘에 무엇이 필요한지(순방향 연결을 되풀이하는 역방향 연결, 별도의 단계)와 어떤 생물학적 근사가 제안되었는지를 다룬다. & \cite{Lillicrap2020, WhittingtonBogacz2019} & 가설 \\
2 & 가중치는 방출 부위 수, 방출 확률, 한 꾸러미에 대한 응답으로 분해된다: $w=N_s\,p\,A_1$~\eqref{eq:quantal} & 방출을 낮춘 개구리 신경근 시냅스: 수신자의 응답은 꾸러미 하나에 대한 자발적 미니 응답의 정수배로 계단식으로 변한다. 스파이크당 꾸러미 수는 푸아송 법칙을 따른다. & \cite{delCastillo1954} & 측정됨 \\
3 & 수신자의 수용체는 동시성 검출기이고, 칼슘이 가중치 변화를 일으킨다 & 생쥐 뉴런의 패치 클램프: 휴지 전위에서는 Mg$^{2+}$ 이온이 글루탐산으로 열린 채널을 막고, 탈분극되면 빠져나간다. 해마 절편: 수신자 안의 칼슘 킬레이터는 장기 강화를 막고, 세포 안에서 빛으로 칼슘을 풀어 주면 그것만으로 전달이 강화된다. & \cite{Nowak1984, Malenka1988calcium} & 측정됨 \\
4 & 결정은 어느 쪽이든 실행한다: 수신자에서는 $A_1$이 커지고, 송신자에서는 $p$가 바뀐다 & 가시 하나 위에서 빛으로 풀어 준 글루탐산은 가시와 그 수용체를 흐르는 전류를 빠르고 지속적으로 키운다. 강화에는 AMPA 수용체의 삽입이 따른다. 수신자의 탈분극은 송신자의 방출을 한동안 억누른다. 이 효과는 틈을 거슬러 가는 엔도카나비노이드가 전달한다(\nt{GABA} 입력에서 확인됨). 단기 고갈은 송신자의 방출 확률에 달려 있다. & \cite{Matsuzaki2004, Malinow2002, WilsonNicoll2001, Tsodyks1997} & 측정됨 \\
5 & 송신자와 수신자가 함께 활동하면 시냅스가 강해진다~\eqref{eq:hebb} & 마취한 토끼: 해마 입력 경로에 짧은 스파이크 열을 준 뒤 응답이 $30$~min에서 $10$~h까지 강화된다. 해마 절편에서 수신자의 스파이크는 같은 순간 활동한 시냅스만 강화한다. & \cite{Bliss1973LTP, Kelso1986hebb} & 측정됨 \\
6 & 규칙~\eqref{eq:oja}은 입력 상관행렬의 주 고유벡터를 찾는다(명제~\ref{prop:oja}) & 정리이며 증명은 본문에 있다. 뉴런이 자기 입력의 주성분을 배운 실험은 없다. 시냅스에서 가중치 증가를 제한하는 메커니즘은 밝혀지지 않았다. & \cite{Oja1982} & 이론 \\
7 & 시각에 따른 헤브 규칙: 창은 약 $\pm20\ms$(그림~\ref{fig:stdp}). 반복되는 순서열에서 사슬이 자란다 & 배양한 해마 뉴런 쌍: 송신자의 스파이크 뒤 $20\ms$ 안의 수신자 스파이크는 지속적 강화를, $20\ms$ 앞의 스파이크는 약화를 준다. 둘 다 NMDA 수용체에 의존한다. 그림의 비율 $A_-=0.6A_+$는 예시다. 이렇게 신경망에서 사슬이 자란다는 것은 직접 측정되지 않았다. & \cite{BiPoo1998} & 측정됨 \\
8 & 적격 흔적~\eqref{eq:threefactor}: 함께 일어난 활동 뒤 $1$--$2\,\text{s}$ 안에 온 도파민은 연결을 강화하고, 더 늦으면 그렇지 않다(명제~\ref{prop:threefactor}) & 선조체 절편, 글루탐산 입력과 도파민 입력의 개별 광유전학 자극: 도파민이 글루탐산 입력 뒤 $0.3$--$2\,\text{s}$ 안에 올 때만 가시가 자란다. 창은 cAMP의 빠른 상승과 분해로 정해진다. 피질에서는 짝짓기가 강화와 약화에 대해 따로따로 흔적을 남기고, 모노아민 수용체가 이것을 몇 초 규모의 창 안에서 장기 변화로 바꾼다. & \cite{Yagishita2014, He2015eligibility} & 측정됨 \\
9 & 몇 초 길이의 창은 도파민 없이도 있다: 세 번째 신호는 수신자의 강한 흥분이다 & 살아 있는 생쥐, CA1 영역: 수신자의 플래토 사건 하나가 그 앞뒤 몇 초 안에 온 입력을 강화하고, 한 번 지나가는 것만으로 장소장을 만든다. & \cite{Bittner2017} & 측정됨 \\
10 & 3요소 규칙의 평균 걸음은 보상의 기울기를 따른다~\eqref{eq:perturbation}(명제~\ref{prop:gradient}) & 무작위 이탈로 학습할 때의 기울기에 관한 정리. 본문에서는 작은 잡음에 대해 유도했다. & \cite{Williams1992} & 이론 \\
11 & 잡음은 탐색이다: 신경망은 자기의 무작위 이탈로부터 배운다 & 어린 금화조: 기저핵 회로의 출력 핵을 끄면 노래의 변동성이 급격하고 가역적으로 사라진다. 성체 금화조: 음높이가 평균의 한쪽에 있는 음절을 부를 때 방해음을 주면 새는 음높이를 빠르게 반대쪽으로 옮긴다. 자기 산포로부터 배우는 것이다. & \cite{Olveczky2005, TumerBrainard2007} & 측정됨 \\
12 & 두 선택지 중 고르기는 $\tau_e=1\,\text{s}$, $\delta=R-\bar R$에서 학습되고, $\tau_e=50\ms$에서는 학습되지 않는다. 기대값을 빼지 않으면 가중치가 한없이 자라고, 학습률이 크면 신경망이 더 나쁜 선택을 굳힐 수 있다 & \texttt{models/dofa\_learning.py}, $\eta=0.05$, 시도 $500$회, 시드 $6$개. $\tau_e=1\,\text{s}$: $A$ 선택 비율 $0.65$--$0.79\to0.96$--$1.00$, 가중치 $0.85$--$1.33$. $\tau_e=50\ms$: $0.38$--$0.55$, 가중치 변화 없음. $\delta=R$: 승자의 가중치 $860$--$1190$. $\delta=R$, $\eta=0.5$, 시드 $3$개: 하나가 $B$에 고정($0.04\to0$). 스크립트가 네 경우를 모두 출력하며, 재계산 결과는 본문과 일치했다. & \texttt{models/\allowbreak dofa\_\allowbreak learning.py} & 모델 \\
13 & 하나의 공통 신호로 학습하는 시간은 잡음을 내는 뉴런 수에 비례해 늘어난다(명제~\ref{prop:broadcastcost}) & 선형 신경망의 학습 곡선: 뉴런 섭동은 정확한 기울기보다 출력 뉴런 수만큼 느리고, 가중치 섭동은 거기에 입력 수만큼 더 느리다. & \cite{Werfel2005} & 이론 \\
14 & \nt{DOPA}의 출력은 무엇보다 행동 선택 영역으로 간다 & 쥐의 흑질선조체 \nt{DOPA} 뉴런 개개의 축삭 완전 재구성: 세포 하나의 가지가 선조체 부피의 $0.45$--$5.7\,\%$(평균 $2.7\,\%$)를 덮고, 선조체 밖에는 가지가 거의 없다. & \cite{Matsuda2009} & 측정됨 \\
15 & 피질은 자기 입력을 예측하도록 배우며, 오차는 그 자리에서 계산된다 & 리뷰: 감각 피질에는 기대한 입력과 실제 입력의 불일치에 대한 응답이 있다. 이것이 피질 학습의 일반 규칙이라는 것은 증명되지 않았다. & \cite{Keller2018} & 가설 \\
16 & 도파민은 오차~\eqref{eq:ctd}처럼 행동한다: 급증, 예고 신호로의 이동, 급감(그림~\ref{fig:ctdcases}) & 원숭이의 \nt{DOPA} 뉴런: 예상 못 한 주스에 급증. 학습 뒤에는 예고 신호에 급증이 생기고 약속된 주스에는 응답이 없다. 주스가 오지 않으면 급감. 연속 형태~\eqref{eq:ctd}는 이론이다. & \cite{Schultz1997, Doya2000} & 측정됨 \\
17 & $\dd V/\dd t$를 계산하고 기대를 빼는 회로 & 생쥐, 배쪽 피개 영역의 기록과 광유전학: \nt{DOPA} 뉴런은 보상에서 기대를 뺀다. 이웃한 \nt{GABA} 뉴런은 보상이 기대될 때 이들을 억제하며, 이 뉴런들을 흥분시키거나 억제하면 오차가 바뀐다. 도함수가 어떻게 계산되는지는 밝혀지지 않았다. & \cite{Eshel2015arithmetic} & 가설 \\
18 & \nt{DOPA} 뉴런은 페이스메이커이며, 급감은 급증보다 얕다 & 절편에서 \nt{DOPA} 뉴런은 느린 페이스메이커 전위 위에서 자발적으로, 매우 규칙적으로 발화한다. 원숭이에서 발화 빈도는 양의 오차를 정량적으로 따르지만 음의 오차는 약하게만 전달한다. & \cite{GraceOnn1989, Bayer2005} & 측정됨 \\
19 & 행위자와 비평자(그림~\ref{fig:actorcritic}) & 알고리즘은 강화 학습 이론에서 왔다. 사람(fMRI)에서 배쪽 선조체는 선택이 있든 없든 비평자처럼 행동하고, 등쪽 선조체는 행동을 골라야 할 때만 그렇다. & \cite{SuttonBarto2018, ODoherty2004actor} & 가설 \\
20 & 두 번째 계열의 수용체를 가진 뉴런에서는 규칙의 $\delta$ 부호가 뒤집힌다 & 선조체 절편: D1 수용체를 가진 뉴런에서는 짝짓기가 D1이 필요한 강화를, D2를 가진 뉴런에서는 D2가 필요한 약화를 준다. & \cite{Shen2008} & 측정됨 \\
\bottomrule
\end{longtable}}

\section{\ref{sec:dofaalt}장. \nt{DOPA}의 다른 이론들}

확장된 그림들은 저마다 도파민의 직접 측정(반전점의 흩어짐, 불쾌한 자극과 새로운 자극에 대한 반응, 완만한 상승, 맛 교체에 대한 반응)에 기대고 있다. 그러나 그것을 설명하는 식들은 이론이며, 그 가운데 두 이론 사이에서는 실험 결과가 아직 서로 모순된다.

{\footnotesize
\setlength{\tabcolsep}{3pt}\renewcommand{\arraystretch}{1.15}
\begin{longtable}{>{\raggedright}p{0.5cm}>{\raggedright}p{4.0cm}>{\raggedright}p{5.6cm}>{\raggedright}p{2.3cm}>{\raggedright\arraybackslash}p{1.7cm}}
\toprule
No. & 명제 & 실험과 측정한 것 & 출처 & 상태 \\
\midrule\endfirsthead
\toprule
No. & 명제 & 실험과 측정한 것 & 출처 & 상태 \\
\midrule\endhead
1 & 비대칭 규칙~\eqref{eq:asymmetric}은 점~\eqref{eq:expectile}으로 수렴한다. $\kappa=1/2$이면 평균이고, 확률 $1/2$의 한 방울이면 $V=\kappa$이다(그림~\ref{fig:expectiles}) & 점~\eqref{eq:expectile}은 익스펙타일, 곧 양의 편차와 음의 편차에 서로 다른 가중치를 준 최소제곱 문제의 해다. 이 장의 예에 대한 유도는 본문에 있다. & \cite{NeweyPowell1987}; 이 장의 유도 & 이론 \\
2 & \nt{DOPA} 뉴런마다 반전점이 다르고, 반전점은 비대칭 순으로 정렬된다(명제~\ref{prop:expectile}) & 생쥐, 배쪽 피개 영역의 \nt{DOPA} 뉴런을 광유전학으로 표지, 크기가 다른 보상: 급증이 급감으로 바뀌는 점이 뉴런마다 다르다. 반응 비대칭이 큰 뉴런일수록 그 점이 높다. 무리 반응으로 보상 분포의 모양이 복원된다. 이것이 흩어짐의 주된 기능이라는 것은 가설이다. & \cite{Dabney2020} & 측정됨 \\
3 & 일부 \nt{DOPA} 뉴런은 불쾌한 자극에 급증으로 반응한다 & 원숭이: 어떤 \nt{DOPA} 뉴런은 보상 예고 신호에 흥분하고 눈에 부는 공기 예고 신호에 억제되며, 다른 뉴런은 둘 다에 흥분한다. 두 무리는 중뇌의 서로 다른 부위에 있다. & \cite{Matsumoto2009, BrombergMartin2010} & 측정됨 \\
4 & 오차의 크기 또는 새로움이 학습률을 정한다 & \nt{DOPA} 뉴런의 중요도 신호가 학습률을 바꾼다는 직접 실험은 인용된 연구에 없다. 리뷰는 “주의, 중요함” 신호의 역할을 제안한다. & \cite{BrombergMartin2010} & 가설 \\
5 & 위험 축을 위한 별도의 전선 & 생쥐: 선조체 꼬리로 가는 \nt{DOPA} 뉴런은 보상이 아니라 새롭고 강한 자극에 반응한다. 이 뉴런을 파괴하면 새로운 물체에 대한 회피가 약해진다. & \cite{Menegas2018} & 측정됨 \\
6 & 최적 행동 시간 $\tau_a^*=\sqrt{C/\bar R}$~\eqref{eq:vigor} & 합 $C/\tau_a+\bar R\tau_a$의 최소. 유도는 본문에 있다. 원래 모델에서 행동 속도는 평균 보상과 같은 시간의 가격이 정한다. & \cite{Niv2007}; 이 장의 유도 & 이론 \\
7 & 느린 도파민 수준은 $\bar R$를 싣고 행동 속도를 정한다(명제~\ref{prop:multiplex}) & 쥐, 측좌핵의 미세투석과 전압전류법: 도파민 수준이 분 단위로 보상 빈도와 행동 속도를 따라간다. 사람에서 L-DOPA는 반응 시간이 평균 보상 빈도에 의존하는 정도를 키운다. 느린 수준이 정확히 $\bar R$라는 것은 증명되지 않았다. & \cite{Hamid2016, Beierholm2013vigor} & 가설 \\
8 & 느린 수준은 두 원천에서 합쳐진다. 방출은 스파이크 발화율이 변하지 않아도 출력 말단에서 달라지지만, 발화율 자체에도 완만한 상승이 있다 & 쥐, 한 과제: 측좌핵의 방출은 변하는 보상 기대를 따라가지만, 식별된 배쪽 피개 영역 \nt{DOPA} 뉴런의 발화율은 그렇지 않다. 가상 복도의 생쥐(10행): 보상에 다가갈 때의 완만한 상승이 식별된 \nt{DOPA} 뉴런의 스파이크에도 있다. & \cite{Mohebi2019, Kim2020} & 측정됨 \\
9 & 동물이 먼 보상에 다가가는 동안 도파민이 완만하게 오른다 & 미로의 쥐, 선조체 전압전류법: 목표에 가까워지면서 도파민 농도가 몇 초에 걸쳐 오르고, 기울기는 거리와 보상 크기에 의존한다. & \cite{Howe2013ramp, Hamid2016} & 측정됨 \\
10 & 상승은 추정이 정확해질 때의 오차 $\delta$다~\eqref{eq:ramp}. 복도에서의 도약은 급증을 준다(그림~\ref{fig:ramp}) & 식과 수치 초당 $\delta=0.75\,V$는 본문의 유도. 실험: 가상 복도의 생쥐를 목표 쪽으로 순간 이동. 세포체의 스파이크, 세포체와 출력 말단의 칼슘, 선조체의 농도가 기대 $V$가 아니라 오차처럼 반응한다. 두 설명 중 어느 것이 모든 출력에서 옳은지는 논쟁 중이다. & \cite{Kim2020} & 측정됨 \\
11 & 뒤돌아보기: 도파민은 인과 연관의 척도~\eqref{eq:retro}를 알린다 & 생쥐, 이 이론과 오차~\eqref{eq:ctd}가 서로 다른 것을 예측하는 실험에서의 측좌핵 도파민 방출: 여러 실험에서 방출이 앞의 이론을 따랐다. & \cite{Jeong2022} & 가설 \\
12 & 학습 중 도파민 반응은 보상에서 예고 신호로 점차 옮겨 간다 & 생쥐, 학습 과정 동안 배쪽 선조체에서 \nt{DOPA} 뉴런 출력 말단의 신호: 급증이 보상 순간에서 예고 신호 순간으로 점차 옮겨 간다. \eqref{eq:ctd}에 의한 학습이 요구하는 그대로다. & \cite{Amo2022} & 측정됨 \\
13 & \nt{DOPA} 뉴런은 가치가 같아도 보상의 맛이 바뀌면 반응한다 & 쥐, 배쪽 피개 영역 뉴런 기록: 주스를 가치가 같은 다른 주스로 바꾸면 오차~\eqref{eq:ctd}가 0인데도 반응이 생긴다. & \cite{Takahashi2017} & 측정됨 \\
14 & 인위적 급증은 두 중립 자극 사이의 연관을 가르친다 & 쥐, 보상 없이 두 자극을 미리 짝짓는 실험: 첫째 자극에서 둘째 자극으로 넘어갈 때 \nt{DOPA} 뉴런을 광유전학으로 켜면 연관이 생기고, 끄면 학습이 방해된다. & \cite{Sharpe2017} & 측정됨 \\
15 & 특징별 오차 벡터~\eqref{eq:vectortd} & 특징 예측 오차 이론. 벡터 형태에 대한 직접 검증은 없다. & \cite{Gardner2018} & 가설 \\
16 & 한 과제에서 \nt{DOPA} 뉴런마다 서로 다른 양을 싣는다 & 가상 복도의 생쥐, \nt{DOPA} 뉴런의 2광자 영상: 어떤 뉴런은 보상과, 어떤 뉴런은 속도와 방향 전환과, 또 어떤 뉴런은 예고 신호와 선택의 정확도와 연관된다. & \cite{Engelhard2019} & 측정됨 \\
17 & 전선 대신 다발(명제~\ref{prop:bundle}) & 2--16행의 요약: 각 분할은 따로따로 측정되었다. 이것들이 모두 한 신호의 부분이라는 통합된 그림은 실험으로 확인되지 않았다. & --- & 가설 \\
\bottomrule
\end{longtable}}

\section{\ref{sec:nora}장. \nt{NORA}}

\nt{NORA} 시스템의 구조, 두 작동 모드, 동공과의 연관(\nt{NORA}만을 통한 것은 아니다), 신호 대 배경 비의 증가, 행동에서의 뒤집힌 U는 직접 측정되었다. 곱셈 이득 $g$는 모델이다. 이득이 선택 모드를 전환하고 역온도로 작동한다는 것은 이 장의 유도 결과다. 이득 조정의 이점은 이 책의 모델로 계산한 것이다. “\nt{NORA}는 탐색의 노브다”와 “\nt{NORA}는 인터럽트다”는 가설이다.

{\footnotesize
\setlength{\tabcolsep}{3pt}\renewcommand{\arraystretch}{1.15}
\begin{longtable}{>{\raggedright}p{0.5cm}>{\raggedright}p{4.0cm}>{\raggedright}p{5.6cm}>{\raggedright}p{2.3cm}>{\raggedright\arraybackslash}p{1.7cm}}
\toprule
No. & 명제 & 실험과 측정한 것 & 출처 & 상태 \\
\midrule\endfirsthead
\toprule
No. & 명제 & 실험과 측정한 것 & 출처 & 상태 \\
\midrule\endhead
1 & 사람의 \nt{NORA} 뉴런은 수만 개이며, 거의 모두 한 무리에 있다 & 사람 뇌줄기의 연속 절편, 티로신 수산화효소 염색: 청반에 $53\,900$개, 이웃한 하위핵에 $6\,260$개의 세포. & \cite{Baker1989LC} & 측정됨 \\
2 & 출력은 거의 뇌 전체로 퍼지지만, 무리의 부분들이 서로 다른 영역을 어느 정도 나누어 맡는다 & 생쥐에서의 바이러스 유전 표지: 청반의 \nt{NORA} 뉴런은 뇌의 넓은 영역에서 입력을 받고, 뇌와 척수의 넓은 영역으로 출력을 보낸다. 쥐: 편도체로 가는 무리는 공포 학습을 강화하고, 전전두피질로 가는 별도의 무리는 공포 소거를 강화한다. & \cite{Schwarz2015LC, Uematsu2017LC, Poe2020} & 측정됨 \\
3 & 배경 모드: 초당 1회 미만에서 몇 회 사이의 고른 스파이크, 발화율은 각성 수준에 따라 변한다 & 쥐, 수면--각성 주기 동안의 청반 뉴런 기록: 발화율은 깨어 있을 때 가장 높고, 서파 수면에서 낮으며, 렘수면에서는 거의 0이다. 발화율 변화가 상태 전환보다 앞선다. & \cite{AstonJonesBloom1981} & 측정됨 \\
4 & 돌발 모드: 과제에 중요한 사건에 대략 결정 순간에 나오는 짧은 버스트 & 원숭이, 드문 표적 찾기 과제: 모든 청반 뉴런이 표적에만 버스트로 응답한다. 표적 뒤 $\sim90\ms$, 손 반응 $\sim200\ms$ 전이다. 수행이 나쁠 때 버스트가 약하다. 양자택일 과제에서 버스트는 자극보다 반응에 더 정확히 맞물리며, 정답과 오답 모두에서 나타난다. & \cite{AstonJones1994vigilance, Clayton2004LC} & 측정됨 \\
5 & 동공 지름은 \nt{NORA}의 부정확한 측정점이다. \nt{NORA}뿐 아니라 \nt{ACET}와 다른 영역도 따른다 & 원숭이: 청반의 스파이크는 세포 하나의 개별 스파이크까지 동공 확장을 예측한다. 비슷하지만 더 늦은 연관이 둔덕과 띠피질에도 있다. 생쥐: 동공 변동은 피질로 가는 노르아드레날린성 출력과 콜린성 출력의 활동을 모두 따른다. & \cite{Joshi2016, Reimer2016} & 측정됨 \\
6 & 노르아드레날린은 배경 활동을 유발 활동보다 더 강하게 억제한다. 곱셈 이득~\eqref{eq:gain}은 이 효과를 전달하는 모델이다 & 깨어 있는 원숭이, 청각 피질, 노르아드레날린 이온영동: 뉴런의 배경 활동이 같은 종의 소리에 대한 응답보다 더 강하게 억제되어 신호 대 잡음비가 커진다. 곱셈 시그모이드~\eqref{eq:gain}는 네트워크 모델에서 가져왔다. 기울기가 문자 그대로 곱해진다는 것은 측정되지 않았다. & \cite{Foote1975NE, ServanSchreiber1990} & 측정됨; $g$는 모델 \\
7 & 이득은 증폭기 뒤의 잡음에 대해서만 $d'$를 높인다~\eqref{eq:gainsnr}(명제~\ref{prop:gainnoise}) & 이 장의 유도. 네트워크 모델에서 이득은 신호 검출을 개선한다. 증폭기 앞과 뒤의 잡음을 분리해 이 차이를 검증한 실험은 없다. & 이 장의 유도; \cite{ServanSchreiber1990} & 이론 \\
8 & 경계 $g\beta=1$은 선택 네트워크를 “숙고 중”에서 “결정함”으로 전환한다~\eqref{eq:wtagain}(명제~\ref{prop:gainwta}, 그림~\ref{fig:pitchfork}) & 서로 억제하는 두 무리의 선형 분석; 이 장의 유도. 뇌에서 이 문턱값을 직접 측정한 것은 없다. & 이 장의 유도 & 이론 \\
9 & 이득은 선택의 역온도다~\eqref{eq:softmax} & 증폭기 뒤의 잡음이 로지스틱이면 선택 확률은 $T=\sigma/g$인 softmax다. 이 장의 유도. & 이 장의 유도 & 이론 \\
10 & \nt{NORA}는 얼마나 탐색할지 정한다. 높은 배경 활동은 작은 유효 $g$(탐색), 결정 순간의 돌발 버스트는 큰 $g$(결정) & 사람: 무작위 선택 직전의 기저 동공은 알려진 최선을 고르기 직전보다 넓다. 동공이 넓은 사람일수록 탐색 성향이 높다. 쥐: 무작위 선택 모드로의 전환은 앞띠피질로 가는 \nt{NORA} 입력이 조절한다. 버스트가 유효 $g$를 높인다는 것은 직접 측정되지 않았다. & \cite{JepmaNieuwenhuis2011, Tervo2014stochastic, AstonJones2005} & 가설 \\
11 & 보상은 $g$에 대해 뒤집힌 U로 의존한다. 빠르게 변하는 세계에서 최선의 $g$가 더 작다(명제~\ref{prop:explore}, 그림~\ref{fig:invertedU}) & \texttt{models/nora\_explore.py}, $4000$번 시도의 실행 $60$회. $1000$번에 한 번 변화: 최선 $g=16$, 비율 $0.976$; $20$번에 한 번: $g=8$, 비율 $0.886$; $g=0.5$에서 $0.77$. 재계산 결과가 본문과 일치했다. & \texttt{models/\allowbreak nora\_\allowbreak explore.py} & 모델 \\
12 & 행동에서의 뒤집힌 U: 배경 활동이 적당하고 버스트가 뚜렷할 때 가장 잘한다 & 원숭이, 4행과 같은 과제: 수행이 좋은 시기에는 배경 발화율이 적당하고 표적에 대한 버스트가 뚜렷하다. 배경 발화율이 높으면 버스트가 없고 오경보가 많다. 배경 활동과 유발 활동의 변화는 수행의 변동을 따라간다. & \cite{Usher1999LC, AstonJones2005} & 측정됨 \\
13 & \nt{NORA}는 예상치 못한 불확실성을, \nt{ACET}는 예상된 불확실성을 알린다 & 두 종류의 불확실성을 다루는 베이즈 추론 이론. 인용한 연구에는 이 신호들을 신경전달물질별로 분리한 직접 실험이 없다. & \cite{YuDayan2005} & 가설 \\
14 & 급격한 변화 뒤 사람은 더 빨리 학습하고 동공이 커진다 & 사람, 갑작스러운 변화가 있는 예측 과제: 동공의 짧은 변화는 변화 확률의 추정을 따르며, 기저 동공과 함께 새 데이터가 추정을 얼마나 바꾸는지(학습률)를 예측한다. 빛이나 과제와 무관한 동공 변화도 이 학습률을 바꾼다. 여기서 동공이 바로 \nt{NORA}를 보여 준다는 것은 5행의 간접 자료에 기댄 추론이다. & \cite{Nassar2012} & 측정됨 \\
15 & 돌발 버스트는 인터럽트로 작동한다. 네트워크가 상태를 초기화한다 & \nt{NORA} 버스트가 네트워크 상태를 초기화한 직접 실험은 없다. 측정된 것은 중요한 사건에 대한 버스트뿐이다(4행). & \cite{BouretSara2005} & 가설 \\
16 & 뜻밖의 일에 따라 $g$나 학습률을 조정해도 최선의 고정 설정보다 거의 낫지 않다 & \texttt{models/nora\_explore.py}, 시기당 $1000$번 시도인 세계($1000$번에 한 번, $20$번에 한 번 변화). 최선의 고정 $g=8$: $0.925$; $g$ 조정: 최대 $0.927$; 학습률 조정: 최대 $0.926$; 고정 학습률 $0.4$: $0.928$. 재계산 결과가 본문과 일치했다. & \texttt{models/\allowbreak nora\_\allowbreak explore.py} & 모델 \\
\bottomrule
\end{longtable}}

\section{\ref{sec:acet}장. \nt{ACET}를 더하다}

아세틸콜린의 세 가지 작용은 절편과 생체 뇌에서 측정되었고, 쓰기와 읽기의 충돌은 이 책의 모델에서 보였다. \nt{ACET}가 쓰기 허용 선으로 쓰이며 예상된 불확실성을 알린다는 것은 실험의 부분적 뒷받침이 있는 가설이다.

{\footnotesize
\setlength{\tabcolsep}{3pt}\renewcommand{\arraystretch}{1.15}
\begin{longtable}{>{\raggedright}p{0.5cm}>{\raggedright}p{4.0cm}>{\raggedright}p{5.6cm}>{\raggedright}p{2.3cm}>{\raggedright\arraybackslash}p{1.7cm}}
\toprule
No. & 명제 & 실험과 측정 내용 & 출처 & 상태 \\
\midrule\endfirsthead
\toprule
No. & 명제 & 실험과 측정 내용 & 출처 & 상태 \\
\midrule\endhead
1 & 뇌의 \nt{ACET} 뉴런은 많지 않고 몇 개의 집단에 모여 있으며, 출력은 피질 전체로 퍼진다. 브로드캐스트 채널이다 & 쥐에서 아세틸콜린 합성 효소에 대한 항체 염색과 역행성 표지: 피질, 해마, 편도체, 후각 망울, 시상은 아세틸콜린을 주로 기저 전뇌와 뇌간의 조밀한 세포 집단 여섯 개에서 받는다 & \cite{Mesulam1983} & 측정됨 \\
2 & 피질의 아세틸콜린 수준은 깨어 있을 때 높고 서파 수면에서 낮다 & 자유롭게 움직이는 고양이의 피질과 해마에서 뇌파 기록과 함께 한 미세투석: 아세틸콜린 방출은 서파 수면보다 조용히 깨어 있을 때 $\sim100\%$, 활발히 깨어 있을 때 $\sim175\%$ 높다. 렘수면에서 피질의 방출은 깨어 있을 때와 같다 & \cite{Marrosu1995,Hasselmo1999} & 측정됨 \\
3 & 신호의 의미는 수용체가 정한다. 아세틸콜린에는 빠른 채널형 수용체와 세포 안 반응을 일으키는 느린 수용체가 있다(명제~\ref{prop:receptor}) & 측정 결과의 총설: 아세틸콜린이 흥분성, 전달, 가소성에 미치는 작용은 방출 장소, 수용체 아형(니코틴성은 이온 채널, 무스카린성은 G 단백질 경유), 수신 세포에 달려 있다 & \cite{Picciotto2012} & 측정됨 \\
4 & 첫째 작용: 외부 입력은 거의 건드리지 않고 내부 연결을 약하게 한다(식~\eqref{eq:acetnet}의 인수 $1-a$) & 쥐 후각 피질 절편: 아세틸콜린 작용제 $100$\,\textmu M에서 같은 세포의 내부 연결(Ib층) 전위는 $60\%$ 넘게, 외부 입력(Ia층) 전위는 $15\%$ 미만으로 줄어든다. 억제는 시냅스 전 작용이다. 해마 절편(CA1): 내부 경로 $89\%$ 대 입력 경로 $40\%$ & \cite{HasselmoBower1992,HasselmoSchnell1994} & 측정됨 \\
5 & 둘째 작용: 입력을 강하게 한다(인수 $\frac{1+a}{2}$) & 신피질 절편: 니코틴성 수용체는 시상에서 오는 입력 시냅스를 강화하고 피질 내부 시냅스에는 작용하지 않는다. 무스카린성 수용체는 두 종류를 모두 약하게 한다. 아세틸콜린은 뉴런의 흥분성을 높인다(총설) & \cite{Gil1997,Hasselmo2006} & 측정됨 \\
6 & 셋째 작용: 가중치 변화를 쉽게 한다(학습률 $\eta a$) & 쥐: 피질 아세틸콜린의 공급원인 기저핵을 전기 자극하면서 소리를 함께 들려주면, 다 자란 동물에서도 청각 피질 지도가 들려준 소리 쪽으로 크게 재편된다 & \cite{KilgardMerzenich1998,Hasselmo2006} & 측정됨 \\
7 & 식~\eqref{eq:acetnet}의 두 번째 식, 즉 드문 패턴용 헤브 규칙은 일부로부터 패턴을 완성하는 연상 메모리를 준다 & 성긴 패턴과 공분산 규칙을 쓰는 네트워크의 용량 계산: 활성 뉴런 비율 $p$가 작으면 네트워크는 $p=1/2$일 때보다 눈에 띄게 많은 패턴을 저장한다 & \cite{Tsodyks1988} & 이론 \\
8 & 낮은 $a$에서의 쓰기는 메모리를 망가뜨린다. 네트워크는 본 것이 아니라 떠올린 것을 기록한다. $a=0.1$에서도 손상은 $a=0$보다 약하지 않다 & \texttt{models/\allowbreak acet\_memory.py}: 뉴런 $200$개, $20$개짜리 패턴 다섯 개, 새 패턴은 그중 하나와 뉴런 $10$개를 공유, 실행 $10$회. $a=0$에서 새 패턴은 기억되지 않고, 옛 패턴은 남의 뉴런 $10$개 중 $5.9$개를 끌고 나온다. $a=0.1$에서 새 패턴은 떠올려지지만($0.97$) 둘이 합쳐진다. 새 패턴은 남의 뉴런 $7.3$개, 옛 패턴은 $7.9$개를 끌고 나온다. $a=0.2$부터는 하나 미만 & 이 장의 유도 & 모델 \\
9 & 명제~\ref{prop:writeread}: 쓰기와 읽기가 똑같이 잘 되는 수준 $a$는 없다(그림~\ref{fig:acetsim}) & 같은 스크립트, 학습률 $\eta a$: 새 패턴은 한 번 제시로 $a\ge0.9$에서 완전히, $a=0.75$에서 $0.8$만큼 기억되고, $a\le0.5$에서는 기억되지 않는다. 절반으로부터의 읽기: $a\le0.2$에서 $1.0$, $a=0.3$에서 $0.96$, $a=0.4$에서 $0.04$ & 이 장의 유도 & 모델 \\
10 & 뇌에서는 브로드캐스트 전선 하나, 즉 \nt{ACET}가 쓰기와 읽기를 전환한다 & 이 발상은 절편 측정을 바탕으로 만든 모델에서 나왔다. 가장 직접적인 실험은 사람에서다. 무스카린성 수용체 차단제 스코폴라민은 새 단어 쌍, 특히 옛 쌍과 겹치는 쌍의 학습을 해치지만, 이미 배운 쌍을 떠올리는 것은 해치지 않는다. 네트워크의 두 모드를 직접 측정한 결과는 없다 & \cite{HasselmoAndersonBower1992,Atri2004} & 가설 \\
11 & 필터~\eqref{eq:kalman}는 최적이다. 입력 가중치와 학습률은 같은 양 $P/R$이고, 정상 상태에서 $P=\sqrt{qR}$, 기억은 $\sqrt{R/q}$이다(명제~\ref{prop:kalman}) & 실험이 필요 없다. 칼만--뷰시 필터의 최적성은 추정 이론의 고전적 결과이고, 정상 상태는 이 장에서 유도했다 & \cite{KalmanBucy1961} & 이론 \\
12 & \nt{ACET}는 $P$와 비슷한 양, 즉 예상된 불확실성을 네트워크에 알린다 & 주의의 확률 모델에서 제안되었다. 아세틸콜린 수준이 추정의 불확실성을 따라간다는 직접 측정은 없다 & \cite{YuDayan2005} & 가설 \\
13 & \nt{ACET} 뉴런의 일부는 보상과 처벌에 20밀리초 남짓 만에, 강화가 뜻밖일수록 더 강하게 응답한다 & 생쥐, 광유전학으로 식별한 기저 전뇌 콜린성 뉴런: 보상과 처벌에 대한 응답은 $18\pm3$\,ms 뒤, 크기는 강화가 뜻밖일수록 커지며, 두 핵의 뉴런에서 응답이 비슷하다 & \cite{Hangya2015} & 측정됨 \\
14 & \nt{NORA}는 세계의 예상치 못한 변화를 알아채고, \nt{ACET}는 네트워크를 쓰기 모드로 유지한다 & 이 분업은 모델에서 제안되었다. 간접 증거: 사람에서 \nt{NORA}와 연관된 동공 지름이 변화와 학습률을 따라간다. \nt{NORA}--\nt{ACET} 쌍을 직접 검증한 실험은 없다 & \cite{YuDayan2005,Nassar2012} & 가설 \\
15 & 서파 수면에서 읽기 모드의 네트워크는 낮에 기록한 것을 재생하고, 이 재생이 그것을 장기 기억으로 옮겨 쓴다 & 쥐 해마 뉴런 앙상블 기록: 낮에 함께 작동한 세포는 이어지는 서파 수면에서 더 자주 함께, 같은 순서로 발화한다(측정됨). 옮겨 쓰기에 관해서는, 학습 뒤 해마 잔물결을 억제하면 공간 기억이 나빠지고 자발적 잔물결을 늘이면 좋아진다. 그 원인이 낮은 \nt{ACET} 수준이라는 것은 증명되지 않았다 & \cite{WilsonMcNaughton1994,Skaggs1996,Girardeau2009,FernandezRuiz2019,Hasselmo1999} & 가설 \\
\bottomrule
\end{longtable}}

\section{\ref{sec:synthesis}장. 기계 전체}

요약 장이므로 새 실험을 소개하지 않는다. 각 행은 그 명제를 다룬 장과 같은 출처에 기댄다. \nt{GLUT} 버스와 \nt{GABA}를 통한 흐름 제어는 확고하게 확립되어 있고, 브로드캐스트 채널의 역할과 그 공동 동작은 대부분 가설이다.

{\footnotesize
\setlength{\tabcolsep}{3pt}\renewcommand{\arraystretch}{1.15}
\begin{longtable}{>{\raggedright}p{0.5cm}>{\raggedright}p{4.0cm}>{\raggedright}p{5.6cm}>{\raggedright}p{2.3cm}>{\raggedright\arraybackslash}p{1.7cm}}
\toprule
No. & 명제 & 실험과 측정한 것 & 출처 & 상태 \\
\midrule\endfirsthead
\toprule
No. & 명제 & 실험과 측정한 것 & 출처 & 상태 \\
\midrule\endhead
1 & 뉴런 $8.6\cdot10^{10}$개에 제어 신호는 네 종류뿐이다~\eqref{eq:machine} & 뇌 균질액에서 세포핵 계수(등방성 분획법): 성인 남성의 뉴런은 $86.1\pm8.1$\,$\times10^9$개 & \cite{Azevedo2009} & 측정됨 \\
2 & 명제~\ref{prop:architecture}의 앞 두 층: 데이터는 \nt{GLUT} 주소 버스를 따라 흐르고, 연결의 부호는 송신자가 정하며, 흐름은 \nt{GABA} 뉴런이 이끌고 정규화한다(\ref{sec:neurontypes}, \ref{sec:gaba}장) & 뉴런당 주된 신경전달물질은 하나(측정 총설); 피드포워드 억제는 피라미드 뉴런이 입력을 합산하는 창을 좁힌다; 전체 활동에 의한 나눗셈은 여러 영역에서 측정됨 & \cite{Strata1999,Pouille2001,Carandini2012} & 측정됨 \\
3 & 소자는 신뢰할 수 없다: 시냅스는 스파이크 대부분을 잃는다(표~\ref{tab:computer}, \ref{sec:code}장) & 해마 절편, 최소 자극: 스파이크의 절반 넘게 CA1에서 응답을 일으키지 못한다; 문턱값과 축삭 전도의 실패는 배제되었고, 원인은 확률적 신경전달물질 방출 & \cite{AllenStevens1994} & 측정됨 \\
4 & 뇌는 $\sim20$와트로 동작하며, 뉴런당 평균 초당 약 한 번 스파이크한다(\ref{sec:code}장); 엔지니어는 아직 이런 기계를 만들지 못했다 & 측정된 비용에서 얻은 추정~\eqref{eq:energy}: 시냅스 작업을 포함해 스파이크당 ATP 분자 약 $10^9$개(설치류 피질); 피질에 대한 상세 추정은 더 낮은 발화율을 준다. 기술 현황은 총설에 따름 & \cite{Attwell2001,Lennie2003,MehonicKenyon2022} & 이론 \\
5 & 대부분 시냅스의 학습은 국소적 적격 흔적과 하나의 공통 수 $\delta$의 곱이다(\eqref{eq:machine}의 둘째 식, \ref{sec:dofa}장) & 원숭이의 \nt{DOPA} 뉴런은 보상 예측 오차에 응답한다; 선조체 절편에서 도파민은 글루탐산 입력 $0.3$--$2$\,s 뒤에 올 때만 가시를 키운다. 적격 흔적이 측정됨 & \cite{Schultz1997,Yagishita2014} & 측정됨 \\
6 & 출력마다 따로 가는 오차 전선은 운동 학습 회로에만 있다(\ref{sec:motorlearning}장) & 소뇌: 등반섬유 신호가 입력과 동시에 오면 그 입력이 약해진다; 원숭이에서 푸르키네 세포의 복합 스파이크는 운동 오차의 방향을 나르고 보정을 일으킨다 & \cite{Ito2001,Herzfeld2018} & 측정됨 \\
7 & 각 브로드캐스트 채널은 몇 개 선로의 다발이다(명제~\ref{prop:bundle}) & \nt{DOPA}: 같은 과제에서 서로 다른 뉴런이 보상, 운동, 자극 특징을 나른다; 빠른 오차와 느린 도파민 수준은 서로 다르다. \nt{SERO}, \nt{NORA}, \nt{ACET}에서는 이런 분해가 덜 연구됨 & \cite{Engelhard2019,Hamid2016,Mohebi2019} & 측정됨 \\
8 & \nt{SERO}: 수준 조절기이자, 가설상 지평 $\tau_\gamma$(\ref{sec:serocomputer}장) & 자기 억제: 자기 5-HT$_{1A}$ 수용체의 작용제가 솔기핵 세로토닌 뉴런을 억제한다. 측정됨. 지평은 이론에서 제안됨; 가장 강한 실험은 생쥐에서 이 뉴런의 광유전학적 활성화가 지연 보상의 기다림을 늘린다는 것 & \cite{Sprouse1987,Doya2002,Miyazaki2014} & 가설 \\
9 & \nt{NORA}: 이득 $g$가 탐색과 결정 사이에서 고른다(\ref{sec:nora}장) & 신호 대 잡음비 증가: 깨어 있는 원숭이의 청각 피질에서 노르아드레날린은 동종의 울음소리에 대한 응답보다 배경 활동을 더 강하게 억제한다. 측정됨; 곱셈적 이득은 모델; 탐색과 결정 사이의 선택에서의 역할은 이론 & \cite{Foote1975NE,ServanSchreiber1990,AstonJones2005} & 가설 \\
10 & \nt{ACET}: 쓰기와 읽기를 전환한다(\ref{sec:acet}장) & 세 작용(내부 연결 약화, 입력 강화, 가소성 촉진)은 측정됨; 모드 전환은 사람에서 스코폴라민 실험으로 간접 지지됨 & \cite{HasselmoBower1992,Gil1997,Atri2004} & 가설 \\
11 & 식~\eqref{eq:machine}은 기계 전체를 기술한다 & \ref{sec:glutcomputer}--\ref{sec:acet}장 모델의 요약이며, 각 부분은 자기 장에 기댄다; 전체로서는 실험으로 검증되지 않음 & 장 안의 유도 & 가설 \\
12 & 노브들은 공통 제어 없이 협조해 동작한다: 오차, 놀람, 쓰기, 복귀(그림~\ref{fig:timeline}) & 실험 없음: 정성적 도식. 개별 고리는 \nt{NORA}와 \nt{ACET}의 분업 이론, 그리고 사람에서 동공과 학습 속도의 관계 & \cite{YuDayan2005,Nassar2012} & 가설 \\
13 & 하나의 공통 신호에 의한 학습은 결과에 영향을 주는 뉴런이 많을수록 느리다(명제~\ref{prop:broadcastcost}) & 선형 네트워크 학습 곡선 계산: 출력 섭동에 의한 학습은 경사 하강법보다 출력 수만큼 느리다 & \cite{Werfel2005} & 이론 \\
14 & 피질이 다층 회로를 어떻게 조정하는지는 알려져 있지 않다; 후보는 스파이크 버스트, 뉴런 가지 안의 계산, 예측에 의한 자기 지도 학습 & 오차 운반자로서의 버스트는 모델; 피질 뉴런 상세 모델의 응답은 $5$--$8$층 네트워크만 재현했다. 모델; 사람 피질 뉴런 가지에서 XOR을 내는 칼슘 스파이크는 절편에서 측정됨; 자기 지도 학습은 증거 총설 & \cite{Payeur2021,Beniaguev2021,Gidon2020,Keller2018} & 가설 \\
\bottomrule
\end{longtable}}

\section{\ref{sec:sensors}장. 기계의 입력}

센서의 구조는 단일 세포 전류 기록, 달팽이관 진동 기록, 망막 세포 계수로 직접 측정되었다. 감도의 한계와 산탄 잡음 식은 물리에서 나온다. 센서의 부호가 최대 정보에 맞춰져 있다는 것은 강한 근거를 가진 가설이다.

{\footnotesize
\setlength{\tabcolsep}{3pt}\renewcommand{\arraystretch}{1.15}
\begin{longtable}{>{\raggedright}p{0.5cm}>{\raggedright}p{4.0cm}>{\raggedright}p{5.6cm}>{\raggedright}p{2.3cm}>{\raggedright\arraybackslash}p{1.7cm}}
\toprule
No. & 명제 & 실험과 측정한 것 & 출처 & 상태 \\
\midrule\endfirsthead
\toprule
No. & 명제 & 실험과 측정한 것 & 출처 & 상태 \\
\midrule\endhead
1 & 센서는 변환기이다: 수용체가 전류를 만들고, 눈과 귀 감각 세포의 출력은 \nt{GLUT} 버스로 가는 글루탐산이다; 첫 세포들은 스파이크를 내지 않는다(그림~\ref{fig:transducer}) & 거북의 막대세포와 원뿔세포는 흥분성 아미노산을 방출한다: 떼어 낸 막 조각의 글루탐산 채널이 이를 잡는다. 쥐 달팽이관의 속털세포: 신경섬유 말단의 전류는 글루탐산 AMPA 수용체를 거치며, 방출 빈도는 세포의 완만한 탈분극과 함께 $150$\,Hz까지 커진다 & \cite{CopenhagenJahr1989,GlowatzkiFuchs2002} & 측정됨 \\
2 & 어둠 속 빛 센서는 광자 하나에 약 1피코암페어의 전류로 응답한다 & 두꺼비 막대세포 바깥 분절에 흡입 전극: 희미한 섬광에 대한 응답이 양자화되어 있고, 사건은 $\approx1$\,pA(암전류의 $5\%$), 퍼짐 $0.2$\,pA, 효율 $0.5$. 사람은 광자 하나가 들어온 것을 우연보다 자주 알아맞힌다 & \cite{Baylor1979,Tinsley2016} & 측정됨 \\
3 & 소리 센서는 1나노미터보다 작은 변위를 알아챈다 & 뫼스바우어 방법, 기니피그 달팽이관 기저막의 $16$--$19$\,kHz 위치: 청신경 응답 문턱에서 막의 속도는 약 $0.04$\,mm/s, 곧 변위 $\approx0.35$\,nm(우리의 환산). 잰 것은 센서 다발이 아니라 막이다 & \cite{Sellick1982,Hudspeth1989} & 측정됨 \\
4 & 어둠 속 정밀도는 광자 산탄 잡음~\eqref{eq:photonnoise}이 제한한다; 약한 빛에서는 적분이 길어진다 & 식은 푸아송 통계에서 나온다. 막대세포에서 희미한 빛의 잡음은 독립적인 양자 사건들의 합과 일치한다; 밝은 배경에서 섬광 응답은 더 작고 짧다. “수십분의 1초”라는 수치는 여기서 따로 검증되지 않음 & \cite{Baylor1979} & 이론 \\
5 & 입력 범위는 빛이 열 자릿수, 소리가 열두 자릿수이고, 스파이크 발화율은 세 자릿수가 안 된다 & 소리: 특성 주파수에서 기저막 응답은 $40$--$80$\,dB 대역에서 기울기 $0.2$\,dB/dB까지로 커지며, 압축은 $100$\,dB 위에서도 보인다. 자릿수 자체는 표준 추정치로, 본문에서는 출처 없이 제시 & \cite{Ruggero1997} & 측정됨 \\
6 & 센서 응답은~\eqref{eq:nakarushton}이며, $\sigma$는 평균 밝기를 따라가 곡선을 로그 축을 따라 이동시킨다(그림~\ref{fig:adaptation}) & 이 식은 처음 물고기 망막의 S 전위에 맞춰졌다. 거북 원뿔세포: 응답은 $V=I V_m/(I+\sigma)$를 따르고 밝기 $\sim3.5$자릿수를 담는다; 배경에 $2$--$3$분 노출하면 곡선은 폭을 유지한 채 수평으로 이동한다. 전체 활동에 의한 나눗셈과의 관계는 총설 & \cite{NakaRushton1966,NormannPerlman1979,Carandini2012} & 측정됨 \\
7 & 센서는 변화를 전하고, 그 부호는 스파이크당 최대 정보에 맞춰져 있다(명제~\ref{prop:sensors}) & 원리는 이론적으로 제안되었다. 가장 강한 증거: 파리에서 첫 층 뉴런의 “대비 $\to$ 응답” 특성이 자연 장면 대비의 누적 분포와 일치한다. 이렇게 해서 최대 정보가 얻어진다 & \cite{Barlow1961,Laughlin1981} & 가설 \\
8 & 켜짐 센서와 꺼짐 센서는 두 전선 부호이다; 이벤트 카메라는 이를 실리콘으로 재현한다 & 실험 총설: 망막의 켜짐 채널과 꺼짐 채널은 첫 시냅스에서 이미 갈라지고 따로 차단된다. 카메라: 프레임 없는 $128\times128$픽셀, 범위 $120$\,dB, 지연 $15\,\mu\text{s}$ & \cite{Schiller1992,Lichtsteiner2008,Gallego2022} & 측정됨 \\
9 & 눈에는 빛 센서 $\sim10^8$개와 출력 뉴런 $\sim10^6$개가 있다: $\sim100$배 압축 & 사람 망막 전체 표본에서의 계수: 평균 막대세포 $92$백만 개, 원뿔세포 $4.6$백만 개; 출력(신경절) 세포는 눈에 따라 $0.7$에서 $1.5$백만 개 & \cite{Curcio1990,CurcioAllen1990} & 측정됨 \\
10 & 생쥐 눈의 출력 채널은 서른 개가 넘는다 & 생쥐: 출력 세포 $11\,000$개 이상의 이광자 칼슘 영상과 응답의 군집화, $30$개보다 뚜렷이 많은 기능 유형. 사람의 수는 측정되지 않음 & \cite{Baden2016} & 측정됨 \\
11 & 기니피그의 눈은 $\sim875\,000$ bit/s를 전송한다; 사람 눈으로 환산하면 $\sim10^7$ bit/s & 기니피그, 다전극 배열, 자연 장면 속 움직임: 세포당 $6$--$13$ bit/s, 출력 세포 $\sim10^5$개에 $\sim875\,000$ bit/s. 사람($\sim10^6$개 세포)의 $\sim10^7$ bit/s는 저자들의 환산 & \cite{Koch2006} & 측정됨 \\
12 & 귀는 거의 로그인 주파수 지도를 가진 대역 통과 필터 뱅크이다(그림~\ref{fig:filterbank}) & 기저막 최대 진동 위치는 주파수에 따라 다르다; “주파수--위치” 함수는 거의 지수적이며, 하나의 형태로 사람과 살아 있는 동물의 달팽이관을 기술한다 & \cite{Greenwood1990,Robles2001} & 측정됨 \\
13 & 공진기는 능동적이다: 센서 일부가 약한 진동을 진폭으로 수천 배($66$--$76$\,dB) 증폭하고, 소리가 커지면 증폭이 줄어든다 & 친칠라 달팽이관 기저부의 레이저 진동 측정: 약한 순음의 특성 주파수에서 등자뼈 대비 이득 $66$--$76$\,dB; 죽으면 감도가 $60$--$81$\,dB(진폭으로 $10^3$--$10^4$배) 떨어지고 압축이 사라진다. 힘의 원천은 바깥털세포 & \cite{Ruggero1997,Robles2001} & 측정됨 \\
14 & 귀의 증폭기는 자기 발진 경계 근처에 있다 & 계산: 호프 분기점 근처에서는 범위 압축, 날카로운 조율, 결합음이 불가피하다. 이것이 알려진 청각의 비선형성 전부이다. 달팽이관이 실제로 그렇게 조율되어 있다는 직접 증명은 없음 & \cite{Eguiluz2000} & 가설 \\
15 & 낮은 주파수에서 스파이크는 소리와 위상을 맞춰 나가며(기니피그에서 위상 고정은 $\sim600$\,Hz부터 약해지고 $\sim3.5$\,kHz까지 드러난다), 이로부터 방향을 찾는다; 높은 주파수에서는 위치 부호만 남는다 & 기니피그 청신경: 위상 고정은 약 $600$\,Hz에서 떨어지기 시작해 $3.5$\,kHz 위에서 사라진다; 고양이와 원숭이에서는 경계가 약 한 옥타브 높다. 두 귀의 도착 시간 차는 총설 & \cite{PalmerRussell1986,Grothe2010} & 측정됨 \\
16 & 나머지 센서(표~\ref{tab:sensors}): 후각은 수용체 유형 $\sim400$개, 뉴런당 하나, 조합 부호; 미각은 일부 ATP와 세로토닌을 통해; 촉각은 빨리, 느리게 적응하는 센서; 따뜻함과 차가움은 별도 & 생쥐: 한 수용체가 많은 냄새를, 한 냄새가 많은 수용체를 인식하며, 냄새마다 집합이 다르다. ATP 수용체가 없으면 미각 신경이 응답하지 않는다; 미뢰는 세로토닌을 방출한다. 사람 손 피부의 단일 섬유 기록: 적응 속도에 따른 네 유형. 냉각 채널은 열 채널과 다르다 & \cite{Mombaerts2004,Malnic1999,Finger2005,Huang2005,JohanssonVallbo1979,McKemy2002} & 측정됨 \\
\bottomrule
\end{longtable}}

\section{\ref{sec:recognition}장. 인식}

인식 속도, 첫 단계들의 구조, 응답의 선형 판독, 시간의 연속성에 의한 학습은 사람과 원숭이에서 측정되었다. 사슬 전체를 두 연산으로 환원하는 것과 영상에 의한 학습은 이 책의 모델로 확인했다. 착시의 설명은 근거가 탄탄한 것부터 논쟁 중인 것까지 있다.

{\footnotesize
\setlength{\tabcolsep}{3pt}\renewcommand{\arraystretch}{1.15}
\begin{longtable}{>{\raggedright}p{0.5cm}>{\raggedright}p{4.0cm}>{\raggedright}p{5.6cm}>{\raggedright}p{2.3cm}>{\raggedright\arraybackslash}p{1.7cm}}
\toprule
No. & 명제 & 실험과 측정한 것 & 출처 & 상태 \\
\midrule\endfirsthead
\toprule
No. & 명제 & 실험과 측정한 것 & 출처 & 상태 \\
\midrule\endhead
1 & 이동이 있을 때 픽셀 단위 견본 비교는 동전보다 낫지 않다; 단계 쌍~\eqref{eq:scstage}는 도형을 어디서나 인식한다 & \texttt{models/\allowbreak recognition\_models.py}, $24$개 점의 “망막”, $4000$번 시행: 뒤집힌 점의 비율 $0$, $0.1$, $0.2$에서 견본은 $0.49$, $0.51$, $0.52$; $S$와 $C$ 쌍은 $1.00$, $0.86$, $0.70$ & 이 장의 유도 & 모델 \\
2 & 그림 속 동물은 $\sim150$\,ms 만에 인식된다. 열 개쯤의 단계를 단계당 $10$--$20$\,ms씩 한 번 통과 & 사람, $20$\,ms 동안 보여 준 사진에서 “동물이 있는가” 과제: 두 답에 대한 유발 전위가 $\sim150$\,ms에 갈라진다. 원숭이: 첫 응답 시간이 단계마다 늘어난다(V1이 가장 먼저, 이어서 V2와 V4). 단계 수와 “단계당 스파이크 0개 또는 1개”는 이 수치에서 나온 추론 & \cite{Thorpe1996,Schmolesky1998} & 측정됨 \\
3 & $S$ 단계는 부분들의 AND, $C$ 단계는 위치들에 대한 최댓값이다(명제~\ref{prop:conveyor}) & 고양이, 일차 시각 피질: 단순 세포는 특정 위치의 특정 방위 가장자리에, 복합 세포는 더 넓은 영역의 같은 방위에 응답한다. 복합 세포의 세포 내 기록: 막대 두 개에 대한 응답은 많은 세포에서 두 응답 중 큰 쪽에 가깝고, 평균적으로 max 연산에 가장 가깝다. 상호 억제에 의한 최댓값은 명제~\ref{prop:wta}, 전체 활동에 의한 나눗셈은 총설 & \cite{HubelWiesel1962,Lampl2004,Carandini2012} & 측정됨 \\
4 & 사슬 위로 갈수록 조합은 커지고 복잡해지며, 응답은 위치에 점점 덜 좌우된다 & 원숭이, 자극을 “결정적 특징”까지 단순화: V4와 뒤쪽 하측두 피질 세포는 작은 수용장에서 형태, 색, 질감의 복잡한 조합에 응답한다; 앞쪽 하측두 피질에서는 수용장이 더 크다. 모델에서 $S$와 $C$의 교대가 이 그림을 재현한다 & \cite{KobatakeTanaka1994,Riesenhuber1999} & 측정됨 \\
5 & 합성곱 신경망은 이 교대를 되풀이하며 윗단계 응답을 가장 잘 예측한다; 물체는 뉴런 무리의 발화율에서 선형으로 읽힌다 & 인식을 학습했지만 뇌에 맞추지 않은 네트워크: 윗층은 원숭이 하측두 피질 뉴런의 응답을, 중간층은 V4의 응답을 예측한다. 무작위로 고른 세포 $\sim100$개의 활동에 대한 선형 분류기가 $12.5$\,ms부터의 창에서 위치와 크기가 달라도 물체와 범주를 알아본다 & \cite{Fukushima1980,Yamins2014,Hung2005} & 측정됨 \\
6 & 흔적을 가진 헵 규칙~\eqref{eq:traceinvariance}은 흔적이 $\sim20$\,ms보다 짧지 않으면 이름표 없는 영상으로 불변성을 가르친다 & 느린 특징과 흔적 규칙의 착상은 이론. \texttt{models/\allowbreak recognition\_models.py}, $C$ 뉴런 두 개, 입력 $16$개, $10$번 실행, 물체 사이 휴지 $0.3$\,s. $\tau_{\text{tr}}=5$\,ms에서 도형과의 일치는 평균 $0.52$, $10$\,ms에서 $0.56$; $20$, $50$, $200$\,ms에서 모든 실행이 $1.00$($30$\,ms에서는 열 번 중 한 번 틀림) & \cite{Foldiak1991,WiskottSejnowski2002} & 모델 \\
7 & 뇌에서 불변성은 시간의 연속성으로 학습된다 & 원숭이: 시야의 특정 위치로 눈이 도약하는 동안 물체를 몰래 바꿔치기했다; 하측두 피질 뉴런의 위치 불변성이 바꿔치기에 따라 변했고, 한 시간 만에 이미 유의했다. 이것이 시각 학습의 주된 규칙이라는 것은 가설 & \cite{LiDiCarlo2008} & 측정됨 \\
8 & 움직임: 방향 검출기, 이어서 넓은 영역에 대한 같은 AND/OR 쌍 & 토끼 망막의 방향 검출기; 원숭이 MT 영역에서 기록한 뉴런 $110$개 모두 방향 선택적이며, 움직임에 대한 응답이 V1보다 강하다 & \cite{Barlow1965,Albright1984} & 측정됨 \\
9 & 윗단계는 아래로 예측을 보내고, 위로는 오차가 올라간다 & 모델은 일차 시각 피질 수용장의 비고전적 효과를 설명한다; 개별 관찰은 부합하지만(총설) 파이프라인의 일반 구조로서는 증명되지 않음 & \cite{RaoBallard1999,Keller2018} & 가설 \\
10 & 어려운 그림은 되먹임과 여러 바퀴를 필요로 한다 & 원숭이: “어려운” 그림에서는 하측두 피질의 판정이 $\sim30$\,ms 늦게 형성된다; 이 늦은 응답은 순방향 네트워크가 잘 예측하지 못하고, 되먹임 네트워크나 더 깊은 네트워크가 더 잘 예측한다 & \cite{Kar2019} & 측정됨 \\
11 & 눈에 띄지 않는 픽셀 위조가 네트워크를 속인다; 사람의 시각은 훨씬 강하지만 무적은 아니다 & 네트워크: 특별히 고른 작은 섭동이 답을 바꾼다; “판다 $\to$ 긴팔원숭이” 예는 두 번째 논문에서. 사람에서는 짧게 보여 주면 네트워크 사이에 잘 전이되는 위조가 선택을 옮긴다 & \cite{Szegedy2014,Goodfellow2015,Elsayed2018} & 측정됨 \\
12 & 기대의 착시: 믿기 어려운 입력은 기대에 밀린다. 흐린 것은 “더 느리게” 움직이고, 드레스는 사람마다 다르게 보인다(\ref{sec:illusions}절) & 대비가 낮은 격자는 더 느리게 움직이는 것처럼 보인다. $1401$명 설문: 드레스는 세 가지 안정된 방식으로 보이고, 조명에 대한 힌트가 색을 바꾼다; 관찰자 $\sim13\,000$명에서 조명에 대한 가정이 결정한다. 느린 움직임을 기대하는 베이즈 모델이 여러 움직임 착시를 설명한다 & \cite{Thompson1982,LaferSousa2015,Wallisch2017,Weiss2002,Purves2011} & 가설 \\
13 & 이득 조절: 폭포, 기울기, 대비; 헤르만 격자는 이웃 억제가 아니다 & 토끼 망막의 방향 검출기는 오래 움직임을 보면 발화율이 떨어지고, 멈춘 뒤에는 배경 아래로 떨어진다. 기울기 착시는 이웃 활동에 의한 나눗셈 모델로 재현된다. 망막 출력 뉴런의 응답을 바꾸지 않는 격자 변형이 점을 없앤다. 이웃 억제 설명은 기각됨 & \cite{BarlowHill1963,Schwartz2009,SchillerCarvey2005} & 측정됨 \\
14 & 지연 보상: 움직이는 물체가 섬광보다 앞서 보인다 & 착시는 측정됨. 정신물리학은 움직임의 전방 연장과 지연 차이 둘 다와 모순된다; 섬광 뒤 $\sim80$\,ms의 사건에 따른 사후 설명이 제안되었다. 논쟁은 끝나지 않음 & \cite{Nijhawan1994,EaglemanSejnowski2000} & 가설 \\
15 & 정지한 “뱀”이 돈다: 지연 차이~\eqref{eq:latency}를 방향 검출기가 읽는다 & 착시는 색 없이도 일어난다; 이웃한 요소 쌍이 하나씩 착시를 만든다; 방향 선택적인 원숭이 피질 뉴런은 이 정지한 쌍에 방향성 응답을 낸다. 사람에서 회전은 미세 눈 도약과 눈 깜박임이 촉발한다 & \cite{Conway2005,OteroMillan2012} & 측정됨 \\
16 & 이중 그림: 상호 억제, 피로, 잡음; 전환은 주로 잡음이 일으킨다 & 경쟁하는 뉴런 무리 모델: 전환은 \emph{잡음}이 일으키며, 잡음이 없으면 전환도 없다; 자극 세기에 따른 전환 빈도 증가를 설명한다. 여기서 피로의 역할은 작다 & \cite{MorenoBote2007} & 가설 \\
17 & 착시의 일부는 기계가 배우는 과제의 결과이다 & 영상의 다음 프레임을 예측하도록 학습한 네트워크는 정지한 고리의 회전을 “보고” 대조 그림에서는 보지 않는다; 이미지의 잡음 제거와 선명화를 위한 네트워크는 사람처럼 색과 밝기 착시에 넘어간다 & \cite{Watanabe2018,GomezVilla2020} & 측정됨 \\
\bottomrule
\end{longtable}}

\section{\ref{sec:muscles}장. 근육으로의 출력}

출력의 구조, 곧 운동 단위, 시냅스의 안전 여유, 크기 순서대로의 동원, 근육 쌍, 신장 루프, 리듬 발생기는 직접 측정되었다. 지연 루프의 안정 조건은 정리이다. 몸의 내부 모델은 근거가 탄탄한 가설이다. 그림의 수치는 이 책의 모델에 따른 계산이다.

{\footnotesize
\setlength{\tabcolsep}{3pt}\renewcommand{\arraystretch}{1.15}
\begin{longtable}{>{\raggedright}p{0.5cm}>{\raggedright}p{4.0cm}>{\raggedright}p{5.6cm}>{\raggedright}p{2.3cm}>{\raggedright\arraybackslash}p{1.7cm}}
\toprule
No. & 명제 & 실험과 측정한 것 & 출처 & 상태 \\
\midrule\endfirsthead
\toprule
No. & 명제 & 실험과 측정한 것 & 출처 & 상태 \\
\midrule\endhead
1 & 운동 단위: \nt{ACET} 뉴런 하나가 수백에서 이천 개 남짓의 섬유 무리를 제어한다; 섬세한 조정을 하는 근육에서는 단위가 더 작다 & 사람 근육, 운동 축삭과 근육 섬유 계수: 뉴런당 섬유는 평균적으로 손의 골간근에서 약 $340$개, 완요골근에서 약 $410$개, 안쪽 장딴지근에서 약 $1930$개. 가장 작은 단위의 정확한 수는 본문에 없다. & \cite{Feinstein1955mu} & 측정됨 \\
2 & 근육 위의 시냅스는 섬유를 동작시키는 데 필요한 양의 몇 배나 되는 신경전달물질을 방출한다 & 신경근 시냅스 측정 총설: 성체 포유류의 안전 여유(방출량 대 문턱 양의 비)는 보통 $3$--$5$; 사람에서는 여유가 방출량보다 수신 측 막의 주름에 더 많이 기댄다. & \cite{WoodSlater2001} & 측정됨 \\
3 & 연축~\eqref{eq:twitch}은 $50\ms$ 정도의 시간 $T_c$ 동안 커진다 & 수의적 힘을 낼 때 사람 손의 단일 운동 단위, 단위의 스파이크로 평균한 힘: 상승 시간 $30$--$100\ms$, 단위의 $80\,\%$가 $70\ms$ 미만. 함수 $(t/T_c)e^{1-t/T_c}$는 단위 풀 모델의 근사. & \cite{MilnerBrown1973rec, Fuglevand1993} & 측정됨 \\
4 & 연축의 합~\eqref{eq:forcesum}: $5$, $15$, $40$ Hz에서 힘 $0.2$--$1.1F_1$, $1.8$--$2.2F_1$, 약 $5.4F_1$(그림~\ref{fig:tetanus}) & \texttt{models/\allowbreak muscle\_\allowbreak models.py}에 따른 계산, $T_c=50\ms$: 마지막 $0.3\,\text{s}$에서 $0.21$--$1.10$, $1.76$--$2.19$, $5.28$--$5.49\,F_1$. 선형 합은 단순화: 실제 근육의 포화는 모델에 없다. & \texttt{models/\allowbreak muscle\_\allowbreak models.py} & 모델 \\
5 & 단위는 크기 순서대로 켜진다; 가장 약한 단위는 가장 강한 단위보다 백 배 약하다 & 고양이 배쪽 뿌리: 반사 입력이 커질 때 뿌리에서 스파이크가 작은 뉴런, 곧 작은 뉴런이 먼저 방전한다. 사람 손의 골간근: 단위의 연축 힘은 $0.1$에서 $10$~g이고, 단위가 켜지는 힘에 거의 선형으로 커진다. & \cite{Henneman1957, MilnerBrown1973rec} & 측정됨 \\
6 & 힘의 단계는 힘에 비례한다, $\Delta F/F\approx q-1$~\eqref{eq:sizeprinciple}: 부동소수점 & 힘의 등비수열에서 본문이 유도. 실험과 부합: 큰 힘에서는 같은 힘 증가에 새로 켜지는 단위가 훨씬 적다. & 이 장의 유도; \cite{MilnerBrown1973rec} & 이론 \\
7 & 힘의 퍼짐은 힘 자체에 비례해 커진다 & 엄지 근육의 수의적 등척성 힘: 힘의 표준편차는 평균 힘에 따라 선형으로 커진다; 전기 자극으로 같은 힘을 낼 때는 퍼짐이 일정하다. 단위 풀 모델: 힘 순서대로의 동원이 선형 증가를 준다. & \cite{Jones2002sdn} & 측정됨 \\
8 & 움직임의 매끄러움은 신호 의존 잡음의 결과이다 & 모델: 이런 잡음에서 끝점 퍼짐을 최소화하는 궤적이 측정된 눈과 팔의 궤적을 재현한다. 이 원인을 다른 원인과 구별하는 직접 실험은 없다. & \cite{HarrisWolpert1998} & 가설 \\
9 & 근육 쌍의 힘 차가 회전력을, 합이 강성을 정한다~\eqref{eq:antagonists} & 공동 긴장에 의한 임피던스 제어 이론. 사람 팔: 작은 밀기로 잰 각 관절의 강성은 그 관절의 회전력에 대략 비례한다; 공동 긴장의 비율이 다르면 손 강성 타원의 모양과 방향이 바뀐다. & \cite{Hogan1984, GomiOsu1998stiff} & 측정됨 \\
10 & 신장 센서는 자기 \nt{ACET} 뉴런을 흥분시키고, 억제성 중개 뉴런을 거쳐 길항근을 끈다(그림~\ref{fig:antagonists}) & 고양이 척수: 신장 센서 섬유로 흥분되는 단일 중개 뉴런이 길항근 \nt{ACET} 뉴런에 단일 시냅스 억제 전위를 만든다, 시냅스 지연 $0.28$--$0.42\ms$. 중개 뉴런의 신경전달물질(글리신)은 이 실험에서 확인하지 않았다. & \cite{Jankowska1972Ia} & 측정됨 \\
11 & 루프 지연: 척수 루프 $25$--$30\ms$, 피질을 거치면 1.5--3배, 눈을 거치면 $100$--$200\ms$ & 사람 팔, 관절 밀기: 근육의 짧은 응답은 $20$--$45\ms$, 긴 응답은 $45$--$105\ms$, 수의적 응답은 $120$--$180\ms$ 뒤. 움직이는 동안 보이는 손가락 위치를 옮기면: 보정은 평균 $160\ms$ 뒤. & \cite{Pruszynski2008rapid, SaundersKnill2003vis} & 측정됨 \\
12 & $\dd x/\dd t=-Gx(t-d)$는 $Gd<\pi/2$일 때 안정하다~\eqref{eq:delaystable}; 경계에서 주파수는 $1/(4d)$ & 지연 방정식의 근에 관한 정리. \texttt{models/\allowbreak muscle\_\allowbreak models.py}로 확인, $d=30\ms$: $3\,\text{s}$까지의 진폭은 초당 $G=40$에서 $3\cdot10^{-6}$, $G=50$에서 $0.13$, $G=55$에서 $38$(경계 $52$). & \cite{Hayes1950} & 이론 \\
13 & 생리적 떨림 $8$--$12$~Hz; 신장 루프는 그 원천 가운데 하나 & 측정 총설: 약 $10$~Hz 대역의 미세 떨림은 건강한 사람에게 있다; 그 기원은 복합적이다. 중추 진동, 단위 방전의 성질, 기계적 공진, 반사 루프의 공진이 있고, 조건에 따라 우세한 것이 다르다. & \cite{McAuleyMarsden2000} & 가설 \\
14 & 뇌는 몸의 내부 모델로 명령의 결과를 예측한다 & 사람이 물체를 손가락으로 집어 들고 팔을 움직인다: 관성, 점성, 탄성 부하에서 쥐는 힘이 루프 지연 없이 부하와 동시에 변한다. 곧 부하가 예측된다. 총설이 이런 실험을 정리한다; 뉴런에서 모델 자체를 직접 관찰한 것은 없다. & \cite{FlanaganWing1997grip, WolpertGhahramani2000} & 가설 \\
15 & 걷기, 숨쉬기, 씹기의 리듬은 일정한 입력에서 국소 네트워크가 만든다 & 실험 총설: 분리된 신경계(척수, 신경절)는 리듬적 입력도 근육의 되먹임도 없이 리듬 운동 패턴을 낸다. & \cite{MarderBucher2001} & 측정됨 \\
16 & 피로를 동반한 상호 억제~\eqref{eq:halfcentre}는 주기 $0.84\,\text{s}\approx2\tau_a$의 리듬을 준다(그림~\ref{fig:halfcentre}) & 정리: 상호 억제와 적응을 가지며 일정한 입력에서 안정 상태가 없는 네트워크는 반드시 진동한다. \texttt{models/\allowbreak muscle\_\allowbreak models.py}에 따른 계산($I=1$, $\beta=2$, $b=2$, $\tau=20\ms$, $\tau_a=0.4\,\text{s}$): 주기 $0.84\,\text{s}$. 실제 발생기가 바로 이렇게 만들어져 있다는 것은 보이지 않았다. & \cite{Matsuoka1985osc} & 모델 \\
\bottomrule
\end{longtable}}

\section{\ref{sec:pain}장. 통증}

통증 센서, 두 전선, 척수의 관문, 하행 음량 노브, 감작, 주의 포획은 직접 측정되었다. 관문과 감작의 식은 이 책의 단순화된 모델이다. 기계가 경보 음량을 고르는 규칙은 가설이다.

{\footnotesize
\setlength{\tabcolsep}{3pt}\renewcommand{\arraystretch}{1.15}
\begin{longtable}{>{\raggedright}p{0.5cm}>{\raggedright}p{4.0cm}>{\raggedright}p{5.6cm}>{\raggedright}p{2.3cm}>{\raggedright\arraybackslash}p{1.7cm}}
\toprule
No. & 명제 & 실험과 측정한 것 & 출처 & 상태 \\
\midrule\endfirsthead
\toprule
No. & 명제 & 실험과 측정한 것 & 출처 & 상태 \\
\midrule\endhead
1 & 높은 문턱값: 통증 센서마다 가열 문턱값은 $41^\circ$에서 $46$--$48^\circ$까지 & 단일 섬유 기록. 원숭이 피부: 느린(C) 센서의 중앙 문턱값 $41^\circ$, 빠른 제2형 센서 $46^\circ$, 제1형은 $53^\circ$ 이상. 사람 피부: 압력에도 응답하는 C 센서 $40.7^\circ$, 압력에 응답하지 않는 것 $48^\circ$. & \cite{Treede1995heat, Weidner1999cfib} & 측정됨 \\
2 & 통증 센서는 촉각 센서보다 훨씬 약하게 적응하고, 가벼운 손상 뒤에는 더 민감해진다; 강한 가열 뒤의 약화는 한 유형에서 측정됨 & 피부의 가벼운 화상($50^\circ$, $100\,\text{s}$): $5$--$10$~분 뒤 사람의 통증 문턱값과 원숭이 C 센서의 문턱값이 $2$--$6^\circ$ 내려간다; 다른 피부 센서에는 이런 이동이 없다. 빠른 제2형 센서는 강한 가열 뒤 응답이 억제된다. 촉각 센서와의 적응 비교는 일반적인 추정으로, 여기서 별도 실험으로 뒷받침되지는 않는다. & \cite{LaMotte1982hyper, Treede1995heat} & 측정됨 \\
3 & 두 전선: 빠른 전선 $5$--$30$~m/s, 느린 전선 $0.5$--$2$~m/s; 도착 시간 $t=L/v$~\eqref{eq:paintime} & 전도 속도: 원숭이의 빠른 통증 센서는 평균 $15$와 $25$~m/s(두 유형); 사람의 C 센서는 $0.8$--$1.0$~m/s. $L=1$~m이면 수십 밀리초와 약 1초. & \cite{Treede1995heat, Weidner1999cfib} & 측정됨 \\
4 & 통증은 두 번 온다: 먼저 빠르고 정확하게, 다음에 둔하고 오래 & 사람 아래팔 가열에 대한 반응 시간: 온도가 오르면 $1100$에서 $400\ms$로; 털 있는 피부의 강한 가열은 이중 통증을 주고, 손바닥은 주지 않는다. 첫 통증은 $6$~m/s보다 빠른 섬유만 나를 수 있다. 잠복기로 보면 손바닥에 없는 빠른 제2형 센서가 이에 해당한다. 역할 분담(“어디”와 “얼마나 나쁜지”)은 해석. & \cite{CampbellLaMotte1983first, Treede1995heat} & 측정됨 \\
5 & 관문: 척수 출력 뉴런이 통증과 촉각을 받고, 억제성 중개 뉴런은 촉각에 흥분한다(그림~\ref{fig:gate}) & 관문 도식은 이론으로 제안되었다. 생쥐, 유전적 표지와 뉴런 집단 제거: 소마토스타틴을 가진 흥분성 뉴런은 기계적 통증에 필요하다; 다이노르핀을 가진 억제성 뉴런은 촉각 센서의 입력이 그 뉴런에 닿지 못하게 한다. 이 뉴런이 없으면 가벼운 접촉이 통증을 일으킨다. & \cite{MelzackWall1965, Duan2014} & 측정됨 \\
6 & 촉각은 통증을 누그러뜨린다 & 사람, 손의 레이저 통증 펄스: 동시에 접촉하면 통증 평가와 펄스 에너지 구별 능력이 낮아진다; 한 피부 분절 안에서 접촉 지점과 통증 지점 사이 거리가 멀수록 효과가 약해진다. & \cite{Mancini2014touch} & 측정됨 \\
7 & 관문 이론의 가정: 강한 통증이 스스로 억제성 중개 뉴런을 꺼서 관문이 활짝 열린다 & 본문에서 원래 이론의 가정으로 표시했다. 생쥐 연구는 관문이 촉각을 통증 채널로 들이지 않는 방식을 기술한다; 통증에 의한 중개 뉴런의 꺼짐은 거기서 보이지 않았다. & \cite{MelzackWall1965} & 가설 \\
8 & 관문~\eqref{eq:gate}: 촉각 없이 문턱값 $0.18$, 촉각이 있으면 $0.86$, $g=2$에서 $0.11$; $\nu=1$에서 촉각은 출력을 $1$에서 $0.25$로 낮추고, $\nu=2$에서는 효과가 없다; 촉각 자체는 통과하지 못한다(그림~\ref{fig:gatecurves}) & 본문의 가중치로 \texttt{models/\allowbreak pain\_\allowbreak gate.py}에 따라 계산하면 이 수치가 모두 재현된다. & \texttt{models/\allowbreak pain\_\allowbreak gate.py} & 모델 \\
9 & 뇌간에서 내려오는 하행 경로는 관문의 이득을 양쪽으로 바꾼다 & 쥐, 연수, 꼬리 가열 중 기록: 어떤 뉴런은 꼬리를 빼기 전에 방전하고(“켜짐”), 다른 뉴런은 조용해진다(“꺼짐”); 이 영역의 약한 자극은 반사를 억제한다. 총설: 같은 회로가 통증 전달을 촉진하기도 억제하기도 하며, 오피오이드는 이를 통해 작용한다. & \cite{Fields1983rvm, Fields2004} & 측정됨 \\
10 & 통증이 줄어들 것이라는 기대는 오피오이드 채널을 통해 통증을 누그러뜨린다 & 급성 통증이 있는 사람들: 기대로 통증이 줄어든 사람에서는 오피오이드 수용체 차단이 통증을 다른 사람들 수준으로 되돌렸다; 다른 사람들에서는 차단이 통증을 바꾸지 않았다. & \cite{Levine1978} & 측정됨 \\
11 & 뇌는 중단의 이득에 따라 경보 음량을 고른다 & 음량 선택 규칙을 측정한 실험은 없다. & --- & 가설 \\
12 & 감작: 손상 뒤 관문 출력이 커지고 문턱값이 떨어지며, 일부는 척수에서 일어난다; 손상 자극이 멈춘 뒤에도 유지된다 & 쥐: 피부 손상 뒤 굽힘 반사의 문턱값이 떨어지고 응답이 커진다; 전기생리학적 분석은 증가의 일부가 척수 자체에서 생김을 보여 준다. 손상 자극은 자극 자체보다 오래가는 감도의 장기 변화를 일으킨다. 정확한 회복 시간은 본문에 없다. & \cite{Woolf1983} & 측정됨 \\
13 & 감작~\eqref{eq:sensitization}은 양의 되먹임이다; 루프가 강하면 경보는 치유 뒤에도 걸린 채 남을 수 있다 & 모델~\eqref{eq:sensitization}에서 나온다(장 끝의 연습문제). 치유 뒤 지속되는 통증이 바로 이런 형태로 걸린 루프임을 보인 실험은 없다. & 장 안의 유도 & 가설 \\
14 & 통증은 음의 보상이며, 그것으로 하는 학습은 시간차 오차를 따른다 & 영상 장치, 사람, 통증 전조의 사슬: 배쪽 선조체와 앞쪽 섬엽 피질의 활동이 시간차 학습 신호와 일치한다. 원숭이: \nt{DOPA} 뉴런 일부는 불쾌한 공기 분사와 그 전조에 흥분하고, 다른 일부는 억제된다. & \cite{Seymour2004, Matsumoto2009} & 측정됨 \\
15 & 통증은 진행 중인 작업을 중단시킨다 & 사람, 전기 통증 자극과 소리 시험: 통증 시작 $100\ms$ 뒤의 시험 응답은 뚜렷이 느려지고, $1.5\,\text{s}$ 뒤에는 대조군과 차이가 없다. 총설은 이런 실험을 주의 포획 모델로 정리한다. & \cite{Crombez1996disrupt, EcclestonCrombez1999} & 측정됨 \\
16 & 회피 반사는 척수에서 닫힌다 & 쥐: 후각 깊은 층의 뉴런이 개별 근육의 회피 반사 공간 패턴을 되풀이한다; 경수에서 역방향으로 흥분시킬 수 없으므로 척수 중개 뉴런이다. & \cite{Schouenborg1995wr} & 측정됨 \\
\bottomrule
\end{longtable}}

\section{\ref{sec:motorlearning}장. 운동 학습}

최소 제곱 규칙과 개별 오차 전선의 이점은 정리이다. 오차 전선 회로의 구조, 겹칠 때의 약화, 지연에 맞춘 적격 흔적, 후효과와 기술의 자발적 복귀는 측정되었다. 회로 전체가 지도 학습으로 동작하며 내부 모델을 만든다는 것은 유력한 가설이다. 두 과정 모델의 수치는 이 책의 모델에 따른 계산이다.

{\footnotesize
\setlength{\tabcolsep}{3pt}\renewcommand{\arraystretch}{1.15}
\begin{longtable}{>{\raggedright}p{0.5cm}>{\raggedright}p{4.0cm}>{\raggedright}p{5.6cm}>{\raggedright}p{2.3cm}>{\raggedright\arraybackslash}p{1.7cm}}
\toprule
No. & 명제 & 실험과 측정한 것 & 출처 & 상태 \\
\midrule\endfirsthead
\toprule
No. & 명제 & 실험과 측정한 것 & 출처 & 상태 \\
\midrule\endhead
1 & 규칙~\eqref{eq:lms}는 평균적으로 오차 제곱의 기울기를 따라 내려간다 & 적응 필터 이론의 결과: 입력과 출력 오차에 비례하는 보정은 평균 제곱 오차의 확률적 경사 하강이다. & \cite{WidrowHoff1960} & 이론 \\
2 & 출력마다 오차 전선이 있으면 속도는 출력 수와 무관하다; 하나의 공통 수로는 잡음 내는 뉴런 수에 따라 떨어진다(명제~\ref{prop:teachers}) & 선형 네트워크의 학습 곡선: 뉴런의 무작위 섭동에 의한 학습은 정확한 기울기보다 섭동받는 뉴런 수만큼 느리다. & \cite{Werfel2005} & 이론 \\
3 & 오차 전선 회로는 지도 학습으로 동작한다(명제~\ref{prop:errorwire}) & 회로 구조에서 유도한 이론. 7--10행의 실험은 이와 부합한다; 회로 전체가 바로 이렇게 동작한다는 것은 증명되지 않았다. & \cite{Marr1969, Albus1971} & 가설 \\
4 & 확장층: 작은 \nt{GLUT} 뉴런 약 $7\cdot10^{10}$개, 뇌의 나머지 전부보다 많다 & 녹인 조직 현탁액에서 핵 계수: 사람 소뇌에 뉴런 $69\cdot10^9$개, 뇌 전체는 $86\cdot10^9$개. 입체해석학: 노년 남성의 소뇌에 작은 뉴런 $101\cdot10^9$개. & \cite{Azevedo2009, Andersen1992cereb} & 측정됨 \\
5 & 입력의 차원을 높이면 원하는 관계가 선형이 된다 & 분할 수에 관한 정리: 문턱값을 가진 입력 $n$개의 가중합은 일반 위치에 있는 벡터 약 $2n$개의 임의의 표지를 구현한다. 소뇌가 바로 이것을 쓴다는 것은 3행 가설의 일부이다. & \cite{Cover1965} & 이론 \\
6 & 출력 \nt{GABA} 뉴런 약 $3\cdot10^7$개, 각각 $10^5$개 정도의 입력 & 사람 소뇌의 입체해석학: 출력 뉴런 $30.5\cdot10^6$개. 전자 현미경, 쥐: 출력 뉴런 하나에 확장층 입력 약 $1.75\cdot10^5$개. 출력 뉴런의 억제성 신경전달물질은 총설. & \cite{Andersen1992cereb, NapperHarvey1988, Ito2001} & 측정됨 \\
7 & 출력 뉴런마다 정확히 하나의 오차 전선, 초당 약 한 번, 매번 강력한 폭발 & 기록 총설: 각 출력 뉴런은 등반하는 \nt{GLUT} 입력 하나를 받으며, 이것이 $1$~Hz 정도의 빈도로 복합 스파이크를 일으킨다. & \cite{Ito2001} & 측정됨 \\
8 & 폭발의 확률은 움직임 오차의 방향에 따라 다르다 & 원숭이, 소뇌의 안구 운동 부위, 도약 안구 운동 적응: 시각 오차의 방향이 복합 스파이크의 확률을, 크기가 그 시각을 정한다; 폭발은 다음 시도의 단순 스파이크를 바꾸고 움직임을 세포의 선호 오차 방향으로 옮긴다. & \cite{Herzfeld2018} & 측정됨 \\
9 & 오차 폭발이 활성 입력과 겹치면 그 입력이 약해진다($-\eta e_i x_j$) & 실험 총설: 확장층 입력과 오차 전선을 함께 자극하면 그 입력이 장기적으로 약해진다; 한 입력만 자극하면 그렇지 않다. & \cite{Ito2001} & 측정됨 \\
10 & 흔적의 길이는 오차 지연에 맞춰져 있다: 지연이 $100\ms$이면 $100\ms$ 앞서 활성이던 입력이 가장 강하게 약해진다 & 절편과 살아 있는 생쥐: 안구 운동 학습을 맡는 부위에서 약화는 입력과 오차 사이 간격 약 $120\ms$에 좁게 맞춰져 있다. 이 과제에서 오차 신호가 오는 데 걸리는 시간이다; 다른 부위에서는 세포마다 다른 간격에 맞춰져 있다. & \cite{Suvrathan2016} & 측정됨 \\
11 & 회로는 몸의 내부 모델을 학습하는 적응 필터~\eqref{eq:adaptivefilter}이다 & 회로 구조에서 유도한 모델. 학습된 필터를 몸의 전달 함수로 측정한 직접 실험은 없다. & \cite{Fujita1982} & 가설 \\
12 & 예측은 자기 움직임의 결과를 빼므로 스스로를 간지럽힐 수 없다 & 사람: 자기 접촉은 같은 남의 접촉보다 덜 간지럽게 느껴진다; 영상 장치에서 체감각 피질은 자기 접촉에 더 약하게 응답하고, 움직임이 피부에 닿을 때 소뇌는 덜 활성이다. 예측을 빼는 설명은 가설. & \cite{Blakemore1998} & 가설 \\
13 & 외부 힘이 있으면 빗나감이 줄고, 힘을 없애면 팔이 반대쪽으로 빗나간다 & 사람이 힘의 장 속에서 로봇 손잡이를 잡고 뻗는다: 궤적이 직선으로 돌아온다; 장을 없애면 후효과는 처음 왜곡의 거의 거울상이다; 훈련하지 않은 공간 영역으로도 전이된다. & \cite{ShadmehrMussaIvaldi1994} & 측정됨 \\
14 & 두 과정~\eqref{eq:tworate}이 학습한다; 반대 섭동 뒤에 기술이 스스로 돌아온다 & 사람, 힘의 장, 이어서 짧은 반대 장과 오차를 고정한 시도: 보정이 스스로 옛 값으로 돌아간다; 두 과정 모델은 이를 재현하고 한 과정 모델은 재현하지 못한다. & \cite{Smith2006} & 측정됨 \\
15 & 그림~\ref{fig:tworate}의 수치: $200$번 움직임에 $0.84$(느린 과정 $0.63$); 반대 $6$번; 빠른 과정 $-0.53$, 느린 과정 $0.46$; $40$번 움직임에 $0.37$까지 복귀 & \texttt{models/\allowbreak motor\_\allowbreak learning.py}에 따른 계산, $A_f=0.92$, $B_f=0.1$, $A_s=0.996$, $B_s=0.02$: $0.840$($0.634$와 $0.205$), 움직임 $6$번, $-0.531$과 $0.457$, $40$번째 움직임에서 최고 $0.371$. & \texttt{models/\allowbreak motor\_\allowbreak learning.py} & 모델 \\
16 & 몸이 여러 속도로 변하므로 두 과정이 필요하다 & 이론: 수명이 여럿인 교란에 대한 최적 추정은 시간 상수가 다른 여러 필터를 주며, 자발적 복귀와 빨라진 재학습을 설명한다. & \cite{Kording2007} & 가설 \\
\bottomrule
\end{longtable}}

\section{\ref{sec:chemoutput}장. 화학적 출력}

심장으로 가는 두 가닥의 전선과 그 속도 차이, 호르몬 캐스케이드의 음의 되먹임, 분량 빈도에 의한 부호, 일주기 발진기와 빛에 의한 맞춤은 직접 측정되었다. $K<8$이라는 경계는 본문에서 유도한 제어 이론의 정리다. 캐스케이드의 되먹임 자체가 맥동을 만든다는 것은 이를 강하게 뒷받침하는 실험이 있는 가설이다.

{\footnotesize
\setlength{\tabcolsep}{3pt}\renewcommand{\arraystretch}{1.15}
\begin{longtable}{>{\raggedright}p{0.5cm}>{\raggedright}p{4.0cm}>{\raggedright}p{5.6cm}>{\raggedright}p{2.3cm}>{\raggedright\arraybackslash}p{1.7cm}}
\toprule
No. & 명제 & 실험과 측정한 것 & 출처 & 상태 \\
\midrule\endfirsthead
\toprule
No. & 명제 & 실험과 측정한 것 & 출처 & 상태 \\
\midrule\endhead
1 & 심장 페이스메이커는 스스로 1분에 약 $100$번 뛴다. 안정 시에는 브레이크가 밟혀 있어 박동수가 보통 $60$--$70$ 정도다 & 사람, 심장으로 가는 두 전선을 모두 차단: 고유 박동수는 안정 시 박동수보다 높고 나이가 들수록 줄며, 성인에서 분당 약 $100$회다. 따라서 안정 시에는 브레이크가 우세하다. 안정 시 박동수 $60$--$70$은 전형적인 값으로 따로 인용하지 않았다. & \cite{JoseCollison1970ihr} & 측정됨 \\
2 & 가속 페달 \nt{NORA}와 브레이크 \nt{ACET}은 저역 통과 필터이며 브레이크가 더 빠르다~\eqref{eq:gasbrake} & 개, 심장의 미주 신경 또는 교감 신경을 주파수 변조 자극하고 심방 박동수까지의 전달 함수를 측정: 두 경로 모두 저역 통과 필터이고, 교감 경로는 차단 주파수가 훨씬 낮으며 약 $1.7\,\text{s}$의 순수 지연이 더해진다. 특성은 평균 긴장도에 따라 달라지므로 선형식~\eqref{eq:gasbrake}은 근사다. & \cite{Berger1989} & 측정됨 \\
3 & 브레이크가 빠른 것은 \nt{ACET}이 2차 전달자 없이 칼륨 채널을 열기 때문이고, \nt{NORA}는 연쇄 반응을 거친다 & 심방 세포의 막 조각: 아세틸콜린이 여는 칼륨 채널은 수용체와 결합한 G 단백질의 $\beta\gamma$ 소단위가 세포 안 2차 전달자 없이 연다. cAMP를 거치는 \nt{NORA} 경로는 여기서 인용하지 않았다. & \cite{Logothetis1987gbg} & 측정됨 \\
4 & 혈액은 약 1분 만에 몸을 한 바퀴 돈다. 혈류 속 \nt{ADRE}는 수용체가 있는 모든 기관에 닿는다 & 일반적인 크기 수준의 추정. 본문에서도 그렇게 서술했고 출처는 인용하지 않았다. & --- & 추정 \\
5 & 호르몬 캐스케이드는 음의 되먹임으로 감싸여 있다. 최종 호르몬이 앞 단계를 억제한다 & 실험 리뷰: 부신 피질 호르몬은 뇌하수체의 ACTH 분비와 시상하부의 방출 호르몬 분비를 빠른, 중간, 느린 세 시간 범위에서 억제한다. & \cite{KellerWood1984fb} & 측정됨 \\
6 & 같은 단 세 개는 $K<8$에서 안정하고~\eqref{eq:cascadestable}, 경계에서의 주기는 $2\pi\tau/\sqrt3\approx3.6\tau$ & 본문에서 유도: 주파수 $\sqrt3/\tau$에서 세 단은 위상을 $180^\circ$ 늦추고 진폭을 $8$분의 1로 줄인다. & 본문에서 유도 & 이론 \\
7 & 농도 오차는 $1/(1+K)$이므로 이런 조절기는 약 $10\,\%$보다 정확하게 유지하지 못한다(명제~\ref{prop:cascade}) & 비례 되먹임에 대한 6행의 따름정리. 실제 축에는 여러 단계와 여러 시간 범위에 되먹임이 있으므로(5행), 이 경계는 모델~\eqref{eq:cascade}에 대한 것이며 측정된 것이 아니다. & 본문에서 유도 & 이론 \\
8 & $K=4$와 $7$에서 농도가 자리 잡고, $K=12$에서는 주기 $3.7$--$3.9\,\tau$의 일정한 맥동(그림~\ref{fig:cascade}) & \texttt{models/\allowbreak chem\_\allowbreak models.py}로 계산: $K=12$에서 연속된 주기 $3.9$, $3.7$, $3.8$, $3.7\,\tau$; $K=4$에서 계산 끝의 진폭 $0.003$. & \texttt{models/\allowbreak chem\_\allowbreak models.py} & 모델 \\
9 & 스트레스 호르몬은 약 1시간에 한 번 맥동으로 분비된다. 맥동은 별도의 발진기 없이 되먹임 자체가 만든다 & 쥐, 방출 호르몬을 일정하게 주입: ACTH와 코르티코스테론이 생리적인, 거의 1시간 주기로 오르내린다. 맥동에 시상하부의 리듬 입력은 필요 없다. 지연된 되먹임이 있는 뇌하수체--부신 모델이 이런 진동을 낸다. & \cite{Walker2012glc, Walker2010} & 가설 \\
10 & 수신자는 농도가 아니라 분량의 빈도를 읽는다 & 생식샘 자극 호르몬 방출 호르몬을 스스로 분비하지 못하는 원숭이: 일정한 주입은 뇌하수체 호르몬 분비를 회복시키지 못하고, 1시간에 한 번 분량으로 주면 회복된다. 다시 일정한 주입으로 바꾸면 반응이 다시 사라진다. 고역 통과 필터로서의 수용체 피로는 해석이다. & \cite{Belchetz1978} & 측정됨 \\
11 & 분량의 빈도가 다르면 같은 샘의 서로 다른 세포에 다르게 작용한다 & 같은 원숭이: 분량을 1시간에 $2$--$5$회로 늘리면 두 뇌하수체 호르몬이 모두 줄고, $3$시간에 1회로 줄이면 하나(LH)는 줄고 다른 하나(FSH)는 늘어난다. 따라서 둘의 비는 빈도가 정한다. & \cite{Wildt1981gnrh} & 측정됨 \\
12 & 일주기 리듬은 몇 시간 지연되는 세포 안의 음의 되먹임이 만든다 & 리뷰: 시계 유전자의 단백질이 지연을 두고 자기 유전자의 전사를 억제한다. 전사와 번역으로 이루어진 루프가 약 하루의 주기를 낸다. & \cite{TakahashiJS2017} & 측정됨 \\
13 & 주 발진기는 수만 개의 뉴런으로 이루어진 무리이며 몸의 나머지를 동기화한다 & 리뷰: 시교차상핵이 주 일주기 발진기다. 배양한 개별 뉴런은 독립된 발진기이고, 네트워크가 이들을 동기화한다. 생쥐에서는 한쪽에 약 $10^4$개의 뉴런이 있다. & \cite{Welsh2010scn} & 측정됨 \\
14 & 사람의 고유 주기는 $24.18$시간 & 하루와 분리된 일정 아래 어두운 빛에서 지낸 사람: 멜라토닌, 체온, 코르티솔 리듬의 주기는 젊은 사람과 노인 모두 평균 $24.18$시간이고 편차가 좁다. & \cite{Czeisler1999} & 측정됨 \\
15 & 빛은 위상 고정 루프처럼 발진기를 맞춘다~\eqref{eq:pll}. 하루 약 $11$~min의 이동이 필요하다 & 사람, 하루 중 여러 시각에 $6.7$시간의 밝은 빛을 한 번 비춤: 진폭 $5$시간의 1형 위상 반응 곡선. 체온 최저점 이전에는 지연, 이후에는 앞당김. 식~\eqref{eq:pll}은 표준 모델이고, $11$~min $=0.18$시간은 산술이다. & \cite{Khalsa2003prc} & 측정됨 \\
16 & 빛이 없으면 맞춤이 없고 리듬은 고유 주기로 흘러간다 & 빛을 전혀 지각하지 못하고 일상 사회에서 사는 전맹인: 약 절반에서 멜라토닌과 코르티솔 리듬이 하루와 맞지 않는 자기 주기로 간다. & \cite{Sack1992blind} & 측정됨 \\
\bottomrule
\end{longtable}}

\section{\ref{sec:debugging}장. 기계 디버깅}

전극이 무엇을 듣는지, 활동의 저차원성, 딱딱한 배열을 쓴 인터페이스의 성능, 뇌와 디코더의 공동 학습, 자극으로 생기는 시각의 점은 동료 심사를 거친 실험에서 측정되었다. 데이터 흐름의 추정은 본문의 산술이다. 사람에게 이식한 Neuralink 장치에 대해서는 확인한 범위에서 주로 회사 자체가 데이터를 발표했다.

{\footnotesize
\setlength{\tabcolsep}{3pt}\renewcommand{\arraystretch}{1.15}
\begin{longtable}{>{\raggedright}p{0.5cm}>{\raggedright}p{4.0cm}>{\raggedright}p{5.6cm}>{\raggedright}p{2.3cm}>{\raggedright\arraybackslash}p{1.7cm}}
\toprule
No. & 명제 & 실험과 측정한 것 & 출처 & 상태 \\
\midrule\endfirsthead
\toprule
No. & 명제 & 실험과 측정한 것 & 출처 & 상태 \\
\midrule\endhead
1 & 전극은 반경 약 100마이크로미터 안의 뉴런, 보통 한 개에서 몇 개의 스파이크를 듣고, 모양으로 구별한다 & 쥐, CA1 영역, 세포 안과 세포 밖 동시 기록: 세포 밖 스파이크의 모양은 세포 안 스파이크의 폭과 진폭을 따른다. 추정: 테트로드는 뉴런 $60$--$100$개를 구별할 수 있지만 실제로는 약 여섯 개를 골라낸다. & \cite{Henze2000extra} & 측정됨 \\
2 & 뇌에는 뉴런 $8.6\cdot10^{10}$개가 있고, 프로브는 대략 1억분의 1을 듣는다 & 녹인 조직의 현탁액에서 핵을 셈: 사람 뇌의 뉴런 $86\cdot10^9$개. 비율은 본문의 산술. & \cite{Azevedo2009} & 측정됨 \\
3 & 운동 피질 뉴런 수백 개의 활동은 공통 변수 열 개 남짓에 들어간다 & 원숭이 운동 피질 기록의 리뷰: 집단 활동은 많은 뉴런이 공유하는 소수의 잠재 변수로 기술된다. & \cite{Gallego2017} & 측정됨 \\
4 & 이웃의 잡음은 상관되어 있고, 뉴런 100개는 1000개와 거의 같은 양을 전달한다(명제~\ref{prop:pooling}) & 원숭이 시각 피질: 이웃 뉴런의 잡음 상관 계수는 약 $0.1$이며, 평균의 이득은 뉴런 100개에서 포화된다. 정보를 제한하는 상관의 이론. & \cite{Zohary1994, MorenoBote2014} & 측정됨 \\
5 & 원시 흐름 $2\cdot10^8$~bit/s 대 스파이크 $8\cdot10^4$~bit/s~\eqref{eq:probedata} & 스파이크 흐름의 용량~\eqref{eq:spikeentropy}에 따른 본문의 산술. 디지털화 매개변수는 실제 값이다: Neuralink 칩은 채널당 $19.3$~kHz, $10$비트. & 이 장의 유도; \cite{Rieke1997, Musk2019} & 이론 \\
6 & Neuralink의 첫 시스템: 칩이 증폭하고 디지털화하며, 원시 흐름은 케이블로 나가고, 스파이크는 외부 프로그램이 골라낸다. 칩 위의 스파이크 검출, 밀폐된 본체, 무선 통신과 무선 전력은 회사 발표에 따르면 사람용 장치에 있다 & 회사 논문: 원시 흐름 전체가 USB-C 케이블로 나가고, 스파이크는 칩 밖의 프로그램이 실시간으로 골라낸다. 탑재 압축, 무선 전력과 전송은 임상 장치의 계획으로 언급되었다. 사람에게 이식한 장치는 회사 발표뿐이다. 칩 위에서 스파이크를 골라내야 한다는 것은 \eqref{eq:probedata}에서 나온 본문의 결론이다. & \cite{Musk2019}; Neuralink 보고, 동료 심사 논문 없음 & 측정됨(회사 논문); 사람은 회사 보고 \\
7 & 실 $96$가닥에 전극 최대 $3072$개, 실마다 $32$개; 실은 로봇이 혈관을 피해 꽂는다 & 회사 논문: 실 $96$가닥에 전극 $3072$개, 실마다 $32$개; 로봇이 1분에 실 $6$가닥(전극 $192$개)을 꽂는다; 배열을 오래 이식한 쥐에서 전극의 최대 $70\,\%$가 스파이크를 듣는다. & \cite{Musk2019} & 측정됨(회사 논문) \\
8 & 사람에서는 실 $64$가닥에 약 1000채널; 실 일부가 밀려났다; 커서는 약 $10$~bit/s & 회사 발표뿐. PubMed 검색에서 사람 대상 결과를 담은 동료 심사 논문을 찾지 못했다. 본문의 단서와 일치한다. & Neuralink 보고; 동료 심사 논문 없음 & 측정됨(회사 보고) \\
9 & 전극 100개 남짓의 배열을 쓴 인터페이스가 커서를 조작한다 & 마비 환자, 일차 운동 피질에 전극 $96$개: 손상 3년 뒤에도 팔을 움직이려는 의도가 스파이크를 바꾼다. 디코더가 “신경 커서”(메일, TV), 의수 손과 로봇 팔 조작을 가능하게 한다. & \cite{Hochberg2006} & 측정됨 \\
10 & “생각 속의 손”으로 글쓰기: 분당 약 $90$자, $8$~bit/s 미만 & 손이 마비된 사람, 손으로 쓰려는 시도, 순환 신경망 디코더: 실시간으로 분당 $90$자, 정확도 $94.1\,\%$, 자동 교정 뒤 $99\,\%$ 이상. 비트 추정은 본문의 산술. & \cite{Willett2021} & 측정됨 \\
11 & 말하려는 시도를 분당 약 $60$단어로 글로 바꾼다 & 또렷한 말을 잃은 사람, 피질의 배열: 분당 $62$단어; 어휘 $50$단어에서 단어 오류 $9.1\,\%$, 어휘 $125\,000$단어에서 $23.8\,\%$. & \cite{Willett2023} & 측정됨 \\
12 & 인터페이스는 용량의 작은 일부만 꺼낸다: 뉴런 하나는 초당 수십 비트, 100개는 수천 비트(명제~\ref{prop:probe}) & 상한~\eqref{eq:spikeentropy}은 이론; 8행, 10행과의 비교는 본문의 결론. 병목이 전선이 아니라 저차원성과 부호에 대한 무지라는 것은 직접 실험으로 확인되지 않았다. & \cite{Rieke1997} & 이론 \\
13 & 뇌와 디코더는 서로에게서 배운다 & 원숭이가 커서를 조작; 디코더가 닫힌 루프에서 맞춰지는 동안 뉴런도 다시 학습한다: 정확도가 오르고, 기술이 유지되며, 자기 팔의 움직임에 덜 방해받는다. 학습 속도를 떼어 놓으라는 조언은 공학 규칙이며, 본문에 별도 실험은 없다. & \cite{Orsborn2014coad} & 측정됨 \\
14 & 전극을 통한 전류는 주변 뉴런과 전선을 유형을 가리지 않고 흥분시킨다 & 생쥐와 고양이 피질, 2광자 칼슘 기록: 약한 전류조차 수십 마이크로미터 지름의 부피에서 축삭을 직접 흥분시켜, 최대 수 밀리미터 떨어진 드문드문한 뉴런을 흥분시킨다. 즉 흥분은 비선택적이지만 빽빽하지는 않다. & \cite{Histed2009stim} & 측정됨 \\
15 & 시각 피질 자극은 점을 켠다. 최대 $16$개 점을 동시에 켜면 몇몇 글자와 물체의 경계를 알아볼 수 있다 & 시각 장애인, 시각 피질에 전극 $96$개를 $6$개월 이식: 전극 하나의 문턱값 $66.8\pm36.5$~\textmu A; 무리를 동시에 자극하면(최대 $16$전극, 자극기의 한계) 문턱값이 낮아지고 구별 가능한 무늬가 생기며, 몇몇 글자와 물체의 경계를 알아본다. & \cite{Fernandez2021} & 측정됨 \\
16 & Neuralink의 두 번째 방향은 시각 피질에 신호를 넣는 장치다 & 개발 중이라는 회사 발표뿐. 작동에 관한 동료 심사 데이터는 찾지 못했다. & Neuralink 보고 & 가설 \\
17 & 브로드캐스트 신경전달물질의 농도는 조직 안의 화학 센서로 잴 수 있다 & 쥐 뇌 절편, 탄소 섬유 고속 순환 전압전류법: 세로토닌의 방출과 재흡수를 측정; 물질이 시냅스 밖으로 나간다. & \cite{Bunin1998} & 측정됨 \\
\bottomrule
\end{longtable}}

\section{\ref{sec:eetypes}장. 계측기로서의 뉴런}

개별 세포의 동작 모드는 전극을 통한 전류와 전압 기록으로 직접 측정되었다. 적분기의 수학과 부류 구분은 고전 이론이다. 뇌가 모드 재구성을 계산에 쓴다는 것은 여전히 가설이다.

{\footnotesize
\setlength{\tabcolsep}{3pt}\renewcommand{\arraystretch}{1.15}
\begin{longtable}{>{\raggedright}p{0.5cm}>{\raggedright}p{4.0cm}>{\raggedright}p{5.6cm}>{\raggedright}p{2.3cm}>{\raggedright\arraybackslash}p{1.7cm}}
\toprule
No. & 명제 & 실험과 측정한 것 & 출처 & 상태 \\
\midrule\endfirsthead
\toprule
No. & 명제 & 실험과 측정한 것 & 출처 & 상태 \\
\midrule\endhead
1 & 공진기: 전압 변화에 맞서는 느린 전류가 뉴런을 몇 Hz에서 수십 Hz에 최대가 있는 대역 통과 필터로 만든다(\ref{sec:resonator}절) & 절편의 \nt{GLUT} 뉴런에 주파수가 서서히 오르는 사인파 전류를 넣고 임피던스 $|Z(f)|$를 측정: $33^\circ$C에서 $2$--$5$~Hz에 최대($38^\circ$C에서 약 $7$~Hz). 공진은 느린 칼륨 전류와 양이온 전류가 만들고 지속성 나트륨 전류가 키운다. 리뷰는 같은 기법을 여러 세포 부류에 적용한다 & \cite{HuVervaekeStorm2002,HutcheonYarom2000} & 측정됨 \\
2 & 느린 전류는 인덕턴스처럼 행동하고, 막의 커패시턴스와 함께 공진 회로를 이룬다 & 전류에 대한 축삭의 문턱 아래 반응을 측정하고 선형화한 컨덕턴스 방정식으로 기술했다. 등가 회로에는 값이 전압에 따라 달라지는 “현상론적” 인덕턴스가 들어 있다 & \cite{Mauro1970} & 이론 \\
3 & 되먹임이 있는 무리의 시상수 $\tau_{\text{eff}}=\tau/(1-w)$~\eqref{eq:perfectint}; $\tau=0.1$~s, $w=0.995$에서 $20$~s & 본문에서 무리의 선형 방정식으로 유도. 눈 위치 유지의 메커니즘으로는 이론 연구에서 제안되었다 & 본문에서 유도; \cite{Seung1996} & 이론 \\
4 & 눈 위치 적분기는 개별 세포의 성질이 아니라 네트워크로 입력을 유지한다 & 물고기에서 눈 위치를 유지하는 핵의 세포 안 기록: 빠르게 눈을 돌린 뒤 뉴런의 발화율이 계단처럼 바뀌어 유지된다. 세포 하나에 짧은 전류를 넣으면 잠깐 이동할 뿐이고, 전압 계단은 문턱 아래에서도 유지된다. 이를 유지하는 것은 시냅스 입력이다 & \cite{Aksay2001} & 측정됨 \\
5 & 무누설 적분기에는 학습에 의한 끊임없는 조정이 필요하다. 어긋나면 잊거나 포화로 치닫는다 & 물고기를 눈 위치의 $\pm$에 비례하는 시각 되먹임으로 훈련: 유지가 불안정해지거나 누설되고, 뉴런의 시상수가 한 자릿수 넘게, $<1$~s까지 떨어졌다. 평범한 시각 되먹임은 안정성을 서서히 되돌렸다 & \cite{Major2004} & 측정됨 \\
6 & 쌍안정 뉴런: 짧은 입력이 오래가는 플래토를 켜고 억제가 끈다. 이 성질은 세로토닌이 켠다(\ref{sec:bistable}절) & 고양이에서 근육을 제어하는 \nt{ACET} 뉴런의 세포 안 기록: 짧은 시냅스 흥분이 뉴런을 오래가는 발화로 옮기고 짧은 억제가 재설정한다. 척수를 절단한 고양이에서는 세로토닌 전구체를 주입한 뒤 이 상태가 나타난다 & \cite{Hounsgaard1988} & 측정됨 \\
7 & 두 부류: 적분기(발화율이 0부터 커진다)와 공진기(문턱에서 곧바로 유한한 발화율) & 부류는 분기 이론에서 유도되었다. 실험: 피질 절편에서 \nt{GLUT} 뉴런은 얼마든지 낮은 발화율로 발화를 시작하고(1형), 빠른 \nt{GABA} 뉴런은 유한한 발화율로 도약한다(2형) & \cite{Izhikevich2007,Tateno2004} & 측정됨 \\
8 & 같은 뉴런이 적분기 또는 동시성 검출기가 된다. 억제가 합산 창을 짧게 한다 & 해마 절편에서 흥분성 입력이 문턱까지 더해지는 창은 $2$~ms보다 짧다. 흥분 직후에 오는 피드포워드 억제가 창을 짧게 한다. 억제가 약한 가지돌기에서는 창이 넓다 & \cite{Pouille2001} & 측정됨 \\
9 & 축삭으로 된 지연선(표~\ref{tab:eetypes}) & 원숭이올빼미에서 동시성 검출기로 가는 축삭의 지연을 측정: 축삭 길이가 다르면 지연이 다르고, 검출기는 소리 도달 시간차에 맞춰져 있다 & \cite{Carr1990} & 측정됨 \\
10 & 깨어 있을 때는 페이스메이커, 잘 때는 조용한 세포 & \nt{SERO} 뉴런은 낮에는 거의 규칙적으로 발화하고, 서파 수면에서 느려지며, 렘수면에서는 거의 조용하다(자세한 내용은 \ref{sec:sleep}장) & \cite{JacobsAzmitia1992,McGintyHarper1976} & 측정됨 \\
11 & 소자는 이온 채널과, 그 채널을 조정하는 브로드캐스트 채널이 정한다(명제~\ref{prop:eetypes}) & 작은 네트워크에서 입증: 조절 물질이 뉴런의 고유 성질과 시냅스 세기를 바꿔, 같은 회로가 서로 다른 리듬을 낸다 & \cite{Marder2012} & 측정됨 \\
12 & 뇌는 모드 재구성을 계산에 쓴다(명제~\ref{prop:eetypes}) & 과제를 푸는 데 세포 모드의 재구성이 필요했다는 직접 실험은 없다. 조절 물질이 회로의 출력을 바꾸는 실험만 있다 & \cite{Marder2012} & 가설 \\
\bottomrule
\end{longtable}}

\section{\ref{sec:dendrites}장. 가지돌기 개수}

부류의 구조와 입력 개수는 해부학과 전기 기록으로 측정되었다. 추정~\eqref{eq:dendriteload}은 커패시터의 간단한 물리이고, 입력 개수에 대한 계산적 설명은 이론과 모델에 기댄다.

{\footnotesize
\setlength{\tabcolsep}{3pt}\renewcommand{\arraystretch}{1.15}
\begin{longtable}{>{\raggedright}p{0.5cm}>{\raggedright}p{4.0cm}>{\raggedright}p{5.6cm}>{\raggedright}p{2.3cm}>{\raggedright\arraybackslash}p{1.7cm}}
\toprule
No. & 명제 & 실험과 측정한 것 & 출처 & 상태 \\
\midrule\endfirsthead
\toprule
No. & 명제 & 실험과 측정한 것 & 출처 & 상태 \\
\midrule\endhead
1 & 막의 단위 면적당 커패시턴스 $c_m\approx1$~\textmu F/cm$^2$ & 뉴런 세포체에서 막을 피펫으로 떼어 내고 전압 계단을 걸어, 용량성 전류의 감쇠로 전체 커패시턴스를 구한 뒤 면적으로 나눴다: 세 부류의 뉴런에서 $0.9$~\textmu F/cm$^2$ & \cite{Gentet2000} & 측정됨 \\
2 & 충전 시간 $t\approx c_mA\Delta V/I$~\eqref{eq:dendriteload}: 큰 뉴런에는 동시 입력 수십 개가 필요하고, 작은 뉴런에는 큰 시냅스 하나면 충분하다 & 커패시터 충전 법칙에 의한 추정. 수치($20$과 $300$~pF, $0.1$과 몇 nA)는 크기 수준 & 본문에서 유도 & 이론 \\
3 & 거대 시냅스 하나는 몇 나노암페어의 전류를 주어 작은 뉴런을 곧바로 문턱값까지 올린다 & 절편에서 말단과 수신 세포를 동시 기록: 시냅스 하나의 전류 $2$--$13$~nA, 입력 스파이크마다 문턱 이상의 반응, $36^\circ$C에서 지연 약 $0.4$~ms; 수신 세포는 \nt{GLYC}를 낸다 & \cite{Borst1995,BorstSoria2012} & 측정됨 \\
4 & 가지돌기 0개: 촉각과 통증의 감각 뉴런에서 스파이크는 세포체를 거치지 않고 센서에서 척수로 간다 & 구조(세포체 옆에서 둘로 갈라지는 축삭)는 해부학으로 확립되었다. 전도에 세포체의 흥분성이 필요 없다는 것은 이런 뉴런의 검증된 모델에서 보였다: 흥분성 세포체가 없어도 스파이크는 분기점을 똑같이 확실하게 통과한다 & \cite{AmirDevor2003} & 이론 \\
5 & 가지돌기 1개: 후각 뉴런은 수용체 한 종류를 지니고 입력 선로 하나를 전달한다 & 수용체 유전자 실험의 리뷰: 후각 뉴런 하나는 큰 유전자군 가운데 수용체 유전자 하나를 발현한다 & \cite{Mombaerts2004} & 측정됨 \\
6 & 가지돌기 2개: 소리 방향 검출기에서 왼쪽 귀와 오른쪽 귀의 입력은 서로 다른 가지돌기에 붙는다 & 리뷰에 모은 포유류의 해부학과 기록: 두 귀의 입력이 반대쪽을 향한 두 가지돌기에 나뉘어 있다 & \cite{Grothe2010} & 측정됨 \\
7 & 입력을 두 가지돌기에 나누면 AND가 더 엄격해진다 & 모델에서 입증: 한 가지돌기의 포화~\eqref{eq:shunt} 때문에 한쪽 입력만으로는 문턱값에 이르지 못하고, 동시성 검출이 개선된다. 입력 배치를 바꾼 직접 실험은 없다 & \cite{AgmonSnir1998} & 이론 \\
8 & 운동 학습 회로에 작은 \nt{GLUT} 뉴런 약 $7\cdot10^{10}$개, 각각 입력 $\sim4$개 & 등방성 분획법에 의한 세포 계수: 사람 뇌의 이 부분에 뉴런 약 $690$억 개. 살아 있는 쥐에서 이런 세포 하나를 기록: 촉각에 소수의 입력에서 오는 이산 전류의 버스트가 도착하고, 이들이 더해지면 세포가 발화한다 & \cite{Azevedo2009,Chadderton2004} & 측정됨 \\
9 & 입력 $n$개인 문턱 소자는 무작위 패턴을 최대 $P_{\max}\approx2n$개 분리한다~\eqref{eq:cover}; 운동 학습 회로의 출력 뉴런에서 이 경계는 $\sim2\cdot10^5$, 현실적인 추정은 $\sim4\cdot10^4$개의 대응 & 경계는 조합 기하학의 고전적 결과. 현실적인 추정은 양의 가중치만 쓰고 잡음 여유를 둔 용량 이론을, 이 뉴런의 측정된 입력 수에 적용한 것이다 & \cite{Cover1965,Brunel2004} & 이론 \\
10 & 입력이 적은 믹서는 입력을 확장하려고 필요하다. “네 개”는 최적에 가깝다 & 이론: 믹서 층이 주는 표현의 차원은 해부학적으로 관찰되는 입력 개수 근처에서 최대가 된다. 확장이라는 원래 가설은 고전 연구에 있다 & \cite{Marr1969,Albus1971,LitwinKumar2017} & 가설 \\
11 & 큰 나무: 입력 $10^4$--$10^5$개 & 전자 현미경 사진으로 시냅스 계수: 쥐의 운동 학습 회로 출력 \nt{GABA} 뉴런 하나에 시냅스 $\approx1.7\cdot10^5$개; 해마 \nt{GLUT} 뉴런에 흥분성 입력 약 $30\,000$개와 억제성 입력 $1\,700$개 & \cite{NapperHarvey1988,Megias2001} & 측정됨 \\
12 & 큰 나무의 가지는 각자 문턱값을 가진 별도의 하위 회로이며, 뉴런은 작은 다층 네트워크다 & 가는 가지돌기의 두 지점을 국소 자극: 같은 가지의 입력은 시그모이드로, 다른 가지의 입력은 선형으로 더해진다. 사람 피질 뉴런의 가지돌기에서 세포 하나가 선형 분리 불가능한 문제를 풀게 하는 등급형 칼슘 스파이크가 발견되었다. 모델: 뉴런 하나를 재현하려면 깊은 네트워크가 필요하다 & \cite{Polsky2004,Gidon2020,Beniaguev2021} & 측정됨 \\
13 & 출력 가지돌기: 망막 \nt{GABA} 뉴런의 가지 하나하나가 별도의 방향 검출기다 & 개별 가지의 2광자 칼슘 기록: 가지의 신호는 세포체에서 끝 쪽으로 가는 움직임에 선택적이지만 세포체의 전압은 그렇지 않다 & \cite{Euler2002} & 측정됨 \\
14 & 입력 개수는 과제를 따른다(명제~\ref{prop:dendrites}) & 부류의 구조는 측정되었다(3--13행). 입력 개수가 속도나 분류를 위해 선택되었다는 것은 개별 부류에 대해 이론과 모델로만 보였다 & \cite{LitwinKumar2017,AgmonSnir1998} & 가설 \\
\bottomrule
\end{longtable}}

\section{\ref{sec:sleep}장. 수면}

수면 모드, 재생, 시냅스 축소는 뉴런 기록, 미세투석, 전자 현미경으로 측정되었다. 수면 일정과 느린 시소는 이 책의 모델로 기술했고, 수면의 목적은 대부분 가설이다.

{\footnotesize
\setlength{\tabcolsep}{3pt}\renewcommand{\arraystretch}{1.15}
\begin{longtable}{>{\raggedright}p{0.5cm}>{\raggedright}p{4.0cm}>{\raggedright}p{5.6cm}>{\raggedright}p{2.3cm}>{\raggedright\arraybackslash}p{1.7cm}}
\toprule
No. & 명제 & 실험과 측정한 것 & 출처 & 상태 \\
\midrule\endfirsthead
\toprule
No. & 명제 & 실험과 측정한 것 & 출처 & 상태 \\
\midrule\endhead
1 & 잠자는 동안 피질 뉴런은 조용하지 않다. 바뀌는 것은 동작 모드다 & 쥐 피질 뉴런의 장시간 기록: 서파 수면에서 발화율은 0이 아니며, 활동 기간과 침묵 기간이 번갈아 온다. 오래 깨어 있은 뒤에는 모든 상태에서 발화율이 높고, 잔 뒤에는 낮다 & \cite{Vyazovskiy2009} & 측정됨 \\
2 & \nt{ACET}은 낮과 렘수면에서 높고 서파 수면에서 낮다(표~\ref{tab:sleepmodes}) & 자유롭게 움직이는 고양이의 피질 미세투석: 각성 시 아세틸콜린 방출이 서파 수면보다 $100$--$175\%$ 높고, 렘수면에서는 조용한 각성과 거의 같은 수준이다 & \cite{Marrosu1995} & 측정됨 \\
3 & \nt{NORA}와 \nt{SERO}는 서파 수면에서 잠잠해지고 렘수면에서 거의 조용하다 & 쥐의 개별 \nt{NORA} 뉴런과 고양이의 \nt{SERO} 뉴런을 수면 단계별로 기록: 발화율은 각성에서 최대, 서파 수면에서 낮고, 렘수면에서 거의 0이다 & \cite{AstonJonesBloom1981,McGintyHarper1976} & 측정됨 \\
4 & 렘수면에서 근육은 \nt{GABA}와 \nt{GLYC}를 통한 억제로 꺼진다 & 쥐에서 씹기 근육의 \nt{ACET} 뉴런 활동을 재고 차단제를 그 뉴런에 직접 주입: 두 유형의 \nt{GABA} 수용체와 \nt{GLYC} 수용체를 모두 차단해야만 마비가 풀린다 & \cite{BrooksPeever2012} & 측정됨 \\
5 & 수면 압력 $S$는 문턱값이 두 개인 커패시터다~\eqref{eq:sleeppressure}, $\tau_r=18.2$~h, $\tau_d=4.2$~h & 2과정 모델. 시상수는 EEG 서파 파워가 각성에 따라 커지고 수면 중에 줄어드는 방식에서, 문턱값의 모양은 자발적으로 깨는 시각에서 얻었다 & \cite{Borbely1982,Daan1984} & 이론 \\
6 & 기계는 스스로 $16.1$~h 각성과 $8.0$~h 수면에 이른다(그림~\ref{fig:twoprocess}) & 문턱값이 오르내리는 \eqref{eq:sleeppressure}에 따른 계산, 스크립트 \texttt{models/sleep\_models.py}: $16.1$~h 각성과 $7.9$--$8.0$~h 수면 & --- & 모델 \\
7 & $40$~h 동안 자지 않은 뒤 $S=0.90$이지만 수면은 $8.8$~h뿐이다. 빚은 길이가 아니라 깊이로 갚는다 & 모델: \texttt{models/sleep\_models.py}가 $S=0.90$과 $8.8$~h를 준다. 실험: $40.5$~h 동안 자지 않은 여덟 명에서 첫 회복 밤의 EEG 서파 파워가 평소보다 유의하게 높다 & \cite{Borbely1981} & 모델 \\
8 & 물리적으로 $S$는 아데노신이다 & 고양이 미세투석: 각성을 조절하는 영역 가운데 하나에서 세포 밖 아데노신이 깨어 있는 동안 오르고, 수면 박탈로 더 오르며, 자는 동안 떨어진다. $S$가 아데노신뿐이라는 것은 입증되지 않았다 & \cite{PorkkaHeiskanen1997,Dunwiddie2001} & 가설 \\
9 & 서파 수면에서 피질은 시소처럼 오간다: 수백 ms의 활동과 침묵, 1초에 한 번보다 드물게($<1$~Hz, 마취 상태에서 대개 $0.3$--$0.4$~Hz) & 마취한 고양이의 세포 안 기록: $0.8$--$1.5$~s의 탈분극과 긴 과분극, 리듬 $<1$~Hz, 대개 $0.3$--$0.4$~Hz. 같은 활동-침묵 교대가 자연 수면에서도 보인다(1행) & \cite{Steriade1993,Vyazovskiy2009} & 측정됨 \\
10 & 시소는 되먹임과 피로가 있는 무리다~\eqref{eq:updown}. $b=5$에서 주기 $1.1$~s(모델의 값), 활동은 시간의 약 $35\%$; $b=1$에서 고른 활동(그림~\ref{fig:updown}) & 계산, 스크립트 \texttt{models/sleep\_models.py}: 주기 $1.11$~s, 활동 시간 비율 $0.35$(정수 개 주기에 대한 평균); $b=1$에서 활동 비율 $1.0$. 모델의 주기는 측정값보다 짧다(9행) & --- & 모델 \\
11 & \nt{ACET}은 시소를 끄고 피질을 고른 작동으로 옮긴다 & 고양이에서 \nt{ACET} 뉴런 핵을 자극하자 피질 뉴런의 $79\%$에서 느린 리듬이 차단되고 고른 작동으로 바뀌었다. 주로 긴 과분극이 사라졌기 때문이다. 아세틸콜린은 피로를 만드는 느린 칼륨 전류를 억누른다 & \cite{SteriadeAmzica1993,McCormick1992} & 측정됨 \\
12 & $7$--$15$~Hz 리듬 폭발은 피질과 중계 핵 사이의 \nt{GLUT}--\nt{GABA} 루프가 만든다 & 페럿의 시각 중계 핵 절편에서, 폭발은 중계 세포가 스스로 흥분시키는 이웃 \nt{GABA} 핵의 억제에 대해 내는 반동 버스트로 생긴다 & \cite{vonKrosigk1993,SteriadeMcCormickSejnowski1993} & 측정됨 \\
13 & 하루의 재생은 빨리 감기로 진행된다: 해마에서 약 $20$배, 피질에서 $6$--$7$배 빠르다 & 서파 수면에서 해마 뉴런의 반복 순서가 시간적으로 압축되어 나타난다. 트랙을 달린 뒤 장소 뉴런의 순서가 같은 순서로, 약 $20$배 압축된 $\sim100$~ms의 짧은 폭발로 반복되었다. 피질에서는 $6$--$7$배 압축 & \cite{Nadasdy1999,LeeWilson2002,Euston2007} & 측정됨 \\
14 & 재생과 피질의 맞물림을 강화하면 기억이 좋아진다(인과 관계) & 쥐에서 학습 뒤 재생에 맞춘 자극으로 재생과 피질의 서파 및 리듬 폭발의 연결을 강화: 다음 날 기억이 높았고, 대조군은 우연 수준이었다 & \cite{Maingret2016} & 측정됨 \\
15 & 재생의 목적은 느린 기억의 섞어 배우기다 & 두 학습 시스템 이론과 인공 네트워크 계산: 순차 학습은 옛것을 망가뜨리고, 섞어 배우기는 그렇지 않다. 재생이 바로 이를 위해 필요하다는 뇌 직접 실험은 없다 & \cite{McClelland1995} & 가설 \\
16 & 잔 뒤 시냅스는 크기에 비례해 줄어들고, 큰 시냅스는 거의 바뀌지 않는다 & 생쥐 피질 시냅스 $6920$개의 3차원 전자 현미경: 잔 뒤 접촉 면적이 깨어 있은 뒤보다 $\sim18\%$ 작고, 감소는 크기에 비례하며, 크고 안정된 시냅스는 바뀌지 않는다 & \cite{deVivo2017} & 측정됨 \\
17 & 수면의 주된 과제는 가중치 재정규화다 & 시냅스 항상성 가설. 1행과 16행은 이와 부합하지만 이것이 주된 기능임을 증명하지는 않는다 & \cite{TononiCirelli2014} & 가설 \\
18 & 렘수면은 세계 모델의 벤치 시험이다 & 모드(표~\ref{tab:sleepmodes})는 측정되었다. 네트워크가 세계 모델을 돌린다는 것은 이론적 추측이며 실험은 없다 & \cite{Hobson2009} & 가설 \\
19 & 수면 중 뇌 세척: 열린 문제 & 한 실험: 잠자는 동안 세포 사이 공간이 $60\%$ 넓고 청소가 빠르다. 측정 방법이 다른 다른 실험: 잠자는 동안과 마취 상태에서 청소가 눈에 띄게 느리다. 결과가 서로 모순된다 & \cite{Xie2013,Miao2024} & 가설 \\
20 & 국소 수면: 깨어 있는 동물의 피질 일부가 잠깐 침묵으로 빠진다 & 오래 깨어 있은 쥐에서 피질 한 영역의 뉴런이 서파 수면에서처럼 잠깐 조용해지고, 국소 EEG에 서파가 나타난다. 그런 순간 쥐는 과제에서 더 자주 실수한다 & \cite{Vyazovskiy2011} & 측정됨 \\
\bottomrule
\end{longtable}}

\section{\ref{sec:livingmachine}장. 살아 있는 뉴런으로 만든 기계}

전극 배열 위 배양의 행동은 기록과 자극을 쓰는 직접 실험으로 측정되었다. 일제 발화의 메커니즘은 시냅스 고갈 모델과 미세 배양 실험에 근거하고, 접시 안의 도파민에 관한 주장은 단 몇 건의 실험에 근거하며, 그중 일부는 저자들 스스로도 시연에 불과하다고 본다.

{\footnotesize
\setlength{\tabcolsep}{3pt}\renewcommand{\arraystretch}{1.15}
\begin{longtable}{>{\raggedright}p{0.5cm}>{\raggedright}p{4.0cm}>{\raggedright}p{5.6cm}>{\raggedright}p{2.3cm}>{\raggedright\arraybackslash}p{1.7cm}}
\toprule
No. & 명제 & 실험과 측정한 것 & 출처 & 상태 \\
\midrule\endfirsthead
\toprule
No. & 명제 & 실험과 측정한 것 & 출처 & 상태 \\
\midrule\endhead
1 & 칩 위의 배양: 대부분 \nt{GLUT}, 표본에 따라 다른 더 적은 비율의 \nt{GABA}. 해마 배양에서는 약 $6\%$ & 전극 배열 위의 뉴런 수천 개 네트워크는 표준 실험이다(3행). 해마 배양의 \nt{GABA} 뉴런 비율은 GABA 염색으로 세었다. 뉴런의 약 $6\%$. 피질 배양에 대해서는 검증된 숫자를 제시하지 않는다 & \cite{Benson1994} & 측정됨 \\
2 & 오가노이드: 인간 줄기세포에서 키운 1밀리미터 정도의 3차원 신경 조직 & 줄기세포에서 여러 뇌 영역을 가진 3차원 오가노이드를 키웠다. 전구세포 층과 성숙한 뉴런이 있는 피질도 포함된다 & \cite{Lancaster2013} & 측정됨 \\
3 & 입력이 없으면 배양은 배양에 따라 1초에서 몇 분 간격으로 일제 발화한다 & 밀도 $3\,000$--$50\,000$ 뉴런의 피질 배양 $58$개를 5주 동안 기록했다(30분 기록 $963$건). 2주 뒤에는 거의 모든 배양에서 네트워크 전체의 일제 발화가 우세하다. 일제 발화 간격은 $1$--$300$~s이며 표본에 크게 의존한다 & \cite{Wagenaar2006} & 측정됨 \\
4 & 일제 발화는 신경전달물질 고갈로 끝나고, 간격은 $T_{\text{burst}}\approx\tau_{\text{rec}}\ln\frac{1}{1-x^*}$~\eqref{eq:burstinterval}. $2$--$4$~s는 $\tau_{\text{rec}}=0.8$~s를 쓴 모델의 추정값이고, 측정된 회복은 더 느리다 & 공식은 본문에서 유도했다. 모델: 고갈되는 시냅스를 가진 네트워크는 스스로 네트워크 전체의 일제 발화를 만든다. 뉴런 $4$--$30$개의 미세 배양 실험: 일제 발화의 길이는 이전 일제 발화 뒤의 시간에 의존하고, 신경전달물질 소포를 고갈시키면 일제 발화가 사라진다. 저자들의 결론은 일제 발화가 신경전달물질 저장량에 의해 제한된다는 것이다. 거기서 저장량 회복은 수십 초이며, 같은 공식으로 수십 초에서 몇 분의 간격이 나온다 & 이 장의 유도; \cite{Tsodyks2000,CohenSegal2011,Tsodyks1997} & 이론 \\
5 & “입을 다무는 교사”: 원하는 응답이 나오면 자극을 끄고, 응답이 굳어진다 & 자극 후 $50\pm10$~ms에 선택한 응답이 나타날 때까지 $0.3$--$1$~Hz로 배양을 자극하고, 그 뒤 $5$~분 동안 자극을 껐다. 몇 주기 뒤에는 원하는 응답이 곧바로 나왔다. 학습 곡선을 그렸다 & \cite{Shahaf2001} & 측정됨 \\
6 & 배양이 테니스를 친다. 세션 안의 유의미한 랠리 증가는 완전한 되먹임일 때만 & 자세한 내용은 \ref{sec:dishbrain}장과 그 부록. 자극 대신 침묵을 주어도 세션 끝의 배양은 되먹임이 없을 때보다 낫지만 효과가 더 작다. 세션 간 학습은 불안정하다 & \cite{Kagan2022} & 측정됨 \\
7 & 배양이 자기 입력을 예측하도록 학습한다는 것은 가설. “지각력”에 관한 과장된 표현은 반박되었다 & 자유 에너지 원리는 이론적 제안이며, 더 단순한 설명과 구별하는 직접 실험은 없다. 반론 논문에서 연구자 30명이 루프 안의 행동 변화는 지능도 감각도 보여 주지 않는다고 지적했다 & \cite{Friston2010,Balci2023} & 가설 \\
8 & 알맞게 고른 자극 패턴으로 배양이 가상의 몸을 목표로 움직인다 & 프로그램이 현재 결과에 따라 훈련 자극 패턴을 스스로 골랐다. 수십 분 만에 네트워크는 지정된 상태에 도달했고, $2$~시간 훈련 뒤 가소성은 $80$~분 동안 유지되었다. 같은 자극을 행동과 무관하게 주면 효과가 없었다 & \cite{Bakkum2008} & 측정됨 \\
9 & 저수지: 학습하는 것은 선형 판독 $y=W_{\text{out}}\mathbf r$~\eqref{eq:reservoir}뿐 & 저수지 컴퓨팅 이론: 감쇠하는 기억을 가진 임의의 비선형 시스템에 학습된 선형 출력을 더한다 & \cite{Maass2002,JaegerHaas2004} & 이론 \\
10 & 저수지로 쓴 오가노이드가 약 $78\%$의 정확도로 화자를 구분한다(저자들의 보고). 오차로 학습하는 것은 판독이고, 오가노이드의 연결은 교사 없이 바뀐다 & 여러 화자의 말소리를 전극 배열의 전류 패턴으로 오가노이드에 넣었다. 학습된 판독은 약 $78\%$의 정확도로 화자를 구분했다고 저자들이 보고한다. 저자들은 훈련 중 오가노이드의 기능적 연결성이 바뀌었다고도 보고하고, 이를 오가노이드 안의 교사 없는 학습이라 부른다 & \cite{Cai2023} & 측정됨 \\
11 & 뉴런 100만 개는 스파이크에 약 $10^{-4}$~W를 쓴다~\eqref{eq:dishpower} & 추정: 피질 에너지 예산~\eqref{eq:energy}에서 얻은 스파이크당 비용 $10^{-10}$~J에 스파이크 수를 곱했다. 배양 전력의 직접 측정은 없다 & 이 장의 유도; \cite{Attwell2001} & 이론 \\
12 & 섬광으로 주는 도파민: 케이지 도파민 $1$~mM, $240$회 반복, 유발 스파이크 수 증가 & 원격 오가노이드 플랫폼에서 광반응성 케이지 속 도파민($1$~mM)을 자극이 더 많은 스파이크를 일으킬 때 자외선($365$~nm, $0.8$~s)으로 방출했다. $240$단계(약 $5$~시간) 동안 스파이크 수는 전체적으로 늘었다. 저자들은 실험 하나로는 도파민과의 관련을 확정할 수 없다고 쓴다. 대조 실험은 없다 & \cite{Jordan2024} & 가설 \\
13 & 살아 있는 세포의 도파민이 유기 트랜지스터의 가중치를 오래, 가역적으로 바꾼다 & 도파민을 분비하는 세포를 유기 뉴로모픽 소자 위에서 키웠다. 전극에서의 도파민 산화가 채널 전도도를 바꾼다. 가중치의 장기 변화와 그 복귀를 보였다 & \cite{Keene2020} & 측정됨 \\
14 & 작동하는 \nt{DOPA} 뉴런을 가진 오가노이드가 존재한다. 그 도파민을 교사로 쓴 적은 없다 & 인간 줄기세포 오가노이드에서 전기적으로 활성인 성숙한 \nt{DOPA} 뉴런과 도파민 생산이 확인되었다. 이 도파민이 다른 네트워크를 가르치는 실험은 문헌에서 찾지 못했다 & \cite{Jo2016} & 측정됨 \\
15 & 막대의 균형을 잡는 오가노이드: 프로그램이 고른 학습 자극이 무작위보다 낫고, 굳히기는 없으며, \nt{GLUT} 시냅스를 통한다 & 생쥐 피질 오가노이드를 “카트 위의 막대” 과제와 닫힌 루프로 연결했다. 대부분의 오가노이드에서 강화 학습이 고른 신호가 무작위 신호나 무신호보다 좋은 결과를 냈다. $45$~분 휴식 뒤에는 향상이 남지 않았다. AMPA 수용체와 NMDA 수용체를 차단하면 향상이 사라졌다 & \cite{Robbins2026} & 측정됨 \\
16 & 제어 레지스터가 없으면 테스트 벤치의 네트워크는 무엇을 성공으로 볼지 스스로 정하지 못한다(명제~\ref{prop:livingmachine}) & 5--15행의 모든 실험에서 학습 신호는 실험자나 프로그램이 정한다. 배양이 스스로 성공 기준을 만든 실험은 없다. 이것은 부재에서 나온 결론이지 실험이 아니다 & --- & 가설 \\
\bottomrule
\end{longtable}}

\section{\ref{sec:dishbrain}장. DishBrain}

테스트 벤치의 구조와 게임 결과는 실험 논문 하나와 그 후속 연구에서 가져왔다. 잡음 공식과 좋은 설정 쪽으로의 표류는 본문에서 유도하고 이 책의 모델로 확인했으며, 접시 안의 학습 메커니즘은 밝혀지지 않았다.

{\footnotesize
\setlength{\tabcolsep}{3pt}\renewcommand{\arraystretch}{1.15}
\begin{longtable}{>{\raggedright}p{0.5cm}>{\raggedright}p{4.0cm}>{\raggedright}p{5.6cm}>{\raggedright}p{2.3cm}>{\raggedright\arraybackslash}p{1.7cm}}
\toprule
번호 & 주장 & 실험과 측정된 것 & 출처 & 상태 \\
\midrule\endfirsthead
\toprule
번호 & 주장 & 실험과 측정된 것 & 출처 & 상태 \\
\midrule\endhead
1 & 테스트 벤치: 칩당 생쥐 피질 세포 약 $800\,000$개 또는 인간 세포 약 $10^6$개. $8$~mm$^2$에 전극 $26\,000$개, 기록 최대 $1024$, 자극은 $8$개로 & 논문의 방법: MaxOne 칩, $8$~mm$^2$에 백금 전극 $26\,000$개, 초당 $20\,000$ 표본으로 $1024$채널 기록. 자극은 기술적으로 $32$개 전극까지 가능하나 독립적으로는 $8$개. 생쥐 배양은 약 $10$일째부터 사용 & \cite{Kagan2022} & 측정됨 \\
2 & $10$~ms 클록의 루프: 운동 영역의 스파이크가 라켓을 움직인다(그림~\ref{fig:dishloop}) & 스파이크를 $10$~ms 창으로 센다. 스파이크에서 응답 자극까지의 지연은 약 $5$~ms로, 장비의 버퍼가 정한다 & \cite{Kagan2022} & 측정됨 \\
3 & 입력: 장소 부호($8$개 전극 중 어느 것)와 초당 $4$--$40$회의 빈도 부호 & 본 실험에서 자극 빈도는 공이 먼 벽에 있을 때 $4$~Hz에서 라켓 앞 $40$~Hz까지 선형으로 오른다. 공 펄스의 진폭은 $75$~mV, 펄스는 직사각형 2상 & \cite{Kagan2022} & 측정됨 \\
4 & 출력 $\Delta p=c\,(n_\uparrow-n_\downarrow)$~\eqref{eq:dishreadout}, 창당 스파이크 수의 잡음 $1/\sqrt{RT}$ & 두 영역의 차이 규칙은 논문에 기술되어 있으나(활동에 따른 영역 가중치 포함) 정확한 공식은 없다. 잡음 추정은 푸아송 계수의 결과이며 본문에서 유도했다 & \cite{Kagan2022}; 본문의 유도 & 이론 \\
5 & 교사: 히트 뒤에는 예측 가능한 자극 $75$~mV, 초당 $100$회, $0.1$~s. 미스 뒤에는 예측 불가능한 자극 $150$~mV, 초당 $5$회, $4$~s 동안 무작위 장소와 시점에, 그 뒤 $4$~s 휴지 & 방법: 히트 뒤 $75$~mV, $100$~Hz, $100$~ms를 $8$개 전극 모두에 동시에. 미스 뒤 $150$~mV, $5$~Hz를 $4$~s 동안 무작위 전극에 무작위 시점으로, 그 뒤 $4$~s 자극 없음. 배양에 도파민은 없다 & \cite{Kagan2022} & 측정됨 \\
6 & 네트워크는 자기 입력이 예측 가능해지도록 행동한다 & 자유 에너지 원리는 저자들이 프로토콜을 이끌어 낸 이론적 제안이다. 실험은 이와 맞지만 더 단순한 메커니즘(7--8행)과 구별하지 못한다 & \cite{Friston2010,Kagan2022} & 가설 \\
7 & 실패가 네트워크를 흔들고 성공은 흔들지 않으면, 설정에 머무는 시간은 $\rho\propto1/P_{\text{miss}}$~\eqref{eq:dishdensity}(명\ref{prop:dishscramble}) & 대칭적 흔들기를 가진 마르코프 연쇄의 흐름 균형에서 나오며, 본문에서 유도했다. 배양에서 직접 확인한 것은 없다 & 본문의 유도 & 이론 \\
8 & `미스 때 뒤섞기' 모델이 학습한다: 히트 비율 $0.21\to0.38$. 모든 랠리 뒤 흔들기는 $\approx0.12$, 흔들기 없음은 $0.19$(그림~\ref{fig:dishmodel}) & 계산, 스크립트 \texttt{models/dishbrain\_model.py}, 모델 배양 $40$개에 랠리 $2000$번씩: $0.21\to0.38$; $0.13\to0.11$; $0.19\to0.19$ & --- & 모델 \\
9 & 랠리가 길어지는 것은 완전한 루프의 살아 있는 배양뿐이고, 대조군은 학습하지 않는다 & $20$~분 세션 $399$번: 생쥐 배양 $9$개(세션 $101$번), 인간 $11$개($138$), 배지만 있는 칩 $6$개($80$), 휴지 배양 $20$개($42$), 모의실험 $3$개($38$). 마지막 $15$~분의 평균 랠리 길이가 처음 $5$~분보다 긴 것은 살아 있는 배양뿐이다. 뉴런이 아닌 세포(HEK293T)와 배지만 있는 칩에서는 효과가 없다 & \cite{Kagan2022} & 측정됨 \\
10 & 세션 안의 유의미한 향상은 완전한 피드백일 때만. 자극 대신 침묵을 주면 세션 끝의 게임은 피드백 없음보다 낫지만 효과가 더 작다 & $3$일 동안 세션 $486$번: `자극'($n=27$), `침묵'($17$), `피드백 없음'($15$). 시간에 따른 향상은 `자극' 조건에서만 유의미하다. 세션 끝에 `침묵' 조건은 더 작은 효과로 `피드백 없음'을 앞섰다 & \cite{Kagan2022} & 측정됨 \\
11 & 효과는 작다. 네트워크가 오래 가는 학습을 하는지는 열린 질문 & 세션 안의 향상은 확실하지만, 저자들 스스로 말하듯 며칠에 걸친 세션 간 학습은 안정적으로 관찰되지 않았다 & \cite{Kagan2022} & 측정됨 \\
12 & 게임 중 네트워크는 임계 상태에 다가가고, 가까울수록 게임을 잘한다 & 같은 배양의 활동 사태 분석: 구조화된 입력은 네트워크를 임계 상태 쪽으로 옮기고, 게임 성적은 임계 상태와의 근접도와 상관된다. 그러나 행동의 결과에 대한 정보 없이 임계성만으로는 학습에 부족하다 & \cite{Habibollahi2023} & 측정됨 \\
13 & 배양의 `지각력'은 보이지 않았다 & 연구자 30명의 반론 논문: 닫힌 루프 안의 행동 변화는 지능도 감각도 증명하지 않는다. 그것을 보여 주는 실험은 없다 & \cite{Balci2023} & 가설 \\
14 & 제안하는 네 번째 대조 실험: 미스 뒤에 같은 길이와 같은 에너지의 예측 가능한 자극(\ref{sec:dishspec}절) & 이런 실험은 수행되지 않았다. 이 책의 제안이다 & --- & 가설 \\
\bottomrule
\end{longtable}}


\begin{thebibliography}{99}
\bibitem{Azevedo2009} F. A. C. Azevedo et al., \emph{Equal numbers of neuronal and nonneuronal cells make the human brain an isometrically scaled-up primate brain}, J. Comp. Neurol. \textbf{513}, 532 (2009).
\bibitem{Schultz1997} W. Schultz, P. Dayan, and P. R. Montague, \emph{A neural substrate of prediction and reward}, Science \textbf{275}, 1593 (1997).
\bibitem{SuttonBarto2018} R. S. Sutton and A. G. Barto, \emph{Reinforcement Learning: An Introduction}, 2nd ed. (MIT Press, Cambridge, MA, 2018).
\bibitem{Doya2002} K. Doya, \emph{Metalearning and neuromodulation}, Neural Networks \textbf{15}, 495 (2002).
\bibitem{Strata1999} P. Strata and R. Harvey, \emph{Dale's principle}, Brain Res. Bull. \textbf{50}, 349 (1999).
\bibitem{vanVreeswijk1996} C. van Vreeswijk and H. Sompolinsky, \emph{Chaos in neuronal networks with balanced excitatory and inhibitory activity}, Science \textbf{274}, 1724 (1996).
\bibitem{AstonJones2005} G. Aston-Jones and J. D. Cohen, \emph{An integrative theory of locus coeruleus-norepinephrine function: adaptive gain and optimal performance}, Annu. Rev. Neurosci. \textbf{28}, 403 (2005).
\bibitem{Beaulieu2011} J.-M. Beaulieu and R. R. Gainetdinov, \emph{The physiology, signaling, and pharmacology of dopamine receptors}, Pharmacol. Rev. \textbf{63}, 182 (2011).
\bibitem{Sparso2001} J. Spars{\o} and S. Furber (eds.), \emph{Principles of Asynchronous Circuit Design: A Systems Perspective} (Kluwer, Boston, 2001).
\bibitem{Carr1990} C. E. Carr and M. Konishi, \emph{A circuit for detection of interaural time differences in the brain stem of the barn owl}, J. Neurosci. \textbf{10}, 3227 (1990).
\bibitem{Wang2001} X.-J. Wang, \emph{Synaptic reverberation underlying mnemonic persistent activity}, Trends Neurosci. \textbf{24}, 455 (2001).
\bibitem{Hahnloser2002} R. H. R. Hahnloser, A. A. Kozhevnikov, and M. S. Fee, \emph{An ultra-sparse code underlies the generation of neural sequences in a songbird}, Nature \textbf{419}, 65 (2002).
\bibitem{Barlow1965} H. B. Barlow and W. R. Levick, \emph{The mechanism of directionally selective units in rabbit's retina}, J. Physiol. \textbf{178}, 477 (1965).
\bibitem{Carandini2012} M. Carandini and D. J. Heeger, \emph{Normalization as a canonical neural computation}, Nat. Rev. Neurosci. \textbf{13}, 51 (2012).
\bibitem{Pouille2001} F. Pouille and M. Scanziani, \emph{Enforcement of temporal fidelity in pyramidal cells by somatic feed-forward inhibition}, Science \textbf{293}, 1159 (2001).
\bibitem{Tsodyks1997isn} M. V. Tsodyks, W. E. Skaggs, T. J. Sejnowski, and B. L. McNaughton, \emph{Paradoxical effects of external modulation of inhibitory interneurons}, J. Neurosci. \textbf{17}, 4382 (1997).
\bibitem{Buzsaki2012} G. Buzs\'aki and X.-J. Wang, \emph{Mechanisms of gamma oscillations}, Annu. Rev. Neurosci. \textbf{35}, 203 (2012).
\bibitem{Rieke1997} F. Rieke, D. Warland, R. de Ruyter van Steveninck, and W. Bialek, \emph{Spikes: Exploring the Neural Code} (MIT Press, Cambridge, MA, 1997).
\bibitem{Mainen1995} Z. F. Mainen and T. J. Sejnowski, \emph{Reliability of spike timing in neocortical neurons}, Science \textbf{268}, 1503 (1995).
\bibitem{Softky1993} W. R. Softky and C. Koch, \emph{The highly irregular firing of cortical cells is inconsistent with temporal integration of random EPSPs}, J. Neurosci. \textbf{13}, 334 (1993).
\bibitem{Thorpe1996} S. Thorpe, D. Fize, and C. Marlot, \emph{Speed of processing in the human visual system}, Nature \textbf{381}, 520 (1996).
\bibitem{Zohary1994} E. Zohary, M. N. Shadlen, and W. T. Newsome, \emph{Correlated neuronal discharge rate and its implications for psychophysical performance}, Nature \textbf{370}, 140 (1994).
\bibitem{MorenoBote2014} R. Moreno-Bote et al., \emph{Information-limiting correlations}, Nat. Neurosci. \textbf{17}, 1410 (2014).
\bibitem{Gollisch2008} T. Gollisch and M. Meister, \emph{Rapid neural coding in the retina with relative spike latencies}, Science \textbf{319}, 1108 (2008).
\bibitem{OKeefe1993} J. O'Keefe and M. L. Recce, \emph{Phase relationship between hippocampal place units and the EEG theta rhythm}, Hippocampus \textbf{3}, 317 (1993).
\bibitem{Destexhe2003} A. Destexhe, M. Rudolph, and D. Par\'e, \emph{The high-conductance state of neocortical neurons in vivo}, Nat. Rev. Neurosci. \textbf{4}, 739 (2003).
\bibitem{Tsodyks1997} M. V. Tsodyks and H. Markram, \emph{The neural code between neocortical pyramidal neurons depends on neurotransmitter release probability}, Proc. Natl. Acad. Sci. USA \textbf{94}, 719 (1997).
\bibitem{Lisman1997} J. E. Lisman, \emph{Bursts as a unit of neural information: making unreliable synapses reliable}, Trends Neurosci. \textbf{20}, 38 (1997).
\bibitem{Attwell2001} D. Attwell and S. B. Laughlin, \emph{An energy budget for signaling in the grey matter of the brain}, J. Cereb. Blood Flow Metab. \textbf{21}, 1133 (2001).
\bibitem{Lennie2003} P. Lennie, \emph{The cost of cortical computation}, Curr. Biol. \textbf{13}, 493 (2003).
\bibitem{Yagishita2014} S. Yagishita et al., \emph{A critical time window for dopamine actions on the structural plasticity of dendritic spines}, Science \textbf{345}, 1616 (2014).
\bibitem{Werfel2005} J. Werfel, X. Xie, and H. S. Seung, \emph{Learning curves for stochastic gradient descent in linear feedforward networks}, Neural Comput. \textbf{17}, 2699 (2005).
\bibitem{Doya2000} K. Doya, \emph{Reinforcement learning in continuous time and space}, Neural Comput. \textbf{12}, 219 (2000).
\bibitem{Oja1982} E. Oja, \emph{Simplified neuron model as a principal component analyzer}, J. Math. Biol. \textbf{15}, 267 (1982).
\bibitem{BiPoo1998} G.-Q. Bi and M.-M. Poo, \emph{Synaptic modifications in cultured hippocampal neurons: dependence on spike timing, synaptic strength, and postsynaptic cell type}, J. Neurosci. \textbf{18}, 10464 (1998).
\bibitem{Williams1992} R. J. Williams, \emph{Simple statistical gradient-following algorithms for connectionist reinforcement learning}, Mach. Learn. \textbf{8}, 229 (1992).
\bibitem{Bayer2005} H. M. Bayer and P. W. Glimcher, \emph{Midbrain dopamine neurons encode a quantitative reward prediction error signal}, Neuron \textbf{47}, 129 (2005).
\bibitem{Shen2008} W. Shen, M. Flajolet, P. Greengard, and D. J. Surmeier, \emph{Dichotomous dopaminergic control of striatal synaptic plasticity}, Science \textbf{321}, 848 (2008).
\bibitem{Dabney2020} W. Dabney, Z. Kurth-Nelson, N. Uchida, C. K. Starkweather, D. Hassabis, R. Munos, and M. Botvinick, \emph{A distributional code for value in dopamine-based reinforcement learning}, Nature \textbf{577}, 671 (2020).
\bibitem{Matsumoto2009} M. Matsumoto and O. Hikosaka, \emph{Two types of dopamine neuron distinctly convey positive and negative motivational signals}, Nature \textbf{459}, 837 (2009).
\bibitem{BrombergMartin2010} E. S. Bromberg-Martin, M. Matsumoto, and O. Hikosaka, \emph{Dopamine in motivational control: rewarding, aversive, and alerting}, Neuron \textbf{68}, 815 (2010).
\bibitem{Menegas2018} W. Menegas, K. Akiti, R. Amo, N. Uchida, and M. Watabe-Uchida, \emph{Dopamine neurons projecting to the posterior striatum reinforce avoidance of threatening stimuli}, Nat. Neurosci. \textbf{21}, 1421 (2018).
\bibitem{Niv2007} Y. Niv, N. D. Daw, D. Joel, and P. Dayan, \emph{Tonic dopamine: opportunity costs and the control of response vigor}, Psychopharmacology \textbf{191}, 507 (2007).
\bibitem{Mohebi2019} A. Mohebi et al., \emph{Dissociable dopamine dynamics for learning and motivation}, Nature \textbf{570}, 65 (2019).
\bibitem{Hamid2016} A. A. Hamid et al., \emph{Mesolimbic dopamine signals the value of work}, Nat. Neurosci. \textbf{19}, 117 (2016).
\bibitem{Kim2020} H. R. Kim, A. N. Malik et al., \emph{A unified framework for dopamine signals across timescales}, Cell \textbf{183}, 1600 (2020).
\bibitem{Jeong2022} H. Jeong et al., \emph{Mesolimbic dopamine release conveys causal associations}, Science \textbf{378}, eabq6740 (2022).
\bibitem{Amo2022} R. Amo et al., \emph{A gradual temporal shift of dopamine responses mirrors the progression of temporal difference error in machine learning}, Nat. Neurosci. \textbf{25}, 1082 (2022).
\bibitem{Takahashi2017} Y. K. Takahashi et al., \emph{Dopamine neurons respond to errors in the prediction of sensory features of expected rewards}, Neuron \textbf{95}, 1395 (2017).
\bibitem{Sharpe2017} M. J. Sharpe et al., \emph{Dopamine transients are sufficient and necessary for acquisition of model-based associations}, Nat. Neurosci. \textbf{20}, 735 (2017).
\bibitem{Gardner2018} M. P. H. Gardner, G. Schoenbaum, and S. J. Gershman, \emph{Rethinking dopamine as generalized prediction error}, Proc. R. Soc. B \textbf{285}, 20181645 (2018).
\bibitem{Engelhard2019} B. Engelhard et al., \emph{Specialized coding of sensory, motor and cognitive variables in VTA dopamine neurons}, Nature \textbf{570}, 509 (2019).
\bibitem{Clements1992} J. D. Clements, R. A. J. Lester, G. Tong, C. E. Jahr, and G. L. Westbrook, \emph{The time course of glutamate in the synaptic cleft}, Science \textbf{258}, 1498 (1992).
\bibitem{Hasselmo2006} M. E. Hasselmo, \emph{The role of acetylcholine in learning and memory}, Curr. Opin. Neurobiol. \textbf{16}, 710 (2006).
\bibitem{JacobsAzmitia1992} B. L. Jacobs and E. C. Azmitia, \emph{Structure and function of the brain serotonin system}, Physiol. Rev. \textbf{72}, 165 (1992).
\bibitem{Schiller1992} P. H. Schiller, \emph{The ON and OFF channels of the visual system}, Trends Neurosci. \textbf{15}, 86 (1992).
\bibitem{Grothe2010} B. Grothe, M. Pecka, and D. McAlpine, \emph{Mechanisms of sound localization in mammals}, Physiol. Rev. \textbf{90}, 983 (2010).
\bibitem{Markram2004} H. Markram et al., \emph{Interneurons of the neocortical inhibitory system}, Nat. Rev. Neurosci. \textbf{5}, 793 (2004).
\bibitem{Joshi2016} S. Joshi, Y. Li, R. M. Kalwani, and J. I. Gold, \emph{Relationships between pupil diameter and neuronal activity in the locus coeruleus, colliculi, and cingulate cortex}, Neuron \textbf{89}, 221 (2016).
\bibitem{ServanSchreiber1990} D. Servan-Schreiber, H. Printz, and J. D. Cohen, \emph{A network model of catecholamine effects: gain, signal-to-noise ratio, and behavior}, Science \textbf{249}, 892 (1990).
\bibitem{YuDayan2005} A. J. Yu and P. Dayan, \emph{Uncertainty, neuromodulation, and attention}, Neuron \textbf{46}, 681 (2005).
\bibitem{Nassar2012} M. R. Nassar et al., \emph{Rational regulation of learning dynamics by pupil-linked arousal systems}, Nat. Neurosci. \textbf{15}, 1040 (2012).
\bibitem{BouretSara2005} S. Bouret and S. J. Sara, \emph{Network reset: a simplified overarching theory of locus coeruleus noradrenaline function}, Trends Neurosci. \textbf{28}, 574 (2005).
\bibitem{Hebb1949} D. O. Hebb, \emph{The Organization of Behavior: A Neuropsychological Theory} (Wiley, New York, 1949).
\bibitem{MacKay1952} D. M. MacKay and W. S. McCulloch, \emph{The limiting information capacity of a neuronal link}, Bull. Math. Biophys. \textbf{14}, 127 (1952).
\bibitem{WilsonCowan1972} H. R. Wilson and J. D. Cowan, \emph{Excitatory and inhibitory interactions in localized populations of model neurons}, Biophys. J. \textbf{12}, 1 (1972).
\bibitem{AllenStevens1994} C. Allen and C. F. Stevens, \emph{An evaluation of causes for unreliability of synaptic transmission}, Proc. Natl. Acad. Sci. USA \textbf{91}, 10380 (1994).
\bibitem{Barlow1961} H. B. Barlow, \emph{Possible principles underlying the transformations of sensory messages}, in \emph{Sensory Communication}, ed. W. A. Rosenblith (MIT Press, Cambridge, MA, 1961), pp.~217--234.
\bibitem{Cardin2009} J. A. Cardin et al., \emph{Driving fast-spiking cells induces gamma rhythm and controls sensory responses}, Nature \textbf{459}, 663 (2009).
\bibitem{Lillicrap2020} T. P. Lillicrap, A. Santoro, L. Marris, C. J. Akerman, and G. Hinton, \emph{Backpropagation and the brain}, Nat. Rev. Neurosci. \textbf{21}, 335 (2020).
\bibitem{Fremaux2016} N. Fr\'emaux and W. Gerstner, \emph{Neuromodulated spike-timing-dependent plasticity, and theory of three-factor learning rules}, Front. Neural Circuits \textbf{9}, 85 (2016).
\bibitem{Haider2006} B. Haider, A. Duque, A. R. Hasenstaub, and D. A. McCormick, \emph{Neocortical network activity in vivo is generated through a dynamic balance of excitation and inhibition}, J. Neurosci. \textbf{26}, 4535 (2006).
\bibitem{HoltKoch1997} G. R. Holt and C. Koch, \emph{Shunting inhibition does not have a divisive effect on firing rates}, Neural Comput. \textbf{9}, 1001 (1997).
\bibitem{Chance2002} F. S. Chance, L. F. Abbott, and A. D. Reyes, \emph{Gain modulation from background synaptic input}, Neuron \textbf{35}, 773 (2002).
\bibitem{Ozeki2009} H. Ozeki, I. M. Finn, E. S. Schaffer, K. D. Miller, and D. Ferster, \emph{Inhibitory stabilization of the cortical network underlies visual surround suppression}, Neuron \textbf{62}, 578 (2009).
\bibitem{HasselmoBower1992} M. E. Hasselmo and J. M. Bower, \emph{Cholinergic suppression specific to intrinsic not afferent fiber synapses in rat piriform (olfactory) cortex}, J. Neurophysiol. \textbf{67}, 1222 (1992).
\bibitem{HasselmoAndersonBower1992} M. E. Hasselmo, B. P. Anderson, and J. M. Bower, \emph{Cholinergic modulation of cortical associative memory function}, J. Neurophysiol. \textbf{67}, 1230 (1992).
\bibitem{Hasselmo1999} M. E. Hasselmo, \emph{Neuromodulation: acetylcholine and memory consolidation}, Trends Cogn. Sci. \textbf{3}, 351 (1999).
\bibitem{Tsodyks1988} M. V. Tsodyks and M. V. Feigel'man, \emph{The enhanced storage capacity in neural networks with low activity level}, Europhys. Lett. \textbf{6}, 101 (1988).
\bibitem{KalmanBucy1961} R. E. Kalman and R. S. Bucy, \emph{New results in linear filtering and prediction theory}, J. Basic Eng. \textbf{83}, 95 (1961).
\bibitem{WilsonMcNaughton1994} M. A. Wilson and B. L. McNaughton, \emph{Reactivation of hippocampal ensemble memories during sleep}, Science \textbf{265}, 676 (1994).
\bibitem{BarnesSharp1999} N. M. Barnes and T. Sharp, \emph{A review of central 5-HT receptors and their function}, Neuropharmacology \textbf{38}, 1083 (1999).
\bibitem{Bunin1998} M. A. Bunin and R. M. Wightman, \emph{Quantitative evaluation of 5-hydroxytryptamine (serotonin) neuronal release and uptake: an investigation of extrasynaptic transmission}, J. Neurosci. \textbf{18}, 4854 (1998).
\bibitem{Sprouse1987} J. S. Sprouse and G. K. Aghajanian, \emph{Electrophysiological responses of serotoninergic dorsal raphe neurons to 5-HT$_{1A}$ and 5-HT$_{1B}$ agonists}, Synapse \textbf{1}, 3 (1987).
\bibitem{Sykova2008} E. Syková and C. Nicholson, \emph{Diffusion in brain extracellular space}, Physiol. Rev. \textbf{88}, 1277 (2008).
\bibitem{WilsonNicoll2001} R. I. Wilson and R. A. Nicoll, \emph{Endogenous cannabinoids mediate retrograde signalling at hippocampal synapses}, Nature \textbf{410}, 588 (2001).
\bibitem{Garthwaite2008} J. Garthwaite, \emph{Concepts of neural nitric oxide-mediated transmission}, Eur. J. Neurosci. \textbf{27}, 2783 (2008).
\bibitem{vandenPol2012} A. N. van den Pol, \emph{Neuropeptide transmission in brain circuits}, Neuron \textbf{76}, 98 (2012).
\bibitem{Dunwiddie2001} T. V. Dunwiddie and S. A. Masino, \emph{The role and regulation of adenosine in the central nervous system}, Annu. Rev. Neurosci. \textbf{24}, 31 (2001).
\bibitem{Skaggs1996} W. E. Skaggs and B. L. McNaughton, \emph{Replay of neuronal firing sequences in rat hippocampus during sleep following spatial experience}, Science \textbf{271}, 1870 (1996).
\bibitem{Baylor1979} D. A. Baylor, T. D. Lamb, and K.-W. Yau, \emph{Responses of retinal rods to single photons}, J. Physiol. \textbf{288}, 613 (1979).
\bibitem{Hudspeth1989} A. J. Hudspeth, \emph{How the ear's works work}, Nature \textbf{341}, 397 (1989).
\bibitem{NakaRushton1966} K. I. Naka and W. A. H. Rushton, \emph{S-potentials from colour units in the retina of fish (Cyprinidae)}, J. Physiol. \textbf{185}, 536 (1966).
\bibitem{Lichtsteiner2008} P. Lichtsteiner, C. Posch, and T. Delbruck, \emph{A $128\times128$ 120 dB 15 $\mu$s latency asynchronous temporal contrast vision sensor}, IEEE J. Solid-State Circuits \textbf{43}, 566 (2008).
\bibitem{Koch2006} K. Koch et al., \emph{How much the eye tells the brain}, Curr. Biol. \textbf{16}, 1428 (2006).
\bibitem{Robles2001} L. Robles and M. A. Ruggero, \emph{Mechanics of the mammalian cochlea}, Physiol. Rev. \textbf{81}, 1305 (2001).
\bibitem{Mombaerts2004} P. Mombaerts, \emph{Genes and ligands for odorant, vomeronasal and taste receptors}, Nat. Rev. Neurosci. \textbf{5}, 263 (2004).
\bibitem{WoodSlater2001} S. J. Wood and C. R. Slater, \emph{Safety factor at the neuromuscular junction}, Prog. Neurobiol. \textbf{64}, 393 (2001).
\bibitem{Henneman1957} E. Henneman, \emph{Relation between size of neurons and their susceptibility to discharge}, Science \textbf{126}, 1345 (1957).
\bibitem{HarrisWolpert1998} C. M. Harris and D. M. Wolpert, \emph{Signal-dependent noise determines motor planning}, Nature \textbf{394}, 780 (1998).
\bibitem{Hogan1984} N. Hogan, \emph{Adaptive control of mechanical impedance by coactivation of antagonist muscles}, IEEE Trans. Autom. Control \textbf{29}, 681 (1984).
\bibitem{McAuleyMarsden2000} J. H. McAuley and C. D. Marsden, \emph{Physiological and pathological tremors and rhythmic central motor control}, Brain \textbf{123}, 1545 (2000).
\bibitem{WolpertGhahramani2000} D. M. Wolpert and Z. Ghahramani, \emph{Computational principles of movement neuroscience}, Nat. Neurosci. \textbf{3}, 1212 (2000).
\bibitem{MarderBucher2001} E. Marder and D. Bucher, \emph{Central pattern generators and the control of rhythmic movements}, Curr. Biol. \textbf{11}, R986 (2001).
\bibitem{Miyazaki2014} K. W. Miyazaki et al., \emph{Optogenetic activation of dorsal raphe serotonin neurons enhances patience for future rewards}, Curr. Biol. \textbf{24}, 2033 (2014).
\bibitem{Fuglevand1993} A. J. Fuglevand, D. A. Winter, and A. E. Patla, \emph{Models of recruitment and rate coding organization in motor-unit pools}, J. Neurophysiol. \textbf{70}, 2470 (1993).
\bibitem{Curcio1990} C. A. Curcio, K. R. Sloan, R. E. Kalina, and A. E. Hendrickson, \emph{Human photoreceptor topography}, J. Comp. Neurol. \textbf{292}, 497 (1990).
\bibitem{Hayes1950} N. D. Hayes, \emph{Roots of the transcendental equation associated with a certain difference-differential equation}, J. London Math. Soc. \textbf{25}, 226 (1950).
\bibitem{MelzackWall1965} R. Melzack and P. D. Wall, \emph{Pain mechanisms: a new theory}, Science \textbf{150}, 971 (1965).
\bibitem{Fields2004} H. Fields, \emph{State-dependent opioid control of pain}, Nat. Rev. Neurosci. \textbf{5}, 565 (2004).
\bibitem{Levine1978} J. D. Levine, N. C. Gordon, and H. L. Fields, \emph{The mechanism of placebo analgesia}, Lancet \textbf{312}, 654 (1978).
\bibitem{Woolf1983} C. J. Woolf, \emph{Evidence for a central component of post-injury pain hypersensitivity}, Nature \textbf{306}, 686 (1983).
\bibitem{Seymour2004} B. Seymour et al., \emph{Temporal difference models describe higher-order learning in humans}, Nature \textbf{429}, 664 (2004).
\bibitem{EcclestonCrombez1999} C. Eccleston and G. Crombez, \emph{Pain demands attention: a cognitive-affective model of the interruptive function of pain}, Psychol. Bull. \textbf{125}, 356 (1999).
\bibitem{WidrowHoff1960} B. Widrow and M. E. Hoff, \emph{Adaptive switching circuits}, IRE WESCON Convention Record, part 4, 96 (1960).
\bibitem{Marr1969} D. Marr, \emph{A theory of cerebellar cortex}, J. Physiol. \textbf{202}, 437 (1969).
\bibitem{Albus1971} J. S. Albus, \emph{A theory of cerebellar function}, Math. Biosci. \textbf{10}, 25 (1971).
\bibitem{Cover1965} T. M. Cover, \emph{Geometrical and statistical properties of systems of linear inequalities with applications in pattern recognition}, IEEE Trans. Electron. Comput. \textbf{EC-14}, 326 (1965).
\bibitem{Ito2001} M. Ito, \emph{Cerebellar long-term depression: characterization, signal transduction, and functional roles}, Physiol. Rev. \textbf{81}, 1143 (2001).
\bibitem{Suvrathan2016} A. Suvrathan, H. L. Payne, and J. L. Raymond, \emph{Timing rules for synaptic plasticity matched to behavioral function}, Neuron \textbf{92}, 959 (2016).
\bibitem{Fujita1982} M. Fujita, \emph{Adaptive filter model of the cerebellum}, Biol. Cybern. \textbf{45}, 195 (1982).
\bibitem{Blakemore1998} S.-J. Blakemore, D. M. Wolpert, and C. D. Frith, \emph{Central cancellation of self-produced tickle sensation}, Nat. Neurosci. \textbf{1}, 635 (1998).
\bibitem{Smith2006} M. A. Smith, A. Ghazizadeh, and R. Shadmehr, \emph{Interacting adaptive processes with different timescales underlie short-term motor learning}, PLoS Biol. \textbf{4}, e179 (2006).
\bibitem{Kording2007} K. P. K\"ording, J. B. Tenenbaum, and R. Shadmehr, \emph{The dynamics of memory as a consequence of optimal adaptation to a changing body}, Nat. Neurosci. \textbf{10}, 779 (2007).
\bibitem{Yao2023} Z. Yao et al., \emph{A high-resolution transcriptomic and spatial atlas of cell types in the whole mouse brain}, Nature \textbf{624}, 317 (2023).
\bibitem{Sanzeni2020} A. Sanzeni, B. Akitake, H. C. Goldbach, C. E. Leedy, N. Brunel, and M. H. Histed, \emph{Inhibition stabilization is a widespread property of cortical networks}, eLife \textbf{9}, e54875 (2020).
\bibitem{Padamsey2022} Z. Padamsey, D. Katsanevaki, N. Dupuy, and N. L. Rochefort, \emph{Neocortex saves energy by reducing coding precision during food scarcity}, Neuron \textbf{110}, 280 (2022).
\bibitem{Stringer2019} C. Stringer, M. Pachitariu, N. Steinmetz, C. B. Reddy, M. Carandini, and K. D. Harris, \emph{Spontaneous behaviors drive multidimensional, brainwide activity}, Science \textbf{364}, eaav7893 (2019).
\bibitem{Bittner2017} K. C. Bittner, A. D. Milstein, C. Grienberger, S. Romani, and J. C. Magee, \emph{Behavioral time scale synaptic plasticity underlies CA1 place fields}, Science \textbf{357}, 1033 (2017).
\bibitem{WhittingtonBogacz2019} J. C. R. Whittington and R. Bogacz, \emph{Theories of error back-propagation in the brain}, Trends Cogn. Sci. \textbf{23}, 235 (2019).
\bibitem{Reimer2016} J. Reimer et al., \emph{Pupil fluctuations track rapid changes in adrenergic and cholinergic activity in cortex}, Nat. Commun. \textbf{7}, 13289 (2016).
\bibitem{Beniaguev2021} D. Beniaguev, I. Segev, and M. London, \emph{Single cortical neurons as deep artificial neural networks}, Neuron \textbf{109}, 2727 (2021).
\bibitem{Gidon2020} A. Gidon et al., \emph{Dendritic action potentials and computation in human layer 2/3 cortical neurons}, Science \textbf{367}, 83 (2020).
\bibitem{Gallego2022} G. Gallego et al., \emph{Event-based vision: a survey}, IEEE Trans. Pattern Anal. Mach. Intell. \textbf{44}, 154 (2022).
\bibitem{Berger1989} R. D. Berger, J. P. Saul, and R. J. Cohen, \emph{Transfer function analysis of autonomic regulation. I. Canine atrial rate response}, Am. J. Physiol. \textbf{256}, H142 (1989).
\bibitem{Walker2010} J. J. Walker, J. R. Terry, and S. L. Lightman, \emph{Origin of ultradian pulsatility in the hypothalamic--pituitary--adrenal axis}, Proc. R. Soc. B \textbf{277}, 1627 (2010).
\bibitem{Belchetz1978} P. E. Belchetz, T. M. Plant, Y. Nakai, E. J. Keogh, and E. Knobil, \emph{Hypophysial responses to continuous and intermittent delivery of hypothalamic gonadotropin-releasing hormone}, Science \textbf{202}, 631 (1978).
\bibitem{Czeisler1999} C. A. Czeisler et al., \emph{Stability, precision, and near-24-hour period of the human circadian pacemaker}, Science \textbf{284}, 2177 (1999).
\bibitem{TakahashiJS2017} J. S. Takahashi, \emph{Transcriptional architecture of the mammalian circadian clock}, Nat. Rev. Genet. \textbf{18}, 164 (2017).
\bibitem{Musk2019} E. Musk and Neuralink, \emph{An integrated brain-machine interface platform with thousands of channels}, J. Med. Internet Res. \textbf{21}, e16194 (2019).
\bibitem{Hochberg2006} L. R. Hochberg et al., \emph{Neuronal ensemble control of prosthetic devices by a human with tetraplegia}, Nature \textbf{442}, 164 (2006).
\bibitem{Willett2021} F. R. Willett et al., \emph{High-performance brain-to-text communication via handwriting}, Nature \textbf{593}, 249 (2021).
\bibitem{Willett2023} F. R. Willett et al., \emph{A high-performance speech neuroprosthesis}, Nature \textbf{620}, 1031 (2023).
\bibitem{Gallego2017} J. A. Gallego, M. G. Perich, L. E. Miller, and S. A. Solla, \emph{Neural manifolds for the control of movement}, Neuron \textbf{94}, 978 (2017).
\bibitem{Fernandez2021} E. Fern\'andez et al., \emph{Visual percepts evoked with an intracortical 96-channel microelectrode array inserted in human occipital cortex}, J. Clin. Invest. \textbf{131}, e151331 (2021).
\bibitem{Tritsch2012} N. X. Tritsch, J. B. Ding, and B. L. Sabatini, \emph{Dopaminergic neurons inhibit striatal output through non-canonical release of GABA}, Nature \textbf{490}, 262 (2012).
\bibitem{Zhang2015} S. Zhang et al., \emph{Dopaminergic and glutamatergic microdomains in a subset of rodent mesoaccumbens axons}, Nat. Neurosci. \textbf{18}, 386 (2015).
\bibitem{Mongillo2008} G. Mongillo, O. Barak, and M. Tsodyks, \emph{Synaptic theory of working memory}, Science \textbf{319}, 1543 (2008).
\bibitem{Stokes2015} M. G. Stokes, \emph{`Activity-silent' working memory in prefrontal cortex: a dynamic coding framework}, Trends Cogn. Sci. \textbf{19}, 394 (2015).
\bibitem{Smith2013} S. L. Smith, I. T. Smith, T. Branco, and M. H\"ausser, \emph{Dendritic spikes enhance stimulus selectivity in cortical neurons in vivo}, Nature \textbf{503}, 115 (2013).
\bibitem{Baden2016} T. Baden et al., \emph{The functional diversity of retinal ganglion cells in the mouse}, Nature \textbf{529}, 345 (2016).
\bibitem{Lottem2018} E. Lottem et al., \emph{Activation of serotonin neurons promotes active persistence in a probabilistic foraging task}, Nat. Commun. \textbf{9}, 1000 (2018).
\bibitem{Payeur2021} A. Payeur, J. Guerguiev, F. Zenke, B. A. Richards, and R. Naud, \emph{Burst-dependent synaptic plasticity can coordinate learning in hierarchical circuits}, Nat. Neurosci. \textbf{24}, 1010 (2021).
\bibitem{MehonicKenyon2022} A. Mehonic and A. J. Kenyon, \emph{Brain-inspired computing needs a master plan}, Nature \textbf{604}, 255 (2022).
\bibitem{Poe2020} G. R. Poe et al., \emph{Locus coeruleus: a new look at the blue spot}, Nat. Rev. Neurosci. \textbf{21}, 644 (2020).
\bibitem{Hangya2015} B. Hangya, S. P. Ranade, M. Lorenc, and A. Kepecs, \emph{Central cholinergic neurons are rapidly recruited by reinforcement feedback}, Cell \textbf{162}, 1155 (2015).
\bibitem{FernandezRuiz2019} A. Fern\'andez-Ruiz et al., \emph{Long-duration hippocampal sharp wave ripples improve memory}, Science \textbf{364}, 1082 (2019).
\bibitem{Herzfeld2018} D. J. Herzfeld, Y. Kojima, R. Soetedjo, and R. Shadmehr, \emph{Encoding of error and learning to correct that error by the Purkinje cells of the cerebellum}, Nat. Neurosci. \textbf{21}, 736 (2018).
\bibitem{Duan2014} B. Duan et al., \emph{Identification of spinal circuits transmitting and gating mechanical pain}, Cell \textbf{159}, 1417 (2014).
\bibitem{Tremblay2016} R. Tremblay, S. Lee, and B. Rudy, \emph{GABAergic interneurons in the neocortex: from cellular properties to circuits}, Neuron \textbf{91}, 260 (2016).
\bibitem{Ren2018} J. Ren et al., \emph{Anatomically defined and functionally distinct dorsal raphe serotonin sub-systems}, Cell \textbf{175}, 472 (2018).
\bibitem{Pfeffer2013} C. K. Pfeffer, M. Xue, M. He, Z. J. Huang, and M. Scanziani, \emph{Inhibition of inhibition in visual cortex: the logic of connections between molecularly distinct interneurons}, Nat. Neurosci. \textbf{16}, 1068 (2013).
\bibitem{Pi2013} H.-J. Pi et al., \emph{Cortical interneurons that specialize in disinhibitory control}, Nature \textbf{503}, 521 (2013).
\bibitem{Keller2018} G. B. Keller and T. D. Mrsic-Flogel, \emph{Predictive processing: a canonical cortical computation}, Neuron \textbf{100}, 424 (2018).
\bibitem{HubelWiesel1962} D. H. Hubel and T. N. Wiesel, \emph{Receptive fields, binocular interaction and functional architecture in the cat's visual cortex}, J. Physiol. \textbf{160}, 106 (1962).
\bibitem{Riesenhuber1999} M. Riesenhuber and T. Poggio, \emph{Hierarchical models of object recognition in cortex}, Nat. Neurosci. \textbf{2}, 1019 (1999).
\bibitem{Fukushima1980} K. Fukushima, \emph{Neocognitron: a self-organizing neural network model for a mechanism of pattern recognition unaffected by shift in position}, Biol. Cybern. \textbf{36}, 193 (1980).
\bibitem{Yamins2014} D. L. K. Yamins et al., \emph{Performance-optimized hierarchical models predict neural responses in higher visual cortex}, Proc. Natl. Acad. Sci. USA \textbf{111}, 8619 (2014).
\bibitem{Hung2005} C. P. Hung, G. Kreiman, T. Poggio, and J. J. DiCarlo, \emph{Fast readout of object identity from macaque inferior temporal cortex}, Science \textbf{310}, 863 (2005).
\bibitem{WiskottSejnowski2002} L. Wiskott and T. J. Sejnowski, \emph{Slow feature analysis: unsupervised learning of invariances}, Neural Comput. \textbf{14}, 715 (2002).
\bibitem{Foldiak1991} P. F\"oldi\'ak, \emph{Learning invariance from transformation sequences}, Neural Comput. \textbf{3}, 194 (1991).
\bibitem{LiDiCarlo2008} N. Li and J. J. DiCarlo, \emph{Unsupervised natural experience rapidly alters invariant object representation in visual cortex}, Science \textbf{321}, 1502 (2008).
\bibitem{RaoBallard1999} R. P. N. Rao and D. H. Ballard, \emph{Predictive coding in the visual cortex: a functional interpretation of some extra-classical receptive-field effects}, Nat. Neurosci. \textbf{2}, 79 (1999).
\bibitem{Kar2019} K. Kar, J. Kubilius, K. Schmidt, E. B. Issa, and J. J. DiCarlo, \emph{Evidence that recurrent circuits are critical to the ventral stream's execution of core object recognition behavior}, Nat. Neurosci. \textbf{22}, 974 (2019).
\bibitem{Szegedy2014} C. Szegedy et al., \emph{Intriguing properties of neural networks}, in \emph{Proc. ICLR} (2014), arXiv:1312.6199.
\bibitem{Weiss2002} Y. Weiss, E. P. Simoncelli, and E. H. Adelson, \emph{Motion illusions as optimal percepts}, Nat. Neurosci. \textbf{5}, 598 (2002).
\bibitem{Purves2011} D. Purves, W. T. Wojtach, and R. B. Lotto, \emph{Understanding vision in wholly empirical terms}, Proc. Natl. Acad. Sci. USA \textbf{108}, Suppl.~3, 15588 (2011).
\bibitem{LaferSousa2015} R. Lafer-Sousa, K. L. Hermann, and B. R. Conway, \emph{Striking individual differences in color perception uncovered by `the dress' photograph}, Curr. Biol. \textbf{25}, R545 (2015).
\bibitem{Wallisch2017} P. Wallisch, \emph{Illumination assumptions account for individual differences in the perceptual interpretation of a profoundly ambiguous stimulus in the color domain: ``The dress''}, J. Vis. \textbf{17}(4), 5 (2017).
\bibitem{Schwartz2009} O. Schwartz, T. J. Sejnowski, and P. Dayan, \emph{Perceptual organization in the tilt illusion}, J. Vis. \textbf{9}(4), 19 (2009).
\bibitem{SchillerCarvey2005} P. H. Schiller and C. E. Carvey, \emph{The Hermann grid illusion revisited}, Perception \textbf{34}, 1375 (2005).
\bibitem{Nijhawan1994} R. Nijhawan, \emph{Motion extrapolation in catching}, Nature \textbf{370}, 256 (1994).
\bibitem{Conway2005} B. R. Conway, A. Kitaoka, A. Yazdanbakhsh, C. C. Pack, and M. S. Livingstone, \emph{Neural basis for a powerful static motion illusion}, J. Neurosci. \textbf{25}, 5651 (2005).
\bibitem{OteroMillan2012} J. Otero-Millan, S. L. Macknik, and S. Martinez-Conde, \emph{Microsaccades and blinks trigger illusory rotation in the ``rotating snakes'' illusion}, J. Neurosci. \textbf{32}, 6043 (2012).
\bibitem{MorenoBote2007} R. Moreno-Bote, J. Rinzel, and N. Rubin, \emph{Noise-induced alternations in an attractor network model of perceptual bistability}, J. Neurophysiol. \textbf{98}, 1125 (2007).
\bibitem{Watanabe2018} E. Watanabe, A. Kitaoka, K. Sakamoto, M. Yasugi, and K. Tanaka, \emph{Illusory motion reproduced by deep neural networks trained for prediction}, Front. Psychol. \textbf{9}, 345 (2018).
\bibitem{GomezVilla2020} A. G\'omez-Villa, A. Mart\'in, J. Vazquez-Corral, M. Bertalm\'io, and J. Malo, \emph{Color illusions also deceive CNNs for low-level vision tasks: analysis and implications}, Vision Res. \textbf{176}, 156 (2020).
\bibitem{delCastillo1954} J. del Castillo and B. Katz, \emph{Quantal components of the end-plate potential}, J. Physiol. \textbf{124}, 560 (1954).
\bibitem{Nowak1984} L. Nowak, P. Bregestovski, P. Ascher, A. Herbet, and A. Prochiantz, \emph{Magnesium gates glutamate-activated channels in mouse central neurones}, Nature \textbf{307}, 462 (1984).
\bibitem{Malinow2002} R. Malinow and R. C. Malenka, \emph{AMPA receptor trafficking and synaptic plasticity}, Annu. Rev. Neurosci. \textbf{25}, 103 (2002).
\bibitem{Matsuzaki2004} M. Matsuzaki, N. Honkura, G. C. R. Ellis-Davies, and H. Kasai, \emph{Structural basis of long-term potentiation in single dendritic spines}, Nature \textbf{429}, 761 (2004).
\bibitem{Rudomin1999} P. Rudomin and R. F. Schmidt, \emph{Presynaptic inhibition in the vertebrate spinal cord revisited}, Exp. Brain Res. \textbf{129}, 1 (1999).
\bibitem{HutcheonYarom2000} B. Hutcheon and Y. Yarom, \emph{Resonance, oscillation and the intrinsic frequency preferences of neurons}, Trends Neurosci. \textbf{23}, 216 (2000).
\bibitem{Seung1996} H. S. Seung, \emph{How the brain keeps the eyes still}, Proc. Natl. Acad. Sci. USA \textbf{93}, 13339 (1996).
\bibitem{Hounsgaard1988} J. Hounsgaard, H. Hultborn, B. Jespersen, and O. Kiehn, \emph{Bistability of alpha-motoneurones in the decerebrate cat and in the acute spinal cat after intravenous 5-hydroxytryptophan}, J. Physiol. \textbf{405}, 345 (1988).
\bibitem{Izhikevich2007} E. M. Izhikevich, \emph{Dynamical Systems in Neuroscience: The Geometry of Excitability and Bursting} (MIT Press, Cambridge, MA, 2007).
\bibitem{AgmonSnir1998} H. Agmon-Snir, C. E. Carr, and J. Rinzel, \emph{The role of dendrites in auditory coincidence detection}, Nature \textbf{393}, 268 (1998).
\bibitem{BorstSoria2012} J. G. G. Borst and J. Soria van Hoeve, \emph{The calyx of Held synapse: from model synapse to auditory relay}, Annu. Rev. Physiol. \textbf{74}, 199 (2012).
\bibitem{Euler2002} T. Euler, P. B. Detwiler, and W. Denk, \emph{Directionally selective calcium signals in dendrites of starburst amacrine cells}, Nature \textbf{418}, 845 (2002).
\bibitem{AstonJonesBloom1981} G. Aston-Jones and F. E. Bloom, \emph{Activity of norepinephrine-containing locus coeruleus neurons in behaving rats anticipates fluctuations in the sleep-waking cycle}, J. Neurosci. \textbf{1}, 876 (1981).
\bibitem{BrooksPeever2012} P. L. Brooks and J. H. Peever, \emph{Identification of the transmitter and receptor mechanisms responsible for REM sleep paralysis}, J. Neurosci. \textbf{32}, 9785 (2012).
\bibitem{Borbely1982} A. A. Borb\'ely, \emph{A two process model of sleep regulation}, Hum. Neurobiol. \textbf{1}, 195 (1982).
\bibitem{PorkkaHeiskanen1997} T. Porkka-Heiskanen et al., \emph{Adenosine: a mediator of the sleep-inducing effects of prolonged wakefulness}, Science \textbf{276}, 1265 (1997).
\bibitem{Steriade1993} M. Steriade, A. N\'u\~nez, and F. Amzica, \emph{A novel slow ($<1$ Hz) oscillation of neocortical neurons in vivo: depolarizing and hyperpolarizing components}, J. Neurosci. \textbf{13}, 3252 (1993).
\bibitem{McCormick1992} D. A. McCormick, \emph{Neurotransmitter actions in the thalamus and cerebral cortex and their role in neuromodulation of thalamocortical activity}, Prog. Neurobiol. \textbf{39}, 337 (1992).
\bibitem{SteriadeMcCormickSejnowski1993} M. Steriade, D. A. McCormick, and T. J. Sejnowski, \emph{Thalamocortical oscillations in the sleeping and aroused brain}, Science \textbf{262}, 679 (1993).
\bibitem{Nadasdy1999} Z. N\'adasdy et al., \emph{Replay and time compression of recurring spike sequences in the hippocampus}, J. Neurosci. \textbf{19}, 9497 (1999).
\bibitem{Maingret2016} N. Maingret et al., \emph{Hippocampo-cortical coupling mediates memory consolidation during sleep}, Nat. Neurosci. \textbf{19}, 959 (2016).
\bibitem{McClelland1995} J. L. McClelland, B. L. McNaughton, and R. C. O'Reilly, \emph{Why there are complementary learning systems in the hippocampus and neocortex}, Psychol. Rev. \textbf{102}, 419 (1995).
\bibitem{TononiCirelli2014} G. Tononi and C. Cirelli, \emph{Sleep and the price of plasticity: from synaptic and cellular homeostasis to memory consolidation and integration}, Neuron \textbf{81}, 12 (2014).
\bibitem{deVivo2017} L. de Vivo et al., \emph{Ultrastructural evidence for synaptic scaling across the wake/sleep cycle}, Science \textbf{355}, 507 (2017).
\bibitem{Xie2013} L. Xie et al., \emph{Sleep drives metabolite clearance from the adult brain}, Science \textbf{342}, 373 (2013).
\bibitem{Miao2024} A. Miao et al., \emph{Brain clearance is reduced during sleep and anesthesia}, Nat. Neurosci. \textbf{27}, 1046 (2024).
\bibitem{Vyazovskiy2011} V. V. Vyazovskiy et al., \emph{Local sleep in awake rats}, Nature \textbf{472}, 443 (2011).
\bibitem{Lancaster2013} M. A. Lancaster et al., \emph{Cerebral organoids model human brain development and microcephaly}, Nature \textbf{501}, 373 (2013).
\bibitem{Jordan2024} F. D. Jordan et al., \emph{Open and remotely accessible Neuroplatform for research in wetware computing}, Front. Artif. Intell. \textbf{7}, 1376042 (2024).
\bibitem{Smirnova2023} L. Smirnova et al., \emph{Organoid intelligence (OI): the new frontier in biocomputing and intelligence-in-a-dish}, Front. Sci. \textbf{1}, 1017235 (2023).
\bibitem{Wagenaar2006} D. A. Wagenaar, J. Pine, and S. M. Potter, \emph{An extremely rich repertoire of bursting patterns during the development of cortical cultures}, BMC Neurosci. \textbf{7}, 11 (2006).
\bibitem{Shahaf2001} G. Shahaf and S. Marom, \emph{Learning in networks of cortical neurons}, J. Neurosci. \textbf{21}, 8782 (2001).
\bibitem{Kagan2022} B. J. Kagan et al., \emph{In vitro neurons learn and exhibit sentience when embodied in a simulated game-world}, Neuron \textbf{110}, 3952 (2022).
\bibitem{Friston2010} K. Friston, \emph{The free-energy principle: a unified brain theory?}, Nat. Rev. Neurosci. \textbf{11}, 127 (2010).
\bibitem{Balci2023} F. Balci et al., \emph{A response to claims of emergent intelligence and sentience in a dish}, Neuron \textbf{111}, 604 (2023).
\bibitem{Bakkum2008} D. J. Bakkum, Z. C. Chao, and S. M. Potter, \emph{Spatio-temporal electrical stimuli shape behavior of an embodied cortical network in a goal-directed learning task}, J. Neural Eng. \textbf{5}, 310 (2008).
\bibitem{Maass2002} W. Maass, T. Natschl\"ager, and H. Markram, \emph{Real-time computing without stable states: a new framework for neural computation based on perturbations}, Neural Comput. \textbf{14}, 2531 (2002).
\bibitem{JaegerHaas2004} H. Jaeger and H. Haas, \emph{Harnessing nonlinearity: predicting chaotic systems and saving energy in wireless communication}, Science \textbf{304}, 78 (2004).
\bibitem{Cai2023} H. Cai et al., \emph{Brain organoid reservoir computing for artificial intelligence}, Nat. Electron. \textbf{6}, 1032 (2023).
\bibitem{Habibollahi2023} F. Habibollahi, B. J. Kagan, A. N. Burkitt, and C. French, \emph{Critical dynamics arise during structured information presentation within embodied in vitro neuronal networks}, Nat. Commun. \textbf{14}, 5287 (2023).
\bibitem{Robbins2026} A. Robbins et al., \emph{Goal-directed learning in cortical organoids}, Cell Rep. \textbf{45}, 116984 (2026).
\bibitem{Keene2020} S. T. Keene et al., \emph{A biohybrid synapse with neurotransmitter-mediated plasticity}, Nat. Mater. \textbf{19}, 969 (2020).
\bibitem{Jo2016} J. Jo et al., \emph{Midbrain-like organoids from human pluripotent stem cells contain functional dopaminergic and neuromelanin-producing neurons}, Cell Stem Cell \textbf{19}, 248 (2016).
\bibitem{Hendry1987} S. H. Hendry, H. D. Schwark, E. G. Jones, and J. Yan, \emph{Numbers and proportions of GABA-immunoreactive neurons in different areas of monkey cerebral cortex}, J. Neurosci. \textbf{7}, 1503 (1987).
\bibitem{Tang2001synapses} Y. Tang, J. R. Nyengaard, D. M. G. De Groot, and H. J. G. Gundersen, \emph{Total regional and global number of synapses in the human brain neocortex}, Synapse \textbf{41}, 258 (2001).
\bibitem{Pakkenberg1997} B. Pakkenberg and H. J. G. Gundersen, \emph{Neocortical neuron number in humans: effect of sex and age}, J. Comp. Neurol. \textbf{384}, 312 (1997).
\bibitem{Matsuda2009} W. Matsuda, T. Furuta, K. C. Nakamura, H. Hioki, F. Fujiyama, R. Arai, and T. Kaneko, \emph{Single nigrostriatal dopaminergic neurons form widely spread and highly dense axonal arborizations in the neostriatum}, J. Neurosci. \textbf{29}, 444 (2009).
\bibitem{Pakkenberg1991} B. Pakkenberg, A. M{\o}ller, H. J. G. Gundersen, A. Mouritzen Dam, and H. Pakkenberg, \emph{The absolute number of nerve cells in substantia nigra in normal subjects and in patients with Parkinson's disease estimated with an unbiased stereological method}, J. Neurol. Neurosurg. Psychiatry \textbf{54}, 30 (1991).
\bibitem{Baker1990} K. G. Baker, G. M. Halliday, and I. T\"ork, \emph{Cytoarchitecture of the human dorsal raphe nucleus}, J. Comp. Neurol. \textbf{301}, 147 (1990).
\bibitem{Baker1991} K. G. Baker, G. M. Halliday, P. Halasz, J.-P. Hornung, L. B. Geffen, R. G. H. Cotton, and I. T\"ork, \emph{Cytoarchitecture of serotonin-synthesizing neurons in the pontine tegmentum of the human brain}, Synapse \textbf{7}, 301 (1991).
\bibitem{Airaksinen1991} M. S. Airaksinen, A. Paetau, L. Palj\"arvi, K. Reinikainen, P. Riekkinen, R. Suomalainen, and P. Panula, \emph{Histamine neurons in human hypothalamus: anatomy in normal and Alzheimer diseased brains}, Neuroscience \textbf{44}, 465 (1991).
\bibitem{Colquhoun1992} D. Colquhoun, P. Jonas, and B. Sakmann, \emph{Action of brief pulses of glutamate on AMPA/kainate receptors in patches from different neurones of rat hippocampal slices}, J. Physiol. \textbf{458}, 261 (1992).
\bibitem{DutarNicoll1988} P. Dutar and R. A. Nicoll, \emph{A physiological role for GABA$_B$ receptors in the central nervous system}, Nature \textbf{332}, 156 (1988).

\bibitem{Rauch2003} A. Rauch, G. La Camera, H.-R. L\"uscher, W. Senn, and S. Fusi, \emph{Neocortical pyramidal cells respond as integrate-and-fire neurons to in vivo-like input currents}, J. Neurophysiol. \textbf{90}, 1598 (2003).
\bibitem{KlumppEady1956} R. G. Klumpp and H. R. Eady, \emph{Some measurements of interaural time difference thresholds}, J. Acoust. Soc. Am. \textbf{28}, 859 (1956).
\bibitem{Brand2002} A. Brand, O. Behrend, T. Marquardt, D. McAlpine, and B. Grothe, \emph{Precise inhibition is essential for microsecond interaural time difference coding}, Nature \textbf{417}, 543 (2002).
\bibitem{Funahashi1989} S. Funahashi, C. J. Bruce, and P. S. Goldman-Rakic, \emph{Mnemonic coding of visual space in the monkey's dorsolateral prefrontal cortex}, J. Neurophysiol. \textbf{61}, 331 (1989).
\bibitem{WangMarkram2006} Y. Wang, H. Markram, P. H. Goodman, T. K. Berger, J. Ma, and P. S. Goldman-Rakic, \emph{Heterogeneity in the pyramidal network of the medial prefrontal cortex}, Nat. Neurosci. \textbf{9}, 534 (2006).

\bibitem{McCullochPitts1943} W. S. McCulloch and W. Pitts, \emph{A logical calculus of the ideas immanent in nervous activity}, Bull. Math. Biophys. \textbf{5}, 115 (1943).
\bibitem{Fried2002} S. I. Fried, T. A. M\"unch, and F. S. Werblin, \emph{Mechanisms and circuitry underlying directional selectivity in the retina}, Nature \textbf{420}, 411 (2002).

\bibitem{Hromadka2008} T. Hrom\'adka, M. R. DeWeese, and A. M. Zador, \emph{Sparse representation of sounds in the unanesthetized auditory cortex}, PLoS Biol. \textbf{6}, e16 (2008).
\bibitem{Abbott1997depression} L. F. Abbott, J. A. Varela, K. Sen, and S. B. Nelson, \emph{Synaptic depression and cortical gain control}, Science \textbf{275}, 221 (1997).
\bibitem{ShadlenNewsome1998} M. N. Shadlen and W. T. Newsome, \emph{The variable discharge of cortical neurons: implications for connectivity, computation, and information coding}, J. Neurosci. \textbf{18}, 3870 (1998).

\bibitem{TrulsonJacobs1979} M. E. Trulson and B. L. Jacobs, \emph{Raphe unit activity in freely moving cats: correlation with level of behavioral arousal}, Brain Res. \textbf{163}, 135 (1979).
\bibitem{GagnonParent2014} D. Gagnon and M. Parent, \emph{Distribution of VGLUT3 in highly collateralized axons from the rat dorsal raphe nucleus as revealed by single-neuron reconstructions}, PLoS One \textbf{9}, e87709 (2014).
\bibitem{VandermaelenAghajanian1983} C. P. Vandermaelen and G. K. Aghajanian, \emph{Electrophysiological and pharmacological characterization of serotonergic dorsal raphe neurons recorded extracellularly and intracellularly in rat brain slices}, Brain Res. \textbf{289}, 109 (1983).
\bibitem{CourtneyFord2016} N. A. Courtney and C. P. Ford, \emph{Mechanisms of 5-HT$_{1A}$ receptor-mediated transmission in dorsal raphe serotonin neurons}, J. Physiol. \textbf{594}, 953 (2016).
\bibitem{Hashemi2012serotonin} P. Hashemi et al., \emph{Brain dopamine and serotonin differ in regulation and its consequences}, Proc. Natl. Acad. Sci. USA \textbf{109}, 11510 (2012).
\bibitem{Black1934feedback} H. S. Black, \emph{Stabilized feedback amplifiers}, Bell Syst. Tech. J. \textbf{13}, 1 (1934).
\bibitem{KobayashiSchultz2008} S. Kobayashi and W. Schultz, \emph{Influence of reward delays on responses of dopamine neurons}, J. Neurosci. \textbf{28}, 7837 (2008).
\bibitem{Schweighofer2008} N. Schweighofer et al., \emph{Low-serotonin levels increase delayed reward discounting in humans}, J. Neurosci. \textbf{28}, 4528 (2008).

\bibitem{Malenka1988calcium} R. C. Malenka, J. A. Kauer, R. S. Zucker, and R. A. Nicoll, \emph{Postsynaptic calcium is sufficient for potentiation of hippocampal synaptic transmission}, Science \textbf{242}, 81 (1988).
\bibitem{Bliss1973LTP} T. V. P. Bliss and T. L{\o}mo, \emph{Long-lasting potentiation of synaptic transmission in the dentate area of the anaesthetized rabbit following stimulation of the perforant path}, J. Physiol. \textbf{232}, 331 (1973).
\bibitem{Kelso1986hebb} S. R. Kelso, A. H. Ganong, and T. H. Brown, \emph{Hebbian synapses in hippocampus}, Proc. Natl. Acad. Sci. USA \textbf{83}, 5326 (1986).
\bibitem{He2015eligibility} K. He et al., \emph{Distinct eligibility traces for LTP and LTD in cortical synapses}, Neuron \textbf{88}, 528 (2015).
\bibitem{Olveczky2005} B. P. \"Olveczky, A. S. Andalman, and M. S. Fee, \emph{Vocal experimentation in the juvenile songbird requires a basal ganglia circuit}, PLoS Biol. \textbf{3}, e153 (2005).
\bibitem{TumerBrainard2007} E. C. Tumer and M. S. Brainard, \emph{Performance variability enables adaptive plasticity of `crystallized' adult birdsong}, Nature \textbf{450}, 1240 (2007).
\bibitem{Eshel2015arithmetic} N. Eshel et al., \emph{Arithmetic and local circuitry underlying dopamine prediction errors}, Nature \textbf{525}, 243 (2015).
\bibitem{GraceOnn1989} A. A. Grace and S.-P. Onn, \emph{Morphology and electrophysiological properties of immunocytochemically identified rat dopamine neurons recorded in vitro}, J. Neurosci. \textbf{9}, 3463 (1989).
\bibitem{ODoherty2004actor} J. O'Doherty et al., \emph{Dissociable roles of ventral and dorsal striatum in instrumental conditioning}, Science \textbf{304}, 452 (2004).

\bibitem{NeweyPowell1987} W. K. Newey and J. L. Powell, \emph{Asymmetric least squares estimation and testing}, Econometrica \textbf{55}, 819 (1987).
\bibitem{Beierholm2013vigor} U. Beierholm et al., \emph{Dopamine modulates reward-related vigor}, Neuropsychopharmacology \textbf{38}, 1495 (2013).
\bibitem{Howe2013ramp} M. W. Howe, P. L. Tierney, S. G. Sandberg, P. E. M. Phillips, and A. M. Graybiel, \emph{Prolonged dopamine signalling in striatum signals proximity and value of distant rewards}, Nature \textbf{500}, 575 (2013).

\bibitem{Baker1989LC} K. G. Baker, I. T\"ork, J.-P. Hornung, and P. Halasz, \emph{The human locus coeruleus complex: an immunohistochemical and three dimensional reconstruction study}, Exp. Brain Res. \textbf{77}, 257 (1989).
\bibitem{Schwarz2015LC} L. A. Schwarz et al., \emph{Viral-genetic tracing of the input--output organization of a central noradrenaline circuit}, Nature \textbf{524}, 88 (2015).
\bibitem{Uematsu2017LC} A. Uematsu et al., \emph{Modular organization of the brainstem noradrenaline system coordinates opposing learning states}, Nat. Neurosci. \textbf{20}, 1602 (2017).
\bibitem{AstonJones1994vigilance} G. Aston-Jones, J. Rajkowski, P. Kubiak, and T. Alexinsky, \emph{Locus coeruleus neurons in monkey are selectively activated by attended cues in a vigilance task}, J. Neurosci. \textbf{14}, 4467 (1994).
\bibitem{Clayton2004LC} E. C. Clayton, J. Rajkowski, J. D. Cohen, and G. Aston-Jones, \emph{Phasic activation of monkey locus ceruleus neurons by simple decisions in a forced-choice task}, J. Neurosci. \textbf{24}, 9914 (2004).
\bibitem{Foote1975NE} S. L. Foote, R. Freedman, and A. P. Oliver, \emph{Effects of putative neurotransmitters on neuronal activity in monkey auditory cortex}, Brain Res. \textbf{86}, 229 (1975).
\bibitem{JepmaNieuwenhuis2011} M. Jepma and S. Nieuwenhuis, \emph{Pupil diameter predicts changes in the exploration--exploitation trade-off: evidence for the adaptive gain theory}, J. Cogn. Neurosci. \textbf{23}, 1587 (2011).
\bibitem{Tervo2014stochastic} D. G. R. Tervo et al., \emph{Behavioral variability through stochastic choice and its gating by anterior cingulate cortex}, Cell \textbf{159}, 21 (2014).
\bibitem{Usher1999LC} M. Usher, J. D. Cohen, D. Servan-Schreiber, J. Rajkowski, and G. Aston-Jones, \emph{The role of locus coeruleus in the regulation of cognitive performance}, Science \textbf{283}, 549 (1999).

\bibitem{Mesulam1983} M.-M. Mesulam, E. J. Mufson, B. H. Wainer, and A. I. Levey, \emph{Central cholinergic pathways in the rat: an overview based on an alternative nomenclature (Ch1--Ch6)}, Neuroscience \textbf{10}, 1185 (1983).
\bibitem{Marrosu1995} F. Marrosu et al., \emph{Microdialysis measurement of cortical and hippocampal acetylcholine release during sleep-wake cycle in freely moving cats}, Brain Res. \textbf{671}, 329 (1995).
\bibitem{Picciotto2012} M. R. Picciotto, M. J. Higley, and Y. S. Mineur, \emph{Acetylcholine as a neuromodulator: cholinergic signaling shapes nervous system function and behavior}, Neuron \textbf{76}, 116 (2012).
\bibitem{HasselmoSchnell1994} M. E. Hasselmo and E. Schnell, \emph{Laminar selectivity of the cholinergic suppression of synaptic transmission in rat hippocampal region CA1: computational modeling and brain slice physiology}, J. Neurosci. \textbf{14}, 3898 (1994).
\bibitem{Gil1997} Z. Gil, B. W. Connors, and Y. Amitai, \emph{Differential regulation of neocortical synapses by neuromodulators and activity}, Neuron \textbf{19}, 679 (1997).
\bibitem{KilgardMerzenich1998} M. P. Kilgard and M. M. Merzenich, \emph{Cortical map reorganization enabled by nucleus basalis activity}, Science \textbf{279}, 1714 (1998).
\bibitem{Atri2004} A. Atri et al., \emph{Blockade of central cholinergic receptors impairs new learning and increases proactive interference in a word paired-associate memory task}, Behav. Neurosci. \textbf{118}, 223 (2004).
\bibitem{Girardeau2009} G. Girardeau, K. Benchenane, S. I. Wiener, G. Buzs\'aki, and M. B. Zugaro, \emph{Selective suppression of hippocampal ripples impairs spatial memory}, Nat. Neurosci. \textbf{12}, 1222 (2009).

\input{supplement/bib_10}
\bibitem{CopenhagenJahr1989} D. R. Copenhagen and C. E. Jahr, \emph{Release of endogenous excitatory amino acids from turtle photoreceptors}, Nature \textbf{341}, 536 (1989).
\bibitem{GlowatzkiFuchs2002} E. Glowatzki and P. A. Fuchs, \emph{Transmitter release at the hair cell ribbon synapse}, Nat. Neurosci. \textbf{5}, 147 (2002).
\bibitem{Tinsley2016} J. N. Tinsley et al., \emph{Direct detection of a single photon by humans}, Nat. Commun. \textbf{7}, 12172 (2016).
\bibitem{Sellick1982} P. M. Sellick, R. Patuzzi, and B. M. Johnstone, \emph{Measurement of basilar membrane motion in the guinea pig using the M\"ossbauer technique}, J. Acoust. Soc. Am. \textbf{72}, 131 (1982).
\bibitem{Ruggero1997} M. A. Ruggero, N. C. Rich, A. Recio, S. S. Narayan, and L. Robles, \emph{Basilar-membrane responses to tones at the base of the chinchilla cochlea}, J. Acoust. Soc. Am. \textbf{101}, 2151 (1997).
\bibitem{NormannPerlman1979} R. A. Normann and I. Perlman, \emph{The effects of background illumination on the photoresponses of red and green cones}, J. Physiol. \textbf{286}, 491 (1979).
\bibitem{Laughlin1981} S. Laughlin, \emph{A simple coding procedure enhances a neuron's information capacity}, Z. Naturforsch. C \textbf{36}, 910 (1981).
\bibitem{CurcioAllen1990} C. A. Curcio and K. A. Allen, \emph{Topography of ganglion cells in human retina}, J. Comp. Neurol. \textbf{300}, 5 (1990).
\bibitem{Greenwood1990} D. D. Greenwood, \emph{A cochlear frequency-position function for several species --- 29 years later}, J. Acoust. Soc. Am. \textbf{87}, 2592 (1990).
\bibitem{Eguiluz2000} V. M. Egu\'iluz, M. Ospeck, Y. Choe, A. J. Hudspeth, and M. O. Magnasco, \emph{Essential nonlinearities in hearing}, Phys. Rev. Lett. \textbf{84}, 5232 (2000).
\bibitem{PalmerRussell1986} A. R. Palmer and I. J. Russell, \emph{Phase-locking in the cochlear nerve of the guinea-pig and its relation to the receptor potential of inner hair-cells}, Hear. Res. \textbf{24}, 1 (1986).
\bibitem{Malnic1999} B. Malnic, J. Hirono, T. Sato, and L. B. Buck, \emph{Combinatorial receptor codes for odors}, Cell \textbf{96}, 713 (1999).
\bibitem{Finger2005} T. E. Finger et al., \emph{ATP signaling is crucial for communication from taste buds to gustatory nerves}, Science \textbf{310}, 1495 (2005).
\bibitem{Huang2005} Y.-J. Huang et al., \emph{Mouse taste buds use serotonin as a neurotransmitter}, J. Neurosci. \textbf{25}, 843 (2005).
\bibitem{JohanssonVallbo1979} R. S. Johansson and A. B. Vallbo, \emph{Tactile sensibility in the human hand: relative and absolute densities of four types of mechanoreceptive units in glabrous skin}, J. Physiol. \textbf{286}, 283 (1979).
\bibitem{McKemy2002} D. D. McKemy, W. M. Neuhausser, and D. Julius, \emph{Identification of a cold receptor reveals a general role for TRP channels in thermosensation}, Nature \textbf{416}, 52 (2002).

\bibitem{Schmolesky1998} M. T. Schmolesky et al., \emph{Signal timing across the macaque visual system}, J. Neurophysiol. \textbf{79}, 3272 (1998).
\bibitem{Lampl2004} I. Lampl, D. Ferster, T. Poggio, and M. Riesenhuber, \emph{Intracellular measurements of spatial integration and the MAX operation in complex cells of the cat primary visual cortex}, J. Neurophysiol. \textbf{92}, 2704 (2004).
\bibitem{KobatakeTanaka1994} E. Kobatake and K. Tanaka, \emph{Neuronal selectivities to complex object features in the ventral visual pathway of the macaque cerebral cortex}, J. Neurophysiol. \textbf{71}, 856 (1994).
\bibitem{Albright1984} T. D. Albright, \emph{Direction and orientation selectivity of neurons in visual area MT of the macaque}, J. Neurophysiol. \textbf{52}, 1106 (1984).
\bibitem{Goodfellow2015} I. J. Goodfellow, J. Shlens, and C. Szegedy, \emph{Explaining and harnessing adversarial examples}, in \emph{Proc. ICLR} (2015), arXiv:1412.6572.
\bibitem{Elsayed2018} G. F. Elsayed et al., \emph{Adversarial examples that fool both computer vision and time-limited humans}, in \emph{Advances in Neural Information Processing Systems 31} (2018), arXiv:1802.08195.
\bibitem{Thompson1982} P. Thompson, \emph{Perceived rate of movement depends on contrast}, Vision Res. \textbf{22}, 377 (1982).
\bibitem{BarlowHill1963} H. B. Barlow and R. M. Hill, \emph{Evidence for a physiological explanation of the waterfall phenomenon and figural after-effects}, Nature \textbf{200}, 1345 (1963).
\bibitem{EaglemanSejnowski2000} D. M. Eagleman and T. J. Sejnowski, \emph{Motion integration and postdiction in visual awareness}, Science \textbf{287}, 2036 (2000).

\bibitem{Feinstein1955mu} B. Feinstein, B. Lindeg{\aa}rd, E. Nyman, and G. Wohlfart, \emph{Morphologic studies of motor units in normal human muscles}, Acta Anat. \textbf{23}, 127 (1955).
\bibitem{MilnerBrown1973rec} H. S. Milner-Brown, R. B. Stein, and R. Yemm, \emph{The orderly recruitment of human motor units during voluntary isometric contractions}, J. Physiol. \textbf{230}, 359 (1973).
\bibitem{Jones2002sdn} K. E. Jones, A. F. de C. Hamilton, and D. M. Wolpert, \emph{Sources of signal-dependent noise during isometric force production}, J. Neurophysiol. \textbf{88}, 1533 (2002).
\bibitem{GomiOsu1998stiff} H. Gomi and R. Osu, \emph{Task-dependent viscoelasticity of human multijoint arm and its spatial characteristics for interaction with environments}, J. Neurosci. \textbf{18}, 8965 (1998).
\bibitem{Jankowska1972Ia} E. Jankowska and W. J. Roberts, \emph{Synaptic actions of single interneurones mediating reciprocal Ia inhibition of motoneurones}, J. Physiol. \textbf{222}, 623 (1972).
\bibitem{Pruszynski2008rapid} J. A. Pruszynski, I. Kurtzer, and S. H. Scott, \emph{Rapid motor responses are appropriately tuned to the metrics of a visuospatial task}, J. Neurophysiol. \textbf{100}, 224 (2008).
\bibitem{SaundersKnill2003vis} J. A. Saunders and D. C. Knill, \emph{Humans use continuous visual feedback from the hand to control fast reaching movements}, Exp. Brain Res. \textbf{152}, 341 (2003).
\bibitem{FlanaganWing1997grip} J. R. Flanagan and A. M. Wing, \emph{The role of internal models in motion planning and control: evidence from grip force adjustments during movements of hand-held loads}, J. Neurosci. \textbf{17}, 1519 (1997).
\bibitem{Matsuoka1985osc} K. Matsuoka, \emph{Sustained oscillations generated by mutually inhibiting neurons with adaptation}, Biol. Cybern. \textbf{52}, 367 (1985).

\bibitem{Treede1995heat} R.-D. Treede, R. A. Meyer, S. N. Raja, and J. N. Campbell, \emph{Evidence for two different heat transduction mechanisms in nociceptive primary afferents innervating monkey skin}, J. Physiol. \textbf{483}, 747 (1995).
\bibitem{Weidner1999cfib} C. Weidner et al., \emph{Functional attributes discriminating mechano-insensitive and mechano-responsive C nociceptors in human skin}, J. Neurosci. \textbf{19}, 10184 (1999).
\bibitem{CampbellLaMotte1983first} J. N. Campbell and R. H. LaMotte, \emph{Latency to detection of first pain}, Brain Res. \textbf{266}, 203 (1983).
\bibitem{LaMotte1982hyper} R. H. LaMotte, J. G. Thalhammer, H. E. Torebj{\"o}rk, and C. J. Robinson, \emph{Peripheral neural mechanisms of cutaneous hyperalgesia following mild injury by heat}, J. Neurosci. \textbf{2}, 765 (1982).
\bibitem{Mancini2014touch} F. Mancini, T. Nash, G. D. Iannetti, and P. Haggard, \emph{Pain relief by touch: a quantitative approach}, Pain \textbf{155}, 635 (2014).
\bibitem{Fields1983rvm} H. L. Fields, J. Bry, I. Hentall, and G. Zorman, \emph{The activity of neurons in the rostral medulla of the rat during withdrawal from noxious heat}, J. Neurosci. \textbf{3}, 2545 (1983).
\bibitem{Crombez1996disrupt} G. Crombez, C. Eccleston, F. Baeyens, and P. Eelen, \emph{The disruptive nature of pain: an experimental investigation}, Behav. Res. Ther. \textbf{34}, 911 (1996).
\bibitem{Schouenborg1995wr} J. Schouenborg, H.-R. Weng, J. Kalliom{\"a}ki, and H. Holmberg, \emph{A survey of spinal dorsal horn neurones encoding the spatial organization of withdrawal reflexes in the rat}, Exp. Brain Res. \textbf{106}, 19 (1995).

\bibitem{Andersen1992cereb} B. B. Andersen, L. Korbo, and B. Pakkenberg, \emph{A quantitative study of the human cerebellum with unbiased stereological techniques}, J. Comp. Neurol. \textbf{326}, 549 (1992).
\bibitem{ShadmehrMussaIvaldi1994} R. Shadmehr and F. A. Mussa-Ivaldi, \emph{Adaptive representation of dynamics during learning of a motor task}, J. Neurosci. \textbf{14}, 3208 (1994).

\bibitem{JoseCollison1970ihr} A. D. Jose and D. Collison, \emph{The normal range and determinants of the intrinsic heart rate in man}, Cardiovasc. Res. \textbf{4}, 160 (1970).
\bibitem{Logothetis1987gbg} D. E. Logothetis, Y. Kurachi, J. Galper, E. J. Neer, and D. E. Clapham, \emph{The $\beta\gamma$ subunits of GTP-binding proteins activate the muscarinic K$^+$ channel in heart}, Nature \textbf{325}, 321 (1987).
\bibitem{KellerWood1984fb} M. E. Keller-Wood and M. F. Dallman, \emph{Corticosteroid inhibition of ACTH secretion}, Endocr. Rev. \textbf{5}, 1 (1984).
\bibitem{Walker2012glc} J. J. Walker et al., \emph{The origin of glucocorticoid hormone oscillations}, PLoS Biol. \textbf{10}, e1001341 (2012).
\bibitem{Wildt1981gnrh} L. Wildt et al., \emph{Frequency and amplitude of gonadotropin-releasing hormone stimulation and gonadotropin secretion in the rhesus monkey}, Endocrinology \textbf{109}, 376 (1981).
\bibitem{Welsh2010scn} D. K. Welsh, J. S. Takahashi, and S. A. Kay, \emph{Suprachiasmatic nucleus: cell autonomy and network properties}, Annu. Rev. Physiol. \textbf{72}, 551 (2010).
\bibitem{Khalsa2003prc} S. B. S. Khalsa, M. E. Jewett, C. Cajochen, and C. A. Czeisler, \emph{A phase response curve to single bright light pulses in human subjects}, J. Physiol. \textbf{549}, 945 (2003).
\bibitem{Sack1992blind} R. L. Sack, A. J. Lewy, M. L. Blood, L. D. Keith, and H. Nakagawa, \emph{Circadian rhythm abnormalities in totally blind people: incidence and clinical significance}, J. Clin. Endocrinol. Metab. \textbf{75}, 127 (1992).

\bibitem{Henze2000extra} D. A. Henze et al., \emph{Intracellular features predicted by extracellular recordings in the hippocampus in vivo}, J. Neurophysiol. \textbf{84}, 390 (2000).
\bibitem{Histed2009stim} M. H. Histed, V. Bonin, and R. C. Reid, \emph{Direct activation of sparse, distributed populations of cortical neurons by electrical microstimulation}, Neuron \textbf{63}, 508 (2009).
\bibitem{Orsborn2014coad} A. L. Orsborn et al., \emph{Closed-loop decoder adaptation shapes neural plasticity for skillful neuroprosthetic control}, Neuron \textbf{82}, 1380 (2014).

\bibitem{HuVervaekeStorm2002} H. Hu, K. Vervaeke, and J. F. Storm, \emph{Two forms of electrical resonance at theta frequencies, generated by M-current, h-current and persistent Na$^+$ current in rat hippocampal pyramidal cells}, J. Physiol. \textbf{545}, 783 (2002).
\bibitem{Mauro1970} A. Mauro, F. Conti, F. Dodge, and R. Schor, \emph{Subthreshold behavior and phenomenological impedance of the squid giant axon}, J. Gen. Physiol. \textbf{55}, 497 (1970).
\bibitem{Aksay2001} E. Aksay et al., \emph{In vivo intracellular recording and perturbation of persistent activity in a neural integrator}, Nat. Neurosci. \textbf{4}, 184 (2001).
\bibitem{Major2004} G. Major et al., \emph{Plasticity and tuning of the time course of analog persistent firing in a neural integrator}, Proc. Natl. Acad. Sci. USA \textbf{101}, 7745 (2004).
\bibitem{Tateno2004} T. Tateno, A. Harsch, and H. P. C. Robinson, \emph{Threshold firing frequency-current relationships of neurons in rat somatosensory cortex: type 1 and type 2 dynamics}, J. Neurophysiol. \textbf{92}, 2283 (2004).
\bibitem{Marder2012} E. Marder, \emph{Neuromodulation of neuronal circuits: back to the future}, Neuron \textbf{76}, 1 (2012).
\bibitem{McGintyHarper1976} D. J. McGinty and R. M. Harper, \emph{Dorsal raphe neurons: depression of firing during sleep in cats}, Brain Res. \textbf{101}, 569 (1976).

\bibitem{Gentet2000} L. J. Gentet, G. J. Stuart, and J. D. Clements, \emph{Direct measurement of specific membrane capacitance in neurons}, Biophys. J. \textbf{79}, 314 (2000).
\bibitem{Borst1995} J. G. G. Borst, F. Helmchen, and B. Sakmann, \emph{Pre- and postsynaptic whole-cell recordings in the medial nucleus of the trapezoid body of the rat}, J. Physiol. \textbf{489}, 825 (1995).
\bibitem{AmirDevor2003} R. Amir and M. Devor, \emph{Electrical excitability of the soma of sensory neurons is required for spike invasion of the soma, but not for through-conduction}, Biophys. J. \textbf{84}, 2181 (2003).
\bibitem{Chadderton2004} P. Chadderton, T. W. Margrie, and M. H\"ausser, \emph{Integration of quanta in cerebellar granule cells during sensory processing}, Nature \textbf{428}, 856 (2004).
\bibitem{LitwinKumar2017} A. Litwin-Kumar et al., \emph{Optimal degrees of synaptic connectivity}, Neuron \textbf{93}, 1153 (2017).
\bibitem{NapperHarvey1988} R. M. A. Napper and R. J. Harvey, \emph{Number of parallel fiber synapses on an individual Purkinje cell in the cerebellum of the rat}, J. Comp. Neurol. \textbf{274}, 168 (1988).
\bibitem{Megias2001} M. Meg\'{\i}as, Z. Emri, T. F. Freund, and A. I. Guly\'as, \emph{Total number and distribution of inhibitory and excitatory synapses on hippocampal CA1 pyramidal cells}, Neuroscience \textbf{102}, 527 (2001).
\bibitem{Polsky2004} A. Polsky, B. W. Mel, and J. Schiller, \emph{Computational subunits in thin dendrites of pyramidal cells}, Nat. Neurosci. \textbf{7}, 621 (2004).
\bibitem{Brunel2004} N. Brunel et al., \emph{Optimal information storage and the distribution of synaptic weights: perceptron versus Purkinje cell}, Neuron \textbf{43}, 745 (2004).

\bibitem{Vyazovskiy2009} V. V. Vyazovskiy et al., \emph{Cortical firing and sleep homeostasis}, Neuron \textbf{63}, 865 (2009).
\bibitem{Daan1984} S. Daan, D. G. Beersma, and A. A. Borb\'ely, \emph{Timing of human sleep: recovery process gated by a circadian pacemaker}, Am. J. Physiol. \textbf{246}, R161 (1984).
\bibitem{Borbely1981} A. A. Borb\'ely et al., \emph{Sleep deprivation: effect on sleep stages and EEG power density in man}, Electroencephalogr. Clin. Neurophysiol. \textbf{51}, 483 (1981).
\bibitem{SteriadeAmzica1993} M. Steriade, F. Amzica, and A. Nu\~nez, \emph{Cholinergic and noradrenergic modulation of the slow ($\approx0.3$ Hz) oscillation in neocortical cells}, J. Neurophysiol. \textbf{70}, 1385 (1993).
\bibitem{vonKrosigk1993} M. von Krosigk, T. Bal, and D. A. McCormick, \emph{Cellular mechanisms of a synchronized oscillation in the thalamus}, Science \textbf{261}, 361 (1993).
\bibitem{LeeWilson2002} A. K. Lee and M. A. Wilson, \emph{Memory of sequential experience in the hippocampus during slow wave sleep}, Neuron \textbf{36}, 1183 (2002).
\bibitem{Euston2007} D. R. Euston, M. Tatsuno, and B. L. McNaughton, \emph{Fast-forward playback of recent memory sequences in prefrontal cortex during sleep}, Science \textbf{318}, 1147 (2007).
\bibitem{Hobson2009} J. A. Hobson, \emph{REM sleep and dreaming: towards a theory of protoconsciousness}, Nat. Rev. Neurosci. \textbf{10}, 803 (2009).

\bibitem{CohenSegal2011} D. Cohen and M. Segal, \emph{Network bursts in hippocampal microcultures are terminated by exhaustion of vesicle pools}, J. Neurophysiol. \textbf{106}, 2314 (2011).
\bibitem{Tsodyks2000} M. Tsodyks, A. Uziel, and H. Markram, \emph{Synchrony generation in recurrent networks with frequency-dependent synapses}, J. Neurosci. \textbf{20}, RC50 (2000).
\bibitem{Benson1994} D. L. Benson, F. H. Watkins, O. Steward, and G. Banker, \emph{Characterization of GABAergic neurons in hippocampal cell cultures}, J. Neurocytol. \textbf{23}, 279 (1994).

\input{supplement/bib_22}
\end{thebibliography}


\end{document}
